% 本文件由 LaTeX 排版 Agent 根据有效 translation.json 自动生成。
% author=Irenaeus of Lyons; work_id=0103; title=驳斥异端

\chapter{驳斥异端}

% structure: p0001 source=h1 target=omitted reason=page-title-already-rendered-by-parent

% structure: p0002 source=h2 target=section reason=relative-heading-rank; base=h2

\section{第一卷}

\Emph{序言}\newline{}\Emph{第一章　瓦伦廷门徒关于其幻想之伊昂的起源、名称、秩序及婚配产出的荒谬观点，以及他们为己见而改编的圣经经文选段}\newline{}\Emph{第二章　元父唯独为独生子所知。索菲亚的野心、骚动及危险；她的无形产物；她由界限恢复。为完成伊昂而产出基督与圣神。耶稣的产出方式。}\newline{}\Emph{第三章　这些异端者用以支持自己意见的圣经经文}\newline{}\Emph{第四章　异端者对阿卡莫特形成的描述；可见世界源于她的骚动}\newline{}\Emph{第五章　德谬哥的形成；对其的描述。他是普雷若马之外万物的创造者}\newline{}\Emph{第六章　这些异端者虚构的三类人：善行于他们无必要，但对他人必要；他们堕落的道德}\newline{}\Emph{第七章　母亲阿卡莫特与其所有种子受成全后，将进入普雷若马，并由那些属灵的人陪伴；德谬哥与属魂的人将进入中间居所；但一切属物质的人将归于朽坏。他们对基督由童贞玛利亚真正降生成人的亵渎观点。他们对预言的看法。德谬哥的愚昧无知。}\newline{}\Emph{第八章　瓦伦廷派如何歪曲圣经以支持自己虔诚的见解}\newline{}\Emph{第九章　驳斥这些异端者的邪恶解释}\newline{}\Emph{第十章　普世教会信仰的统一}\newline{}\Emph{第十一章　瓦伦廷及其门徒与其他人的见解}\newline{}\Emph{第十二章　托勒密与克罗巴苏斯追随者的教义}\newline{}\Emph{第十三章　马可的欺诈手段与邪恶行径}\newline{}\Emph{第十四章　马可与其他人种种假说。关于字母与音节的学说}\newline{}\Emph{第十五章　塞吉向马可启示二十四元素与耶稣的诞生。揭露这些荒谬。}\newline{}\Emph{第十六章　马可派荒谬的解释}\newline{}\Emph{第十七章　马可派理论，认为受造之物是按不可见之物的肖像而造}\newline{}\Emph{第十八章　梅瑟书中的段落，异端者歪曲以支持其假说}\newline{}\Emph{第十九章　他们用以证明至圣父在基督来前未为人知的圣经段落}\newline{}\Emph{第二十章　马可派的外典与伪经，以及他们歪曲的福音段落}\newline{}\Emph{第二十一章　这些异端者关于救赎的看法}\newline{}\Emph{第二十二章　异端者偏离真理}\newline{}\Emph{第二十三章　西蒙·马古与梅南德的教义与实践}\newline{}\Emph{第二十四章　撒土尼鲁与巴西里德的教义}\newline{}\Emph{第二十五章　卡波克拉特的教义}\newline{}\Emph{第二十六章　塞林特、厄俾翁派与尼苛劳派的教义}\newline{}\Emph{第二十七章　切尔多与马西昂的教义}\newline{}\Emph{第二十八章　塔提安、禁欲派及其他人的教义}\newline{}\Emph{第二十九章　其他各种诺斯替派别的教义，尤其是巴比洛斯派或波波利安派}\newline{}\Emph{第三十章　奥菲特派与塞特派的教义}\newline{}\Emph{第三十一章　该隐派的教义}\par

% structure: p0004 source=h2 target=section reason=relative-heading-rank; base=h2

\section{第二卷}

序言\newline{}第一章 唯一天主：不可能有其他情形\newline{}第二章 世界不是由天使或任何其他存在违背至高天主的旨意而形成的，而是由圣父藉圣言所创造\newline{}第三章 瓦伦廷派的深渊与圆满，以及马西昂的神，均被证明是荒谬的；世界实际上是由那位构想出它的同一存有所创造，并非缺陷或无知的产物\newline{}第四章 异端者所谓的虚空与缺陷的荒谬性得到证明\newline{}第五章 这个世界并非由圣父所包含的领域内的任何其他存在所形成\newline{}第六章 天使和世界的创造者不可能不知道至高天主\newline{}第七章 受造物不是那些在圆满之内的永世的肖像\newline{}第八章 受造物并非圆满的阴影\newline{}第九章 世界的创造者只有一位，即天主圣父：这是教会恒常的信念\newline{}第十章 异端者对圣经的曲解：天主从无中创造了万有，而非从先在的物质\newline{}第十一章 异端者因不信真理而陷入谬误的深渊：研究其体系的理由\newline{}第十二章 异端者的三十之数既因缺陷也因过剩而犯错：智慧绝不可能离开她的配偶产生任何东西；圣言与沉默不可能同时存在\newline{}第十三章 异端者所持的初始生产秩序是完全不可辩护的\newline{}第十四章 瓦伦廷及其追随者从其体系的原则来源于异教徒；仅名称有所改变\newline{}第十五章 对这些生产物无法给出任何说明\newline{}第十六章 世界的创造者要么自发产生了将要制造之物的影像，要么圆满是按照某个先前体系的影像形成的，如此类推以至无穷\newline{}第十七章 对永世生产的探讨：无论其假定的性质如何，它在各方面都不一致；并且根据异端者的假设，就连心灵和圣父本身也将被无知所玷污\newline{}第十八章 智慧从未真正处于无知或激情中；她的意念不可能与她分离，或表现出其特有的倾向\newline{}第十九章 异端者关于自身起源的荒谬看法：他们关于造物主的意见被证明同样站不住脚和可笑\newline{}第二十章 从比喻、犹达斯的背叛和救主的受难中引出证明第十二永世受苦的论据之无效\newline{}第二十一章 十二宗徒不是永世的预像\newline{}第二十二章 三十永世并非由基督在三十岁时受洗这一事实所预表：他并非在受洗后第十二个月受苦，而是去世时已逾五十岁\newline{}第二十三章 患血漏的妇人并非受苦永世的预像\newline{}第二十四章 异端者从数字、字母和音节中引出的论据之愚蠢\newline{}第二十五章 不应该借助字母、音节和数字来寻求天主；这类研究中谦虚的必要性\newline{}第二十六章 \Quote{知识使人自高自大，惟有爱德能建立人}\newline{}第二十七章 解释比喻和圣经晦涩段落的正确方式\newline{}第二十八章 在今世无法达到圆满的知识：许多问题必须谦卑地交托在天主手中\newline{}第二十九章 驳斥异端者关于灵魂和身体未来命运的观点\newline{}第三十章 他们自称属灵，而造物主却被宣布为属魂的荒谬性\newline{}第三十一章 对前述论点的总结与应用\newline{}第三十二章 进一步揭露异端者为邪恶亵渎的教义\newline{}第三十三章 灵魂轮回说之荒谬\newline{}第三十四章 灵魂在分离状态中可被辨认，并且虽有一开始却是不朽的\newline{}第三十五章 驳斥巴西里德，以及认为先知受不同神祇默感而发预言的看法\par

% structure: p0006 source=h2 target=section reason=relative-heading-rank; base=h2

\section{第三卷}

\Emph{前言}\newline{}\Emph{第一章} 宗徒们直到被赐予圣神的恩赐和能力之后，才开始宣讲福音或记录任何事物。他们只宣讲独一的天主，天地的创造者。\newline{}\Emph{第二章} 异端者既不遵从圣经，也不遵从圣传。\newline{}\Emph{第三章} 从各教会中始终保持宗徒继承的事实，驳斥异端者。\newline{}\Emph{第四章} 真理只能在天主教会中找到，她是宗徒教义的唯一宝库。异端是近期形成的，无法追溯其起源至宗徒。\newline{}\Emph{第五章} 基督和他的宗徒们毫无欺诈、欺骗或虚伪，宣讲独一的天主圣父是万物的创造者。他们没有使他们的教义迎合听众的偏见。\newline{}\Emph{第六章} 圣神在整部旧约圣经中，除了那位真实的天主之外，没有提及任何其他神或主。\newline{}\Emph{第七章} 回应基于圣保禄的话\BibleRef{格后 4:5}的反对意见。\Person{圣保禄}偶尔使用不合语法顺序的词语。\newline{}\Emph{第八章} 回应基于基督的话\BibleRef{玛 6:24}的反对意见。只有天主应该真正被称为天主和主，因为他是无始无终的。\newline{}\Emph{第九章} 同一位天主，天地的创造者，就是先知们所预言的、福音所宣告的那位。一开始就从\BibleBook{玛窦}{福音}证明这一点。\newline{}\Emph{第十章} 从\BibleBook{马尔谷}{福音}和\BibleBook{路加}{福音}引出上述的证明。\newline{}\Emph{第十一章} 继续从\BibleBook{若望}{福音}引出证明。福音共有四部，不多不少。其神秘原因。\newline{}\Emph{第十二章} 其余宗徒的教义。\newline{}\Emph{第十三章} 驳斥认为\Person{保禄}是唯一认识真理的宗徒的意见。\newline{}\Emph{第十四章} 如果\Person{保禄}知道任何其他宗徒未被启示的奥秘，那么他忠实的同伴和旅伴\Person{路加}不可能不知道它们；真理也不可能隐藏在他那里，因为通过他我们才得知福音历史的许多最重要细节。\newline{}\Emph{第十五章} 从\Person{圣路加}的作品（必须整体接受）驳斥贬低\Person{圣保禄}权威的\OriginalTerm{厄俾翁派}{Ebionites}。揭露\OriginalTerm{诺斯替派}{Gnostics}的虚伪、欺骗和骄傲。宗徒及其门徒认识并宣讲独一的天主，世界的创造者。\newline{}\Emph{第十六章} 从宗徒著作证明耶稣基督是同一位，天主的独生子，完美的天主和完美的人。\newline{}\Emph{第十七章} 宗徒们教导说，降临在耶稣身上的既不是基督，也不是\OriginalTerm{救主}{Saviour}，而是圣神。这降临的原因。\newline{}\Emph{第十八章} 延续上述论证。从\Person{圣保禄}的著作和主的话证明，基督和耶稣不能被视为不同的存在；也不能声称天主子仅仅是外表上成为人，而是真实地成了人。\newline{}\Emph{第十九章} 耶稣基督不是按自然常理由若瑟所生的凡人，而是由至高圣父所生的真天主，由童贞女所生的真人。\newline{}\Emph{第二十章} 天主通过人的堕落显示出他的忍耐、慈爱、怜悯和拯救的大能。因此，人如果忘记自己的处境和赐予他的恩惠，不承认天主的恩宠，就是最忘恩负义的。\newline{}\Emph{第二十一章} 为依撒意亚的预言\BibleRef{依 7:14}辩护，驳斥\OriginalTerm{提奥多提安}{Theodotion}、\OriginalTerm{阿奎拉}{Aquila}、\OriginalTerm{厄俾翁派}{Ebionites}和犹太人的错误解释。\OriginalTerm{七十贤士译本}{Septuagint}的权威。证明基督由童贞女所生的论据。\newline{}\Emph{第二十二章} 基督取真实血肉，由童贞女怀孕并出生。\newline{}\Emph{第二十三章} 反对\OriginalTerm{塔提安}{Tatian}的论据，表明符合天主的公义和仁慈，使第一个亚当首先分享基督提供给所有人的救恩。\newline{}\Emph{第二十四章} 总结针对\OriginalTerm{诺斯替派}{Gnostic}不虔诚的各种论据。异端者被各种教义之风吹动，而教会始终如一的教导则反对他们，这教导永不改变且自身一致。\newline{}\Emph{第二十五章} 这个世界由一位天主的眷顾所统治，他既有无限的正义惩罚恶人，又有无限的良善祝福虔诚者，并赐予他们救恩。\par

% structure: p0008 source=h2 target=section reason=relative-heading-rank; base=h2

\section{\Emph{第四卷}}

\Emph{序言}\newline{}\Emph{第一章}主承认只有一位天主和圣父。\newline{}\Emph{第二章}从梅瑟和其他先知明确的见证（他们的话就是基督的话）证明只有一位天主，世界的创造者，也就是我们的主所宣讲并称为圣父的那一位。\newline{}\Emph{第三章}答诺斯替派的刁难。我们不应以为真实的天主会改变或终结，因为诸天是他的宝座，大地是他的脚凳，将要消逝。\newline{}\Emph{第四章}答另一异议，显示耶路撒冷——大君的城邑——的毁灭并未减损天主至高的尊威和能力，因为这一毁灭是由同一天主最明智的安排执行的。\newline{}\Emph{第五章}作者回到先前的论证，指明律法和先知所宣布的只有一位天主，基督承认他为圣父，他借着他的圣言——与他同为永生的一位天主——在双重盟约中将自己启示给人。\newline{}\Emph{第六章}解释基督的话：\Quote{除了圣子，没有人认识圣父}等等；异端者误解了这些话。证明：藉着圣父启示圣子，也藉着圣子被启示，圣父从未是未知的。\newline{}\Emph{第七章}重述先前的论证，显示亚巴郎藉着圣言的启示认识圣父和天主圣子的来临。为此，他欢喜看见基督的日子，那时对他所作的许诺将要应验。这欢喜的果实流到了后世，即那些有分于亚巴郎信德的人，而不是拒绝天主圣言的犹太人。\newline{}\Emph{第八章}马西昂及其追随者的徒劳尝试，他们排除亚巴郎于基督所赐的救恩之外；基督非但没有毁坏律法，反而在安息日治病时成全了律法，从而解放了亚巴郎和亚巴郎的后裔。\newline{}\Emph{第九章}双重盟约只有一位作者和同一个终向。\newline{}\Emph{第十章}旧约圣经，特别是梅瑟所写的，处处提及天主圣子，并预言他的降临和受难。由此事实可知它们是由同一的天主所默感。\newline{}\Emph{第十一章}古先知和义人事先知道基督的降临，切望看见并听见他；他藉着圣神在圣经中启示自己，自身毫无改变，日日以恩惠丰富人，但赐予后世比前代更丰盛。\newline{}\Emph{第十二章}从基督谴责与旧律相悖的传统和习俗，同时确认其最重要的规诫，并教导他自己就是梅瑟律法的终向，清楚地显出新旧律法只有一位作者。\newline{}\Emph{第十三章}基督并没有废除律法的自然规诫，反而成全并扩展了它们。他移除了旧律的轭和束缚，使人类现今得释放，能以适合子女的孝爱虔诚事奉天主。\newline{}\Emph{第十四章}如果天主要求人服从，如果他造了人，召叫了他，把他置于法律之下，那完全是为了人的益处；并非天主需要人，而是他仁慈地以各种方式将恩惠赐予人。\newline{}\Emph{第十五章}起初，天主认为将自然律或十诫刻在人心上就已足够；但后来，他发觉有必要用梅瑟律法的轭来约束那些滥用自由的犹太人的欲望，甚至因为他们的心硬而增加一些特别的诫命。\newline{}\Emph{第十六章}完美的义既不是由割损，也不是由任何其他法律仪式所赐予的。然而，十诫并未被基督废除，而是永远有效：人从未被免除其诫命。\newline{}\Emph{第十七章}证明天主设立肋未制度并非为了自己，也不是因为需要这种服务；因为实际上他不需要人任何东西。\newline{}\Emph{第十八章}论祭献和供物，以及真正献上它们的人。\newline{}\Emph{第十九章}地上的事物可以是天上的预象，但天上的事物不能是其他更崇高未知之物的预象；除非彻底疯狂，我们也不能主张天主仅仅被我们认知为某个仍然未知的更高存在的预象。\newline{}\Emph{第二十章}一位天主藉着圣言和圣神创造了世间万物：并且虽然在此生他为我们是无形不可理解的，但他并非未知；因为他的工程宣告他，他的圣言也表明他可以在多种方式下被看见和认识。\newline{}\Emph{第二十一章}亚巴郎的信德与我们的相同；这信德由古圣祖的言行所预象。\newline{}\Emph{第二十二章}基督并非只为某一时代的人而来，而是为所有正义虔诚生活、相信他的人，也为那些将要相信的人。\newline{}\Emph{第二十三章}古圣祖和先知通过指出基督的降临，从而加强了后裔通往基督信德的道路；因此宗徒的辛劳得以减轻，因为他们收获了他人辛劳的果实。\newline{}\Emph{第二十四章}外邦人的归化比犹太人的更困难；因此从事前一项任务的宗徒的辛劳大于承担后一项的。\newline{}\Emph{第二十五章}双重盟约都在亚巴郎和塔玛尔的劳作中预象；然而，每个盟约只有同一位天主。\newline{}\Emph{第二十六章}隐藏在圣经中的宝藏就是基督；圣经的真确解释只有在教会内才能找到。\newline{}\Emph{第二十七章}古时之人的罪过，曾招致天主的责怒，被天主眷顾写下来，使我们能由此获得教训，不致充满骄傲。因此，我们不可推论有另一位不同于基督所宣讲的天主；我们更应该敬畏，免得同一位曾惩罚古人的天主，加给我们更重的刑罚。\newline{}\Emph{第二十八章}那些过度渲染基督的仁慈，却对审判缄默不语，只注视新约更丰盛的恩宠，却忘记新约要求我们更高的成全，从而试图证明有一位超越创造世界的天主的另一位神的人，显明自己是无知的。\newline{}\Emph{第二十九章}驳斥马西昂派的论据，他们企图显示天主是罪的作者，因为他使法老和他的臣仆眼瞎。\newline{}\Emph{第三十章}驳斥马西昂派提出的另一论据：天主吩咐希伯来人夺取埃及人的财物。\newline{}\Emph{第三十一章}我们不应草率地将圣经未谴责的行为归咎为古人的罪行，而应从中寻求预象：以罗特所犯的乱伦为例。\newline{}\Emph{第三十二章}一位天主是双重盟约的作者，这由一位曾受教于宗徒的长老之权威所确认。\newline{}\Emph{第三十三章}凡承认一位天主是双重盟约的作者，并与教会中的长老一同勤读圣经的人，就是真正的属灵门徒；他将正确地理解和解释先知关于基督和新约自由的一切宣告。\newline{}\Emph{第三十四章}驳斥马西昂派的证明：先知的所有预言都指向我们的基督。\newline{}\Emph{第三十五章}驳斥那些声称先知一部分预言是由至高者默感说出，另一部分来自\OriginalTerm{德谬哥}{Demiurge}的人。华伦提努派内部对于这些预言亦互有分歧。\newline{}\Emph{第三十六章}先知是被同一位父派遣的，圣子也是从他那里受派遣。\newline{}\Emph{第三十七章}人拥有自由意志，并被赋予选择的能力。因此，说一些人生来为善，另一些人生来为恶，是不正确的。\newline{}\Emph{第三十八章}为何人起初没有被造成完美的。\newline{}\Emph{第三十九章}人被赋予辨别善恶的能力；因此，在没有强迫的情况下，他有能力凭自己的意志和选择去履行天主的诫命，藉此避免为叛逆者预备的灾祸。\newline{}\Emph{第四十章}同一天主圣父惩罚被弃者，并赏赐蒙选者。\newline{}\Emph{第四十一章}那些不信天主、悖逆的人，是魔鬼的天使和子女，不是出于本性，而是出于模仿。本书结束，及续篇的范围。\par

% structure: p0010 source=h2 target=section reason=relative-heading-rank; base=h2

\section{\WorkTitle{卷五}}

\Emph{Preface}\newline{}\Emph{Chapter 1}基督独自能教导神圣之事并救赎我们：祂同一者由童贞玛利亚取了肉身，非仅外表，而是实际，藉圣神的运作，为了更新我们。对瓦伦蒂诺及厄彼翁之妄想的评论。\newline{}\Emph{Chapter 2}当基督以祂的恩宠探访我们时，祂并没有来到不属于祂的地方；而且，藉着为我们倾流祂真实的血，并在圣体圣事中向我们显示祂真实的肉，祂赋予我们的肉获得救恩的能力。\newline{}\Emph{Chapter 3}天主的德能与荣耀在人类肉体的软弱中闪耀，因为祂将使我们的身体参与复活与不朽，尽管祂是从地上的尘土形成它；祂也将赐予它享受不朽，正如祂赐予它与灵魂共享这短暂的生命。\newline{}\Emph{Chapter 4}那些臆想另一位圣父（不同于世界的创造主）的人是被欺骗的；因为如果祂不能或不愿意将永生延伸至我们的身体，祂必是软弱无用，或是恶毒且充满嫉妒。\newline{}\Emph{Chapter 5}古人的长寿，厄里亚和哈诺客以自身身体被提升天，以及约纳、沙得辣客、默沙客、阿贝得乃哥在极端危险中的保全，都是天主能复活我们身体得永生的清楚证明。\newline{}\Emph{Chapter 6}天主将把救恩赐予整个人性，即身体与灵魂紧密联合，因为圣言取了这人性，并饰以圣神的恩赐，我们的身体是（并被称为）圣神的宫殿。\newline{}\Emph{Chapter 7}既然基督在我们的肉中复活，我们也将同样复活；因为应许给我们的复活不应指自然不朽的灵魂，而应指本身可朽的身体。\newline{}\Emph{Chapter 8}我们所领受的圣神恩赐预备我们进入不朽，使我们成为属神的，并将我们与属血肉的人分开。这两类在律法时期的洁净与不洁净动物中得到了预表。\newline{}\Emph{Chapter 9}说明宗徒那节被异端者曲解的经文（\BibleQuote{血肉不能承受天主的国}）应如何理解\newline{}\Emph{Chapter 10}通过从野橄榄树（其品质而非本性因嫁接而改变）所作的比喻，他证明了更重要的事；他也指出，人若没有圣神，就不能结出果实，也不能继承天主的国。\newline{}\Emph{Chapter 11}论属血肉之人和属神之人的行为；也论属神的洁净并不涉及我们身体的实质，而是涉及我们先前生活的方式。\newline{}\Emph{Chapter 12}生命与死亡的区别；生命的气息与赋予生命的圣神：身体的实质如何使曾经死去的复活。\newline{}\Emph{Chapter 13}在基督所复活的人身上，我们拥有复活的最大证明；我们的心能接受天主的圣神，因此证明它们能承受永生。\newline{}\Emph{Chapter 14}除非肉体要得救，圣言就不会取了与我们相同实质的肉体；由此推论，我们也不会由祂而得和好。\newline{}\Emph{Chapter 15}从依撒意亚和厄则克耳引用的复活证据：创造我们的同一天主也将复活我们。\newline{}\Emph{Chapter 16}既然我们的身体归于尘土，它们便从尘土得到实质；并且，由于圣言的降临，我们内天主的肖像以更清晰的光显现。\newline{}\Emph{Chapter 17}只有一位主和一位天主，即万有的圣父和创造者，祂在基督内爱了我们，赐给我们诫命，并赦免了我们的罪；基督在赦免我们罪时证明了自己是祂的圣子和圣言。\newline{}\Emph{Chapter 18}天主圣父和祂的圣言藉其本身的能力和智慧形成了所有受造物（他们使用这些受造物），而非出于缺陷或无知。天主子从圣父领受了全部权能，否则祂绝不会取肉体。\newline{}\Emph{Chapter 19}在悖逆犯罪的厄娃与她的主保童贞玛利亚之间进行了一个比较。提及各种分歧的异端。\newline{}\Emph{Chapter 20}应当听从宗徒们托付教会、并持守同一得救教义的牧者；另一方面，应当避开异端者。对于信仰的奥秘，我们必须持守清醒的思考。\newline{}\Emph{Chapter 21}基督是已提及万物的头。祂理应由万有的创造者圣父派遣，取了人性，受撒殚试探，以完成应许，并获得光荣而完美的胜利。\newline{}\Emph{Chapter 22}法律宣明真实的主和唯一的天主，并\newline{}\Emph{Chapter 23}魔鬼惯于说谎，亚当因此被引入歧途，在创造的第六日犯了罪，在此日他也被基督更新。\newline{}\Emph{Chapter 24}论魔鬼的不断谎言，以及世界上的权力和治理，我们应当服从它们，因为它们是由天主（而非魔鬼）设立的。\newline{}\Emph{Chapter 25}论假基督的欺诈、骄傲和暴虐的统治，正如\newline{}\Emph{Chapter 26}若望和达尼尔预言了罗马帝国的瓦解和荒凉，这将在世界末日和基督永恒国度之前发生。驳斥诺斯替派（撒殚的工具），他们捏造了不同于创造主的另一位父。\newline{}\Emph{Chapter 27}基督将来的审判。共融与分离\newline{}\Emph{Chapter 28}义人与恶人之间的区别。假基督时期的将来背叛，以及世界的终结。\newline{}\Emph{Chapter 29}万物被创造为人服务。假基督的欺骗、邪恶和背教权势。这在洪水时被预表，后来由沙得辣客、默沙客、阿贝得乃哥所受的迫害所预表。\newline{}\Emph{Chapter 30}虽然对假基督名字的数目有一定把握，但我们不应对该名字本身作出草率结论，因为这个数目能适合许多名字。圣神保留此点的理由。假基督的统治和死亡。\newline{}\Emph{Chapter 31}我们身体的保存由基督的复活和升天所证实：圣人们的灵魂在中间时期处于期望的状态，等待他们接受完美而圆满光荣的时刻。\newline{}\Emph{Chapter 32}圣人们在其中遭受如此多苦难的肉体，将接受他们劳苦的果实；尤其因为一切受造物都在等待这事，且天主向亚巴郎和他的后裔应许了这事。\newline{}\Emph{Chapter 33}同一命题的进一步证明，来自基督所作的应许，当时祂宣称要在祂父的国里与祂的门徒喝葡萄树的果实，同时祂应许赏报他们百倍，并使他们参加宴席。雅各伯的祝福已经指出了这点，正如帕皮亚和长老们所解释的。\newline{}\Emph{Chapter 34}他用依撒意亚、厄则克耳、耶肋米亚和达尼尔的各样见证，以及主人应许要服侍的醒寤仆人的比喻，来加强他对圣人们复活后暂时和地上国度的意见。\newline{}\Emph{Chapter 35}他坚持说这些已引用的见证不能寓意式地理解为天上的祝福，而将在假基督来临之后、复活之后，在尘世的耶路撒冷中实现。在先前的预言之后，他又附加了来自依撒意亚、耶肋米亚和\BibleBook{若望}{默示录}的其他预言。\newline{}\Emph{Chapter 36}人们将实际复活：世界不会被消灭；但圣人们将有各等的住所，按照每个人被指定的等级。万物将服从天主圣父，这样祂将成为万有中的万有。\par

\SourceRecord{Translated by Alexander Roberts and William Rambaut.；Ante-Nicene Fathers；Vol. 1.；Edited by Alexander Roberts, James Donaldson, and A. Cleveland Coxe.；Buffalo, NY: Christian Literature Publishing Co.,；1885.；<http://www.newadvent.org/fathers/0103.htm>.}

% structure: p0001 source=h1 target=section reason=page-title

\section{驳斥异端（第一卷，前言）}

1. 既然有些人置真理于不顾，引进虚谎的言语和无聊的族谱，正如宗徒所说，这些事\Quote{引发争议，而非信德中合乎敬虔的造益}，并藉其巧妙编造的似是而非之辞，引诱缺乏经验者的心思，将其掳去，[我深感有责任，亲爱的朋友，撰写以下论著，以揭露并抵制他们的阴谋。]这些人篡改天主的圣言，证明自己是那美好启示之言的邪恶诠释者。他们还颠覆许多人的信德，藉着假装拥有[更高]知识，将人从那创造并装饰宇宙者那里引开，仿佛他们能揭示出某种比那创造天地及其中万物的天主更卓越、更崇高的事物。藉着动听而貌似有理的言辞，他们狡猾地诱使头脑简单的人探究他们的体系；然而，当他们将这些简单人引入其关于\OriginalTerm{德谬哥}{Demiurge}的亵渎而不虔的见解时，却又笨拙地毁灭了他们；而这些人连在这样的问题上也无法分辨谬误与真理。\par

2. 其实，谬误从不以其赤裸的丑陋面目出现，以免一旦暴露就立刻被识破。但它被狡猾地装扮上一件诱人的外衣，以便藉其外在形式，在缺乏经验的人面前显得（说来可笑）比真理本身更真实。一位远比我高明的人曾就此恰当地说：\Quote{精妙的玻璃仿制品仿佛藐视那珍贵宝石绿宝石（某些人极为珍视），除非它受到能检验并揭露赝品之人的审视。或者，当黄铜与银混合时，缺乏经验的人岂能轻易察觉黄铜的存在？}因此，唯恐因我的疏忽，一些人被狼掠走，如同羊群，而他们却未察觉这些人的真面目——因为他们外表披着羊皮（主曾告诫我们当提防他们，见：\BibleRef{玛 7:15}），且他们的言语与我们相似，而思想却大相径庭——我认为有责任（在阅读了瓦伦廷门徒们所谓的一些\WorkTitle{注释}，并亲自与其中一些人交往而了解其教义之后）向你，我的朋友，揭示这些怪诞深奥的奥秘，它们并非每个人的理智都能领会，因为并非所有人都充分净化了自己的头脑。我这样做，是为了使你认识这些事，进而向所有与你相关的人解释，并劝勉他们避开这种疯狂和亵渎基督的深渊。因此，我打算尽己所能，简明清晰地陈述那些现今传播异端之人的见解。我特别指托勒密的门徒，他们的学派可说是瓦伦廷学派的一个分支。我也将尽力，按我微薄的能力，提供推翻他们的方法，表明他们的陈述是多么荒谬且与真理不符。我并非在作文或雄辩上训练有素，但我的友爱之情促使我向你及你所有的同伴显明那些至今被隐藏，却终因天主的仁慈而得以显明的教义。\BibleQuote{因为掩盖的事没有不显露出来的，隐秘的事没有不被人知道的。}\BibleRef{玛 10:26}\par

3. 你不可期望，我这个居住在\OriginalTerm{凯尔特人}{Keltæ}中、大多使用蛮族语言的人，会展示出任何我从未学过的修辞技巧，或任何我从未练习过的写作造诣，或任何我自诩的风格之美与说服力。但你会以善意接受我同样以诚朴之心、直率而通俗地写给你的话；而你自己（因比我更有能力）会扩展我寄给你的那些仅如种子般的观念；并凭你理解的广大，将我所简略提及的要点充分展开，以便将我所软弱陈述的，有力地陈明在你同伴面前。总之，正如我（为满足你长期渴望了解这些人教义的心愿）不辞劳苦，不仅使你认识这些学说，也提供证明其虚假的方法；同样，你要按主赐给你的恩宠，向他人作热忱有效的服务者，使人们不再被这些异端者的似是而非体系所迷惑。——我现在便开始描述。\par

\SourceRecord{Translated by Alexander Roberts and William Rambaut.；Ante-Nicene Fathers；Vol. 1.；Edited by Alexander Roberts, James Donaldson, and A. Cleveland Coxe.；Buffalo, NY: Christian Literature Publishing Co.,；1885.；<http://www.newadvent.org/fathers/0103100.htm>.}

% structure: p0001 source=h1 target=section reason=page-title

\section{\WorkTitle{驳斥异端（第一卷 第一章）}}

1. 他们于是主张，在不可见、不可言说的至高之处，存在着一个完美、先存的\OriginalTerm{永世}{Æon}，他们称之为\OriginalTerm{普罗阿耳刻}{Proarche}、\OriginalTerm{普罗帕托耳}{Propator}和\OriginalTerm{彼托斯}{Bythus}，并描述为不可见、不可理解的。他是永恒、非受生的，在无数世代的长河中处于深沉的宁静与寂静之中。与他同在的还有\OriginalTerm{恩诺亚}{Ennœa}，他们也称她为\OriginalTerm{卡里斯}{Charis}和\OriginalTerm{西革}{Sige}。最后，这位彼托斯决定从他自身发出万有的开端，并将这一他决意产生的产物植入与他同时的西革中，如同种子植入子宫。她于是接受了这粒种子，怀孕后生出了\OriginalTerm{努斯}{Nous}，他既像他、又与他相等，是唯独能领悟他父的伟大。他们称这位努斯为\OriginalTerm{默诺革内斯}{Monogenes}，也称为\OriginalTerm{父}{Father}和万有之始。与他一同产生的还有\OriginalTerm{阿勒忒亚}{Aletheia}；这四位构成了第一且首生的毕达哥拉斯式\OriginalTerm{四元组}{Tetrad}，他们也将此称为万有之根。因为首先有彼托斯和西革，然后有努斯和阿勒忒亚。默诺革内斯领悟了他被产生的目的，自己也发出了\OriginalTerm{逻各斯}{Logos}和\OriginalTerm{佐厄}{Zoe}，成为他之后所有将要产生者的父，以及整个\OriginalTerm{普累若麻}{Pleroma}的开端和形成。通过逻各斯和佐厄的结合，产生了\OriginalTerm{安特罗波斯}{Anthropos}和\OriginalTerm{艾克利西亚}{Ecclesia}；这样就形成了首生的\OriginalTerm{俄格多阿德}{Ogdoad}，即万有之根和实体，在他们中被四个名字称呼，即彼托斯、努斯、逻各斯和安特罗波斯。因为每一个都是雄雌同体的，如下：普罗帕托耳通过结合与他的恩诺亚相连；然后是默诺革内斯，即努斯，与阿勒忒亚；逻各斯与佐厄；安特罗波斯与艾克利西亚。\par

2. 这些\OriginalTerm{永世}{Æon}为父的光荣而产生，并希望通过自身努力实现这一目标，便通过结合发出流溢。\OriginalTerm{逻各斯}{Logos}和\OriginalTerm{佐厄}{Zoe}在产生\OriginalTerm{安特罗波斯}{Anthropos}和\OriginalTerm{艾克利西亚}{Ecclesia}之后，又发出了其他十个永世，其名称如下：\OriginalTerm{彼提乌斯}{Bythius}和\OriginalTerm{米克西斯}{Mixis}，\OriginalTerm{阿革拉托斯}{Ageratos}和\OriginalTerm{赫诺西斯}{Henosis}，\OriginalTerm{奥托菲厄斯}{Autophyes}和\OriginalTerm{赫多内}{Hedone}，\OriginalTerm{阿基内托斯}{Acinetos}和\OriginalTerm{辛克拉西斯}{Syncrasis}，\OriginalTerm{默诺革内斯}{Monogenes}和\OriginalTerm{马卡里亚}{Macaria}。这就是他们宣称由逻各斯和佐厄所产生的十个永世。然后他们补充说，安特罗波斯本人与艾克利西亚一起产生了十二个永世，他们给这些永世起名如下：\OriginalTerm{帕拉克莱托斯}{Paracletus}和\OriginalTerm{皮斯蒂斯}{Pistis}，\OriginalTerm{帕特里科斯}{Patricos}和\OriginalTerm{厄尔皮斯}{Elpis}，\OriginalTerm{迈特里科斯}{Metricos}和\OriginalTerm{阿加佩}{Agape}，\OriginalTerm{艾诺斯}{Ainos}和\OriginalTerm{西内西斯}{Synesis}，\OriginalTerm{艾克利西亚斯提科斯}{Ecclesiasticus}和\OriginalTerm{马卡里奥特斯}{Macariotes}，\OriginalTerm{忒勒托斯}{Theletos}和\OriginalTerm{索菲亚}{Sophia}。\par

3. 这些就是这些人的错误体系中的三十个\OriginalTerm{永世}{Æon}；他们描述说，这些永世可以说是包裹在沉默中，不为任何人所知〔除这些自称的教师外〕。此外，他们宣称这不可见、属灵的\OriginalTerm{普累若麻}{Pleroma}是三分组成的，分为一个\OriginalTerm{俄格多阿德}{Ogdoad}、一个\OriginalTerm{十元组}{Decad}和一个\OriginalTerm{十二元组}{Duodecad}。为此，他们断言\Quote{救主}——因为他们不愿意称祂为\Quote{主}——在三十年内没有公开行事\BibleRef{路 3:23}，以此表明这些永世的奥秘。他们还主张，这三十个永世在葡萄园工人的比喻\BibleRef{玛 20:1}中得到了最清楚的指示。因为有些人在约第一时辰被派去，有些人在第三时辰，有些人在第六时辰，有些人在第九时辰，有些人在第十一时辰。现在，如果我们把这里提到的时辰数字加起来，总和将是三十：因为一、三、六、九和十一相加等于三十。他们坚持认为，通过这些时辰，永世被指出来；同时他们主张，这些是他们特别职责要揭示的伟大、奇妙、迄今为止无法言说的奥秘；因此，每当他们在圣经所包含的大量事物中找到任何可以采纳并适应其无根据推测的内容时，他们就这样做。\par

\SourceRecord{Translated by Alexander Roberts and William Rambaut.；Ante-Nicene Fathers；Vol. 1.；Edited by Alexander Roberts, James Donaldson, and A. Cleveland Coxe.；Buffalo, NY: Christian Literature Publishing Co.,；1885.；<http://www.newadvent.org/fathers/0103101.htm>.}

% structure: p0001 source=h1 target=section reason=page-title

\section{\WorkTitle{驳斥异端（卷一，第二章）}}

1. 他们接着告诉我们，他们体系中的\OriginalTerm{元父}{Propator}只被从他而出的\OriginalTerm{独生子}{Monogenes}所认识，换言之，只被\Person{Nous}所认识，而对于其他一切\OriginalTerm{永世}{Æon}来说，他是不可见、不可理解的。而且，据他们说，唯有\Person{Nous}乐意瞻仰圣父，并因思量其不可测量的伟大而欢欣；同时，他也思量如何将圣父的伟大传达给其余的永世，向他们启示圣父是何等浩瀚、大能，如何无始无终——超乎理解，完全无法看见。但是，按照圣父的旨意，\OriginalTerm{沉默}{Sige}约束了他，因为圣父的旨意是要引导他们全都认识前述的元父，并在他们内心激起探究其本性的渴望。同样，其余的永世也以某种安静的方式，渴望看见他们存在的根源，并默观那无始的第一因。\par

2. 但有一个永世远远跑在其余之前，他是从\Person{人}和\Person{教会}而生的\OriginalTerm{十二永世}{Duodecad}中最年幼的一个，即\Person{索菲亚}；她离开了她的配偶\Person{忒勒托斯}的拥抱而遭受了情欲。这情欲最初确实发生在那些与\Person{Nous}和\Person{阿勒忒娅}相连的永世中，但却像传染病一样传给了这个堕落的永世；她以爱为伪装行事，实则受轻率的影响，因为她没有像\Person{Nous}那样享受与完美圣父的共融。他们说，这情欲在于渴望探究圣父的本性；因为她渴望理解他的伟大。当她未能达到目的，因为她追求不可能之事，而陷入了极度的精神痛苦时，既因圣父的无限深奥和不可探究的本性，也因她对圣父的爱，她不断向前伸展，就处于被他的甘甜吸收并消解于他的绝对本质的危险之中，除非她遇到那支持万物并将它们保持在不可言说的伟大之外的能力。他们称这能力为\OriginalTerm{界限}{Horos}；他们说，她被这能力所约束和支持；然后，她好不容易才恢复自我，确信圣父是不可理解的，于是放弃了原初的设计，连同那因她惊叹的压倒性影响而从她内心产生的情欲。\par

3. 但 他们中的另一些人以神话方式描述\Person{索菲亚}的情欲与恢复如下：他们说，她从事了一件不可能且不可行的事，就生出了一个无形体的实体，正如她的女性本性所能产生的。当她观看这实体时，她最初的感受是因产生的不完美而悲伤，随后是恐惧，生怕这会导致她自身的存在终结。接着，她仿佛失去了自制，陷入极大的困惑，竭力想发现这一切的原因，以及如何隐藏所发生的事。被这些情欲大为困扰后，她终于改变了想法，试图重新回到圣父那里。然而，当她稍作尝试时，力量不济，她成了圣父的哀求者。其余的永世，尤其是\Person{Nous}，也同她一起呈上祈求。因此，他们宣称物质实体起源于无知、悲伤、恐惧和混乱。\par

4. 随后，圣父通过\Person{独生子}，按自己的形象产生了上述的\OriginalTerm{界限}{Horos}，没有结合，是男女同体的。因为他们主张，有时圣父与\Person{沉默}结合行动，但有时他显示自己独立于男女。他们称这界限为\OriginalTerm{十字架}{Stauros}、\OriginalTerm{救赎者}{Lytrotes}、\OriginalTerm{切分者}{Carpistes}、\OriginalTerm{立界者}{Horothetes}和\OriginalTerm{引导者}{Metagoges}。他们宣称，\Person{索菲亚}通过这界限得到了净化并稳固，同时被恢复到与她相配的结合。因为她的\OriginalTerm{内生意念}{Enthymesis}连同随之而来的情欲被从她身上取走后，她自己当然留在\OriginalTerm{圆满界}{Pleroma}内；但她的内生意念连同其情欲被界限分离、围栏并驱逐出那个圆圈。这内生意念无疑是一种精神实体，拥有永世的一些自然倾向，但同时无定形、无形式，因为它什么也没有领受。因此，他们说它是一个软弱而女性化的产物。\par

5. 在这实体被置于永世的圆满界之外，其母亲恢复到与她相配的结合之后，他们告诉我们，\Person{独生子}按照圣父审慎的先见，产生了另一对配偶，即\Person{基督}和\OriginalTerm{圣神}{Holy Spirit}（以免任何永世陷入类似\Person{索菲亚}的灾祸），为的是坚固和增强圆满界，同时也完成了永世的数目。然后，\Person{基督}教导他们关于他们结合的本性，并教导那些拥有对无始者的理解的人自足。他也向他们宣告了与认识圣父有关的事——即圣父不能被理解、不能领悟，甚至不能被看见或听见，除非仅被\Person{独生子}认识。至于其余的永世拥有永续存在的原因，在于圣父本性中不可理解的部分；但他们起源和形成的原因，则在于圣父可被理解的部分，即在圣子内。因此，刚刚产生的\Person{基督}在他们中间成就了这些事。\par

6. 但是，圣神教导他们，在彼此之间得以完全平等时，献上感恩，并引导他们进入真正安息之境。于是，他们告诉我们，众\OriginalTerm{永世}{Æon}在形态与情感上被造得彼此平等，以致全体都变成了如\OriginalTerm{诺斯}{Nous}、\OriginalTerm{逻各斯}{Logos}、\OriginalTerm{人}{Anthropos}与\OriginalTerm{基督}{Christus}一般。同样，女性永世全体都变成了如\OriginalTerm{真理}{Aletheia}、\OriginalTerm{生命}{Zoe}、\OriginalTerm{灵}{Spiritus}与\OriginalTerm{教会}{Ecclesia}一般。一切既如此安定，并被带入完美安息之境后，他们接着告诉我们，这些存有满怀大喜乐向\OriginalTerm{太父}{Propator}高唱赞美诗，而他本人也同享丰盛的欢跃。然后，为感恩所赐予的大恩，众永世的全体\OriginalTerm{圆满}{Pleroma}，怀着同一意愿与渴望，并得\OriginalTerm{基督}{Christ}与圣神的同意，他们的父也以其认可为他们的行为盖印，将每位存有自身所拥有的最美、最贵重的一切聚集起来，并把这些贡献巧妙地融合成一个整体，为了\OriginalTerm{深渊}{Bythus}的荣耀与尊荣，产生了一个至完美之美的存有，即圆满之星，那完美的果实——即是耶稣。他们也称他为\Quote{救主}、\Quote{基督}，并依父名称\OriginalTerm{逻各斯}{Logos}和\Quote{万有}，因为他是由所有存有的贡献所形成。然后我们被告知，为示尊崇，与他自己相同本性的天使们同时被产生，作为他的护卫。\par

\SourceRecord{Translated by Alexander Roberts and William Rambaut.；Ante-Nicene Fathers；Vol. 1.；Edited by Alexander Roberts, James Donaldson, and A. Cleveland Coxe.；Buffalo, NY: Christian Literature Publishing Co.,；1885.；<http://www.newadvent.org/fathers/0103102.htm>.}

% structure: p0001 source=h1 target=section reason=page-title

\section{驳斥异端（第一卷，第三章）}

1. 这就是他们对\OriginalTerm{普累若麻}{Pleroma}内所发生之事所作的述说；这就是从那已经命名的、几乎因好奇探究圣父而陷入被普遍实体吸收而灭亡的\OriginalTerm{永世}{Æon}之苦难中流出的灾祸；这就是通过\OriginalTerm{Horos}{Horos}、\OriginalTerm{Stauros}{Stauros}、\OriginalTerm{Lytrotes}{Lytrotes}、\OriginalTerm{Carpistes}{Carpistes}、\OriginalTerm{Horothetes}{Horothetes}和\OriginalTerm{Metagoges}{Metagoges}使那永世从痛苦状态中得以巩固[的过程]。这也是对后来永世起源的述说，即第一个基督和圣神，两者都是由圣父在\OriginalTerm{索非亚}{Sophia}悔改之后产生的；以及第二个基督（他们也称之为救主），他是由[永世们]的共同贡献而产生的。然而，他们告诉我们，这一知识并未公开泄露，因为并非所有人都能接受它，而是由救主通过比喻神秘地启示给那些有资格理解它的人。具体如下。三十个永世由（正如我们已经指出的）救主三十年间没有公开活动以及葡萄园工人的比喻所暗示。他们也断言，保禄非常清楚且频繁地提到这些永世，甚至保留了它们的顺序，当他说：\Quote{直到永世之永世的万代。}不仅如此，我们在感恩礼上念出\Quote{至于永世之永世}（永远永远）时，也陈述了这些永世。最后，无论何处出现\Quote{永世}或\Quote{永世们}这个词，他们立即将它们指涉到这些存在者。\par

2. 再者，永世之十二组的产生，由主在与律法教师辩论时是\Emph{十二}岁\BibleRef{路 2:42}以及宗徒的拣选所暗示，因为宗徒是十二个\BibleRef{路 6:13}。另外十八个永世以这种方式显现：[他们声称]主复活后与门徒交谈了十八个月。他们也断言，这十八个永世由祂名字[\OriginalTerm{耶稣}{\GreekText{Ἰησοῦς}}]的头两个字母，即\Emph{Iota}和\Emph{Eta}所显著暗示。同样，他们断言，十个永世由祂名字开始的字母\Emph{Iota}所指示；出于同样的原因，他们告诉我们救主说：\Quote{一\Emph{Iota}，或一画，决不会过去，直到一切都实现。}\BibleRef{玛 5:18}\par

3. 他们进一步主张，发生在第十二位永世身上的苦难由第十二位宗徒犹达斯的背弃所暗示，也由基督在第十二个月受苦所暗示。因为他们认为，祂受洗后只传教一年。同样的事情也由患血漏病的妇人所最清楚地表明。因为她遭受此病十二年之后，在触摸救主衣边时被救主的到来治愈；因此救主说：\Quote{谁摸了我？}\BibleRef{谷 5:31}——教导祂的门徒们发生在永世中的奥秘，以及那位遭受苦难的永世的治愈。因为患病十二年的妇人代表了那种力量，按照他们的叙述，这种力量的本质在伸展并流入无限；除非她触摸了圣子的衣服，即第一\OriginalTerm{泰特拉得}{Tetrad}的\OriginalTerm{阿勒忒亚}{Aletheia}，由所说的衣边表示，她就会溶解于[她所参与的]普遍实体中。然而，她停住了，不再受苦。因为从圣子发出的力量（他们称此力量为Horos）治愈了她，并将苦难与她分离。\par

4. 他们还断言，救主被表明来自所有的永世，并且祂自己就是\Emph{万物}，通过以下经文：\Quote{凡开胎首生的男性。}\BibleRef{出 13:2} 因为祂，作为万物，开启了那受苦永世的\OriginalTerm{恩图弥西斯}{enthymesis}（当它被逐出普累若麻时）的子宫。他们也称此为第二\OriginalTerm{八元组}{Ogdoad}，我们将稍后讨论。他们声称，保禄显然因此说：\Quote{祂自己就是万有；}\BibleRef{哥 3:11} 又：\Quote{万有归于祂，万有出自于祂；}\BibleRef{罗 11:36} 再：\Quote{在祂内住有整个天主性圆满；}\BibleRef{哥 2:9} 又：\Quote{万有在天主内藉基督而总归。}\BibleRef{弗 1:10} 他们就这样解释圣经中这些和任何类似的段落。\par

5. 他们进一步表明，他们的那个Horos（他们以各种名称称呼它）有两种能力：一是支撑的能力，一是分离的能力；就其支撑和维持而言，它是Stauros，就其划分和分离而言，它是Horos。然后他们表明救主指示了这种双重能力：首先，支撑的能力，当祂说：\Quote{谁不背自己的十字架（Stauros）跟随我，不能做我的门徒；}又：\Quote{背起十字架，跟随我；}\BibleRef{玛 10:21} 但分离的能力，当祂说：\Quote{我来不是送和平，而是送刀剑。}\BibleRef{玛 10:34} 他们还主张若翰指示了同一件事，当他说：\Quote{簸箕在祂手里，祂要彻底扬净禾场，将麦子收进仓里，但将糠秕用不灭的火烧掉。}\BibleRef{路 3:17} 这一声明陈述了Horos的能力。因为他们解释那簸箕就是十字架（Stauros），它无疑燃烧一切物质事物，如火烧糠秕，但它洁净所有得救的人，如簸箕洁净麦子。此外，他们断言宗徒保禄本人用以下言语提到了这个十字架：\Quote{十字架的道理对灭亡的人是愚妄，但对我们得救的人，却是天主的能力。}\BibleRef{格前 1:18} 又：\Quote{我绝不自夸，只夸我们主耶稣基督的十字架，藉着祂，世界于我已被钉在十字架上，我于世界也被钉在十字架上。}\par

6. 他们所有人就是这样描述他们的\OriginalTerm{圆满}{Pleroma}和宇宙的形成，他们竭力将启示的良言曲解以适应他们自己的邪恶发明。他们不仅从福音作者们和宗徒们的著作中，通过歪曲的解释和欺骗性的说明，努力为自己的意见寻找根据；他们也同样对待法律和先知书，这些书卷包含许多比喻和寓言，往往可以按照所采用的解释方法被曲解为多种含义。他们中还有其他人，以极大的诡诈，将圣经的某些部分改编以适应他们自己的虚构，从而俘获那些不能坚守对唯一的天主、全能圣父，以及对唯一的主耶稣基督、天主圣子的坚定信德的人，使他们离开真理。\par

\SourceRecord{Translated by Alexander Roberts and William Rambaut.；Ante-Nicene Fathers；Vol. 1.；Edited by Alexander Roberts, James Donaldson, and A. Cleveland Coxe.；Buffalo, NY: Christian Literature Publishing Co.,；1885.；<http://www.newadvent.org/fathers/0103103.htm>.}

% structure: p0001 source=h1 target=section reason=page-title

\section{\WorkTitle{驳斥异端（第一卷·第四章）}}

1. 以下是他们叙述的发生在\OriginalTerm{普累若麻}{Pleroma}之外的事迹：那位居于高处的\OriginalTerm{索菲亚}{Sophia}（他们又称之为\OriginalTerm{亚卡莫特}{Achamoth}）的\OriginalTerm{思念}{enthymesis}，连同她的\OriginalTerm{情欲}{passion}，从普累若麻中被移去，他们述说，她自然就在那些黑暗与空虚之处（她曾被放逐到那里）变得极其激动。因为她被排除在光明与普累若麻之外，并且没有形状或形态，如同一个未成熟的产物，因为她没有从男性父亲那里得到任何东西。但居于高处的基督怜悯了她；他通过并超越\OriginalTerm{斯陶罗斯}{Stauros}延伸自己，赐予她一个形状，但仅仅是就实质而言，并非赋予悟性。完成后，他撤回自己的影响并返回，留下亚卡莫特独自一人，好使她意识到自己因与普累若麻分离而受苦，从而受到对更好事物的渴望的影响，同时她拥有一种由基督和圣神留在她里面的不朽的芬芳。因此她有两个名字——按她的父亲称为索菲亚（因为索菲亚被认为是她的父亲），以及按与基督同在的那位圣神称为圣神。在获得了一个连同悟性的形状之后，立即被那曾无形地与她同在的\OriginalTerm{逻各斯}{Logos}——即基督——所抛弃，她竭力寻找那遗弃了她的光明，但未能实现目标，因为被\OriginalTerm{浩罗斯}{Horos}阻止。当浩罗斯如此阻碍她进一步前进时，他喊道：\Emph{雅欧}，他们说，由此产生了雅欧这个名字。当她因所陷入的情欲而不能通过浩罗斯，又因独自被留在外面时，她便任凭自己陷入那种她所遭受的多种多样的激情状态；这样她一方面因未得到她渴望的对象而悲伤，另一方面恐惧，唯恐生命本身会离她而去，如同光明已经离去一样，此外，她还处于极大的困惑中。所有这些情感都与无知相关联。她的这种无知不同于她母亲——第一位索菲亚（一位\OriginalTerm{爱昂}{Æon}）——的那种因情欲而堕落的无知，而是因为与生俱来的对知识的对立。此外，另一种情欲落在她（亚卡莫特）身上，即渴望回到那赐予她生命者那里。\par

2. 他们宣称，这一系列情欲就是构成这世界由以形成的质料的实体。因为从她的返回渴望中，属于这世界的每个灵魂以及\OriginalTerm{德穆革}{Demiurge}本人的灵魂都获得了起源。其他一切事物都源自她的恐惧和悲伤。因为从她的眼泪形成了所有液体性质的东西；从她的微笑形成了所有发亮的东西；从她的悲伤和困惑形成了世上所有有形体的元素。因为他们断言，她有时会在黑暗和空虚中因独自被遗弃而哭泣哀伤；有时想到遗弃了她的光明，又会充满喜悦并大笑；然后又陷入恐惧；或者陷入惊愕和迷失。\par

3. 这一切导致了什么？这并非一出轻悲剧，如他们中每一个人都浮夸地解释的那样，一人一种方式，另一个人另一种方式，从何种情欲和何种元素衍生其起源。在我看来，他们有充分的理由不愿公开教导这些事给所有人，而只给那些能付出高价以获得如此深奥奥秘知识的人。因为这些教义绝不同于我们的主所说的：\Quote{你们白白得来的，也要白白分施。}\BibleRef{玛 10:8}相反，它们是晦涩、怪异且深奥的奥秘，只有那些喜爱谎言的人才能经过巨大努力获得。因为谁会不惜花费一切所有，只要他能在回报中学到：从那个陷入情欲的\OriginalTerm{爱昂}{Æon}的\OriginalTerm{思念}{enthymesis}的眼泪中，海洋、泉源、河流以及每种液体物质获得了起源；光明从她的微笑中迸发而出；并且从她的困惑和惊愕中，世界的有形元素得以形成？\par

4. 我本人有些倾向于为发展他们的体系提供一些提示。因为我注意到，水有一部分是淡水，如泉源、河流、雨水等等，有一部分是咸水，如海水，我暗自思忖，所有这些水都不可能来自她的眼泪，因为眼泪只有咸味。因此，显然只有咸水来自她的眼泪。但很可能在剧烈的痛苦和困惑中，她全身出汗。因此，按照他们的观点，我们可以设想泉源、河流以及世上所有淡水都来自这个来源。因为很难相信，我们知道所有眼泪都是同一种性质，咸水和淡水都来自它们。更合理的假设是，有些来自她的眼泪，有些来自她的汗水。既然世上还有某些性质上热且酸的水，你必须猜测它们的起源，如何以及来自何处。这些就是他们假设的一些结果。\par

5. 他们进而陈述说：当母亲\Person{阿卡莫特}经历了各种激情，并好不容易从中脱身后，她转而向那曾离弃她的光明，即基督，恳求。然而，基督已返回\OriginalTerm{普累若麻}{Pleroma}，或许不愿再次降下，便派遣\OriginalTerm{帕拉克莱特}{Paraclete}，即救世主，到那里。这个存在由圣父赋有一切权能，圣父将万有置于他权下，永世亦然，以致\BibleQuote{藉着他，一切有形可见的与无形不可见的，都受了造：诸如王位、神权、主权。}\BibleRef{哥 1:16}然后他被派遣到她那里，同他同时代的天使一起。他们又述说，阿卡莫特充满敬畏，起初因羞怯而遮掩自己，但不久，当她看见他带着全部恩赐，并由他的面容获得力量后，就奔跑上前迎接他。他随后赋予她理智方面的形式，并为她的激情带来治愈，将它们从她身上分离，但并非完全驱逐出思想。因为不能使它们如先前那样化为乌有，因为它们已扎根并获得了力量（以致拥有不可毁灭的存在）。他所能做的只是将它们分离并分别放置，然后掺合、凝聚，以便将无形的激情转化为未成形的质料。他藉此过程赋予它们一种适宜性和本性，使之成为凝结物和形体结构，以便形成两种实体：一是由激情产生的恶的实体，二是由她的转变而来的、虽易受苦但源自她的实体。为此（即由于这种观念质料的实体化），他们说救世主实际上创造了世界。但当阿卡莫特从激情中解脱后，她狂喜地凝视与他同在的天使们的壮丽景象；在神魂超拔中，她因他们而怀孕，据他们告诉我们，她生出了新的存在，一部分是按照她自己的肖像，一部分是按照救世主随从们的肖像的属灵后代。\par

\SourceRecord{Translated by Alexander Roberts and William Rambaut.；Ante-Nicene Fathers；Vol. 1.；Edited by Alexander Roberts, James Donaldson, and A. Cleveland Coxe.；Buffalo, NY: Christian Literature Publishing Co.,；1885.；<http://www.newadvent.org/fathers/0103104.htm>.}

% structure: p0001 source=h1 target=section reason=page-title

\section{驳斥异端（第一卷，第五章）}

1. 因此，按照他们的说法，现在形成了三种存在：一种来自激情，即物质；第二种来自皈依，即属魂的；第三种是她（阿卡摩特）自己产生的，即属神的。接着，她着手赋予这些存在以形式。但她无法成功赋予属神的存在以形式，因为它与她同质。于是她致力于赋予从她自己皈依而出的属魂实体以形式，并启示救主所教导的内容。他们说，她首先从属魂实体中形成了万有之父与王，这万有既包括与他同质的存在，即他们所谓的属魂的（也称为右手的），也包括源于激情和物质的存在（他们称为左手的）。因为他们断言，他在其母的暗中推动下，形成了在他之后产生的一切事物。因此，他们称他为\OriginalTerm{母父}{Metropator}、\OriginalTerm{无父}{Apator}、\OriginalTerm{德谬戈}{Demiurge}和父，说他是右方实体（即属魂的）之父，左方实体（即物质的）的德谬戈，同时又是万有的王。他们说，这个\OriginalTerm{思想}{Enthymesis}渴望为荣耀\OriginalTerm{永世}{Æons}而创造万物，便形成了它们的肖像，或者更确切地说，是救主通过她作成了这些。她、作为不可见之父的肖像，对德谬戈隐藏了自己。而他则是独生子的肖像，由他所造的众天使和总领天使则是其余永世的肖像。\par

2. 因此，他们断言，他被立为普累若麻之外一切事物的父与天主，是一切属魂和属物质实体的创造者。因为他将此前混杂的这两种存在区分开来，从非物质的实体中造出物质的实体，塑造了天上的和地上的事物，成为物质的和属魂的、右手的和左手的、轻的和重的、向上趋向的和向下趋向的事物的\OriginalTerm{工匠}{Demiurge}。他还创造了七重天，他们说德谬戈存在于这七重天之上。因此，他们称他为\OriginalTerm{七元}{Hebdomas}，称他的母亲阿卡摩特为\OriginalTerm{八元}{Ogdoads}，以此保存作为普累若麻的首生且原初的八元之数。此外，他们断言这七重天是有智力的，并称它们为天使，而德谬戈本身则被说成是类似天主的众天使之一。同样，他们声称位于第三层天之上的乐园，是一个拥有权能的第四位天使，亚当与他交谈时从他那里获得了某些特性。\par

3. 他们接着说，德谬戈以为这一切都是他自己创造的，但实际上他是与阿卡摩特的生殖能力共同完成的。他造了天，却不知天；他造了人，却不识人；他使地显现，却对地一无所知。同样，他们声称他对所造的一切形式都不了解，甚至不知道他母亲的存在，反而以为他自己就是一切。他们还断言，他母亲使他产生了这种想法，因为她希望把他造成这样一种性格：他应成为自己本性的源头和首脑，并成为一切（后来进行的）运作的绝对主宰。他们称这位母亲为\OriginalTerm{八元}{Ogdoad}、\OriginalTerm{索菲亚}{Sophia}、\OriginalTerm{地}{Terra}、\OriginalTerm{耶路撒冷}{Jerusalem}、\OriginalTerm{圣神}{Holy Spirit}，并以阳性意义称她为\OriginalTerm{主}{Lord}。她的居所是一个中间地带，位于德谬戈之上，但在普累若麻之下和之外，直至终结。\par

4. 因此，既然他们主张一切物质实体均由三种激情（即恐惧、悲伤和困惑）形成，他们的叙述如下：属魂的实体源于恐惧和皈依；德谬戈也被描述为源于皈依；但所有其他属魂实体（如非理性动物的灵魂、野兽和人的灵魂）的存在，他们都归于恐惧。为此，德谬戈由于不能认识任何属神的本质，便以为自己是独一的天主，并借先知宣告说：\Quote{我是天主，除我以外没有别的。}\BibleRef{依 45:5}他们还教导说，邪恶之灵源于悲伤。因此，魔鬼（他们也称他为\OriginalTerm{宇宙统治者}{Cosmocrator}）、众恶魔和众天使，以及一切存在的邪恶属灵存在，都源于此。他们认为德谬戈是他们母亲（阿卡摩特）的儿子，而宇宙统治者则是德谬戈的造物。宇宙统治者认识在他之上的事物，因为他是邪恶的\Emph{灵}；但德谬戈对这些事无知，因为他仅仅是\Emph{属魂的}。他们的母亲居住在天之上的地方，即中间住所；德谬戈居住在天上的地方，即七元之中；而宇宙统治者则居住在我们这个世界中。世界的物质元素，正如我们之前所说，源于困惑和迷惑，如同来自一个更卑劣的来源。因此，大地源于她的昏沉；水源于她恐惧引起的激动；空气源于她悲伤的固结；而火，产生死亡与腐朽，内在于所有这些元素中，正如他们教导说，无知也隐藏在这三种激情中。\par

5. 这样造了世界之后，他（\OriginalTerm{德谬哥}{Demiurge}）也创造了属土的人，不是取此干燥的尘土，而是取一种由可熔性与流体物质构成的不可见实体。然后，按照他们所定义的过程，向他吹入了其本性中属魂的部分。正是后者（属魂的部分）是按他的肖像和模样造的。其实，属物质的部分，就肖像而言，非常接近天主，但与祂并非同一实体。而属魂的部分，则在模样方面也是如此，因此其实体被称为生命之灵，因为它源自一种属灵的流溢。这一切之后，他们说，他又被一层皮肤包裹，他们所指的就是外在可感觉的肉体。\par

6. 他们还进一步断言，德谬哥本人并不知道他母亲\OriginalTerm{阿卡莫特}{Achamoth}的后裔，这后裔是因她默观那些服侍救主的众天使而产生的，且与她一样，属于属灵的本性。她利用这种无知，将她的产物暗中安置在他里面，使他毫不察觉，为的是藉着他的工具性作用，被注入到从他本身发出的那属魂灵魂中，并像在子宫中那样被承载于这物质身体内，当它逐渐增强力量时，最终能适于接受完全的理性。因此，按照他们的说法，就这样发生了：在德谬哥毫不知情的情况下，由他的吹气所造的人，同时通过一种不可言喻的天主眷顾，因从\OriginalTerm{索菲亚}{Sophia}接受的同步吹气而成了属灵的人。因为，正如他不知道自己的母亲，他也没有认出她的后裔。这个后裔他们宣称就是\OriginalTerm{厄克勒西亚}{Ecclesia}，是天上厄克勒西亚的象征。这就是他们所设想的人：他的属魂灵魂来自德谬哥，他的身体来自尘土，他的肉体来自物质，而他的属灵的人来自母亲阿卡莫特。\par

\SourceRecord{Translated by Alexander Roberts and William Rambaut.；Ante-Nicene Fathers；Vol. 1.；Edited by Alexander Roberts, James Donaldson, and A. Cleveland Coxe.；Buffalo, NY: Christian Literature Publishing Co.,；1885.；<http://www.newadvent.org/fathers/0103105.htm>.}

% structure: p0001 source=h1 target=section reason=page-title

\section{驳斥异端（第一卷，第六章）}

1. 既然有三种物质，他们宣称一切物质（他们也称之为\Quote{左手的}）必然会朽坏，因为它不能接受任何不朽的\OriginalTerm{灵感}{afflatus}。至于一切属魂的存在（他们也称之为\Quote{右手的}），他们认为，既然它介于属灵的和属物质之间，它便倾向于它所偏向的一方。属灵的物质，则被描述为为此目的而被差遣，即在此与属魂的相结合，使其成形，两者同时受到同样的训练。他们宣称这就是\BibleQuote{盐}\BibleRef{玛 5:13}和\BibleQuote{世界的光}。因为属魂的物质需要藉着外在感官的训练；为此他们肯定世界被创造，且救主来到属魂的物质（它拥有自由意志），以便确保它的救恩。因为他们肯定，祂从\Person{阿卡莫特}领取了那些祂要拯救的初果，即属灵的，同时祂被\OriginalTerm{德缪哥}{Demiurge}赋予属魂的基督，并被一种特殊的安排包裹着一个具有属魂本性的身体，但以不可言喻的技巧构造，使之可见、可触，并能忍受苦难。同时，他们否认祂将任何物质（本性）纳入自身，因为物质不能得救。他们还认为，万物的终结将在所有属灵的事物通过\OriginalTerm{诺斯}{Gnosis}被形成和成全时发生；他们所说的属灵的人是指那些对天主有圆满认识、并被\Person{阿卡莫特}引入这些奥秘的人。他们自称就是这些人。\par

2. 属魂的人受教于属魂的事；这些人是指那些藉着他们的行为和单纯的信德而立的人，却没有圆满的知识。他们声称我们教会人士就是这些人。因此他们还主张善工对我们来说是必要的，因为否则我们不可能得救。但至于他们自己，他们认为自己将完全且无疑地得救，不是藉着行为，而是因为他们生来就是属灵的。因为，正如物质不可能有分于救恩（他们确信它无法接受救恩），同样，属灵的物质（他们指自己）也不可能陷入腐朽，无论他们沉溺于何种行为。因为正如金子浸在污秽中并不因此失去它的美丽，而是保持其本质特质，污秽无力损害金子，所以他们肯定他们无论如何不会受到伤害，也不会失去属灵的物质，无论他们可能卷入什么样的物质行为。\par

3. 因此，在他们中\Quote{\OriginalTerm{最完美的}{most perfect}}的人毫无畏惧地沉溺于所有那些禁行之事，圣经向我们保证\BibleQuote{行这样事的人不得继承天主的国}\BibleRef{迦 5:21}。例如，他们毫不犹豫地吃祭偶像的肉，认为这样可以不被玷污。然后，在每一个为偶像举行的异教节庆上，这些人首先聚集；甚至有些人不远离那令天主和人厌恶的血腥场面，在那里角斗士与野兽搏斗或一对一厮杀。其他人则极其贪婪地屈服于肉欲，认为属肉体的事应让给肉体本性，而属神的事则供给属神的人。此外，其中有些人习惯玷污那些他们曾教导上述教义的妇女，正如那些被他们中某些人引入歧途的妇女，在回到天主的教会并承认这一错误连同其他错误时，经常承认的。还有一些人，公开且毫无羞耻地，对某些妇女产生强烈爱恋，诱骗她们离开丈夫，并与她们自己结婚。还有其他人，起初假装与她们如姐妹般同住，持守一切端正，但随着时间的推移，当那位姐妹被她假称的弟兄发现怀孕时，真相败露。\par

他们又犯了许多其他可憎和不虔敬的事，却贬低我们（我们因敬畏天主，连思想言语上也谨防犯罪）为极可鄙、无知的人，同时他们却极力抬高自己，声称自己是完美的，是特选的种子。因为他们宣称，我们只是为使用而领受恩宠，因此它也会再次从我们身上被夺去；但他们自己却拥有恩宠，作为他们自己特有的产业，这恩宠是通过一种不可言喻、不可描述的\OriginalTerm{结合}{conjunction}从上降下的；因此，还要赐给他们更多。\BibleRef{路 19:26}因此，他们主张，无论如何，他们总有必要实践结合的奥秘。为了说服无思想的人相信这一点，他们惯常使用这些话：\Quote{凡在此世\Emph{内}不如此爱一女人而获得她的，不属于真理，也不能达到真理。但凡\Emph{属}此世的，与女人交合，不能达到真理，因为他在情欲之权下如此行动了。}因此，他们告诉我们，我们这些他们称为\Emph{属魂的人}，并描述为\Emph{属}此世的，有必要实践节制和善工，以便借此最终达到中间住所；但对于那些被称为\Quote{属神而完美者}的人，这样的行为举止是完全不必要的。因为并非任何行为能引入\OriginalTerm{圆满}{Pleroma}，而是从那里发出、处于软弱未成熟状态、并在此时此地得以成全的种子。\par

\SourceRecord{Translated by Alexander Roberts and William Rambaut.；Ante-Nicene Fathers；Vol. 1.；Edited by Alexander Roberts, James Donaldson, and A. Cleveland Coxe.；Buffalo, NY: Christian Literature Publishing Co.,；1885.；<http://www.newadvent.org/fathers/0103106.htm>.}

% structure: p0001 source=h1 target=section reason=page-title

\section{驳斥异端（卷一，第七章）}

1. 当所有的种子达到完美时，他们声称，那时他们的母亲\OriginalTerm{阿卡莫特}{Achamoth}将从中间之处离去，进入\OriginalTerm{圆满}{Pleroma}之内，并将接受那出自所有\OriginalTerm{永世}{Æons}的\OriginalTerm{救主}{Saviour}作为她的配偶，这样就在救主与\OriginalTerm{索菲亚}{Sophia}（即阿卡莫特）之间形成一个结合。这些，就是新郎和新娘，而新房是整个圆满的领域。再者，属灵的种子，脱去他们的动物灵魂，成为理智的灵，将以不可抗拒且不可见的方式进入圆满之内，并被作为新娘赐给那些侍奉救主的天使。\OriginalTerm{德穆革}{Demiurge}自己将进入他母亲索菲亚的地方，即中间住所。在这个中间之处，义人的灵魂也将安息；但没有任何动物性质的东西被允许进入圆满。当这些事如所述发生之后，那隐藏在世界中的火将爆发燃烧；在毁灭所有物质的同时，也将与物质一同熄灭，不再存在。他们断言德穆革在救主来临之前不知道这些事中的任何一件。\par

2. 也有一些人对这观点持守：德穆革也产生了基督作为他自己的固有儿子，但本性质为\OriginalTerm{动物的}{animal}，并且先知们曾提及他。这位基督通过玛利亚，如同水通过管子；当祂受圣洗时，属于圆满的救主（由其中所有居民联合形成）以鸽子的形式降临于祂。在祂内也存在那出自阿卡莫特的属灵种子。据此，他们主张，我们的主，在保存首生与首有\OriginalTerm{四元组}{tetrad}的原型时，由这四种实体组成——属灵的（就祂出自阿卡莫特而言），属动物的（就祂出自德穆革，藉由一种特殊安排，因为祂被以不可言喻的技巧形成身体），以及救主（就那降临于祂的鸽子而言）。祂也持续免受一切苦难，因为那不可理解且不可见的祂不可能受苦。为此原因，当祂被带到比拉多面前时，被安置在祂内的基督之灵被取走。他们进一步主张，甚至祂从母亲[阿卡莫特]接受的种子也不受苦；它也是不受苦难的，因为是属灵的，并且对德穆革自己也是不可见的。因此，按照他们，动物的基督和那藉由特殊安排神秘形成的，遭受了苦难，以便母亲可以藉着他展示上面基督的预像，即那通过\OriginalTerm{斯塔乌若斯}{Stauros}伸展自身、并赋予阿卡莫特形状（就实质而言）的那一位。因为他们宣称所有这些事件都是上面发生之事的对应。\par

3. 他们进一步主张，那些拥有\OriginalTerm{阿卡莫特}{Achamoth}种子的灵魂优于其余，并比其它灵魂更被\OriginalTerm{德穆革}{Demiurge}所珍爱，而他不知其真正原因，却想象它们之所以如此是因他的宠爱。因此，他们也说他将它们分配给先知、司铎和君王；并且他们宣称，许多话是由这种子通过先知说出的，因为它具有超凡崇高的本性。母亲也，他们说，说了许多关于上面的事，并且既通过他也通过被他形成的灵魂。然后，他们将预言分为[不同类别]，坚持一部分由母亲说出，第二部分由她的种子说出，第三部分由德穆革说出。同样，他们主张耶稣说了某些话是受救主的影响，某些是受母亲的影响，还有些是受德穆革的影响，正如我们将在本书后面展示的。\par

4. \OriginalTerm{德穆革}{Demiurge}虽然不知道那些高于自己的事，确实被[通过先知]宣布的事所激动，但轻视它们，有时归于一个原因，有时归于另一个原因；或归于预言之灵（它本身拥有自我激动之力），或归于[单纯的]人，或仅是下层[和较卑微的人]的狡猾计谋。他保持无知直到主的显现。但他们叙述，当\OriginalTerm{救主}{Saviour}来了，德穆革从祂学到一切，并喜悦地以他全部权力归附于祂。他们主张他是福音中提到的百夫长，他对救主说了这些话：\Quote{因为我自己也是有权下的人，我有士兵和仆人；我吩咐什么，他们就做什么。}\BibleRef{玛 8:9} 他们进一步主张，只要合适和需要，他将持续管理世事，特别使他能照顾教会；同时他被为他准备的赏报的知识所影响，即他可能达到他母亲的住所。\par

5. 他们因而设想三类人：\OriginalTerm{属灵的}{spiritual}、\OriginalTerm{物质的}{material}、\OriginalTerm{动物的}{animal}，由\OriginalTerm{加音}{Cain}、\OriginalTerm{亚伯尔}{Abel}和\OriginalTerm{舍特}{Seth}代表。这三种本性不再在一个人中发现，而是构成不同种类[的人]。物质的当然归于朽坏。动物的，如果它选择较好的部分，就在中间之处安息；但如果选择较坏的，它也必归于毁灭。但他们断言，由\OriginalTerm{阿卡莫特}{Achamoth}撒下的属灵原则，从那时直到现在在义人的灵魂中被管教和滋养（因为当被她给出时它们还很软弱），最终达到完美，将被赐给\OriginalTerm{救主}{Saviour}的天使作为新娘，而它们的动物灵魂必然永远与\OriginalTerm{德穆革}{Demiurge}在中间之处安息。又进一步细分动物灵魂本身，他们说有些本性善，有些本性恶。善的是那些能够接受[属灵]种子的；本性恶的是那些永远不能接受那种子的。\par

\SourceRecord{Translated by Alexander Roberts and William Rambaut.；Ante-Nicene Fathers；Vol. 1.；Edited by Alexander Roberts, James Donaldson, and A. Cleveland Coxe.；Buffalo, NY: Christian Literature Publishing Co.,；1885.；<http://www.newadvent.org/fathers/0103107.htm>.}

% structure: p0001 source=h1 target=section reason=page-title

\section{驳斥异端（卷一，第八章）}

1. 他们的体系就是如此，既非先知所预言，亦非主所教导，更非宗徒所传授，但他们却夸口说，自己比所有他人对此有更完美的认识。他们的观点并非来自圣经，而是从别处搜罗；并且，正如一句常见的谚语所说，他们努力编织沙绳，同时试图以似乎合理的口吻，将主的比喻、先知的话、宗徒的言辞牵强附会于自己独特的说法，以使自己的体系不至于显得毫无根据。然而，这样做时，他们忽视了圣经的次序与关联，就他们所能而言，肢解并摧毁了真理。通过移接经文、重新装扮、张冠李戴，他们成功地以邪恶的技艺将主的圣言适配于自己的意见，迷惑了许多人。他们的行径就像这样：一位巧匠用珍贵的宝石制作了一幅美丽的国王肖像，然后有人将这幅肖像全部拆散，重新排列宝石，拼凑成狗或狐狸的形状，甚至做得粗劣不堪；然后他却坚持并宣称，\Emph{这}就是那位巧匠制作的美丽国王肖像，指着那些由第一位匠人巧妙地组装成国王形象的宝石，却被后一人糟蹋成狗的形状；他这样展示宝石，欺骗那些不知国王形象为何物的无知之人，并说服他们：那幅可怜的狐狸像实际上就是美丽的国王肖像。同样，这些人拼凑老妇的荒诞故事，然后强行将所遇见的言词、表述和比喻从它们恰当的脉络中扯出来，以适配他们毫无根据的虚构。我们已经谈过他们在\OriginalTerm{普累若麻}{Pleroma}内部如何行事。\par

2. 再者，关于他们普累若麻之外的事物，以下是他们试图从圣经中牵强附会于自己意见的一些例子。他们宣称，主在世界的末期来临，忍受苦难，为的是指示发生在最后一位\OriginalTerm{移涌}{Æons}身上的受难，并借自己的结局宣告在移涌中间所兴起之骚动的止息。他们进一步主张，会堂长十二岁的女儿\BibleRef{路 8:41}，主走近她并使她从死者中复活，乃是\OriginalTerm{阿卡莫特}{Achamoth}的预像；他们的基督通过延伸自己，赋予她形状，并引她重新认识那已离弃她的光。至于救主在她如流产般躺在普累若麻之外时向她显现，他们断言保禄在致格林多人书中：\BibleQuote{最后，也显现了我这个像流产儿的人。}\BibleRef{格前 15:8}同样，救主与其随从降临阿卡莫特，保禄在同一书信中如此宣告：\BibleQuote{女人在头上应有遮盖，因为天使的缘故。}至于阿卡莫特在救主来临时因羞耻而蒙上帕子，梅瑟将帕子蒙在脸上时便已表明。此外，他们还说，她所承受的激情由主在十字架上指示出来。例如，当他说：\BibleQuote{我的天主，我的天主，你为什么舍弃了我？}\BibleRef{玛 27:46}这不过表明\OriginalTerm{索菲亚}{Sophia}被光明遗弃，受到\OriginalTerm{霍洛斯}{Horos}的拦阻而不能前进。她的痛苦再次由他的话指示：\BibleQuote{我的心灵忧闷得要死。}\BibleRef{玛 26:38}她的恐惧由这话：\BibleQuote{父啊！如果可能，就让这杯离开我吧！}\BibleRef{玛 26:39}而她的困惑则由他说：\Quote{我该说什么，我不知道。}\par

3. 他们又教导说，主以如下方式指出了三种人：\Emph{属物质的}，当祂对那问祂的人说：\BibleQuote{你愿跟随我吗？}\BibleRef{路 9:57}\BibleQuote{人子没有枕头的地方}\BibleRef{路 9:58}；\Emph{属魂的}，当祂对那声称\BibleQuote{我要跟随你，但请容许我先告别家里的人}的人说：\BibleQuote{手扶着犁向后看的，不适合天主的国}\BibleRef{路 9:61}（因为他们宣称这个人属于中间类别，正如另一个虽然声称行了大量的义，却拒绝跟随祂，并被财富的爱所胜，以致从未达到成全的人一样——他们乐意将这个人置于属魂的类别）；\Emph{属灵的}，当祂说：\BibleQuote{任凭死人埋葬他们的死人，你去传扬天主的国}\BibleRef{路 9:60}，以及当祂对税吏匝凯说：\BibleQuote{快下来！今天我应当住在你家里}\BibleRef{路 19:5}——因为他们宣称这些属于属灵的类别。又如那女人将酵母藏在三斗面里的比喻，他们宣称揭示了这三个类别。因为按照他们的教导，女人代表\OriginalTerm{索菲亚}{Sophia}；三斗面代表三种人——属灵的、属魂的和属物质的；而酵母则代表救主自己。保禄也很清楚地陈明了属物质的、属魂的和属灵的，在一处说：\BibleQuote{那属于土的怎样，凡属于土的也怎样}\BibleRef{格前 15:48}；在另一处说：\BibleQuote{然而属魂的人不接受圣神的事}\BibleRef{格前 2:14}；又说：\BibleQuote{属灵的人判断万有}\BibleRef{格前 2:15}。他们断言\BibleQuote{属魂的人不接受圣神的事}这句话是针对\OriginalTerm{德谬哥}{Demiurge}说的，因为他既是属魂的，就不认识他那属灵的母亲，也不认识她的种子，也不认识\OriginalTerm{圆满}{Pleroma}中的永世。至于救主收取了祂所要拯救之人的初熟之果，保禄在说：\BibleQuote{如果初熟之果是圣的，面团也是圣的}\BibleRef{罗 11:16}时予以宣布，教导说\Quote{初熟之果}一词所指是属灵的，而\Quote{面团}意指我们，即属魂的教会，那个面团，他们说是祂所取用并融合于自身中的，因为祂就是\Quote{酵母}。\par

4. 此外，他们宣称\OriginalTerm{阿卡莫特}{Achamoth}流浪到\OriginalTerm{圆满}{Pleroma}之外，并从基督那里获得形状，被救主寻觅，主在说祂来寻找那只迷羊时予以指明\BibleRef{路 15:4}。因为他们解释迷羊意指他们的母亲，教会就是由她播种的。流浪本身表示她在圆满之外处于各种情感状态，他们主张物质由此产生。再者，那打扫房屋找到银币的女人，他们宣称代表上面的索菲亚，她失去了自己的\OriginalTerm{思维}{enthymesis}，后来因救主的来临万物得以净化而重新找到。因此，这种实体按他们所说也被重新安置在圆满中。他们也说，西默盎\BibleQuote{把基督抱在怀里，赞美天主说：主啊！现在可照祢的话，让祢的仆人平安去吧}\BibleRef{路 2:28}是德谬哥的预像，他因救主的到来得知了自己的地位改变，而向\OriginalTerm{比特斯}{Bythus}感谢。他们还断言，福音中提到的女先知亚纳\BibleRef{路 2:36}，她与丈夫生活了七年，此后终生守寡，直到看见救主并认出祂，向众人谈论祂，这最清楚地指明了阿卡莫特，她曾短暂地与救主及其同伴相处，此后一直住在中间地带，等候祂再次来临，使她回到自己真正的配偶。她的名字也被救主所指明，当祂说：\BibleQuote{然而智慧由其子女得称为义}\BibleRef{路 7:35}。保禄也在这些话中如此做了：\BibleQuote{我们在成全的人中讲论智慧}\BibleRef{格前 2:6}。他们还宣称，保禄提到了圆满内部的结合，透过一个结合来显明它们；因为论及今生的婚姻结合时，他这样说：\BibleQuote{这是极大的奥秘，但我是指基督和教会说的}\BibleRef{弗 5:32}。\par

5. 此外，他们教导说，主的门徒\Person{若望}指出了第一\OriginalTerm{八元体}{Ogdoad}，并用以下的话表达自己：\Person{若望}，主的门徒，为陈述万有的起源，以说明圣父如何产生了整体，便奠立了一个元始——即天主首生的那一位，他将那一位称为\OriginalTerm{独生子}{Monogenes}和\Person{天主}，圣父以种子的方式在祂内产生了万物。圣言是由祂产生的，众\OriginalTerm{永世}{Æon}的全部实体也在祂内，圣言后来亲自赋以形式。因此，既然他论述万物的首要起源，他理所当然地从元始开始，即从天主和圣言开始。他这样表达自己：\BibleQuote{在元始已有圣言，圣言与天主同在，圣言就是天主；那圣言在元始就与天主同在。}\BibleRef{若 1:1}他在首先区分了这三者——天主、元始和圣言——之后，又将它们合而为一，为显示它们各自的产生，即圣子和圣言的生产，同时显示它们彼此之间以及与圣父的合一。因为\Quote{元始}曾在圣父内，并出自圣父；而\Quote{圣言}曾在元始内，并出自元始。因此，他恰当地说：\BibleQuote{在元始已有圣言，}因为祂在圣子内；\BibleQuote{圣言与天主同在，}因为祂就是元始；\BibleQuote{圣言就是天主，}当然，因为凡由天主所生的就是天主。\BibleQuote{那圣言在元始就与天主同在}——这一句揭示了产生的次序。\BibleQuote{万物是藉着祂而造的；凡受造的，没有一样不是藉着祂造的；}\BibleRef{若 1:3}因为圣言是一切在祂之后出现的众永世的形体与元始的创造者。但\Quote{在祂内有生命，}若望说，\Quote{这生命就是人的光。}这里他又指出了结合；因为他说，万物都是\Emph{藉着}祂造的，但\Emph{在}祂内却有生命。那么，这存在于祂内的，比那些仅藉着祂造的更加紧密地与祂相连，因为它与祂同在，并由祂展开。当他再补充说：\BibleQuote{这生命就是人的光，}他在提及\OriginalTerm{人}{Anthropos}时，也藉那一表述暗示了\OriginalTerm{教会}{Ecclesia}，为藉只用一个名称，显明他们因结合而彼此相通。因为\OriginalTerm{人}{Anthropos}和\OriginalTerm{教会}{Ecclesia}出自\OriginalTerm{圣言}{Logos}和\OriginalTerm{生命}{Zoe}。此外，他将\OriginalTerm{生命}{Zoe}称为人的光，因为他们被她光照，即被塑造并显明。\Person{保禄}也在以下的话中声明这一点：\BibleQuote{因为凡显明出来的，就是光。}\BibleRef{弗 5:13}既然\OriginalTerm{生命}{Zoe}显明并产生了\OriginalTerm{人}{Anthropos}和\OriginalTerm{教会}{Ecclesia}，她就被称为他们的光。如此，若望藉这些话既揭示了其他事，也揭示了第二\OriginalTerm{四元体}{Tetrad}：圣言和生命，人和教会。而且，他还指出了第一\OriginalTerm{四元体}{Tetrad}。因为，在谈论救主并宣告圆满之外的一切事物都由祂赋以形式时，他说祂是整个圆满的果实。因为他称祂为一束\BibleQuote{在黑暗中照耀的光，黑暗却没有理解祂，}\BibleRef{若 1:5}因为当祂给所有源于情欲的事物赋以形式时，黑暗却不认识祂。他还称祂为子、\OriginalTerm{真理}{Aletheia}、\OriginalTerm{生命}{Zoe}，以及\Quote{成了血肉的圣言，}他说道，\InnerQuote{我们看见了祂的光荣；祂的光荣恰如独生者（由圣父所赐）的光荣，满溢恩宠和真理。}\BibleRef{若 1:14}（但\Person{若望}真正所说的却是这样：\BibleQuote{圣言成了血肉，居住在我们中间；我们看见了祂的光荣，正如圣父独生者的光荣，满溢恩宠和真理。}）这样，他（按照他们的说法）在提及圣父、\OriginalTerm{恩宠}{Charis}、\OriginalTerm{独生子}{Monogenes}和\OriginalTerm{真理}{Aletheia}时，清楚地陈述了第一\OriginalTerm{四元体}{Tetrad}。如此，\Person{若望}也讲述了第一\OriginalTerm{八元体}{Ogdoad}，即一切众永世之母。因为他提及了圣父、\OriginalTerm{恩宠}{Charis}、\OriginalTerm{独生子}{Monogenes}、\OriginalTerm{真理}{Aletheia}、\OriginalTerm{圣言}{Logos}、\OriginalTerm{生命}{Zoe}、\OriginalTerm{人}{Anthropos}和\OriginalTerm{教会}{Ecclesia}。这就是\Person{托勒密}的观点。\par

\SourceRecord{Translated by Alexander Roberts and William Rambaut.；Ante-Nicene Fathers；Vol. 1.；Edited by Alexander Roberts, James Donaldson, and A. Cleveland Coxe.；Buffalo, NY: Christian Literature Publishing Co.,；1885.；<http://www.newadvent.org/fathers/0103108.htm>.}

% structure: p0001 source=h1 target=section reason=page-title

\section{驳斥异端（卷一，第九章）}

朋友，你看到了这些人为欺骗自己所采用的方法，他们滥用圣经，试图从中支持自己的体系。为此，我列举了他们的表达方式，好让你明白他们做法的欺骗性和错误的邪恶。因为，首先，如果若望的目的是要提出那至高的\OriginalTerm{八元体}{Ogdoad}，他当然会保持其产生的顺序，并且无疑会先放置首要的\OriginalTerm{四元体}{Tetrad}——按照他们的说法，这是最可敬的——然后附加上第二个\OriginalTerm{四元体}{Tetrad}，这样通过名字的顺序，\OriginalTerm{八元体}{Ogdoad}的次序就能显示出来；而不会在如此长的间隔之后，仿佛一时遗忘又再想起，最后才提到首要的\OriginalTerm{四元体}{Tetrad}。其次，如果他想要指出它们的结合，他肯定不会省略\OriginalTerm{厄克勒西亚}{Ecclesia}的名字；而对于其他的结合，他要么满足于只提男性的\OriginalTerm{永世}{Aeon}（因为其他的，如\OriginalTerm{厄克勒西亚}{Ecclesia}，可以理解），以保持前后一致；要么，如果他列举了其余的结合，他也会宣布\OriginalTerm{安特罗波斯}{Anthropos}的配偶，而不会让我们靠猜测去找出她的名字。\par

这种解释的谬误是显而易见的。因为若望宣告唯一天主、全能者，和唯一耶稣基督、独生子，万物藉祂而造成，他宣告这位是天主子，这位是独生子，这位是万物的形成者，这位是真光，照耀每个人，这位是世界的创造者，这位是来到自己领域的那位，这位是成了血肉并居住在我们中间的那位——这些人却用一种似是而非的解释，歪曲这些陈述，主张有另一位\OriginalTerm{莫诺革涅斯}{Monogenes}，按产生而言，他们也称他为\OriginalTerm{阿尔刻}{Arche}。他们还主张有另一位救主，另一位\OriginalTerm{逻各斯}{Logos}，\OriginalTerm{莫诺革涅斯}{Monogenes}之子，以及另一位基督，为重建\OriginalTerm{圆满体}{Pleroma}而产生。就这样，他们从真理中曲解每一个被引用的表达，滥用这些名字，将它们转移到自己的体系中；因此，按照他们的说法，在所有这些术语中，若望根本没有提及主耶稣基督。因为如果他提到了圣父、\OriginalTerm{卡里斯}{Charis}、\OriginalTerm{莫诺革涅斯}{Monogenes}、\OriginalTerm{阿勒忒亚}{Aletheia}、\OriginalTerm{逻各斯}{Logos}、\OriginalTerm{佐厄}{Zoe}、\OriginalTerm{安特罗波斯}{Anthropos}和\OriginalTerm{厄克勒西亚}{Ecclesia}，按他们的假设，他这样说话指的是首要的\OriginalTerm{八元体}{Ogdoad}，其中还没有耶稣，也没有基督，即若望的老师。但宗徒所说的并不是关于他们的结合，而是关于我们的主耶稣基督，他也承认基督是天主的圣言，这一点他自己已表明。因为总结他先前关于圣言的论述，他进一步宣告：\BibleQuote{圣言成了血肉，居住在我们中间}。但按照他们的假设，圣言根本没有成为血肉，因为祂从未离开过\OriginalTerm{圆满体}{Pleroma}，而是那位救主（由特殊安排从所有\OriginalTerm{永世}{Aeon}中形成，比圣言晚出现）成了血肉。\par

因此，愚昧的人哪，你们要知道：为我们受苦并居住在我们中间的耶稣，就是天主的圣言。因为如果有任何其他\OriginalTerm{永世}{Aeon}为了我们的救恩而成了血肉，那么宗徒很可能是在说另一位。但如果圣父的圣言下降，也就是那上升的一位，即唯一天主的独生子，为了人的缘故，按圣父的旨意成了血肉，那么宗徒当然不是在说任何其他，也不是关于任何\OriginalTerm{八元体}{Ogdoad}，而是关于我们的主耶稣基督。因为按照他们的说法，圣言原本并没有成为血肉。因为他们主张，救主取了一个动物性的身体，是由一种不可言喻的上智安排按特殊造化而形成的，以便成为可见可触的。但\Quote{血肉}是指古时天主用尘土为亚当形成的那个，而若望宣布天主的圣言成了这血肉。这样，他们的首要和首生的\OriginalTerm{八元体}{Ogdoad}就归于无有。因为既然\OriginalTerm{逻各斯}{Logos}、\OriginalTerm{莫诺革涅斯}{Monogenes}、\OriginalTerm{佐厄}{Zoe}、\OriginalTerm{佛斯}{Phōs}、\OriginalTerm{索忒耳}{Soter}、\OriginalTerm{基督}{Christus}、天主子，以及为我们成了肉身的那一位，已经被证明是同一的，他们所建立的\OriginalTerm{八元体}{Ogdoad}就立刻瓦解了。当这个被摧毁时，他们的整个体系就陷入废墟——一个他们虚假地梦想出来的体系，从而伤害圣经，同时建立自己的假设。\par

然后，他们收集散见于圣经各处的一组表达和名字，如我们所说，将它们从自然的意义扭曲为不自然的意义。他们这样做，就像那些提出任何他们幻想的假设，然后试图从\Person{荷马}的诗歌中支持它的人一样，以致无知的人以为\Person{荷马}实际上创作了那些符合这个假设的诗句，而这个假设其实是新编造的；许多人被这些诗句的规律性序列所引导，以致怀疑\Person{荷马}是否真的创作了它们。例如下面这段，有人描述\Person{赫拉克勒斯}被\Person{欧律斯透斯}派往地狱之犬那里，通过以下荷马诗行来进行——我们引用这些作为例子并无不妥，因为两边的尝试是类似的：——\par

\Quote{如此说着，他从家中出来，深深叹息。}— \WorkTitle{Od.}, x. 76.\newline{}\Quote{英雄赫拉克勒斯，与伟大的事迹相熟。}— \WorkTitle{Od.}, xxi. 26.\newline{}\Quote{欧律斯透斯，斯忒涅卢斯之子，珀耳修斯之后裔。}— \WorkTitle{Il.}, xix. 123.\newline{}\Quote{为使他从厄瑞玻斯带来阴森的普路托的犬。}— \WorkTitle{Il.}, viii. 368.\newline{}\Quote{他前进如山中养育的狮子，充满力量的自信。}— \WorkTitle{Od.}, vi. 130.\newline{}\Quote{迅速穿过城市，所有朋友跟随其后。}— \WorkTitle{Il.}, xxiv. 327.\newline{}\Quote{少女与青年，以及忍受苦难的老者。}— \WorkTitle{Od.}, xi. 38.\newline{}\Quote{为他悲恸，如同走向死亡之人。}— \WorkTitle{Il.}, xxiv. 328.\newline{}\Quote{但墨丘利与蓝眼明眸的弥涅耳瓦引导着他。}— \WorkTitle{Od.}, xi. 626.\newline{}\Quote{因她知道她兄弟的心思，如何因忧愁而劳苦。}— \WorkTitle{Il.}, ii. 409.\par

试问，有哪个心思单纯的人不会被这样的诗句所迷惑，以为荷马确实是根据所指的主题编排了它们？但熟悉荷马著作的人会辨认出这些诗句，却不会识别它们所应用的主题，因为他知道其中一些是关于尤利西斯的，另一些是关于赫拉克勒斯本人的，还有一些是关于普里阿摩的，另外一些则是关于墨涅拉俄斯和阿伽门农的。然而，若有人把这些诗句取来，将每一句放回它恰当的位置，他立刻就会摧毁这个叙述。同样，那藉着圣洗在心中持守真理的准则、毫不动摇的人，无疑能认出取自圣经的名词、语句和比喻，却绝不会认可这些人所做的亵渎性运用。因为，虽然他会认出宝石，但他绝不会接受狐狸来替代君王的肖像。但当他把所引用的每个语句都恢复到恰当的位置，使之与真理的整体相适合时，他就会揭露并证明这些异端者的虚构是毫无根据的。\par

5. 但由于这场表演还缺少致命一击，使得任何人都可以在追随其闹剧至终后，随即附加一个推翻它的论证，我们认为最好先指出这个寓言本身的原创者们彼此之间在哪些方面存在分歧，仿佛他们受了不同错误之灵的感动。因为这一事实本身就构成了一个\Emph{先验}的证明：教会所宣告的真理是不可动摇的，而这些人的理论不过是一堆谎言的编织。\par

\SourceRecord{Translated by Alexander Roberts and William Rambaut.；Ante-Nicene Fathers；Vol. 1.；Edited by Alexander Roberts, James Donaldson, and A. Cleveland Coxe.；Buffalo, NY: Christian Literature Publishing Co.,；1885.；<http://www.newadvent.org/fathers/0103109.htm>.}

% structure: p0001 source=h1 target=section reason=page-title

\section{驳斥异端（卷一，第十章）}

1. 教会虽然分散于普世，直至地极，但从\OriginalTerm{宗徒}{apostles}及其门徒接受了这\OriginalTerm{信德}{faith}：她信唯一的天主，全能的圣父，天地的创造者，海洋及其中万物的创造者；并信唯一的基督耶稣，天主子，为了我们的\OriginalTerm{救恩}{salvation}而\OriginalTerm{降生成人}{incarnate}；并信\OriginalTerm{圣神}{Holy Spirit}，祂曾借着先知们宣讲天主的诸般分施、来临、由童贞女所生、受难、从死者中\OriginalTerm{复活}{resurrection}、带着肉身升入天堂，以及所爱的基督耶稣，我们的主，将来在圣父的光荣中从天上显现，\BibleQuote{使万有归于一体，\BibleRef{弗 1:10}} 并重新兴起全人类的全体肉身，为使对基督耶稣，我们的主、天主、救主、君王，按照无形圣父的旨意，\BibleQuote{一切在天上的、地上的和地下的，无不屈膝叩拜，一切唇舌无不承认，\BibleRef{斐 2:10}} 祂并向众人施行公正的审判；祂要将\BibleQuote{邪恶的灵，\BibleRef{弗 6:12}} 以及那些悖逆堕落的天使，连同人类中不敬虔、不义、邪恶、亵渎的人，投入永远的火里；但祂要运用祂的\OriginalTerm{恩宠}{grace}，赐予义人、圣洁的人、遵守祂诫命、持守祂的爱（有些从起初，有些从悔改之时起）的人以不朽，并以永远的光荣围绕他们。\par

2. 正如我已观察到的，教会接受了这宣讲和这信德，虽然分散于普世，却好似住在一所房子里，谨慎地保存它。她也相信这些道理，就好像她只有一个灵魂、同一颗心，并以完全的和谐宣讲、教导、传授它们，仿佛她只有一张口。因为尽管世界的语言各不相同，传承的要旨却是一致相同的。因为在\Place{日耳曼}建立的教会，并不相信或传授任何不同的东西，在\Place{西班牙}、\Place{高卢}、\Place{东方}、\Place{埃及}、\Place{利比亚}，以及在世界中央区域建立的教会，也是如此。但正如天主的受造物——太阳，在普世是同一的，真理的宣讲也到处照耀，并光照所有愿意认识真理的人。教会中的首领，无论他在口才上多么有天赋，也不会教导异于这些的道理（因为没有人大于师傅）；另一方面，口才不足的人也不会损害传承。因为信德始终是同一的，能长篇大论讲论它的人，不会对它增添什么；言语寥寥的人，也不会减少它。\par

3. 不能因为人们拥有不同程度的智慧，就因此改变信仰本身的实质，去设想在这宇宙的塑造者、创造者和保存者之外另有别的神（仿佛祂对他们不够），或另有一个基督、另一个\OriginalTerm{独生子}{Only-begotten}。这里所指的仅仅是：一个人可能比其他人在更准确地阐述那些用比喻所说的事物的含义，并使它们适应信仰的整体计划；以及特别清晰地解释与人类救恩相关的天主的作为和\OriginalTerm{分施}{dispensation}；并显示天主在对待悖逆堕落的天使以及人类的悖逆时，显出了恒久忍耐；并阐明为何同一个天主使一些事物暂时、一些永远，一些属天、一些属地；并理解为什么天主虽不可见，却向先知们显现，且不是以同一形象，而是对不同的人以不同的形式；并显示为何多次\OriginalTerm{盟约}{covenants}赐予人类；并教导每一盟约的特殊性质；并探求为何\BibleQuote{天主\BibleRef{罗 11:32}将众人都禁锢于不信之中，为要怜悯众人}；并感恩地描述\OriginalTerm{圣言}{Word of God}为何降生成人并受苦；并讲述天主子为何在这末世，即终结时，而非在起初来临；并展开圣经中关于终结和未来之事的含意；并不可缄默于天主如何使外邦人（他们的救恩本已绝望）同为承继者、同为一体、同享圣人的福分；并论述\BibleQuote{这必朽的身体如何穿着不朽，这必朽坏的如何穿着不朽坏\BibleRef{格前 15:54}}，并宣告天主在何种意义上说：\BibleQuote{那本来不是我子民的，我要称为我的子民；本来不是蒙爱的，我要称为蒙爱的\BibleRef{欧 2:23}}，以及在何种意义上祂说：\BibleQuote{不怀孕、不生养的，你要欢呼；因为弃妇的儿女，比有丈夫的儿女更多\BibleRef{依 54:1}}。因为关于这些点和类似的事，宗徒喊道：\BibleQuote{啊！天主的智慧和知识的富饶何其深奥！祂的审判何其难测，祂的道路何其难寻\BibleRef{罗 11:33}}。但这里所指的高超技巧并不在于：有人竟在创造者和塑造者之外，设想一个迷途\OriginalTerm{永世}{Æon}的\OriginalTerm{思想}{Enthymesis}，即他们共同之母，并以此达到如此亵渎的地步；也不在于：有人又虚妄地想象，在这之上还有一个\OriginalTerm{圆满}{Pleroma}，有时被认为包含三十个，有时又是无数个永世，正如这些缺乏真正神性智慧的教师所主张的；而\OriginalTerm{天主教会}{Catholic Church}却在普世拥有同一的信德，如我们已说过的。\par

\SourceRecord{Translated by Alexander Roberts and William Rambaut.；Ante-Nicene Fathers；Vol. 1.；Edited by Alexander Roberts, James Donaldson, and A. Cleveland Coxe.；Buffalo, NY: Christian Literature Publishing Co.,；1885.；<http://www.newadvent.org/fathers/0103110.htm>.}

% structure: p0001 source=h1 target=section reason=page-title

\section{\WorkTitle{驳斥异端（第一卷，第十一章）}}

1. 现在让我们来看看那些异端者（大约有两三种）自相矛盾的观点，他们在处理相同问题时彼此不一致，反而在事物和名称上同样提出相互抵触的见解。他们中的第一人，\Person{瓦伦廷}，将被称为\Quote{\OriginalTerm{诺斯替}{Gnostic}}的异端原则应用于其学派的独特特征，教导如下：他主张存在某种不可用任何名称表达的\OriginalTerm{二元}{Dyad}（双重存有），其中一部分应称为\OriginalTerm{不可言者}{Arrhetus}，另一部分称为\OriginalTerm{沉默}{Sige}。从这二元中产生了第二个二元，一部分他命名为\OriginalTerm{父}{Pater}，另一部分命名为\OriginalTerm{真理}{Aletheia}。从这个\OriginalTerm{四元体}{Tetrad}又生出\OriginalTerm{逻各斯}{Logos}和\OriginalTerm{生命}{Zoe}、\OriginalTerm{人}{Anthropos}和\OriginalTerm{教会}{Ecclesia}。这些构成最初的\OriginalTerm{八元体}{Ogdoad}。他接着说，从逻各斯和生命产生了十种能力，正如我们之前所提的。但从人和教会产生了十二种，其中一种分离出去，脱离其原初状态，产生了其余一切受造之物。他还设想了两种名为\OriginalTerm{界限}{Horos}的存在，一个位于\OriginalTerm{深渊}{Bythus}与\OriginalTerm{圆满}{Pleroma}其余部分之间，分隔受造的\OriginalTerm{永恒者}{Æons}与非受造的圣父；另一个则将他们的母亲与圆满分开。\Person{基督}也不是从圆满内的永恒者产生的，而是由被排除在圆满之外的母亲在回忆更美好事物的作用下产生的，但并非没有某种阴影。他作为男性，将阴影从自身分离后回到了圆满；但他的母亲留在阴影中，被剥夺了灵性本质，又生了另一个儿子，即\OriginalTerm{工匠神}{Demiurge}，他也称其为一切隶属之物的至高统治者。他还断言，与工匠神一同产生了左手的能力，在这一点上他同意那些假称为\OriginalTerm{诺斯替}{Gnostic}的人，我们稍后将谈到他们。有时他又主张\Person{耶稣}是从那位与母亲分离并与其余部分联合者（即\OriginalTerm{意愿}{Theletus}）产生的；有时说从那位回到圆满者（即基督）产生；其他时候又说来自人和教会。他宣称\OriginalTerm{圣神}{Holy Spirit}是由真理为检视和使永恒者结果实而生的，他无形地进入他们中，这样永恒者就生出了真理的植物。\par

2. \Person{塞孔杜斯}则断言最初的\OriginalTerm{八元体}{Ogdoad}由右手和左手的两个\OriginalTerm{四元体}{Tetrad}组成，并教导其中之一称为\OriginalTerm{光明}{light}，另一称\OriginalTerm{黑暗}{darkness}。但他主张那分离出去并堕落的能力并非直接来自三十个\OriginalTerm{永恒者}{Æons}，而是来自他们的果实。\par

3. 另有一位在他们中著名的大师，竭力追求某种更崇高者，欲达到更高知识，这样解释了最初的\OriginalTerm{四元体}{Tetrad}：存在一种先于万物的\OriginalTerm{原始开端}{Proarche}，超越一切思想、言说和名称，我称之为\OriginalTerm{单一}{Monotes}。与这单一并存一种能力，我又称之为\OriginalTerm{统一}{Henotes}。这统一和单一成为一体，产生了（但并未分离自身而产出）万物的开端，一个理智、非受造、不可见的存在，语言称之为\Quote{\OriginalTerm{单一者}{Monad}}。与这单一者共存一种同质的能力，我又称之为\OriginalTerm{一}{Hen}。这些能力——单一、统一、单一者、一——产生了其余所有的\OriginalTerm{永恒者}{Æons}。\par

4. 唉，唉！呜呼，呜呼！——我们完全有理由在这些悲剧性的感叹中，为他在编造名称方面所显示的这般厚颜无耻而惊叹，他为自己的谬误体系设计了一套命名法，毫无愧色。因为他宣称：存在一个先于万物的\OriginalTerm{原始开端}{Proarche}，超越一切思想，我称之为\OriginalTerm{单一}{Monotes}；又与这单一并存一种能力，我也称之为\OriginalTerm{统一}{Henotes}——这最清楚地表明他所言皆是他自己的发明，他亲自为他的事物体系命名，而这些名称此前从未有人提出。同样明显的是，正是他本人有足够的胆量杜撰了这些名称；因此，除非\Emph{他}出现在世上，真理否则将始终无名。但既然如此，在处理同一主题时，任何人都可以毫无阻碍地按如下方式赋予名称：存在一个王家的\OriginalTerm{原始开端}{Proarche}，超越一切思想，一种先于一切其他实体的能力，向四面八方扩展于空间中。但与它并存一种能力，我称之为\OriginalTerm{瓜}{Gourd}；与这瓜并存一种能力，我又称之为\OriginalTerm{虚空}{Utter-Emptiness}。这瓜和虚空成为一体，产生了（但并非简单地产生，以致分离自身）一个果实，随处可见，可食且美味，果实语言称之为\OriginalTerm{黄瓜}{Cucumber}。与这黄瓜并存一种同质的能力，我又称之为\OriginalTerm{甜瓜}{Melon}。这些能力——\OriginalTerm{瓜}{Gourd}、\OriginalTerm{虚空}{Utter-Emptiness}、\OriginalTerm{黄瓜}{Cucumber}和\OriginalTerm{甜瓜}{Melon}——生出了瓦伦廷的其余狂乱甜瓜之群。因为如果用于描述宇宙的语言可以转用于首要的四元体，并且任何人都可以随意命名，谁又能阻止我们采用这些名字呢？它们比那些名字更可信，更通用，且为人所理解。\par

5. 然而，另一些人却称他们的首要和首生的\OriginalTerm{奥格多亚德}{Ogdoad}为以下名称：第一，\OriginalTerm{普罗阿尔刻}{Proarche}；第二，\OriginalTerm{阿内诺厄托斯}{Anennoetos}；第三，\OriginalTerm{阿瑞托斯}{Arrhetos}；第四，\OriginalTerm{阿奥拉托斯}{Aoratos}。然后，从第一\OriginalTerm{普罗阿尔刻}{Proarche}，产生了第一和第五位的\OriginalTerm{阿尔刻}{Arche}；从\OriginalTerm{阿内诺厄托斯}{Anennoetos}，产生了第二和第六位的\OriginalTerm{阿卡塔勒普托斯}{Acataleptos}；从\OriginalTerm{阿瑞托斯}{Arrhetos}，产生了第三和第七位的\OriginalTerm{阿诺诺马斯托斯}{Anonomastos}；从\OriginalTerm{阿奥拉托斯}{Aoratos}，产生了第四和第八位的\OriginalTerm{阿格内托斯}{Agennetos}。这就是第一个\OriginalTerm{奥格多亚德}{Ogdoad}的\OriginalTerm{普雷若玛}{Pleroma}。他们主张这些权能先于\OriginalTerm{彼托斯}{Bythus}和\OriginalTerm{西革}{Sige}，为的是显得比完美者更完美，比那些诺斯替派本人更博学！对此，人可正当地喊道：\Quote{你们这些微不足道的诡辩家！}因为，即使关于\OriginalTerm{彼托斯}{Bythus}本身，他们中间也有许多相互矛盾的见解。因为有人宣称他没有配偶，既非男亦非女，事实上什么都不是；另一些人则断言他是阴阳同体，赋予他雌雄同体的本性；还有一些人则将\OriginalTerm{西革}{Sige}许配给他为妻，以便形成第一个结合。\par

\SourceRecord{Translated by Alexander Roberts and William Rambaut.；Ante-Nicene Fathers；Vol. 1.；Edited by Alexander Roberts, James Donaldson, and A. Cleveland Coxe.；Buffalo, NY: Christian Literature Publishing Co.,；1885.；<http://www.newadvent.org/fathers/0103111.htm>.}

% structure: p0001 source=h1 target=section reason=page-title

\section{\WorkTitle{驳斥异端}（卷一，第十二章）}

1. 但是托勒密的追随者说，他（\Person{比托斯}）有两个配偶，他们也称之为\OriginalTerm{意趣}{Diatheses}（affections），即\OriginalTerm{意念}{Ennœa}和\OriginalTerm{意志}{Thelesis}。因为他们断言，他首先构思了产生某物的思想，然后以此意愿行动。因此，这两个情感或能力，即意念和意志，仿佛彼此结合，\OriginalTerm{独生者}{Monogenes}和\OriginalTerm{真理}{Aletheia}便据此结合而产生。这两个作为父两个情感的典型和肖像出现——那些不可见之物的可见代表——\OriginalTerm{努斯}{Nous}（即独生者）出自意志，而真理出自意念，因此由意志产生的形象是男性的，而由意念产生的则是女性的。因此，\OriginalTerm{意志}{Thelesis}（will）仿佛成了\OriginalTerm{意念}{Ennœa}（thought）的一种能力。因为意念不断渴望后代；但凭自己无法生出所渴望之物。但当意志的能力（意志的机能）降临于她时，她便生出了她所孕育的那一位。\par

2. 这些幻想的存在（如同荷马笔下的\Person{宙斯}，他被描绘成焦虑不眠之夜，谋划着如何荣耀\Person{阿基里斯}并毁灭许多希腊人），亲爱的朋友，在你看来，不会比宇宙的天主拥有更多的知识。祂一旦思考，也同时完成祂所意愿的；祂一旦意愿，也同时思考祂所意愿的；祂意时即思，思时即意，由于祂完全是思想，[全是意志，全是理智，全是光明]，全是眼睛，全是耳朵，是一切美善事物唯一的完整源泉。\par

3. 那些被认为比刚才提到的人更有技巧的人说，第一个\OriginalTerm{八度}{Ogdoad}并非逐渐产生，一个\OriginalTerm{移涌}{Æon}由另一个发出，而是所有移涌由\OriginalTerm{原父}{Propator}及其\OriginalTerm{意念}{Ennœa}同时产生。\Person{科罗巴苏斯}（Colorbasus）对此断言如同他亲临其出生一样肯定。因此，他和他的追随者主张，\OriginalTerm{人}{Anthropos}和\OriginalTerm{教会}{Ecclesia}并非如其他人所认为的那样来自\OriginalTerm{圣言}{Logos}和\OriginalTerm{生命}{Zoe}；相反，\OriginalTerm{圣言}{Logos}和\OriginalTerm{生命}{Zoe}来自\OriginalTerm{人}{Anthropos}和\OriginalTerm{教会}{Ecclesia}。但他们以另一种形式表达如下：当原父构思产生某物的思想时，他得到了\Emph{父}的名称。但因为他所产生的是\Emph{真实的}，便得名为\OriginalTerm{真理}{Aletheia}。同样，当他希望启示自己时，这被称为\OriginalTerm{人}{Anthropos}。最后，当他产生了那些他先前所思想的存在时，这些便被称为\OriginalTerm{教会}{Ecclesia}。\OriginalTerm{人}{Anthropos}通过言语形成了\OriginalTerm{圣言}{Logos}：这是首生子。但\OriginalTerm{生命}{Zoe}紧随\OriginalTerm{圣言}{Logos}之后；这样第一个\OriginalTerm{八度}{Ogdoad}便完成了。\par

4. 他们对救主也有很大争论。有些人认为他是从所有移涌中形成的，因此被称为\OriginalTerm{意悦者}{Eudocetos}，因为整个\OriginalTerm{圆满}{Pleroma}通过他\Emph{满意}于荣耀父。但其他人断言他仅由出自\OriginalTerm{圣言}{Logos}和\OriginalTerm{生命}{Zoe}的十位移涌产生，因此被称为\OriginalTerm{圣言}{Logos}和\OriginalTerm{生命}{Zoe}，保留了祖名。又有人主张他出自从\OriginalTerm{人}{Anthropos}和\OriginalTerm{教会}{Ecclesia}生出的十二位移涌，因此他自称为\Quote{人子}。还有人断言他由\OriginalTerm{基督}{Christ}和\OriginalTerm{圣神}{Holy Spirit}产生（他们是为圆满的安全而产生的），因此被称为\OriginalTerm{基督}{Christ}，保留了产生他的父的名称。还有一些人宣称整个\OriginalTerm{原父}{Propator}、\OriginalTerm{初始元始}{Proarche}、\OriginalTerm{前无概念者}{Proanennoetos}被称为\OriginalTerm{人}{Anthropos}；这就是那伟大而深奥的奥秘，即超越一切、包容万有的能力被称为\OriginalTerm{人}{Anthropos}；因此救主自称为\Quote{人子}。\par

\SourceRecord{Translated by Alexander Roberts and William Rambaut.；Ante-Nicene Fathers；Vol. 1.；Edited by Alexander Roberts, James Donaldson, and A. Cleveland Coxe.；Buffalo, NY: Christian Literature Publishing Co.,；1885.；<http://www.newadvent.org/fathers/0103112.htm>.}

% structure: p0001 source=h1 target=section reason=page-title

\section{驳斥异端（第一卷，第十三章）}

1. 但在这些异端者中另有一人，名叫\OriginalTerm{马库斯}{Marcus}，他自夸改进了其老师的教义。此人精通魔术骗术，借此吸引了大批男子，也有不少妇女，诱使他们加入自己，视他为拥有最高知识和成全之人，并从不可见、不可言说的上界领受了至高权能。如此看来，他似乎是\OriginalTerm{假基督}{Antichrist}的真正先驱。因为他将\Person{阿那克西劳斯}的把戏与所谓的\OriginalTerm{术士}{magi}的诡诈结合起来，被他那些无知狂妄的追随者视为借此施行奇迹。\par

2. 他假装祝圣搀了酒的杯，并刻意延长呼求之词，设法使杯呈现紫红色，好让人以为那超越万有的\OriginalTerm{卡里斯}{Charis}藉他的呼求将自己的血滴入杯中，从而令在场的人欢喜地品尝那杯，以为这样，这术士所呈现的\OriginalTerm{卡里斯}{Charis}也会流入他们体内。他又将搀了酒的杯递给妇女，命她们在他面前祝圣。这事之后，他亲自取出另一只远大于那被迷惑的妇女所祝圣的杯，将那较小的杯（由妇女祝圣）倾入他自己带来的大杯中，同时念道：\Quote{愿那超越万有、超越一切知识与言语的\OriginalTerm{卡里斯}{Charis}充满你们的内在人，并将在你们内加增她的知识，如同将芥菜种子撒在好土中一样。}接着又说了一些类似的话，以此刺激那可怜的女人，随后大杯竟从小杯盛满，甚至因所得而溢出，他便显得行了一件奇事。他还做了许多类似的事，完全欺骗了许多人，并将他们引离了正道。\par

3. 此人似乎确实有一个邪魔作为他的护符，藉此他好像能说预言，也能使那些他自己认为配得上分享他的\OriginalTerm{卡里斯}{Charis}的人说预言。他特别致力于妇女，尤其是那些出身高贵、衣着雅致、家资丰厚者，常用以下诱人的言语勾引她们：\Quote{我切望使你成为我的\OriginalTerm{卡里斯}{Charis}的分享者，因为万有之父常在祂面容前注视你的天使。你的天使的位置在我们中間：我们应当成为一体。首先从我和通过我接受\OriginalTerm{卡里斯}{Charis}的恩赐。你当装扮自己，如同期待新郎的新娘，好使你成为我所是的，我成为你所是的。在你们的洞房中建立光的种子。从我这里接受一个配偶，成为他的收容所，同时被他收容。看，\OriginalTerm{卡里斯}{Charis}已降临到你身上；开口说预言吧。}那女人回答说：\Quote{我从未说过预言，也不知道如何说预言；}于是他再次行某种呼求，以惊骇那受骗的牺牲品，对她说：\Quote{开口，说出你心中所想，你必要说预言。}她于是被这些话空虚无谓地膨胀得意，灵魂因期待自己将说预言而极度激动，心跳剧烈，达到所需的胆量，信口开河，不知羞耻地说出一些偶然想到的胡言乱语，正如一个被虚空气体激动的人所可能说的。（关于此事，一位比我高明的人曾说过：灵魂被虚空气体激动时，既狂妄又无耻。）从此她自以为是女先知，并感谢马库斯使她分享了他自己的\OriginalTerm{卡里斯}{Charis}。然后她努力回报他，不仅捐出自己的财产（他借此聚集了巨额财富），而且献出自己的身体，渴望以各种方式与他结合，好能完全与他成为一体。\par

4. 但已经有几位最忠信、持守对天主的敬畏且未被欺骗的妇女（他仍尽力像引诱其余人一样，命她们说预言以诱惑她们），憎恶并咒骂他，从这无耻的狂欢团体中退出。她们这样做，是因为清楚地知道，预言之恩并非由术士马库斯赐予人，唯有那些从天主那里领受了祂从上而来的恩宠的人，才拥有天主所赐的说预言之能；并且他们只在天主所愿的时间和地点说话，而不是在马库斯命令他们的时候。因为命令者比受命者更大、更有权威，前者管辖，后者则处于服从地位。因此，如果马库斯或任何其他人发出命令——如同这些人常在他们宴会上玩抽签游戏，并按照签令彼此命令说预言，说出符合自己欲望的神谕——那么，命令者（虽然不过是一个人）就会比预言之神更大、更有权威，这是不可能的。但这些受他们命令、在他们愿意时讲话的灵，是属地的、软弱的、狂妄无耻的，由撒殚派来，为诱惑和毁灭那些不持守起初从教会所领受的坚固信德的人。\par

5. 此外，这位马可调配春药和爱情魔药，用以侮辱某些妇女（即便不是全部），那些已经回到天主的教会——此事经常发生——的妇女已承认此事，也告明自己被他玷污，并对他的情欲炽烈。可悲的例子发生在某个亚细亚人身上，他是我教会的执事之一，曾接待马可到家中。他的妻子，一位容貌出众的女子，身心皆沦为这术士的猎物，并与他同行多时。最后，在弟兄们费尽周折才使她归正后，她终日公开告解，为从这术士所受的玷污哀哭悲叹。\par

6. 他的一些门徒也沉迷于同样的行径，欺骗了许多愚昧的妇女并玷污了她们。他们自称为\Quote{成全者}，以致无人在知识的广大上能与他们相比，即使你提及保禄、伯多禄或任何其他宗徒也不行。他们声称自己比所有人知道得更多，唯独他们吸收了那不可言说之能的伟大知识。他们还认为自己已超越一切权能之上，因此在各方面都可以任意而行，无所畏惧。因为他们断言，因着\Quote{救赎}，他们既不能被审判者捕拿，甚至不能被看见。但假如审判者偶然抓住了他们，他们只需站在审判者面前，连同\Quote{救赎}一起说出这些话：\InnerQuote{坐在天主旁边者，奥妙的永恒\OriginalTerm{沉寂}{Sige}，藉着你那权能的天使们（大能者）不断瞻仰圣父的容貌，以你为向导和引介者，从高处获得形体，而她在伟大的胆识中，因着\OriginalTerm{父}{Propator}的良善而赋予心智，生出了我们作为她的肖像——她其时如在梦中般专注于高处的事——看，审判者已临近，传令官命我答辩。但你既通晓双方之事，请向审判者陈明我们双方之案，因为实际上这只是一个案子。}当母亲听到这些话，便将荷马史诗中冥王的隐身头盔戴在他们头上，使他们能隐形逃脱审判者。然后她立刻抓住他们，领他们进入洞房，交给他们的配偶。\par

7. 在我们隆河流域，他们正是以这样的言行欺骗了许多妇女，她们的良心如同被热铁烙过。\BibleRef{弟后 3:6} 她们中有些人公开告明自己的罪过；但另一些人羞于如此，在沉默中绝望于获得天主的生命，有人完全背弃了信仰；还有人在这两种态度间犹豫不决，陷于俗语所说的\Quote{既不在外也不在内}，从这\Quote{知识之子}的种子中结出这样的果实。\par

\SourceRecord{Translated by Alexander Roberts and William Rambaut.；Ante-Nicene Fathers；Vol. 1.；Edited by Alexander Roberts, James Donaldson, and A. Cleveland Coxe.；Buffalo, NY: Christian Literature Publishing Co.,；1885.；<http://www.newadvent.org/fathers/0103113.htm>.}

% structure: p0001 source=h1 target=section reason=page-title

\section{\WorkTitle{驳斥异端（卷一，第十四章）}}

1. 这位玛尔谷于是宣称，唯有他是科洛巴索斯之\OriginalTerm{沉默}{Sige}的孕床与容器，因为他既是独生子，便以类似如下方式将那由缺陷的\OriginalTerm{思念}{Enthymesis}托付给他的东西产生出来。他宣称，那无限崇高的\OriginalTerm{四元体}{Tetrad}以女性的形态（因为世界无法承受它以男性形态来临）从不可见、不可言说之处降临到他身上，独向他解释了它自身的本质和万物的起源——这是它先前从未向任何神或人启示过的。此事以如下方式发生：当那无始、不可理解、无形无质、既非男亦非女的圣父愿意生出那对祂而言不可言说者，并使那不可见者具有形象时，祂张开自己的口，发出与祂相似的圣言；圣言站在祂近旁，向祂显示了祂自己是什么，因为祂已以那不可见者的形态显现出来。此外，祂的名字的发音如下：——祂说出名字的第一个字，那是（其余一切的）开始，那话语由四个字母组成。祂加上第二个字，这也由四个字母组成。接着祂说出第三个字，这又包含十个字母。最后祂说出第四个字，由十二个字母组成。如此，整个名字的宣示得以完成，共三十个字母，四个不同的发音。这些要素中的每一个都有自己独特的字母、特征、发音、形式和图像，没有哪一个能感知它是其组成部分的那个（话语）的形状。也没有任何一个认识它自己，也不认识它邻伴的发音，而是每一个都以为凭借自己的发音便确实说出了整个名字。因为尽管每一个都是整体的一部分，它却以为自己的声音就是整个名字，并且不停息地发出声音，直到凭借自己的发音达到了每个要素的最后一个字母为止。这位师傅宣称，当所有这些字母混合为一个字母，发出同一个声音时，万物的复原将发生。他认为这个发音的象征见于我们齐声所说的\Emph{阿们}。\BibleRef{格前14:16} 他又说，那些不同的声音就是赋予那无形无质、无生的\OriginalTerm{永世}{Æon}以形式的；而这些又是主称为\Quote{天使}的那些形式，他们\Quote{常见圣父的面}。\BibleRef{玛18:10}\par

2. 那些可说且常见的元素的名称，他称之为永世、话语、根、种子、圆满和果实。他断言，这些名称中的每一个以及每个名称所特有的一切，都应理解为包含在\OriginalTerm{教会}{Ecclesia}这一名称之内。这些元素中，最后那个元素的最后一个字母发出了声音，这声音下去，生出了它自己的元素，按（其他）元素的形象，他声称，下面的事物由此被安排到它们所占据的秩序中，而先前的事物则被召叫而存在。他还主张，那个字母本身——它的声音随下面的声音而来——又被它所属的音节收回，为的是整体的完成，但声音如同被抛在外面，留在了下面。但那本身有其特殊发音下降到下面的字母的元素，他断言由三十个字母组成，而这些字母中的每一个又在其自身内包含其他字母，藉以表达该字母的名称。如此，其他字母又由其他字母命名，其余以此类推，以至字母的众多性膨胀至无限。你可以通过以下例子更清楚地理解我的意思：——\Emph{Delta}这个词包含五个字母，即D、E、L、T、A；这些字母又由其他字母写出，其他字母又由其他字母写出。那么，如果\Emph{Delta}这个词（如此分析下去）的整个构成趋于无限，字母不断生出其他字母，前后相继，那么这一个词的字母海洋该比它本身浩瀚多少呢！即使一个字母就如此无限，请想想整个名称中字母的浩大；玛尔谷的\OriginalTerm{沉默}{Sige}教导我们，\OriginalTerm{先祖}{Propator}即由这些字母组成。为此之故，圣父知道自己本质的不可理解性，便将每个字母发出自己发音的（能力）赋予那些他称为永世的元素，因为没有一个能凭自己发出整个发音。\par

3. 此外，那\OriginalTerm{四元体}{Tetrad}更充分地向他解释这些事，说：——我希望向你显示\OriginalTerm{真理}{Aletheia}本身；因为我已将她从上界带下来，好使你毫无遮蔽地看见她，理解她的美——也使你听见她发言，钦佩她的智慧。那么，请看她的头部在高处：\Emph{Alpha}和\Emph{Omega}；她的颈项：\Emph{Beta}和\Emph{Psi}；她的肩膀连同双手：\Emph{Gamma}和\Emph{Chi}；她的胸脯：\Emph{Delta}和\Emph{Phi}；她的横膈膜：\Emph{Epsilon}和\Emph{Upsilon}；她的背：\Emph{Zeta}和\Emph{Tau}；她的腹部：\Emph{Eta}和\Emph{Sigma}；她的大腿：\Emph{Theta}和\Emph{Rho}；她的膝盖：\Emph{Iota}和\Emph{Pi}；她的腿：\Emph{Kappa}和\Emph{Omicron}；她的脚踝：\Emph{Lambda}和\Emph{Xi}；她的脚：\Emph{Mu}和\Emph{Nu}。这就是真理的身体，依照这位术士的说法，这就是元素的形象，这就是字母的特征。他称这个元素为\OriginalTerm{人}{Anthropos}，并说这是所有言语的泉源，一切声音的开始，一切不可言说者的表达，以及沉默\OriginalTerm{沉默}{Sige}的口。这确实是真理的身体。但是，你当高扬你心思的意念，从真理的口中倾听那自生的圣言，祂也是圣父恩惠的分施者。\par

4. 当\Quote{四元体}说出这些之后，\Quote{真理}注视着他，开口说出一个词。那词是一个名号，而这名号就是我们所知并所说的那一个，即\Quote{耶稣基督}。她说完这名号后，立即又陷入沉默。马古斯正期待她会说更多，这时\Quote{四元体}再次上前说：——\Quote{你竟把你从\Quote{真理}口中听到的那个词视为可鄙的。你所知并似乎拥有的，并非一个古老的名号。因为你仅拥有它的声音，却对其能力一无所知。因为\Quote{耶稣}（\OriginalTerm{耶稣}{\GreekText{Ἰησοῦς}}）是一个具有数字象征意义的名号，由六个字母组成，且为所有属于蒙召者的人所知。但\Quote{圆满}Eon中的那一位则由许多部分组成，具有另一种形式和形状，且为那些与祂有亲密关系的（天使）所知，他们的形状（权能）常与祂同在。}\par

5. 那么，要知道你所拥有的二十四字母，是那包含全部更高元素数字的三个权能的象征性流溢。你当这样推算——那九个哑音字母是\Quote{父}与\Quote{真理}的肖像，因为它们无声，即其性质不可言说或发音。半元音则代表\Quote{圣言}与\Quote{生命}，因为它们似乎介于辅音与元音之间，兼具两者的性质。元音则代表\Quote{人}与\Quote{教会}，因为从\Quote{人}发出的声音赋予它们存在；声音的声调赋予它们形式。如此，\Quote{圣言}与\Quote{生命}拥有八个（字母）；\Quote{人}与\Quote{教会}拥有七个；\Quote{父}与\Quote{真理}拥有九个。但因各份的数目不等，那位存在于父内的便降下，特由祂所从出的那一位派遣，为纠正所发生的事，好使圆满的统一体因得了平等而在一切中发展出那从一切流出的唯一权能。如此，那仅有七个字母的部分得了八的权能，三组在数目上变得相似，全成了八元组；这三者合起来便构成二十四这个数目。那三个要素（他宣称它们与三个权能结合，从而形成那从其中流出二十四字母的六）因不可言说的\Quote{四元体}之言而四倍化，产生与它们相同的数目；他又主张这些要素属于那不可名状者。这些要素又由三个权能赋予了那看不见者的相似性。他还说那些我们称为双重的字母，是这些要素肖像的肖像；若将它们加于二十四字母，按类比之力，便形成三十之数。\par

6. 他断言，这个安排与类比的果实已显现在一个肖像的相似中，即那位六天后与另外三人上山\BibleRef{玛 17:7}，之后成为六人中的第六位者，祂以此身份降下并包含在七元组内，因祂是光辉的八元组，并在自身内包含了全部要素的数目——这由鸽子（她是阿尔法与奥梅伽）在祂前来受洗时降临所清楚显明，因为鸽子的数目是八百零一。为此，梅瑟宣布人在第六日被造；而根据安排，又在第六日，即预备日，最后一人出现，为再生第一人。这一安排的开始和结束都在第六时辰，即祂被钉于树时。因为那完美的\Quote{灵智}知道六这个数目具有形成与再生的能力，便向光明之子宣告了那由祂以标示数显现的再生。因此他也宣称，双重字母含有标示数之数；因为此标示数加入二十四个要素后，便完成了三十字母的名号。\par

7. 他（马古斯）正如马古斯的\Quote{静默}所宣称的那样，使用了七个字母的权能为器皿，以使（阿卡莫特）独立意志的果实得以彰显。\Quote{看看这当前的标示数，}她说——\Quote{这位按着（原初）标示数所形成者，似乎被分割或切割为两部分，且留在外面；祂凭自己的权能与智慧，借着自身所产生的，给这由七个权能构成的世界赋与生命，如同七元组权能的肖像，并如此形成它，使它成为一切可见之物的灵魂。祂确实如同凭自己自由意志所形成的那样使用此作品；但其余的，作为那不能完全模仿之物的肖像，则从属于母亲的思想。第一重天发出\Quote{阿尔法}之音，第二重天发出\Quote{艾普西隆}，第三重天发出\Quote{伊塔}，第四重天（亦在七重之中）发出\Quote{约塔}之音，第五重天发出\Quote{奥米克隆}，第六重天发出\Quote{宇普西隆}，第七重天（亦是从中间数的第四重）发出优美的\Quote{奥梅伽}，}——正如马古斯的\Quote{静默}大谈废话却不说真理那样自信地断言。\Quote{而这些权能，}她补充道，\Quote{全部同时彼此拥抱，奏出那产生它们者的荣耀；那荣耀之声被传至上方的\Quote{首父}。}她还断言，\Quote{这赞美之声飘扬到地上，便成了地上万物的塑造者与父母。}\par

8. 他以此证明，列举刚出生的婴儿，他们一出母腹时的哭声，与这些元素中每一个的声音都一致。他说，正如七种权力光荣圣言，婴儿的哀怨灵魂也是如此。为此，达味说：\BibleQuote{由赤子乳儿的口中，你取得了完美的赞颂；} 又说：\BibleQuote{高天陈述天主的光荣。} 也因此，当灵魂陷入困难和痛苦时，为了自我安慰，它呼喊\Quote{啊}（\OriginalTerm{欧米伽}{\GreekText{Ω}}），以尊崇这个字母，以便它在上界的同源灵魂能认出[它的痛苦]，并赐下安慰。\par

9. 因此，关于由三十个字母组成的整个名字，以及从这[名字]的字母获得增长的\OriginalTerm{深渊}{Bythus}，还有由十二个肢体（每个肢体由两个字母组成）构成的\OriginalTerm{真理}{Aletheia}的身体，以及她根本没有说出却发出的声音，并且关于那个无法用言语表达的名字的分析，以及世界和人的灵魂——按照它们拥有那种依照[上界事物]形象的结构——他说出了他那些荒谬的见解。接下来我要叙述\OriginalTerm{四元组}{Tetrad}如何从名字中向他显示了一个数目相等的权能；这样，我的朋友，我所听到的他所说的一切，就不会对你有所隐瞒；这样，你多次向我提出的请求就得以实现。\par

\SourceRecord{Translated by Alexander Roberts and William Rambaut.；Ante-Nicene Fathers；Vol. 1.；Edited by Alexander Roberts, James Donaldson, and A. Cleveland Coxe.；Buffalo, NY: Christian Literature Publishing Co.,；1885.；<http://www.newadvent.org/fathers/0103114.htm>.}

% structure: p0001 source=h1 target=section reason=page-title

\section{\WorkTitle{驳斥异端（第一卷，第十五章）}}

1. 全智的\OriginalTerm{西热}{Sige}遂向他宣告那二十四元素的产生如下：——与\OriginalTerm{莫诺特斯}{Monotes}一同并存的，有\OriginalTerm{赫诺特斯}{Henotes}，由之生出两个产物，正如我们以上所论，即\OriginalTerm{莫那斯}{Monas}与\OriginalTerm{亨}{Hen}，此二者加于前二者，合而为四，因为二乘二得四。再者，二与四相加，显出六之数。进而，此六数乘以四，便生出二十四形式。这第一个\OriginalTerm{泰特拉德}{Tetrad}的名号，被认为至为神圣，非言语所能表达，惟圣子知之，圣父亦知其为何。其他按此人之言须以敬意、信德与敬畏诵念的名号，是：\OriginalTerm{阿勒托斯}{Arrhetos}与\OriginalTerm{西热}{Sige}、\OriginalTerm{帕特尔}{Pater}与\OriginalTerm{阿勒忒亚}{Aletheia}。如今这泰特拉德的总数达二十四字母；因为阿勒托斯一名本身包含七字母，西热五字母，帕特尔五字母，阿勒忒亚七字母。若将所有这些相加——两次五，两次七——便完成二十四之数。同样，第二个泰特拉德，即\OriginalTerm{逻各斯}{Logos}与\OriginalTerm{佐厄}{Zoe}、\OriginalTerm{安特罗波斯}{Anthropos}与\OriginalTerm{艾克勒西亚}{Ecclesia}，也显示相同数目的元素。此外，那堪可诵念的救主名号，即\OriginalTerm{耶稣}{\GreekText{᾿Ιησοῦς}}，由六字母构成，但其不可言说之名则包含二十四字母。\Emph{基督圣子}（\OriginalTerm{基督圣子}{\GreekText{υἱὸς} \GreekText{Χρειστός}}）一名包含十二字母，但基督中不可言说之名则有三十字母。为此，他宣称祂是\Emph{阿尔法}与\Emph{敖默加}，以指明鸽子，因那鸟的名号有此数目。\par

2. 但他断言，耶稣有如下不可言说的起源。从万有之母，即第一个泰特拉德，生出第二个泰特拉德，有如女儿；如此形成\OriginalTerm{奥格多阿德}{Ogdoad}，从它又生\OriginalTerm{德卡德}{Decad}：于是产出德卡德与奥格多阿德。德卡德与奥格多阿德结合，并十倍乘之，便生出八十之数；再以八十乘十，则产生八百之数。如此，从奥格多阿德［乘］入德卡德而来的全部字母总数，便是八百八十八。这就是耶稣之名；因为此名，若照字母数值计算，即为八百八十八。因此，你们对其关于超天界耶稣之起源的意见，便有了清楚的陈述。为此，希腊字母包含八个\OriginalTerm{莫纳斯}{Monad}、八个德卡德和八个\OriginalTerm{赫卡塔德斯}{Hecatad}，表示八百八十八之数，即\OriginalTerm{耶稣}{Jesus}，祂由所有数字形成，因此被称为\Emph{阿尔法}与\Emph{敖默加}，指明祂起源自万有。再者，他们如此阐述：若第一个泰特拉德按数字级数相加，则出现十之数。因为一、二、三、四相加，得十；而按他们之意，这便是耶稣。此外，\OriginalTerm{克瑞斯多斯}{Chreistus}，他说，为一八字母词，指示第一个奥格多阿德，此数乘以十，便生出耶稣（888）。基督圣子，他说，也被言及，即\OriginalTerm{杜德卡德}{Duodecad}。因为\OriginalTerm{圣子}{\GreekText{υἰός}}一名包含四字母，基督（克瑞斯多斯）含八字母，两者相结合，指明杜德卡德的伟大。但他声称，在此名的\Emph{\OriginalTerm{埃皮塞蒙}{Episemon}}出现之前，即耶稣圣子，人类深陷于巨大无知与错误中。但当这六字母之名被显明（那具有此名者披上肉身，以便为人感官所理解，且自身含有此六与二十四字母），人类得以认识祂，便摆脱无知，从死亡进入永生，此名作为他们导向真理之父的向导。因为万有之父决意终止无知，摧毁死亡。而废除无知正是对祂的认识。因此，那\OriginalTerm{人}{Anthropos}是按祂的旨意被拣选，按上界相应权能的肖像而被塑造的。\par

3. 至于永世，它们来自泰特拉德，而在那泰特拉德中，有安特罗波斯与艾克勒西亚、逻各斯与佐厄。他宣称，由这些发出的权能，生出了那位在世上显现的耶稣。天使加俾额尔占据逻各斯之位，圣神占据佐厄之位，至高者的权能占据安特罗波斯之位，而童贞女则指明艾克勒西亚之位。这样，藉着特别的处置，由祂借着玛利亚，生出了那个人，万有之父在祂经过母胎时，拣选了他，为要藉着圣言获得对祂自己的认识。当祂来到水［圣洗］时，有那一位曾上升高天、并完成第十二数目者，以鸽子形状降于祂身，其中含有与祂同时产生的那些人的种子，并与祂一同下降、上升。此外，他主张那下降的权能是圣父的种子，本身既含有圣父与圣子，也含有藉二者可知但不可言说的西热之权能，以及所有永世。这便是那曾藉耶稣的口说话、且承认自己是人子并启示了圣父的那一位灵，祂降于耶稣内，与祂成为一体。他说，由特别处置形成的救主确已消灭死亡，但基督启示了圣父。因此，他主张耶稣是由特别处置形成的那个人之名，且他是按那将要降于他身上的［天上］安特罗波斯的肖像与形状而形成的。在祂接受了那永世之后，祂就拥有了安特罗波斯本身、逻各斯本身、帕特尔、阿勒托斯、西热、阿勒忒亚、艾克勒西亚和佐厄。\par

4. 这样的狂言，我们现在完全可以称之为超越\Emph{Iu, Iu, Pheu, Pheu}以及一切悲剧的感叹或哀号。因为当看到真理被\Person{马库斯}变成一个纯粹的形象，而且这个形象还被字母表戳得千疮百孔时，谁不会憎恶这个如此胆大包天的荒谬谎言的可怜制造者呢？希腊人承认他们最初是从\Person{卡德摩斯}那里获得了十六个字母，而且那是相对于起源而言不久之前的事（正如俗语所说的\Quote{昨天与前天}）；后来，随着时间的推移，他们自己又在某个时期发明了送气音，在另一个时期发明了双字母，最后，据说\Person{帕拉墨德斯}在原有字母的基础上加上了长元音。难道在希腊人发生这些事之前，真理就不存在吗？因为按照你\Person{马库斯}的说法，真理的形体比\Person{卡德摩斯}及其前辈更晚——也比那些添加其他字母的人更晚——甚至比你自己还晚！因为唯独你把你所谓的真理塑造成了一个（外在的、可见的）形象。\par

5. 但谁能容忍你那无聊的\OriginalTerm{沉默}{Sige}，它称呼那不可称呼者，阐述那无法言说者的本性，探求那不可探究者，并宣称你所认为无形无体的那一位张开了口，发出了圣言，仿佛祂被包含在有机体之中；而且祂的圣言，虽然相似于祂的创造者，并承载着不可见者的形象，却由三十个元素和四个音节组成？那么，按照你的理论，万有的圣父，根据圣言的相似性，也将由三十个元素和四个音节组成！或者，谁又能容忍你摆弄形式和数字——一会儿三十，一会儿二十四，一会儿又只有六——而把万有的建立者、塑造者和创造者——天主的圣言——禁锢其中；然后，又把他切成四音节和三十元素，把那创立诸天的万有的主贬低为八百八十八这个数字，使他与字母表相似；并把那不可容纳却容纳万有的圣父再分成\OriginalTerm{四元组}{Tetrad}、\OriginalTerm{八元组}{Ogdoad}、\OriginalTerm{十元组}{Decad}和\OriginalTerm{十二元组}{Duodecad}；通过这样的倍增，来阐述圣父那不可言说、不可思议的本性，正如你自己所宣称的那样？而你自己则像一个精通邪恶发明的\Person{代达罗斯}，是一位至高权能的邪恶建筑师，你用大量彼此衍生的字母，为那被你称为无形无体的那一位构建了一种本性和实质。你宣称那不可分割的权能，却又把它分成辅音、元音和半元音；你错误地将那些无声的字母归于万有的圣父和祂的\OriginalTerm{思想}{Ennœa}，从而把所有信赖你的人都推向了亵渎的顶峰和极大的不虔敬。\par

6. 因此，那位神圣的长老和真理的宣讲者，针对你的鲁莽尝试，非常恰当和合宜地以诗句爆发如下：\par

\Quote{\Person{马库斯}，塑造偶像者，预兆的探查者，\newline{}擅长占星术，精于黑暗的魔法，\newline{}总是用这样的伎俩巩固错误的教义，\newline{}为你所迷惑的人提供征兆，\newline{}那是完全与天主隔绝、背弃的权能奇迹，\newline{}这是\Person{撒殚}，你的真父，使你仍然能够完成的，\newline{}通过\Person{阿匝则耳}，那个堕落而仍大能的天使——\newline{}使你成为他自己的不虔敬行为的先驱。}\par

这就是那位圣善长者的言辞。我将努力简明地陈述他们那极为冗长的神秘体系的其余部分，并将长久以来隐藏的事物揭露到光中。因为这样，这些事情将容易被所有人识破。\par

\SourceRecord{Translated by Alexander Roberts and William Rambaut.；Ante-Nicene Fathers；Vol. 1.；Edited by Alexander Roberts, James Donaldson, and A. Cleveland Coxe.；Buffalo, NY: Christian Literature Publishing Co.,；1885.；<http://www.newadvent.org/fathers/0103115.htm>.}

% structure: p0001 source=h1 target=section reason=page-title

\section{驳异端（第一卷，第十六章）}

1. 这些人在自己的\OriginalTerm{永世}{Æons}的生成中，融合了福音中提及的迷途和寻回那只羊\BibleRef{路 15:4}，试图以更神秘的方式阐述事物，同时将一切都归于数字，主张宇宙是由一个\OriginalTerm{单子}{Monad}和一个\OriginalTerm{双元}{Dyad}形成的。然后，从一数到四，他们便生出了\OriginalTerm{十元}{Decad}。因为一、二、三、四相加，产生了十永世的数目。再者，双元从自身（按两倍）推进到六——二、四、六——产生了\OriginalTerm{十二元}{Duodecad}。同样，如果我们以同样的方式数到十，就会出现数目三十，其中包含八、十和十二。因此，他们称十二元为\OriginalTerm{受难}{passion}，因为它包含了\OriginalTerm{附标}{Episemon}，并且附标（可以这么说）依附于它。由于这个原因，因为与第十二个数字相关的错误发生了，那只羊跳跃着逃脱了，并且迷途了；因为他们断言从十二元发生了一次偏离。同样，他们郑重地声称，有一个能力也从十二元离开而消失了；这由那丢失了一个\OriginalTerm{德拉克马}{drachma}并且点上灯重新找到它的女人所象征\BibleRef{路 15:8}。因此，剩余的数字，即就钱币而言的九和就羊而言的十一，相乘产生了九十九，因为九乘以十一是九十九。所以他们也主张\Quote{阿们}这个词包含这个数目。\par

2. 然而，我不会不厌其烦地叙述他们的其他解释，以便您随处看到结果。例如，他们主张字母\Emph{\OriginalTerm{伊塔}{\GreekText{η}}}连同\Emph{\OriginalTerm{附标}{\GreekText{ς}}}构成了一个\OriginalTerm{八元}{Ogdoad}，因为它占据从第一个字母起的第八位。然后，同样，不带附标，计算字母的数目，一直加到伊塔，他们得出\OriginalTerm{三十元}{Triacontad}。因为如果从阿尔法开始，结束于伊塔，略去附标，并依次把各字母的数值相加，他将发现它们的总数共达三十。因为到\OriginalTerm{伊普西龙}{\GreekText{ε}}为止，形成了十五；然后加上七，总和达到二十二。接着，加上伊塔，其值为八，最奇妙的三十元就完成了。因此他们宣称八元是三十永世的母亲。既然数目三十由三个能力（八元、十元和十二元）组成，乘以三，就产生九十，因为三乘三十是九十。同样，这个\OriginalTerm{三元}{Triad}自乘，产生九。这样，八元通过这些方式生出了九十九。既然第十二永世因她的偏离在高层留下了十一，他们便主张字母的位置是他们计算方法的真实坐标（因为\OriginalTerm{兰布达}{\GreekText{λ}}在字母表中是第十一位，代表数目三十），也构成了上面事物布局的象征，因为从阿尔法开始，略去附标，直到兰布达，按照字母的连续数值相加，包括兰布达本身，总和是九十九；但这个兰布达，作为第十一位，下降寻找一个与自己相等的，以便完成十二个字母的数目，当它找到这样一个时，数目就完成了，这从字母的形态本身可以明显看出；因为兰布达可以说在寻找一个与自己相似的，找到后，把它抱紧，从而填充了第十二位的位置，字母\OriginalTerm{缪}{\GreekText{Μ}}由两个\OriginalTerm{兰布达}{\GreekText{ΛΛ}}组成。因此，藉着他们的\Quote{知识}，他们避开了九十九的位置，即偏离——左手的象征——但努力争取多一个，加在九十九上，使他们的计数转向右手。\par

3. 我深知，亲爱的朋友，当你读完这一切，你必会为他们那膨胀的愚昧智慧而开怀大笑！但这些人真的值得哀悼，他们传播这样一种宗教，如此冷酷而乖僻地拆毁那真正不可言喻之能力的伟大，以及天主本身那般引人注目的计划，借助阿尔法和贝塔，并依靠数字的帮助。但是，所有离开教会并听从这类老妇荒言的人，真是自我定罪的；保禄命令我们\BibleQuote{在第一次和第二次劝诫之后，就要避开他们。}\BibleRef{铎 3:10}主耶稣的门徒若望加重了他们的定罪，他要求我们甚至不要向他们致\BibleQuote{平安}的问候；因为他说：\BibleQuote{那向他们问安的人，就分享他们的恶行。}而这有道理，因为主说：\BibleQuote{恶人没有平安。}\BibleRef{依 48:22}这些人确实是不虔敬至极，他们竟声称天地的创造者、唯一全能的天主、除他以外没有别的神，是由一个缺陷产生的，而这个缺陷本身又源自另一个缺陷，因此按照他们，他是第三个缺陷的产物。这样的观点我们应憎恶和咒骂，并应远远逃离持有此观点的人；并且，他们越是猛烈地坚持并喜悦于他们虚构的教义，我们就越该相信他们是被\OriginalTerm{奥格多阿德}{Ogdoad}的邪恶神灵所影响——正如那些陷入狂乱的人，他们笑得越厉害，自认为健康，做一切事仿佛身心（身体和心灵）都健全，甚至有些事做得比真正健康的人还好，但这只是显示他们病得更重。同样，这些人越是在智慧上似乎超越他人，并拉紧弓弦耗尽力气，就越是显出他们是更大的愚人。因为当愚昧的邪灵出去后，发现他们不事奉天主，而忙于世俗的问题时，就\BibleQuote{带了七个比自己更恶的灵来}\BibleRef{玛 12:43}，并用以为自己能构想出超越天主之物的观念膨胀这些人的心思，并且恰当地预备他们接受欺骗，就在他们心中植入愚昧邪恶之灵的奥格多阿德。\par

\SourceRecord{Translated by Alexander Roberts and William Rambaut.；Ante-Nicene Fathers；Vol. 1.；Edited by Alexander Roberts, James Donaldson, and A. Cleveland Coxe.；Buffalo, NY: Christian Literature Publishing Co.,；1885.；<http://www.newadvent.org/fathers/0103116.htm>.}

% structure: p0001 source=h1 target=section reason=page-title

\section{\WorkTitle{驳斥异端}（卷一，第十七章）}

1. 我也愿意向你们解释他们关于受造界本身如何通过母亲由\OriginalTerm{德谬哥}{Demiurge}（似乎是在他不知情的情况下）按照无形之物的形象而形成的理论。他们声称，首先，四种元素——火、水、土、气——是按照上界\OriginalTerm{四重体}{Tetrad}的形象而产生的；然后，加上它们的性能，即热、冷、干、湿，就呈现出一个与\OriginalTerm{八重体}{Ogdoad}完全相似的形像。接着，他们以下列方式推算出十种权能：有七个球状体，他们也称之为天；然后是一个包容这些球状体的球状体，他们称之为第八重天；此外，还有太阳和月亮。他们宣称，这十个实体是无形\OriginalTerm{十重体}{Decad}的预像，这十重体出自\OriginalTerm{圣言}{Logos}和\OriginalTerm{生命}{Zoe}。至于\OriginalTerm{十二重体}{Duodecad}，则由所谓的黄道十二宫所表示；因为他们断言，十二个星宿最明显地预表了\OriginalTerm{十二重体}{Duodecad}，它是\OriginalTerm{人}{Anthropos}和\OriginalTerm{教会}{Ecclesia}的女儿。由于最高天击打在最外层球体上，并因自身重力与整个体系的极快进动相连，作为一种制衡与平衡，使它在三十年内完成从一宫到另一宫的循环——他们说这是\OriginalTerm{何鲁斯}{Horus}的形像，环绕他们那位有三十名号的母亲。再者，月亮在三十天内行经所分配的天区，他们认为这三十天表达了三十个\OriginalTerm{永世}{Æons}的数目。太阳也在十二个月内运行完其轨道，然后回到圆上同一点，因此十二个月显明了\OriginalTerm{十二重体}{Duodecad}；而一天由十二小时计量，也是无形\OriginalTerm{十二重体}{Duodecad}的预像。此外，他们宣称，一小时是日的十二分之一，由三十个部分构成，以表明\OriginalTerm{三十重体}{Triacontad}的形象。同样，黄道十二宫本身的圆周包含三百六十度（每宫含三十度）；因此他们也断言，通过此圆保存了十二与三十之间连接关系的形象。更进一步，他们声称地球分为十二带，每一带根据太阳的垂直照射从天上接受能力，产生与降下影响的能力相应的产物；他们坚持这是\OriginalTerm{十二重体}{Duodecad}及其后裔最明显的预像。\par

2. 此外，他们宣称，\OriginalTerm{德谬哥}{Demiurge}渴望模仿上界\OriginalTerm{八重体}{Ogdoad}的无限、永恒、浩瀚和不受时间度量限制的特性，但作为缺陷的果实，他无法表达其恒久与永恒，于是采取将它的永恒延伸为时间、季节和漫长年岁的权宜之计，以为通过如此大量时间可以模仿其浩瀚。他们进一步宣称，真理逃避了他，他追随了虚假，因此当时间期满时，他的工程必将毁灭。\par

\SourceRecord{Translated by Alexander Roberts and William Rambaut.；Ante-Nicene Fathers；Vol. 1.；Edited by Alexander Roberts, James Donaldson, and A. Cleveland Coxe.；Buffalo, NY: Christian Literature Publishing Co.,；1885.；<http://www.newadvent.org/fathers/0103117.htm>.}

% structure: p0001 source=h1 target=section reason=page-title

\section{驳斥异端（第一卷，第十八章）}

1. 当他们关于创造提出这样的主张时，他们中的每个人都根据自己的能力每天产生一些新东西；因为若不在他们中间编造出一些伟大的虚构，就没有人被认为是\Quote{完美的}。因此，首先有必要指出他们从先知著作中改编（为己所用）了哪些事物，然后加以驳斥。于是，他们声称，\Person{梅瑟}在开始创造记述时，就指出了\OriginalTerm{万有之母}{the mother of all things}，他说：\BibleQuote{在起初天主创造了天地}\BibleRef{创 1:1}。因为他们主张，通过提及这四者——天主、起初、天、地——他展示了他们的\OriginalTerm{四元组}{Tetrad}。\BibleQuote{大地是看不见的、未成形的}\BibleRef{创 1:2}，他们也坚持认为，他通过提及\OriginalTerm{深渊}{abyss}和\OriginalTerm{黑暗}{darkness}（其中也有水，以及运行在水面上的圣神），这样论及第一个四元组所产生的第二个四元组。然后，他继续提及\OriginalTerm{十元组}{Decad}，提到了光、昼、夜、穹苍、晚、晨、旱地、海、植物，第十位是树木。这样，通过这十个名称，他指出了十个\OriginalTerm{永世}{Æons}。再者，\OriginalTerm{十二元组}{Duodecad}的能力被他这样预示：他提到了太阳、月亮、星辰、季节、年份、鲸鱼、鱼类、爬行动物、鸟类、四足动物、野兽，最后，第十二位是人。这样，他们教导说，\OriginalTerm{三十元组}{Triacontad}是借圣神通过\Person{梅瑟}所说的。此外，人也是按照上面那个能力的肖像形成的，他自身就拥有那源自唯一泉源的能力。这能力位于大脑区域，从中产生四种官能，如同上面四元组的肖像，它们被称为：第一种是\Emph{视觉}，第二种是\Emph{听觉}，第三种是\Emph{嗅觉}，第四种是\Emph{味觉}。他们说，\OriginalTerm{八元组}{Ogdoad}由人这样指示：人有两只耳朵、同样数量的眼睛、两个鼻孔，以及双重的味觉，即苦和甜。此外，他们教导说，整个人包含三十元组的完整肖像，如下：在他手上，通过他的手指，他承载着十元组；而在他的整个身体上承载着十二元组，因为他的身体分为十二部分；因为他们那样分配，如同\OriginalTerm{真理之体}{the body of Truth}被他们分割——这一点我们已经谈过。但八元组，由于是不可言说和不可见的，被理解为隐藏在脏腑之内。\par

2. 再者，他们断言，太阳这个大光源是在第四天形成的，参照了四元组的数字。同样，根据他们，\Person{梅瑟}建造的帐幕的庭院，由细麻、蓝色、紫色、朱红色构成，也指向同一肖像。此外，他们主张，司铎垂到脚上的长袍，饰有四行宝石\BibleRef{出 28:17}，指示四元组；并且，如果在圣经中有任何其他可能被勉强纳入数字四的事物，他们就宣称这些事物是为了四元组而存在的。八元组又如下显示：——他们断言，人是在第八天形成的，因为有时他们主张人是在第六天造的，有时在第八天，除非他们或许意指他的属土部分是在第六天形成的，而他的\OriginalTerm{属血肉的部分}{the fleshly part}在第八天，因为他们区分了这两者。他们中的一些人还认为，一个人是按照天主的肖像和模样形成的，是\OriginalTerm{两性兼备的}{masculo-feminine}，就是\OriginalTerm{属神的人}{the spiritual man}；另一个人是由泥土形成的。\par

3. 此外，他们宣称，关于洪水时方舟的安排，通过它八个人得救\BibleRef{创 6:18}，最清楚地指示了带来救恩的八元组。\Person{达味}也显示了同样的，他在兄弟中按年龄排在第八位\BibleRef{撒上 16:10}。再者，发生在第八天的割礼\BibleRef{创 17:12}代表了上面八元组的割礼。总之，无论他们在圣经中找到什么可以关联到数字八的事物，他们都宣称实现了八元组的奥秘。再者，关于十元组，他们主张它由天主应许给\Person{亚巴郎}为产业的十个民族所指示\BibleRef{创 15:19}。\Person{撒辣}在十年之后将自己的婢女\Person{哈加尔}给他\BibleRef{创 16:2}，好借着她得一个儿子，这一安排也显示了同样的事。此外，\Person{亚巴郎}的仆人被派往\Person{黎贝加}那里，在井旁送给她十只金手镯，她的兄弟们留他住了十天；\Person{雅洛贝罕}也接受了十个权杖（支派）\BibleRef{列上 11:31}，以及帐幕的十个帷幔\BibleRef{出 26:1}和十肘高的柱子\BibleRef{出 36:21}，\Person{雅各伯}最初被派往埃及买粮的十个儿子\BibleRef{创 42:3}，以及主复活后向他显现的十位宗徒——\Person{多默}不在其中\BibleRef{若 20:24}——根据他们，代表了\OriginalTerm{不可见的十元组}{the invisible Decad}。\par

4. 关于\OriginalTerm{十二集}{Duodecad}——与之相关的是缺陷所受之苦的奥迹，他们认为一切有形之物皆由此苦而形成——他们断言此数在圣经中随处可见且明显彰显。他们宣称，雅各伯的十二个儿子\BibleRef{创 35:22}，由此产生十二支派；大司祭的胸牌上有十二块宝石和十二个小铃铛；梅瑟放在山脚的十二块石头\BibleRef{出 24:4}；若苏厄放在河中的同样数目的石头\BibleRef{苏 4:3}，以及河对岸抬约柜的人\BibleRef{苏 3:12}；厄里亚在献牛犊为全燔祭时所立的石头\BibleRef{列上 18:31}；宗徒的数目；总之，凡包含\Emph{十二}这个数字的一切事件——都表征他们的十二集。而所有这些的联合，即所谓的\OriginalTerm{三十集}{Triacontad}，他们竭力借诺厄方舟来证明，其高度为三十肘\BibleRef{创 6:15}；撒慕尔的事迹中，他让撒乌耳坐于三十位宾客的首位\BibleRef{撒上 9:22}；达味在田野中藏匿三十天\BibleRef{撒上 20:5}；以及那些与他一同进入山洞的人；还有圣帐幕的长度（高度）为三十肘；此外，若遇到其他类似的数字，他们都将其应用于他们的三十集。\par

\SourceRecord{Translated by Alexander Roberts and William Rambaut.；Ante-Nicene Fathers；Vol. 1.；Edited by Alexander Roberts, James Donaldson, and A. Cleveland Coxe.；Buffalo, NY: Christian Literature Publishing Co.,；1885.；<http://www.newadvent.org/fathers/0103118.htm>.}

% structure: p0001 source=h1 target=section reason=page-title

\section{反对异端（卷一，第十九章）}

1. 我认为还需要补充以下内容：他们如何通过篡改圣经段落，试图说服我们相信他们的\OriginalTerm{元父}{Propator}——这位元父在基督来临之前是无人知晓的。他们的目的是要表明，我们的主宣告了另一位父，而非这个宇宙的造物主；他们不虔地宣称，这造物主（正如我们先前所说）是缺陷的产物。例如，当\Person{依撒意亚}先知说：\BibleQuote{以色列却不认识我，我的百姓也不明白我}\BibleRef{依 1:3}，他们歪曲他的话，将其解释为对无形\OriginalTerm{深渊}{Bythus}的无知。\Person{欧瑟亚}所说的话：\BibleQuote{在他们内没有真理，也不认识天主}\BibleRef{欧 4:1}，他们也竭力作同样的解释。至于\BibleQuote{没有明白的，没有寻求天主的；他们都偏离了正路，一同变为无用}，他们坚持认为这是指对\OriginalTerm{深渊}{Bythus}的无知。此外，\Person{梅瑟}所说的话：\BibleQuote{没有人看见了我还能活着}\BibleRef{出 33:20}，他们也试图让我们相信是同样的所指。\par

2. 因为他们错误地认为，造物主曾被先知们看见。但\BibleQuote{没有人看见了我还能活着}这段话，他们解释为是对祂那不可见、不为任何人所知之伟大的论述；事实上，我们都清楚，\BibleQuote{没有人看见天主}这句话是指那不可见的圣父、宇宙的造物主说的；然而，这些话并非指他们所虚构的那个\OriginalTerm{深渊}{Bythus}，而是指造物主（祂就是那不可见的天主），这一点我们将在下文证明。他们还声称，\Person{达尼尔}也表达了同样的事，因为他曾请求天使解释比喻，表明他自己也不明白。但天使向他隐瞒了\OriginalTerm{深渊}{Bythus}的极大奥秘，对他说：\Quote{\Person{达尼尔}，你赶快去吧，因为这些话是封闭的，直到那些有智慧的人能明白，那些\Emph{洁白}的人变成\Emph{洁白}。}此外，他们夸耀自己就是那些\Emph{洁白}和\Emph{有智慧}的人。\par

\SourceRecord{Translated by Alexander Roberts and William Rambaut.；Ante-Nicene Fathers；Vol. 1.；Edited by Alexander Roberts, James Donaldson, and A. Cleveland Coxe.；Buffalo, NY: Christian Literature Publishing Co.,；1885.；<http://www.newadvent.org/fathers/0103119.htm>.}

% structure: p0001 source=h1 target=section reason=page-title

\section{驳斥异端（第一卷，第二十章）}

除了上述的[歪曲]之外，他们还引用了无数自己编造的伪经和伪造著作，以迷惑愚昧之人和不谙真理圣经者。他们尤其提出那个虚假而邪恶的故事：我们的主在孩童时期学习字母时，老师照例对他说：\Quote{念 \InnerQuote{阿尔法}，} 主便回答：\Quote{阿尔法。} 但当老师再次命他说 \Quote{贝塔} 时，主却答道：\Quote{先告诉我什么是 \InnerQuote{阿尔法}，然后我就告诉你什么是 \InnerQuote{贝塔}。} 他们将此解释为：唯独祂知道那不可知的（神），并以 \Quote{阿尔法} 为象征将其启示出来。\par

同样，福音中的一些段落也被他们作类似渲染，例如祂十二岁时对母亲所说的话：\BibleQuote{你们不知道我必须在我父的事上吗？} \BibleRef{路 2:49} 他们说，祂以此向他们宣告了他们所不知的父。为此，祂曾派遣门徒到十二支派去，向他们宣讲那未知之神。对那称祂为\BibleQuote{善师}的人，\BibleRef{谷 10:17} 祂承认那真正为善的天主，说：\BibleQuote{你为什么称我为善？除了天主一位外，没有善的，就是那在诸天之上的父；} \BibleRef{路 18:18} 他们断言此处的\OriginalTerm{诸天}{heavens}乃指\OriginalTerm{永世}{Æons}。此外，祂没有回答那些对祂说\BibleQuote{你凭什么权柄做这些事}的人，\BibleRef{玛 21:23} 反而用反问使他们完全混乱；按他们的解释，祂以此显示父的不可言说性。再者，当祂说\Quote{我多次愿意听这些话中的一句，却没有人能说出来}时，他们认为祂以\Quote{一句}指明了他们所不知的独一真神。还有，当祂临近耶路撒冷，为它哭泣说\Quote{你若知道，至少在这日子，你明白了关乎你平安的事，但这事如今在你面前是隐藏的}时，祂以\Quote{隐藏}一词示明了\OriginalTerm{深渊}{Bythus}的隐晦本质。另外，当祂说\BibleQuote{凡劳苦和负重担的，你们都到我这里来，我要使你们安息。你们跟我学吧，} \BibleRef{玛 11:28} 祂便宣告了真理之父。因为他们所不知的，这些人声称祂应许要教导他们。\par

他们更提出以下段落作为最高见证，甚至可说是其体系的冠冕：\BibleQuote{父啊，天地的主宰，我称谢祢，因为祢将这些事瞒住了智慧及明达的人，而启示给了小孩子。是的，父啊，因为祢这样喜欢。我父将一切交给了我；除了父，没有人认识子；除了子和子所愿意启示的人，没有人认识父。} \BibleRef{玛 11:25} 他们声称，这些话清楚表明，他们所炮制的真理之父，在祂来临之前无人知晓。他们试图歪曲这段经文，似乎认为创造及形成[世界]者一向为众人所知，而主所说的话是指那位无人知晓的父，即他们如今所宣讲的。\par

\SourceRecord{Translated by Alexander Roberts and William Rambaut.；Ante-Nicene Fathers；Vol. 1.；Edited by Alexander Roberts, James Donaldson, and A. Cleveland Coxe.；Buffalo, NY: Christian Literature Publishing Co.,；1885.；<http://www.newadvent.org/fathers/0103120.htm>.}

% structure: p0001 source=h1 target=section reason=page-title

\section{\WorkTitle{驳斥异端（第一卷，第二十一章）}}

1. 他们的关于\OriginalTerm{救赎}{redemption}的传统是无形且不可理解的，因为它是不可理解且无形之事的母亲；因此，由于它是变化的，不可能简单且一下子揭示其性质，因为其中每一个人都凭自己的倾向传下来。这样，有多少位这些神秘意见的教师，就有多少种\Quote{救赎}方案。当我们驳斥他们时，我们将在适当的地方指出，这类人已被撒殚煽动，否认那使人重生归于天主的圣洗，从而弃绝了整个（基督徒）信仰。\par

2. 他们主张，那些达到圆满知识的人必须重生入那高于万有的能力。否则不可能进入\OriginalTerm{圆满}{Pleroma}，因为这种（重生）带领他们进入\OriginalTerm{深渊}{Bythus}的深处。因为可见的耶稣所建立的圣洗是为赦免罪过，但降临在他身上的基督带来的救赎是为成全；他们声称前者是属生灵的，后者是属神的。若翰的洗礼是为悔改而宣告的，但耶稣的救赎是为成全而带来的。对此他引用说：\Quote{我还有另一种洗礼要受，而且我急切地渴望它。}此外，他们肯定主将这种救赎加给了载伯德的儿子们，当他们的母亲请求他们在他的王国中坐在他左右时，他说：\BibleQuote{你们能受我将受的洗礼吗？}\BibleRef{谷 10:38}他们还宣称，保禄常以明确的言辞论述了那在基督耶稣内的救赎；而这正是他们以如此多样且不一致的形式所传下来的。\par

3. 因为他们中的一些人准备了一张婚床，并对那些受入门礼者举行一种神秘的仪式（念诵某些词句），宣称这是按照上面的结合所举行的属神婚姻。另一些人带领他们到有水的地方，用这些话为他们付洗：\Quote{进入宇宙不可知之父的名内——进入真理，万有之母——进入那降临在耶稣身上的那位——进入结合、救赎以及与诸能的共融。}还有一些人重复某些希伯来语词汇，以便更彻底地迷惑那些受入门礼者，如下：\Quote{巴塞玛，卡莫斯，巴奥诺拉，米斯塔迪亚，鲁阿达，库斯塔，巴巴弗尔，卡拉赫泰。}这些词的解释如下：\Quote{我呼求那高于父的一切能力的，称为光、善神和生命的那一位，因为你曾统御在身体内。}另一些人如此陈明救赎：那隐藏于任何神祇、统治和真理之外的名字，即纳匝肋耶稣在基督之光的生命中所穿戴的——基督，他因圣神而生活，为天使的救赎。重建的名字如下：梅西阿，乌法雷格，纳门普索曼，卡尔多奥，莫索梅迪亚，阿克弗拉诺，普萨瓦，耶稣纳扎里亚。这些词的解释如下：\Quote{我不分割基督的神，也不分割心，也不分割那仁慈的超天之力；愿我享受你的名，啊，真理的救主！}施礼者的话如此；但受入门者答复说：\Quote{我被建立，我被救赎；我赎回我的灵魂脱离这个世代（世界），以及与之相关的一切，因伊奥的名，他赎回了自己的灵魂进入在基督内的救赎，基督生活着。}然后旁观者加上这些话：\Quote{愿平安归于所有这名字所降临的人。}此后他们用香膏傅抹受入门者；因为他们声称这油膏是那超乎万有的芬芳之气的预象。\par

4. 但其中有些人认为带人到水边是多余的，而是将油和水混合，把这混合物放在受入门者的头上，使用我们已经提到的某些类似词句。他们主张这就是救赎。他们也惯常用香膏傅抹。其他一些人则拒绝所有这些做法，主张那不可言说且无形之能力的奥迹，不应由可见且可朽的受造物施行；那不可设想、无形体且超越感觉的（存有）的奥迹，也不应由可感觉且具形体的事物施行。他们认为，对那不可言说的伟大的知识本身就是圆满的救赎。因为既然缺陷和情欲都是从无知而来，那么由无知所形成的一切本质便因知识而消灭；因此知识是内在的人的救赎。然而这并非属身体的本质，因为身体是可朽的；也非属生灵的，因为生灵之魂是缺陷的果实，并且仿佛是灵的居所。因此救赎必须是属神的本质；因为他们肯定内在的、属神的人藉知识得以救赎，并且他们获得了对万有的知识后，就不再需要别的了。这便是真正的救赎。\par

5. 另有一些人，甚至直到死亡的时刻，仍继续施行救赎，他们将油和水，或前面提到的油膏与水，放在人头上，同时念诵上述的呼求，使这些人变得不能被掌权者与权势者所抓住或看见，并使他们的内在之人以一种不可见的方式升到高处，仿佛他们的身体被留在这世界受造物中，而他们的灵魂则被送到德缪格那里。他们又教导这些人，当他们抵达掌权者与权势者时，要使用这些话：\Quote{我是父亲的一个儿子——那先在的父亲，以及他里面一个先在的儿子。我来是为了观看一切，既属于我自己，也属于别人的事物，虽然严格来说，它们并不属于别人，而是属于那本性为女的阿卡莫特，她为自己造了这些东西。因为我从那先在者获得存在，我又回到我原来的地方，即我从此之处出来的地方。} 他们断言，说了这些话后，他便逃脱了那些权势。然后他进到德缪格的同伴那里，如此对他们说：\Quote{我是一个比那形成你们的女性更宝贵的器皿。你们的母亲若不知道自己的根源，我却认识自己，也知道我从何处来，并且我呼求那不朽的索菲亚，她在父亲内，是你们母亲的母亲，她没有父亲，也没有任何男性配偶；但一个出自女性的女性形成了你们，她对自己的母亲一无所知，以为唯独她自己存在；但我呼求她的母亲。} 他们宣称，当德缪格的同伴听到这些话时，大为骚动，并谩骂他们的起源和母亲的种族。但他则进入自己的地方，已经丢掉了他的锁链，即他的兽性本性。以上便是我们关于\Quote{救赎}所了解到的细节。但由于他们无论是在教义还是传统上都彼此分歧，并且那些被认为是最现代的人每天努力发明一些新意见，提出前人从未想到过的东西，所以描述他们全部的意见是一件困难的事。\par

\SourceRecord{Translated by Alexander Roberts and William Rambaut.；Ante-Nicene Fathers；Vol. 1.；Edited by Alexander Roberts, James Donaldson, and A. Cleveland Coxe.；Buffalo, NY: Christian Literature Publishing Co.,；1885.；<http://www.newadvent.org/fathers/0103121.htm>.}

% structure: p0001 source=h1 target=section reason=page-title

\section{\WorkTitle{駁斥異端（卷一，第二十二章）}}

我們所持守的真理準則是：有一位全能的天主，祂藉著祂的聖言創造了萬物，並從無中塑造並形成了所有存在的事物。聖經如此說：\BibleQuote{諸天因上主的話而創造，萬象藉祂口中的氣而形成。}又說：\BibleQuote{萬物是藉著祂而造成的；凡受造的，沒有一樣不是由祂而造成的。}\BibleRef{若 1:3}沒有例外或削減；但聖父藉著祂創造了萬物，無論是可見的還是不可見的，感官的還是理智的，因賦予它們某種特性而受時間限制的，或是永恆的；這些永恆的事物，祂不是藉著天使，也不是藉著任何從祂的\OriginalTerm{意念}{Ennœa}分離的權能創造的。因為天主不需要這一切，祂是那藉著祂的聖言和聖神創造、安排、治理萬物，並命令一切事物存在的那一位——祂塑造了世界（因為世界是萬有的）——祂塑造了人——祂是亞巴郎的天主、依撒格的天主、雅各伯的天主，在祂之上沒有別的神，沒有元始原則，沒有權能，沒有\OriginalTerm{圓滿}{Pleroma}——祂是我們的主耶穌基督之父，我們將要證明這一點。因此，持守這個準則，我們將輕易地顯示，儘管他們的意見繁多且多樣，這些人都已偏離了真理；因為幾乎所有不同派別的異端者都承認有一位天主；但然後，藉著他們毒害的教義，他們將這真理變為錯誤，正如外邦人藉著偶像崇拜所做的——從而證明自己對創造他們的天主忘恩負義。此外，他們輕視天主的創造，說話反對自己的救恩，成為自己最嚴厲的控告者，並成為（反對自己的）假證人。儘管他們不情願，這些人終有一天要在肉身內復活，承認那使他們從死者中復活者的能力；但因著他們的不信，他們不會被列入義人之中。\par

因此，既然識別並定罪所有異端者是一項複雜且多樣的任務，而我們的意圖是按照他們各自的特點來回應他們，我們認為有必要首先說明他們的源頭和根源，以便藉著了解他們那至高的\OriginalTerm{深淵}{Bythus}，你就可以明白結出如此果實的樹木的性質。\par

\SourceRecord{Translated by Alexander Roberts and William Rambaut.；Ante-Nicene Fathers；Vol. 1.；Edited by Alexander Roberts, James Donaldson, and A. Cleveland Coxe.；Buffalo, NY: Christian Literature Publishing Co.,；1885.；<http://www.newadvent.org/fathers/0103122.htm>.}

% structure: p0001 source=h1 target=section reason=page-title

\section{驳斥异端（第一卷，第二十三章）}

1. 撒玛黎雅人西满乃是那术士，路加——宗徒的门徒与追随者——记载说：\Quote{有一个人名叫西满，向来在那城里行邪术，迷惑撒玛黎雅的百姓，自称是个大人物，众人从最小到最大的都听从他，说：\BibleQuote{这人就是那称为大能的天主的能力。}他们都依从他，因为他长时间用邪术使他们疯狂。}\BibleRef{宗 8:9} 这西满假装相信，以为宗徒们行医治是靠邪术而非天主的德能；又见宗徒们藉覆手使那些因他们所传的基督耶稣而信从天主的人充满圣神，他便怀疑这也是出于更高明的邪术知识，于是向宗徒们献银，以为自己也能获得随意赐予圣神的能力。伯多禄对他说：\Quote{愿你的银子与你一同丧亡，因为你以为天主的恩赐可以用钱买得。你在这事上无份无关，因为你在天主面前心怀不正。我看出你正在苦胆之中，在不义的束缚里。}他因此更不信从天主，反而热心与宗徒们对抗，好使自己看似非凡人物，并更勤奋地钻研全部邪术，以便更有效地迷惑和制伏众人。他在克劳狄乌斯皇帝执政时如此行事，据说因其邪术而被立像尊崇。此人被许多人奉若神明，他教导说：在犹太人中间出现的是子，在撒玛黎雅下降的是父，而来到其他民族时则是圣神的形态。总之，他自称是至高权能，即万有之父，并任由人以任何称呼来称呼他。\par

2. 这撒玛黎雅人西满乃是一切异端的根源，他以下列材料建立其派别：他从腓尼基的提洛城赎出一名名为赫肋纳的妇女，并常带她同行，宣称这女子是他思想的首个概念、万有之母，他在起初心中构想创造天使与总领天使。这\OriginalTerm{恩诺亚}{Ennœa}从他发出，领悟了父的意愿，降于下界，产生了天使与权能，世界便是由这些天使与权能所造。但她生下它们后，因嫉妒而被它们扣留，因为它们不愿被视为由任何其他存有所出。它们对西满一无所知，但他的恩诺亚却被她所生的权能与天使扣留。她遭受它们种种侮辱，以致不能返回父处，甚至被囚于人的身体中，历经世代，从一个女性的身体转移到另一个，如同从器皿到器皿。例如，她曾在引发特洛伊战争的赫肋纳身上；为此斯特西科罗斯因在诗中咒骂她而失明，但后来悔改，写了所谓\Emph{翻案诗}赞美她，才恢复视力。她就这样从一个身体到另一个身体，在每个身体中受辱，最终沦为娼妓，而这正是那迷失的羊所指的。\BibleRef{玛 18:12}\par

3. 为此，西满降临，为要首先赢得她，使她脱离奴役，同时藉着向人彰显自己而赐予人救恩。因为天使们统治世界不善，各自贪图主权，所以他前来改善局面，并下降，变形并同化于权能、率领者与天使，使自己在人间看似是人，其实并不是人；如此，人们以为他在犹太地区受苦，其实他并未受苦。此外，先知们是在造世界之天使的默感下发出预言，因此那些信从他与赫肋纳的人不再看重先知，而是自由地随意生活；因为人是藉他的恩宠得救，并非因自己的义行。因为此类行为在本性上并非义行，只是偶然如此，正如造世界的那些天使认为应当如此设立，藉着此类诫命使人受捆绑。为此，他承诺世界将被解散，属于他的人将从造世界者的管辖下得释放。\par

4. 于是，此派别的神秘司铎们过着放荡的生活，并各自尽其所能行邪术。他们使用驱魔咒与符咒。爱情灵药、护符，以及所谓\OriginalTerm{伴灵}{Paredri}和\OriginalTerm{遣梦者}{Oniropompi}，以及其他一切可用的稀奇法术，都被热切应用。他们还铸有西满的像，仿照朱庇特，另有赫肋纳的像，形似弥涅耳瓦，并加以敬拜。总之，他们从西满——这些最不敬教义的创始人——得名，称为西满派；从他们那里，\BibleQuote{伪称的知识}\BibleRef{弟前 6:20} 开始了，即便从他们自己的宣称也可得知。\par

5. 此人的继任者是\Person{默南德尔}，他也是撒玛黎雅人出身，并且同样是魔术实践的高手。他断言\OriginalTerm{首能}{primary Power}始终不为众人所知，而他自己就是从不可见的存在者当中被派遣来作救主，以救赎人类。世界是由天使造的，如\Person{西蒙}一样，他主张这些天使是由\OriginalTerm{厄诺娅}{Ennœa}产生的。他还声称，藉由他所教导的那魔术，他给予这样的知识：人可以胜过那些创造世界的天使；因为他的门徒藉着受洗归于他而得著\Emph{复活}，不再能死亡，反而永葆不朽的青春。\par

\SourceRecord{Translated by Alexander Roberts and William Rambaut.；Ante-Nicene Fathers；Vol. 1.；Edited by Alexander Roberts, James Donaldson, and A. Cleveland Coxe.；Buffalo, NY: Christian Literature Publishing Co.,；1885.；<http://www.newadvent.org/fathers/0103123.htm>.}

% structure: p0001 source=h1 target=section reason=page-title

\section{驳斥异端（卷一，第二十四章）}

1. 在这些人物中，出身于靠近达夫内的安提约基雅的撒督尼诺，以及巴西里德，乘机而起，分别推行了不同的教义体系——前者在叙利亚，后者在亚历山大里亚。撒督尼诺如同默南德尔一样，主张有一位众人皆不知的父，他创造了天使、总领天使、掌权者与有能者。世界及其中万物是由一组七位天使所造。人也是天使的工艺品，是一个从至高权能面前迸发而出的光辉形像，而后下界。他说，当天使们无法抓住这个形像（因为它立刻又向上冲去）时，他们便互相鼓励说：\Quote{让我们按我们的形像和样式造人。}\BibleRef{创 1:26} 于是人被造出，却因天使无力将那种能力传递给他而无法直立，只能像虫子一样在地上蠕动。至高权能怜悯他，因为他是按自己的样式造的，便赐下一线生命火花，使人类得以直立，接合关节，并使之存活。因此他宣称，人死之后，这生命火花便回归同类的本质，而身体的其余部分则分解为原始元素。\par

2. 他还断言救主是无生、无体、无形的，只是假想为可见之人。他主张犹太人的天主是众天使之一。因此，由于所有权势都想要消灭他的父，基督便来毁灭犹太人的天主，而拯救那些信从他的人——即那些拥有他生命火花的人。这位异端者是第一个声称天使造了两种人——一种邪恶，另一种善良的。由于魔鬼协助最邪恶之人，救主便来毁灭恶人与魔鬼，而拯救善人。他们还宣称婚姻与生育出自撒殚。\BibleRef{弟前 4:3} 他的许多追随者也戒除肉食，并以此种假装节制的行为吸引大众。此外，他们认为某些预言是由创造世界的天使们发出的，有些则出自撒殚；撒督尼诺视撒殚本身为天使，与世界的创造者敌对，尤其与犹太人的天主敌对。\par

3. 另一方面，巴西里德为了显得发现了某种更高深且更似有理的学说，对其教义作了极大的发挥。他主张从无生的父首先生出了\OriginalTerm{努斯}{Nous}，从他再度生出了\OriginalTerm{逻各斯}{Logos}，从逻各斯生出了\OriginalTerm{智慧之念}{Phronesis}，从智慧之念生出了\OriginalTerm{智慧}{Sophia}与\OriginalTerm{能力}{Dynamis}，从能力与智慧生出了众权能、众首领以及众天使——他称之为\Emph{第一个}；并由他们造了第一层天。然后其他权能由这些权能流溢而出，造了另一层与第一层相似的天；同样，当他们又由这些权能流溢而出，恰好与上界对应时，他们又造了第三层天；接着从这第三层向下，有第四代后裔；如此类推，依同样方式，他们宣称形成了越来越多的首领和天使，以及三百六十五层天。因此，一年中的天数与这些天的数目相同。\par

4. 那些占据最低一层天（即我们可见的那一层）的天使，造了世界中的一切事物，并彼此分配了大地及地上的万国。他们的首领被认为是犹太人的天主。由于他想要使其他民族臣服于自己的百姓（即犹太人），所有其他首领都抵抗并反对他。因此所有其他民族都与他的民族为敌。然而那位无生无名的父，察觉他们将被毁灭，便派遣自己的首生\OriginalTerm{努斯}{Nous}（他即被称为基督），来解救那些信从他的人，脱离世界创造者的权势。于是他作为人，向这些权能的诸民族显现，并行了奇迹。因此他本人并未受死，而是一个名叫西满的基勒乃人被迫代他背了十字架；如此，后者被他变形，令人以为是耶稣，便在无知与错误中被钉死，而耶稣自己取了西满的形象，站在一旁嘲笑他们。因为他乃是无形体的权能，即无生之父的努斯（理智），能随意变形，并如此升到了派遣他者那里，嘲笑他们，因为他不能被捉住，且对所有人不可见。因此，知道这些事的人已脱离了创造世界的首领；所以我们无需宣认那位被钉者，而应宣认那位以人形显现、被认为被钉、被称为耶稣、由父所派遣的那一位，祂藉此安排毁灭世界创造者的工程。因此他宣称，若有人宣认被钉者，那人仍是奴隶，且处于造我们身体之权能的统治之下；而那否认他的人，则已脱离这些存在，并得知无生之父的安排。\par

5. 救恩只属于灵魂，因为身体在本性上是可朽坏的。他也宣称，预言来自那些创造世界的权能者，但法律是由他们的首领——他曾带领百姓离开埃及地——特别颁布的。他对祭偶像的肉毫不介意，认为无关紧要，毫无犹豫地食用；他也认为使用其他东西以及实践各种情欲是完全无关紧要的事。此外，这些人行魔术，使用偶像、咒语、呼求以及各种奇术。他们还编造一些名字，仿佛它们是天使的名字，宣称其中一些属于第一层天，另一些属于第二层天；然后他们试图陈述那三百六十五个想象出来的层天的名字、原则、天使和权能。他们还断言救主上升和下降所用的那个蛮族名字是\OriginalTerm{考劳靠}{Caulacau}。\par

6. 那么，学习了这些、并认识所有天使及其根源的人，对天使和一切权能而言，就成了看不见和不可理解的，就如\OriginalTerm{考劳靠}{Caulacau}一样。又如子不为任何人所知，同样他们也不该为任何人所知；但他们认识一切、通过一切，自身却对一切保持看不见和未知。他们说：\Quote{你们要认识一切，却不要让人认识你们。}因此，有这种信念的人也随时准备放弃自己的主张；甚至，他们不可能仅仅为一个名称而受苦，因为他们与所有人相似。然而，大众不能理解这些事，只有千分之一或万分之一的人才能理解。他们声称自己不再是犹太人，也还不是基督徒；并且认为完全不适宜公开谈论他们的奥秘，而应保持沉默加以保密。\par

7. 他们以与数学家相同的方式，确定那三百六十五层天的位置。因为他们接受了数学家的定理，并将其转移到自己的教义类型中。他们主张自己的首领是\Emph{\OriginalTerm{阿布拉克萨}{Abraxas}}；因此，这个词本身就包含了三百六十五这个数目。\par

\SourceRecord{Translated by Alexander Roberts and William Rambaut.；Ante-Nicene Fathers；Vol. 1.；Edited by Alexander Roberts, James Donaldson, and A. Cleveland Coxe.；Buffalo, NY: Christian Literature Publishing Co.,；1885.；<http://www.newadvent.org/fathers/0103124.htm>.}

% structure: p0001 source=h1 target=section reason=page-title

\section{\WorkTitle{驳斥异端}（第一卷，第二十五章）}

1. \Person{卡尔波克拉特}及其追随者又主张，世界及其中万物是由远逊于\OriginalTerm{无始圣父}{unbegotten Father}的天使所造。他们还认为\Person{耶稣}是\Person{若瑟}之子，与其他凡人无异，唯有一点不同：因祂的灵魂坚定纯洁，完美地纪念了在\OriginalTerm{无始天主}{unbegotten God}领域内所见之事。为此，从圣父那里有權能降于祂，使祂能逃脱世界的创造者；他们说这權能穿越所有境界，完全自由地存留后，又升回到圣父及那些同样拥抱类似之物的權能那里。他们更宣称，\Person{耶稣}的灵魂虽受犹太人习俗熏陶，却蔑视这些习俗，因此祂被赋予能力，摧毁了作为惩罚居于人内的那些情欲。\par

2. 因此，那类似基督的灵魂能藐视作为世界创造者的掌权者，并同样获得成就相同效果的能力。这种思想使他们骄傲至极，以致有人宣称自己与\Person{耶稣}相似；而另一些更强大者则自称优于祂的门徒，如\Person{伯多禄}、\Person{保禄}及其他宗徒，他们认为这些宗徒毫不逊于\Person{耶稣}。因为他们的灵魂既来自同一领域，也同样藐视世界的创造者，便被视为配得同等的權能，并再回到同一地方。若有人比祂更轻视今世之物，便显得胜于祂。\par

3. 他们也行巫术与咒语，配制春药与爱情魔药；又求问邪灵、招梦之魔鬼及其他可憎之事，宣称他们现在就能统治这世界的君侯与创造者，不单他们，连其中万物皆可统治。这些人如外邦人一般，被\Person{撒殚}差遣来玷污教会，以致人们听到他们所说的话，以为我们全都是这样的人，便掩耳不听真理的宣讲；或见到他们所行之事，便诽谤我们全体——其实我们在教义、道德、日常行为上与他们毫无关系。但他们过着放荡的生活，为掩盖不敬的教义，妄用基督之名来隐藏邪恶，以致\Quote{他们受审判是公义的}\BibleRef{罗 3:8}，因为他们要从天主那里按自己的行为得相应的报应。\par

4. 他们的狂妄如此无度，竟宣称自己有权行一切不敬与邪恶之事，并可随意为之；因为他们主张事物之善恶仅在于人的意见。因此他们认为，灵魂必须借由从一身体转世到另一身体，经历每一种生活方式和每一种行为——除非一次降生就能防止其他需要，一次性完全地行出那些我们既不敢说也不敢听、甚至不能思想、若在同胞中提及也不可相信的事——这样，如他们著作所言，灵魂体验了各种生活，在离去时不缺任何经验。必须坚持这一点，免得为得救还缺一事，被迫再次投胎。他们宣称\Person{耶稣}为此说了以下比喻：\BibleQuote{当你和你的对头还在路上时，要尽力与他了结，免得他把你交给判官，判官交给差役，把你投在狱里。我实在告诉你：非等你还清最后一分钱，你绝不能从那里出来。}\BibleRef{玛 5:25} 他们也说\Quote{对头}是世上天使之一，他们称之为魔鬼，认为他受造就是为了把那些从世界灭亡的灵魂带到至高统治者那里。他们描述他是在世界创造者中为首的，并说他将那些灵魂交给另一位服侍他的天使，将他关进其他身体；因为他们宣称身体就是\Quote{监狱}。他们又解释\BibleQuote{非等你还清最后一分钱，你绝不能从那里出来}这句话，意思是无人能逃脱那些造世界天使的权势，必须从一身体转到另一身体，直到经历这世上可行的一切行为；当什么都不缺时，释放的灵魂才升到那超越造物天使的天主那里。如此，所有灵魂都得救——无论是他们自己的灵魂，若在一生中警惕不延迟而参与各种行为；还是那些借由转世，完成所入每种生活方式所需之事而获得释放的，最终不再被关在身体内。\par

5. 因此，若在他们中间有不敬、非法、禁行之事，我不能再相信他们是那样的人。在他们的著作中我们读到如下的解释，他们宣称\Person{耶稣}是在奥秘中私下对祂的门徒和宗徒说的，而门徒和宗徒请求并获准将所教之事传给其他配得且相信的人。我们确实是因信德和爱德而得救；但其他一切事物，其本性是无关紧要的，仅凭人的意见——有些被认作善，有些被认作恶，实际上没有事物在本性上是恶的。\par

6. 他们中的另一些人在外表上使用记号，在他们门徒的右耳垂内烙印。从这些人中也兴起了\Person{玛策琳纳}，她在\Person{阿尼塞特}主教任内来到\Place{罗马}，并持守这些学说，迷惑了大批群众。他们自称\OriginalTerm{诺斯替派}{Gnostics}。他们还拥有图像，有些是绘制的，有些则由不同材料制成；并且他们断言，在\Person{耶稣}生活在他们中间时，\Person{比拉多}曾制作了基督的肖像。他们给这些图像加冕，并将它们与世界上的哲学家——即\Person{毕达哥拉斯}、\Person{柏拉图}、\Person{亚里士多德}等人的图像——并排摆放。他们还有其他尊崇这些图像的方式，如同外邦人一样。\par

\SourceRecord{Translated by Alexander Roberts and William Rambaut.；Ante-Nicene Fathers；Vol. 1.；Edited by Alexander Roberts, James Donaldson, and A. Cleveland Coxe.；Buffalo, NY: Christian Literature Publishing Co.,；1885.；<http://www.newadvent.org/fathers/0103125.htm>.}

% structure: p0001 source=h1 target=section reason=page-title

\section{反异端（第一卷，第二十六章）}

1. \Person{塞林托}，又是一个在埃及人的智慧中受教育的人，教导说世界不是由原始的天主创造的，而是由某种远隔他且远离那超越宇宙的至上元首的能力所创造，并且不认识那超越万有的天主。他声称耶稣不是由童贞女所生，而是依照人类生育的通常过程为\Person{若瑟}和\Person{玛利亚}的儿子，但他在义德、谨慎和智慧上却比其他人更优越。此外，在他受圣洗后，基督以鸽子的形状从至高主宰那里降临在他身上，并且那时他宣报了未知的圣父，并行了奇迹。但最后基督离开了耶稣，然后耶稣受难并复活了，而基督却保持了不受苦，因为他是个属神的存在。\par

2. 那些被称为厄彼翁派的人同意世界是由天主创造的；但他们对于主的意见与\Person{塞林托}和\Person{卡波克拉特斯}相似。他们只使用\BibleBook{玛窦}{福音}，并否认宗徒 \Person{保禄}，声称他是法律的背教者。关于先知著作，他们试图以某种奇特的方式解释；他们实行割损，坚持遵守法律所命令的习俗，他们的生活方式如此犹太化，以致他们甚至崇敬\Place{耶路撒冷}，仿佛它是天主的住所。\par

3. 尼苛劳派是那个\Person{尼苛劳}的追随者，他是宗徒们最初祝圣的七位执事之一。他们过着放纵无度的生活。这些人的品行在\BibleBook{若望}{默示录}中非常清楚地被指出来，［他们］被描述为教导说，行奸淫和吃祭偶像之物是无关紧要的。因此，圣言也这样论及他们说：\BibleQuote{然而，你有一件可取的事，就是你恨恶尼苛劳派的作为，我也恨恶。}\BibleRef{默 2:6}\par

\SourceRecord{Translated by Alexander Roberts and William Rambaut.；Ante-Nicene Fathers；Vol. 1.；Edited by Alexander Roberts, James Donaldson, and A. Cleveland Coxe.；Buffalo, NY: Christian Literature Publishing Co.,；1885.；<http://www.newadvent.org/fathers/0103126.htm>.}

% structure: p0001 source=h1 target=section reason=page-title

\section{反异端 第一卷 第二十七章}

1. \Person{切尔东}从\Person{西满}的追随者中采纳其体系，来到\Place{罗马}，时值\Person{希吉诺}在位，他在宗徒以降的主教传承中位居第九。他教导说：法律和先知所宣告的天主，并非我主耶稣基督的圣父。因为前者是已知的，后者却是未知的；而且前者是正义的，后者却是仁慈的。\par

2. \Place{本都}的\Person{马尔基昂}继承了他，并发展了他的学说。他借此对法律和先知所宣告为天主的那一位，提出了最狂妄的亵渎，宣称祂是万恶之源，喜爱战争，意志薄弱，甚至自相矛盾。但耶稣出自那高于创造世界的天主的父，在\Person{提庇留凯撒}的代理官\Person{本丢彼拉多}总督的时代，进入\Place{犹太}，以人的形状显现给在犹太的人，废除法律和先知，以及那创造世界的天主的一切作为——他也称这位天主为\OriginalTerm{宇宙统治者}{Cosmocrator}。此外，他割裂了根据\Person{路加}所写的福音，删去所有关于主诞生的记载，并摒弃了大量主的教导，在这些教导中，主被记载为极其明确地承认宇宙的创造者是祂的圣父。他同样说服他的门徒，他自己比那些将福音传给我们的人更值得信赖，他给他们的不是福音，而仅仅是福音的一部分。同样地，他也肢解了\Person{保禄}的书信，删去宗徒所说关于那创造世界的天主的一切，即祂是我们主耶稣基督的圣父，也删去宗徒引用先知著作的段落，那些段落是为了教导我们，他们预先宣布了主的来临。\par

3. 只有那些学习了其学说的灵魂才能获得救恩；而身体，由于取自尘土，不能分享救恩。除了对天主本身的亵渎之外，他还提出了这一点，真正是像用魔鬼的口说话，说出一切与真理直接相反的话——就是说：\Person{加音}及其同类、索多玛人、埃及人和其他像他们的人，以及一切行各种可憎之事的外邦人，当主降入\Place{阴府}时，他们都跑到主那里去了，因此得救，并迎接主进入他们的王国。但住在马尔基昂内的毒蛇却宣称：\Person{亚伯尔}、\Person{哈诺客}、\Person{诺厄}、以及从圣祖\Person{亚巴郎}所出的其他义人，连同所有先知和那些取悦天主的人，都没有分享救恩。因为这些人，他说，知道他们的天主常常试探他们，所以现在他们也怀疑祂是在试探他们，就没有跑到耶稣那里去，也不相信祂的宣告：因此他宣称他们的灵魂仍留在阴府。\par

4. 但是，既然此人是唯一敢于公开割裂圣经，并比所有其他人更无耻地抨击天主的人，我特意要驳斥他，从他的著作中定他的罪；并且，依靠天主的助佑，我将用主和宗徒们那些对他有权威、他所引用的言论来推翻他。然而现在，我只是被引导提及他，好让你们知道，所有以任何方式败坏真理、损害教会宣讲的人，都是\Place{撒玛黎雅}的\Person{西满}术士的门徒和继承者。虽然他们不承认他们师傅的名字，以便更多地诱惑他人，但他们确实教导他的学说。的确，他们将基督耶稣的名字作为一种诱饵提出，但以各种方式引入西满的不敬神之言论；他们就这样毁灭众人，邪恶地利用一个好名声来散布自己的教义，并通过其甜蜜与美好，将毒蛇——那背道的大源头的苦涩恶毒的毒液——传播给听众。\BibleRef{默 12:9}\par

\SourceRecord{Translated by Alexander Roberts and William Rambaut.；Ante-Nicene Fathers；Vol. 1.；Edited by Alexander Roberts, James Donaldson, and A. Cleveland Coxe.；Buffalo, NY: Christian Literature Publishing Co.,；1885.；<http://www.newadvent.org/fathers/0103127.htm>.}

% structure: p0001 source=h1 target=section reason=page-title

\section{反异端论（卷一，第二十八章）}

1. 从我们已经描述的那些异端者中，已经产生了许多异端的枝节。这源于他们中的许多人——实际上，我们可以说所有人——都渴望自己成为教师，并脱离他们曾经卷入的特定异端。他们从一个完全不同的意见体系中形成一套教义，然后又从其他体系形成另一套，坚持教导某种新东西，宣称自己是他们所能构想出的任何意见的发明者。举例来说：从撒托尔里诺和玛西昂衍生出那些被称为\OriginalTerm{刻苦己身者}{Encratites}的人，他们反对婚姻，从而废弃了天主起初的创造，间接责备了那造男造女以繁衍人类的天主。其中有些人还引入了禁食肉类，从而证明自己对那造了万物的天主忘恩负义。他们也否认那最初受造的人的救恩。然而，这种意见只是最近才在他们中间发明出来的。一个名叫塔提安的人首先引入了这种亵渎。他原是犹斯定的听众，只要他继续与犹斯定在一起，他就没有表达过这样的观点；但在犹斯定殉道后，他脱离了教会，因自认为是教师而沾沾自喜、自高自大，好像他比别人优越，于是创作了自己特有的教义类型。他像瓦伦廷的追随者一样，发明了一套关于某些无形永世的体系；同时，像玛西昂和撒托尔里诺一样，他宣称婚姻无非是败坏和淫乱。但他否认亚当的救恩的意见则完全是他自己的独创。\par

2. 另有一些人，追随巴西里得和卡尔波克拉特，引入了乱交和多个妻子，并对于吃祭偶像的肉毫不在意，声称天主并不十分在意这些事。但何必继续呢？因为要一一列举所有以某种方式偏离真理的人，是不可能的尝试。\par

\SourceRecord{Translated by Alexander Roberts and William Rambaut.；Ante-Nicene Fathers；Vol. 1.；Edited by Alexander Roberts, James Donaldson, and A. Cleveland Coxe.；Buffalo, NY: Christian Literature Publishing Co.,；1885.；<http://www.newadvent.org/fathers/0103128.htm>.}

% structure: p0001 source=h1 target=section reason=page-title

\section{驳斥异端（第一卷，第二十九章）}

1. 然而，除了这些我们已提及的西蒙派异端者之外，还有一大群诺斯替派如同蘑菇从地中长出，涌现出来并为人所知。我现在继续描述他们所持的主要意见。其中有些人提出一个永不衰老的永世，存在于童贞神灵中；他们称他为巴贝罗。他们宣称，在某处有一个不可名的圣父，他渴望向这巴贝罗启示自己。于是这思想上前，站在他面前，向他要求预知。预知出来以后，这两者又要求不朽，不朽也出来了，随后是永生。巴贝罗为此荣耀，并默观其伟大，在构思中因这伟大而喜乐，便生出了相似的光。他们宣称这既是光的开始，也是万有生成的开端；圣父看见这光，就用他自己的良善傅抹它，使之完善。此外，他们主张这就是基督；基督又按他们的说法，要求将理智赐给他作为助手；理智便随之而出。此外，圣父又发出了圣言。这样，思想与圣言、不朽与基督便结为结合；永生与意志结合，理智与预知结合。这些便荣耀了那大光和巴贝罗。\par

2. 他们还肯定，后来从思想和圣言发出了自生者，作为大光的肖像，并大受尊荣，万有都服从于他。与他同发出了真理，自生者与真理形成结合。但他们宣称，从光（即基督）和不朽发出了四个光体，围绕自生者；又从意志和永生发出了另外四个放射物，来服侍这四个光体；他们称这些为恩宠、意愿、理解和明智。其中，恩宠与第一个大光体相连；他们将他描绘为救主，并称为阿摩革内。意愿与第二个光体相连，他们又称之为拉古埃尔；理解与第三个光体相连，他们称之为达味；明智与第四个光体相连，他们称之为厄勒勒特。\par

3. 这一切既然这样安排好了，自生者又产生了一个完美而真实的人，他们称之为阿达玛斯，因为他自己从未被征服，他的根源也未被征服；他连同第一道光一起，从阿摩革内分离出来。此外，自生者还与人一同发出了完美的知识，并与那人结合；因此那人获得了对至高者的认识。童贞神灵也将不可战胜的能力赋予了他；万有便在他内安息，歌颂那伟大的永世。由此他们也宣称母亲、父亲、儿子显明出来；从人和知识产生了那棵树，他们又称之为知识本身。\par

4. 其次，他们主张，从靠近独生者的第一位天使那里，发出了圣神，他们又称之为索菲亚和普律尼科。他见其他一切都有伴侣，而自己却没有，便寻找一个可与他结合的存在；没有找到，他便竭力伸展自己，向下注视底层区域，期望在那里找到一个伴侣；仍然没有找到，他便在极大的急躁中跳跃出来，因为他没有得到父亲的同意就作了尝试。后来，在单纯和善良的影响下，他产生了一个作品，其中存有无知和鲁莽。他们宣称这个作品是原初元首，即这低层受造物的创造者。但他们叙述说，一股大能将他从他母亲那里带走，他远离她定居在低层区域，并形成了天上的穹苍，他们也肯定他居住在其中。在无知中，他形成了那些低于他的权能——天使、穹苍和一切地上的事物。他们肯定，他与胆大妄为结合，产生了邪恶、嫉妒、妒忌、愤怒和情欲。当这些被生出时，母亲索菲亚深感悲痛，逃走，升到上层区域，成为八元体中最末的一位（由上往下数）。当她这样离去时，他以为自己是唯一存在的；因此他宣称：\BibleQuote{我是忌邪的天主，在我以外没有别的神。}\BibleRef{出 20:5} 这些人所捏造的谎言就是如此。\par

\SourceRecord{Translated by Alexander Roberts and William Rambaut.；Ante-Nicene Fathers；Vol. 1.；Edited by Alexander Roberts, James Donaldson, and A. Cleveland Coxe.；Buffalo, NY: Christian Literature Publishing Co.,；1885.；<http://www.newadvent.org/fathers/0103129.htm>.}

% structure: p0001 source=h1 target=section reason=page-title

\section{反异端（第一卷，第三十章）}

1. 另有一些人离奇地宣称，在\OriginalTerm{比斯}{Bythus}的权能中，存在着一种原初的光，是可祝福的、不朽坏的、无限的；这就是万有之父，被称为第一个人。他们也主张，从他的\OriginalTerm{思想}{Ennœa}发出，产生了一个儿子，这就是人子——第二个人。在这二者之下，又有圣神；在这至高的神灵之下，元素彼此分离，即水、黑暗、深渊、混沌；他们宣称圣神运行于这混沌之上，称他为第一个女人。之后，他们宣称，第一个人与他的儿子，因那神灵（即那女人）的美貌而欢欣，并将光照耀在她身上，由她生出了一个不朽坏的光，即第三个男性，他们称之为基督——他是第一人与第二人之子，也是圣神、第一个女人的儿子。\par

2. 父亲与儿子都与那女人（他们也称她为永生之母）交合。然而，当她不能承受也不能容纳那光的伟大时，他们宣称她充满了以至溢出，在左边沸腾；于是，他们的独生子基督，属于右边，常趋向更高之处，立即与他的母亲一同被提，形成一个不朽坏的\OriginalTerm{永世}{Æon}。这就是真实而神圣的教会，它成为了万有之父、第一人、儿子、第二人、他们的儿子基督，以及所提及的那女人的称谓、汇集与合一。\par

3. 他们教导说，从那女人因沸腾而发出的权能，被洒上了光，从它起源之处向下坠落，却因自己的意志而拥有那光的洒落；他们称之为\OriginalTerm{左方（辛尼斯特拉）}{Sinistra}、\OriginalTerm{普鲁尼库斯}{Prunicus}、\OriginalTerm{索菲亚}{Sophia}，以及\OriginalTerm{阴阳同体}{masculo-feminine}。这一存在，出于其单纯，降入尚处于静止状态的水中，也赋予了它们运动，任意地作用于它们直至最深处，并从它们得了一个身体。因为他们断言，万物都向那光的洒落奔涌并依附于它，从各方包围它。若非它拥有那洒落，它或许会完全被物质性所吸收和压倒。因此，被一个由物质构成的身体所束缚并大大压累，这一存在后悔它所行的道路，试图逃离水上到它的母亲那里；但因其身体覆盖并包裹着它，它未能成功。它感到十分不安，至少试图隐藏那从上而来的光，生怕它也被下层元素损害，如同自己所遭遇的。当它从所拥有的那光的洒落中获得力量时，它又跳回并升到高处；在高处，它伸展自己，覆盖了一片空间，并用它的身体形成了这个可见的天；但它仍留在它所造的天之下，因为仍具有水体的形式。当它渴望上面的光，并通过万物获得力量后，它放下这个身体，从中解脱。他们称这权能所抛下的这个身体为来自女性的女性。\par

4. 他们还宣称，她的儿子自己也有从他母亲那里留下的一丝不朽的气，并借此运作；他变得有力，他自己，如他们所说，也从水中生出一个无母的儿子；因为他们不允许他认识一个母亲。他的儿子，又效法他的父亲，生了另一个儿子。这第三个又生了一个第四个；第四个也生了一个儿子；他们主张第五个又生了一个儿子；第六个也生了第七个。这样，根据他们，\OriginalTerm{七元}{Hebdomad}就完成了，母亲占据第八位；正如在他们世代中的情况，在尊严和能力方面，他们也轮流优先。\par

5. 他们还在其虚妄的体系中为各个实体起了名字，如下：那从母亲首先降生的被称为\Person{Ialdabaoth}；从他降生的被称为\Person{Iao}；从这一位降生的被称为\Person{Sabaoth}；第四个被称为\Person{Adoneus}；第五个，\Person{Eloeus}；第六个，\Person{Oreus}；第七个，即最后一个，\Person{Astanphæus}。此外，他们描绘这些天、掌权者、德能、天使和创造者，按其世代有序地坐在天上，无形地统治属天和属地的事物。他们中的第一位，即Ialdabaoth，藐视他的母亲，因为他未得任何人许可就生了儿子和孙子，甚至天使、总领天使、德能、掌权者和宰制者。这些事发生后，他的儿子们转而与他争吵，争夺至高权柄——这行为使Ialdabaoth深感忧伤，驱使他陷入绝望。在此情况下，他将目光投向下面的物质渣滓，并将欲望固定其上，他们宣称他的儿子即由此而来。这个儿子就是\OriginalTerm{努斯}{Nous}本身，扭曲成蛇的形状；从此生出灵、魂和一切尘世之物；由此也产生了所有遗忘、邪恶、竞争、嫉妒和死亡。他们宣称，当这个蛇形扭曲的努斯与他们父亲在天上及乐园中时，父亲将这努斯扭曲得更加严重。\par

为此，\OriginalTerm{雅达巴沃特}{Ialdabaoth}心高气傲，自夸超越其下一切，并喊道：\Quote{我是父，是天主，在我之上没有谁。}但他的母亲听见他这样说，就向他喊道：\Quote{不要撒谎，雅达巴沃特！因为万有之父、第一位\OriginalTerm{人}{\GreekText{Ανθρωπος}}在你之上，而\OriginalTerm{人子}{\GreekText{Ανθρωπος} \GreekText{υἱὸς}}也在你之上。}于是，当众者因这新奇之声和未曾预料的宣告而惊惶，正查问那声音从哪里来时，他们断言雅达巴沃特喊道：\Quote{让我们按我们的肖像造人。}\BibleRef{创 1:26}六位权能听见这话，他们的母亲给他们灌输了一个人的意念（为藉着他耗尽他们原有的能力），便共同形成了一个极其庞大的人，无论宽度还是长度。但他只能在地上扭动，于是他们把他带到他们的父面前；\OriginalTerm{索菲亚}{\GreekText{Σοφία}}在这事上如此劳作，为要耗尽他（雅达巴沃特）所沾染的光明，使他虽仍有能力，却不能再向上面的权能自高。他们宣称，那时藉着向人吹入生命之气息，他的能力被秘密耗尽了；于是人成了拥有\OriginalTerm{理智}{\GreekText{νοῦς}}和\OriginalTerm{思想}{\GreekText{ἐνθύμησις}}的存在；他们断言这些是参与救恩的能力。他立刻向第一位人感谢，并离弃了那些创造他的人。\par

但\OriginalTerm{雅达巴沃特}{Ialdabaoth}对此心怀嫉妒，乐意设计藉着女人再次耗尽人，便从他的\OriginalTerm{思想}{\GreekText{ἐνθύμησις}}中产生一个女人，那位\OriginalTerm{普鲁尼库斯}{\GreekText{Προύνικος}}捉住她，秘密耗尽了她的能力。但其他者前来，赞叹她的美丽，给她起名\OriginalTerm{厄娃}{\GreekText{Εὔα}}，并与她坠入情网，与她生下儿子，他们也称这些儿子为天使。但他们的母亲\OriginalTerm{索菲亚}{\GreekText{Σοφία}}狡猾地策划了一个计谋，藉着蛇引诱\OriginalTerm{厄娃}{\GreekText{Εὔα}}和\OriginalTerm{亚当}{\GreekText{Ἀδάμ}}违背雅达巴沃特的命令。厄娃听从了这话，仿佛它是来自天主之子，便轻易相信了。她也说服亚当去吃那棵天主曾说不可以吃的树上的果子。他们宣称，当他们这样吃了以后，就获得了那超越万有的权能的知识，并离弃了那些创造他们的人。当\OriginalTerm{普鲁尼库斯}{\GreekText{Προύνικος}}察觉权能们被他们自己的受造物挫败时，大大欢喜，又喊道：\Quote{既然父是不朽的，那从前自称父的雅达巴沃特是个撒谎者；并且，既然人和第一位女人（圣神）先前存在，这一个（厄娃）犯了奸淫罪。}\par

然而，\OriginalTerm{雅达巴沃特}{Ialdabaoth}由于所陷入的遗忘，没有理会这些事，将\OriginalTerm{亚当}{\GreekText{Ἀδάμ}}和\OriginalTerm{厄娃}{\GreekText{Εὔα}}逐出\OriginalTerm{乐园}{\GreekText{Παράδεισος}}，因为他们违背了他的诫命。因为他想藉着厄娃生子，但未实现他的愿望，因为他的母亲在每一点上反对他，并秘密耗尽了亚当和厄娃所沾染的光明，为使那来自至高权能的圣神既不参与咒诅也不受责难（因违命所致）。他们还教导说，他们如此被剥夺了神圣本质后，被他咒诅，并从天上被抛到这世界。但那曾反对父的蛇也被他抛到这下界；然而，他使这里的权能屈从于他的权势，并生了六个儿子，他自己构成第七位，仿效那围绕父的\OriginalTerm{七元体}{\GreekText{Ἑβδομάς}}。他们进一步宣称，这些是七个\OriginalTerm{世俗魔鬼}{\GreekText{κοσμικοὶ} \GreekText{δαίμονες}}，他们总是反对并抵挡人类，因为人类是他们的父被抛到这下界的原因。\par

\OriginalTerm{亚当}{\GreekText{Ἀδάμ}}和\OriginalTerm{厄娃}{\GreekText{Εὔα}}先前有光明、清澈、仿佛灵性的身体，正如他们受造时那样；但来到这世界后，这些变成更晦暗、更粗糙、更迟钝的身体。他们的灵魂也软弱无力，因为他们从造物主那里仅仅得到属世的气息。这持续到\OriginalTerm{普鲁尼库斯}{\GreekText{Προύνικος}}出于对他们的怜悯，将光明的甜美香气归还给他们，藉此他们记起了自己，并知道自己是赤身的，也知道身体是物质实体，从而认识到自己带着死亡。于是他们变得忍耐，知道只在短时间内被包裹在身体里。他们还找到了食物，藉着\OriginalTerm{索菲亚}{\GreekText{Σοφία}}的引导；当他们满足后，便有了肉体的认识，生下\OriginalTerm{加音}{\GreekText{Κάϊν}}，那曾连同其儿子被抛下的蛇立刻捉住他，用世俗的遗忘充满他，催使他陷入愚妄和鲁莽，以致他杀了弟弟\OriginalTerm{亚伯尔}{\GreekText{Ἄβελ}}，从而首先显露出嫉妒和死亡。此后，他们断言，藉着普鲁尼库斯的预见，生下了\OriginalTerm{舍特}{\GreekText{Σήθ}}，然后生下了\OriginalTerm{诺利亚}{\GreekText{Νωρία}}，他们声称所有其余人类都是从她而生的。他们被低级的\OriginalTerm{七元体}{\GreekText{Ἑβδομάς}}驱使行各种邪恶，被高级的圣洁\OriginalTerm{七元体}{\GreekText{Ἑβδομάς}}驱使行背道、偶像崇拜和普遍藐视一切，因为母亲总是暗中反对他们，并小心保存她自己的东西，即光明的沾染。他们还主张，圣洁七元体是七颗星，他们称之为行星；他们断言被抛下的蛇有两个名字：\OriginalTerm{弥额尔}{\GreekText{Μιχαήλ}}和\OriginalTerm{撒摩耳}{\GreekText{Σαμαήλ}}。\par

10. 再者，雅达巴沃特因世人不尊崇他、不奉他为父和天主而恼怒，遂降洪水灭尽世人。但索菲亚在此事上亦反对他，诺厄及其家人藉着从她发出的那光的洒布而在方舟中得救，世界又藉此充满人类。雅达巴沃特亲自从这些人中拣选了一个名叫亚巴郎的人，与他订立盟约，若他的后裔继续侍奉他，他就将土地赐给他们为产业。后来，他藉着梅瑟将亚巴郎的后裔从埃及领出，赐给他们法律，使他们成为犹太人。在这民族中，他选择了七天，他们也称之为圣七联。每一位都领受自己的使者，为荣耀并宣扬天主；这样，当其余人听见这些赞美时，他们也会侍奉那些由先知们宣告为神明的。\par

11. 此外，他们如下分配先知：梅瑟、农的儿子若苏厄、亚毛斯、哈巴谷属于雅达巴沃特；撒慕尔、纳堂、约纳、米该亚属于雅乌；厄里亚、约厄尔、匝加利亚属于撒巴沃特；依撒意亚、厄则克耳、耶肋米亚、达内尔属于阿多奈；多俾亚和哈盖属于厄洛伊；米该雅和纳鸿属于奥勒乌；厄斯德拉和索福尼亚属于阿斯唐斐。如此，他们各人荣耀自己的父和天主，并声称索菲亚自己也藉着他们说了许多关于第一个人和那超越的基督的事，以劝诫和提醒人们那不朽之光、第一个人以及基督的降临。诸权力因这些事惊惧，并对先知所宣布的新奇之事感到诧异，普吕尼库斯便藉着雅达巴沃特（他不知自己所做的是什么）使两个人生育降临，一个来自荒胎的依撒伯尔，另一个来自童贞玛利亚。\par

12. 由于她自己（即下界的索菲亚）在天上或地上都不得安宁，便呼求她的母亲在她困苦中协助她。于是，她的母亲，即第一个女人，因女儿悔改而动了怜悯之心，向第一个人祈求，让基督被派来援助她；基督被派来后，降临到他的妹妹那里，也降临到那光的洒布中。当他认出她（即下界的索菲亚）时，她的兄弟降临到她那里，藉着若翰宣告他的来临，并预备了悔改的洗礼，预先接纳了耶稣，为的是当基督降临时，他能找到一个纯洁的器皿，并且藉着雅达巴沃特的儿子，那妇人得以被基督宣告。他们进一步宣称，基督穿过七个天，取了它们儿子的形像，逐步剥夺了它们的能力。因为他们声称，整个光的洒布都涌向他，基督降临此世后，首先将之披在他妹妹索菲亚身上，然后两者在彼此的交往中欢欣互慰：他们将此场景描述为新郎和新娘。但耶稣，因他由天主的能力从童贞女所生，比所有人更智慧、更纯洁、更正义：基督与索菲亚联合，降临到他内，于是产生了耶稣基督。\par

13. 他们声称，许多门徒并不知道基督降临在他内；但当基督降临在耶稣身上时，他才开始行奇迹、治病、宣布那未知的父，并公开承认自己是第一个人之子。诸权力和耶稣的父对这些事发怒，竭力要消灭他；当他为此被带去时，他们说基督本人连同索菲亚离开了他，进入不朽的永世，而耶稣被钉在十字架上。然而，基督没有忘记他的耶稣，而是从上头将一种能力遣入他内，使他在那身体中再次复活，他们称这身体既是魂魄性的，也是属灵的；因为他将属世的部分又送回世界。当门徒看见他复活了，他们并没有认出他——甚至耶稣本人（藉着他从死者中复活的那位）也没有被认出。他们声称在门徒中流传着这种极大的错谬，即他们以为他是在一个属世的身体中复活的，不知道\Quote{血肉不能承受天主的国}。\par

14. 他们试图以基督在受洗前和从死者中复活后，门徒并未记载他行过什么大能的事来确立基督的降临和升天，不知道耶稣是与基督联合，不朽的永世是与七联联合；他们宣称他的属世身体与动物的本性相同。但他复活后在地上逗留了十八个月，然后从上天降临的知识教导了清楚的事。他将这些事教导给少数门徒，他知道他们能领悟如此伟大的奥秘，然后被接升天，基督坐在他父雅达巴沃特的右边，为要收纳那些认识这些事的灵魂（在它们脱去属世肉体之后），这样他便在不被他父亲知道或察觉的情况下充实自己；因此，耶稣以圣善灵魂充实自己多少，他父亲就相应地遭受损失和减少，被这些灵魂夺去自己的权力。因为他现在不会再拥有圣善灵魂将他们再派到世界，除非那些属于他本体的灵魂，即那些他吹入气息的灵魂。但万物的终结将在整个光的灵之洒布被聚集并带走形成一个不朽的永世时发生。\par

15. 这些便是流行于这些人中的见解；从\Person{瓦伦廷}学派中，像\OriginalTerm{勒耳那水蛇}{Lernæan hydra}一般，一个多头怪兽被生了出来。因为他们中有些人宣称\Person{索菲亚}自己变成了蛇；因此她敌视亚当的创造者，并将知识植入人类，故此蛇被称为比一切其他更聪明。此外，通过我们肠子的位置（食物正是通过它们输送），以及它们具有如此形状的事实，我们体内蛇形的构造揭示了隐藏的\OriginalTerm{生成者}{generatrix}。\par

\SourceRecord{Translated by Alexander Roberts and William Rambaut.；Ante-Nicene Fathers；Vol. 1.；Edited by Alexander Roberts, James Donaldson, and A. Cleveland Coxe.；Buffalo, NY: Christian Literature Publishing Co.,；1885.；<http://www.newadvent.org/fathers/0103130.htm>.}

% structure: p0001 source=h1 target=section reason=page-title

\section{驳斥异端（卷一，第三十一章）}

1. 另有一些人宣称加音是从上方的权能获得存有，并承认厄撒乌、科辣黑、索多玛人以及所有这类人物都与他们有关。为此，他们补充说，他们曾受造物主攻击，但其中没有一人受到伤害。因为索菲亚习惯于将他们所有之物从他们那里取归自己。他们宣称叛徒犹大完全了解这些事，并且唯有他，像没有人那样知道真理，完成了出卖的奥迹；藉着他，一切事物，无论属地的属天的，都陷入了混乱。他们炮制了这样一种虚构的历史，称之为\WorkTitle{犹大福音}。\par

2. 我也收集了他们的著作，在其中他们主张废除赫斯特拉的行为。此外，他们称这赫斯特拉为天地的创造者。他们也像卡尔波克拉特那样主张，人除非经历各种经验，否则不能得救。他们坚持认为，在他们每一次罪恶可憎的行为中，都有一位天使陪伴着他们，并促使他们冒险行恶、招致玷污。无论行为性质如何，他们都宣称是奉那位天使的名而行，说：\Quote{哦，天使，我使用你的工作；哦，权能，我完成你的运作！}他们维护说这就是\Quote{圆满的知识}，毫不畏缩地冲入那些甚至不可提及的行为中。\par

3. 有必要清楚证明，正如他们的见解和规条所显示的那样，那些属于瓦伦提诺学派的人起源于这样的母亲、父亲和祖先，并且也要提出他们的教义，希望他们中或许有人会悔改，回归唯一的造物主和宇宙的创造者天主，从而获得救恩，并且使其他人从今以后不再被他们邪恶而似是而非的劝说所迷惑，以为会从他们那里得到一些更伟大更崇高的奥迹的知识。但愿他们反而从我们这里有效地学习这些人的邪恶主张，鄙弃他们的教义，同时怜悯那些仍然固守这些可怜且毫无根据的寓言的人，他们已狂妄到因这种知识——或更确切地说，无知——而自以为超越他人。他们现在已被完全揭露；仅仅展示他们的观点，就是战胜他们。\par

4. 因此我努力提出并清楚地揭示这只可怜小狐狸极其病态的尸体。\BibleRef{歌 2:15}因为当他们的教义体系已向众人显明时，无需多言便能推翻它。正如一只野兽隐藏于树林，常常冲出残害多人，有人绕树搜寻并彻底探查，迫使它暴露，并不试图捕获它，因为它实为凶残之兽；但旁观者可以留意并躲避它的攻击，从四面八方投掷标枪，刺伤它，最终击杀这有害的畜生。同样，我们既已将他们彼此间缄默隐藏的奥迹带到光天化日之下，现在便无需多言摧毁他们的意见体系。因为你们和所有同道现在能够熟悉已说的话，推翻他们邪恶而不成熟的教义，并确立符合真理的教义。既然如此，我将按照承诺，尽我所能，在下一卷书中通过驳倒他们所有人来努力推翻他们。即使仅仅讲述这些事也是冗长乏味的，如你们所见。但我要提供推翻他们的方法，按照已描述的顺序回应他们所有的意见，不仅使野兽暴露于众，而且从四面八方予以创伤。\par

\SourceRecord{Translated by Alexander Roberts and William Rambaut.；Ante-Nicene Fathers；Vol. 1.；Edited by Alexander Roberts, James Donaldson, and A. Cleveland Coxe.；Buffalo, NY: Christian Literature Publishing Co.,；1885.；<http://www.newadvent.org/fathers/0103131.htm>.}

% structure: p0001 source=h1 target=section reason=page-title

\section{\WorkTitle{驳斥异端（第二卷，前言）}}

1. 在紧接此卷的第一卷书中，我揭露了那些属于\Person{瓦伦提诺}学派的人所制定的、以多种对立方式构成的整个体系是虚假且毫无根据的，正如\BibleRef{弟前 6:20}所言，是\Quote{伪知识}。我亲爱的朋友，我已向你展示这一切。我还陈述了其前辈的教义，证明他们不仅彼此分歧，而且早已背离了真理本身。此外，我详尽解释了魔术师\Person{马可}的教义与实践，因他也属于这些人；并仔细指出了他们为了适应自己的虚构而从圣经中曲解的段落。再者，我详细叙述了他们如何通过数字和字母表中的二十四个字母，大胆地试图建立（他们视为）真理。我也讲述了他们如何思考并教导万有是按照他们那不可见的\OriginalTerm{圆满}{Pleroma}的形象而形成的，以及他们关于\OriginalTerm{得缪格}{Demiurge}的观点，同时声明了他们的始祖，撒玛黎雅的\Person{西满}魔术师以及所有继承他之人的教义。我还提到了从他衍生的众多\OriginalTerm{诺斯替派}{Gnostics}，并指出他们之间的差异、各自的教义以及承继次序，同时陈述了由他们起源的所有这些异端。此外，我表明所有这些异端者，源于\Person{西满}，已将不敬神和不信神的教义引入此生；并解释了他们的\Quote{救赎}的本质，以及他们使那些成为\Quote{成全}者的人们入教的方法，连同他们的祈祷和奥秘。我也证明了天主是唯一的创造者，祂并非任何缺陷的产物，也没有任何事物在祂之上或在祂之后。\newline{}2. 在本卷书中，我将在时间允许的范围内，确立那些符合我计划的观点，并通过按不同标题进行的详细论述来推翻他们的整个体系；因此，既然这是一项对其意见的揭露和颠覆，我便这样命名了这部作品的创作。因为通过对其结合体的清晰揭露和颠覆，来终结这些隐藏的联合并终结\OriginalTerm{深渊}{Bythus}本身，从而证明他无论在过去还是现在都从未存在过，这是合宜的。\par

\SourceRecord{Translated by Alexander Roberts and William Rambaut.；Ante-Nicene Fathers；Vol. 1.；Edited by Alexander Roberts, James Donaldson, and A. Cleveland Coxe.；Buffalo, NY: Christian Literature Publishing Co.,；1885.；<http://www.newadvent.org/fathers/0103200.htm>.}

% structure: p0001 source=h1 target=section reason=page-title

\section{\WorkTitle{反异端（第二卷，第一章）}}

1. 那么，我应当从首要和最重要的论点开始，即创造主天主，祂创造了天、地及其中万物（这些人亵渎地称之为缺陷的产物），并证明在祂之上或之后没有任何存在；祂并非受任何一位影响，而是凭自己的自由意志创造了万物，因为祂是唯一的天主，唯一的主，唯一的创造者，唯一的圣父，独自包含万有，并亲自命令万有存在。\par

2. 因为在祂之上怎么可能有任何其他的\OriginalTerm{圆满}{Pleroma}、元始、能力或天主呢？既然天主，即这一切的\OriginalTerm{圆满}{Pleroma}，必然以其无限包含万有，而不被任何一位所包含。但如果在祂之外有某物，那么祂就不是万有的\OriginalTerm{圆满}{Pleroma}，也不包含万有。因为那些人所宣称的在祂之外的东西，将缺席于\OriginalTerm{圆满}{Pleroma}，或者缺席于那超越万有的天主。而那缺席的、以任何方式有欠缺的，就不是万有的\OriginalTerm{圆满}{Pleroma}。在这种情况下，祂就会相对于那些在祂之外者而有开始、中间和结束。如果祂相对于下面的事物有结束，那么祂相对于上面的事物也有开始。同样，祂绝对必然在其他所有点上也经历同样的事情，并被那些在祂之外的存在所限制、界定和包围。因为那个向下结束的存在，必然限制和包围着在其内结束的那一位。因此，按照他们的说法，万有圣父（即他们称之为\OriginalTerm{先存者}{Proön}和\OriginalTerm{元始}{Proarche}的那一位），连同他们的\OriginalTerm{圆满}{Pleroma}，以及\Person{马西昂}的善天主，被建立在并封闭于另一位之中，并从外部被另一个强大的存在所包围，这个存在必然更伟大，因为包含者大于被包含者。但那个更伟大的，也是更强有力的，并在更大程度上是主；而那个更伟大、更强有力、在更大程度上是主的，必定是天主。\par

3. 既然，按照他们的说法，还存在某种他们宣称在\OriginalTerm{圆满}{Pleroma}之外的东西，并且他们还主张那个迷失的高等能力曾下降到其中，那么\OriginalTerm{圆满}{Pleroma}必然要么包含那外在的东西，却又被包含（否则它就不会在\OriginalTerm{圆满}{Pleroma}之外；因为如果有什么在\OriginalTerm{圆满}{Pleroma}之外，就会有一个\OriginalTerm{圆满}{Pleroma}在这个\OriginalTerm{圆满}{Pleroma}之内，即他们宣称在\OriginalTerm{圆满}{Pleroma}之外的那个\OriginalTerm{圆满}{Pleroma}，而\OriginalTerm{圆满}{Pleroma}会被那外在的东西所包含；并且与\OriginalTerm{圆满}{Pleroma}一起被理解的还有第一个天主）；要么，他们必须彼此相隔无限的距离——我是指\OriginalTerm{圆满}{Pleroma}和它之外的东西。但如果他们坚持这一点，那么就会有第三种存在，它以无限分隔\OriginalTerm{圆满}{Pleroma}和那之外的东西。这第三种存在因此将界定和包含另外两者，并将大于\OriginalTerm{圆满}{Pleroma}和那之外的东西，因为它将两者包含在它的怀抱中。如此，关于被包含者和包含者的谈论可以永远持续下去。因为如果这第三种存在有上的开始和下的结束，它绝对必然也在两侧被界定，或在某些其他点开始或结束（那里有新的存在开始）。这些，以及上面和下面的其他存在，将在某些其他点有其开始，如此\OriginalTerm{以至无穷}{ad infinitum}；这样，他们的思想永远不会止于一位天主，而是由于追求多于存在的东西，会漂流向不存在的东西，远离真实的天主。\par

4. 这些评论同样适用于反驳\Person{马西昂}的追随者。因为他的两位神也将被分隔他们彼此的无限间距所包含和限定。但于是必须设想无数位神，彼此被无限距离从各侧分隔，彼此开始，彼此结束。因此，正是根据他们用来教导在天地创造主之上有某个\OriginalTerm{圆满}{Pleroma}或天主的同一推理过程，任何选择使用它的人都可以主张在\OriginalTerm{圆满}{Pleroma}之上有另一个\OriginalTerm{圆满}{Pleroma}，在它之上又有另一个，在\OriginalTerm{深渊}{Bythus}之上另有神性之汪洋，同样地，在两侧也有相同的连续序列；如此，他们的学说漫延至无限，将总是有必要设想其他的\OriginalTerm{圆满}{Pleroma}和其他的\OriginalTerm{深渊}{Bythus}，以致永远无法停止，而总是继续寻找除已提及者之外的其他存在。此外，我们设想的东西是在下面，还是实际就是上面的事物，将是不确定的；同样，关于他们所说的上面的事物，它们真的是上面还是下面，也是不确定的；因此我们的意见不会有固定的结论或确定性，而必然漂流向没有界限的世界和无法计数的神。\par

5. 这些事既然如此，每个神祇必会满足于自己的拥有，不会因好奇而干预他人的事务；否则，他便是不义、贪婪，并且不再成其为天主。每个受造物也必会荣耀自己的创造者，并满足于他，不认识任何其他者；否则，它极有理由被其余一切视为背教者，并受到应有的惩罚。因为，要么存在一位包含万有的存有，按自己的意愿在自己的领域内创造了所有受造之物；要么，又存在无数无限的创造者和神祇，他们彼此起始，又彼此终结于四面八方；那么，就必须承认，其余一切都被某个更伟大的存有从外部包含，并且他们各自被封在自己的领域内，居留其中。因此，他们之中没有一个能是天主。因为每一神祇都将缺乏许多，与其余一切相比，他只拥有极小部分。全能者的名号就此终结，这样的看法必然堕入不敬神之境地。\par

\SourceRecord{Translated by Alexander Roberts and William Rambaut.；Ante-Nicene Fathers；Vol. 1.；Edited by Alexander Roberts, James Donaldson, and A. Cleveland Coxe.；Buffalo, NY: Christian Literature Publishing Co.,；1885.；<http://www.newadvent.org/fathers/0103201.htm>.}

% structure: p0001 source=h1 target=section reason=page-title

\section{\WorkTitle{驳斥异端（第二卷，第二章）}}

1. 此外，那些说世界是由天使或任何其他造物者违背至高圣父的意愿而形成的人，首先在这一点上就错了：他们断言天使们违背了至高天主的意愿，创造了如此宏大而强有力的受造界。这暗示天使比天主更有能力；若不然，则说明天主或是漫不经心、或是低劣、或根本不关心自己产业中所发生的事——无论结果好坏，祂既不能驱散并阻止坏事，也不能赞美并喜悦好事。若有人甚至不愿将这样的行为归给一个稍有才干的人，更何况归给天主呢？\par

2. 其次，请他们告诉我们：这些事物是在祂所包含的界限内、在祂自己的领域内形成的，还是在属于别人、位于祂之外的区域中形成的？若他们说这些事是在祂之外发生的，那么前面提到的一切荒谬之处就会临到他们，并且至高天主将被那在祂之外的东西所包围，而祂也必在其中找到自己的终点。若这些事是在祂自己的领域内发生的，那么说世界在祂自己的领域内、违背祂的意愿、由那些本身在祂权下的天使或任何其他存在如此形成，就十分荒谬了——仿佛祂自己并未看见自己产业中所发生的一切，或不知道天使们将做的事。\par

3. 然而，若这些事并非违背祂的意愿，而是得到祂的同意和知晓——如同这些人中的某些人所认为的那样——那么天使或世界的塑造者（无论他是谁）就不再是那形成的因，而天主的意愿才是。因为若\Emph{祂}是世界的塑造者，祂也创造了天使，或至少是天使受造的原因；而那位预备了形成之因的，便应被视为创造了世界。尽管这些人主张天使是由一连串长久的下传而产生，或世界的塑造者源自至高圣父（如\Person{巴西里德}所断言），但那造成这些受造之物的原因仍会追溯至那位创始如此传承者。正如战争的成功归于那位预备了胜利之因的君王；同样，任何城邦或工程的创立，都归于那位为后来成就之事预备了材料的人。因此，我们不说斧子砍了木头，也不说锯子分了木头，而应恰当地说，是那位为此目的造了斧子和锯子的人，并且更早地造了制造斧子和锯子本身所需的一切工具。因此，按照类似的推理，万有的圣父必被宣告为这个世界的塑造者，而不是天使或任何其他所谓的世界的塑造者——除了那位作为世界创始者、最初就已为这类受造预备了原因者之外。\par

4. 这种说法对于那些不认识天主、将祂比作有需要的人、比作那些不能立即且无需帮助而创造任何事物、却需要许多工具来实现其意图的人，也许看似合理或有说服力。但对于那些知道天主一无所缺、祂凭自己的圣言创造并造成了万有的人而言，这说法绝不可信。祂既不需要天使帮助祂生产那些受造之物，也不需要任何远逊于祂、不认识父的权能，也不需要任何缺陷或无知，为叫那认识祂的人成为人。祂自己在自身内，以一种我们既无法描述也无法想象的方式，预定了万有，并按照祂的喜悦形成了它们，将和谐赋予万有，为它们指定了各自的位置和受造的开始。如此，祂将属灵的、不可见的本性赋予了属灵之物，将属天的本性赋予了超天之物，将天使的本性赋予了天使，将动物的本性赋予了动物，将适于水的本性赋予了水生之物，将适于陆地的本性赋予了陆地生灵——总之，祂为一切赋予了适于所指定生命性质的本性；祂凭那永不疲倦的圣言造成了所有受造之物。\par

5. 因为这是天主卓越性的一个特征：祂不需要其他工具来创造那些被召唤而存在的事物。祂自己的圣言适于并足以形成万有，正如主的门徒\Person{若望}论及祂时所说：\BibleQuote{万物是藉着祂造的；凡被造的，没有一样不是藉着祂造的。}\BibleRef{若 1:3}在这\Quote{万物}中，我们的世界必然被包括在内。因此，世界也是藉着祂的圣言所造，正如圣经在\BibleBook{}{创世纪}中告诉我们，祂藉着祂的圣言造了与我们世界有关的一切。\Person{达味}也表达了同样的真理，他说：\Quote{因为祂说了，它们就造成；祂命令了，它们就创造。}那么，关于世界的创造，我们应当相信谁呢？是那些在此事上胡言乱语、自相矛盾的上述异端者，还是主的门徒以及\Person{梅瑟}——他既是天主忠信的仆人，又是先知？\Person{梅瑟}最初这样叙述世界的形成：\BibleQuote{在起初天主创造了天地}\BibleRef{创 1:1}，以及其他一切依次受造之物；其中既没有神也没有天使参与。\par

现在，这位天主就是我们的主耶稣基督的圣父，宗徒保禄也曾宣告说：\BibleQuote{只有一个天主，圣父，他在万有之上，贯穿万有，并在我们众人之内。}我确已证明只有一个天主；但我还要从宗徒们自身、从主的言论进一步证明这一点。因为我们若离弃先知、主和宗徒的言论，而去听从这些不说一句明智话的人，那会是什么样的行为呢？\par

\SourceRecord{Translated by Alexander Roberts and William Rambaut.；Ante-Nicene Fathers；Vol. 1.；Edited by Alexander Roberts, James Donaldson, and A. Cleveland Coxe.；Buffalo, NY: Christian Literature Publishing Co.,；1885.；<http://www.newadvent.org/fathers/0103202.htm>.}

% structure: p0001 source=h1 target=section reason=page-title

\section{驳斥异端（第二卷，第三章）}

1. 因此，他们所设想的深渊及其圆满，以及马西昂的天主，是前后矛盾的。倘若真如他们所说，深渊在自身之下和超越自身有一个他物，他们称之为虚空与阴影，那么此虚空便被证明比他们的圆满更大。但即使声称以下说法也是矛盾的：当深渊将万有包含于自身内时，造化却由另一位形成。因为他们绝对必须承认（在属灵的圆满之下）存在某种虚空而混沌的存在，此宇宙就在其中形成，并且元父特意将这混沌原样留下，或是预知其中将要发生的事，或是对其无知。若他真无知，那么天主便不会预知万有。但即便如此，他们也无法提出理由，说明为何他在如此漫长的时期中任由这处虚空存在。反之，若他预知，并在心意中默想了那将要在此处存在的造化，那么正是他自己创造了它，也是他事先（在理念上）在自己内形成了它。\par

2. 因此，他们应当停止断言世界是由另一位造的；因为天主一旦在心意中形成观念，那心意中所构思的便已成就。因为不可能是一位在心智中形成观念，而另一位实际产生心智所构思的事物。但按这些异端者的说法，天主或是在心智中构思了一个永恒的世界，或是一个暂时的世界，\Emph{两者}不可能同时为真。然而，若他将之构思为永恒、属灵且可见的，它也会被造成那样。但若它被造成如今的模样，那么正是\Emph{祂}构思了它如此，并使之如此；或是他按照心智的构思，在圣父的理念中意愿它存在，正如它现今的模样：复合、可变而短暂的。既然它恰如圣父在与自身商议中（在理念上）形成的样式，它必是与圣父相称的。但声称万有之父在心意中所构思并预造的，正如实际被造的那样，乃是缺陷的果实和无知的产物，这是极大的亵渎。因为按照他们的说法，万有之父由此便被视作在自己胸中，按自己的心智构想，产生了缺陷的流出和无知的果实，因为他心智中所构思的事物的确被产生了。\par

\SourceRecord{Translated by Alexander Roberts and William Rambaut.；Ante-Nicene Fathers；Vol. 1.；Edited by Alexander Roberts, James Donaldson, and A. Cleveland Coxe.；Buffalo, NY: Christian Literature Publishing Co.,；1885.；<http://www.newadvent.org/fathers/0103203.htm>.}

% structure: p0001 source=h1 target=section reason=page-title

\section{\WorkTitle{驳斥异端}（第二卷，第四章）}

1. 因此，应当探求天主如此安排的缘由，但世界的形成不可归因于任何其他存在。应当说，万物都被天主预先安排，使之得以成为它们所是的样式；但不可凭空捏造出阴影和空虚。然而，请问，他们所说的这个\Quote{空虚}从何而来？如果它确实是由那位他们称为万有之父和创作者所产生，那么它就在尊荣上与其他\OriginalTerm{永世}{Æons}相等且相关，甚至可能比它们更古老。此外，如果它来自同一源头，那么它在本质上必然类似于产生它的那一位，也类似于与它一同被产生的那些。因此，必然会有这样的绝对结果：他们所说的\OriginalTerm{深渊}{Bythus}连同\OriginalTerm{沉默}{Sige}在本质上类似于空虚，也就是说，他实际上是一个空虚；而其余\Quote{永世}，既然是空虚的兄弟，也应当缺乏实体。另一方面，如果它不是这样产生的，那么它必定是自我生成、自我产生的，在这种情况下，它在年岁上将与那位他们称为万有之父的\Quote{深渊}相等；这样，空虚将与那据他们说是宇宙之父者同性质、同尊荣。因为它要么是由某位产生的，要么是自我生成、自我产生的。但是，如果实际上空虚是被产生的，那么它的产生者\Person{瓦伦廷}也是一个空虚，他的追随者也是如此。反之，如果它不是被产生的，而是自我生成的，那么那个真正的空虚就与\Person{瓦伦廷}所宣告的那位父相似、是兄弟、同尊荣；而且它比\Person{托勒密}本人、\Person{赫拉克利翁}以及所有持相同见解者的其余\Quote{永世}更古老、更早存在、更尊贵。\par

2. 但是，如果他们对此感到绝望，并承认万有之父包含一切，在\OriginalTerm{圆满}{Pleroma}之外绝无任何存在（因为若有任何之外的，它必然被比它更大的东西所界定和限定），并且他们说\Emph{外面}和\Emph{里面}是就知识与无知而言，而非就空间距离而言；但在\Quote{圆满}中，或在父所包含的事物中，我们所知被造的全部受造界，由\OriginalTerm{造物主}{Demiurge}或天使所造，都被那不可言说的伟大所包含，如同圆心在圆内，或斑点在一件衣服上——那么，首先，那个\Quote{深渊}是怎样的存在，他竟然允许在自己的怀中有一个污点，并允许另一存在在他的领域内违背他自己的意愿创造或产生东西？这样的行为方式确实会招致对整个\Quote{圆满}的退化指控，因为它本来可以一开始就切断那个缺陷及其产生的流溢，而不应同意允许在无知、情欲或缺陷中形成受造界。因为那些后来能够纠正缺陷、仿佛洗去污点的人，本可以更早地确保在他的所有物中起初就不存在这样的污点。或者，如果他起初就允许所造之物如此（因为它们实际上无法以其他方式形成），那么它们必然永远保持同样的状态。因为那些最初无法得到纠正的事物，后来怎么可能得到纠正呢？或者，人们怎能说他们被召叫达至完美，而正是那些作为人起源的原因的存在——无论是造物主本人还是天使——却被宣称处于缺陷之中？而且，如果正如所主张的那样，至高的存在因其仁慈最终怜悯了人，并赐予他们完美，那么他本应首先怜悯那些创造人者，并赐予他们完美。这样，人也会真正分享到他的同情，由完美者造就为完美。因为如果他怜悯这些存在的\Emph{作品}，他早就该怜悯\Emph{他们本身}，而不应让他们陷入如此可怕的盲目。\par

3. 他们关于阴影和空虚的谈论——他们声称我们所关心的受造界是在其中形成的——也将归于虚无，如果所提及的事物是在父所包含的领域内被造的话。因为如果他们持守他们父的光是那样充满他内部的一切并光照一切，那么在被\OriginalTerm{圆满}{Pleroma}和父之光所包含的领域内，怎么可能存在任何空虚或阴影？因为在这种情况下，他们必须指出在\OriginalTerm{元父}{Propator}或\Quote{圆满}之内某个未被光照、也未被任何人占有的地方，而天使或造物主在那里随意形成了他们所愿的。而且，要设想如此巨大受造界的形成，所需的空间不会很小。因此，必然会有这样的绝对结果：在他们所说的\Quote{圆满}或父之内，他们必须设想某个地方是空虚、无形式、充满黑暗的，其中形成了那些已形成之物。然而，这种假设会使其父之光蒙受指责，仿佛他不能光照并充满他自身之内的事物。这样，当他们主张这些事物是缺陷的果实和错误的作品时，他们实际上把缺陷和错误引入了\Quote{圆满}和父的怀中。\par

\SourceRecord{Translated by Alexander Roberts and William Rambaut.；Ante-Nicene Fathers；Vol. 1.；Edited by Alexander Roberts, James Donaldson, and A. Cleveland Coxe.；Buffalo, NY: Christian Literature Publishing Co.,；1885.；<http://www.newadvent.org/fathers/0103204.htm>.}

% structure: p0001 source=h1 target=section reason=page-title

\section{\WorkTitle{驳斥异端}（卷二，第五章）}

1. 因此，我刚才所作的评论适合于回答那些断言这个世界的形成是在圆满之外，或在一位\Quote{善天主}之下的人；而这样的人，连同他们所谈论的圣父，将被完全隔绝于圆满之外，而他们最终也必然停留在那里。再者，回答那些主张这个世界是由某些其他存在者在圣父所包含的领域内形成的人，现在已经注意到的所有那些论点都将呈现出其荒谬和不连贯之处；他们将被迫要么承认圣父内的一切都是光明、充满和有力的，要么指责圣父的光明，仿佛祂不能光照万物；或者，作为他们圆满的一部分，必须承认整个圆满都是空虚、混沌、充满黑暗的。他们指责所有其他受造物，仿佛这些仅仅是暂时的，或者（至多）若是永恒的，也只是物质的。但这些（永世）应被认为超出此类指责的范围，因为他们是在圆满之内的；否则，所指责的问题也将同样落在整个圆满上；因此，他们所谈论的基督被发现是无知的始作俑者。因为按照他们的说法，当祂给予他们所设想的那位母亲以形式（就实体而言）时，祂将她抛出了圆满；也就是说，祂将她与知识隔绝了。因此，那将她与知识分离的，实际上在她内产生了无知。那么，同一个人怎么能既将知识恩赐给其余那些比祂（在被产生方面）更早的永世，却又成为祂母亲的无知之源？因为祂将她置于知识之外，当祂将她抛出圆满时。\par

2. 此外，如果他们像他们中某些人那样，将圆满之内和之外解释为分别指知识和无知（因为拥有知识的人是在那知道者之内），那么他们必然要承认救主自己（他们称之为\Emph{万有}）是处于无知状态的。因为他们主张，当祂来到圆满之外时，祂将形式赋予了他们的母亲（阿卡莫特）。那么，如果他们断言，凡是在圆满之外的都一无所知，而救主出去是为了赋予他们的母亲以形式，那么祂就置身于一切知识之外；也就是说，祂处于无知中。那么，当祂自己处于知识之外时，怎么能将知识传达给她呢？因为我们也处于圆满之外，他们说，因为我们处于他们所拥有的知识之外。而且再一次：如果救主真正走出圆满去寻找那只迷失的羊，但圆满与知识是等同的，那么祂就将自己置于知识之外，也就是说，在无知中。因为他们必须承认，要么圆满之外是地方意义上的，在这种情况下，先前所作的所有评论都将反对他们；要么如果他们谈论内在方面指知识，外在方面指无知，那么他们的救主，以及远在他之前的基督，必定是在无知中形成的，因为他们走出了圆满，即知识之外，以便将形式赋予他们的母亲。\par

3. 这些论证同样可以适用于所有那些以任何方式主张这个世界是由天使或任何其他不同于真天主者形成的人。因为他们对德穆革以及那些被造成物质和暂时之物所提出的指控，实际上将落回圣父身上；如果确实在圆满怀中形成的事物，最后竟因圣父的许可和善意而开始解体。那么，这位（直接）创造者并不是这作品的（真正）作者，祂以为自己造得非常美好，但那位允许并认可缺陷的产物和谬误的作品存在于祂自己领域之内，并允许暂时与永恒、可朽与不朽、以及那些分享谬误的与那些属于真理的混杂在一起者才是。但如果这些事物是在没有万有圣父的许可或认可下形成的，那么那位在其（圣父）本应属己的领域内造了这些事物而未经祂许可的存在者，就必定更强大、更有力、更具君王性。再次，如果如某些人所说，他们的圣父许可了这些事物但并不认可，那么祂是因某种必要而给予许可，要么能够阻止（这一过程），要么不能。但如果祂不能（阻止），那么祂是软弱无力的；而如果祂能，祂就是一个引诱者、一个伪善者、一个必然性的奴隶，因为祂不同意（这样的进程），却又允许它，仿佛祂同意似的。并且祂允许谬误起初产生并逐渐增长，却在后来的时候试图摧毁它，而此时已有许多人因那（最初的）缺陷而悲惨地丧亡了。\par

然而，说统驭万有的天主是必然的奴隶，因为祂是自由而独立的，是不合适的；或者说任何事都是经祂许可却违背祂意愿而发生的；否则，他们就会使必然比天主更伟大、更具王权，因为最有能力的事物超越一切。祂本应在最初就斩断（想象的）必然的根源，而不应该让自己陷入屈服于那必然的境地，通过许可任何与祂相宜之事以外的事。因为，从最初就斩断这类必然的原则，远胜过后来仿佛因后悔而试图在必然的结果发展到如此地步时将其根除，这样做要好得多、一致得多，也更像天主。如果万有的圣父是必然的奴隶，必须屈服于命运，同时祂不情愿地容忍所发生的事，却又无力反对必然和命运（如同荷马笔下的宙斯，论到必然时说：\Quote{我心甘情愿地给了你，却心不甘愿不情愿}），那么，按照这种推理，他们所说的\OriginalTerm{深渊}{Bythus}就会被认为是必然和命运的奴隶。\par

\SourceRecord{Translated by Alexander Roberts and William Rambaut.；Ante-Nicene Fathers；Vol. 1.；Edited by Alexander Roberts, James Donaldson, and A. Cleveland Coxe.；Buffalo, NY: Christian Literature Publishing Co.,；1885.；<http://www.newadvent.org/fathers/0103205.htm>.}

% structure: p0001 source=h1 target=section reason=page-title

\section{Against Heresies (Book II, Chapter 6)}

1. 再者，天使或世界的创造者怎么可能不认识至高天主呢？他们原是祂的财产、祂的受造物，并被祂所包含。祂确实可能因超越性而对他们不可见，但祂绝不可能因眷顾而为他们所不知。因为尽管如他们所言，他们因本性的低下而与祂相距甚远，然而，既然祂的统治延伸至他们全体，他们理应认识自己的统治者，并特别意识到：创造他们的那一位是万有之主。因为祂那不可见的本质是大能的，它赐予众人一种深邃的内在直觉和感知，对祂那最强大、甚至全能伟大的感知。因此，尽管\Quote{除了圣子，没有人认识圣父，除了圣父，也没有人认识圣子，以及圣子愿意启示的人}，\BibleRef{玛 11:27}，然而所有存在物至少知道这一事实：因为植入他们理智中的理性推动他们，并向他们启示有一位天主，万有之主。\par

2. 因此，万有都（普遍一致地）被置于那位被称为至高者、全能者的权下。藉着呼求祂，甚至在我们主来临之前，人们就从最邪恶的灵、各类魔鬼及各种背逆权势中被拯救。这并非因为地上的灵或魔鬼看见了祂，而是因为他们知道那统御万有的天主存在，在呼求祂时他们战栗，如同在祂治理下的每一个受造物、统治权、权势及一切具有活力的存在都战栗一样。类比而言，那些生活在罗马帝国统治下的人，虽然从未见过皇帝，且与他相隔海陆，岂不因经历他的统治而很清楚是谁拥有国家的主权吗？那么，那些在我们之上（本性上）的天使，甚至他们称为世界创造者的那一位，怎么可能不认识全能者呢？连哑巴牲畜在呼求祂的名时也战栗屈服。而且，虽然他们未曾见过祂，但万有都臣服于我们主的名，同样也必臣服于那藉祂的圣言创造并建立万有的那一位，因为正是祂形成了世界。为此，犹太人至今也藉着这种呼求来驱赶魔鬼，因为一切受造物都畏惧那创造它们之主的呼求。\par

3. 那么，如果他们不愿承认天使比哑巴牲畜更无理性，他们就会发现，天使理应——尽管未见过那统御万有的天主——认识祂的能力与主权。因为如果他们主张，他们这些住在地上的人尚且认识从未见过的统御万有的天主，却不允许那据他们所说形成了他们及整个世界、且住在高处、在诸天之上的那一位，知道他们这些住在地上的人所知道的事，那就显得极为可笑。除非他们或许主张\OriginalTerm{深渊}{Bythus}住在地下的\OriginalTerm{塔耳塔罗斯}{Tartarus}，因此他们得以在那些住高处的天使之前认识祂。他们就这样陷入疯狂的深渊，竟宣称世界的创造者没有理解力。他们实在可怜，竟以如此荒唐的愚昧断言，祂（世界的创造者）既不认识祂的母亲，也不认识她的后裔，不认识\OriginalTerm{圆满}{Pleroma}的\OriginalTerm{永世}{Æons}，不认识\OriginalTerm{原始父}{Propator}，也不知道祂所造之物是什么；而是说这些是圆满之内那些事物的影像，救主暗中工作，使它们被（无意识的\OriginalTerm{德穆革}{Demiurge}）如此形成，以荣耀上面的事物。\par

\SourceRecord{Translated by Alexander Roberts and William Rambaut.；Ante-Nicene Fathers；Vol. 1.；Edited by Alexander Roberts, James Donaldson, and A. Cleveland Coxe.；Buffalo, NY: Christian Literature Publishing Co.,；1885.；<http://www.newadvent.org/fathers/0103206.htm>.}

% structure: p0001 source=h1 target=section reason=page-title

\section{\WorkTitle{驳斥异端}（第二卷，第七章）}

1. 正当\OriginalTerm{德谬哥}{Demiurge}对万物如此无知时，他们告诉我们，救主通过他的母亲所招致的受造界，将尊荣归于\OriginalTerm{圆满}{Pleroma}，因为他产生了上界事物的相似与肖像。但我已经证明，在圆满之外（他们告诉我们，在这个外部区域，制作了圆满内事物的肖像）不可能有任何东西存在，或者这个世界是由至高的天主以外的任何一位形成的。但是，从各方面推翻他们并证明他们是谎言的贩卖者，乃是一件乐事；让我们反对他们说：如果这些事物是由救主按照上界事物的肖像所造，为尊崇那些上界事物，那么这些事物就应该永远存续，以使那些受到尊崇的事物得以永恒地保持在尊荣中。但如果它们实际上消逝了，那么这段极其短暂的尊荣（一种曾经不存在、且将再次归于无有的尊荣）有什么用呢？那样的话，我将证明救主与其说是尊崇上界事物，不如说是追求虚荣。因为那些暂时的事物能给永恒持久的事物带来什么尊荣呢？那些消逝的事物能给恒存的事物带来什么尊荣呢？那些可朽的事物能给不朽的事物带来什么尊荣呢？——因为即使在终有一死的人类中，人们也不看重迅速消失的尊荣，而是看重尽可能持久的尊荣。但是那些刚造出来就终结的事物，与其说它们是为了尊崇那些被认为受它们尊崇的事物而形成的，不如说是为了轻蔑它们；当永恒之物的肖像被败坏和瓦解时，那永恒之物受到了侮辱。但如果他们的母亲没有哭泣、欢笑并陷入绝望呢？那样救主就没有任何办法尊崇\OriginalTerm{圆满}{Fulness}，因为她最后的混乱状态本身没有实体可以用来尊崇\OriginalTerm{太初之父}{Propator}。\par

2. 唉，这虚荣的尊荣即刻消逝，不再显现！将会有某个\OriginalTerm{移涌}{Æon}，对于它而言，这样的尊荣根本不会被认为曾经存在过，然后上界事物将不受尊崇；或者必须再次产生另一个哭泣、绝望的母亲，以尊崇圆满。何等不同且同时亵渎的肖像！你们告诉我，世界的创造者产生了\OriginalTerm{独生子}{Only-begotten}的肖像，而你们又希望这创造者被认为是\OriginalTerm{万有之父}{Father of all}的\OriginalTerm{理智}{Nous}（即心智），并且坚持这肖像对自己无知，对受造界无知，对母亲无知，对一切存在之物以及由它所造之物无知；而你们在反对自己的同时，竟将无知归给独生子本身，难道不感到羞耻吗？因为如果这些[下界]事物是由救主按照上界事物的肖像所造，而那按此肖像所造的他（德谬哥）竟处于如此大的无知中，那么由此必然得出：在他的周围，并按照他（即这位无知者是被按照他的肖像所形成的那一位），那种无知在灵性上存在。因为不可能——两者都是灵性产生的，既非塑造也非构造——在有些中肖像得以保存，而在另一些中肖像的像受到破坏，这在此产生的肖像本来是为了与上界产生的肖像类似。但如果它不相似，那么过失就在于救主，它产生了一个不同的肖像——即可以说是一个无能的工匠。因为他们无法断言救主没有生产能力，因为他们称他为\Emph{\OriginalTerm{万有}{All Things}}。那么，如果肖像不同，他就是拙劣的工匠，根据他们的假说，责任在救主。另一方面，如果相似，那么同样的无知将在他们\OriginalTerm{太初之父}{Propator}的\OriginalTerm{理智}{Nous}中，也就是在\OriginalTerm{独生子}{Only-begotten}中被发现。\OriginalTerm{太初之父}{Propator}的\OriginalTerm{理智}{Nous}，那样的话，对自身无知；对太初之父无知；而且对由祂所形成的那些事物无知。但如果你蒙\Emph{祂}知识，那么必然也得出：由救主按祂肖像所形成的那一位，应该知道相似的事物；这样，按照他们自己的原则，他们可憎的亵渎就被推翻了。\par

3. 此外，属于受造界的万物——它们种类繁多、千变万化、不可胜数——如何能是那三十个在\OriginalTerm{圆满}{Pleroma}内的\OriginalTerm{永世}{Æons}的肖像？这些人所规定的那些永世的名字，我已在前面一卷书中陈述过。他们不仅无法将整个受造界的广阔多样性适应于其圆满的相对狭小，甚至无法将任何一部分——无论是[属于]天上的、地上的，还是水中的——与之对应。因为他们自己作证，说他们的圆满由三十个永世构成；但任何人都会指出，在受造物中，仅一个门类里，他们算出的不是三十种，而是成千上万种。那么，构成如此多样性的受造物——它们彼此本性对立、互相冲突、彼此毁灭——如何能是那三十个圆满的永世的肖像和模样呢？如果他们宣称这些永世具有同一本性，拥有相等相似的性质，毫无差异？如果这些事物是那些永世的肖像，既然他们宣告有些人生来邪恶，另一些人则生来良善，那么他们也就必须指出那些永世中也有同样的差异，并坚持有些永世天生良善，有些天生邪恶，好使事物相似性的假设与永世相合。再者，世界上有些受造物温和，有些凶猛；有些无害，有些则有害并毁灭其余；有些居住在地上，有些在水中，有些在空中，有些在天上；同样，他们也必须证明永世拥有这样的属性——如果真是彼此为肖像的话。此外，\BibleQuote{永火，就是圣父为魔鬼和他的使者所预备的}\BibleRef{玛 25:41}——他们应当指出这是上面哪个永世的肖像；因为火也被视为受造界的一部分。\par

4. 然而，如果他们说这些事物是那个陷于情欲的\OriginalTerm{永世}{Æon}的\OriginalTerm{思念}{Enthymesis}的肖像，那么首先，他们便是亵渎自己的母亲，宣称她是邪恶和朽坏肖像的第一原因。再者，那些繁多、不同且本性相悖的事物，如何能是同一个存有的肖像？如果他们说\OriginalTerm{圆满}{Pleroma}的天使众多，而众多的事物是这些天使的肖像——即便如此，他们的解释也不令人满意。因为，首先，他们就必须指出圆满的天使之间彼此对立的不同之处，正如下面的肖像是彼此矛盾的本性一样。再者，既然有众多、甚至无数天使围绕在造物主周围，正如众先知所承认的——\Quote{万万千千侍立在祂面前，无数万万侍奉祂}——那么，按照他们的说法，圆满的天使便以造物主的天使为肖像，而整个受造界仍处于圆满的肖像之中，但这样一来，三十个永世就不再对应于受造界的多样变化了。\par

5. 更进一步，如果这些[下面的]事物是照着那些[上面的]事物相似造成的，那么那些事物又是照着谁的相似造成的呢？因为如果世界的造物主不是直接源于自己的思想而创造这些事物，而是像一个无能的建筑师或初学的小童，照着他人提供的原型摹仿，那么他们的\OriginalTerm{深渊}{Bythus}从何处获得那起初受造界的样式呢？显然，他必须从一位高于他的存有接受模型，而那位又从另一位接受。如此，若我们不定睛于唯一创造者和唯一天主——祂凭自身形成了所有受造之物——那么关于肖像的谈论，就像关于众神的谈论一样，将无穷无尽。难道在对待普通人的时候，我们允许他们靠自己发明生活所需之物，却不允许那位创造世界的天主靠自己创造那些被造之物的样式，并赋予其有序的安排吗？\par

6. 但是，这些[下面的]事物如何能是那些[上面的]事物的肖像呢？既然它们实际上彼此对立，毫无相通之处。因为相互对立的事物固然会毁灭与之对立的事物，却绝不能成为它们的肖像——例如，水与火；又如，光明与黑暗，以及其他类似的事物，永远不能互为肖像。同样，那些可朽坏、属地上、复合且短暂的事物，也不能是那些人所谓属灵之物的肖像；除非这些事物本身也被认为是复合的、受空间限制、有确定形状的，因而不再是属灵的、弥漫的、扩散至广阔范围且不可理解的。因为若要成为真实的肖像，它们必须具有确定的形状，并被局限在一定界限内；那么它们就不是属灵的。然而，如果这些人坚持认为那些事物是属灵的、弥漫的、不可理解的，那么那些有形状、受限制的事物，如何能成为无形状、不可理解之物的肖像呢？\par

7. 此外，若他们声称，上面的事物既非按形状也非按样式，而是按数目和产生的顺序成为下面这些事物的影像，那么首先，下面这些事物就不应被称为上面那些\OriginalTerm{永世}{Æons}的影像和相似。因为那些既没有上面事物的样式又没有形状的东西，如何能成为它们的影像呢？其次，他们应当调整上面永世的数目和产生，使它们与属于受造界者相同且相似。但现在，既然他们只提及三十个永世，并宣称受造界所包含的无数事物是那三十个永世的影像，我们就有充分的理由谴责他们完全缺乏理智。\par

\SourceRecord{Translated by Alexander Roberts and William Rambaut.；Ante-Nicene Fathers；Vol. 1.；Edited by Alexander Roberts, James Donaldson, and A. Cleveland Coxe.；Buffalo, NY: Christian Literature Publishing Co.,；1885.；<http://www.newadvent.org/fathers/0103207.htm>.}

% structure: p0001 source=h1 target=section reason=page-title

\section{\WorkTitle{反异端论}第二卷第八章}

1. 如果他们又声称这些[下面的]事物是那些[上面的]事物的\OriginalTerm{阴影}{shadow}（因为其中有些人胆敢如此主张，以至于在这方面它们就是肖像），那么，他们就必须承认那些在上者拥有身体。因为那些在上者的身体会投下阴影，但属神的实体却不会，因为它们绝不能使其他事物黯淡。然而，即使我们容许他们这一点（实际上这是不可能的），即那些属神的光明的本质拥有一个阴影，而他们声称他们的母亲降入了其中；既然那些[在上者]是永恒的，并且它们所投射的阴影也永远持续，那么[下面的]这些事物也不是暂时的，而是与那些将阴影投射在它们之上的事物一同持续。另一方面，如果这些[下面的]事物是暂时的，那么那些[在上者]——这些事物是它们的阴影——也必然随之消逝；反之，如果它们持续，它们的阴影也同样持续。\par

2. 然而，如果他们主张所谈论的阴影并非由[那些在上者]的投影产生，而仅仅是在于这一点，即[下面的]事物与那些[在上者]远远分离，那么他们就会控告他们父之光的软弱和无能，仿佛它不能延伸至这些事物，反而未能充满那虚空并驱散阴影，而且没有任何人设置障碍。因为，按照他们的说法，他们父之光将变成黑暗，被埋没于幽暗之中，并在那些以空虚为特征的场所终结，因为它不能穿透并充满一切。那么，他们不要再声称他们的\OriginalTerm{深渊}{Bythus}是万有的圆满，如果它确实既未充满也未照亮那\OriginalTerm{虚空}{vacuum}和阴影；或者，另一方面，如果他们父之光的确充满一切，就让他们停止谈论\OriginalTerm{虚空}{vacuum}和\OriginalTerm{阴影}{shadow}。\par

3. 那么，在首要的父——即统御万有的天主——之外，既不可能有任何\OriginalTerm{圆满}{Pleroma}，而他们声称那遭受情欲的移涌的\OriginalTerm{思想}{Enthymesis}降入了其中（因此圆满本身，或首要的天主，不应被那在外者所限制和界定，从而实际上被它包容）；也不可能有\OriginalTerm{虚空}{vacuum}或\OriginalTerm{阴影}{shadow}的存在，因为父先在，所以他的光不会失败，也不会在\OriginalTerm{虚空}{vacuum}中终结。此外，设想一个场所，在其中那位按照他们来说是\OriginalTerm{先父}{Propator}和\OriginalTerm{起源}{Proarche}以及万有和这圆满之父的存在终止并终结，这是不合理且不虔敬的。再者，基于已经陈述的理由，也不可容许声称某个其他存有在父的怀抱中形成了如此庞大的受造界，无论得到他的同意与否。因为宣称如此庞大的受造界是由天使或某个特定的产物——不认识真天主——在他自己的领域中形成的，同样是既不虔敬又愚昧的。那些属于尘世和物质的事物不可能在他们的圆满之内形成，因为圆满完全是属神的。此外，那些属于多种形态的受造界、并以相互对立之性质形成的事物，也不可能按照上面事物的肖像被创造，因为这些（即移涌）据说是少数的、同样形态的、同质的。他们关于\Emph{\OriginalTerm{虚空}{kenoma}}——即虚空——的阴影的谈论，在各方面都已证明是虚假的。因此，他们的虚构[无论从哪个方面看]已被证明毫无根据，他们的学说站不住脚。那些听从他们的人也同样是虚空的，并且确实正在坠入丧亡的深渊。\par

\SourceRecord{Translated by Alexander Roberts and William Rambaut.；Ante-Nicene Fathers；Vol. 1.；Edited by Alexander Roberts, James Donaldson, and A. Cleveland Coxe.；Buffalo, NY: Christian Literature Publishing Co.,；1885.；<http://www.newadvent.org/fathers/0103208.htm>.}

% structure: p0001 source=h1 target=section reason=page-title

\section{驳斥异端（第二卷，第九章）}

1. 甚至连那些在许多方面说话反对祂、却又承认祂、称祂为造物主和天使的人，都接受天主是世界的创造者，更不用说所有圣经都这样呼喊，主也教导我们这位在天上的圣父，而非其他，正如我将在本著作的后续部分所表明的。但目前，从那些倡导与我们相反教义的人那里获得的证据本身就足够了——事实上，所有人都同意这一真理：古人特别从最初受造的人的传统小心地保存这一信念，同时他们颂扬一位天主、天地的创造者的赞美；另一些人，在他们之后，由天主的先知提醒这一事实，而外邦人则从创造本身学到它。因为创造本身显明了创造它的那一位，所造的工程暗示了制造它的那一位，世界显示了安排它的那一位。此外，普世教会通过全世界从宗徒们接受了这一传统。\par

2. 这位天主，正如我所说，被承认，并且从所有人那里得到祂存在的见证，那么他们虚构出来的那位圣父无疑是站不住脚的，并且没有证人[证明祂的存在]。\Person{西满·玛古斯}是第一个说他自己就是超越万有的天主，并且世界是由他的天使们形成的。然后那些继承他的人，正如我在第一卷中所表明的，通过他们各自的观点，更严重地败坏了他的教导，用他们反对造物主的不敬神和不虔敬的教义。这里提到的这些异端者，是那些人的门徒，使得认同他们的人比外邦人更坏。因为前者\Quote{事奉受造物而不事奉造物主}\BibleRef{罗 1:25}，并且\Quote{那些不是神的}\BibleRef{迦 4:8}，尽管他们将神性中的首位归于那位制造这个宇宙的天主。但后者坚持那位（即这个世界的造物主）是缺陷的产物，并描述祂具有动物性的本性，并且不知道那高于祂的大能，同时祂还宣称：\Quote{我是天主，除我以外没有别的神}\BibleRef{依 46:9}。他们断言祂说谎，他们自己就是说谎者，将各种邪恶归给祂；并且设想一个实际并不存在的、高于这一位的神，因此他们自己的观点证明了他们亵渎那真实存在的神，同时他们虚构一个不存在的神来定自己的罪。因此，那些自称\Quote{成全的}，并拥有万有知识的人，被发现比外邦人更坏，并且对他们的造物主持有更亵渎的观点。\par

\SourceRecord{Translated by Alexander Roberts and William Rambaut.；Ante-Nicene Fathers；Vol. 1.；Edited by Alexander Roberts, James Donaldson, and A. Cleveland Coxe.；Buffalo, NY: Christian Literature Publishing Co.,；1885.；<http://www.newadvent.org/fathers/0103209.htm>.}

% structure: p0001 source=h1 target=section reason=page-title

\section{驳斥异端（第二卷，第十章）}

因此，这是极其不合理的：我们竟然不理会那真正是天主、且从万有获得见证的天主，而去探究在祂之上是否还有那另一位[存在者]，他实际上并不存在，也从未被任何人宣告过。因为关于祂并没有任何清楚的言论，他们自己就提供了证据；由于他们以可怜的成功，将他们所设想的那个存在者套用到圣经中的比喻上——这些比喻无论以什么形式说出，都是为了[这个目的]而被寻求的——显然他们现在是生出了另一位神，一位从未被寻求过的神。因为他们这样努力解释圣经中模糊的段落（这些段落之所以模糊，并非指涉及另一位神，而是指涉及真天主的安排），他们就构造了另一位神，如我先前所说，编织沙绳，并将更重要的问题依附于较不重要的问题。因为任何问题都不能通过另一个本身有待解决的问题来解决；而且，在有见识的人看来，歧义不能靠另一个歧义来解释，谜也不能靠另一个更大的谜来解释；但这类事物是从那些明显、一致且清晰的事物中得到解决的。\par

2. 但这些异端者，在努力解释圣经段落和比喻的同时，提出了另一个更重要且确实不虔敬的问题，大意是：\Quote{是否在创造世界的天主之上真的另有另一位神？}他们并不是在解决他们所提出的问题；因为他们怎么能找到解决的方法呢？但他们将一个重要的问题附加在一个次要的问题上，从而将无法解决的困难插入他们的思辨中。为了让他们认识\Quote{知识}本身——却不了解这一事实：主在三十岁时来接受真理的圣洗——他们不虔敬地藐视那位创造者天主，正是祂派遣了祂（主）为人的救恩。而且，为了让人以为他们能告知我们物质的本质从何而来，他们不相信天主按祂的喜悦，运用祂自己的意志和能力，从不存在的事物中形成了万有（使现存的事物得以存在），从而收集了众多空谈。他们因此真正显露出了他们的不信；他们不相信真正存在的事物，却堕落于实际上不存在的事物的信仰中。\par

3. 因为，当他们告诉我们，一切湿润的物质来自\Person{阿卡摩特}的眼泪，一切明亮的物质来自她的微笑，一切固态的物质来自她的忧伤，一切流动的物质来自她的恐惧，并且因此他们拥有崇高的知识，借此他们优于他人——这些事如何能不被视为可鄙且真正可笑的呢？他们不相信天主（全能的、富于一切资源）创造了物质本身，因为他们不知道一个精神性和神性的本质能成就多少。但他们却相信他们的母亲——他们称之为女性所生的女性——从上述她的情感中产生了如此庞大的受造物质。他们也询问，受造物的物质是从何处供给创造者的；但他们却不询问，他们的母亲（他们称她为迷失的永世的\OriginalTerm{恩提米西斯}{Enthymesis}和冲动）从哪里获得如此大量的泪水、汗水、忧伤，或产生其余物质的东西。\par

4. 因为，将受造物的物质归于万有之天主的能力和意志，是既可信又可接受的。这也合乎理性，并且关于这样的信仰，可以恰当地说：\BibleQuote{在人不可能的事，在天主却是可能的。}\BibleRef{路 18:27}事实上，人不能从无中造出任何东西，只能从已有的物质中造出；然而天主在这点上远比人优越，因为祂自己将祂创造物的物质召唤到存在中，而此前它并不存在。但声称物质是从迷失的永世的\OriginalTerm{恩提米西斯}{Enthymesis}产生的，那永世与她的恩提米西斯远远分离，而且她的情感和感觉脱离她自身而成为物质——这是不可信的、荒谬的、不可能的、站不住脚的。\par

\SourceRecord{Translated by Alexander Roberts and William Rambaut.；Ante-Nicene Fathers；Vol. 1.；Edited by Alexander Roberts, James Donaldson, and A. Cleveland Coxe.；Buffalo, NY: Christian Literature Publishing Co.,；1885.；<http://www.newadvent.org/fathers/0103210.htm>.}

% structure: p0001 source=h1 target=section reason=page-title

\section{\WorkTitle{Against Heresies (Book II, Chapter 11)}}

1. 他们不相信那万有之上的天主，曾藉其圣言，在自己的领域内，按祂所喜悦的，创造了现有各种不同的受造之物，因为祂是万物的塑造者，如同智慧的建筑师，又像最有权能的君主。反之，他们相信是天使，或是某个与天主分离、对祂一无所知的能力，形成了这个宇宙。因此，他们不归依真理，反而沉溺于虚谎，丧失了真实生命的食粮，陷入了空虚与阴影的深渊。他们就像伊索寓言中的狗，丢掉了面包，却试图攫取它的影子，从而失去了真正的食物。从主亲口的话中很容易证明，祂承认那由法律和先知所宣讲的、创造世界并塑造人的独一圣父与造物主，祂不认识别的神，且这一位实在是万有之上的天主；祂教导说，那属于圣父的义子名分（即永生）是藉着祂自己而赐予的，祂将这永生赐给所有义人。\par

2. 但由于这些人喜欢攻击我们，并以他们真正的吹毛求疵者身份，用一些确实与我们无关的论点来反对我们，向我们提出许多比喻和刁钻的问题，我认为最好反过来，首先就他们自己的学说向他们提出以下质询，以显示其不合理性，并终止他们的狂妄。在此之后，我计划引述主的言论，使他们不仅失去攻击我们的手段，而且因无法合理回应所提出的问题，看到他们的论证体系被摧毁；这样，他们要么回归真理，谦卑自下，停止他们五花八门的幻想，为他们所出的亵渎之言向天主求情，并获得救恩；要么，如果他们仍执迷于那占据他们思想的虚荣体系，至少他们必须改变反对我们的那种论调。\par

\SourceRecord{Translated by Alexander Roberts and William Rambaut.；Ante-Nicene Fathers；Vol. 1.；Edited by Alexander Roberts, James Donaldson, and A. Cleveland Coxe.；Buffalo, NY: Christian Literature Publishing Co.,；1885.；<http://www.newadvent.org/fathers/0103211.htm>.}

% structure: p0001 source=h1 target=section reason=page-title

\section{\WorkTitle{驳斥异端（卷二，第十二章）}}

1. 首先，关于他们的\OriginalTerm{三十}{Triacontad}，我们可以指出，它全部奇妙地坍塌于两面，即就欠缺与过多而言。他们说，主在三十岁时受洗，以此作为暗示。但这一论断实际上明显推翻了他们的全部论证。就欠缺而言，情况如下：首先，因为他们将\OriginalTerm{原父}{Propatōr}算在其他的\OriginalTerm{永世}{Æons}之中。因为万有之父不应与其他受生之物同列；祂不是受生者，而是非受生者；祂无人能理解，是那被祂理解者的理解者，因此祂本身是不可理解的；祂无形无像，却与有确定形状的同列。因为祂超越其余一切，故不应与他们同列，更何况那不动情、不犯错的一位，竟与一个会动情、实际犯错的永世同列。因为我在紧接此卷的前一卷中已经表明，从\OriginalTerm{深渊}{Bythus}算起，他们一直数到\OriginalTerm{索菲亚}{Sophia}——他们描述为犯错之永世——共三十；我也在那里列出了他们的永世名称；但若不将原父算在内，按照他们自己的说法，就不再是三十个永世受生，而只有二十九个了。\par

2. 其次，关于第一个受生物\OriginalTerm{意念}{Ennœa}——他们也称其为\OriginalTerm{沉默}{Sige}——他们说\OriginalTerm{理智}{Nous}和\OriginalTerm{真理}{Aletheia}又由其发出；他们在两个方面都错了。因为不可能有任何人的意念或沉默被理解为离开他自身，且超出他而被发出，并拥有一个特殊的自身形状。但如果他们断言意念不是超出原父而发出，而是与原父合一，那么为何又将她与其他永世——那些不与父合一、因而不知父伟大者——同列？如果她是与父合一的（让我们也考虑这一点），则必然从这合一不可分的结合——构成一个存有——产生一个不可分合一的产物，使之与发出者相似。但若如此，正如深渊与沉默，理智与真理也将构成同一存有，永远互相依附。正如一方不能脱离另一方而被设想，正如水不能脱离湿润、火不能脱离热、石头不能脱离坚硬（因为这些事物互相联结，一方不能与另一方分离，而是始终共存），同样，深渊应与意念、理智与真理如此结合。\OriginalTerm{圣言}{Logos}和\OriginalTerm{生命}{Zoe}又由那些合一者发出，它们自身也应合一，构成一个存有。依照这样的推理过程，\OriginalTerm{人}{Homo}和\OriginalTerm{教会}{Ecclesia}，以及所有其余永世的对偶，都应合一，永远共存。因为按照他们的意见，必有一女性永世与一男性永世并存，因为她可谓是那男性情感的发出。\par

3. 事既如此，且他们宣扬这些观点，却又不知羞耻地教导说，十二系列中较年幼的永世——他们称其为\OriginalTerm{索菲亚}{Sophia}——远离与她的配偶——他们称其为\OriginalTerm{愿望}{Theletus}——的结合，遭受了动情，并单独、无其助而生育了一个产物，他们称之为\Quote{由女性所生的女性}。他们如此陷入狂乱，竟对同一要点形成两个截然相反的见解。因为若深渊永远与沉默合一，理智与真理、圣言与生命等等，索菲亚怎可能不与她的配偶结合而遭受或产生任何事物？再者，若她真的离他而遭受动情，则其他对偶之间也必然容许分裂与分离——我已证明这不可能。因此索菲亚不可能离开愿望而遭受动情；如此，他们的整个体系又被推翻。因为他们又从那动情——他们声称她远离与配偶的结合而经历——引出了所有其余物质实体，如同悲剧的组成一般。\par

4. 然而，如果他们厚颜地坚持——为挽救他们虚妄的想象免于毁灭——其他所有对偶也因这最后对偶的缘故而彼此分离断裂，那么我回答说：第一，他们立足于不可能之事。因为怎能将原父与其意念分离，或将理智与真理、圣言与生命分离，等等？他们又如何能自称这些对偶又趋向合一，且事实上全都合一，如果这些在\OriginalTerm{圆满}{Pleroma}之内的对偶并不保持合一，反而彼此分离；以至于到如此地步，它们不但遭受动情，而且彼此不结合便行生育之工，就像母鸡不与公鸡交配一样。\par

5. 其次，他们的首生\OriginalTerm{八元}{Ogdoad}亦将以下列方式被推翻：他们必须承认，\OriginalTerm{深渊}{Bythus}与\OriginalTerm{缄默}{Sige}、\OriginalTerm{理智}{Nous}与\OriginalTerm{真理}{Aletheia}、\OriginalTerm{圣言}{Logos}与\OriginalTerm{生命}{Zoe}、\OriginalTerm{人}{Anthropos}与\OriginalTerm{教会}{Ecclesia}各自同居同一\OriginalTerm{圆满}{Pleroma}中。但\OriginalTerm{缄默}{Sige}不可能存在于\OriginalTerm{圣言}{Logos}面前，反之，\OriginalTerm{圣言}{Logos}亦不可能在\OriginalTerm{缄默}{Sige}面前彰显自己。因为这两者互相摧毁，正如光明与黑暗决不可能共处一处：若光明得胜，便无黑暗；若黑暗，便无光明，因为光明出现，黑暗便被驱逐。同样，有\OriginalTerm{缄默}{Sige}之处，便无\OriginalTerm{圣言}{Logos}；有\OriginalTerm{圣言}{Logos}之处，定无\OriginalTerm{缄默}{Sige}。但若他们说\OriginalTerm{圣言}{Logos}仅内在于（未表达），则\OriginalTerm{缄默}{Sige}亦将内在于，且并不更少地被内在于的\OriginalTerm{圣言}{Logos}摧毁。然而，他并非只是心智中的概念，他们的\OriginalTerm{永世}{Æons}产生的顺序本身便已显明。\par

6. 那么，他们不要再宣称首生\OriginalTerm{八元}{Ogdoad}由\OriginalTerm{圣言}{Logos}与\OriginalTerm{缄默}{Sige}组成，而应（必然地）排除\OriginalTerm{缄默}{Sige}或\OriginalTerm{圣言}{Logos}；如此，他们的首生\OriginalTerm{八元}{Ogdoad}便告终结。因为若他们描述\OriginalTerm{永世}{Æons}的结合是联合的，则其整个论点便瓦解。因为若它们联合，\OriginalTerm{索菲亚}{Sophia}如何能在不与她的配偶结合的情况下生出一个缺陷？反之，若他们坚持，如同在实际产生中，每个\OriginalTerm{永世}{Æons}拥有其独特实体，那么\OriginalTerm{缄默}{Sige}与\OriginalTerm{圣言}{Logos}又如何能在同一地方彰显？至此，关于缺陷的讨论已足。\par

7. 然而，他们的\OriginalTerm{三十元}{Triacontad}又因过度而被以下考虑推翻。他们描述\OriginalTerm{界限}{Horos}（他们以多种名号称之，我在前书中已提及）是由\OriginalTerm{独生者}{Monogenes}产生，如同其他\OriginalTerm{永世}{Æons}一样。他们中有些人主张此\OriginalTerm{界限}{Horos}由\OriginalTerm{独生者}{Monogenes}产生，另些人则确认祂是由\OriginalTerm{先祖}{Propator}按自己的肖像所发出。他们进一步确认，由\OriginalTerm{独生者}{Monogenes}产生了一个产出——\OriginalTerm{基督}{Christ}与\OriginalTerm{圣神}{Holy Spirit}；他们并未将这些算在\OriginalTerm{圆满}{Pleroma}之数中，也未将\OriginalTerm{救主}{Saviour}算在内，而他们亦宣告祂为\Emph{Totum}（万有）。现在，即使瞎子也清楚，不仅如他们所主张的三十种产出被发出，而是另有四种加在这三十种之上。因为他们将\OriginalTerm{先祖}{Propator}本身算在\OriginalTerm{圆满}{Pleroma}内，亦将那些依次由彼此所生者算在内。那么，为何那些（其他存在）不与被算在同一个\OriginalTerm{圆满}{Pleroma}中，既然它们以同样方式产生？他们能有何正当理由，不将\OriginalTerm{基督}{Christ}（他们描述为按父旨由\OriginalTerm{独生者}{Monogenes}所生）、或\OriginalTerm{圣神}{Holy Spirit}、或\OriginalTerm{界限}{Horos}（他们亦称为\OriginalTerm{救主}{Soter}）、甚至\OriginalTerm{救主}{Saviour}本身（祂来协助并赋形于他们的母亲）与其他\OriginalTerm{永世}{Æons}一同计算？这岂不是好像这些后者比前者更弱，因而配不上\OriginalTerm{永世}{Æons}之名或列入其数？抑或它们更为优越卓越？但如何可能更弱，既然它们是为建立并纠正别的而产？再者，它们不可能优于首生的\OriginalTerm{四元}{Tetrad}，因它们亦由其产生；因为\OriginalTerm{四元}{Tetrad}本身亦被列于上述数目中。因此，这些后者也应被计入\OriginalTerm{永世}{Æons}的\OriginalTerm{圆满}{Pleroma}中，否则那些享有此称谓（\OriginalTerm{四元}{Tetrad}）的\OriginalTerm{永世}{Æons}应被剥夺其尊荣。\par

8. 因此，既然他们的\OriginalTerm{三十元}{Triacontad}如此归于虚无，如我所显示的，无论从缺陷或过度而言（因为处理这样一个数目，任何程度上的过多或过少都会使这个数目难以成立，何况如此大的变种？），则他们关于\OriginalTerm{八元}{Ogdoad}与\OriginalTerm{十二元}{Duodecad}的所言纯属无稽之谈，无法成立。他们的整个体系亦倒塌，当其根基被摧毁并消解于\OriginalTerm{深渊}{Bythus}，即消解于虚无之中。让他们从此寻求其他理由，解释为何主在三十岁时受洗，以及宗徒的\OriginalTerm{十二元}{Duodecad}，以及关于患血漏的妇人所陈述的事实，以及其他诸点，他们在这些事上徒然疯狂劳碌。\par

\SourceRecord{Translated by Alexander Roberts and William Rambaut.；Ante-Nicene Fathers；Vol. 1.；Edited by Alexander Roberts, James Donaldson, and A. Cleveland Coxe.；Buffalo, NY: Christian Literature Publishing Co.,；1885.；<http://www.newadvent.org/fathers/0103212.htm>.}

% structure: p0001 source=h1 target=section reason=page-title

\section{\WorkTitle{驳斥异端}（卷二，第十三章）}

1. 我现在着手证明，他们构想的第一个产生次序必须被摒弃。因为他们主张，\OriginalTerm{理智}{Nous}和\OriginalTerm{真理}{Aletheia}是从\OriginalTerm{深渊}{Bythus}及其\OriginalTerm{意念}{Ennœa}产生的，而这被证明是矛盾的。因为理智本身就是首要、最高，且简直是所有理解的原则和源泉。再者，从它生出的意念，则是关于任何主题的某种情感。因此，理智不可能由深渊和意念产生；他们更应主张意念是作为\OriginalTerm{元父}{Propator}和这个理智的女儿而产生的，这更接近真理。因为意念不是理智的女儿，如他们所断言，而是理智成为意念的父亲。因为当理智占据了那存在于元父内隐秘不可见情感的首要和首要位置时，它如何能被元父产生呢？藉由这情感，感觉、意念和\OriginalTerm{深思}{Enthymesis}以及其他事物被产生，它们都只是理智本身的同义词。正如我已说过的，它们仅仅是那同一个能力关于某一特定主题的某些确定的思维活动。我们理解这些〔若干〕术语是根据它们意义的长度和广度，而不是根据任何〔根本〕的〔变化〕；而这些〔各种思维活动〕被〔同一〕知识的范围所限制，并由〔同一〕术语一起表达，这〔完全相同〕的感觉留在内部，并创造、管理，并凭自身能力自由地支配那些先前提到的事物，如其所愿。\par

2. 因为那〔能力〕关于任何事物的首次活动，被称为\OriginalTerm{意念}{Ennœa}；但当它持续、增强并占据整个灵魂时，被称为\OriginalTerm{深思}{Enthymesis}。这深思，当它长时间在同一问题上活动，并如同经过考验时，称为\OriginalTerm{感觉}{Sensation}。这感觉，当它大大发展时，成为\OriginalTerm{谋略}{Counsel}。谋略的进一步增长和极大发展，成为\OriginalTerm{思维审断}{Examination of thought}（判断）；而它存留于心灵中时，最恰当地称为\OriginalTerm{逻各斯}{Logos}（理性），而从它产生出言语的逻各斯（话语）。但所有已提及的〔思维活动〕根本上是一体的，都从理智获得起源，并根据其增长获得不同名称。正如人的身体，有时年轻，有时壮年，有时老年，根据增长和持续而获得不同名称，但不是由于实体变化，也不是由于任何身体真实损失；那些〔思维活动〕也是一样。因为当人〔在思想上〕注视任何事物时，他也思考它；当他思考它时，他也拥有关于它的知识；当他认识它时，他也考虑它；当他考虑它时，他也〔在思想〕操作它；当他操作它时，他也说出它。但是，正如我已经说过的，是理智统治所有这些〔思维过程〕，祂本身是不可见的，并通过那些已提及的过程，如同通过〔发自祂的〕光芒来说话，但祂自身并不由任何其他事物发出。\par

3. 这些说法可说适用于人，因为人本性是复合的，由身体和灵魂组成。但那些断言\OriginalTerm{意念}{Ennœa}从天主发出，\OriginalTerm{理智}{Nous}从意念发出，然后依次，\OriginalTerm{逻各斯}{Logos}从这些发出的人，首先应受指责，因为他们不当使用这些产物；其次，他们描述人的情感、激情和心灵倾向，却〔由此证明自己〕对天主无知。借他们的说话方式，他们将适用于人的事物归于万有之父（他们又宣称祂是所有人不可知的），并否认祂亲自创造了世界，以免将无能归于祂；同时却将人的情感和激情赋予祂。但如果他们认识圣经，并受真理教导，他们必无疑会知道，天主不像人，祂的思想也不像人的思想。\BibleRef{依 55:8}因为万有之父远超越那些在人类中运作的情感和激情。祂是单纯、非复合的存在，没有多样部分，全然相似且等同自身，因为祂完全是领悟、完全是灵、完全是思想、完全是理智、完全是理性、完全是听见、完全是看见、完全是光，完全是全部美善的源头——正如虔诚和敬神的人惯于谈论天主的。\par

4. 然而，祂超越〔所有〕这些属性，因此不可言说。因为祂完全可以并恰当地被称为领悟万物的理解，但祂并不〔因此〕像人的理解；祂最恰当地被称为光，但祂毫不像我们所熟知的光。同样，在其他所有方面，万有之父丝毫不似人的软弱。祂按〔我们对祂的〕爱而被这样描述；但就伟大而言，我们对祂的思想超越这些表述。如果即使在人的情况，理解本身并不从发出产生，产生其他事物的理智也不与活人分离，而它的运动和情感却显现出来，更何况天主的理智（祂是全部的理解）绝不会与自己分离；〔在祂的情形〕，也没有任何事物能如同从一个不同存在者那样产生出来。\par

5. 因为如果祂产生了理智，那么按照他们的观点，那产生理智的一位必须被理解为复合的和有形体的存有；这样，发出那理智的天主与理智是分离的，而那被发出的理智也与祂分离。但若他们断言理智从理智发出，那么他们就将天主的理智割裂，并分成了部分。它到哪里去了？是从哪里被发出的？因为任何从一个地方发出的东西，必然进入另一个地方。但有什么存在比天主的理智更古老，他们竟声称理智被发到那里？那必须是一个多么广阔的区域，能够接收和容纳天主的理智！然而，他们若肯定这发射如同光线从太阳射出，那么，正如接收光线的下层空气必须先于光线存在，同样，他们也将指出有某种存在，天主的理智被发射到其中，能够容纳它，并且比它自身更古老。由此，我们必须认为，正如我们看到太阳——比万物都小——从自身发出光线到很远的地方，同样我们也说\OriginalTerm{原始根源}{Propator}从自身发出光线到远方，远离自己。但有什么能被认为在远方或远离天主，祂将这道光线发射进去呢？\par

6. 再者，若他们断言那理智并非被发到圣父之外，而是在圣父自身之内，那么，首先，说它被发出就是多余的。因为如果它仍留在圣父内，又如何能被发出呢？因为发射是所发射之物在发射者之外的显现。其次，这理智被发出，那么那从祂生出的\OriginalTerm{逻各斯}{Logos}仍将在圣父内，将来从\OriginalTerm{逻各斯}{Logos}发出的发射也将如此。那么，它们就不能在这种情况下对圣父无知，因为它们都在祂内；而且，既然都被圣父等同地包围，就没有哪一个能因发射的降序而更少认识祂。它们也必同样保持不受感受，因为它们存在于圣父的怀中，没有一个会陷入败坏或降格的状态。因为在圣父内没有腐败，除非可能如同大圆中包含小圆，小圆中再包含更小的圆；或者除非他们论圣父，说祂以球体或方体的方式，从各方面在自身内包含着球体的相似，或以方体的形式产生其余\OriginalTerm{永世}{Æon}，每一个都被比自己更大的那一位所包围，而又包围比自己更小的那一位；因此，那最小的和最末的，位于中心，这样远离圣父，确实对\OriginalTerm{原始根源}{Propator}无知。但如果他们坚持这种假说，他们就必须将他们的\OriginalTerm{深渊}{Bythus}封闭在确定的形式和空间内，而祂既包围其他，又被它们包围；因为他们必然承认有某种东西在祂之外包围祂。并且关于包围者和被包围者的谈论将无限延伸；所有\OriginalTerm{永世}{Æon}都将清楚显现为彼此包容的形体。\par

7. 此外，他们还必须承认，要么祂只是虚空，要么整个宇宙都在祂内；那样的话，所有都将同等地分有圣父。正如，若有人在水中画圆，或画圆形或方形，所有这些都将同等分有水；同样，那些形成于空气中的必然分有空气，那些形成于光中的分有光；同样，那些在祂内的也必将同等分有圣父，无知在他们中没有位置。那么，这充满万有的圣父的分有在哪里呢？诚然，若祂充满了万物，他们中间就不会有无知。因此，他们所设想的败坏的工作就归于无有，以及物质的产生和世界其余部分的形成；他们声称这些东西来自激情和无知。反之，若他们承认祂是虚空，那么他们就陷入最大的亵渎；他们否认祂的属神本性。因为祂怎能是属神的存在，若祂甚至不能充满那些在祂内的事物呢？\par

8. 如今，关于\OriginalTerm{理智}{Nous}流溢所说的话，同样适用于反对巴西里得学派的人，以及反对其余诺斯替派的人，这些（瓦伦廷派）从后者接受了关于流溢的观念，并在第一卷中已被驳斥。但是我现在已经清楚地表明，他们所说的理智的首次产生，是一种站不住脚、不可能的观点。且让我们看看关于其余（伊雍）的情况如何。因为他们声称，\OriginalTerm{圣言}{Logos}和\OriginalTerm{生命}{Zoe}由他（即理智）发出，作为这个\OriginalTerm{圆满}{Pleroma}的塑造者；同时他们以人类情感的类比来构想圣言的流溢，即圣言，并鲁莽地对天主形成猜想，仿佛他们在断言圣言由理智产生时发现了什么奇妙之事。所有人都清楚感知到，这在人类方面可以合乎逻辑地肯定。但在那位高于万有的天主内，既然祂全是理智，全是圣言，如同我先前说过的，祂内在没有任何比其他更古老或更晚近的，也没有彼此不一致的，而是完全平等、相似、同质的，那么就不再有任何理由按照上述次序构想这样的产生。正如那个人没有错，他宣称天主全是视觉，全是听觉（因为祂以何种方式看见，也同样以何种方式听见；祂以何种方式听见，也同样以何种方式看见），同样，谁肯定祂全是理智，全是圣言，并且无论从哪方面看祂是理智，从哪方面看祂也是圣言，而这理智就是祂的圣言，那么他确实仍然只对万有之父有不充分的观念，但要比那些把人类发出的话语的产生转移到天主的永恒圣言上、为（圣言）指定一个开始和产生过程的人，怀有远为合宜的（关于祂的思想）。如果天主圣言——甚至天主自己，因为祂就是圣言——遵循与人的言语相同的次序和产生过程，那么祂与人的言语又有什么区别呢？\par

9. 他们也在\OriginalTerm{生命}{Zoe}上陷入了错误，因为他们主张她是在第六个位置被产生的，而她理应居于所有（其余）之先，因为天主是生命、不朽坏和真理。这些以及诸如此类的属性并不是按照逐步的下降等级产生的，而是那些永远存在于天主内的完美之名的名称，是按照人可能且适宜地听到和谈论天主的方式。因为与天主的名称和谐的有以下词语：\OriginalTerm{理智}{Nous}、圣言、生命、不朽坏、真理、智慧、良善，以及诸如此类。没有人能够主张理智比生命更古老，因为理智本身即是生命；也不能主张生命比理智更晚，以致万有的理智，即天主，曾在某一时刻缺乏生命。但如果他们肯定生命确实（先前）就在圣父内，但为了使\OriginalTerm{圣言}{Logos}能生活，而在第六个位置被产生，那么按照这种推理，生命理应早已在第四位被发出，以使理智拥有生命；而且更进一步，甚至在理智之先，它应与\OriginalTerm{深渊}{Bythus}同在，以使他们的深渊拥有生命。因为把\OriginalTerm{缄默}{Sige}同他们的\OriginalTerm{元父}{Propator}算在一起，并把她配给祂作为祂的配偶，却不同时将生命列入数目——这岂非超越所有其他疯狂？\par

10. 再者，关于从这些（已经提到的伊雍）发出的第二个产生——即\OriginalTerm{人}{Homo}和\OriginalTerm{教会}{Ecclesia}——他们的父辈，即那些虚假地自称诺斯替派的人，彼此争论，每个人都试图维护自己的见解，从而证明自己是邪恶的窃贼。他们主张，对于产生（理论）来说——事实上，这似乎是真实的——\OriginalTerm{圣言}{Logos}是由人产生的，而不是人由圣言产生的；并且人存在于圣言之先，而这位人其实就是万有之上的天主。因此，正如我先前指出的，他们把人类的一切情感、心理活动、意向的形成、言语的表达，以某种貌似合理的方式堆砌在一起，却毫无根据地说谎反对天主。因为当他们把发生在人身上的事，以及他们自己所体验到的一切，都归之于神圣的理性时，他们在那些不认识天主的人看来，似乎说得头头是道。借着这些人类的情欲，他们误导自己的理智，当他们描述天主圣言在第五位的起源和产生时，他们断言如此就教导了奇妙的奥秘，是不可言说、崇高无比的，除他们以外无人知晓。（他们声称，）主所说的正是关于这些：\BibleQuote{你们求，就得着}\BibleRef{玛 7:7}，也就是说，他们应当探究：\OriginalTerm{理智}{Nous}和\OriginalTerm{真理}{Aletheia}如何从\OriginalTerm{深渊}{Bythus}和\OriginalTerm{缄默}{Sige}中发出；圣言和生命是否又从这些再次衍生；然后，\OriginalTerm{人}{Anthropos}和教会是否从圣言和生命发出。\par

\SourceRecord{Translated by Alexander Roberts and William Rambaut.；Ante-Nicene Fathers；Vol. 1.；Edited by Alexander Roberts, James Donaldson, and A. Cleveland Coxe.；Buffalo, NY: Christian Literature Publishing Co.,；1885.；<http://www.newadvent.org/fathers/0103213.htm>.}

% structure: p0001 source=h1 target=section reason=page-title

\section{\WorkTitle{驳斥异端（第二卷，第十四章）}}

1. 古代喜剧诗人之一\Person{安提法奈斯}在其\WorkTitle{神谱}中关于万物起源的描述，远比这更接近真理，也更令人愉悦。因为他讲述从黑夜和沉默中产生了混沌；接着说爱产生了，它来自混沌和黑夜；从爱又产生了光明；根据他的看法，最原始的神族其余一切皆由此而来。此后，他接着介绍第二代神族以及世界的创造；然后他叙述了人类是由第二代神祇所造。这些人（异端者）将这则寓言采纳为自己的，并将他们的意见围绕它排列，仿佛出于某种自然过程，只改变了所指事物的名称，并以完全相同的方式阐述万物生成及其产生的开端。他们用\OriginalTerm{深渊}{Bythus}和\OriginalTerm{沉默}{Sige}取代黑夜和沉默；他们用\OriginalTerm{心灵}{Nous}取代混沌；他们提出\OriginalTerm{圣言}{Word}取代爱（喜剧诗人称，万物皆由爱安排）；他们制造了\OriginalTerm{移涌}{Æons}取代首要和最大的众神；至于次等神祇，他们告诉我们，那是他们的母亲在\OriginalTerm{圆满}{Pleroma}之外创造的，称之为第二\OriginalTerm{奥格多阿德}{Ogdoad}。他们向我们宣称，像上述作者一样，世界的创造和人的形成出于此（奥格多阿德），并坚持只有他们才知晓这些不可言喻和未知的奥秘。喜剧演员在剧场中以最清晰的声音到处表演的那些东西，他们搬到了自己的体系中，无疑是通过同样的论点教导它们，仅仅改了名字。\par

2. 这些人不仅被证实将喜剧诗人中可找到的那些东西当作自己的（原始观念）提出，而且还汇集了所有不认识天主、被称为哲学家的人所说的话；他们将这些碎布片缝合成一件五彩斑斓的衣服，凭其巧妙的表达方式，为自己披上了一件本不属于自己的外衣。的确，他们引入了一种新教义，因为这种教义已通过一种新艺术被取代（旧教义）。然而它实际上既陈旧又无用，因为这些意见是从充满无知和不虔敬的古老教义中拼凑而成的。例如，\Person{米利都的泰勒斯}断言水是万物的生成和本原原则。现在，无论我们说\Emph{水}还是\OriginalTerm{深渊}{Bythus}，都是一回事。\Person{荷马}又认为，\Person{俄刻阿诺斯}与母亲\Person{特堤斯}是众神的起源：这些人将这一观念转移到了深渊和\OriginalTerm{沉默}{Sige}。\Person{阿那克西曼德}规定无限是万物的第一原则，它内在具有万物的种子性生成，并由此宣称形成众多世界（现存者）：他们再次重新包装了这一观念，并归之于深渊和他们的\OriginalTerm{移涌}{Æons}。\Person{阿那克萨戈拉}，他又被称为\Quote{无神论者}，声称动物是从天上降到地上的种子形成的。这一思想，这些人又转移到他们母亲的\Quote{种子}，他们坚持自己就是该种子；因此，在有识者看来，他们立即承认自己就是不信神的\Person{阿那克萨戈拉}的后裔。\par

3. 再者，他们采纳了\Person{德谟克利特}和\Person{伊壁鸠鲁}的阴影和虚空（观念），并将这些应用于自己的观点，追随那些早已大谈虚空和原子的教师；他们把其中之一称为\Emph{存在}，另一个称为\Emph{非存在}。同样，这些人将\OriginalTerm{圆满}{Pleroma}之内的事物称为真实存在，正如那些哲学家看待原子一样；而他们坚持认为圆满之外的事物没有真实存在，正如那些人对待虚空一样。他们因此将自己放逐于这个世界（因为他们在此处于圆满之外）的一个没有存在的地方。再者，当他们坚持下面的事物是上面真实存在的影像时，他们又一次极其明显地重复了\Person{德谟克利特}和\Person{柏拉图}的学说。因为\Person{德谟克利特}是第一个主张无数不同的形象被印上（上方事物的）形式，并从宇宙空间降临到这世界的人。但\Person{柏拉图}则谈论质料、模型和天主。这些人追随这些区分，将柏拉图所谓的理念、模型称为那些在上之物的\Emph{影像}；而仅仅通过改名，他们就自夸是这类虚构幻想的发现者和设计者。\par

4. 他们还持有一种意见，认为创造者是从先前存在的质料中形成世界的，\Person{阿那克萨戈拉}、\Person{恩培多克勒}和\Person{柏拉图}在他们之前都表达过这种看法；而且，我们得知他们也在其母亲的感召下这样做。再者，关于万物必然归回到它们据以构成的那些东西，以及天主是这种必然性的奴隶，以至于祂不能将不朽赋予可朽者，也不能将不败坏赋予可败坏者，而是每一事物都归于与其本性相似的实体，这种意见，那些被称为廊下派的斯多葛主义者（由\OriginalTerm{廊下}{\GreekText{στοὰ}}得名），以及所有不认识天主的人，无论诗人还是历史学家，都同样认定。那些持同样不信体系的（异端者）无疑已将属灵存在归入其自身领域——即\OriginalTerm{圆满}{Pleroma}之内，而将属魂存在归入中间空间，将属肉体的归入物质方面。他们断言天主自己也无法做别的事，而是每一种（不同种类的）实体都归于与其相同本性的东西。\par

5. 此外，关于他们说救主是由所有永世所形成，每个永世都在他内，可以说是存入了自己特别的花朵，他们所提出的并无新意，无非是可以在赫西俄德的潘多拉中找到的。因为赫西俄德关于她所说的，这些人暗示关于救主，将祂呈现为我们如同潘多洛斯（万赐者），仿佛每个永世都将自己最完美的东西赐予了祂。同样，他们关于食用肉类和其他行为的漠然态度，以及他们自认为出于其高贵的本性，无论吃什么或做什么都丝毫不会沾染污秽的观点，乃是从犬儒学派学来的，因为他们实际上与这些哲学家属于同一类。他们还努力将对信德问题的处理引入那种吹毛求疵和微妙的方式，这实际上是模仿亚里士多德。\par

6. 同样，关于他们表现出的将整个宇宙归于数字的倾向，他们是从毕达哥拉斯学派学来的。因为这些人最先提出数字是万物的本原，并描述他们的本原既相等又不相等，由此他们设想，可感之物和非物质之物都源于此（两个属性）。他们认为，一组本原产生了事物的质料，另一组产生了它们的形式。他们断言，万物都由这些本原形成，正如雕像由金属和特定形状构成一样。如今，异端者将此应用于圆满之外的事物。毕达哥拉斯学派主张理智的本原与能量成比例，心智作为可理解之物的接受者，追求其探究，直到疲惫，最终消溶于不可分的一中。他们还断言，\OriginalTerm{太一}{\GreekText{Ἕν}}，即一，是万物的第一本原，是一切被形成之物的实体。由此又产生了\OriginalTerm{二元}{\GreekText{Δυάς}}、\OriginalTerm{四元}{\GreekText{Τετράς}}、\OriginalTerm{五元}{\GreekText{Πεντάς}}，以及其他众多的衍生。这些事，异端者逐字重复，指向他们的\OriginalTerm{圆满}{\GreekText{Πλήρωμα}}和\OriginalTerm{深渊}{\GreekText{Βυθός}}。同样，他们也力图将那些从统一中产生的结合引入流行。\Person{马尔库斯}以这些观点自夸，仿佛这些是他自己的，并似乎被看作发现了比他人更新颖的东西，然而他不过是把毕达哥拉斯的四元组作为万物的本原和母亲陈述出来。\par

7. 但我只对这些说一句话：——所有那些被提及的、与你们在表达上被证实一致的人，他们知道真理还是不知道？如果他们知道，那么救主降临此世便是多余的。因为祂为何降临？难道是祂要把那已经为认识真理者所知的真理带给认识真理的人吗？另一方面，如果这些人\Emph{不}知道，那么，既然你们用与那些不知道真理者相同的术语表达自己，你们又怎能夸口唯独你们拥有那超越万物的知识，尽管那些对天主无知的人也拥有它？因此，通过彻底颠倒语言，他们将无知的真理称为知识：保禄说得对，他们\Quote{虚假知识的标新立异之词}。因为他们的知识确实被证明是虚假的。然而，如果这些人厚颜无耻地就这些点声称，人们确实不知道真理，但是他们的母亲，即圣父的种子，透过这样的人（正如通过先知一样）宣告了真理的奥秘，而\OriginalTerm{造物主}{\GreekText{Δημιουργός}}对此一无所知，那么我首先回答，所预言的事情并非无法被任何人理解；因为那些人自己知道他们在说什么，他们的门徒也知道，再传的门徒也知道。其次，如果母亲或她的种子知道并宣告了属于真理的事（而圣父就是真理），那么根据他们的理论，救主说话就是虚假的，因为祂曾说：\BibleQuote{除了圣子，无人认识圣父}\BibleRef{玛 11:27}，除非他们坚持认为他们的种子或母亲是\Emph{无人}。\par

8. 至此，他们藉著【将人类情感归于其\OriginalTerm{永世}{Æon}】以及他们的话语与许多不认识天主的人高度吻合的事实，似乎颇有道理地引诱了一些人【离开真理】。他们利用那些人们熟悉的【表达】，引导他们进入那种论述万物的言谈，展示\OriginalTerm{天主圣言}{Logos}、\OriginalTerm{生命}{Zoe}、\OriginalTerm{心灵}{Nous}的产出，仿佛将神性的【连续】流溢带入世间。他们提出的观点，既无合理性也无掩饰，从头到尾纯粹是谎言。这正如那些为了引诱和捕捉某种动物的人，把它们的惯常食物放在前面，藉着熟悉的饵食逐渐引它们上钩，直到终于捉住它们；但一旦捕获，便使它们遭受最苦的束缚，并暴力地将它们拖到任何他们想去的地方。同样，这些人也藉著貌似合理的言辞，逐渐温和地劝服【他人】接受前面提到的流溢，然后提出不一致的东西，以及其余流溢的形态，并非人们所期望的那样。例如，他们宣称【十】个永世由\OriginalTerm{圣言}{Logos}和\OriginalTerm{生命}{Zoe}发出，而从\OriginalTerm{人}{Anthropos}和\OriginalTerm{教会}{Ecclesia}发出了十二个，尽管他们既无证据、无见证、无或然性、也无任何此类性质的东西【支持这些断言】；并以同样愚蠢和大胆，希望人们相信从圣言和生命（二者是永世）发出了\OriginalTerm{深渊}{Bythus}、\OriginalTerm{混合}{Mixis}、\OriginalTerm{不朽}{Ageratos}、\OriginalTerm{合一}{Henosis}、\OriginalTerm{自生}{Autophyes}、\OriginalTerm{快乐}{Hedone}、\OriginalTerm{不动}{Acinetos}、\OriginalTerm{融合}{Syncrasis}、\OriginalTerm{独生}{Monogenes}、\OriginalTerm{真福}{Macaria}。此外，【如他们所断言，】以类似方式，从人和教会（二者是永世）发出了\OriginalTerm{护慰者}{Paracletus}、\OriginalTerm{信德}{Pistis}、\OriginalTerm{父性}{Patricos}、\OriginalTerm{望德}{Elpis}、\OriginalTerm{母性}{Metricos}、\OriginalTerm{爱德}{Agape}、\OriginalTerm{赞美}{Ainos}、\OriginalTerm{理解}{Synesis}、\OriginalTerm{教会者}{Ecclesiasticus}、\OriginalTerm{真福者}{Macariotes}、\OriginalTerm{意志}{Theletos}、\OriginalTerm{智慧}{Sophia}。\par

9. 关于这位\OriginalTerm{智慧}{Sophia}的激情与错误，以及她如何因探究圣父的【本性】而几乎丧命，按他们的叙述，以及在\OriginalTerm{圆满}{Pleroma}之外发生了什么事，他们教导说世界的创造者是由何种缺陷产生的，这些我已在前面一卷中陈述过，在其中我详尽描述了这些异端者的意见。【我也详细说明了他们】关于\OriginalTerm{基督}{Christ}的观点，他们描述基督是在这一切之后产生的，以及关于\OriginalTerm{救主}{Soter}，【根据他们说，】他的存来自那些在圆满内形成的永世。但我现在不得不提及他们的名称，以便从这些名称中显明他们谎言的荒谬，以及他们所发明术语的混乱性质。因为他们自身因这类众多的名称而减损了其永世的【尊荣】。他们向外邦人提出貌似合理可信的名称，【类似】那些被称为十二神明的名称，甚至他们还要将这些神明视为其十二永世的肖像。但【所谓的】肖像能够产生【自己的】名称，这些名称通过词源更能表明神性，【比他们幻想中的原型】更为适宜和有力。\par

\SourceRecord{Translated by Alexander Roberts and William Rambaut.；Ante-Nicene Fathers；Vol. 1.；Edited by Alexander Roberts, James Donaldson, and A. Cleveland Coxe.；Buffalo, NY: Christian Literature Publishing Co.,；1885.；<http://www.newadvent.org/fathers/0103214.htm>.}

% structure: p0001 source=h1 target=section reason=page-title

\section{\WorkTitle{驳斥异端（第二卷，第十五章）}}

1. 然而，让我们回到之前关于[\OriginalTerm{永世}{Æons}]产生的那个问题。首先，请他们告诉我们，永世的产生为何是这样的：它们不与任何属于受造界的事物接触。因为他们主张，那些[在上者]并非因受造界而造，而是受造界因它们而造；前者也不是后者的影像，而是后者是前者的影像。因此，既然他们为影像提供理由，说月亮有三十天是因为三十个永世，白天有十二个小时，一年有十二个月，是因为\OriginalTerm{圆满}{Pleroma}内的十二个永世，以及其他类似的胡说，那么现在也请他们告诉我们，永世的那种产生为何具有这样的性质：为何第一个和首生的\Emph{\OriginalTerm{八元体}{Ogdoad}}被发施出来，而不是五元体，或三元体，或七元体，或其他由不同数目所限定的东西？此外，为何从\OriginalTerm{圣言和生命}{Logos and Zoe}发施出十个永世，不多不少？而从\OriginalTerm{人和教会}{Anthropos and Ecclesia}又发出十二个，而它们本来可能更多或更少？\par

2. 然后，再次关于整个\OriginalTerm{圆满}{Pleroma}，有什么理由将它分为这三个——\OriginalTerm{八元体}{Ogdoad}、\OriginalTerm{十元体}{Decad}和\OriginalTerm{十二元体}{Duodecad}——而不是分成其他不同的数目？此外，关于这个划分本身，为何分成\Emph{三}部分，而不是四、五、六，或某些与受造界数目无关的其他数目？因为他们描述那些[在上者]比这些[在下者]更古老，它们应当在自身内拥有它们[存在]的原理，这个原理在创造之前就已存在，而不是按照受造界的样式，所有点都完全一致。\par

3. 我们所给出的关于受造界的解释是与[世界中]那种规律顺序一致的，因为这个方案适应于那些[实际]被造的事物；但他们必然无法为那些[创造]之前、自身就完备的存在物提供属于事物本身的理由，从而陷入最大的困境。因为关于他们质问我们不知受造界之事的问题，当他们反过来被质问关于\OriginalTerm{圆满}{Pleroma}时，要么只提及纯粹的人类情感，要么诉诸那种只关乎受造界中可观察到的和谐的说法，不恰当地给我们回答关于次要事物，而不是关于他们声称是首要的事物。因为我们问他们的，不是关于属于受造界的和谐，也不是关于人类情感；而是因为他们必须承认，关于他们的\OriginalTerm{八形、十形和十二形的圆满}{octiform, deciform, and duodeciform Pleroma}（他们宣称受造界是它的影像），他们的父空幻而轻率地将它塑造成那个形状，并且若他没有理由地造任何东西，就必须将畸形归于他。或者，如果他们说，圆满是按照父的预见这样产生的，为了受造界，好像他这样对称地安排它的本质本身，那么圆满就不再被认为是因自身而形成，而是为了那个将要成为它的影像、拥有其相似者的[受造界]（正如粘土模型不是为了自身而塑造，而是为了即将用铜、金或银塑造的雕像），那么受造界将比圆满更受尊崇，如果为了它那些[在上者]被产生。\par

\SourceRecord{Translated by Alexander Roberts and William Rambaut.；Ante-Nicene Fathers；Vol. 1.；Edited by Alexander Roberts, James Donaldson, and A. Cleveland Coxe.；Buffalo, NY: Christian Literature Publishing Co.,；1885.；<http://www.newadvent.org/fathers/0103215.htm>.}

% structure: p0001 source=h1 target=section reason=page-title

\section{驳斥异端（第二卷，第十六章）}

1. 但如果他们不同意这些结论中的任何一点，因为他们会被我们证明无法为他们的\OriginalTerm{圆满}{Pleroma}的这种产生提供任何理由，那么他们必然会被迫承认：在他们的\OriginalTerm{圆满}{Pleroma}之上，还有另一个更属灵、更有力的体系，他们的圆满是依照其形象形成的。因为如果\OriginalTerm{造物主}{Demiurge}并没有自己构建那存在的受造物的形象，而是按照上面事物的形式制造了它，那么他们的\OriginalTerm{深渊}{Bythus}——它无疑使圆满具有这种配置——是从谁那里接受了那些在他之前存在的事物的形象呢？因为必须要么是（创造的）意图存在于那个创造世界的天主之中，以便他凭自己的力量、从自身获得其形成模式的模型；要么，如果偏离了这个存在，那么就会产生不断地追问上面那个存在者从何处得到了已造之物的配置、生产的数量以及模型本身的实质等问题。然而，如果深渊有能力凭自己将这种配置赋予圆满，那么为什么造物主就不能凭自己形成这样存在的世界呢？再者，如果受造物是上面事物的形象，我们为什么不能断言那些（上面的事物）又是它们上面其他事物的形象，而这些上面的事物又是其他事物的形象，如此推断无穷无尽的形象形象呢？\par

2. 巴西里德在完全偏离真理之后，就遇到了这个困难，他设想通过彼此形成的那些存有者的无限系列来逃避这种困境。当他宣告通过彼此相继和相似形成了三百六十五层天，并且这些天的明显证据存在于一年的天数中，如我先前所说；并且在这些之上有一个力量，他们也称之为\OriginalTerm{不可名者}{Unnameable}，及其管理体系——他即使这样也没有逃脱这种困境。因为，当被问及那最高之天的配置形象从何而来——他希望其余的天是通过系列由它形成——他将说，来自那属于\OriginalTerm{不可名者}{Unnameable}的管理体系。那么他必须说，要么是\OriginalTerm{不可言者}{Unspeakable}凭自己形成了它，要么他不得不承认在这个存有者之上还有另一个力量，他的不可名者正是从它那里得到了如此大量的配置，正如他所说的存在的那样。\par

3. 那么，立即承认真理是何等更安全、更准确的做法：即这个天主，创造者，他形成了世界，是唯一的天主，除他以外没有别的天主——他自己从自身接受了所造之物的模型和形象——而不是在使我们对如此不敬而曲折的描述感到厌倦之后，被迫在某个点上把心思固定在某个一上，并承认受造之物的配置是从他而来的。\par

4. 至于瓦伦提努斯的追随者对我们的指责，他们声称我们停留在下面的\OriginalTerm{七元体}{Hebdomad}中，好像我们不能将心智提升到高处，也不能理解上面的事物，因为我们不接受他们怪诞的主张：这同样的指责，巴西里德的追随者又反过来指责他们，因为（瓦伦提努斯派）还在围绕下面的事物，直到第一和第二\OriginalTerm{八元体}{Ogdoad}，并且他们不熟练地想象，在三十个\OriginalTerm{永世}{Æons}之后，他们就发现了超越万物的圣父，没有在思想上追踪他们的研究到那超过三百六十五层天的\OriginalTerm{圆满}{Pleroma}，即超过四十五个八元体。而任何一个人，如果想象四千三百八十层天或永世（因为一年中的天数包含那么多小时），也可以对他们提出同样的指责。如果再加上夜晚，从而将前述的小时加倍，想象他已经发现了大量的八元体和无穷无尽的永世，并因此反对那超越万物的圣父，自以为比所有人都更完美，那么他也会对所有人提出同样的指责，因为他们不能上升到他所宣布的如此众多的天或永世的概念，而是要么欠缺，停留在下面的事物中，要么停留在中间的空间。\par

\SourceRecord{Translated by Alexander Roberts and William Rambaut.；Ante-Nicene Fathers；Vol. 1.；Edited by Alexander Roberts, James Donaldson, and A. Cleveland Coxe.；Buffalo, NY: Christian Literature Publishing Co.,；1885.；<http://www.newadvent.org/fathers/0103216.htm>.}

% structure: p0001 source=h1 target=section reason=page-title

\section{\WorkTitle{反驳异端（卷二，第十七章）}}

1. 那么，那个关于他们\OriginalTerm{圆满}{Pleroma}的体系，尤其是其中涉及首要\OriginalTerm{八元体}{Ogdoad}的部分，既已背负如此众多的矛盾和困惑，现在让我继续考察他们体系的其他部分。我这样做，是因他们的狂妄，不得不查问一些并不真实存在的事物；但这样做是必要的，因为这项任务已托付于我，我愿众人认识真理，也因你亲自向我请求，要得到推翻这些人的（观点）的充分而完全的方法。\par

2. 那么，我问：其余的\OriginalTerm{永世}{Æon}是以何种方式产生的？是与其生产者结合，如同太阳光线与太阳结合？还是实际而分离地，使每个永世都拥有独立的存在和自己独特的形状，如同一个人从另一个人、一群牲畜从另一群牲畜那样？或是像发芽一样，如同树枝从树长出？它们与生产者是同质的，还是从别的某种实体获得其实体？还有，它们是同时产生，因而同时代，还是按一定顺序，有些较年长，有些较年轻？再者，它们是纯一而单一的，彼此完全相等和相似，如同灵和光那样产生，还是复合而不同的，在各自的组成部分上不相像？\par

3. 如果每个永世都是如同人那样实际地并按各自世代产生的，那么，由父所生的这些永世，要么与父同质，与作者相似；要么如果它们显得不相似，就必须承认它们是由某种不同的实体形成的。现在，如果父所生的存在者与作者相似，那么被产生的就必须永远不受苦，如同那生产者一样；但如果它们是由一种不同的、能受苦的实体形成的，那么这种不同的实体何以能在不朽的圆满内找到位置？而且，根据这个原则，每个永世都必须被理解为完全彼此分离，如同人彼此不相混合也不相联，各自有自己独特的形象和确定的活动领域，而且每个永世也有特定的尺寸——这些都是身体的特征，而非灵的特征。所以，请他们不要再称圆满为\Emph{属灵的}，或自称\Quote{\InnerQuote{属灵的}}，如果他们的永世确实与父一同坐席，如同人一样，而父本身也具有那些由祂所生的永世所揭示的形貌。\par

4. 再者，如果永世是从\OriginalTerm{逻各斯}{Logos}而来，逻各斯从\OriginalTerm{努斯}{Nous}而来，努斯从\OriginalTerm{深渊}{Bythus}而来，正如光从光点燃——例如火把从火把点燃——那么它们确实在生成和大小上彼此不同；但由于它们与它们的生产者同质，它们要么全都永远不受苦，要么它们的父本身也必须受苦。因为后来点燃的火把，不能拥有与先前不同的光。因此，它们的光混合为一时就回归原初的同一性，因为那从起初就存在的光被形成为一体。但对于光本身，我们不能说部分较新，部分较旧（因为整体只是一个光）；对于接受了光的火把，我们也不能这样说（因为它们就质料而言是同时的，火把的质料是同一的），而只能就点燃的时间而言，因为一个不久前点燃，另一个刚刚点燃。\par

5. 因此，那个与无知相关的受苦的缺陷，要么也会附着于他们的整个圆满，因为（其所有成员）是同质的；而\OriginalTerm{原始父}{Propator}也将分享这个无知的缺陷——即，他将不认识自己；要么另一方面，圆满内的所有光将永远不受苦。那么，最小的永世的受苦从何而来，如果父的光是形成所有其他光的光，并且本性不受苦？而且，在永世之间，如何能说有一个较年轻或较年长，既然整个圆满内只有一个光？如果有人称它们为星辰，它们也全都显示分享同样的本性。因为如果\BibleQuote{星辰与星辰的光辉虽有差别}\BibleRef{格前 15:41}，但它们在性质、实体或受苦与否方面并无差别；那么，所有这些光，既然都从父的光而来，要么本性不受苦、不变异，要么它们与父的光共同受苦，并能够经历败坏的变化。\par

6. 同样的结论也会得出，即使他们声称永世是从逻各斯产生，如同树枝从树长出，因为逻各斯是从他们的父而生。因为所有（永世）都由与父相同的实体形成，彼此只在大小而非在本性上不同，并充满父的广大，如同手指完成手。因此，如果父存在于受苦和无知中，那么由祂所生的永世也必如此。但如果将无知和受苦归于万有之父是不敬，他们怎能描述由祂所生的永世为能受苦？当他们将同样的不敬归于天主的\OriginalTerm{智慧}{Sophia}本身时，他们又怎能自称为虔诚的人？\par

7. 此外，倘若他们宣称他们的\OriginalTerm{永世}{Æons}被发出，正如光芒从太阳发出一样，那么，既然所有都出于同一实体并源于同一源头，所有便要么都能与那产生它们者一同\OriginalTerm{受难}{passion}，要么都将永远\OriginalTerm{不受苦}{impassible}。因为他们再不能坚持说，在所产生之物中，有些是不受苦的，有些是可受难的。那么，如果他们宣称所有都不受苦，他们便自毁其论据。因为如果所有都不受苦，最年幼的\OriginalTerm{永世}{Æon}如何能遭受受难？另一方面，如果他们宣称所有都分享了这受难（正如他们中有些人竟敢主张的那样），那么，既然这受难始于\OriginalTerm{圣言}{Logos}，而涌向\OriginalTerm{智慧}{Sophia}，他们便会被证实将受难追溯至圣言，亦即这\OriginalTerm{元父}{Propator}的\OriginalTerm{心灵}{Nous}，从而承认元父的心灵及圣父本身也经历了受难。因为万有之父不应被视为一种复合的存在，能与其心灵分离，正如我已表明的；实际上，心灵就是圣父，圣父就是心灵。因此，必然得出：那从祂而生作为圣言者，或者确切地说，心灵本身（既然祂是圣言）既是完美且不受苦的，而由祂所出的那些产物，既然与祂同一实体，也应是完美且不受苦的，并应永远保持与那产生它们者相似。\par

8. 因此，再不能如这些人所教导的那样，认为\OriginalTerm{圣言}{Logos}在生成中占据第三位置而对父无知。此类事在人类的诞生中或许可视为可能，因为子女常常不知其父母；但在父的圣言身上则完全不可能。因为如果圣言存在于父内，祂知道祂所存在的那一位——即祂不无知于自己——那么，从祂发出的那些产物，既作为祂的能力且常与祂同在，便不会对那发出它们者无知，正如光芒之于太阳。因此，在\OriginalTerm{圆满}{Pleroma}内的天主的\OriginalTerm{智慧}{Sophia}，既然是以这样的方式被产生，便不可能受\OriginalTerm{受难}{passion}的影响并孕育如此的无知。但属于\Person{瓦伦提诺}体系的智慧，既然她是魔鬼的产物，则可能陷入各种受难并显示出最深的无知。因为当他们自己作证说，他们的母亲是一位漂泊\OriginalTerm{永世}{Æon}的后裔时，我们便无需再探究如此母亲的子女为何总在无知的深渊中游荡。\par

9. 我不知道，除了这些已提及的产物，他们还能谈及任何别的；事实上，我虽与他们就此形式多有讨论，却未曾知晓他们曾提出任何其他特殊种类的存有如所述方式被产生。他们只坚持这一点：这些中的每一个都被产生得\Emph{如此被产生}，以致仅知道那产生它的一位，而对那直接先于它的一位\OriginalTerm{无知}{ignorant}。但在这一问题上，他们并未就其产生方式或此种事如何在属灵存在中发生，提出任何论证。因为无论他们选择如何推进，他们都将感到自己不得不（而就真理而言，他们完全偏离了正当理性）进展到坚持说，他们的\OriginalTerm{圣言}{Word}（亦即那从\OriginalTerm{元父}{Propator}的\OriginalTerm{心灵}{Nous}生出者）——我说，不得不坚持说祂是在一种堕落状态中被产生的。因为他们认为，由完美的\OriginalTerm{深渊}{Bythus}先产生的完美心灵，并不能使由祂而出的产物完美，而只能使之对父的知识和伟大完全盲目。他们还坚持说，\OriginalTerm{救主}{Saviour}在那生来失明的人身上展示了这奥秘的象征\BibleRef{若 9:1}，因为\OriginalTerm{永世}{Æon}也是被\OriginalTerm{独生者}{Monogenes}以这种盲目方式产生的，即处于无知中，从而将无知和盲目妄加给天主的圣言——按他们的理论，圣言占据从元父而产生的第二位置。可钦佩的诡辩家！未知之父崇高境界的探索者！以及那些超天界奥秘的宣读者——\BibleQuote{就是天使也都切愿窥探这些奥秘}\BibleRef{伯前 1:12}——好让他们得知：从那位超越万有之父的心灵中，圣言被产生为\Emph{盲目的}，即对那产生祂的父无知！\par

10. 但是，你们这些可怜的诡辩家，父的理智，甚至父自己既然是理智且在一切事上圆满，如何能生出祂自己的圣言成为一个不完全、盲目的永世呢？祂本来也能连同圣言一起产生对父的认识。既然你们断言基督是在其余之后生成的，却又宣称祂被生为完全的，那么圣言比祂年长，更应被同一个理智生为毫无疑问的完全，而非盲目；祂也不会再产生比祂自己更盲目的永世，直到最后你们的索菲亚，始终完全盲目，生出了如此庞大的一堆邪恶。你们的父是这一切祸患的原因；因为你们宣称你们父的伟大和能力是无知的原因，把祂比作深渊，并将这个名归于不可言说的父。但无知是恶，你们宣称一切恶都从它获得力量，同时你们又主张父的伟大和能力是无知的原因，这样你们确实把祂说成是\Quote{一切}祸患的肇始者。因为你们陈述恶的原因是\Quote{没有人}能默观祂的伟大。但如果父从一开始实在不可能使祂所造的那些存在认识祂，那么祂在这事上应被视为无可指责，因为祂\Emph{不能}消除随祂之后者们的无知。但若在后来，当祂愿意时，祂\Emph{能够}消除那随着相继的受造物而增长、因而深深扎根于永世中的无知，那么，若祂愿意，祂更早就能阻止那尚不存在的无知产生。\par

11. 既然因此，一旦祂愿意，祂不仅使永世认识祂，也使这些生活在末世的世人认识祂；但既然祂从起初不愿意被认识，祂便保持不认识——那么，按你们所说，无知的原因是父的旨意。因为若祂预知这些事将来会这样发生，祂为何不在无知在这些存在中获得立足之前就防止它，反而在事后，仿佛出于后悔，通过产生基督来处理它呢？因为那通过基督传达给众人的认识，祂本可更早地通过圣言（祂也是独生子的首生者）来传授。或者，若祂预先知道却愿意这些事如此发生，那么无知的作为必永远存续，永不消逝。因为按照你们最先之父的旨意所造成的事物，必与那制定旨意者的旨意一起存续；或若它们消逝，那定意使它们存在之人的旨意也必随它们一同消逝。永世们为何在得知\Quote{父是完全不可理解的}之后才得安息并达到圆满的认识呢？它们本可在陷入激情之前就拥有这认识；因为父的伟大从起初并未减少，以致它们能知道祂是完全不可理解的。若因祂的无限伟大，祂一直不被认识，那么也因祂的无限慈爱，祂就该使祂所产生的那些存在保持无激情，因为没有什么阻碍，且合宜性反而要求它们从起初就知道父是完全不可理解的。\par

\SourceRecord{Translated by Alexander Roberts and William Rambaut.；Ante-Nicene Fathers；Vol. 1.；Edited by Alexander Roberts, James Donaldson, and A. Cleveland Coxe.；Buffalo, NY: Christian Literature Publishing Co.,；1885.；<http://www.newadvent.org/fathers/0103217.htm>.}

% structure: p0001 source=h1 target=section reason=page-title

\section{驳斥异端（第二卷，第十八章）}

1. 她们还宣称这个\OriginalTerm{索菲亚}{Sophia}曾陷入无知、堕落与情欲之中，这除了说荒谬之外，还能怎么看待呢？因为这些事情与智慧是相异且相反的，绝不能成为其属性。凡是没有先见、不知何者有益的地方，智慧就不存在。因此，她们不要再称这位受苦的\OriginalTerm{永世}{Æon}为\OriginalTerm{索菲亚}{Sophia}，而应当要么放弃她的名字，要么放弃她的苦难。而且，她们也不该称整个\OriginalTerm{圆满}{Pleroma}为属灵的，如果这位\OriginalTerm{永世}{Æon}在经历如此激烈的情欲骚动时还位于它之中的话。因为就连一个强有力的灵魂（更不用说属灵实体）也不会经历任何此类事情。\par

2. 再者，她的\OriginalTerm{内思}{Enthymesis}怎能与情欲一同出去而成为一个独立的存在呢？因为\OriginalTerm{内思}{Enthymesis}总是与某个人相关联，绝不可能自行孤立存在。因为一个坏的\OriginalTerm{内思}{Enthymesis}会被一个好的\OriginalTerm{内思}{Enthymesis}所消灭和吸收，一如疾病状态被健康所吸收。那么，在情欲的\OriginalTerm{内思}{Enthymesis}之前，先有的是哪一种\OriginalTerm{内思}{Enthymesis}呢？就是探究父的本性和思想他的伟大。但她后来相信了什么而恢复了健康呢？就是：父是不可理解的，是无法探测的。因此，她想认识父并非一种恰当的情感，并因此成为可受情欲的；但当她知道他是不可探测的，她就恢复了健康。就连\OriginalTerm{努斯}{Nous}本人，他先前探究父的本性，据她们说，当他得知父是不可理解的时，也停止了继续研究。\par

3. 那么，\OriginalTerm{内思}{Enthymesis}怎么能单独构想出情欲，而这些情欲本身就是它的感受呢？因为感受必然与个体相连：它不能独立产生或存在。然而，她们的这种见解不仅站不住脚，而且与主所说的相悖：\BibleQuote{你们求，就会找到。}\BibleRef{玛 7:7} 因为主使祂的门徒通过寻找并找到父而变得完美；而她们那位在高天之上的基督，却因命令\OriginalTerm{永世}{Æon}们不要寻求父，使它们变得完美，说服它们即使努力劳作也找不到祂。她们宣称自己完美，是因为她们坚持找到了她们的\OriginalTerm{深渊}{Bythus}；而\OriginalTerm{永世}{Æon}们则因为她们所寻求的祂是不可探测的而变得完美。\par

4. 既然\OriginalTerm{内思}{Enthymesis}本身不能与\OriginalTerm{永世}{Æon}分离而独立存在，那么她们进一步把她的情欲分开、分离出来，并宣称那是物质的实体，这显然加大了谎言。好像天主不是光，好像没有圣言能定罪他们、推翻他们的邪恶。因为确实如此：\OriginalTerm{永世}{Æon}所想的就是她所经历的，她所经历的就是她所想的。据她们所说，她的\OriginalTerm{内思}{Enthymesis}无非就是一个思考如何把握不可把握者的情欲。因此，\OriginalTerm{内思}{Enthymesis}就是情欲，因为她思考的是不可能的事。那么，感受和情欲怎么能与\OriginalTerm{内思}{Enthymesis}分离并分离开，以至于成为如此庞大物质世界的实体呢？既然\OriginalTerm{内思}{Enthymesis}本身就是情欲，情欲就是\OriginalTerm{内思}{Enthymesis}？因此，\OriginalTerm{内思}{Enthymesis}不能离开\OriginalTerm{永世}{Æon}而独立具有实体，感受也不能离开\OriginalTerm{内思}{Enthymesis}而独立具有实体；这样，她们的体系又一次崩溃瓦解了。\par

5. 但是，\OriginalTerm{永世}{Æon}怎么会解体（化为组成部分）并成为情欲的奴隶呢？她无疑与\OriginalTerm{圆满}{Pleroma}具有相同的实体；而整个\OriginalTerm{圆满}{Pleroma}来自父。任何实体，当与相似本性接触时，不会化为无有，也不会面临毁灭的危险，反而会继续存在并增长，如火遇火，灵遇灵，水遇水；而彼此本性相反的实体相遇时，会受苦、改变、毁灭。同样，如果有光的产物，它不会在相似的光中受苦或遇到危险，反而会更加明亮并增长，如同太阳的光芒使白昼增长；因为他们主张\OriginalTerm{深渊}{Bythus}本身就是父亲\OriginalTerm{索菲亚}{Sophia}的形象。凡是习性相异、彼此陌生、本性相反的动物，相遇时会陷入危险并毁灭；而另一方面，那些彼此熟悉、性情和谐的动物，在同一地方相处没有危险，反而因此获得安全和生命。因此，如果这个\OriginalTerm{永世}{Æon}是由与整体相同实体的\OriginalTerm{圆满}{Pleroma}所产生，她就不可能经历变化，因为她与相似和熟悉的存在在一起，属灵本质与属灵者在一起。因为恐惧、惊慌、情欲、解体等，可能发生在我们这些有身体的人因对立冲突而产生的过程中；但在属灵的存在以及有光照耀他们的存在中，不可能发生此类灾祸。但在我看来，这些人赋予他们的\OriginalTerm{永世}{Æon}以那种情欲，如同喜剧诗人米南德笔下的角色，其人深陷爱河却遭人憎恨。因为发明这些见解的人，与其说具有属灵和神圣实体的观念和想象，不如说具有人类中某个不幸的情人的观念和想象。\par

6. 此外，默想如何探究完美圣父的（本性），渴望存在于祂内，并理解祂的（伟大），并不能给一个属神的永世带来无知或情欲的玷污；反而会（产生）完美、无情感和真理。因为他们并不声称，即便是他们，虽然是凡人，通过默想那在他们之前的那位——而现在仿佛理解完美者，并被置于对祂的认识中——就会陷入困惑的情欲，反而达到对真理的认识和理解。因为他们断言，救主曾对祂的门徒说：\Quote{求则得之}，正是为此，要他们寻求那由想象而被他们视为在万有的创造者之上的那位——不可言说的深渊；并且他们自己渴望被视为\Quote{完美者}，因为他们尚在世上时就已经寻求并找到了完美者。然而他们却宣称，那在圆满之内的永世，一个完全属神的存在，因寻求先父，并竭力要进入祂的伟大之中，且渴望理解圣父的真理，就堕入了（承受）情欲，且是那样的情欲，若非她遇见了那维持万有的能力，她就会消融于（众永世的）普遍实体中，从而结束她（个别的）存在。\par

7. 如此假设是荒谬的，真正是完全缺乏真理之人的意见。因为，按照他们自己的体系，他们承认这个永世比他们自己更高、更古老，当他们断言自己正是那经历情欲的永世之深思的果实，如此这永世是他们母亲的父亲，也就是他们自己的祖父。对于他们这些后来的孙辈，对圣父的寻求，如他们所持守的，带来了真理、完美、稳固、从不稳定物质中的解放以及与圣父的和好；然而对于他们的祖父，同样的寻求却招致了无知、情欲、恐惧和困惑，他们也宣称物质的实体正是从这些（扰动）中形成的。因此，若说对完美圣父的寻求与探究，以及与祂共融和结合的渴望，对他们是相当有益的，但对那位连他们自己的起源也来自其中的永世，这些事却是消融与毁灭的原因，那么这样的主张除了被视为全然矛盾、愚蠢和非理性之外，还能如何看呢？那些听从这些师傅的人，真是盲目自己，而他们拥有的盲目的向导，也（理当）与他们一同坠入下面那无知的深渊。\par

\SourceRecord{Translated by Alexander Roberts and William Rambaut.；Ante-Nicene Fathers；Vol. 1.；Edited by Alexander Roberts, James Donaldson, and A. Cleveland Coxe.；Buffalo, NY: Christian Literature Publishing Co.,；1885.；<http://www.newadvent.org/fathers/0103218.htm>.}

% structure: p0001 source=h1 target=section reason=page-title

\section{\WorkTitle{Against Heresies (Book II, Chapter 19)}}

1. 但是，关于他们的种子的这种说法——即它是由母亲按照那些侍候救主的天使的形象而受孕的，无形无状、无定形且不完美；并且它在德穆革不知情的情况下被存放在他内，以便通过他的中介，在他所拥有的、[可以这么说]被充满种子的那个灵魂中达到完美和形状——这又是什么言论呢？这首先是肯定那些侍候他们救主的天使是不完美的，且没有形象或形状；如果那按照他们的模样受孕的产物确实生成了任何此类存在[如所述]的话。\par

2. 其次，关于他们所说的造物主对于发生在他内的种子存放一无所知，以及他对由他所作的给与人的种子传递也一无所知，这些话是空洞无益的，并且丝毫不能被证明。因为如果那种子具有任何实体和特殊属性，他怎能对此无知呢？反之，如果它没有实体和性质，因此实际上是虚无，那么他当然对此无知。因为那些具有自身特定运动和性质的事物，无论是热、快、甜，还是在光辉中与其它事物不同，都不会逃避人类，因为它们混合在人类活动范围内；更不可能[向]天主、这个宇宙的创造者[隐藏]。然而，有理由[说]他们的种子不为他所知，因为它没有任何普遍有用的性质，也没有任何行动所需的实体，实际上是一个纯粹的非实体。在我看来，主确实针对这类意见表达了如下话语：\BibleQuote{凡人们所说的每一句废话，在审判之日，都要交账。}因为所有像这样的导师，用废话充满人的耳朵，当他们站在审判宝座前时，将为他们徒然想象并虚假地反对主而言论交账，因为他们已经到了如此胆大妄为的地步，竟然宣称自己由于种子的实体而认识属神的普累若玛，因为那内住的人向他们启示了真正的父；因为属魂的本性需要藉着感官来受管教。但[他们认为]德穆革，虽然将整个这种子接受到自己内，但由于它被母亲存放给他，却仍然对所有事情完全无知，对任何与普累若玛相关的事毫不理解。\par

3. 并且他们声称自己是真正\Quote{属神的}，因为宇宙之父的某一部分已被存放在他们的灵魂中，因为按照他们的断言，他们拥有的灵魂与德穆革本身由同样的质料形成，然而他，虽然从母亲一次性地接受了全部[神圣的]种子，并拥有它在自己内，却仍然保持属魂的本性，对高于自己的事没有丝毫理解，而这些事他们吹嘘自己仍在世时就理解了——这不是达到一切可能荒谬之最吗？因为想象完全相同的种子将这些人的灵魂传递了认识和完美，却使创造他们的天主产生无知，这种意见只有那些完全疯狂、完全缺乏常识的人才会持有。\par

4. 再者，他们说种子通过这样被存放而得以成形、增长，并因此准备好接受完全的理性，这也是极其荒谬且毫无根据的。因为其中将会有物质——那种他们认为来自无知和缺失的实体；[而这将证明自身]比他们父的光更适宜和有用，因为如果那[光]按照那[光]的观照而生出时，它是无形无状的，但从此[物质]获得了形状、外貌、增长和完美。因为如果来自普累若玛的那光是一个属神存在的原因，使之既无形状也无外貌且无自身特殊的体积，而它下降到这个世界却增添了所有这些事物，并使之达到完美，那么在此世居留（他们也称之为黑暗）似乎比他们父的光更有功效和有用。但怎么不能被视为荒谬呢？他们肯定说，他们的母亲几乎陷于物质之中，几乎被物质消灭，若不是她当时好不容易向外伸展并仿佛跳出自身，从父得到援助；但她的种子在这同一物质中增长了，接受了形状，并变得适合接受完全的理性；而且，这是在与其不相似且不熟悉的实体中\Quote{涌现出来}，按照他们自己的宣称，属神之物与属魂之物对立，属魂之物与属神之物对立。既然如此，他们所说的\Quote{小粒子}怎能在与自己对立且不熟悉的实体中增长、获得形状并达到完美呢？\par

5. 但是，进一步，除上述以外，还有一个问题：他们的母亲，当她看见众天使时，是同时一下子生出种子，还是一个一个[依次]生出？如果她是同时一下子全部生出，那么所生出的东西现在不可能是幼稚的；因此，它降入到现今存在的人们中必定是多余的。但如果是一个一个，那么她并不是按照她所看见的那些天使的形象而受孕；因为她同时一次性地观想它们，以便藉着它们受孕，她本该一次性生出那来自于其形象她曾一次性受孕的后代。\par

6. 再者，为何她在看见天使与\OriginalTerm{救主}{Saviour}同在时，确实怀了\Emph{他们的}形像，却没有怀那比她更美丽的\Emph{救主}的形像呢？难道救主不合她的心意，因此她没有按他的肖像怀孕吗？再者，他们称为动物性的\OriginalTerm{造物主}{Demiurge}，如他们所主张，拥有自己特定的体积与形状，为何在实体上已完美地产生，而那属灵的——它本应比属动物的更有力量——却被发出时不完美，且需要降入一个灵魂中，在其中获得形状，然后变得完美，才适合接受完美的理性呢？如果它仅在属地上、属人的凡人中得形状，那么它便不能再说是按被称为光的那些天使的肖像，而是按这些地上之人的肖像。因为在那种情况下，它所拥有的不是天使的肖像与外观，而是那些灵魂的肖像与外观，它也在其中获得形状；正如水倒入器皿时，取那器皿的形状，若偶尔在其中凝结，便获得它所冻结的器皿的形状；因为灵魂本身拥有其所居住身体的外形，它们自己已适应于其所在的器皿，如我先前所说。因此，若那（所谓的）\OriginalTerm{种子}{seed}在此固化并形成确定形状，它将拥有人的形状，而非天使的形状。那么，那种子既然按人的肖像获得形状，又如何能是天使的肖像呢？再者，既然它属于属灵的本性，又何需降入肉身呢？因为属血肉的若欲得救，确实需要属灵的，以便在其中被圣化、清除一切不洁，并使可朽的被不朽吞没；但属灵的本性完全不需要这些下界之物。因为不是我们使之受益，而是它使我们改善。\par

7. 他们关于那种子的言论更明显地显为虚假，且对众人显而易见，因为他们宣称那些从\OriginalTerm{母亲}{Mother}接受种子的灵魂超越所有其他灵魂；因此它们也被造物主尊荣，被立为首领、君王和司祭。若果真如此，大司祭盖法、亚纳斯，以及其他司祭长、经师和民众的统治者本应首先信从主，因他们与那亲属关系相符；甚至在他们之前，黑落德王也应信从。但既然他、司祭长、统治者以及民众中的显贵都没有（以信德）归向祂，反而是那些在路旁乞讨的、耳聋的、瞎眼的归向了祂，而祂被其他人弃绝轻视，正符合保禄所宣称的：\BibleQuote{弟兄们！你们看看你们是怎样蒙召的：按肉眼来看，你们中并不多有智慧者，并不多有权贵者，并不多有显赫者；天主却拣选了世上愚拙的。}因此，这样的灵魂并非因其中的种子而优于其他，也并非因此被造物主尊荣。\par

8. 关于他们的体系软弱、站不住脚且完全荒诞不经这一点，已经说得够多了。因为无需用一句俗语说，想要知道海水是咸的，就得把大海喝干。但正如一个用泥土塑成的雕像，外表涂上颜色，使人以为它是金的，而实际上它是泥的；有人从它上面取下一小块，因而将它打开，露出泥土，便可使寻求真理的人从错误的看法中解脱出来；同样，我（不仅揭露了其体系的一小部分，而是揭露了其中最重要的几条要点）已向那些不愿意故意被误导的人指明了\OriginalTerm{瓦伦廷派}{Valentinians}和所有其他异端者的邪恶意念，他们散布关于造物主——即这个宇宙的工匠和形成者，而事实上祂是唯一真天主——的恶毒意见，（正如我所做的那样，）显明了他们的观点多么容易被推翻。\par

9. 因为凡稍有智力、拥有少许真理的人，怎能容忍他们肯定说在创造主之上另有其他神？又说另有另一位\OriginalTerm{独生者}{Monogenes}，以及另有一位天主的\OriginalTerm{圣言}{Word}，他们还将祂描述为在（一种）退化状态中产生；另有另一位\OriginalTerm{基督}{Christ}，他们断言祂与\OriginalTerm{圣神}{Holy Spirit}一同形成，晚于其他\OriginalTerm{永世}{Æons}；另有另一位\OriginalTerm{救主}{Saviour}，他们说祂并非出自万有之父，而是那些在（一种）退化状态中形成的永世的一种共同产物，并且祂是由于这退化本身而被必然产生的？因此他们认为，除非永世处于无知和退化状态，否则基督、圣神、\OriginalTerm{界限}{Horos}、救主、天使、他们的母亲、她的种子以及其余世界的架构根本不会产生；宇宙将会是荒芜的，并且没有其中存在的许多美善之物。因此，他们不仅犯有对创造主的不虔敬，宣称祂是缺陷的果实，而且也对基督和圣神不虔敬，肯定他们是因那缺陷而产生的；同样，救主是在那缺陷之后（产生的）。谁能容忍他们其余的空谈呢？他们狡猾地试图将空谈附会到比喻中，从而既使自己、也使那些相信他们的人陷入了最深的不虔敬之中。\par

\SourceRecord{Translated by Alexander Roberts and William Rambaut.；Ante-Nicene Fathers；Vol. 1.；Edited by Alexander Roberts, James Donaldson, and A. Cleveland Coxe.；Buffalo, NY: Christian Literature Publishing Co.,；1885.；<http://www.newadvent.org/fathers/0103219.htm>.}

% structure: p0001 source=h1 target=section reason=page-title

\section{驳斥异端（第二卷，第二十章）}

1. 他们不恰当且不合逻辑地将主的比喻和作为应用于他们虚假编造的系统，我证明如下：例如，他们试图通过以下事实来证明他们所说的发生在第十二永世身上的受难：救主的受难是由第十二位宗徒引起的，并且发生在第十二个月。因为他们认为祂在受洗后只传道一年。他们还主张，同一件事在患血漏的妇人身上清楚地预示了。因为那妇人患病十二年，通过触摸救主的衣边，被从救主身上发出的能力治愈，他们断言那能力先前就已存在。因为那受难的能力正向外伸展并流向无限，以致她有被溶解到永世的普遍实体中的危险；但随后，她触及了由衣边所预表的首四元组，就被止住，停止了她的受难。\par

2. 再者，至于他们断言第十二永世的受难是通过犹达斯的行为证明的，犹达斯怎可能与她相比，作为她的预像呢？——他是从十二位中被驱逐，从未恢复其地位的人。因为那永世——他们宣称犹达斯是她的预表——在与她的思念分离后，被恢复或召回她原来的位置；而犹达斯被剥夺了职务，被驱逐，而玛弟亚按所记载的被任命代替他：\BibleQuote{他的监督职位，让别人取得}\BibleRef{宗 1:20}。因此，他们应当主张第十二永世被逐出圆满，并且产生了另一位来填补她的位置；如果她确实被犹达斯所指向的话。此外，他们告诉我们，是永世自己受难，而犹达斯是出卖者。连他们自己也承认是受难的基督，而非犹达斯，前来承受苦难。那么，出卖那位为了我们的救恩而必须受难之人的犹达斯，怎么会是那受难永世的预表和形象呢？\par

3. 但事实上，基督的受难既与永世的受难不相似，也非发生在类似的情况下。因为永世经历了一种瓦解和毁灭的受难，以致那受难者也有被毁灭的危险。但是主，我们的基督，经历了一种有效的、而非偶然的受难；祂自己不仅没有毁灭的危险，反而用自己的力量坚固了堕落的人，并召叫他们回归不朽。再者，永世在寻求圣父时受难，且未能找到祂；但主受难，是为了将那些偏离圣父的人领回认识祂并与祂共融。对圣父伟大的寻求对她成了一种导致毁灭的受难；但主受难并赐予对圣父的认识，反而为我们带来了救恩。她的受难，如他们所说，产生了一个女性后裔，软弱、病态、未成形且无效；但祂的受难产生了力量和权能。因为主通过受难，\BibleQuote{上升高天，掳掠了俘虏，恩赐予人}\BibleRef{厄 4:8}，并赐给那些信靠祂的人权柄\BibleQuote{践踏蛇和蝎子，并制胜仇敌的一切势力}\BibleRef{路 10:19}，即那背道领袖的势力。我们的主也通过受难毁灭了死亡，驱散了谬误，终结了腐朽，消除了无知，同时显明了生命，启示了真理，并赐予了不朽的恩赐。但他们的永世，当她受难时，却建立了无知，并产生了一个无形质体，从中产生了所有物质性作品——死亡、腐朽、谬误以及诸如此类。\par

4. 那么，在门徒中位列\Emph{第十二}的犹达斯并不是那受难永世的预表，主的受难也不是；因为这两件事已经被证明在各方面互不相似且不协调。这不仅就我已提及的各点而言如此，就数字本身而言也是如此。因为出卖者犹达斯位列\Emph{第十二}，这是所有人同意的，因为在福音中提到了十二位宗徒的名字。但是这位永世不是\Emph{第十二}，而是\Emph{第三十}；因为按照所考虑的观点，由圣父的意志所产的永世不是十二位，她也不是在\Emph{第十二}位被派出；相反，他们认为她是在\Emph{第三十}位产生的。那么，位列\Emph{第十二}的犹达斯，怎能是那占据\Emph{第三十}位的永世的预表和形象呢？\par

5. 但如果他们说丧亡的犹达斯是她的\OriginalTerm{思念}{Enthymesis}的肖像，那么即使如此，这肖像与那[假设]与之对应的真理也没有任何类似之处。因为\OriginalTerm{思念}{Enthymesis}从\OriginalTerm{永世}{Æon}中分离出来，后来从基督那里获得形态，又由救主分有理智，并按照\OriginalTerm{圆满}{Pleroma}内的事物的肖像，形成了\OriginalTerm{圆满}{Pleroma}之外的一切事物，最终据说被它们接纳进入\OriginalTerm{圆满}{Pleroma}，并按照[结合]的原则，与那由所有事物形成的救主联合。但是犹达斯既已永远被抛弃，就绝不会再回到门徒的数目中；否则就不会另选一人填补他的位置了。此外，主也论及他说：\BibleQuote{那出卖人子的人有祸了！}\BibleRef{玛 26:24}，并且：\BibleQuote{那人若没有生下来，为他倒好。}\BibleRef{谷 14:21}，他被主称为\BibleQuote{丧亡之子}\BibleRef{若 17:12}。然而，如果他们说犹达斯是\OriginalTerm{思念}{Enthymesis}的预象，不是指与\OriginalTerm{永世}{Æon}分离的思念，而是指与她纠缠的\OriginalTerm{苦难}{passion}，那么这样十二这个数也不能被视为三这个数的[合适]预象。因为在前一种情况下，犹达斯被抛弃，玛弟亚被选立代替他；但在后一种情况下，\OriginalTerm{永世}{Æon}据说处于瓦解和毁灭的危险中，并且还有她的\OriginalTerm{思念}{Enthymesis}和\OriginalTerm{苦难}{passion}：因为他们明显区分\OriginalTerm{思念}{Enthymesis}和\OriginalTerm{苦难}{passion}；他们描述\OriginalTerm{永世}{Æon}得到复原，\OriginalTerm{思念}{Enthymesis}获得形态，但\OriginalTerm{苦难}{passion}在与此两者分离后，成为物质。因此，既然有这三个，即\OriginalTerm{永世}{Æon}、她的\OriginalTerm{思念}{Enthymesis}和她的\OriginalTerm{苦难}{passion}，而犹达斯和玛弟亚仅有两人，就不能作为它们的预象。\par

\SourceRecord{Translated by Alexander Roberts and William Rambaut.；Ante-Nicene Fathers；Vol. 1.；Edited by Alexander Roberts, James Donaldson, and A. Cleveland Coxe.；Buffalo, NY: Christian Literature Publishing Co.,；1885.；<http://www.newadvent.org/fathers/0103220.htm>.}

% structure: p0001 source=h1 target=section reason=page-title

\section{驳斥异端（第二卷，第二十一章）}

1. 如果，再者，他们主张十二宗徒只是那由\OriginalTerm{人}{Anthropos}与\OriginalTerm{教会}{Ecclesia}结合所产生的十二\OriginalTerm{永世}{Æons}集团的预像，那么让他们另外举出十个宗徒作为另外十个\OriginalTerm{永世}{Æons}的预像，据他们宣称，这些\OriginalTerm{永世}{Æons}是由\OriginalTerm{圣言}{Logos}和\OriginalTerm{生命}{Zoe}产生的。因为不合理的是：较低级因而较卑微的\OriginalTerm{永世}{Æons}由救主通过选举宗徒来预示，而较高级因而更优越的\OriginalTerm{永世}{Æons}却没有这样被预示；因为救主（如果，也就是说，祂选择宗徒是出于这个目的，即通过他们来显示在\OriginalTerm{圆满}{Pleroma}中的\OriginalTerm{永世}{Æons}）本可以选择另外十个宗徒，同样还可以在这之前选择另外八个，以此预示最初原初的\OriginalTerm{八元体}{Ogdoad}。祂不能通过宗徒数目（已经）构成预像而显示第二个\OriginalTerm{二元十元体}{Duo Decad}的任何象征。因为\EditorialNote{祂没有选择这样的其他数目的门徒；但}在十二宗徒之后，我们发现主又差遣了其他七十人先祂而行。\BibleRef{路 10:1}现在\Emph{七十}不可能是一个\OriginalTerm{八元体}{Ogdoad}、\OriginalTerm{十元体}{Decad}或\OriginalTerm{三十元体}{Triacontad}的预像。那么，为什么较低级的\OriginalTerm{永世}{Æons}，如我所说，通过宗徒被代表，而较高级的，前者从它们获得存在，却完全没有被预表呢？但如果十二宗徒被选是为了通过他们表明十二\OriginalTerm{永世}{Æons}的数目，那么七十人也应该被选为七十\OriginalTerm{永世}{Æons}的预像；那样的话，他们必须断言\OriginalTerm{永世}{Æons}不再是三十，而是八十二个。因为那选择宗徒使之成为在\OriginalTerm{圆满}{Pleroma}中的那些\OriginalTerm{永世}{Æons}的预像的那一位，绝不会使宗徒成为一些\OriginalTerm{永世}{Æons}的预像而不成为另一些的；祂会通过宗徒努力保存一个形象并展示一个在\OriginalTerm{圆满}{Pleroma}中那些\OriginalTerm{永世}{Æons}的预像。\par

2. 再者，关于\Person{保禄}我们不可缄默，而要问他们那位宗徒是根据何种\OriginalTerm{永世}{Æon}的预像传给我们的，除非也许\EditorialNote{他们断言他是}那由所有\OriginalTerm{永世}{Æons}复合而成的救主（他来自整体收集的恩赐，他们称之为\OriginalTerm{万物}{All Things}，因为是由它们全体形成）的代表。关于这位存在，诗人\Person{赫西俄德}曾惊人地表达，称他为\OriginalTerm{潘多拉}{Pandora}——即\Quote{万有的礼物}——因为全体的最好恩赐都集中于他。在描述这些恩赐时，有如下记载：\Person{赫耳墨斯}（希腊语如此称呼），\OriginalTerm{将欺诈和欺骗的言语以及偷窃的习性植入他们}{\GreekText{Αἱμυλίους} \GreekText{τε} \GreekText{λόγους} \GreekText{καὶ} \GreekText{ἐπίκλοπον} \GreekText{ἦθος} \GreekText{αὐτοὺς} \GreekText{Κάτθετο}}（或用我们的话说），\Quote{将欺诈和欺骗的言语以及偷窃的习性植入他们}，目的是引导愚昧者误入歧途，使他们相信这些假话。因为他们的母亲——即\OriginalTerm{勒托}{Leto}——秘密地激动他们（因此她被称为勒托，根据希腊词义，因为她\Emph{秘密地}激动人），在\OriginalTerm{造物主}{Demiurge}不知情的情况下，向发痒的耳朵灌输深奥不可言喻的奥秘。\BibleRef{弟后 4:3}而且，不仅是他们的母亲通过\Person{赫西俄德}宣布了这一奥秘；而且非常巧妙地，通过抒情诗人\Person{品达}，当他向\OriginalTerm{造物主}{Demiurge}描述\OriginalTerm{珀罗普斯}{Pelops}的案例——其肉被父亲切成碎片，然后被所有神明收集、带拢并重新结合——她就这样暗示了\OriginalTerm{潘多拉}{Pandora}，而且这些人被她的烙铁烙过良心，宣称同样的事情，据他们声称，\EditorialNote{被证明}与其他人同出一族一气。\par

\SourceRecord{Translated by Alexander Roberts and William Rambaut.；Ante-Nicene Fathers；Vol. 1.；Edited by Alexander Roberts, James Donaldson, and A. Cleveland Coxe.；Buffalo, NY: Christian Literature Publishing Co.,；1885.；<http://www.newadvent.org/fathers/0103221.htm>.}

% structure: p0001 source=h1 target=section reason=page-title

\section{反驳异端（第二卷，第二十二章）}

1. 我已表明，\OriginalTerm{三十}{thirty}这个数字在各方面都不适用于他们；按他们的说法，有时在\OriginalTerm{圆满}{Pleroma}内找到的\OriginalTerm{永世}{Æons}太少，有时又太多（无法与那个数字对应）。因此，并没有三十个永世，救主也不是在三十岁时来受洗，为的是彰显他们体系中的三十个沉默的永世；否则，他们必须首先从万有之圆满中将救主自己分离并逐出。此外，他们声称祂在第十二个月受苦，这样祂在受洗后只传教了一年；他们试图从先知那里确立这一点（因为经上记载：\BibleQuote{宣扬上主恩慈之年，和我们天主报仇的日子} \BibleRef{依 61:2}），他们实在是瞎眼，竟声称发现了\OriginalTerm{深渊}{Bythus}的奥秘，却不明白依撒意亚所说的上主恩慈之年，也不明白报仇的日子。因为先知既不是说一个包含十二时辰的日子，也不是说一个长度为十二个月的年份。因为他们自己也承认，先知们常常用比喻和寓言来表达自己，[并且]不应[按照]字面的声音来理解。\par

2. 那么，那被称为报仇的日子，就是主要照各人的行为报答各人的日子——即审判。再者，上主恩慈之年乃是现今的时期，在这时期内，那些信祂的人蒙祂召唤，成为天主所悦纳的——即从祂来临直到（万物的）终结的整个时期，在这时期内，祂为自己获得那些得救的人，作为（慈悲计划的）果实。因为按先知的用语，报仇的日子紧随（恩慈）之年；如果主只传教一年，而先知说到它，那么先知将被证明是说谎者。因为报仇的日子在哪里呢？因为那年已经过去，报仇的日子尚未到来；但祂仍然\BibleQuote{叫日头照好人，也照歹人；降雨给义人，也给不义的人} \BibleRef{玛 5:45}。而义人受迫害、被磨难、被杀害，罪人却拥有丰裕，\BibleQuote{弹琴鼓瑟唱消闲的歌曲，却不留意上主的作为} \BibleRef{依 5:12}。但是，按（先知所用的）说法，它们应当结合在一起，报仇的日子应紧随（恩慈）之年。因为经文说：\BibleQuote{宣扬上主恩慈之年，和我们天主报仇的日子}。因此，现今这个时期，人们被主召唤并得救，适当地被理解为\Quote{上主恩慈之年}所指；紧随其后的乃是\Quote{报仇的日子}，即审判。而这里所指的时期不仅被称为\Emph{一年}，也被先知和保禄称为\Emph{一日}；这宗徒想起圣经，在写给罗马人的书信中说：\BibleQuote{正如经上所载：我们为了你整天被置于死地，被视为待宰的羊} \BibleRef{罗 8:36}。但这里\Quote{整天}一词被用来指整个这段时期，在此期间我们受迫害，被宰杀如同羊。因此，正如这\Emph{日}不是指一个由十二时辰组成的日子，而是指基督信徒为祂受苦被杀的整个时期，同样那里提到的\Emph{年}也不是指一个由十二个月组成的年份，而是指整个信德的时期，在此期间人们听闻并相信福音的宣讲，那些与祂合一的人成为天主所悦纳的。\par

3. 但令人大为惊异的是，这些人声称发现了天主的奥秘，却没有查考福音，去确证主领受圣洗后，按犹太人每年从各地聚集到耶路撒冷庆祝逾越节的惯例，曾多少次在逾越节期间上耶路撒冷。首先，祂在加里肋亚的\Place{加纳}变水为酒后，便上到耶路撒冷过逾越节庆日，当时记载：\BibleQuote{许多人看见祂所行的神迹，就信了祂，}\BibleRef{若 2:23}，主之门徒\Person{若望}如此记录。随后，祂〔从\Place{犹太}〕退去，出现在\Place{撒玛黎雅}；当时祂与撒玛黎雅妇人交谈，并远距离用一句话治愈了百夫长的儿子，说：\BibleQuote{你去吧，你的儿子活了。}\BibleRef{若 4:50}。后来祂第二次上耶路撒冷守逾越节庆日；当时祂治愈了那在池旁躺了三十八年的瘫子，命他起来，拿起床榻离去。再者，祂从那里退到\Place{提庇哩亚}海的对岸\BibleRef{若 6:1}，见有一大群人跟随祂，便用五个饼使所有群众吃饱，剩下的碎块装满了十二筐。后来，当祂复活了\Person{拉匝禄}，法利塞人图谋害祂时，祂退到一座名叫\Place{厄弗辣因}的城；从那里，如经上记载：\BibleQuote{逾越节前六天，祂来到\Place{伯达尼}，}\BibleRef{若 11:54}，并从\Place{伯达尼}上耶路撒冷，在那里吃了逾越节晚餐，次日受难。现在，任何人都必须承认，这三个逾越节不可能发生在同一年。而且，庆祝逾越节以及主受难的那个特定月份不是第十二个月，而是第一个月；那些自夸通晓一切的人，若不知此事，可从\Person{梅瑟}学习。因此，他们关于年份和第十二个月的解释已被证明是虚假的；他们应当要么摒弃自己的解释，要么摒弃福音；否则，这个无可回避的问题就迫使他们面对：主怎么可能只传道一年呢？\par

4. 祂领受圣洗时已三十岁，已具有师傅的完全年龄，祂来到耶路撒冷，为能恰当地被众人认作师傅。因为祂并非表里不一，像那些声称祂只是外表是人的所断言的那样；而是祂是什么，就显出是什么。所以祂既是师傅，也拥有师傅的年龄，不轻视或逃避任何人性处境，也不在自己身上废除祂为人性所制定的法律，而是藉着与自己相应的每一时期，圣化每一个年龄。因为祂来是要藉着自己拯救所有的人——我说的是所有藉着祂重生归于天主的人——婴孩、幼童、少年、青年和老人。因此，祂经历了每一个年龄：为婴孩而成为婴孩，以此圣化婴孩；为幼童而成为幼童，以此圣化这个年龄的人，同时为他们成为虔敬、公义和顺服的榜样；为少年而成为少年，成为少年的榜样，从而为他们圣化归于主。同样，祂也为老人而成为老人，为要成为众人的完善的师傅，不仅在于阐明真理，也在于年龄，同时圣化年长者，也成为他们的榜样。最后，祂终于走向死亡本身，为要成为\BibleQuote{死者中的首生者，使祂在一切事上居首位，}\BibleRef{哥 1:18}成为\OriginalTerm{生命之主}{Prince of life}\BibleRef{宗 3:15}，祂先于万有，且走在万有之前。\par

5. 然而，他们为了支持自己关于经上所记\BibleQuote{宣布上主恩慈之年}的虚假见解，就主张主只传道一年，然后在第十二个月受难。他们这样说，是于己不利地忘却了，摧毁了祂的全部工作，剥夺了祂比任何年龄都更必要、更可敬的那个年龄；我指的是那个更成熟的年龄，在此期间，祂作为教师也超越了所有其他人。因为如果祂不教导，怎么能有门徒呢？如果祂没有达到师傅的年龄，又怎么能教导呢？因为祂来领受圣洗时，尚未满三十岁，而是开始约三十岁（路加记载了祂的岁数，如此表达：\BibleQuote{耶稣开始传教时，大约三十岁，}\BibleRef{路 3:23}当祂来领受圣洗时）；并且，〔照他们所说，〕从领受圣洗算起，祂只传道一年。祂在满三十岁时受难，实际上仍是一个年轻人，根本没有达到成熟的年龄。现在，每个人都承认，人生的第一阶段包括三十年，并延续到第四十年；但从四十岁到五十岁，人开始向老年衰迈，而我们的主在履行教师职务时，已经处于这个年龄，正如福音和所有长者所见证的；那些在\Place{亚细亚}与主之门徒\Person{若望}交往的人，〔证实〕\Person{若望}向他们传达了这一信息。而且，他一直与他们在一起，直到\Person{图拉真}时期。他们中的一些人不仅见过\Person{若望}，也见过其他宗徒，并从他们那里听到同样的记述，为这陈述作证。那么，我们应当相信谁呢？是这些人，还是从未见过宗徒、甚至做梦也从未触及丝毫宗徒踪影的\Person{仆托肋买}？\par

6. 然而，除此之外，当时与主耶稣基督辩论的那些犹太人自己也最清楚地表明了同样的事。因为主对他们说：\BibleQuote{你们的祖宗亚巴郎欢欣喜乐地要看见我的日子，他看见了，且高兴起来。}他们就回答他说：\BibleQuote{你还没有五十岁，就见过亚巴郎吗？}\BibleRef{若 8:56} 这样的话正适用于一位已过四十岁、尚未到达五十岁，但离后者不远的人。但若有人只三十岁，无疑会被说：\Quote{你还没有四十岁。}因为那些想要证明他说谎的人，绝不会把他所达到的年龄远远超过他们所见到的；而是提了一个接近他实际年龄的数目，无论他们是确实从官方登记册得知，还是仅凭观察推测他四十多岁，且肯定绝非三十岁。因为认为他们错了二十年，想要证明他比亚巴郎的时代更年轻，是完全不合理的。因为他们所见，他们也说了；他们看见的并非幻影，而是真实血肉之躯。他当时离五十岁并不远；因此他们对他说：\BibleQuote{你还没有五十岁，就见过亚巴郎吗？}所以，他传道不止一年，也不是在那年的第十二个月受难。因为从三十岁到五十岁之间的时期，决不能被视为\Emph{一}年，除非在他们那些\OriginalTerm{永世}{Æons}中，对于坐在\OriginalTerm{圆满}{Pleroma}中与\OriginalTerm{深渊}{Bythus}同列者，分配了这么长的年岁；诗人荷马曾提到这些受造物，无疑受了他们错误体系的\Quote{母亲}的启示：—\par

\OriginalTerm{}{\GreekText{Οἱ} \GreekText{δὲ} \GreekText{θεοὶ} \GreekText{πὰρ} \GreekText{Ζηνὶ} \GreekText{καθήμενοι} \GreekText{ἠγορόωντο}}\newline{}\OriginalTerm{}{\GreekText{Χρυσέῳ} \GreekText{ἐν} \GreekText{δαπέδῳ}:}\par

我们可以如下意译：\par

\Quote{众神围绕，朱庇特主持会议，\newline{}在金碧辉煌的地板上交谈。}\par

\SourceRecord{Translated by Alexander Roberts and William Rambaut.；Ante-Nicene Fathers；Vol. 1.；Edited by Alexander Roberts, James Donaldson, and A. Cleveland Coxe.；Buffalo, NY: Christian Literature Publishing Co.,；1885.；<http://www.newadvent.org/fathers/0103222.htm>.}

% structure: p0001 source=h1 target=section reason=page-title

\section{\WorkTitle{驳斥异端}（第二卷，第二十三章）}

1. 此外，关于那位患血漏的妇女，她触摸了主的衣边而得以痊愈的事例，他们的无知昭然若揭。因为他们声称，通过此妇人所展现的是那承受苦难、流向无限的第十二能力，即第十二永世\OriginalTerm{永世}{Æon}。首先，正如我所示，按照他们自己的体系，那并非第十二永世。即使暂且认可这一点，即共有十二永世，其中十一个据说一直未受苦难，而第十二个承受了苦难；但另一方面，那妇女在第十二年得到治愈，显然她曾患病十一年，而在第十二年痊愈。如果他们说，十一个永世都经历了苦难，而第十二个得到治愈，那么说这妇女是这些永世的预像或许合理。然而，她患病十一年，未得任何医治，而在第十二年痊愈，这样她如何能成为第十二永世的预像呢？因为按照假设，十一个永世根本未受苦，只有第十二个参与了苦难。诚然，预像与象征有时在质料和实体上与被象征的真理有所不同，但它应当在总体形式和特征上保持相似，借此以现有事物预示那将要到来的事物。\par

2. 不仅在这妇人的事例中提及了她病患的年数（他们断言这符合他们的虚构），而且，看！还有另一位妇人，同样患病十八年后被治愈；主论及她说：\Quote{难道不应该让这位亚巴郎的女儿，撒殚捆绑了她十八年，在安息日得以解脱吗？}\BibleRef{路 13:16} 如果前者是受苦的第十二永世的预像，那么后者也应该是受苦的第十八永世的预像。但他们无法坚持这一点；否则他们的首要和原始的\OriginalTerm{八元}{Ogdoad}将被包括在共同受苦的永世数目中。此外，还有另一个人，主治愈了他，他患病三十八年\BibleRef{若 5:5}；因此，他们应当断言那占据第三十八位的永世受苦了。因为如果他们声称主所行的事是\OriginalTerm{圆满}{Pleroma}内所发生之事的预像，那么这预像应当贯穿始终。但他们既不能将医治十八年后痊愈的妇人的事例，也不能将医治三十八年后的男子的病例，纳入他们虚构的体系。在任何方面，宣称救主在某些情况下保持了预像，而在另一些情况下却没有，都是荒谬且自相矛盾的。因此，那个妇人的预像已被证明与他们的永世体系毫无类比关系。\par

\SourceRecord{Translated by Alexander Roberts and William Rambaut.；Ante-Nicene Fathers；Vol. 1.；Edited by Alexander Roberts, James Donaldson, and A. Cleveland Coxe.；Buffalo, NY: Christian Literature Publishing Co.,；1885.；<http://www.newadvent.org/fathers/0103223.htm>.}

% structure: p0001 source=h1 target=section reason=page-title

\section{驳斥异端（第二卷，第二十四章）}

1. 这一点更充分证明他们的观点虚假，其虚构体系站不住脚：他们试图通过数字和名字的音节，有时也通过音节的字母，甚至通过希腊人惯例中不同字母所含的数字，来为这体系提出证据——〔我说，〕这一点最清楚地表明他们的失败或混淆，以及他们〔自称〕的知识的虚妄和乖谬。因为他们将属于另一种语言的名字\Emph{耶稣}引入希腊人的数字体系，有时称它为\Quote{厄庇塞蒙}，因为它有六个字母；有时又称它为\Quote{奥格多阿斯的圆满}，因为它包含八百八十八之数。但他的对应希腊名\Quote{索忒耳}，即\Emph{救主}，却因无论在数值上还是字母上都不合他们的体系而被他们默然忽略。然而，如果他们果真按照圣父的预旨，视主的名通过数值和字母指示圆满中的数字，那么\Emph{\OriginalTerm{索忒耳}{Soter}}作为希腊名，理应通过其字母和数值（因它是希腊文）展现圆满的奥秘。但事实并非如此，因为它是一个五字母词，数值为一千四百零八。这些与他们的圆满毫不相符；因此，他们对圆满中运作的叙述不可能是真实的。\par

2. 此外，\Emph{耶稣}是希伯来本族语中的词，据他们中的学者所言，包含两个半字母，意指那包含天地的\Emph{主}；因为在古希伯来语中，\Emph{耶稣}意为\Quote{天}，而\Quote{地}则由\Emph{\OriginalTerm{苏拉乌塞尔}{sura usser}}一词表达。因此，包含天地的词就是\Emph{耶稣}。那么，他们对\Emph{\OriginalTerm{厄庇塞蒙}{Episemon}}的解释是虚假的，其数字计算也明显被推翻。因为在他们自己的语言中，\Emph{\OriginalTerm{索忒耳}{Soter}}是一个五字母的希腊词；但在希伯来语中，\Emph{耶稣}只包含两个半字母。他们算出的总数即八百八十八便落空。而且，希伯来字母与希腊字母在数字上并不对应，尽管希伯来字母作为更古老、不变的字母，本应支持与名字相关的计算。因为这些古老、原始、通常被称为\Emph{神圣的}希伯来字母共有十个（但通过十五个来书写），最后一个字母与第一个相连。这样，有些字母按自然顺序书写，一如我们；但其他的则反向书写，从右向左，即倒写字母。\Emph{基督}这名也应能与他们圆满的众多永世相符，因为据他们称，祂是为了建立和纠正他们的圆满而产生的。同样，\Person{圣父}也应通过字母和数值包含祂所产生的那些永世的数目；\OriginalTerm{深渊}{Bythus}亦然，\OriginalTerm{莫诺格尼斯}{Monogenes}也不应例外；但最首要的是那名超越一切的名，即天主在希伯来语中所称的\Emph{\OriginalTerm{巴路克}{Baruch}}，它也包含两个半字母。因此，从这些在希伯来语和希腊语中更重要的名字无论是字母数还是计算出的数值都不符合他们的体系这一事实，他们关于其余名字的强行计算便显而易见。\par

3. 因为，他们从法律中挑选出与他们体系所采纳的数字相符的事物，就这样强行地试图获取其有效性的证据。但是，如果他们的\OriginalTerm{母亲}{Mother}或\OriginalTerm{救主}{Saviour}真的有意通过\OriginalTerm{造物主}{Demiurge}来呈现那些在\OriginalTerm{圆满}{Pleroma}中的事物的预象，那么他们本该注意让预象出现在更为精确对应且更为神圣的事物中；特别是，在约柜的事上——因这柜，整个会幕才被建造起来——约柜的尺寸是这样：它的长\BibleRef{出 25:10}是二肘半，宽一肘半，高一肘半。然而这样的肘数完全不符合他们的体系，但通过它，预象本应比任何其他事物都更清楚地被呈现出来。赎罪盖\BibleRef{出 25:17}同样也完全不符合他们的解释。此外，供饼桌\BibleRef{出 25:23}长二肘，高一肘半。这些物件放置在至圣所前，然而其中没有一个数字大到足以包含\OriginalTerm{四元组}{Tetrad}或\OriginalTerm{八元组}{Ogdoad}或他们\OriginalTerm{圆满}{Pleroma}其余部分的指示。那么灯台\BibleRef{出 25:31}呢？它本有七枝、七盏灯。然而，若这些是按照预象制造的，它本该有八枝和相应数量的灯，以符合原初\OriginalTerm{八元组}{Ogdoad}的预象——这八元组在\OriginalTerm{永世}{Æons}中照耀得最为卓越，并光照整个\OriginalTerm{圆满}{Pleroma}。他们仔细列举了幕幔\BibleRef{出 26:1}共十幅，宣称这些是十个\OriginalTerm{永世}{Æons}的预象；却忘记了计算那十一张\BibleRef{出 26:7}皮制的罩棚。他们也没有量这些幕幔的尺寸：每幅幕幔\BibleRef{出 26:2}长二十八肘。他们提出柱子的长度为十肘，以指向\OriginalTerm{十元组}{Decad}的\OriginalTerm{永世}{Æons}。\Quote{但每根柱子的宽度是一肘半；}\BibleRef{出 26:16}这一点他们却不解释，正如他们也不解释柱子和横闩的总数一样，因为那不符合他们的论证。但是，那祝圣了整个会幕的傅油\BibleRef{出 26:26}又如何呢？或许是救主忽略了它，或是当他们的母亲睡觉时，造物主自行指示了其重量；正因如此，它与他们的\OriginalTerm{圆满}{Pleroma}不一致：这傅油\BibleRef{出 30:23}含有五百舍客勒的没药、五百舍客勒的肉桂、二百五十舍客勒的桂枝、二百五十舍客勒的菖蒲，外加橄榄油，因此它由五种成分组成。同样地，香\BibleRef{出 30:34}由苏合香、香螺、白松香、薄荷和乳香组成，所有这些，无论是其混合物还是重量，都与他们的论点毫不相符。因此，声称预象并没有保存在法律中那些崇高且更为威严的诫命中，却在其他点上，当任何数字与他们的主张相符时，便断言那是\OriginalTerm{圆满}{Pleroma}中事物的预象，这是不合理且完全荒谬的。而事实是，每一个数字在圣经中都极多种类地出现，以至于任何人若愿意，他不仅可以从中形成一个\OriginalTerm{八元组}{Ogdoad}、一个\OriginalTerm{十元组}{Decad}和一个\OriginalTerm{十二元组}{Duodecad}，而且可以形成任意数目的数字，然后声称这是他自己所虚构的谬误体系的预象。\par

4. 但如下的圣经经文将证明这一点是真实的，即那被称为\Emph{五}的数字，虽然与他们的论证毫不相符，不与他们的体系调和，也不适合作为\OriginalTerm{圆满}{Pleroma}中事物的预像，[却广泛存在]。\OriginalTerm{救主}{Soter}是一个由五个字母组成的名字；\OriginalTerm{父}{Pater}也包含五个字母；\OriginalTerm{爱}{Agape}也由五个字母构成。我们的主在祝福了五个饼后，用它们喂饱了五千人。五个童女\BibleRef{玛 25:2}被主称为明智的；同样，五个被称为糊涂的。再者，当主从父那里获得见证时\BibleRef{玛 17:1}，据说有五人与主同在，即伯多禄、雅各伯、若望、梅瑟和厄里亚。主也作为第五个人进入了已死女孩的房间，并再次使她复活；因为[圣经]说：\Quote{祂不许任何人进去，只有伯多禄、雅各伯、以及那女孩的父母。}\BibleRef{路 8:51}地狱中的富人\BibleRef{路 16:28}自称有五个兄弟，他希望从死者中复活的一位到他们那里去。主命令那瘫痪的人回家去的池子，有五个廊檐。十字架的形状本身也有五个端点：两个长度方向，两个宽度方向，一个在中间，被钉子钉住的人就靠在那[最后一个]上面。我们每只手有五个手指；我们也有五种感官；我们的内脏也可以算作五个：即心脏、肝脏、肺、脾和肾。此外，整个人也可以分为这个[部分]数目——头、胸、腹、大腿和脚。人类经历五个年龄：首先婴孩期，然后童年，然后青年，然后壮年，然后老年。梅瑟将法律传给百姓，共五卷书。他从天主领受的每块石版上包含五条诫命。遮盖至圣所的幔子\BibleRef{出 26:37}有五个柱子。全燔祭台宽度也是五肘。在旷野中选定了五位司祭——即亚郎\BibleRef{出 28:1}、纳达布、阿彼胡、厄肋阿匝尔、依塔玛尔。厄弗得、胸牌以及其他司祭衣饰是由五种材料\BibleRef{出 28:5}制成的；因为它们包含了金、蓝色、紫色、朱红色和细麻。还有五位阿摩黎人的王\BibleRef{苏 10:17}被农的儿子若苏厄关在洞里，并吩咐百姓踩踏他们的头。事实上，任何人都可以收集数以千计的其他相同事物，无论是关于这个数字还是他选择确定的任何其他数字，无论是从圣经中，还是从他观察到的自然作品中。但是，尽管如此，我们并不因此断言在\OriginalTerm{造物主}{Demiurge}之上有五个\OriginalTerm{永世}{Æons}；我们也不奉\OriginalTerm{五重}{Pentad}为神圣之物；我们也不试图通过那种徒劳的劳作去建立站不住脚的东西或[他们所沉溺的]狂言；我们也不悖逆地迫使一个由天主充分适应[所要服务的目的]的受造物改变自身，成为并不真实存在之物的预像；我们也不寻求提出不敬且可憎的教义，其识破和推翻对所有有才智的人来说都是容易的。\par

5. 因为谁能同意他们所说的，一年只有三百六十五天，以便有十二个月，每月三十天，仿照十二个\OriginalTerm{永世}{Æons}的样式，而实际上这样式与[原型]完全不符？因为在前一种情况中，每个永世是整个\OriginalTerm{圆满}{Pleroma}的三十分之一，而在后一种情况中，他们宣称一个月是一年的十二分之一。如果确实一年被分为三十部分，一个月被分为十二部分，那么就可能被认为为他们的虚构体系找到了合适的样式。但相反，实际情况是，他们的圆满被分为三十部分，其中一部分又被分为十二；而整个一年被分为十二部分，其中一部分又被分为三十。因此，救主在使一个月成为整个圆满的样式，而使一年只成为圆满中那\OriginalTerm{十二重}{Duodecad}的样式时，是不明智的；因为更适宜的做法是将一年分为三十部分，正如整个圆满被分，而将一个月分为十二部分，正如永世在他们的圆满中一样。此外，他们将整个圆满分为三部分——即\OriginalTerm{八重}{Ogdoad}、\OriginalTerm{十重}{Decad}和十二重。但我们的年分为四部分——即春、夏、秋、冬。此外，他们主张是\OriginalTerm{三十重}{Triacontad}样式的月份也并非精确地由三十天组成，而是有的多些有的少些，因为还有五天作为剩余。白天也并非总是精确地由十二小时组成，而是从九小时增加到十五小时，然后又从十五小时减到九小时。因此，不能认为每月三十天的月份是为了[预像]永世而形成的；因为那样的话，它们就应该精确地由三十天组成；同样，月份的日数也不能是为了通过十二小时来象征十二个永世；因为那样的话，它们就应该总是精确地由十二小时组成。\par

6. 此外，至于他们称物质实体为\Quote{左边}，并主张这些在左边的事物必然陷入败坏，而他们又断言救主来到迷失的羊那里，是为了将其转移到\Quote{右边}，即那九十九只安然无恙、未失丧、仍留在圈内的羊，但这些羊却属于左边；如此则他们必须承认，安息的享受并不意味着救恩。而那个不同样拥有相同数目的事物，他们将被强迫承认为属于左边，即属于败坏。这个希腊词\OriginalTerm{爱}{Agape}，按照希腊字母（他们以此进行推算）具有数值\Emph{九十三}，同样被归为左边的安息之所。\OriginalTerm{真理}{Aletheia}也同样按照上述原则具有数值\Emph{六十四}，存在于物质实体中。最终，他们将被迫承认，所有那些达不到一百数值、仅包含左边手指所计之数的神圣名号，都是可朽坏的和物质的。\par

\SourceRecord{Translated by Alexander Roberts and William Rambaut.；Ante-Nicene Fathers；Vol. 1.；Edited by Alexander Roberts, James Donaldson, and A. Cleveland Coxe.；Buffalo, NY: Christian Literature Publishing Co.,；1885.；<http://www.newadvent.org/fathers/0103224.htm>.}

% structure: p0001 source=h1 target=section reason=page-title

\section{《驳斥异端》（第二卷，第二十五章）}

1. 如果有人对这些事回答说：那么，名字的安排、宗徒的拣选、主的作为、受造物的秩序，岂非毫无意义、出于偶然吗？——我们回答他们：当然不是；相反，一切都是天主以大智慧、大细心，明确地造成，各按其用而配合、预备妥当。祂的圣言形成了古代的事和最近时期的事。人不应该把那些事与\Emph{三十}这个数字相连，而应该使它们与实际存在或正确理性相和谐。也不应该通过数字、音节和字母来寻求探究有关天主的事。因为这是一种不确定的方法，因为数字系统多种多样、各不相同，如今任何人都可以以类似方式编造各种假设，以致他们可以从这些理论本身推导出反对真理的论据，因为这些理论可以在许多不同方向上被扭转。相反，他们应该使数字本身以及已经被造的事物适应于摆在面前的真正理论。因为系统并非出自数字，而是数字出自系统；天主并非源于受造之物，而是受造之物源于天主。因为万有皆出自同一天主。\par

2. 然而，既然受造之物多样性且为数众多，它们确实整个地良好配合、适应整个受造界；然而个别来看，它们又相互对立、不协调，就像竖琴的声音，由许多对立音符组成，通过每个音符之间的间隔，产生出一首连贯的旋律。因此，爱好真理的人不应被每个音符之间的间隔所欺骗，也不应以爲一个音符属于一位艺术家和作者，另一个属于另一位，或者一个调配高音弦，另一个低音弦，第三个中音弦；而应坚持是同一人形成了整体，以显明整个作品和智慧的判断、良善和技艺。那些聆听旋律的人，应该赞美并颂扬艺术家，钦佩某些音符的张力，注意另一些音符的柔和，感受介乎两者之间的其他音符的声响，并考虑另一些音符的特殊性质，以便探究每个音符的目的及其多样性的原因，始终不放弃我们的准则，既不放弃那位艺术家，也不放弃形成万有的唯一天主的信仰，也不亵渎我们的造物主。\par

3. 然而，如果有人无法发现所有成为探究对象的事物的原因，让他思量：人无限低于天主；他仅部分地领受了恩宠，尚未与他的造物主相等或相似；而且，他不能像天主那样体验或形成对万有的概念；正如那今日才被形成、才被创造开始的人，低于那非受造的、永恒不变者一样，同样地，他在知识和探究万有原因的能力上低于那创造他的。因为，人啊，你并非非受造者，也未曾与天主永远共存，像祂的圣言那样；而是如今，因着祂卓越的良善，你领受了你受造的开始，你逐渐从圣言学习那创造你的天主的管理计划。\par

4. 因此，要保守你知识的正确秩序，不要因为不认识真正良善的事物而妄图超越天主本身，因为祂不能被超越；也不要在造物主之上追求任何一位，因为你找不到这样的。因为你的创造者不能被限制在界限之内；尽管你应测量这整个宇宙，穿越祂的全部受造物，并考虑其深度、高度和长度，你也不能在圣父本身之上想象任何别的。因为你不能完全想出祂，反而沉溺于与你的本性相反的思绪中，你就显明自己是愚昧的；如果你坚持这样下去，你会陷入彻底疯狂，同时自以为比你的造物主更高更大，并幻想你能穿透到祂的领域之外。\par

\SourceRecord{Translated by Alexander Roberts and William Rambaut.；Ante-Nicene Fathers；Vol. 1.；Edited by Alexander Roberts, James Donaldson, and A. Cleveland Coxe.；Buffalo, NY: Christian Literature Publishing Co.,；1885.；<http://www.newadvent.org/fathers/0103225.htm>.}

% structure: p0001 source=h1 target=section reason=page-title

\section{\WorkTitle{驳斥异端（第二卷 第二十六章）}}

1. 因此，做一个单纯的、未受过教育的人，藉着爱德而接近天主，比自认为有学问、有技巧，却成为亵渎自己所信天主的人（因为他们捏造了另一位神为父），更好，也更有益。为此，保禄曾喊道：\Quote{知识使人自高自大，爱德却立人：}\BibleRef{格前 8:1} 他这样说，并不是要反对对天主的真知识，因为那样他就是在指责自己；而是因为他知道有些人，因着自以为有知识而自高自大，远离了天主的爱德，并以为自己已是成全的，为此他们提出一个不全成的创造主，以期终止这种因知识而生的骄傲。他说：\Quote{知识使人自高自大，爱德却立人。}再没有比这更大的狂妄，即一个人以为自己比那创造、形成他，并赐给他生命气息，且命令这存在本身的那一位更优越、更成全。因此，如我所说，与其因这种知识而自高自大，远离那作为人之生命的爱德，倒不如对受造物中任何一件东西被造的缘由毫无所知，而只信靠天主，并居住在祂的爱德中；只寻求那为我们被钉十字架的天主圣子耶稣基督的知识，而不藉着刁钻的问题和咬文嚼字陷入不敬虔。\par

2. 如果一个人，因着这类尝试而逐渐自高自大，因为主曾说：\Quote{连你们的头发，也都一一数过了。}\BibleRef{玛 10:30} 就着手探究每个人头上的头发数目，并试图找出一个人有这么多、另一个人有那么多（因为并非所有人都有相等的数目，而是成千上万的人有着千差万别的数目，因为有的人头大，有的人头小，有的人头发浓密，有的人稀疏，有的人几乎无发）的原因，然后那些自以为发现了头发数目的人，便尽力将他们所臆想的内容用于赞扬自己的派别，那将会怎样？再者，如果某人因为福音中这句话：\Quote{两只麻雀不是卖一个铜钱吗？然而没有一只落在地上，不经过你们父的许可。}\BibleRef{玛 10:29} 就趁机计算每日捕捉的麻雀数目，无论是在全世界还是在某个特定地区，并探究为什么昨天捕捉了这么多，前天这么多，今天又是这么多，然后将麻雀的数目与他（自己的）假设联系起来，他岂不彻底误入歧途，并将那些赞同他的人驱向完全的疯狂吗？因为人总想在这样的事上被认为是发现了比自己的老师更非凡的东西。\par

3. 但若有人问我们，是否一切已造和受造之物的每一个数目都为天主所知，并且每一个这样的数目都依照祂的眷顾而获得了它所包含的特定量；当我们同意并承认：任何已存在、正存在或将要存在的事物，无一能逃脱天主的认知，而是藉着祂的眷顾，每一事物都获得了它的本性、等级、数目和特定的量，没有任何事物是徒然或偶然地存在或产生的，而是具有极其适切性，并运用了超越的知识，并且能够区分并产生如此体系之适当原因，是一种令人赞叹且真正神圣的理智——那么，如果某人，我再说，在我们应允并同意这一点之后，着手计算地上的沙粒和卵石，甚至海中的波浪和天上的星辰，并试图想出其自以为发现的数目的原因，他的劳作岂不徒劳？这样的一个人岂不理所当然地被一切有常识的人宣布为疯狂和缺乏理性？他越是比他人专注于这类问题，越是自以为比他人发现得多，称他人为无技巧、无知、属血气的，因为他们不从事他这种无用的劳作，他就越是（实际上）疯狂、愚钝，仿佛被雷击一般，因为他丝毫没有承认自己低于天主；相反，藉着他自以为发现的知识，他改变了天主本身，并将自己的意见抬高到创造主的伟大之上。\par

\SourceRecord{Translated by Alexander Roberts and William Rambaut.；Ante-Nicene Fathers；Vol. 1.；Edited by Alexander Roberts, James Donaldson, and A. Cleveland Coxe.；Buffalo, NY: Christian Literature Publishing Co.,；1885.；<http://www.newadvent.org/fathers/0103226.htm>.}

% structure: p0001 source=h1 target=section reason=page-title

\section{\WorkTitle{反对异端（第二卷，第二十七章）}}

1. 一个健全的头脑，一个不使其拥有者陷入危险、致力于虔诚与热爱真理的头脑，会热切默想那些天主置于人类能力之内、交给我们认识的事物，并通过每日研习而在对它们的认识上取得进步，使对它们的认识变得容易。这些事物就是清楚呈现在我们眼前、在圣经中以明确无误的措辞清晰阐述的那些事物。因此，比喻不应被牵强附会地解释成含糊的表述。因为，若不如此，解释者便可以无风险地解释它们，比喻将会得到所有人同样的解释，真理的整体保持完整，其各部分和谐配合，而无任何冲突。但是，将不清晰或不明显的表述应用于比喻的解释，就像每个人凭自己的好恶所发现的那样，\Quote{是荒谬的。}因为这样一来，没有人将拥有真理的准则；相反，随着解释比喻的人数多寡，将会出现各种彼此对立的真理体系，宣扬相互冲突的教义，就像外邦哲学家中间流行的那些问题一样。\par

2. 按照这种程序，人将永远探求而永无发现，因为他已抛弃了发现的方法。当新郎\BibleRef{玛 25:5}来到时，那油灯没有修剪、没有以稳定光明的亮度燃烧的人，就被列在那些使比喻的解释变得晦暗的人当中；他们离弃了那位藉着明白的宣告白白赐予一切来者恩赐的主，并被排除在新郎的婚房之外。既然全部圣经、先知书和福音，虽然并非所有人都相信，却能被所有人清楚、明确、和谐地理解，而且它们宣告只有一位天主，排除所有其他，祂藉着祂的圣言形成了万物，无论可见不可见、天上地上、水中或地下——正如我从圣经的原文所证明的——并且我们所从属的这受造界本身，通过呈现在我们眼前的事物，见证了一位造物主和治理者——那么，那些人对如此清楚的证明视而不见，不肯看见那宣告的光明，而给自己戴上镣铐，每个人都以为藉着他们对比喻的晦暗解释，找到了一位他自己的天主，这样的人可真显得愚昧。因为，在任何部分的圣经中，没有什么公开、明确、毫无争议地提及那些持相反意见者所设想的那位圣父——他们自己证明了这一点，因为他们坚持说救主私下教导这些事，不是教给所有人，而只教给祂的某些门徒，这些门徒能理解祂，并能通过论证、谜语和比喻来领悟祂的意图。他们最终得出这样的结论：他们坚持有一位存在被宣告为天主，另一位则为圣父，那位圣父是通过比喻和谜语而被揭示的。\par

3. 但是，既然比喻容许多种解释，哪个热爱真理的人不会承认，他们放弃确定无疑、不可辩驳、真实的事物，而主张该从比喻中寻找天主，这是那些甘愿冒险、行事如同失去理智的人的行为？这样的行为不正是没有将房屋建造在稳固、坚强、位于开阔处的磐石\BibleRef{玛 7:25}上，而是建造在流沙上吗？因此，这样的建筑很容易倒塌。\par

\SourceRecord{Translated by Alexander Roberts and William Rambaut.；Ante-Nicene Fathers；Vol. 1.；Edited by Alexander Roberts, James Donaldson, and A. Cleveland Coxe.；Buffalo, NY: Christian Literature Publishing Co.,；1885.；<http://www.newadvent.org/fathers/0103227.htm>.}

% structure: p0001 source=h1 target=section reason=page-title

\section{\WorkTitle{驳斥异端（第二卷，第二十八章）}}

1. 我们既以真理本身为准则，又清楚地拥有关于天主的见证，就不应当为了追求对问题作出众多不同的回答，而抛弃对天主稳固而真实的知识。更为适宜的是：我们应当以这样的方式来指导我们的探究，在追查生活天主的奥迹和管理时操练自己，并增加对那为我们成就了如此伟大之事的他的爱德，却绝不从这信仰中坠落；这信仰最清楚地宣告：唯有这一位才是真正的天主和圣父，他创造了这个世界，塑造了人，赋予受造物增长的能力，并把人从较小的事物召往在他面前那更大的事物，正如他把受孕在胎中的婴儿带到阳光之下，又在赐予麦穗以饱满的力量后，把麦子收进仓里。原是同一个造物主塑造了子宫并创造了太阳；同一个主养育了麦秆，增加并繁衍了麦子，并预备了谷仓。\par

2. 然而，如果我们不能为圣经中所有那些作为探究主题的事找到解释，我们也不可因此寻求其他神祇，离弃那真实存在的一位。因为这是极不虔敬的事。我们应当把这类性质的事留给那创造了我们的天主，并且坚信圣经确实是完美的，因为它是借天主圣言和他的圣神所说出的；而我们，既然劣于且后于天主圣言和他的圣神，因而欠缺对他奥迹的知识。如果我们在属灵和属天的事上，以及那些需要藉启示向我们显明的事上如此，那并不足为奇，因为即使是我们脚下许多的事（我说的是那些属于这个世界、我们亲手触摸、亲眼看见、密切接触的事）也超越我们的知识，以至连这些我们也必须留给天主。因为天主应在一切知识上超越我们。例如，我们若试图解释\Place{尼罗河}泛滥的原因，我们可以就这个问题说很多话，不论合理与否；但关于它的真实、确定且无可争辩的知识，只属于天主。再者，鸟类的栖所——我指的是那些春天飞到我们这里、秋天又飞走的鸟——虽然是与这个世界有关的事，却超出我们的知识。我们又能如何解释海洋的潮汐涨落呢？尽管每个人都承认这些现象必有某个原因。对于海洋之外那些事物的性质，我们又能说什么？此外，关于雨、闪电、雷、云彩的聚集、水汽、风的爆发以及诸如此类的事，我们又能说什么？关于雪、冰雹以及其他类似之物的仓库，我们又能知道什么？关于云彩形成的条件，或天空中水汽的真实性质，我们知道什么？关于月亮圆缺的原因，或各种水、金属、石头等物性质不同的原因，我们又知道什么？在这一切问题上，我们固然可以在探究其原因时说很多话，但唯独那创造它们的天主能宣告它们的事实。\par

3. 因此，如果连对于受造之物，有些事的知识只属于天主，有些则属于我们自己的知识范围，那么，对于我们在圣经（是全然属灵的）中探究的那些事，我们能藉天主的恩宠解释其中一些，而将其余的交在天主手中——不仅在此世，也在来世，使天主永远教导，人永远学习天主所教导的事——我们有什么可抱怨的呢？正如宗徒在此事上所说：当其他事物被废除时，这三样却存留，即\Quote{信德、望德和爱德}。\BibleRef{格前 13:13} 因为那关于我们主的信德，坚定不移地存留，肯定只有一位真天主，并且我们应当永远真实地爱他，因为他唯独是我们的圣父；同时我们期望不断从天主那里领受更多，并向他学习，因为他美善，拥有无尽的财富、无终的王国和取之不竭的教导。因此，倘若我们按我已陈述的准则，将一些问题留在天主手中，我们就能保全信德不受损伤，并安全无虞；而天主所赐予我们的一切圣经，在我们看来将完全一致；比喻将与那些完全清晰的段落和谐；意义明了的陈述将用以解释比喻；藉圣经中众多不同的言语，我们内里将听到一首和谐的旋律，赞美那创造万有的天主。例如，如果有人问：\Quote{天主在创造世界之前做什么？}我们回答说，这类问题的答案在于天主自己。因为圣经教导我们，这个世界在天主手中完美地形成，在时间中有一个开始；但圣经没有启示我们天主在这之前做了什么。因此，那问题的答案仍在天主那里，我们不宜试图提出愚昧、鲁莽和亵渎的假设作为回答；这样，某人若自以为发现了物质的起源，实际上却否定了那创造万有的天主本身。\par

4. 试想，你们这些编造此类意见的人：既然只有圣父本身被称为天主，祂真实存在，而你们却称祂为\OriginalTerm{德穆革}{Demiurge}；此外，圣经也只承认祂是天主；再者，主也亲口承认祂是自己的唯一父亲，并不认识别的，正如我将从祂的话中表明的那样——当你们称这一位为缺陷之果、无知之裔，并描述祂对超越自身之事一无所知，以及你们对祂所作的各种其他指摘时——请想想你们对那位真实的天主所犯的可怕亵渎。你们似乎郑重而诚实地声称信天主，但既然你们完全不能揭示任何别的天主，你们就宣称你们所信奉的这一位是缺陷之果和无知之裔。这种盲目和愚昧之言源于你们不为天主保留任何奥秘，反而想宣告天主本身及其思想、圣言、生命和基督的诞生与产生。你们从人的经验中构想这些，不明白（正如我先前所说）对于人是复合体而言，可以这样谈论人的理智和思想：说\OriginalTerm{思想}{Ennœa}来自\OriginalTerm{理智}{sensus}，\OriginalTerm{意图}{enthymesis}又来自思想，\OriginalTerm{言语}{logos}来自意图——但这是哪种言语？因为希腊人有一种是思维的理性原则，另一种是表达思想的工具——并且人有时安静沉默，有时说话行动。但天主是全部理智、全部理性、全部行动的神、全部光明，永远同一不变（正如我们默想天主所当有的，并从圣经中所学到的），这样的情感和分割不能适当地归于祂。因为我们的舌头是属血肉的，不足以服务于人类理智的敏捷（因理智是属精神的），因此我们的言语被约束在我们内，不能像被理智构思的那样立刻表达出来，而是通过连续的努力发出，正如舌头所能服务的那样。\par

5. 但天主是全部理智、全部圣言，祂完全说出祂所想，也完全想祂所说。因为祂的思想是圣言，圣言是理智，而包含万有的理智就是圣父本身。因此，凡谈论天主理智并为其赋予独立起源的人，就是宣称祂是复合体，仿佛天主是一回事，原始理智是另一回事。同样，关于圣言，若有人将祂置于从圣父而生的第三位，他就不知道祂的伟大，从而将圣言远远地与天主分离。先知论及祂说：\BibleQuote{谁能述说祂的世代？}\BibleRef{依 53:8}。而你们却企图阐述祂从圣父而生的世代，并将人的言语（通过舌头产生）的产生方式转移到天主的圣言上，因此你们自己被正义地揭露为既不知人事也不识神事的人。\par

6. 但你们因自己的智慧而过度膨胀，狂妄地声称你们知晓天主不可言喻的奥秘；而连主——天主子本身——也承认唯独圣父知道审判的那日那时，祂明确地说：\BibleQuote{至于那日子和那时辰，没有人知道，连子也不知道，只有父知道。}如果连子都不以将那个日子的知识唯独归于父为耻，而是宣告了事实，那么我们也不应以将可能遇到的更重大难题保留给天主为耻。因为\BibleQuote{没有人高过自己的师傅。}\BibleRef{玛 10:24} 因此，若有人对我们说：\Quote{那么子是如何由父所生的？}我们回答他：没有人理解那个生产、或生育、或呼召、或启示，或任何可能用来描述祂的世代（这世代是完全不可描述的）的名称。无论是\OriginalTerm{瓦伦提努}{Valentinus}、\OriginalTerm{玛尔基翁}{Marcion}、\OriginalTerm{撒杜尔尼努}{Saturninus}、\OriginalTerm{巴西里得}{Basilides}，还是天使、总领天使、宰制者、掌权者，都不知道这个奥秘，唯有生者之父和被生之子知道。既然祂的世代是不可言喻的，那些企图阐述世代和生产的人必定神志不清，因为他们承担了描述不可描述之事。因为人人都明白：言语是由思想和理智发出的。因此，那些编造\Quote{流溢}说的人，并没有发现什么伟大的事，也没有揭示任何深奥的奥秘，只不过是把人人理解的东西转移到了天主的独生圣言身上；他们称祂为不可言说和不可名状的，却又自以为协助了祂的诞生，从而阐发了祂最初世代的产生和形成，将其比作由流溢形成的人类的言语。\par

7. 然而，如果我们同样肯定关于物质的本体，即天主产生了它，我们也不会错。因为我们从圣经中得知，天主掌管万有。但祂在何处或以何种方式产生了它，圣经从未在任何地方明示；我们也不宜猜测，以至于根据自己的意见，对天主作无尽的臆测，而应将这种知识交在天主自己手中。同样，我们也应将以下原因交托于天主：为什么万物虽由天主造成，祂的某些受造物却犯罪并背离了顺服天主的境界，而另一些（实际上是绝大多数）却坚持并仍坚持（意愿地）服从于那形成了他们的那一位，以及那些犯罪的是什么样的受造物，那些坚持的又是什么样的——我们应将这些事情的原因交托于天主和祂的圣言，唯有祂曾对圣言说：\Quote{你坐在我右边，直到我使你的仇敌作你的脚凳。}而我们还在地上生活，尚未坐在祂的宝座上。因为虽然那在祂内的救主的圣神\Quote{洞察一切，甚至洞察天主的深奥}，\BibleRef{格前 2:10}但因我们有\Quote{各种恩赐、各种职分、各种功效}的区别；我们在地上，正如保禄也宣称的，是\Quote{以部分知道，以部分说预言}。\BibleRef{格前 13:9}因此，既然我们仅知道一部分，我们应将各种[困难]问题交托于那一位，祂曾在某种程度上（仅此而已）赐予我们恩宠。那永恒之火已为罪人预备，主已清楚宣告，其余圣经也证明。天主预知此事的发生，圣经同样证明，因为祂从起初就为那些后来要违背祂诫命的人准备了永恒之火；但这类犯罪者的本性之原因本身，没有任何经文告知我们，也没有宗徒告诉我们，主也未曾教导我们。因此，我们应将对此事的认识交托于天主，正如主对待审判的时日那样，而不应贸然陷入如此危险的极端，以至于在天主手中不留任何事，尽管我们在这世界上只领受了有限的恩宠。但当我们探究那些超越我们、且我们无法获得满足的问题时，我们却表现出如此极端的僭妄，好像我们已经通过关于\OriginalTerm{流溢}{emissions}的空谈发现了天主本身、万有的创造者，并断言祂的实体源自背道与无知，从而构建一个亵渎的反天主假设——这是荒谬的。\par

8. 此外，他们对自己的体系毫无证据，这体系是由他们最近发明的，有时依赖于某些数字，有时依赖于音节，有时又依赖于名字；有时还通过字母中包含的字母、未被正确解释的比喻、或某些[无根据的]臆测，他们竭力证明自己编造的那些荒诞说法。因为如果有人询问为什么那位与圣子在一切事上共融的圣父，唯独被主宣告知道审判的时日，他目前找不到比这更合适、更恰当或更稳妥的理由（因为主是唯一的真正导师），即我们应透过祂学习到圣父超越万有。因为主说：\Quote{圣父比我大。}\BibleRef{若 14:28}因此，我们的主已宣告圣父在知识上超越；为此原因，我们只要还处于这个世界的事物秩序中，就应将完美的知识和此类问题交托于天主，而不可因寻求探究圣父的崇高本性而陷入提出\Quote{在天主之上是否有另一位天主？}这一问题的危险。\par

9. 但若有爱争论者反驳我所说的话，以及宗徒所肯定的\Quote{我们以部分知道，以部分说预言}，\BibleRef{格前 13:9}并想象他已获得对一切存在者的全部知识而非部分知识——这样的人就像\Person{瓦伦提努}、\Person{托勒密}或\Person{巴西里德}，或任何其他宣称自己已探究天主深奥的人——他不要（以虚妄的荣耀装饰自己）夸耀自己比别人更了解那些不可见或无法观察的事物；而应殷勤探究并从圣父那里获取信息，告诉我们（我们所不知道的）这个世界中事物的原因——例如他自己头上的头发数目，以及每日被捕获的麻雀，以及我们事先不知道的其他此类要点——这样我们才能也相信他关于更重要之事的说法。但如果那些\Emph{完美者}还不理解手中、脚下、眼前、地上的这些事物，尤其是关于自己头上头发的规律，我们怎能相信他们关于属神之事、超天之事以及他们虚妄自信地断言超越天主之事呢？关于数字、名字、音节、此类超越我们理解力的问题，以及他们对比喻的不恰当解释，我就说这么多；因为你自己可以加以发挥。\par

\SourceRecord{Translated by Alexander Roberts and William Rambaut.；Ante-Nicene Fathers；Vol. 1.；Edited by Alexander Roberts, James Donaldson, and A. Cleveland Coxe.；Buffalo, NY: Christian Literature Publishing Co.,；1885.；<http://www.newadvent.org/fathers/0103228.htm>.}

% structure: p0001 source=h1 target=section reason=page-title

\section{驳斥异端（卷二，第29章）}

1. 让我们回到他们体系中其余的观点。他们宣称，在万物终结时，他们的母亲将重新进入\OriginalTerm{圆满}{Pleroma}，并接受\OriginalTerm{救主}{Saviour}作为她的配偶；他们自己，作为属灵者，在摆脱属魂的灵魂并成为理智的灵之后，将成为\OriginalTerm{属灵的天使}{spiritual angels}的配偶；但他们称之为属魂的\OriginalTerm{造物主}{Demiurge}将进入母亲的位置；义人的灵魂将在中间地带得到属魂的安息——当他们宣称同类将归于同类，属灵的归于属灵的，而物质的东西继续存在于物质的东西之中时，他们实际上自相矛盾，因为他们不再坚持灵魂因其本性而进入中间地带与类似自己的实体结合，而是因（身体中）所行的行为而进入，因为他们断言义人的灵魂进入那居所，而不敬虔者的灵魂则继续在火中。因为如果所有灵魂因其本性而达到享乐之地，并且所有灵魂仅因为它们是灵魂而属于中间地带（因为与它本性相同），那么信仰完全是多余的，\OriginalTerm{救主}{Saviour}的降临（到这个世界）也是多余的。另一方面，如果是因为他们的义（他们才达到这样的安息之地），那么就不是因为他们是\Emph{灵魂}，而是因为他们是\Emph{义人}。但如果灵魂除非成为义人就会灭亡，那么义也必定有能力拯救（这些灵魂所居住的）身体；因为为什么不拯救它们呢？既然它们也参与了义？因为如果本性和实体是得救的手段，那么所有灵魂都将得救；但如果义和信仰是得救的手段，为什么它们不拯救那些与灵魂同样将要进入不朽的身体呢？因为在这类事情上，义要么显得无力，要么显得不义，如果它确实通过参与它而拯救了一些实体，却不拯救其他实体。\par

2. 显然，那些被视为义的行为是在身体中进行的。因此，要么所有灵魂必然进入中间地带，并且永远不会有审判；要么，身体也参与了义，将与同样参与义的灵魂一同达到享乐之地，如果义确实有足够的能力将那些参与它的实体带到那里的话。然后，我们所相信的关于身体复活的教义将从他们的体系中显现为真实和确定的；因为（如我们所持守的）\OriginalTerm{天主}{God}，当祂使保持义的我们必死的身体复活时，将使它们成为不朽的和不死的。因为\OriginalTerm{天主}{God}超越本性，并且在祂自身内具有（施恩的）倾向，因为祂是善的；也有能力这样做，因为祂是强大的；也有完全实现祂目的的才能，因为祂是丰富和完美的。\par

3. 但这些人在所有方面都与自身不一致，当他们断定并非所有灵魂都进入中间地带，而只有义人的灵魂才能进入时。因为他们主张，根据本性和实体，\OriginalTerm{母亲}{Mother}产生了三种（存在）：第一种出于困惑、厌倦和恐惧——即物质实体；第二种出于冲动——即属魂实体；但她在看见那些侍奉\OriginalTerm{基督}{Christ}的天使之后所生的，是属灵实体。那么，如果她所生的那个实体因为是属灵的而必定进入\OriginalTerm{圆满}{Pleroma}，而物质实体因为是物质的而留在下面，并会被其中燃烧的火完全烧毁，那么为什么全部的属魂实体不进入中间地带呢？他们也把\OriginalTerm{造物主}{Demiurge}送到那里。但什么会进入他们的\OriginalTerm{圆满}{Pleroma}呢？因为他们主张灵魂会留在中间地带，而身体因具有物质实体，当它们分解为物质时，将被存在于其中的火烧毁；但身体这样被毁灭，灵魂留在中间地带，人的任何部分都不会再进入\OriginalTerm{圆满}{Pleroma}内。因为人的理智——他的心智、思想、精神意向及诸如此类——无非是他的灵魂；而灵魂本身的情绪和运作没有独立于灵魂的实体。那么，他们的哪一部分还会留下来进入\OriginalTerm{圆满}{Pleroma}呢？因为他们自己，就他们是灵魂而言，留在中间地带；就他们是身体而言，将与其他物质一起被烧毁。\par

\SourceRecord{Translated by Alexander Roberts and William Rambaut.；Ante-Nicene Fathers；Vol. 1.；Edited by Alexander Roberts, James Donaldson, and A. Cleveland Coxe.；Buffalo, NY: Christian Literature Publishing Co.,；1885.；<http://www.newadvent.org/fathers/0103229.htm>.}

% structure: p0001 source=h1 target=section reason=page-title

\section{驳斥异端（第二卷，第三十章）}

1. 既然如此，这些迷失心智的人宣称他们凌驾于造物主（\OriginalTerm{德缪革}{Demiurge}）之上；并且，鉴于他们自认为比那位创造并装饰了诸天、大地及其中的万物的神更为优越，声称自己\OriginalTerm{属神的}{spiritual}，而事实上由于他们如此大的不敬而可耻地属血肉——他们断言：那位使祂的天使成为神灵、以光为衣、手握大地之环、其上的居民在祂眼中如同蝗虫、且是一切属神实体的造物主及主宰，是\OriginalTerm{属魂的}{animal}——他们无疑确实暴露了自己的疯狂；而且，仿佛真的被雷击中，甚至比异教神话中的那些巨人更甚，他们带着虚妄的傲慢和不稳的荣耀高举自己的意见反对天主——这些人，世上所有的菉藜芦都不足以净化他们，使他们摆脱极度的愚妄。\par

2. 优越的人应由其行为证明。那么，他们怎能显示自己优于造物主呢？（同样，由于当前论证的需要，我也要降到他们不敬的程度，在天主与愚蠢之人之间进行比较，并通过降级到他们的论证，常常以他们自己的教义驳斥他们；但在这样做时，愿天主怜悯我，因为我敢于说出这些话，并非为了将祂与他们比较，而是为了谴责并推翻他们疯狂的见解。）——那些为许多愚人所如此钦佩的人，好像他们能从这些人那里学到比真理本身更宝贵的东西！他们解释圣经中的这句话：\BibleQuote{你们求，就会得到}\BibleRef{玛 7:7}，认为是为此目的说的：他们应发现自己在造物主之上，自称为比天主更大更好，称自己为\OriginalTerm{属神的}{spiritual}，而称造物主为\OriginalTerm{属魂的}{animal}；并断言因此他们向上高升于天主之上，因为他们进入了\OriginalTerm{普累若麻}{Pleroma}内，而祂则留在中间之地。那么，让他们通过自己的行为证明自己优于造物主吧；因为优越的人不应由言词证明，而应由实在的事物证明。\par

3. 那么，他们将指出由救主或他们的母亲通过他们完成了什么工作，比那位安排我们周围万物的造物主所产生的工作更伟大、更光荣、或更具智慧？他们建立了什么诸天？奠定了什么大地？召叫了什么星辰？或使什么天上之光闪耀？又用什么界限将它们约束？或者，他们给大地带来了什么雨、霜、雪，各适季节和特殊气候？此外，作为对立面，他们又设置了什么炎热或干燥？他们使什么河流流动？涌出了什么泉源？他们用什么样的花朵和树木装饰了这个月下世界？他们造了什么样的动物，有些有理性，有些无理性，但都饰以美丽？谁能逐一举出由天主的德能所建立、并由祂的智慧所管理的其余所有事物？谁能探究创造它们的那个天主的伟大？关于那些高于诸天、永不消逝的存在，如\OriginalTerm{天使}{Angels}、\OriginalTerm{总领天使}{Archangels}、\OriginalTerm{上座者}{Thrones}、\OriginalTerm{宰制者}{Dominions}、\OriginalTerm{掌权者}{Powers}，不可胜数，又能述说什么呢？那么，他们反对这些工作中的哪一件？他们有什么相似的东西可以展示，作为通过自己或由自己造成的，既然他们自己也是这位造物主的工艺品和受造物？因为无论救主还是他们的母亲（用他们自己的说法，通过他们使用的术语本身证明他们是错误的），正如他们所主张，用这个存在（德缪革）来塑造那些在\OriginalTerm{普累若麻}{Pleroma}内的事物以及她所见到的侍奉救主的一切存在的象，她使用他（德缪革）作为（某种意义上的）比她优越、并更适合通过他的工具来完成她的目的的个体；因为她绝不会通过一个较劣的、而是通过一个较优的工具来形成如此重要存在的象。\par

4. 因为，请注意观察，按照他们自己的说法，他们当时作为\OriginalTerm{属神的}{spiritual}受孕存在，是由于默观那些如同卫星般围绕\Person{潘多拉}排列的存在。他们确实一直无用，母亲没有通过他们完成任何事情——一个闲散的受孕，归因于救主，并且一无所用，因为看不出任何事是由他们做的。而按照他们的说法，那位被产生的神，虽然他们争辩说祂低于他们（因为他们坚持说祂是属魂的），却是万物中主动的行动者，有效且适合完成工作，以致通过祂，万物的象都被造了；不仅那些可见之物由祂形成，而且所有不可见之物，\OriginalTerm{天使}{Angels}、\OriginalTerm{总领天使}{Archangels}、\OriginalTerm{宰制者}{Dominations}、\OriginalTerm{掌权者}{Powers}、\OriginalTerm{德能者}{Virtues}——（我说，由祂）作为较优者，且能够服务于她的愿望。但看来母亲完全没有通过他们做任何事，正如他们自己所承认；所以，人们可以公正地认为他们是母亲痛苦分娩产生的流产。因为没有产婆来为母亲接生，因此他们被驱逐出来，如同流产，一无所用，被形成却未能完成母亲的任何工作。然而他们却描述自己优于那位完成并安排了如此伟大和可钦佩之工作的神，尽管根据他们自己的推理，他们被证明是如此可怜地低劣！\par

5. 这好比有两件铁制工具或器具，其中一件经常在工匠手中，不断使用，他用这件工具随意做出他所喜欢的东西，展示他的技艺和技巧；而另一件则闲置无用，从未被使用，工匠似乎从未用它做出任何东西，也不在任何工作中使用它。然后有人主张，这件无用、闲置、未被使用的工具，在本性和价值上优于工匠在工作中使用、并借以赢得名声的那件工具。若有这样的人，他理应被看作愚昧、精神不正常。那些自称是属神的、更优越的，而说造物主具有属魂的本性，并主张因此他们将升到高处，进入\OriginalTerm{圆满}{Pleroma}内，到达他们自己的配偶那里——因为根据他们自己的说法，他们自己是女性的——而天主（造物主）本性较低，因此停留在中间领域的人，也应受到同样的判断。然而他们对自己的主张提不出任何证据：因为更好的人由他的工作显明，而所有工作都是由造物主完成的；但他们自己却提不出任何值得理性称道的东西是由他们产生的，他们正遭受着最严重、最不可救药的疯狂。\par

6. 然而，如果他们竭力主张：虽然一切物质事物，如天和天下存在的整个世界，确实是由\OriginalTerm{得穆革}{Demiurge}形成的，但比这些更属神的一切事物——即那些在天之上的，如宰制者、大能者、天使、总领天使、上座者、异力者——是由一种属神的生育过程产生的（他们声称自己就是这样），那么，首先，我们从权威的圣经证明：所有提到的事物，有形无形的，都是由一位天主所造。因为这些人并不比圣经更可靠；我们也不应放弃主、梅瑟和其余先知——他们宣告了真理——的声明，而去相信这些人，他们确实没有说出任何明智的话，而是胡言乱语地坚持站不住脚的意见。其次，如果那些在天之上的事物真的是通过他们造成的，那么请他们告诉我们无形之物的本性，数算天使的数目和总领天使的品级，揭示上座者的奥秘，教导我们上座者、宰制者、大能者和异力者之间的区别。但他们对此说不出什么；因此，这些存在不是由他们造的。另一方面，如果这些是由造物主造的（实际正是如此），并且是属神的、圣洁的，那么产生属神之物的那位本身就不是属魂的本性，这样他们那可怕的亵渎体系就被推翻了。\par

7. 因为在天上确有属神的存在，所有圣经都大声宣告；保禄（\Person{保禄}）明确见证有属神的事物，当他宣称他被提到第三层天，又被带到乐园，听见不可言传的话，那是人不能说出的。但这对他有什么益处呢？无论是他进入乐园，还是他被提到第三层天，既然这些都仍在\OriginalTerm{得穆革}{Demiurge}的权下，如果——像有些人敢于主张的那样——他已经开始看见和听见那些被断言在得穆革之上的奥秘？因为如果他真的开始认识那超越得穆革的秩序，他就绝不会留在得穆革的领域内，甚至不会彻底探究这些领域（因为按照他们的说法，他面前还有四层天，如果他接近得穆革，就会看见全部七层天在他下面）；但他也许会被允许进入中间领域，即到\OriginalTerm{母亲}{Mother}面前，好从她那里接受关于\OriginalTerm{圆满}{Pleroma}内之事的教导。因为那在他里面的\OriginalTerm{内在的人}{inner man}，并在他内发言（如他们所说），虽然不可见，不仅能达到第三层天，甚至能到达他们的母亲面前。因为如果他们主张他们自己——即他们的（内在的）人——立刻升到得穆革之上，并去到母亲那里，那么宗徒的（内在的）人更应如此；因为得穆革不会阻止他，按照他们的断言，得穆革自己已服从于救主。但如果他试图阻止他，那努力也必归于徒然。因为得穆革不可能比圣父的眷顾更强，而且当内在的人甚至对得穆革也是不可见的时候。但既然他（保禄）把那次被提到第三层天描述为伟大和卓越的事，那么这些人不可能升到第七层天以上，因为他们肯定不比宗徒更优越。如果他们竟主张自己比他更卓越，那就让他们用工作证明自己吧，因为他们从未假装有过（像他描述的发生在他身上的）任何事。为此他补充说：\BibleQuote{或在身内，或在身外，天主知道}，这样身体就不会被认为参与了那异象，仿佛它能参与它所看见和听见的那些事；也不至于有人说，由于身体的重量他没有被提得更高；但因此允许人甚至在身体之外看见属神的奥秘，这些是天主的作为，他造了天和地，造了人，并把他安置在乐园里，使那些像宗徒一样在爱天主方面达到高度成全的人成为这些事的旁观者。\par

8. 因此，这一存在也创造了属神的事物，关于这些，宗徒曾被提升到第三层天，成为目睹者，并听到了不可言传的话语，那是人不可能讲出的，因为它们是属神的；祂亲自按自己的意愿将\OriginalTerm{恩赐}{gifts}赐予配得者，因为乐园是属于祂的；祂确实是天主的圣神，而不是一个属动物的德谬哥，否则祂绝不会创造属神的事物。但如果祂确实具有属动物的本性，那么就让他们告诉我们，属神的事物是由谁造的。他们无法提出任何证据证明这是通过他们母亲的产痛完成的，而他们自称就是那母亲。因为，且不说属神的事物，这些人连一只苍蝇、一只蚊子或任何其他微小而不重要的动物都不能创造，而不遵守从起初就由天主自然产生的动物所遵循的规律——即通过种子在同一种类中繁衍。没有一样东西是由母亲单独形成的；[因为]他们说这个德谬哥是由她产生的，并且\Emph{他}是万有的\OriginalTerm{主}{Lord}（创造者）。他们坚持认为，那创造并主宰万有的他具有属动物的本性，而他们自己却是属神的——他们既不是任何一件作品的创造者或主宰者，不仅对外部事物如此，甚至连他们自己的身体也不是！此外，这些自称属神并高于创造者的人，却常常违背自己的意愿，遭受许多身体上的痛苦。\par

9. 因此，我们公正地谴责他们已远远偏离了真理。因为如果救主通过他（德谬哥）形成了受造物，那么他（救主）在这种情况下就被证明并不低于他们，反而高于他们，因为他被发现甚至是他们自身的形成者；因为连他们自己也在受造物中占有一席之地。那么，怎么能说这些人确实是属神的，而创造他们的他却是属动物的呢？或者，另一方面，如果（这确实是唯一真实的假设，正如我通过大量最清晰的论证所表明的那样）他（造物主）自由地创造了一切，并凭自己的能力安排和完成了它们，并且他的意志是万有的实体，那么他就会被发现是唯一的、创造万有的天主，唯独他是全能的，并且是唯一一位形成并塑造万有的圣父，包括可见的和不可见的，那些可以为我们感官所感知的和那些不能感知的，天上的和地上的，\Quote{借祂大能的言语；} \BibleRef{希 1:3} 并且他凭自己的智慧安排了一切，同时他容纳万有，但他自己不能被任何人容纳：他是形成者，他是建造者，他是发现者，他是创造者，他是万有的主；在他以外或在他之上没有别的，他也没有母亲，正如他们虚假地归于他的那样；也没有第二位神，如玛尔强所想象的；也没有三十个永世的\OriginalTerm{圆满}{Pleroma}，这已被证明是一个虚妄的假设；也没有像\OriginalTerm{深渊}{Bythus}或\OriginalTerm{元始}{Proarche}这样的存在；也没有一系列的诸天；也没有一个童贞之光，或一个不可名的\OriginalTerm{永世}{Æon}，事实上，没有任何一件这些异端者以及所有异端者疯狂梦想的事物。但只有一位天主，造物主——他超越一切率领者、掌权者、宰制者和异能者：他是圣父，他是天主，他是奠基者，他是制造者，他是创造者，他凭自己，即通过他的圣言和他的智慧，创造了这一切——天、地、海和其中的万物：他是公义的；他是良善的；是他形成了人，种植了乐园，创造了世界，引发了洪水，拯救了诺厄；他是亚巴郎的天主，依撒格的天主，雅各伯的天主，活人的天主：他是法律所宣告的，先知所预言的，基督所启示的，宗徒们所教导我们的，以及教会所相信的。他是我们的主耶稣基督的父：通过他的圣言，即他的圣子，通过他，他向所有蒙他启示的人启示和显明了自己；因为只有那些圣子向他启示的人才知道他。但是圣子，与圣父永恒共存，从古至今，是的，从起初，总是向天使、总领天使、掌权者、异能者以及所有他愿意他们认识天主的人启示圣父。\par

\SourceRecord{Translated by Alexander Roberts and William Rambaut.；Ante-Nicene Fathers；Vol. 1.；Edited by Alexander Roberts, James Donaldson, and A. Cleveland Coxe.；Buffalo, NY: Christian Literature Publishing Co.,；1885.；<http://www.newadvent.org/fathers/0103230.htm>.}

% structure: p0001 source=h1 target=section reason=page-title

\section{\WorkTitle{驳斥异端}（第二卷，第三十一章）}

1. 凡属于华伦提努学派的人既被推翻，实际上所有异端群体也都被颠覆了。因为我对他们的\OriginalTerm{普累若麻}{Pleroma}以及超越它的事物的所有论证，指明了万有之父如何被超越祂的事物所限制和圈定（如果确实有超越祂的事物的话），又如何[按他们的理论]绝对必然地在每一方面设想出许多父、许多普累若麻和许多世界的创造，从一个系列开始、以另一个系列结束；而且所有[被提及的存在者]都留在自己的领域内，并不好奇地干涉其他领域，因为彼此之间既无共同利益也无任何交往；并且没有另一位万有之神，只有全能者才拥有这个名号——[所有这些论证，我说]同样适用于那些属于马尔西昂、西蒙、梅安德学派的人，或任何其他像他们一样将我们所关联的受造界与圣父割裂的人。再者，我用来反对那些人的论证——那些人声称万有之父无疑包含一切，但我们所属的受造界并非由祂所造，而是由某种其他能力或由对\OriginalTerm{原父}{Propator}毫无所知的天使所造，这原父如同被宇宙的广袤范围所环绕的中心点，正如污点被[环绕的]斗篷所环绕一样——当我表明，除了万有之父外，任何其他存在形成我们所属的受造界并非一种可能的假设时，这些相同的论证将适用于撒图尼努斯、巴西里德、卡波克拉底斯以及表达类似观点的其余诺斯替派信徒。再者，关于流溢、\OriginalTerm{永世}{Æons}、[所谓的]堕落状态、以及她们母亲的不稳定性格的陈述，同样推翻了巴西里德和所有假称为诺斯替派的人，他们实际上只是以不同名称重复相同观点，但却比前者更大幅度地将真理之外的事物转移到他们自己的教义体系中。而我对数字的评论也将适用于所有滥用真理之物来支持此类体系的人。所有关于\OriginalTerm{造物主}{Demiurge}的论述——表明祂唯独是天主和万有之父——以及在随后各卷中可能做出的任何评论，我都将其应用于所有异端者。你们可以转化并说服其中较为温和合理的人，使他们不再亵渎自己的创造者、制造者、维持者和主，也不将祂的起源归于缺陷和无知；但那些凶猛、可怕且无理的人，你们要远远驱离，以免再忍受他们空洞的饶舌。\par

2. 此外，属于西蒙和卡波克拉底斯的人，以及任何其他据说行异能的人——他们行异能既非藉天主的德能，也非与真理相连，亦不为人的福祉，而是为了毁灭和误导人类，藉着魔法欺骗和普遍诡计，从而对相信他们的人造成比益处更大的危害，在他们误导人的那一方面——也将如此被驳倒。因为他们既不能使瞎子看见，不能使聋子听见，也不能驱除各类魔鬼——[确实，]除非是他们自己送到别人身上的那些，如果连这点他们也能做到的话。他们也不能医治软弱、瘸腿、瘫痪或身体其他部分受苦的人，正如在身体疾患方面常常发生的那样。他们也不能为可能发生的外在事故提供有效补救。至于他们使死人复活的能力，更是远不能及；主曾使死人复活，宗徒藉着祈祷也曾做到，而因着某种需要，在弟兄团体中常常发生——整个地方教会以多日禁食祈祷哀求[恩惠]，死者的灵魂返回，他因圣人的祈祷而获得恩赐——他们甚至不相信这是可能做到的，[并且主张]从死者中复活仅仅是对他们所宣告的真理的一种认知。\par

3. 因此，既然在他们中间存在错误和误导影响，以及邪恶地在人眼前行出的魔法幻象；但在教会中，同情、怜悯、坚定和真理，为了人类的援助和鼓励，不仅无偿地彰显出来，而且我们自己也为了他人的益处付出自己的财物；并且，被治愈的人常常缺乏所需之物，就从我们这里获得——[既然如此]，这些人无疑被证明完全与神圣本性、天主的恩惠和一切属灵美德无缘。但他们却充满了各种欺骗、背信的灵感、魔鬼的行为和偶像崇拜的幻象，实际上是那条龙的前驱\BibleRef{默 12:14}，这条龙将藉着同样类型的欺骗，用它的尾巴使三分之一星辰从天坠落，并将它们投到地上。我们必须像逃避它一样逃避他们；他们据说[行奇迹]时显示的奇观越大，我们就越应谨慎提防他们，因为他们被赋予了更强大的邪恶之灵。如果有人思考所引用的预言和这些人的日常作为，就会发现他们的行为方式与魔鬼如出一辙。\par

\SourceRecord{Translated by Alexander Roberts and William Rambaut.；Ante-Nicene Fathers；Vol. 1.；Edited by Alexander Roberts, James Donaldson, and A. Cleveland Coxe.；Buffalo, NY: Christian Literature Publishing Co.,；1885.；<http://www.newadvent.org/fathers/0103231.htm>.}

% structure: p0001 source=h1 target=section reason=page-title

\section{\WorkTitle{驳斥异端}（第二卷，第三十二章）}

1. 此外，他们这关于行为的亵渎意见——即认为他们应当经历各种行为，甚至最可憎的——已为主的教导所驳斥。在主那里，不仅犯奸淫者被弃绝，就连有意犯奸淫的人也被弃绝\BibleRef{玛 5:21}；不仅实际杀人的凶手被判定为有罪而招致自己的丧亡，就连无缘无故向弟兄发怒的人也是如此。主命令祂的门徒不仅不可恨人，还要爱仇敌；不仅不可发虚誓，甚至完全不发誓；不仅不可说邻人的坏话，甚至不可称任何人为\Quote{拉加}和\Quote{愚人}；否则有陷入地狱之火的危险；不仅不可打人，甚至自己被打时，应转另一面颊给那些虐待他们的人；不仅不可拒绝放弃他人的财物，甚至自己的财物被夺去，也不可从夺走的人那里再索回；不仅不可伤害邻人，也不可向他们行任何恶事，甚至自己受恶待时，要忍耐，并善待那些伤害他们的人，为他们祈祷，使能借着悔改而得救。这样，我们在任何事上都不效法别人的傲慢、情欲和骄傲。因此，既然这些人自夸的师傅——他们断言他的灵魂远比其他人的更高超、更卓越——确实非常恳切地命令某些事当行，作为美好和卓越的，另一些事不仅在实际行出时当戒绝，就连引向这些行为的思念也当戒绝，视为邪恶、有害、可憎；那么，当他们声称这样一位师傅的精神更高超、比别人更好，却显然传授完全违背祂教导的道理时，他们怎能不陷于混乱呢？再者，如果根本就没有所谓的善与恶，某些事被视为正义，某些事被视为不义，仅仅出于人的意见，祂就绝不会在教导中这样说：\Quote{义人将在他们圣父的国里发光如太阳}\BibleRef{玛 13:43}；但祂要把不义的人和不行正义事的人投入\Quote{永不熄灭的火，那里他们的虫子不死，火也不灭}\BibleRef{玛 25:41}。\par

2. 他们进一步主张，必须经历各种工作和行为，以便若有可能，在一次现世显现中完成一切，从而立即转入成全的境界。然而，他们却丝毫没有努力去从事那些属于美德、辛劳、光荣和技巧的事，这些事也普遍被认为是好的。因为如果必须经历每一种工作和每一种操作，他们首先应当学习所有的技艺——无论是理论还是实践，无论是通过克制获得，还是通过劳动、练习和毅力掌握的；例如，各种音乐、算术、几何、天文学，以及所有涉及智力的学科；其次，全部医学知识、植物学知识，以便了解那些为人类健康而预备的植物；绘画和雕刻艺术、铜器和石工，以及类似的艺术；此外，还必须学习各种农活、兽医学、畜牧工作、各种手艺，这些据说遍及人类活动的全部领域；还有与航海生活相关的、体操、狩猎、军事和统治活动，以及其他许多可能存在的，他们即使付出最大努力，终其一生也无法学到其中的十分之一，甚至千分之一。事实上，他们根本不努力学这些，尽管他们声称必须经历各种工作；但他们转而追求享乐、情欲和可憎的行为，因此当被自己的教义审判时，他们已自证有罪。因为他们缺乏上述所有美德，就必然要进入火的毁灭。这些人夸耀耶稣是他们的师傅，实际上却效仿伊壁鸠鲁的哲学和犬儒学派的漠然，他们称耶稣为他们的师傅，而耶稣不仅使祂的门徒远离恶行，甚至远离恶言恶念——正如我已经指出的。\par

3. 再者，当他们声称自己的灵魂与耶稣来自同一领域，并像祂一样，有时甚至声称自己超越祂；又宣称他们像祂一样受造，为要成就有益于人类利益和建立人类的工作，却发现他们并未行任何与祂相同或类似的事，也没有任何可与之比拟的事。如果他们确实借着魔术行过什么令人瞩目的事，他们也不过是欺骗性地误导愚昧人，因为他们并未给那些他们宣称施展超自然能力的人带来任何真正的益处或祝福；反而带来小男孩作为练习对象，欺骗他们的视觉，显示一些立刻消失、连片刻也不持久的幻象，借此证明他们不像我们的主耶稣，而像西满术士。同样，从主在第三天从死者中复活，向祂的门徒显现，并在他们眼前被接升天的事实来看，这些人既然死去，不再复活，也不向任何人显现，他们就被证明拥有与耶稣毫不相似的灵魂。\par

4. 然而，若他们主张主也仅在外表上行了这些事，我们就将他们引到先知著作，并从其中证明：所有关于祂的事都已被预言，且确实发生了，而祂是天主的独生子。因此，那些真正作祂门徒的人，从祂领受恩宠，因祂的名施行[奇迹]，以促进他人的福祉，各人按从祂所领受的恩赐。因为有些人确实真实地驱逐魔鬼，以致那些这样从恶神中被洁净的人常会相信基督，并加入教会。另有些人预知未来之事；他们看见神视，说出先知性言词。还有些人藉覆手在病人身上治好他们，他们就痊愈了。不仅如此，如我曾说，甚至死人复活了，并留在我们中间多年。我还有什么可说的？无法数算那散布全世界的教会，从天主那里，因着那在\Person{般雀比辣多}下被钉的\Person{耶稣基督}之名所领受的神恩，并且她天天施行这些神恩以造福外邦人，既不欺骗任何人，也不从他们那里收取任何报酬\BibleRef{宗 8:9}[基于这样的奇迹干预]。因为她白白地从天主领受了\BibleRef{玛 10:8}，也白白地服务[他人]。\par

5. 她也不藉天使的呼求、或咒语、或任何其他邪恶的奇术行事；而是以纯洁、真诚、正直的精神，将祈祷导向那创造万有的主，并呼求我们的主\Person{耶稣基督}之名，惯常施行奇迹以造福人类，而非引人走入歧途。因此，若我们的主\Person{耶稣基督}之名如今仍然赐予恩惠[于人类]，并彻底有效地治愈所有任何地方相信祂的人，而非\Person{西蒙}、\Person{默南德}、\Person{卡尔波克拉特}或任何其他人的名，那么显然，当祂成为人时，祂与祂自己的受造界共融，并按照万有父的旨意，真真实实地藉天主的权能行了这一切，正如先知们所预言的。但这些事是什么，将在处理先知著作中的证明时加以描述。\par

\SourceRecord{Translated by Alexander Roberts and William Rambaut.；Ante-Nicene Fathers；Vol. 1.；Edited by Alexander Roberts, James Donaldson, and A. Cleveland Coxe.；Buffalo, NY: Christian Literature Publishing Co.,；1885.；<http://www.newadvent.org/fathers/0103232.htm>.}

% structure: p0001 source=h1 target=section reason=page-title

\section{驳斥异端（第二卷，第三十三章）}

1. 我们可凭以下事实推翻他们关于灵魂从一身体转入另一身体的学说：灵魂对于前生所发生的任何事件毫无记忆。因为，如果灵魂被派遣到世上是为了体验各种行为，那么它必然要记住先前已完成的事，以便弥补尚缺的经历，而不至于无休止地反复进行同样的追求，徒然劳苦（因为身体与灵魂的结合并不能完全抹去对先前经历的记忆与沉思），尤其灵魂来到世间本就是为此目的。正如身体安眠休息时，灵魂独自看见并在异象中所行的一切，它记住其中许多事，并将它们传达给身体；又如人醒来后，或许过了很久，还能讲述梦中所见；同样，灵魂无疑也会记得它进入这个特定身体之前所行的事。因为，若仅在极短时间所见、或仅在梦境中由灵魂单独幻象所感知的事，在灵魂再次与身体结合、分散于各肢体后仍能被记住，那么，灵魂在先前整个一生中长久停留之事，岂不更应被牢记？\par

2. 针对这些反驳，那位古雅典人\Person{柏拉图}——他亦是首先引入此见解者——因无法驳倒它们，便发明了\Quote{遗忘之杯}的概念，以为如此便可逃避此类难题。他并未为他的假设提出任何证明，只是武断地回答：灵魂进入此生时，在进入所分配的身体之前，由守护其入口的魔鬼令其饮下\Quote{遗忘之杯}。他未曾察觉，这样说反而陷入了更大的困境。因为，若饮下\Quote{遗忘之杯}能抹去一切所作之事的记忆，那么，\Person{柏拉图}啊，你如何得知此事（因你的灵魂如今在身体内），即灵魂进入身体前曾被魔鬼令饮下\Quote{致遗忘的药水}？你若记得魔鬼、杯盏及进入此生之事，就也应知道其他事；但若你不知这些，那么关于魔鬼和精心调制的\Quote{遗忘之杯}的故事便不真实。\par

3. 反之，对于那些声称身体本身就是\Quote{遗忘之药}的人，可提出如下观察：为何灵魂在身体被动时，无论通过梦境、反思或全神贯注的精神努力，凭自身所见的一切，它都能记住并告诉邻人？再者，若身体本身是遗忘的原因，那么灵魂寓居身体内时，甚至无法记住不久前通过眼或耳所感知的事；但一旦眼睛从所视之物移开，其记忆也必定立即消失。因为灵魂寓居在遗忘之因本身中，除了当下所见之物外，不能知晓任何其他东西。既然如此，灵魂又如何能认识神圣事物并在身体内保持记忆——既然他们主张身体本身就是遗忘之因？此外，先知们在世时，当他们在异象中灵性上看见或听见天上的事物后，恢复平常心神时也会记住这些，并将之告诉他人。因此，身体并不使灵魂忘记灵性所见之事；相反，灵魂教导身体，并与身体分享它所享有的灵性异象。\par

4. 因为身体并不比灵魂更有能力——实际上，身体由灵魂赋予生气、赋予生命、成长并维系；而灵魂拥有并统治身体。无疑，灵魂的速度因身体参与其运动而相应减缓，但它从不失去本属于自己的知识。身体可比作乐器，灵魂则拥有艺术家的理智。因此，正如艺术家心中迅速涌现作品的构想，但因工具不够柔顺而只能缓慢实现，其心智的敏捷与工具的迟缓相混合，便产生适中的运动；同样，灵魂与其所属的身体相混合，也在一定程度上受到阻碍，其敏捷与身体的迟缓相混合。然而，它并未完全丧失自身特有的能力；当它仿佛与身体共享生命时，它自身并不停止存活。同样，当它将其他事物传达给身体时，它既不丧失对这些事物的知识，也不丧失对所见证之事的记忆。\par

5. 因此，如果灵魂不记得前生所发生的一切，只对现世的事物有感知，那么，它便从未存在于其他身体中，也从未做过它所不知的事，也从未【一度】知晓它无法【现在从心智上】思虑的事。然而，正如我们每人因天主巧妙的运作而获得身体，同样，人也拥有自己的灵魂。因为天主并非如此贫乏或资源匮乏，以至于不能赐予每个身体它本有的灵魂，正如祂也赋予其特有的性格。因此，当【预定】的数目完成时——就是祂在自己计划中预定的数目——所有已被登记于【永生】的人都要复活，拥有他们自己的身体，也拥有他们自己的灵魂和精神，他们曾藉以悦乐天主。另一方面，那些应受惩罚的人，将进入惩罚中，他们也拥有自己的灵魂和身体，他们曾藉以远离天主的恩宠。这两类人届时将不再生育或受生，也不再嫁娶；如此，相应于天主预定的数目的人类，既已完成，便能圆满实现圣父所制定的计划。\par

\SourceRecord{Translated by Alexander Roberts and William Rambaut.；Ante-Nicene Fathers；Vol. 1.；Edited by Alexander Roberts, James Donaldson, and A. Cleveland Coxe.；Buffalo, NY: Christian Literature Publishing Co.,；1885.；<http://www.newadvent.org/fathers/0103233.htm>.}

% structure: p0001 source=h1 target=section reason=page-title

\section{\WorkTitle{驳斥异端（卷二，第三十四章）}}

1. 主已极其完满地教导我们：灵魂不仅继续存在，并非从一身体转入另一身体，而且它们保持与它们所相配的身体相同的形状[在分离状态中]，并记念它们在现世中所行的事——这些事如今已停止。这在关于那位富人和那在亚巴郎怀中安息的拉匝禄的记载中显明了。在此叙述中，祂宣称\BibleRef{路 16:19}，\OriginalTerm{富人}{Dives}死后认出了拉匝禄，亚巴郎也同样认出，并且他们每个人都保持在自己适当的位置上，而[富人]请求拉匝禄被派来救济他——拉匝禄，他从前甚至不给他桌上掉下的碎屑。[祂也告诉我们]亚巴郎的回答，亚巴郎不仅知道关于他自己的事，也知道关于富人的事，并且嘱咐那些不愿进入那痛苦之地的人，要相信梅瑟和先知们，并接受那将从死者中复活者的宣讲。因此，这些事清楚地宣告：灵魂继续存在，它们不从一身体转入另一身体，它们具有人的形状，以便被认出，并保留对现世之事的记忆；此外，亚巴郎拥有预言的恩赐，并且每一类[灵魂]在审判之前就已领受了所应得的住所。\par

2. 但若有人此时主张：那些刚刚开始存在的灵魂不能长久存留；相反，它们要么必须是非受生的，以便不朽；要么如果它们通过出生而有开始，就应与身体一同死亡——让他们知道：唯有天主，万有的主，是无始无终的，真实且永远同一，始终如一，保持不变。但一切由祂而出的事物，凡受造的和正在受造的，的确有其生成的开始，因此低于那形成它们的一位，因为它们不是非受生的。然而，它们依照造物主天主的旨意存留，并将其存在延长至漫长的岁月序列，使得祂恩准它们起初如此形成，此后如此存在。\par

3. 因为正如我们上面的穹苍、太阳、月亮、其余星辰及其全部宏伟，虽然它们先前并不存在，却因天主的旨意而被召唤存在，并经过漫长的岁月继续存在；同样，任何这样思考灵魂、精神以及实际上一切受造之物的人，必不会偏离正轨，因为一切受造之物在形成时有一个开始，但依照天主的旨意，它们存留多久，天主就愿意它们存在多久。预言之神对此观点作证，祂宣告：\BibleQuote{因为祂说话，它们就成就；祂命令，它们就受造。祂使它们永远、永远立定。} 再者，关于人的救恩，祂这样说：\BibleQuote{祂向祢求生命，祢便将长寿赐给祂，永远、永远。} 表明是万有的父将永远的存续赐给那些得救的人。因为生命并非出于我们，也非出于我们自己的本性，而是按天主的恩宠赐予的。因此，那保存所受的生命并感谢赐予者的人，也必得着永远、永远的长寿。但那拒绝它并显明对造物主忘恩的人——既然他被造，却不认识施恩者——就自绝于永远、永远的存续。为此，主对那些对祂忘恩的人说：\BibleQuote{如果你们在小事上不忠信，谁会把大事交给你们呢？} 指出那些在此短暂现世生活中对施恩者显明忘恩的人，理应从祂那里得不到永远、永远的长寿。\par

4. 但动物的身体本身当然不是灵魂，然而只要天主愿意，它就与灵魂有共融；同样，灵魂本身不是生命，而是分享天主赐予她的生命。因此先知的话论及最初受造的人说：\BibleQuote{他成了一个活灵，} \BibleRef{创 2:7} 教导我们：因着生命的分享，灵魂成为活的；所以灵魂及其拥有的生命应被理解为各自独立的存在。因此当天主赐予生命和永恒延续时，就发生这样的事：甚至连先前不存在的灵魂也从此存留[永远]，因为天主既愿意它们存在，也愿意它们继续存在。因为天主的旨意应当统御并统治万有，而其他一切事物都向祂让步、服从并事奉祂。至此，让我就灵魂的受造和持续存留这样说吧。\par

\SourceRecord{Translated by Alexander Roberts and William Rambaut.；Ante-Nicene Fathers；Vol. 1.；Edited by Alexander Roberts, James Donaldson, and A. Cleveland Coxe.；Buffalo, NY: Christian Literature Publishing Co.,；1885.；<http://www.newadvent.org/fathers/0103234.htm>.}

% structure: p0001 source=h1 target=section reason=page-title

\section{驳斥异端（第二卷，第三十五章）}

此外，除上述所言之外，\Person{巴西里得}自己根据其原则，必然不仅要主张有三百六十五重天依次由彼此生成，而且还要主张有无数不可胜数的天一直在生成中、正在生成中、并将继续生成，以至于此类天的形成永不止息。因为如果从第一重天的流出中，第二重天按其肖像被造，第三重天按第二重天的肖像，如此类推所有后续的天，那么，作为必然的结果，从我们的天（他称之为最后一重）的流出中，将形成另一个相似的天，再从那个天形成第三个；如此，从已经生成的天中流出，以及[新]天的制造，都永无止息，而这一运作必须持续\Emph{无限进行}，并产生完全不确定数量的天。\par

其余那些被称为\Quote{诺斯替派}的假知识者，他们主张先知们在不同神祇的默感下发出预言，将会被这一事实轻易推翻：所有先知都宣告了同一天主和主，即天与地及其中万物的创造者本身；此外，他们还预告了祂圣子的来临，正如我将在后续卷册中从圣经本身所证明的那样。\par

因为\OriginalTerm{Eloë}{Eloë}一词在犹太语言中意指天主，而\OriginalTerm{Elōeim}{Elōeim}和\OriginalTerm{Eleōuth}{Eleōuth}在希伯来语中意为\Quote{ \Emph{包容万有者} }。至于\OriginalTerm{Adonai}{Adonai}这一称谓，有时它意指\Emph{可名状}和\Emph{可钦佩}者；但另一些时候，当其中的\OriginalTerm{Daleth}{Daleth}字母被双写，且该词获得一个初始的喉音——即\OriginalTerm{Addonai}{Addonai}——[它意指]，\Quote{一位限定并分开陆地与水者，}使得水不再淹没陆地。同样地，\OriginalTerm{Sabaoth}{Sabaoth}，当最后一个音节用希腊字母Omega拼写时[\OriginalTerm{Sabaōth}{Sabaōth}]，意指\Quote{ \Emph{自愿的行动者} }；但当用希腊字母Omicron拼写时——例如\OriginalTerm{Sabaŏth}{Sabaŏth}——它表达\Quote{ \Emph{第一重天} }。同样的，\OriginalTerm{Jaōth}{Jaōth}一词，当最后一个音节被拉长并送气时，意指\Quote{ \Emph{预定的尺度} }；但当用希腊字母Omicron写短时，即\OriginalTerm{Jaŏth}{Jaŏth}，它意指\Quote{ \Emph{驱逐邪恶者} }。所有其他表达同样彰显同一位存在者的名号；例如，\Emph{万军之主、万有之父、全能天主、至高者、创造者、造物主}等等。这些并非一系列不同存在者的名字和称号，而是同一位的名号，藉此揭示了独一天主和父，祂包含万有，并赐予万有存在的恩赐。\par

现在，宗徒们的宣讲、主的权威教导、先知们的宣告、宗徒们的受默感言论、以及法律的施行——所有这些都赞美同一位存在者，即万有的天主和圣父，而非许多不同的存在者，也不是一个从不同神祇或能力中衍生其本体的存在者，而是[宣告]万物由同一位圣父所形成（祂却使[祂的工作]适应所处理材料的本性和倾向），有形可见与不可见之物，简言之，一切受造之物[被创造]既不是由天使，也不是由任何其他能力，而是唯独由天主圣父——这些都与我们的陈述和谐一致，我认为这已充分证明；而通过这些有力的论证，已经表明只有一个天主，万物的创造者。但为避免有人认为我回避从主的圣经中得出的那系列证据（事实上，这些圣经更清楚明晰地宣告了这一点），我将，至少为了那些不带败坏心智来阅读的人，专门写一卷书来讨论所引用的圣经，将公平地追循它们[并解释它们]，并从这些神圣圣经中清晰地提出证明，以满足所有真理爱好者。\par

\SourceRecord{Translated by Alexander Roberts and William Rambaut.；Ante-Nicene Fathers；Vol. 1.；Edited by Alexander Roberts, James Donaldson, and A. Cleveland Coxe.；Buffalo, NY: Christian Literature Publishing Co.,；1885.；<http://www.newadvent.org/fathers/0103235.htm>.}

% structure: p0001 source=h1 target=section reason=page-title

\section{\WorkTitle{驳斥异端（第三卷，序言）}}

你实在嘱咐了我，我至亲爱的朋友，要把瓦伦廷教义——正如其信奉者所想象的那样隐藏着——揭示出来；要展示其多样性，并撰写一篇驳斥它们的论文。因此，我承担了这项工作——表明它们源自所有异端者的始祖西满——既要展示他们的教义与传承，也要提出反对所有他们的论据。因此，由于对这些人的定罪和揭露在许多方面只是一项工作，我已寄给你[某些]书卷，其中第一卷包含了所有这些人的意见，并展示了他们的习俗和行为特征。在第二卷中，再次将他们乖谬的教导摧毁和推翻，并按其真实面目暴露于光天化日之下。但在这一卷，即第三卷中，我将从圣经中提出证据，以便在你所嘱咐的事上毫无欠缺；是的，超过你所期望的，你从我这里得到与那些以各种方式散布谎言的人争战并战胜他们的方法。因为天主的爱，既丰富又不吝啬，赐予恳求者超过他能所求的。那么，请回想我在前两卷中所陈述的事，并将这些与它们联系起来，你将从我这里得到对所有异端者非常充分的驳斥；并且你要忠实地、奋力地抵抗他们，护卫那唯一真实而赋予生命的信仰，就是教会从宗徒们所领受并传授给她的子女的。因为万有的主将福音的能力赐给了祂的宗徒们，借着他们我们也认识了真理，即天主圣子的教导；主也曾向他们宣告：\BibleQuote{听从你们的，就是听从我；拒绝你们的，就是拒绝我，并拒绝那派遣我的。}\BibleRef{路 10:16}\par

\SourceRecord{Translated by Alexander Roberts and William Rambaut.；Ante-Nicene Fathers；Vol. 1.；Edited by Alexander Roberts, James Donaldson, and A. Cleveland Coxe.；Buffalo, NY: Christian Literature Publishing Co.,；1885.；<http://www.newadvent.org/fathers/0103300.htm>.}

% structure: p0001 source=h1 target=section reason=page-title

\section{\WorkTitle{反异端}（第三卷第一章）}

1. 我们救恩的计划，并非从别人，而是从那些藉着他们将福音传给我们的人学来的；他们曾一度公开宣讲，后来，按天主的旨意，以圣经形式传授给我们，成为我们信德的根基和支柱。因为，断言他们在拥有\Quote{圆满的知识}之前就宣讲，这是不合法的，正如有些人甚至胆敢说，自夸为宗徒的改进者。因为，我们的主从死者中复活后，[宗徒们] 被赋予来自高天的权能，当圣神降临 [在他们身上] 时，他们充满了 [祂所有的恩赐]，并拥有了圆满的知识：他们前往地极，宣讲从天主赐予我们的美事的喜讯，并向人宣告天上的平安，这些人确实都同样并个别地拥有天主的福音。\Person{玛窦}也在希伯来人中用他们自己的方言发表了一部书写的福音，那时\Person{伯多禄}和\Person{保禄}正在\Place{罗马}宣讲，并奠定教会的根基。在他们离去后，\Person{马尔谷}，\Person{伯多禄}的门徒和翻译者，也将\Person{伯多禄}所宣讲的以书面形式传授给我们。同样，\Person{路加}，\Person{保禄}的同伴，将\Person{保禄}所宣讲的福音记录在一本书中。随后，主的门徒\Person{若望}，就是曾倚靠在祂胸膛的那位，在\Place{亚细亚}的\Place{厄弗所}居住期间，亲自发表了一部福音。\par

2. 所有这些人都向我们宣告：有一位天主，天地的创造者，由法律和先知所宣告；并有一位基督，天主圣子。若有人不同意这些真理，他就藐视主的同伴；不仅如此，他还藐视基督自己，即主；是的，他也藐视圣父，并自定己罪，抵抗并反对自己的救恩，正如所有异端者的情形一样。\par

\SourceRecord{Translated by Alexander Roberts and William Rambaut.；Ante-Nicene Fathers；Vol. 1.；Edited by Alexander Roberts, James Donaldson, and A. Cleveland Coxe.；Buffalo, NY: Christian Literature Publishing Co.,；1885.；<http://www.newadvent.org/fathers/0103301.htm>.}

% structure: p0001 source=h1 target=section reason=page-title

\section{驳斥异端（第三卷，第二章）}

1. 然而，当他们被圣经驳倒时，他们转而指责这些圣经本身，仿佛它们不正确、没有权威，并断言它们含混不清，无知于传统的人无法从中提取真理。因为他们声称，真理并非通过书面文件传递，而是\OriginalTerm{口传}{viv\={a} voce}。因此保禄也宣告说：\Quote{\BibleQuote{但在成熟的人中，我们宣讲智慧，却不是这世界的智慧。}}\BibleRef{格前 2:6}他们每个人声称这种智慧是他自己捏造的虚构，诚然！因此，按照他们的想法，真理时而驻留在\Person{瓦伦提诺}身上，时而驻留在\Person{马尔基昂}身上，时而驻留在\Person{塞林图}身上，其后又在\Person{巴西里德斯}身上，或者甚至无差别地存在于任何其他不能讲论任何有关救恩之事的反对者身上。因为这些人中的每一个，心地完全乖僻，败坏了真理的系统，却毫不羞耻地宣扬自己。\par

2. 但是，当我们把他们引向那源自宗徒、并通过教会中长老的继承而保存的传统时，他们又反对传统，声称他们自己不仅比那些长老，甚至比宗徒们更有智慧，因为他们发现了未掺杂的真理。因为他们主张，宗徒们将法律的事与救主的话混杂在一起；而且不仅是宗徒们，甚至主自己有时从\OriginalTerm{德缪哥}{Demiurge}说话，有时从中间之地说话，有时又从\OriginalTerm{圆满}{Pleroma}说话；而他们自己则无可置疑地、纯洁无误地知晓那隐藏的奥秘。这实在是无耻至极地亵渎他们的创造者！因此，这些人现在既不同意圣经，也不同意传统。\par

3. 我极亲爱的朋友，我们所应对的对手就是如此，他们像滑溜的蛇一样试图从各个角落逃脱。因此，必须在各个方面反对他们，或许通过切断他们的退路，我们能够使他们回转归向真理。因为，虽然被谬误所影响的灵魂悔改不是一件容易的事，但另一方面，当真理被带到谬误旁边时，逃脱谬误也并非完全不可能。\par

\SourceRecord{Translated by Alexander Roberts and William Rambaut.；Ante-Nicene Fathers；Vol. 1.；Edited by Alexander Roberts, James Donaldson, and A. Cleveland Coxe.；Buffalo, NY: Christian Literature Publishing Co.,；1885.；<http://www.newadvent.org/fathers/0103302.htm>.}

% structure: p0001 source=h1 target=section reason=page-title

\section{\WorkTitle{驳斥异端}（第三卷，第三章）}

因此，在任何教会中，凡愿看到真理的人，都能清楚默观那传遍全世界的宗徒传统；而且我们能够历数那些由宗徒在各教会中设立为主教的人，并［证明］这些人的继承直到我们当代；这些人既没有教导，也不知道这些异端者所胡言乱语的任何类似之事。因为，如果宗徒知道隐秘的奥秘，并且惯于私下和秘密地将之传授给\Quote{完美者}，那么他们就会特别将这些奥秘交给那些他们也将教会本身托付的人。因为他们渴望这些人成为非常完美、在一切事上无可指摘的人，正是这些人他们留下作为自己的继承者，将自己的治理权交付他们；这些人若忠诚履行其职责，将成为［教会］的巨大福分，但若他们堕落，则会造成最可怕的灾祸。\par

然而，在这样一卷书中历数所有教会的继承将十分冗长，因此，我们使那些以任何方式——或出于邪恶的自满、虚荣、盲目和偏执的见解——非法聚会的人陷入混乱；［我们这样做］是指明那来自宗徒的传统，即由两位最光荣的宗徒——\Person{伯多禄}和\Person{保禄}——在\Place{罗马}建立并组织的非常伟大、非常古老且举世闻名的教会；同时也［指明］那向人宣讲的信德，它借着主教的继承传到我们时代。因为所有教会，即各处忠信之人，都必须与此教会一致，因其卓越的权柄，既然传统已由各处那些忠信之人持续保存。\par

蒙福的宗徒，在建立并坚固了教会之后，将主教职务托付给\Person{理诺}。保禄在\BibleBook{弟茂德}{书}中提到了这位\Person{理诺}。继他之后的是\Person{安纳克肋托}；在他之后，第三位自宗徒起，\Person{格肋孟}被指派担任主教职务。此人既曾见过蒙福的宗徒，并与他们交往，可谓宗徒的宣讲仍回荡［在他耳中］，他们的传统仍在眼前。他并非独自［如此］，因为还有许多仍存留者曾接受宗徒的训诲。在\Person{格肋孟}的时代，\Place{格林多}的弟兄之间发生了不小的纷争，\Place{罗马}的教会寄了一封极有力量的\WorkTitle{书信}给\Place{格林多}人，劝勉他们和好，更新他们的信德，并宣告它最近从宗徒领受的传统，宣扬唯一的天主、全能者，天地\OriginalTerm{创造者}{Maker}，人的\OriginalTerm{造物主}{Creator}，祂曾带来洪水，召叫了\Person{亚巴郎}，带领子民离开\Place{埃及}地，与\Person{梅瑟}说话，颁布法律，派遣先知，并为魔鬼及其天使预备了火。凡愿意者，可以从这\WorkTitle{书信}中得知，我们的主耶稣基督的圣父曾被各教会宣讲，也可明白教会的传统，因为这封\WorkTitle{书信}比现今传播谎言并另造一个超越创造者及万有之创造者的神明的人更古老。这位\Person{格肋孟}之后继任的是\Person{恩华利斯笃}；\Person{亚历山大}继任\Person{恩华利斯笃}；然后第六位自宗徒起，\Person{西斯笃}被任命；此后是\Person{泰利福禄}，他光荣地殉道；然后\Person{希奇诺}；此后是\Person{比约}；然后在他之后是\Person{安尼塞托}。\Person{索泰尔}继任\Person{安尼塞托}，现在\Person{厄留提若}，自宗徒起第十二位，持有主教继承。按此秩序，并由此继承，教会自宗徒的传统和真理的宣讲传到了我们这里。这便是最充分的证据，证明有同一个赋予生命的信德，自宗徒直到如今在教会中保存，并真实地传递下来。\par

4. 然而，\Person{波利卡普}不仅由宗徒们教导，并与许多见过基督的人交谈，而且由宗徒们在\Place{亚细亚}被任命为\Place{斯米纳}教会的\OriginalTerm{主教}{bishop}，我在幼年时也见过他，因为他在地上逗留了很长时间，当他在极其年迈时，光荣而最高贵地殉道，离开此世，他始终教导他从宗徒们所学到的、教会所传承的、以及唯独真实的事物。所有\Place{亚细亚}教会都为此作证，那些继承\Person{波利卡普}直到如今的人也是如此——他是一个远比\Person{瓦伦蒂努斯}、\Person{马尔基翁}以及其他异端者更有分量、更坚定的真理见证者。正是他在\Person{阿尼塞特}时代来到\Place{罗马}，使许多人从上述异端者转向天主的教会，宣告他从宗徒们领受了这唯一且独一的真理——即教会所传承的真理。还有一些人从他那里听到，主的门徒\Person{若望}，在\Place{厄弗所}去沐浴时，看见\Person{塞林特}在里面，便未沐浴就冲出浴室，喊道：\Quote{我们快逃，免得浴室倒塌，因为真理的敌人\Person{塞林特}在里面。}\Person{波利卡普}本人也曾回答遇见的\Person{马尔基翁}，\Person{马尔基翁}说：\Quote{你认识我吗？}\Person{波利卡普}说：\Quote{我认识你，你这\Person{撒殚}的长子。}宗徒们及其门徒对于与任何真理败坏者甚至言语交流都怀有如此憎恶；正如\Person{保禄}也说：\BibleQuote{异端的人，第一次和第二次警戒以后，就该弃绝他；因为你知道这样的人已经背道，自己定自己的罪。}\BibleRef{铎 3:10}还有\Person{波利卡普}写给\Place{斐理伯}人的一封非常有力量的书信，那些愿意并且关心自己得救的人，可以从这封信中了解他信德的特质和真理的宣讲。此外，由\Person{保禄}建立、并有\Person{若望}在他们中间一直居住到\Person{图拉真}时代的\Place{厄弗所}教会，是宗徒传统的一个真实见证。\par

\SourceRecord{Translated by Alexander Roberts and William Rambaut.；Ante-Nicene Fathers；Vol. 1.；Edited by Alexander Roberts, James Donaldson, and A. Cleveland Coxe.；Buffalo, NY: Christian Literature Publishing Co.,；1885.；<http://www.newadvent.org/fathers/0103303.htm>.}

% structure: p0001 source=h1 target=section reason=page-title

\section{\WorkTitle{反异端（第三卷，第四章）}}

1. 既然我们有了这样的证据，就不必在其他人那里寻求真理，因为从教会容易获得；因为宗徒们如同一个富人在银行里存款一样，把一切属于真理的事物丰丰富富地交在她手中，以便凡愿意的人，都可以从她那里汲取生命的水。\BibleRef{默 22:17} 因为她就是生命的入口；其他一切都是贼和强盗。因此，我们应当避开\Emph{他们}，而要极其谨慎地选择属于教会的事物，并紧紧抓住真理的传统。因为情况如何？假如在我们中间对某个重要问题产生了争议，难道我们不应该求助于与宗徒们经常交往的最古老的教会，从他们那里得知关于当前问题明确而清楚的事吗？假如宗徒们自己并没有给我们留下著作，那又该怎么办？难道那时，我们不是必须遵循他们传授给那些他们曾托付教会者的传统吗？\par

2. 许多信奉基督的未开化民族同意这一做法，他们有圣神写在心里的救恩，没有纸和墨，又小心地保存着古老的传统，相信唯一的天主，天地的创造者，以及其中万物的创造者，藉着天主子耶稣基督；祂因对受造物超乎寻常的爱，屈尊由童贞女所生，祂亲自将人通过自己与天主结合，并在本丢·比拉多执政时受苦，复活，在光荣中被接升天，将在光荣中来临，作得救者的救主，和受审判者的审判者，并将那些歪曲真理、藐视祂的父和祂的来临的人投入永火。这些人在没有文字的情况下相信了这信仰，就我们的语言而言是未开化的，但在教义、生活方式和处世态度上，他们因信德而极其智慧，并以一切正义、贞洁和智慧来安排他们的交往，他们确实中悦天主。如果有人用他们自己的语言向这些人宣讲异端者的发明，他们会立刻掩耳，尽可能逃离，甚至无法忍受听那亵渎的言论。因此，藉着宗徒的那古老传统，他们不允许自己的思想接受这些教师怪诞言论中的任何[教义]，因为在这些教师中，既从未建立过教会，也从未建立过教义。\par

3. 在瓦伦提诺之前，那些追随瓦伦提诺的人并不存在；在马尔基翁之前，那些来自马尔基翁的人也不存在；简言之，我上面列举的那些心怀恶意的任何人，在他们那邪说的创始者和发明者之前，都不存在。因为瓦伦提诺在希吉诺时期来到罗马，在庇护下兴盛，并留到阿尼塞托时期。同样，马尔基翁的前辈切尔东本人也在希吉诺时期来到罗马，希吉诺是第九位主教。他经常进入教会，公开宣认，因此如此留下，有时秘密教导，然后再次公开宣认；但最后，因败坏教导被指控，他从弟兄的集会中被逐出。随后，马尔基翁继任，在阿尼塞托时期兴盛，阿尼塞托是第十位主教。但其余被称为诺斯替派的人，源自西满的门徒梅南德，如我已表明的；他们每一位都似乎是那教义的父亲和司祭长，且已被引入其中。但所有这些（马尔科西昂派）都是在更晚的时候，甚至在教会的中间时期，爆发了他们的背道。\par

\SourceRecord{Translated by Alexander Roberts and William Rambaut.；Ante-Nicene Fathers；Vol. 1.；Edited by Alexander Roberts, James Donaldson, and A. Cleveland Coxe.；Buffalo, NY: Christian Literature Publishing Co.,；1885.；<http://www.newadvent.org/fathers/0103304.htm>.}

% structure: p0001 source=h1 target=section reason=page-title

\section{驳斥异端（第三卷第五章）}

1. 所以，既然宗徒的传授在教会中如此存在，并在我们中间持久不变，让我们回到那些也写了福音的宗徒所提供的圣经证据吧。他们在福音中记录了关于天主的道理，指明我们的主耶稣基督就是真理\BibleRef{若 14:6}，祂内毫无谎言。正如达味也预言了祂由童贞女所生以及从死者中复活，说：\Quote{真理从地上生出。}宗徒们同样身为真理的门徒，超越一切谎言；因为谎言与真理不相交，正如黑暗与光明不相交，但一方的出现排斥另一方。所以，我们的主既是真理，就不说谎言；祂知道那些从缺陷中产生的，绝不会承认其为天主，即万有的天主，至高的君王，也是祂自己的圣父——将不完美的当作完美的，属魂的当作属神的，将在\OriginalTerm{圆满}{Pleroma}之外的当作在圆满之内的。祂的门徒也没有提及任何别的神，或称任何别的为主，只有那真正是万有的天主和主的一位，正如这些最虚妄的诡辩家所声称的——宗徒们确实以虚伪的方式按听者的能力来教导他们的教义，并根据提问者的意见来回答问题——为盲人编造盲目的东西，按照他们的盲目；为迟钝者编造迟钝的东西；为陷入错误者编造符合他们错误的东西。对于那些想象\OriginalTerm{德缪革}{Demiurge}是独一天主的人，他们传讲他；但对于那些能够领悟不可言说的圣父的人，他们通过比喻和谜语来宣布那不可言说的奥秘：这样，主和宗徒执行教师的职务，不是为了促进真理的事业，反而是出于虚伪，并按照各人所能接受的！\par

2. 这样的做法不属于那些医治或赐予生命的人；反而属于那些带来疾病、增加无知的人；而法律会比这些人更真实，它宣告凡是使盲人在路上走迷的，都是可诅咒的。因为宗徒们被派遣去寻找迷失的人，为眼睛看不见的人作光明，为软弱的人作医药，他们肯定不是按他们当时的意见来对他们说话，而是按照启示的真理。因为如果一个人劝告即将跌入悬崖的盲人继续走那条最危险的道路，好像那是正确的道路，好像他们可以安全地前行，那么任何这类人都不会合宜地行事。或者，有哪个渴望治好病人的医生，会按照病人的任性来开处方，而不是按照所需的药物呢？但主来作病人的医生，祂亲自宣告说：\Quote{健康的人不需要医生，而需要医生的是有病的人；我来不是召义人，而是召罪人悔改。}\BibleRef{路 5:31} 那么，病人如何得坚强，罪人如何悔改呢？是靠坚持同样的行为吗？或者相反，是靠经历一个巨大的改变和彻底扭转他们以前的生活方式——他们由此给自己带来了不少的病痛和许多的罪？而无知，这一切之母，被知识驱除。所以主常把知识赐给祂的门徒，祂也习惯借此医治那些受苦的人，并阻止罪人犯罪。因此，祂不是按照他们原有的观念对他们说话，也不是按照提问者的意见回答他们，而是按照导向救恩的道理，毫无虚伪或偏私。\par

3. 这也从主的话中清楚显明，祂真实地向受割礼的人启示了天主子——就是先知所预言的基督；也就是说，祂亲自显明自己，恢复了人的自由，并赐予他们不朽的产业。此外，宗徒们教导外邦人，要离开他们想像为神的虚妄的木石，敬拜那位创造并造了整个人类家庭的真实天主，并藉着祂的受造物养育、增长、坚强并保存他们的存在；并且他们要期待祂的儿子耶稣基督，祂用自己的血救赎了我们脱离背道，使我们成为圣洁的子民——祂将要带着祂父的权能从天降下，审判所有人，并要将天主的美好事物自由地赐给那些遵守祂诫命的人。祂在这末世出现，作房角石，将远处的和近处的聚集为一，并联合起来\BibleRef{弗 2:17}，就是那受割礼和未受割礼的，使耶斐特扩张，住在闪的帐棚中。\BibleRef{创 9:27}\par

\SourceRecord{Translated by Alexander Roberts and William Rambaut.；Ante-Nicene Fathers；Vol. 1.；Edited by Alexander Roberts, James Donaldson, and A. Cleveland Coxe.；Buffalo, NY: Christian Literature Publishing Co.,；1885.；<http://www.newadvent.org/fathers/0103305.htm>.}

% structure: p0001 source=h1 target=section reason=page-title

\section{反对异端（第三卷，第六章）}

1. 因此，主、圣神和宗徒们绝不会明确地、绝对地称那并非天主者为天主，除非祂确实是天主；他们也不会称任何祂本身为主，除非是天主圣父统御万有，以及祂的圣子从圣父那里接受了统御一切受造物的权柄，正如这段经文所说：\BibleQuote{主对我主说：你坐在我右边，等我使你的仇敌作你的脚凳。}此处（\WorkTitle{圣经}）向我们表明圣父对圣子说话；祂将万民的产业赐给了圣子，并使所有仇敌服从于祂。既然圣父是真正的主，圣子也是真正的主，圣神恰当地用\Quote{主}这一称号来指称他们。此外，论及索多玛人的毁灭，\WorkTitle{圣经}说：\BibleQuote{那时，主从天上从主那里降下硫磺与火，落在索多玛和哥摩辣身上。} \BibleRef{创 19:24} 因为这里指明，那曾与亚巴郎交谈的圣子，已领受了权柄，因他们的邪恶审判索多玛人。以下的经文也宣告同样的真理：\BibleQuote{天主，你的宝座永远长存；你王国的权杖是正直的权杖。你喜爱正义，憎恨不义；因此天主，你的天主，用喜乐油膏了你。} 因为圣神用\Quote{天主}这一称号称呼两者——被膏者作为圣子，以及施膏者，即圣父。又说：\BibleQuote{天主站在诸神的会中，在诸神中施行审判。} 祂（在此）指圣父、圣子以及那些领受了义子身份的人；这些人就是教会。因为她是天主的会众，是天主——即圣子本身——亲自聚集的。关于他们，祂又说：\BibleQuote{万神之神，上主发言，召唤大地。} 这里的天主是谁？就是祂曾说的：\BibleQuote{我们的天主来临，绝不缄默；} 即圣子，祂曾向人们显现，说：\BibleQuote{我向那不寻找我的人显现。} \BibleRef{依 65:1} 但这诸神又是谁？是那些祂曾对他们说：\BibleQuote{我曾说：你们是神，都是至高者的儿子。} 无疑是指那些领受了\Quote{义子}之恩的人，\BibleQuote{藉此我们呼喊：阿爸，父啊！} \BibleRef{罗 8:15}\par

2. 因此，正如我已说过的，唯有那万有之主、万有之神，祂曾对梅瑟说：\BibleQuote{我是自有者。你要这样对以色列子民说：那自有者派遣我到你们这里来；} \BibleRef{出 3:14} 以及祂的圣子耶稣基督我们的主，祂使那些信祂名的人成为天主的子女。此外，当圣子对梅瑟说话时，祂说：\BibleQuote{我下来是要拯救这百姓。} \BibleRef{出 3:8} 因为正是祂为了人的救恩下降又上升。因此，天主藉着圣子而启示出来，圣子在圣父内，圣父在圣子内——祂是自有者，圣父为圣子作证，圣子宣告圣父。——正如依撒意亚所说：\BibleQuote{你们是我的见证，} 祂宣告，\BibleQuote{上主天主说，并是我所拣选的圣子，为叫你们认识、相信并明了我是自有者。} \BibleRef{依 43:10}\par

3. 然而，当\WorkTitle{圣经}称那些不是神的事物为\Quote{神}时，它并非如我已指出的那样，以绝对意义称它们为神，而是带有某种附加和限定，表明它们根本不是神。如达味所说：\BibleQuote{外邦人的神都是魔鬼的偶像；} 以及\BibleQuote{你们不可追随别的神。} 因为祂说\Quote{外邦人的神}——但外邦人不认识真天主——又称它们为\Quote{别的神}，这就拒绝了它们被视为神的资格。至于它们本身是什么，达味论及它们说：\BibleQuote{因为它们是魔鬼的偶像，他说。} 依撒意亚也说：\BibleQuote{愿一切亵渎天主、雕刻无用之物的人都蒙羞；我是见证，上主说。} \BibleRef{依 44:9} 祂将它们从神的类别中移除，但使用了这个词，目的是让我们知道祂在说谁。耶肋米亚也说了同样的话：\BibleQuote{那未曾创造天地的神，愿它们从天下地上消灭。} \BibleRef{耶 10:11} 因为祂附加了它们的毁灭，这就表明它们根本不是神。厄里亚在加尔默耳山召集所有以色列人时，为使他们脱离偶像崇拜，对他们说：\BibleQuote{你们摇摆在两个意见之间要到几时呢？如果上主是天主，你们就跟随祂。} \BibleRef{列上 18:21} 此外，在燔祭时，祂对偶像司祭说：\BibleQuote{你们呼求你们神的名，我呼求上主我的天主的名；那以火回应的上主，祂就是天主。} 先知的这些话证明了那些被人们视为神的东西根本不是神。祂引导他们归向祂所信的那位天主，即真正的天主；祂呼求说：\BibleQuote{上主，亚巴郎的天主，依撒格的天主，雅各伯的天主，求你今天应允我，使这百姓知道你是以色列的天主。} \BibleRef{列上 18:36}\par

4. 因此，我也呼求你，上主，亚巴郎的天主，依撒格的天主，雅各伯和以色列的天主，你是我主耶稣基督的圣父，你因你丰厚的仁慈恩待了我们，使我们认识你；你创造了天地，统御万有，是唯一真实的天主，在你之上没有别的天主；求你藉我们的主耶稣基督，赐予圣神的治理权能；求你赐给本书的每位读者认识你，你独自是天主，并在你内得以坚固，远离一切异端、无神和邪恶的教义。\par

5. 宗徒保禄也说：\BibleQuote{你们从前服事那些本来不是神的，如今却认识天主，更确切地说，是被天主所认识} \BibleRef{迦 4:8}，祂在那些不是神者与那是神者之间做了区分。再者，论及\OriginalTerm{假基督}{Antichrist}，他说：\BibleQuote{那敌对者、高举自己在一切称为神或受人敬拜者以上} \BibleRef{得后 2:4}。祂在此指出那些被不认识天主的人称为\OriginalTerm{众神}{gods}的，即\OriginalTerm{偶像}{idols}。因为万有的圣父被称为天主，也确实是；假基督被高举，不是高于祂，而是高于那些确实被称为神、却并非神者。保禄自己说这是真的：\BibleQuote{我们知道偶像算不得什么，也只有一个天主。因为即使有称为神的，或在天上，或在地上；但为我们只有一个天主，就是圣父，万物出于祂，我们也归于祂；也只有一个主耶稣基督，万物藉祂而有，我们也藉祂而有。} \BibleRef{格前 8:4} 因为他做了区分，将那些确实被称为神、却并非神的，与那唯一的圣父天主分开；万物出于祂，而他也以最明确的方式亲身宣认了唯一的主耶稣基督。但在\Quote{\InnerQuote{或在天上，或在地上}}这一语句中，他并非指\OriginalTerm{世界的塑造者}{formers of the world}，像这些教师所解释的那样；他的意思与梅瑟所说的相似：\BibleQuote{你不可为自己制造任何天主的形象，凡是天上、地上、水中之物的形象。} \BibleRef{申 5:8} 他也如此解释天上之物所指为何：\BibleQuote{恐怕你，}他说，\BibleQuote{举目望天，看见太阳、月亮、星辰、天上的万象，迷惑了，便朝拜事奉它们。} \BibleRef{申 4:19} 梅瑟本人，作为天主的人，在法朗面前确实被立为神；\BibleRef{出 7:1} 但他并不适当地被称为主，先知们也不称他为天主，而是被圣神称为\BibleQuote{梅瑟，忠信的管家和天主的仆人} \BibleRef{希 3:5}，这确实是他。\par

\SourceRecord{Translated by Alexander Roberts and William Rambaut.；Ante-Nicene Fathers；Vol. 1.；Edited by Alexander Roberts, James Donaldson, and A. Cleveland Coxe.；Buffalo, NY: Christian Literature Publishing Co.,；1885.；<http://www.newadvent.org/fathers/0103306.htm>.}

% structure: p0001 source=h1 target=section reason=page-title

\section{驳斥异端（第三卷，第七章）}

1. 至于他们声称，\Person{保禄}在《\WorkTitle{格林多后书}》中明确说：\BibleQuote{在他们身上，这世界的神明蒙蔽了不信者的心意}\BibleRef{格后 4:4}，并坚持确实有一位这世界的神，但另有一位超越一切元首、元始和权能的神，那么，如果那些自称知道天主之外的奥秘的人，竟不知道如何阅读\Person{保禄}，这并不归咎于我们。因为若有人按\Person{保禄}的习惯——正如我在别处用许多例子所表明的，他常用倒装词序——如此读这段经文：\Quote{在他们身上，天主}，然后停顿一下，稍作间歇，同时将剩余部分读作一个子句：\Quote{蒙蔽了这世界不相信之人的心意}，他就会发现真义，即这句话包含的意思是：\Quote{天主蒙蔽了这世界不信者的心意}。这一点正由那小小的间歇所表明。因为\Person{保禄}并没有说\Quote{这世界的神明}，好像承认在他之外另有别神；而是承认天主确实是天主。他说\Quote{这世界的不信者}，因为他们不会继承未来不朽的世代。我将在本著作中从\Person{保禄}本人那里表明，天主是如何蒙蔽了不信者的心意，以便我们现在不因广泛游移而分散对当前问题的注意力。\par

2. 从许多其他例子中，我们也可以发现，宗徒由于他讲论的急促和他内里圣神的推动，常在他的句子中使用倒装词序。例如在《\WorkTitle{迦拉达书}》中，他这样表达：\BibleQuote{那么，为什么有行为的律法呢？它是加添的，直到那应许的后裔来到；由天使藉着中保之手所制定的。}\BibleRef{迦 3:19} 因为词序是这样：\Quote{那么，为什么有行为的律法呢？由天使藉着中保之手所制定的，它是加添的，直到那应许的后裔来到}——人这样发问，圣神作出回答。此外，在《\WorkTitle{得撒洛尼后书}》中，论及假基督，他说：\BibleQuote{那时，那无法无天的人将要显露出来，主耶稣将以口中的气息杀死他，并以他的来临的显现消灭他；他的来临是依照撒殚的运作，具有一切异能、征兆和欺骗的奇迹。}\BibleRef{得后 2:8} 现在在这些句子中，词序是这样：\Quote{那时，那无法无天的人将要显露出来，他的来临是依照撒殚的运作，具有一切异能、征兆和欺骗的奇迹，主耶稣将以口中的气息杀死他，并以他的来临的显现消灭他。} 因为他不是说主的来临是依照撒殚的运作，而是那无法无天之人的来临，我们也称他为假基督。因此，若有人不注意经文的正确读法，不按呼吸的间歇停顿，那么不仅会产生矛盾，而且读的时候还会说出亵渎的话，仿佛主的来临可能依照撒殚的运作。所以，在这类经文中，必须通过朗读来展现\Emph{倒装法}，并保持宗徒的意思连贯；因此我们在那段经文中所读的不是\Quote{这世界的神明}，而是\Quote{天主}，我们确实称他为天主；我们听到的是说这世界的不信者和被蒙蔽者，他们将不会继承那将要来临的生命之世。\par

\SourceRecord{Translated by Alexander Roberts and William Rambaut.；Ante-Nicene Fathers；Vol. 1.；Edited by Alexander Roberts, James Donaldson, and A. Cleveland Coxe.；Buffalo, NY: Christian Literature Publishing Co.,；1885.；<http://www.newadvent.org/fathers/0103307.htm>.}

% structure: p0001 source=h1 target=section reason=page-title

\section{驳斥异端（第三卷，第八章）}

1. 既然这些人的诽谤已被驳倒，便清楚证明：先知和宗徒从未称另一位为天主，或呼他为主，除非是那位真实而唯一的天主。更不用说主自己了，祂也曾指示我们\BibleQuote{凯撒的归凯撒，天主的归天主}；\BibleRef{玛 22:21}，明确称凯撒为凯撒，却宣认天主为天主。同样，那说\BibleQuote{你不能事奉两个主人}，\BibleRef{玛 6:24} 的经文，主自己也解释，说\BibleQuote{你不能事奉天主又事奉\OriginalTerm{玛门}{mammon}}。祂确实承认天主为天主，但提到玛门，这是一种存在之物。当祂说\BibleQuote{你不能事奉两个主人}时，祂并未称玛门为主，而是教导那些事奉天主的门徒，不要受玛门的支配，也不要被它统治。因为祂说\BibleQuote{犯罪的是罪恶的奴隶。}\BibleRef{若 8:34} 既然祂称那些事奉罪恶的人为\BibleQuote{罪恶的奴隶}，却当然没有称罪恶本身为天主，同样，祂也称那些事奉玛门的人为\Quote{玛门的奴隶}，而不称玛门为天主。因为玛门，按犹太人的用语（撒玛黎雅人也使用），是一个\OriginalTerm{贪恋者}{covetous}，即一个想要拥有超过应有之物的人。但按希伯来语，它通过添加一个音节（\Emph{附属音节}）称为\OriginalTerm{玛末尔}{Mamuel}，意为\OriginalTerm{贪食者}{gulosum}，即一个喉咙贪得无厌的人。因此，根据这两点所示，我们不能事奉天主又事奉玛门。\par

2. 但是，当祂说起魔鬼是强者时，并非绝对地强，而是相对于我们而言，主在各方面都显示并真实地是那位强者，祂说：\BibleQuote{除非先捆绑那强者，否则没有人能抢夺强者的财物，然后才能洗劫他的家}。\BibleRef{玛 12:29} 我们曾是这位强者的器物和家，因为我们处于背逆状态；他随意使用我们，不洁的灵居住在我们内。因为他并不强于那捆绑他、洗劫他家者；而是相对于那些是他的工具的人，因为他使他们思想偏离天主：主从他的手爪中夺回了他们。正如耶肋米亚也宣告：\BibleQuote{上主赎回了雅各伯，从比他更强者的手中把他夺回。}\BibleRef{耶 31:11} 那么，如果祂没有指出那捆绑并抢夺他财物者，而只是说他为强者，那强者就应是不被制服的。但祂又加上了那夺取并保持占有者；因为捆绑者制胜，而被捆绑者受制。祂这样做毫无比较可言，使这背逆的奴仆不能与主相比：因为不仅是他，连一切受造和隶属之物，都永远不能与天主圣言——我们的主耶稣基督——相比，万物都是藉着祂造成的。\par

3. 因为万物——无论是天使、总领天使、宝座、宰制——都是那超越万有的天主藉着祂的圣言而建立和创造的，若望这样指明了。当他说到天主的圣言曾在圣父内时，他补充说：\BibleQuote{万物是藉着祂造成的；凡受造的，没有一样不是由祂造成的。}\BibleRef{若 1:3} 达味在列举赞美之后，也一一提到了我所提及的一切，包括诸天和其中的权能：\BibleQuote{因祂命令，它们被创造；祂发言，它们就被形成。}那么，祂命令了谁？无疑是圣言，\BibleQuote{藉着祂，诸天被建立，以及它们的一切军威因祂口中的气息而成。}但祂曾自由地、随己意创造万物，达味又说：\BibleQuote{但我们的天主在天上，在世上；祂创造了凡祂所喜爱的一切。}然而，被建立的事物与建立者不同，受造之物与创造者不同。因为祂自身非受造，无始无终，无所匮乏。祂自足，并且还将存在这一事实赐予其他一切；但由祂所造的事物有一个开始。但凡有开始、易于解体、依赖并需要那创造者的事物，必定在各方面有一个不同的名称（即使对仅有中等辨别力的人也是如此）；这样，那创造万物的唯一者偕同祂的圣言才能被适当地称为天主和主；但受造之物不能被赋予这个称号，也不应合理地占有属于造物主的名称。\par

\SourceRecord{Translated by Alexander Roberts and William Rambaut.；Ante-Nicene Fathers；Vol. 1.；Edited by Alexander Roberts, James Donaldson, and A. Cleveland Coxe.；Buffalo, NY: Christian Literature Publishing Co.,；1885.；<http://www.newadvent.org/fathers/0103308.htm>.}

% structure: p0001 source=h1 target=section reason=page-title

\section{驳斥异端（第三卷，第九章）}

1. 既然这一点已在此清楚证明（且将来还会更清楚地证明），即先知、宗徒、主基督亲自都不承认任何其他的主或神，只承认至高的天主与主：先知和宗徒明认圣父与圣子，但不称任何别的为神，也不承认任何别的为主；而主亲自传授给祂的门徒：圣父是唯一的天主和主，祂独自是万物之天主和统治者。那么，我们若真是祂的门徒，就必须遵循他们在这方面的见证。宗徒玛窦知道那向亚巴郎应许使他的后裔如天上星辰的\BibleRef{创 15:5}同一位天主，也正是那通过祂的圣子基督耶稣，从崇拜石头中召叫我们认识祂自己的天主，使那非子民成为子民，那不被怜爱的成为蒙怜爱的\BibleRef{罗 9:25}。他宣布：若翰在预备基督的道路时，对那夸耀按血统与亚巴郎有关系、却心沾染各种邪恶的人，宣讲那要引领他们脱离恶行的悔改，说：\BibleQuote{毒蛇的种类！谁指示你们逃避那将要来的忿怒？你们要结出与悔改相称的果实。不要心里说：我们有亚巴郎为父。我告诉你们：天主能从这些石头中给亚巴郎兴起子孙来。}\BibleRef{玛 3:7}因此，他向他们宣讲脱离邪恶的悔改，但他没有向他们宣告另一位神，除了那向亚巴郎作出应许的神。他是基督的前驱，关于他，玛窦又记载，路加也同样记载：\BibleQuote{这人就是那藉先知依撒意亚所说的：在旷野有人声喊着说：预备主的道，修直我们天主的路径。一切山洼都要填满，大小山岗都要削平；弯弯曲曲的地方要改为正直，高高低低的道路要改为平坦。凡有血气的，都要看见天主的救恩。}\BibleRef{玛 3:3}因此，只有一位且同一位天主，就是我们主的父，祂也曾藉先知应许要派遣祂的前驱；并且祂使祂的救恩——即祂的圣言——给一切有血肉的显现，\OriginalTerm{圣言}{Word}亲自降生成人，使万有的君王在一切事上显明出来。因为受审判者必须看见审判者，并认识他们从谁领受审判；同样，那些将要进入光荣的人也应该认识那赐予他们光荣恩赐的那一位。\par

2. 再者，玛窦在论及天使时说：\BibleQuote{主的天使在梦中向若瑟显现。}\BibleRef{玛 1:20}他亲自解释这是哪位主：\BibleQuote{这是为应验主藉先知所说的话：我从埃及召叫了我的儿子。}\BibleRef{玛 2:15}\BibleQuote{看，贞女将怀孕生子，人要称祂的名字为厄玛奴耳，翻出来就是：天主与我们同在。}\BibleRef{玛 1:23}达味也论到那由贞女而生的厄玛奴耳：\BibleQuote{求祢不要转脸不顾祢的受膏者。主向达味起了誓，决不反悔：我要从你身所生的后裔，立在你宝座上。}\BibleRef{咏 132:10-11}又说：\BibleQuote{在天主被认识于犹大地，祂的居所设立在平安中，祂的住处设在熙雍。}\BibleRef{咏 76:2-3}因此，只有一位且同一位天主，曾由先知所预言并由福音所宣告；祂的儿子，就是达味身所生的后裔，即出自达味家族的贞女，就是厄玛奴耳。关于祂的星，巴郎也曾预言：\BibleQuote{有星要出于雅各伯，有杖要兴起于以色列。}\BibleRef{户 24:17}玛窦说，从东方来的\Person{贤士}喊道：\BibleQuote{我们在东方看见了祂的星，特来朝拜祂。}\BibleRef{玛 2:2}他们被那星引领进入雅各伯的家，到了厄玛奴耳面前，藉着所献的礼物表明所朝拜的是谁：\Emph{没药}，因为祂是那将为可朽的人类而死并被埋葬的；\Emph{黄金}，因为祂是君王，\BibleQuote{祂的王国无穷无尽}\BibleRef{路 1:33}；\Emph{乳香}，因为祂是天主，\BibleQuote{曾在犹大地被认识}，并\BibleQuote{向那些不寻求祂的人显明}。\par

3. 接着，论及祂的圣洗，\Person{玛窦}说：\BibleQuote{天开了，祂看见天主圣神如同鸽子降在祂身上；并且有声音从天上说：这是我的爱子，我所喜悦的。}\BibleRef{玛 3:16}因为那时基督并没有降到耶稣身上，也不是基督有一位而耶稣有另一位：而是天主圣言——祂是万民的救主，天地的主宰，祂就是耶稣，正如我已经指出过的，祂又取了肉身，并由圣父藉圣神受傅——成为耶稣基督，正如\Person{依撒意亚}也说：\BibleQuote{从\Person{叶瑟}的根上将发出一个嫩枝，从他的根上将开出一朵花；天主圣神将安息在祂身上：智慧和聪敏之神，超见和刚毅之神，明达和孝爱之神，以及敬畏上主之神将充满祂。祂不按外表审判，也不按言语责斥；却要为谦卑的人伸张正义，责斥地上的高傲者。}\BibleRef{依 11:1}\Person{依撒意亚}又提前指明祂的傅油以及祂受傅的理由，亲自说：\BibleQuote{天主圣神临到我身上，因为祂给我傅了油；祂派遣我向卑微者传报福音，医治心灵破碎的人，向俘虏宣告释放，向盲者宣告复明；宣布上主恩慈之年和报仇之日；安慰一切哀恸的人。}\BibleRef{依 61:1}因为天主圣言是来自\Person{叶瑟}之根的人，是\Person{亚巴郎}之子，在此意义上，天主圣神安息在祂身上，并傅油于祂，向卑微者传报福音。但既然祂是天主，祂就不按外表审判，也不按言语责斥。因为\BibleQuote{祂不需要任何人给祂作证关于人，因为祂自己知道在人内有什么。}\BibleRef{若 2:25}祂召叫一切哀恸的人；并赐予那些被自己的罪引入俘虏的人宽恕，解开他们的锁链，关于这些人，\Person{撒罗满}说：\BibleQuote{每人必被自己的罪绳索束缚。}\BibleRef{箴 5:22}因此，天主圣神——那位曾藉先知们许诺要傅油于祂的神——降临到祂身上，好使我们从祂丰富的傅油中领受，得以得救。\Person{玛窦}的［见证］就是这样。\par

\SourceRecord{Translated by Alexander Roberts and William Rambaut.；Ante-Nicene Fathers；Vol. 1.；Edited by Alexander Roberts, James Donaldson, and A. Cleveland Coxe.；Buffalo, NY: Christian Literature Publishing Co.,；1885.；<http://www.newadvent.org/fathers/0103309.htm>.}

% structure: p0001 source=h1 target=section reason=page-title

\section{驳斥异端（第三卷，第十章）}

1. 路加也是宗徒的追随者和门徒，提及匝加利亚和依撒伯尔，按恩许，若翰正是由他们而生，他说：\Quote{他们二人在天主面前是义人，遵行主的一切诫命和典章，无可指摘。}\BibleRef{路 1:6} 又说及匝加利亚：\Quote{当他按着班次，在天主面前执行司祭职务时，按照司祭职的惯例，拈阄烧香；}然后他前来献祭，\Quote{进入主的圣殿。}\BibleRef{路 1:8} 那位天使加俾额尔，既是显赫地站在主面前的，便单纯、绝对而坚定地以自己的身份宣认为天主和主的那一位，即拣选了耶路撒冷并设立了司祭职务的那一位。因为他知道没有别位在他以上；因为，如果他认识任何另一位比他更完美的天主和主，他就绝不会——正如我已指出的——把那位他知道是缺陷的果实的，绝对而完全地宣认为天主和主。然后，论及若翰，他这样说：\Quote{他在主面前将要为大，许多以色列子民要使他们归向主、他们的天主。他要以厄里亚的精神和能力走在主前面，为主预备一个完善的百姓。}\BibleRef{路 1:15} 那么，他预备百姓是为谁，他在哪位主面前为大呢？确实是那位曾说若翰甚至\Quote{比先知还大，}\BibleRef{玛 11:9}并且\Quote{在妇人所生者中，没有比洗者若翰更大的；}的那一位；他也为主的来临预备百姓，警告他的同仆们，向他们宣讲悔改，使他们得以在主来临时从祂那里获得赦免，他们曾因罪恶和过犯而与主疏离，如今回转归向祂。正如达味也说：\Quote{离弃的人从母腹便是罪人：他们一出母胎就走错了路。}正是为此，他引领他们归向他们的主，以厄里亚的精神和能力，为主预备了一个完善的百姓。\par

2. 再者，提及天使时，他说：\BibleQuote{那时，天使加俾额尔受天主派遣，向童贞女说：玛利亚，不要害怕，因为你在天主面前蒙受了恩宠。}\BibleRef{路 1:26} 他论及主说：\BibleQuote{他将为上主至高者的儿子，上主天主要把他祖先达味的御座赐给他。他要为王统治雅各伯家，直到永远；他的王权没有终结。}\BibleRef{路 1:32} 因为除了我们的主耶稣基督、至高天主之子以外，还有谁能永远不间断地统治雅各伯家呢？祂曾藉法律和先知应许，要使祂的救恩在万民面前显现；为此，祂成为人子，好使人也能成为天主子。玛利亚为此欢跃，代表教会大声预言说：\BibleQuote{我的灵魂颂扬上主，我的心神欢跃于天主我的救主。因为他扶助了他的仆人以色列，纪念他的仁慈，正如他向我们的祖先亚巴郎和他的后裔所说过的，直到永远。}\BibleRef{路 1:46} 藉着这些以及类似的话，福音指出：是那位曾与祖先发言的天主；是那位藉梅瑟设立了法律制度的天主；藉这法律的颁赐，我们知道他曾经与祖先说话。这同一天主，在他极大的美善之后，将他的怜悯倾注在我们身上，藉此怜悯，\BibleQuote{旭日由高天照耀我们，照亮那些坐在黑暗和死影中的人，并引领我们的脚步走上和平的道路。}\BibleRef{路 1:78} 同样，匝加利亚从他因不信所患的哑症中恢复后，充满了新的精神，以新的方式赞美了天主。因为一切事物都进入了新的阶段，圣言以新的方式安排了肉身的降临，好将那离开了天主的\OriginalTerm{人性}{hominem}重新带回天主面前；因此，人们被教导以新的方式敬拜天主，而非另一位神，因为实际上只有\BibleQuote{一位天主，他要因信德使受割礼的成义，也要藉信德使未受割礼的成义。}\BibleRef{罗 3:30} 但匝加利亚预言时，高声说：\BibleQuote{上主，以色列的天主应受赞美，因他眷顾救赎了自己的百姓，并在他的仆人达味家中，为我们兴起了大能的救主，正如他借历代诸圣先知的口所说过的，拯救我们脱离敌人和仇恨我们者的手。他向我们的祖先施行仁慈，记忆起他的神圣盟约，就是他向我们的父亲亚巴郎所起的誓，恩赐我们，使我们从敌人手中被救出后，无惧地在他面前，以圣善和正义事奉他，终生不渝。}\BibleRef{路 1:68} 然后他向若翰说：\BibleQuote{孩子，你要称为至高者的先知，因为你要走在上主前面，预备他的道路，使他的百姓认识救恩，以获得他们罪恶的赦免。}\BibleRef{路 1:76} 这就是他们欠缺的救恩知识，即关于天主子的知识，若翰宣讲说：\BibleQuote{请看，天主的羔羊，除免世罪者！这就是我曾说过的那位：在我以后要来的一个人，却成了在我以前的，因为他本来先我而有；从他的圆满中我们众人都领受了。}\BibleRef{若 1:29} 因此，这就是救恩的知识；但这并不在于另一位神，另一位父，或\OriginalTerm{深渊}{Bythus}，或三十永世的\OriginalTerm{圆满}{Pleroma}，或（下层）\OriginalTerm{八质}{Ogdoad}之母；而在于认识天主子，他被称作且实际就是救恩、救主和赋予救恩者。祂是救恩，正如经上说：\BibleQuote{上主，我期待你的救恩。}\BibleRef{创 49:18} 祂是救主：\BibleQuote{看啊！天主是我的救主，我要依靠他。}\BibleRef{依 12:2} 但祂也是赋予救恩者：\BibleQuote{上主在万民眼前显示了他的\OriginalTerm{救恩}{salutare}。} 因为他确实是救主，作为天主子与圣言；但又是赋予救恩者，因为他是圣神；如经上说：\Quote{我们鼻中的气息，基督上主。} 但祂也是救恩，因为是肉身：因为\BibleQuote{圣言成了血肉，居住在我们中间。}\BibleRef{若 1:14} 因此，若翰将救恩的知识传授给那些悔改并相信除免世罪的天主羔羊的人。\par

3. 他说，主的使者也向牧人们显现，向他们宣告喜乐：\Quote{因为今天在达味城中，为你们诞生了一位救主，就是基督主。}\BibleRef{路 2:11}随后出现了大队天军，讚美天主说：\Quote{光荣在天上归于天主，平安在地上归于善意的人。}那些自称\OriginalTerm{诺斯替派}{Gnostics}的人说，这些天使来自\OriginalTerm{第八层}{Ogdoad}，并显示了至高基督的降临。但他们又错了，他们说，那位来自高天的基督和救主并非诞生，而是在分施性耶稣受洗之后，他（\OriginalTerm{普累若麻}{Pleroma}的基督）像鸽子一样降在他身上。因此，按这些人的说法，第八层的天使撒谎了，因为他们说：\Quote{今天在达味城，为你们诞生了一位救主，就是基督主。}因为按他们的说法，那时既无基督也无救主诞生；那只是分施性耶稣，属于世界的创造者（\OriginalTerm{德穆革}{Demiurge}），而在他的洗礼之后，也就是三十年之后，他们主张那来自高天的救主降临在他身上。但天使为什么加上\Quote{在达味城}呢？如果他们不是宣布天主向达味所许承诺的应验，即从达味的后裔中要兴起一位永恒的君王？因为整个宇宙的创造者（\OriginalTerm{德穆革}{Demiurge}）曾向达味许诺，正如达味自己所宣告的：\Quote{我的救援来自天主，祂创造了天地。}又说：\Quote{祂手中掌握大地四极，山岳的高峰也属于祂。海洋是祂的，祂亲手创造了它；祂的手奠定了旱地。来，让我们朝拜，向祂俯伏，在创造我们的主面前哭泣，因为祂是主我们的天主。}圣神显然藉达味向那些听从他的人宣告，将会有人藐视那造了我们、且是唯一天主的祂。因此，他先前说了那些话，意思是：你们要谨慎，不要犯错；除了祂以外，没有别的神，你们不应向别的神伸开双手，从而使我们敬虔、感恩于那位创造、建立并养育我们的祂。那么，那些对造物主说了如此多亵渎话的人，将会怎样呢？天使们也宣告了这同一真理。因为当他们高呼\Quote{光荣在天上归于天主，平安在地上}时，他们用这些话荣耀了那位至高者（即天上诸事）的创造者，以及地上万物的建立者：祂从天上将祂救恩的祝福赐给了祂自己的手工作品，即人类。因此，他补充道：\Quote{牧人们回去了，为所有他们听到和看见的、正如天使所说的事，光荣天主。}\BibleRef{路 2:20}因为以色列的牧人们并非光荣另一位神，而是那位由法律和先知所宣告的、万有的创造者，天使们也光荣了祂。但如果那些来自第八层的天使习惯于光荣另一位不同于牧人们所朝拜的神，那么这些来自第八层的天使带给他们的就是错误，而非真理。\par

4. 路加进一步论及主说：\Quote{到了取洁的日子，他们带着祂上耶路撒冷去，把祂献给主，正如主的法律上所记载的：凡开胎首生的男子，应称为圣，奉献给主；并按照主的法律上所说的，献上一对斑鸠或两只雏鸽。}\BibleRef{路 2:22}他本人最清楚地称那制定了法律规条的那一位为主。但他又说：\Quote{西默盎讚美天主说：主啊！现在可照祢的话，放祢的仆人平安去吧！因为我亲眼看见了祢的救援，即祢在万民面前所准备的：为启示外邦人的光明，和祢的子民以色列的荣耀。}\BibleRef{路 2:29}又说：\Quote{女先知亚纳}\BibleRef{路 2:38}同样，她看见了基督，就光荣天主，\Quote{并向所有期待耶路撒冷得救赎的人讲论祂。}这一切都昭示了同一天主，祂藉着祂圣子的新来临，向人类揭示了自由的新安排、新盟约。\par

5. 因此，马尔谷，伯多禄的译员和跟随者，这样开始他的福音叙述：\Quote{天主之子耶稣基督福音的开始，正如先知书上所记载的：看，我派遣我的使者在你面前，他要在你前面预备你的道路。旷野中有呼喊者的声音：预备主的道路，修直我们天主的途径。}\BibleRef{路 1:17}福音的开端明确引用圣先知的话，并立刻指明他们承认为天主和主的那一位；即我们的主耶稣基督的父，祂曾向祂许诺，要派遣祂的使者在祂前面，就是若翰，在旷野中呼喊，带着\Quote{厄里亚的精神和能力}\BibleRef{路 1:17}\Quote{预备主的道路，修直我们天主的途径。}因为先知们并非宣告一位又一位神，而是同一位；只是以不同的方面和许多名称。因为父是多方面且富于属性的，正如我在前书中已表明的；我在本卷后续部分也要从先知本身证明这同一真理。同样，在他的福音结尾处，马尔谷说：\Quote{主耶稣对他们说了话以后，就被接升天，坐在天主的右边。}\BibleRef{谷 16:19}这就证实了先知所说的：\Quote{上主对我主说：你坐在我右边，等我把你的仇敌放在你脚下。}因此，天主和父确实是同一位；祂是由先知们所宣告、并由真福音所传承的；我们基督徒全心全意朝拜和爱慕祂，作为天地及其中万物的创造者。\par

\SourceRecord{Translated by Alexander Roberts and William Rambaut.；Ante-Nicene Fathers；Vol. 1.；Edited by Alexander Roberts, James Donaldson, and A. Cleveland Coxe.；Buffalo, NY: Christian Literature Publishing Co.,；1885.；<http://www.newadvent.org/fathers/0103310.htm>.}

% structure: p0001 source=h1 target=section reason=page-title

\section{驳斥异端（第三卷，第十一章）}

1. \Person{若望}，主的门徒，宣讲这信德，并藉着宣讲福音，试图消除那个由\Person{塞林特}在人中间散布的、以及更早由所谓\Person{尼苛劳党人}（他们是那错误称为\Quote{知识}的一个分支）所传播的谬误，为要驳倒他们，说服他们只有一位天主，祂藉其圣言创造了万物；而非如他们所声称：造物主是一位，而主的父是另一位；并且造物主的儿子（诚然）是一位，而从上而来的基督是另一位，祂也是不受苦难的，降临在造物主之子\Person{耶稣}身上，又飞回祂的\OriginalTerm{圆满}{Pleroma}中；以及\OriginalTerm{独生子}{Monogenes}是元始，而\OriginalTerm{圣言}{Logos}是独生子的真子；并且我们所属的这个世界并非由首要的天主所造，而是由某种远低于祂、且与不可见和不可言说之事隔绝的能力所造。因此，主的门徒为终结所有这类学说，并在教会中确立真理的准则，即有一位全能的天主，祂藉其圣言创造了万物，无论可见的与不可见的；同时指明，天主藉以创造万物的那圣言，也赐予了包含在创造中的人类救恩；如此，他在福音中开始教导说：\BibleQuote{在起初已有圣言，圣言与天主同在，圣言就是天主。这一位在起初与天主同在。万物是藉着他而造成的；凡受造的，没有一样不是由他而造成的。在他内有生命，这生命是人的光。光在黑暗中照耀，黑暗决不能胜过他。}\BibleRef{若 1:1} \Quote{万物}，他说道，\Quote{是藉着他而造成的；}因此，在\Quote{万物}中，我们这个世界也包含在内，因为我们不能向这些人让步说\Quote{万物}这个词是指他们圆满中的那些事物。因为如果他们的圆满确实包含这些，那么这个世界，既然是这个世界，就不是在它之外，正如我在前一卷中所证明的；但如果他们在圆满之外（这似乎是不可能的），那么他们的圆满就不能称为\Quote{万物}：因此这个浩大的世界并不在圆满之外。\par

2. 然而，\Person{若望}亲自消除了我们这边的所有争议，当他说：\BibleQuote{他在世界上，世界原是藉着他造成的，但世界却不认识他。他来到自己的领域，自己的人却没有接受他。}\BibleRef{若 1:10} 但是按照\Person{马尔西昂}和像他那样的人，世界并非藉着他造成的；他也不是来到自己的领域，而是来到别人的领域。而且，按照某些\Person{诺斯底人}，这个世界是由天使造成的，而非由天主的圣言。但按照\Person{华伦提努}的追随者，世界并非由他造成，而是由\OriginalTerm{德穆革}{Demiurge}造成的。因为\OriginalTerm{救主}{Soter}按照他们的说法，促使了这些相似物按照上面模型的样式被造；但德穆革完成了创造的工程。因为他们说，他（这个创造计划的领主和创造者，他们认为这个世界是藉着他造的）是从\OriginalTerm{母亲}{Mother}产生出来的；而福音明确宣称，万物是由那在起初与天主同在的圣言造成的，那圣言，他说，\BibleQuote{成了血肉，住在我们中间。}\BibleRef{若 1:14}\par

3. 但是，按照这些人，圣言并没有成为血肉，基督也没有，从所有\OriginalTerm{永世}{Æons}联合的贡献中产生出来的\OriginalTerm{救主}{Soter}也没有。因为他们主张圣言和基督从未进入这个世界；救主也没有降生成人，也没有受苦，而是像鸽子一样降临在\OriginalTerm{经世耶稣}{dispensational Jesus}身上；并且一旦宣告了那未知的父，他便再次升入圆满。有些人却断言这位经世耶稣确实降生成人并受苦了，他们描述他像水通过管子一样穿过\Person{玛利亚}；而另一些人则声称他是德穆革的儿子，经世耶稣降临在他身上；又有些人说\Person{耶稣}是从\Person{若瑟}和\Person{玛利亚}生的，而那从上而来的基督——没有肉身、不受苦难——降临在他身上。但是，按照任何一个异端者的意见，天主的圣言都没有成为血肉。因为如果任何人仔细检查他们所有人的体系，就会发现天主的圣言被他们所有人都引入为没有成为血肉\OriginalTerm{无肉}{sine carne}且不受苦难的，从上而来的基督也是如此。另一些人认为他显现为一个变形的人；但他们坚持他既非受生也非成为血肉；而另一些人（认为）他根本没有采取人的形象，而是像鸽子一样降临在那由\Person{玛利亚}所生的\Person{耶稣}身上。因此，主的门徒指出他们都是假见证，说：\BibleQuote{圣言成了血肉，住在我们中间。}\BibleRef{若 1:14}\par

4. 为了使我们不必追问：圣言是成了哪一位天主的血肉？他自己先教导我们说：\Quote{曾有一人，由天主派遣，名叫若翰。这人来，是为作证，给那光作证。他不是那光，而是来为光作证的。} \BibleRef{若 1:6} 那么，作为前驱的若翰，他见证那光，是由哪一位天主派遣而来的呢？确实是由那一位，加俾额尔是祂的天使，这位天使也曾宣布他诞生的喜讯：（就是那一位）天主，祂曾藉先知们应许，要在祂儿子面前派遣祂的使者，\BibleRef{拉 3:1} 预备祂的道路，即他要以厄里亚的精神和能力，为那光作证。\BibleRef{路 1:17} 可是，厄里亚又是哪一位天主的仆人和先知呢？是那位创造天地者，正如他自己所承认的。那么，既然若翰是由这个世界的创造者和建立者所派遣，他又如何能为那从不可言说、不可见之处降临的光作证呢？因为所有异端者都断定，\OriginalTerm{神工者}{Demiurge}不认识在他之上的那一位权势，而若翰正是那位权势的见证者和先驱。因此，主说祂认为他\Quote{比先知还大}。\BibleRef{玛 11:9} 因为所有其他先知都预言了父光的来临，并渴望配得看见他们所预言的这一位；但若翰既如其他人一样预先宣布了（祂的来临），又在祂来临时确实看见了祂，并指出祂，使许多人相信祂，以致他本人同时扮演了先知和宗徒的角色。因为这正是比先知还大，因为\Quote{第一是宗徒，第二是先知}；\BibleRef{格前 12:28} 但这一切都出自同一天主。\par

5. 那由天主在葡萄园中所产、且首先被饮用的酒是好的。\BibleRef{若 2:3} 饮用这酒的人无一挑剔；主也饮了它。然而那由圣言从水立刻变成的酒更好，纯粹是为那些被召赴婚筵的人所用。因为虽然主有能力不依靠任何受造物而为赴宴者提供酒，并使饥饿者饱足，但他没有这样做；而是拿取大地所产的饼，祝谢后，\BibleRef{若 6:11} 另一次则将水变成酒，使那些斜靠（在桌前）的人满足，并给那些被请赴婚筵的人饮酒，以此表明：那创造大地、命令它结果实，建立诸水、引出泉源的天主，正是那在末世藉其圣子，将食物的祝福和饮料的恩赐赐予人类的天主；\OriginalTerm{不可理解者}{Incomprehensible}藉由可理解者行动，不可见者藉由可见者行动；因为除祂之外没有别的，祂存在於圣父怀中。\par

6. 因为他说：\Quote{从来没有人见过天主，}唯有\Quote{天主的独生子，在圣父怀里的，祂所启示的。} \BibleRef{若 1:18} 因为那在圣父怀里的圣子，向众人启示那不可见的圣父。所以，圣子向谁启示圣父，\Emph{他们}就认识祂；同样，圣父藉圣子使那些爱祂的人认识祂的圣子。连纳塔乃耳也由此受教而认识（祂），主也为他作证说，他是个\Quote{真以色列人，在他内毫无诡诈}。\BibleRef{若 1:47} 那以色列人认识了他的君王，于是向他喊说：\Quote{辣彼，祢是天主子，祢是以色列的君王。} 连伯多禄也由此受教，承认基督是永生天主之子，当时（天主）说：\Quote{看哪，我的爱子，我所喜悦的：我要把我的神放在祂身上，祂要向外邦人宣扬公理。祂不争辩，也不呼喊，街上也没有人听见祂的声音。压伤的芦苇祂不折断，将残的灯火祂不吹灭，直到祂使公理得胜；外邦人都要仰望祂的名。}\par

7. 那么，福音的首要原则便是：有一位天主，是这宇宙的创造者；祂也是先知们所预言的那一位，并由梅瑟规定了法律的经济：这些原则宣告了我们的主耶稣基督的圣父，并否认有除祂之外的任何其他神或父。这些福音所依据的基础是如此稳固，以致连异端者自身也为它们作证，并且每个人都从这些（文献）出发试图建立自己的独特教义。因为那些只使用玛窦福音的\OriginalTerm{厄彼翁派}{Ebionites}，被这同一福音所驳斥，他们对主作了错误的假设。\Person{马尔基昂}则删减了路加福音，藉由他仍然保留的（段落）证明自己是唯一存在天主的亵渎者。那些将耶稣与基督分开、声称基督是不可受难、而受苦的是耶稣的人，偏爱马尔谷福音，若以爱真理之心阅读它，他们的错误可得纠正。此外，那些追随\Person{瓦伦提努斯}的人，大量使用若望福音来图解他们的配对，藉由这同一福音将被证明完全错误，正如我在第一卷中所指出的。既然我们的对手为我们作证并使用这些（文献），那么我们从它们得出的证明是坚实而真实的。\par

8. 福音书的数目不可能比现在更多或更少。因为，既然我们所居住的世界有四个区域，有四种主要的风，而教会遍布整个世界，教会的\BibleQuote{柱石和基础}\BibleRef{弟前 3:15}是福音和生命之神，那么她拥有四根柱石，从四方吹出不朽之气，重新赋予人生命，是合宜的。由此明显可见，圣言，万物的创造者，那居于革鲁宾之上，包含万有，并向人显示自己的那一位，以四种形式将福音赐给了我们，但由同一圣神联结在一起。正如\Person{达味}在恳求祂显现时所说：\Quote{你坐在革鲁宾之上，请显现。}因为革鲁宾也是四面的，他们的面容是圣子之救恩计划的图像。因为，（如圣经）说：\BibleQuote{第一个活物像狮子}\BibleRef{默 4:7}，象征祂有效的工作、祂的领导地位和君王权能；第二个活物像牛犊，象征祂的祭献和司祭职；\BibleQuote{第三个活物有面如同人的脸面}——明显地描述了祂作为人的来临；\BibleQuote{第四个活物像飞鹰}，指出圣神之恩赐以翅膀覆盖教会。因此，福音与这些事一致，而基督耶稣就坐于其中。因为按照\Person{若望}的福音，叙述祂从圣父而来的原始、有效和荣耀的出生，如此宣告：\BibleQuote{在起初已有圣言，圣言与天主同在，圣言就是天主。}\BibleRef{若 1:1}又说：\BibleQuote{万物是藉着祂造成的，凡受造的，没有一样不是藉着祂造成的。}\BibleRef{若 1:3}为此，那部福音充满一切信心，因为祂的人格正是如此。但是按照\Person{路加}的福音，采取司祭的职分，以司祭\Person{匝加利亚}向天主献祭开始。因为肥牛犊已预备好，将要为寻找回幼子而被宰杀。\Person{玛窦}则叙述祂作为人的出生，说：\BibleQuote{亚巴郎之子，达味之子耶稣基督的族谱}\BibleRef{玛 1:1}；又说：\BibleQuote{耶稣基督的诞生是这样的。}\BibleRef{玛 1:18}那么，这就是祂人性的福音；为此，在整个福音中保持了一个谦卑温顺的人的特征。相反，\Person{玛尔谷}则以先知之神从高天降临到人开始，说：\BibleQuote{天主子耶稣基督福音的开始，正如依撒意亚先知所记载的}\BibleRef{谷 1:1}——指出福音有翅膀的特征；因此他做了一个简略而匆忙的叙述，因为先知的特征正是如此。而天主圣言自己在梅瑟之前的先祖时代，曾按照祂的神性和光荣与他们交谈；但对于那些在法律之下的，祂设立了司祭和礼仪服务。后来，为了我们而成为人，祂将天上圣神的恩赐派遣到整个大地，用祂的翅膀保护我们。因此，正如圣子所行进的途径，活物的形状也是如此；又如活物的形状，福音的特征也是如此。因为活物是四形的，福音也是四形的，如同主的途径也是四形的。为此，人类获得了四个\OriginalTerm{主要的}{\GreekText{καθολικαί}}盟约：第一个，在洪水之前，在\Person{亚当}之下；第二个，在洪水之后，在\Person{诺厄}之下；第三个，法律的颁布，在\Person{梅瑟}之下；第四个，那更新人，并以福音将万有总归于自身，将人提升并背负到天国的。\par

9. 既然如此，所有破坏福音形式的人都是虚妄、无知且狂妄的；我指的是那些认为福音的幅度比上述数目更多或更少的人。前者这样做，是为了显得他们发现了超越真理的东西；后者这样做，是为了废弃天主的救世计划。因为\Person{马尔西翁}（Marcion）既摒弃整个福音，更确切地说，他自绝于福音，却夸耀自己有分于福音的福分。另一些人（\OriginalTerm{孟他努派}{Montanists}）则为了藐视圣神的恩赐——这恩赐在末后时期，因圣父的美意沛降于人类——不接受\BibleBook{若望}{福音}所呈现的那一\Emph{幅度}（福音救世计划的），其中主应许要派遣\OriginalTerm{护慰者}{Paraclete}；\BibleRef{若 14:16} 但他们同时废弃了福音和预言之神。可怜的人！他们确实希望成为假先知，却将预言的恩赐从教会中剔除；他们的行为就像那些因虚伪之徒而回避弟兄共融的人（例如\OriginalTerm{恩克拉提派}{Encratitæ}）。此外，我们必须断定，这些人也不能接纳宗徒保禄。因为在他所写的\WorkTitle{致格林多人书}中，\BibleRef{格前 11:4} 他明确提及预言恩赐，并承认在教会中男女都可以说预言。因此，他们在所有这些事上得罪了天主圣神，\BibleRef{玛 12:31} 陷于不可赦免的罪。另一方面，来自\Person{华伦提诺}（Valentinus）的人则全然鲁莽，他们提出自己的著作，夸耀自己拥有比实际更多的福音。实际上，他们竟大胆到将他们近期著作命名为\OriginalTerm{真理福音}{Gospel of Truth}，但这著作与宗徒们的福音毫无一致之处，以致他们根本没有不充满亵渎的福音。因为如果他们出版的著作是真理的福音，却又与从宗徒传下来的那些福音全然不同，那么任何愿意学习的人都可以从圣经本身看出，从宗徒传下来的东西不能再算作真理的福音。但我已通过如此众多且有力的论据证明，唯有这些福音是真实可靠的，既不增加也不减少上述数目。因为既然天主按适当的比例和适应创造了万物，那么福音的外在形式也应当妥善安排与协调。因此，在调查了那些将福音传给我们的先辈的意见（从其源头开始）之后，让我们也继续考察其余宗徒，探究他们关于天主的教义；然后，按适当顺序，我们将聆听主亲口所说的话。\par

\SourceRecord{Translated by Alexander Roberts and William Rambaut.；Ante-Nicene Fathers；Vol. 1.；Edited by Alexander Roberts, James Donaldson, and A. Cleveland Coxe.；Buffalo, NY: Christian Literature Publishing Co.,；1885.；<http://www.newadvent.org/fathers/0103311.htm>.}

% structure: p0001 source=h1 target=section reason=page-title

\section{驳斥异端（第三卷，第十二章）}

1. 宗徒伯多禄，因此，在主复活并升天之后，渴望补全十二宗徒的数目，并在选举一位由天主所拣选的替代者来代替犹大时，对在场的人这样说：\Quote{人们和弟兄们，这圣经必须应验，圣神藉达味的口预先论及犹大，他成了引导那些捉拿耶稣的人的向导。他原是我们中的一员，\BibleRef{宗 1:16}……愿他的居所变为荒场，无人居住；并且，愿他的监督职分让别人取代。}——这样，根据达味所说的话，就完成了宗徒的数目。再次，当圣神降临在门徒身上，使他们都说预言和说方言时，有些人嘲笑他们，以为他们被新酒灌醉了。伯多禄说他们并没有醉，因为那时是第三时辰；而是那藉先知所说的：\Quote{在末日，天主说，我要将我的圣神倾注在一切血肉之人身上，他们就要说预言。}\BibleRef{岳 2:28}因此，那曾藉先知应许要赐给全人类祂的圣神的天主，正是那赐下圣神的天主；而伯多禄宣称天主自己成就了祂的应许。\par

2. 因为伯多禄说：\Quote{以色列人哪，请听我的话：纳匝肋人耶稣，是天主藉着祂在你们中间所行的奇能、奇事和征兆，向你们所证明的人，正如你们自己也知道：祂按天主的定旨和预知被交付，你们藉不法之人的手把祂钉在十字架上杀了。天主却将祂复活了，解除了死亡的痛苦，因为祂不能被死亡拘禁。因为达味论到祂说：我常将主摆在我眼前，因祂在我右边，我便不致摇动。因此，我心欢喜，我的舌快乐，并且我的肉身也要安息于希望中，因为祢必不将我的灵魂撇在阴间，也不让祢的圣者见到朽坏。}\BibleRef{宗 2:22}接着他大胆地对他们论到先祖达味，说他死了，也葬了，并且他的坟墓直到今日还在他们那里。他说：\Quote{但达味既是先知，也知道天主曾向他起誓，要从他腰中的果子立一位坐在他的宝座上；他预先看见了，就讲论基督的复活，说祂没有被撇在阴间，祂的肉身也没有见到朽坏。这位耶稣，天主已经复活了，我们都是此事的见证人。祂被高举到天主的右边，从圣父领受了所应许的圣神，就倾注了你们现在所看见所听见的这恩赐。因为达味并没有升到天上，但他自己说：主对我主说：你坐在我右边，等我使你的仇敌作你的脚凳。因此，以色列全家当确实地知道，你们所钉死的这位耶稣，天主已经立祂为主、为基督了。}\BibleRef{宗 2:30}当群众喊说：\Quote{我们该作什么？}伯多禄对他们说：\Quote{你们悔改，各人要以耶稣基督的名受洗，使罪得赦，就必领受圣神的恩赐。}\BibleRef{宗 2:37}这样，宗徒们并没有宣讲另一位神，或另一个\OriginalTerm{圆满}{Fulness}；也没有宣讲那受苦又复活的基督是一位，而那升到高处、不受苦难的基督是另一位；而是宣讲同一位天主父，和从死者中复活的基督耶稣；他们向那些不信天主之子的人宣讲信德，并劝他们从先知的话中得知：天主所应许要派遣的基督，祂就在耶稣身上派遣了，就是他们钉死而天主复活的那位。\par

3. 此外，当伯多禄同若望一起，看见那个生来瘸腿的人坐在那称为丽门的殿门口，求人施舍，就对他说：\Quote{金银我都没有，只把我所有的给你：因纳匝肋人耶稣基督的名，起来行走！}\BibleRef{宗 3:6} 那时，群众因这意外的事从各处聚集，伯多禄就对他们说：\Quote{以色列人哪，为什么因这事惊奇？为什么定睛看我们，好像我们凭自己的能力和虔诚使这人行走？亚巴郎的天主、依撒格的天主、雅各伯的天主、我们祖先的天主，光荣了自己的圣子，他就是你们所解送，并在比拉多前所否认的，虽然比拉多原意要释放他。你们否认了圣洁和正义的人，反倒要求把凶手赐给你们；你们杀害了生命的元首，天主却从死人中复活了他，我们就是这事的见证人。因信他的名，他的名就强壮了你们所看见所认识的这人；那由他而来的信德，在你们众人面前赐给了他完全的康复。如今，弟兄们！我知道你们所行的是出于无知；但天主借着众先知的口，预言他的基督当受难的事，就这样成就了。所以，你们悔改，并要回心转意，好消除你们的罪过，为使那安舒的时期从主面前来到；他必派遣预先为你们准备好的耶稣基督。基督必须留在天上，直到万物复兴的时期，就是天主借他圣先知的口所说的一切。梅瑟曾向我们的祖先说过：你们的天主上主要从你们弟兄中给你们兴起一位像我一样的先知；凡他所说的一切，你们都要听从。将来无论谁，若不听从那位先知，必从民间铲除。从撒慕尔起，以及所有随后发言的先知，都预报了这些日子。你们是先知和盟约的子孙，这盟约是天主与你们的祖先所订立的，他曾对亚伯郎说：地上万民都要因你的后裔蒙受祝福。天主先给你们兴起他的圣子，派遣他来祝福你们，使你们各人回头，离开你们的邪恶。}\BibleRef{宗 3:12} 伯多禄同若望一起，向他们宣讲了这喜讯的明白信息，即天主对祖先的应许已由耶稣实现了；他们当然不是宣扬另一位神，而是宣扬天主子，他也成了人，并受了苦；如此引导以色列人认识，并借着耶稣宣讲死人复活，\BibleRef{宗 4:2} 并指明，先知们所预言关于基督受苦的一切，天主都已实现了。\par

4. 为此，当司祭长聚集时，伯多禄充满勇气对他们说：\Quote{百姓的领袖和以色列的长老！如果今天我们因对那瘸子所行的善事，并因他怎样得到痊愈，受你们的查问，那么，你们所有的人和全以色列百姓都该知道：站在你们面前这人痊愈了，是因纳匝肋人耶稣基督的名字，他就是你们钉死，天主从死人中复活的耶稣基督。这石头是你们建筑者所弃而不用的，反而成了屋角的基石。[除他以外，别无救恩；]因为在天下人间，没有赐下别的名字，使我们赖以得救的。}\BibleRef{宗 4:8} 这样，宗徒们并没有改变天主，而是向百姓宣讲：耶稣就是被钉死的基督，那曾派遣先知的天主，就是天主自己，把他复活了，并在他内赐给人救恩。\par

5. 因此，他们既因这个治愈的例子（\BibleQuote{那得治愈的人，年纪已在四十岁以上}\BibleRef{宗 4:22}），又因宗徒们的教导和先知们的讲解而感到困惑。当时司祭长已打发伯多禄和若望离开。［他们二人］回到其余同为宗徒和主的门徒那里，即教会那里，报告了所发生的事，以及他们如何因耶稣的名字勇敢行事。然后，据说全体教会\BibleQuote{他们听了这事，就同心合意高声向天主说：\InnerQuote{主啊！你是天主，你创造了天、地、海和其中的万物；你曾藉着圣神，用你仆人、我们的祖先达味的口说过：\InnerQuote{外邦为什么嚣张？万民为什么妄想？地上的君王起来，首领聚集，反对主，反对他的基督。}因为，在这城里，确实的黑落德和般雀比拉多，同外邦人和以色列人，一起聚集，反对你的圣子耶稣，你所傅油的，所作的一切，都是你的手和你的旨意所预定的。}}\BibleRef{宗 4:24}这就是那产生一切教会之教会的声音；这就是新盟约子民之京城的声音；这就是宗徒们的声音；这就是主的门徒们，那些真正成全的人的声音，他们在主升天之后，由圣神得以成全，呼求那创造天、地、海的天主——就是先知们所预言的那一位——以及耶稣基督、他的圣子，就是天主所傅油的，他们不认识别的［神］。因为在那时、那地，既没有华伦提努，也没有玛尔强，也没有其他这些（真理的）颠覆者及其追随者。因此，万有的创造者天主垂听了他们。因为经上记载：\BibleQuote{他们聚集的地方震动起来，众人都充满了圣神，大胆地宣讲天主的话}\BibleRef{宗 4:31}，向每个愿意相信的人。又补充说：\BibleQuote{并且以极大的能力，宗徒们为主耶稣的复活作见证}\BibleRef{宗 4:33}，对他们说：\BibleQuote{我们祖先的天主，举扬了耶稣，就是你们所捉拿、所杀害，悬挂在木架上的；天主用右手高举了他，使他成为元首和救主，为将悔改和赦罪赐给以色列。我们就是这些事的见证人，圣神也是如此，就是天主赐给顺从他的人的。}\BibleRef{宗 5:30}又说：\BibleQuote{每天，在圣殿里，并且挨家挨户，他们从不停止教导和宣讲基督耶稣}\BibleRef{宗 5:42}，即天主子。因为这就是救恩的知识，它使那些承认他圣子来临的人，在天主面前成为成全的。\par

6. 但若这些人中有些无耻地断言，宗徒们在犹太人中间宣讲时，不能向他们宣报另一个神，即除了他们所信仰的神以外的神，我们要对他们说：如果宗徒们过去是按照人们心中固有的意见对人说话，那么就没有人从他们那里学到真理，更早的时候也没有从主那里学到；因为那些人说，主自己也以同样的方式说话。因此，这些人自己也不认识真理；但既然他们对天主的看法就是如此，他们只是照他们所能听的接受了教训。那么，按照这样的说话方式，真理的准则就不能存在于任何人身上；但所有学习者都会把这种做法归给所有的［教师］，即每个人如何思想，按他的能力所及，向他所说的话也是如此。但主的来临将显得多余而无用，如果他真的来是为了容忍和保留各人心中从古以来对天主的看法。此外，还有一件更艰巨的任务，就是那位被犹太人视为人、并被钉在十字架上的，竟被传为基督、天主子、他们永远的君王。既然事实如此，他们确实没有按照他们旧有的信仰对他们说话。因为那些人，当面告诉他们是杀死主的人，他们自己也会更勇敢地宣讲那高于神工者的父，而不是每个人要自己相信的［关于天主的］东西；而且罪过要小得多，如果他们确实没有把更优越的救主（他们本应升到他那里）钉在十字架上，因为他是不能受苦的。因为，他们向外邦人说话时，并没有迎合他们的观念，而是大胆地告诉他们，他们的神不是神，而是魔鬼的偶像；同样，如果他们知道另一位更伟大、更完美的父，他们也会以同样的方式向犹太人宣讲，而不是滋养或加强这些人对天主的错误看法。此外，在摧毁外邦人的错误、把他们从他们的神那里带走的同时，他们当然没有把另一种错误引到他们身上；相反，移除了那些不是神的，他们指出了那唯一的天主和真正的父。\par

7. 因此，从伯多禄在凯撒勒雅对百夫长科尔乃略和与他同在的外邦人（他们首先听到天主的圣言）所说的话，我们可以了解宗徒们当时宣讲的内容、他们宣讲的性质以及他们对天主的观念。因为这位科尔乃略，据记载，是\Quote{一个虔诚的人，他和他全家都敬畏天主，多多周济百姓，常常向天主祈祷。大约在第九时辰，他看见一位天主的天使进来对他说：你的施舍已在天主前成为纪念。所以，你要打发人往西满那里去，请那称呼伯多禄的。}\BibleRef{宗 10:1}但伯多禄在异象中听见从天上传来的声音说：\Quote{天主所洁净的，你不可称为污秽。}\BibleRef{宗 10:15}这事发生，是为教导他：那位曾藉法律区分洁与不洁的天主，正是那藉祂圣子的血洁净了外邦人的那位——也是科尔乃略所敬拜的那位。伯多禄进去后对他说：\Quote{我真正看出天主是不看情面的；在各民族中，凡敬畏祂而行正义的，都为祂所悦纳。}\BibleRef{宗 10:34}他清楚表明：科尔乃略先前所敬畏、并曾透过法律和先知听见过、且因此也施舍的那位，确实是天主。然而，他尚缺乏对圣子的认识；故此，伯多禄补充说：\Quote{你们知道：在若翰宣讲洗礼以后，从加里肋亚开始，传遍整个犹太的这道圣言，就是关于纳匝肋人耶稣，天主怎样以圣神和德能傅了祂；祂巡行各处，施恩行善，治好所有受魔鬼压制的人，因为天主与祂同在。我们就是祂在犹太地和耶路撒冷所行一切事的证人；他们却把祂悬在木架上，杀死了祂。天主在第三日使祂复活，并且显现出来，不是给所有百姓，而是给天主预先拣选的证人，就是我们这些在祂从死者中复活后，与祂同食共饮的人。祂命令我们向百姓宣讲，并证明祂就是天主所立定的，审判生者死者的那一位。众先知都为祂作证：凡信祂的人，赖祂的名字，都获得罪的赦免。}\BibleRef{宗 10:37}因此，宗徒们所宣讲的，是人所不认识的天主圣子，以及祂的来临，对象是已受教认识天主的人；但他们并没有引入另一位神。因为如果伯多禄知道任何这类事，他一定会放胆向外邦人宣讲：犹太人的天主是一位，而基督徒的天主是另一位；而且他们无疑都会因天使显现而畏惧，相信他所说的一切。但从伯多禄的话明显可知，他确实仍持守他们已认识的那位天主；同时他也向他们见证：耶稣基督是天主圣子，是审判生者死者的那一位，并命令他们为得罪的赦免而受洗归于祂；不仅如此，他还见证耶稣本身就是天主圣子，祂曾受圣神傅油，被称为耶稣基督。祂就是由玛利亚所生的那一位，正如伯多禄的见证所暗示的。难道伯多禄那时尚未拥有这些人后来发现的完满知识吗？按照他们的说法，伯多禄是不完美的，其余的宗徒也是不完美的；因此，他们应当复活，成为这些人的门徒，以便自己也得以成全。但这实在可笑。事实上，这些人被证明不是宗徒的门徒，而是他们自己邪恶思想的门徒。这也导致他们中间存在各种意见分歧，因为每个人都按其所能接受了谬误。但遍布全世界的教会，其根源稳固来自宗徒，在天主和祂圣子的事上，持守着同一见解。\par

8. 再者：斐理伯向埃塞俄比亚女王的总管（他正从耶路撒冷返回，诵读依撒意亚先知书）所宣讲的，当他们两人独处时，岂不是那位先知所说：\BibleQuote{祂被牵去宰杀，如同一只绵羊；又像羔羊在剪毛者前缄默，同样祂不开口。}\BibleQuote{谁能述说祂的世代？因为祂的生命从地上被夺去。}\BibleRef{宗 8:32}的那位吗？斐理伯宣告这位就是耶稣，并且圣经在祂身上应验了；那位信了的总管也如此相信：他立即请求受洗，说：\BibleQuote{我信耶稣基督是天主圣子。}\BibleRef{宗 8:37}这人也被派往埃塞俄比亚地区，宣讲他自己所相信的：只有一位天主，是先知们所宣讲的，但这一位天主的圣子已经\OriginalTerm{按人性}{secundum hominem}显现，并被牵去宰杀如同一只绵羊；以及先知们论及祂所说的其他一切话。\par

9. \Person{保禄}本人也同样——主从天上向他说话，指示他：迫害祂的门徒，就是迫害他自己的主，并派遣\Person{阿纳尼雅}到他那里，使他复明并受洗——经上记载说：\Quote{宣讲}，\Quote{耶稣在大马士革的会堂中，放胆讲论，说这一位就是天主子、基督。}\BibleRef{宗 9:20}这就是他说由启示所知晓的奥秘：那位在\Person{般雀·比拉多}手下受难的，就是万有的主、王、天主和审判者，从万有的天主那里接受了权能，因为祂成了\Quote{听命至死，且死在十字架上}\BibleRef{斐 2:8}。既然这是真的，当他在\Place{阿勒约帕哥}向雅典人宣讲时——那里没有犹太人，他得以放胆宣讲天主——他对他们说：\Quote{创造世界及其中的万有的天主，祂既是天地的主宰，不住人手所造的殿宇，也不受人手的服侍，好像需要什么，反而将生命、气息和万有赐给众人。祂从一人血脉造出万族，住在地上，预先定准他们的年限和居住的疆界，要他们寻求神性，或者可以揣摩而得，其实祂离我们各人不远。因为我们生活、行动、存在都在祂内，正如你们自己的诗人说的：\InnerQuote{我们也是祂的子孙}。因此，既然我们是天主的子孙，就不该以为神性像金、银或石头，由人工艺所雕刻。所以，天主忽略世人蒙昧无知的时期，如今却命令各处的人都要悔改；因为祂已定了日子，要藉着祂所立的人按正义审判天下，且叫祂从死人中复活，给万人作可靠的凭证。}\BibleRef{宗 17:24}如今在这段经文中，他不仅向他们宣告天主是世界的创造者（那里没有犹太人），而且祂也造了一族人类住在地上；正如\Person{梅瑟}所宣告的：\Quote{至高者划分万民时，分散\Person{亚当}的子孙，按天主天使的数目，定下了万民的疆界；}但相信天主的民族如今不在天使的权下，而是在主的权下。\Quote{因为\Person{雅各伯}是主的子民，\Place{以色列}是祂产业的份。}\BibleRef{申 32:9}再者，在\Place{吕考尼雅的吕斯特辣}，当\Person{保禄}与\Person{巴尔纳伯}一起，因我们的主\Person{耶稣基督}的名使一个生来瘸腿的人行走，群众因这奇事要尊他们为神时，他对他们说：\Quote{我们也是像你们一样的人，向你们宣讲天主，叫你们离开这些虚假的神，归向永生的天主，祂创造了天、地、海和其中的万物。祂在过去的世代容忍万民各行其道，然而祂并不是没有为自己留下证据，常行善事，从天上降雨，赐给你们丰年，使你们饮食饱足，满心喜乐。}\BibleRef{宗 14:15}但是，他的一切书信与这些宣告是一致的，我将在解释宗徒时，从书信本身在适当的地方表明。然而，当我藉这些证据引出圣经的真理，并扼要地陈明以各种方式陈述的事情时，你们也要耐心地留意，不要认为冗长；要考虑到，包含在圣经里的事物的证据，除非从圣经本身，是无法显明的。\par

10. 更进一步，\Person{斯德望}——由宗徒们选为首位执事，且最先跟随主殉道的足迹，为承认基督而被杀的第一个人——在民众中放胆谈论教训他们说：\Quote{荣耀的天主显现给我们的祖先\Person{亚巴郎}……对他说：你离开你的家乡和你的亲戚，到我指示你的地方去；……于是天主把他迁到你们现在所住的这地方。在这地方，天主没有留给他产业，连立足之地也没有；却应许要将他承受此地为产业，并赐给他的后裔。……天主这样说：他的后裔必在外邦寄居，被奴役虐待四百年；那奴役他们的民族，我要惩罚，主说。以后他们要出来，在这地方事奉我。并赐给他割损的盟约；这样，[亚巴郎]生了\Person{依撒格}。}\BibleRef{宗 7:2}随后的经文也宣告了同一位天主，祂与\Person{若瑟}和先祖们同在，并与\Person{梅瑟}说话。\par

11. 宗徒教义的全部内容宣扬同一天主，祂曾迁徙\Person{亚巴郎}，向他作出继承的许诺，在适当的时候赐给他割损之约，将他的后裔从埃及领出，藉割损（祂赐此作为标记，使他们不与埃及人相似）在外部保守他们；祂是万有的创造者，祂是主耶稣基督的父，祂是荣耀的天主——凡愿意的人可以从宗徒的言行本身学习，并默观这个事实：天主是唯一的，在祂之上没有别的天主。即使有另一个神明在祂之上，我们也会说，在比较（各自工作的）量以后，后者优于前者。因为行为显明更好的人，正如我已经说过的；而且，既然这些人没有他们父的工作可以提出，后者就被证明是唯一的天主。可是，若有任何人\BibleQuote{好辩成性}\BibleRef{弟前 6:4}，认为宗徒所宣告关于天主的事应当寓意化，就请他考虑我之前的陈述，在其中我提出了唯一的天主为万有的创建者和创造者，并摧毁和揭露了他们的主张；他将发现这些陈述与宗徒的教义一致，从而坚持他们过去教导并确信的：有一位天主，万有的创造者。当他将思想从这样的错误及其包含的对天主的亵渎中解脱出来时，他将自行找到理由承认：\Person{梅瑟}法律和新盟约的恩宠，正如两者都适合（赐予的）时代，是由同一天主为了人类的益处而赐下的。\par

12. 因为所有那些存有乖谬心思的人，与\Person{梅瑟}立法对立，认为它与福音教义不同且相反，却没有致力于考察两个盟约差异的原因。因此，他们既被父爱抛弃，又被\Person{撒殚}鼓胀，被引向\Person{西满魔术士}的教义，就在意见上背弃了天主，并想象自己比宗徒发现了更多，找到了另一个神明；并且（主张）宗徒宣讲福音仍然多少受犹太意见的影响，而他们自己在（教义上）更纯正、更聪明，超过宗徒。因此，\Person{马尔西翁}及其追随者开始窜改圣经，根本不承认某些书卷；并且删节\BibleBook{路加}{福音}和\BibleBook{保禄}{书信}，断言只有这些（被他们自行缩短的）才是可靠的。然而，在另一部著作中，若天主赐我（力量），我将从他们仍然保留的这些经书出发驳斥他们。但所有其余的人，膨胀于\Quote{知识}的虚名下，确实承认圣经；但他们曲解解释，正如我在第一卷中所指出的。的确，\Person{马尔西翁}的追随者直接亵渎造物主，声称祂是邪恶的创造者，（但）对其来源持有一种较可容忍的理论，（并）主张有两个本性为神明的存在，彼此不同——一个为善，另一个为恶。然而，\Person{华伦提努}的追随者，虽然使用更尊贵的名称，并提出那位造物主既是父、又是主、又是天主，却（仍然）使他们的理论或派别更加亵渎，因为他们主张祂并非出自\OriginalTerm{圆满}{Pleroma}之内的任何\OriginalTerm{永世}{Æons}，而是出自被驱逐到圆满之外的那个缺陷。对圣经和天主经世之无知给它们带来了这一切。在这部著作的过程中，我将一方面讨论盟约差异的原因，另一方面讨论它们的统一与和谐。\par

13. 但是，宗徒及其门徒既如此按照教会所宣讲的教导，并如此教导得以成全，因此他们也被召叫到达那成全之境——\Person{斯德望}，当他还在世上教导这些真理时，看见了天主的荣耀，并耶稣站在祂右边，就大声说：\BibleQuote{看，我见天开了，并见人子站在天主右边。}\BibleRef{宗 7:56}他说了这些话，就被石头砸死；这样他就完成了成全的教义，在各方面效法殉道元首，并为那些杀害他的人祈祷，说了这些话：\BibleQuote{主，不要将这罪归到他们身上。}这样，那些认识同一天主、从始至终以各种经世与人类同在的人得以成全。正如\Person{欧瑟亚}先知所宣告：\BibleQuote{我已增加异象，并藉先知的手使用比喻。}\BibleRef{欧 12:10}因此，那些为基督的福音将自己的灵魂交付到死的人——他们怎能照旧有的意见向人说话呢？如果他们采取了这种路线，他们就不应受苦；但是，因为他们宣讲了与那些不赞同真理之人相反的事，所以他们受苦。因此，很明显，他们没有放弃真理，而是放胆向犹太人和希腊人宣讲。向犹太人，他们（宣告）曾被他们钉死的耶稣是天主子，审判生者死者的，并已从祂的父那里接受了在以色列的永恒王国，如我所指出的；向希腊人，他们宣讲一位创造万有的天主，和祂的圣子耶稣基督。\par

这从宗徒们的书信中得到更清楚的证明：他们不是写给犹太人，也不是写给希腊人，而是写给那些从外邦人中相信基督的人，以坚定他们的信德。因为有人从犹太下到安提约基雅——在那里，主的门徒因信基督，首先被称为基督徒——并试图说服那些信了主的人受割礼，并遵守法律的其它规条。保禄和巴尔纳伯为这个问题上耶路撒冷去见宗徒，整个教会聚集在一起，伯多禄就这样对他们说：\Quote{诸位仁人弟兄，你们知道，从最初天主在你们中间拣选了我，使外邦人从我口中听到福音的话而相信。知道人心的天主为他们作了证，赐给他们圣神，如同赐给我们一样；在我们和他们之间没有作任何区别，因信德净化了他们的心。既然如此，现在你们为什么试探天主，要把一个我们的祖先和我们自己都负担不了的轭，放在门徒的颈上？但是，我们相信，我们得救是藉着主耶稣的恩宠，和它们一样。} \BibleRef{宗 15:15}此后，雅各伯发言如下：\Quote{诸位仁人弟兄，西满述说了天主当初怎样眷顾外邦人，从他们中选取一个百姓，归于自己的名下。先知的话也与此符合，正如经上所载：\InnerQuote{此后我要回来，重建达味已倒塌的帐幕；我要重建它的废墟，并要把它竖立起来，使剩余的人寻求主，并所有呼求我名的外邦人——行这些事的主说。}\BibleRef{亚 9:11}天主从永远就知道他的工作。因此，我判定，不要难为那归向天主的外邦人；但要写信给他们，叫他们戒避偶像的虚妄、奸淫和血，并他们不愿意别人对自己做的事，也不要对别人做。} \BibleRef{宗 15:14}这些话说完，众人都同意，他们就写信这样说：\Quote{宗徒和长老弟兄们，给在安提约基雅、叙利亚和基里基雅的外邦众弟兄请安。我们听说有人从我们这里出去，用言语搅乱了你们，颠覆你们的灵魂，说你们必须受割礼并遵守法律；我们并没有吩咐过他们。我们同心合意，决定拣选几个人，同我们亲爱的巴尔纳伯和保禄，打发到你们那里去；他们曾为我们主耶稣基督的名，交出了自己的性命。所以我们派了犹达和西拉，他们也要亲口报告同样的事。因为圣神和我们决定，不再将别的重担加在你们身上，唯有这几件事是必要的：就是戒避祭偶像的肉、血、奸淫，并你们不愿意别人对自己做的事，也不要对别人做。你们如果自己谨守，不犯这些，就好了。愿你们在圣神内行走。}由此可见，他们并没有教导另一位父的存在，而是将自由的新盟约赐给那些藉着圣神刚刚相信天主的人。从他们所辩论的问题（即门徒是否仍须受割礼）的性质，可以清楚看出，他们并不知道另一位神。\par

否则，他们也不会对第一个盟约持这样的态度，甚至不愿意与外邦人一起吃饭。因为，即使伯多禄曾被派遣去教训他们，并因一个神视而被迫这样做，但他仍然有些犹豫地对他们说：\Quote{你们知道，一个犹太人是不可以同外邦人来往或接近的；但天主指示给我，不可把任何人看作凡俗或不洁的。所以，我一被邀请，便毫不推辞地来了。} \BibleRef{宗 10:28}这些话表明，若不是奉命，他是不会去的。同样，若不是他听见他们当圣神降临时说先知话，他也不会如此轻易地给他们付洗。因此他喊道：\Quote{谁能禁止用水给他们施洗呢？他们既和我们一样领受了圣神。} \BibleRef{宗 10:47}同时，他说服了那些与他同在的人，并指出，除非圣神降临在他们身上，恐怕会有人反对给他们付洗。而雅各伯同来的宗徒们允许外邦人自由行动，把我们交托给天主圣神。但他们自己，虽然认识同一位天主，却仍遵守古时的规矩；因此，连伯多禄也怕受他们责备，尽管从前因神视和圣神降临而在外邦人中间吃饭，但当有人从雅各伯那里来的时候，就退避，不与他们同食。保禄说，连巴尔纳伯也照样行了。\BibleRef{迦 2:12}因此，主立为一切行动和一切道理的见证的宗徒们——因为我们常常看到伯多禄、雅各伯和若望与主同在——审慎地遵行梅瑟法律的安排，表明这法律出于同一位天主；如果他们从主那里学到另有别的父（我先前已经说过），他们决不会这样做。\par

\SourceRecord{Translated by Alexander Roberts and William Rambaut.；Ante-Nicene Fathers；Vol. 1.；Edited by Alexander Roberts, James Donaldson, and A. Cleveland Coxe.；Buffalo, NY: Christian Literature Publishing Co.,；1885.；<http://www.newadvent.org/fathers/0103312.htm>.}

% structure: p0001 source=h1 target=section reason=page-title

\section{\WorkTitle{驳斥异端}（第三卷，第十三章）}

1. 至于那些（马尔基翁派）声称唯独\Person{保禄}认识真理，并且奥秘是由启示向他显示的人，让\Person{保禄}亲自定他们的罪，因为他说，同一的天主在\Person{伯多禄}身上为受割礼的人施行了宗徒职务，在我身上则为外邦人施行了宗徒职务。\BibleRef{迦 2:8}因此，\Person{伯多禄}正是那一位的宗徒，\Person{保禄}也是属于那一位；\Person{伯多禄}在受割礼的人中宣告为天主的，并同样地宣告为天主子，\Person{保禄}也在外邦人中宣告了。因为我们的主从来不是仅仅为了拯救\Person{保禄}而来，天主也不是如此受限，以致祂只有一位知道祂圣子救恩计划的宗徒。再者，当\Person{保禄}说：\BibleQuote{那传布美好事物、宣讲和平福音的人的脚步，是多么美丽啊！}\BibleRef{罗 10:15}，他清楚地表明，不是仅仅一个人，而是有许多人曾经宣讲真理。再者，在\WorkTitle{格林多前书}中，当他列举了所有在复活后看见过天主的人之后，他接着说道：\BibleQuote{无论是我或是他们，我们如此宣讲，你们也如此相信了。}\BibleRef{格前 15:11}，承认所有在死者复活后看见过天主的人的宣讲是同一的。\par

2. 再者，主答复了愿意见父的\Person{斐理伯}：\BibleQuote{我同你们在一起这么长久，你还不认识我吗？\Person{斐理伯}！谁看见了我，就是看见了父；你怎么说：把父显示给我们呢？因为我在父内，父也在我内；从今以后，你们认识祂，并且已经看见祂了。}因此，主为这些人作证，说他们在祂身上已经认识并看见了圣父（而圣父就是真理）。因此，断言这些人不认识真理，就是扮演假见证的角色，并与基督的教导疏远的人。因为如果这些人不认识真理，主为什么派遣十二宗徒到以色列家迷失的羊那里去呢？\BibleRef{玛 10:6}那七十个人又如何宣讲，除非他们自己预先知道了所宣讲的真理？或者，\Person{伯多禄}怎么可能无知呢？主为他作证，说不是肉和血启示了他，而是在天上的父。\BibleRef{玛 16:17}正如\BibleQuote{\Person{保禄}做宗徒，不是出于人，也不是借着人，而是借着耶稣基督和天主父}\BibleRef{迦 1:1}［其余的也是如此；］圣子确实带领他们到圣父那里，而圣父则向他们启示圣子。\par

3. 保禄同意那些因所提出的问题而召唤他到宗徒那里的人，并与\Person{巴尔纳伯}一同上到耶路撒冷，这并非没有理由，而是为使外邦人的自由由他们得以确认，他本人在\WorkTitle{迦拉达书}中这样说：\BibleQuote{过了十四年，我同\Person{巴尔纳伯}再上耶路撒冷去，也带了\Person{弟铎}同去。但我是因启示上去的，向他们陈述了我在外邦人中所宣讲的福音。}\BibleRef{迦 2:1}再者他说：\BibleQuote{我们暂时顺服了，是为使福音的真理在你们中间得以保存。}\BibleRef{迦 2:5}如果任何人从\WorkTitle{宗徒大事录}仔细考察那记载他因前述问题而上耶路撒冷的时间，他将会发现\Person{保禄}提到的那些年份与之吻合。这样，\Person{保禄}的陈述与\Person{路加}关于宗徒的见证是和谐一致的，可以说是同一的。\par

\SourceRecord{Translated by Alexander Roberts and William Rambaut.；Ante-Nicene Fathers；Vol. 1.；Edited by Alexander Roberts, James Donaldson, and A. Cleveland Coxe.；Buffalo, NY: Christian Literature Publishing Co.,；1885.；<http://www.newadvent.org/fathers/0103313.htm>.}

% structure: p0001 source=h1 target=section reason=page-title

\section{\WorkTitle{反对异端（卷三，第十四章）}}

1. 但这路加与保禄不可分离，且是他福音上的同工，他自己清楚表明，并非出于夸耀，而是因真理本身所迫。因他说，当巴尔纳伯与号称马尔谷的若望离开保禄，航行至塞浦路斯时，\BibleQuote{我们来到特洛阿}；\BibleRef{宗 16:8} 当保禄在梦中看见一个马其顿人，说：\BibleQuote{请来马其顿，保禄，帮助我们}，\BibleQuote{我们立即}，他说，\BibleQuote{设法往马其顿去，明白是主召我们给他们传福音。于是，从特洛阿开船，我们直航撒摩辣刻。} 然后他仔细说明其余全部旅程，直到斐理伯，以及他们如何初次宣讲：\BibleQuote{因为坐下来}，他说，\BibleQuote{我们向聚集的妇女讲道}；\BibleRef{宗 16:13} 有些人信了，甚至很多。他又说：\BibleQuote{过了无酵节，我们从斐理伯开船，五天到了特洛阿，在那里住了七天。} \BibleRef{宗 20:5} 他叙述了与保禄同行的其余一切细节，详尽地指出地点、城市和天数，直到他们上耶路撒冷；以及保禄在那里所遭遇的，\EditorialNote{参阅《宗徒大事录》第21章} 他如何被捆绑解送罗马；押送他的百夫长的名字；\EditorialNote{参阅《宗徒大事录》第27章} 船的标记，以及他们如何遭遇船难；\BibleRef{宗 28:11} 他们脱险的海岛，在那里如何受到善待，保禄治愈了岛上的首领；又如何从那里航行至仆特约利，再由此抵达罗马；以及在罗马居留了多久。因路加亲历这一切，他仔细地以文字记录下来，以致他不能被指为虚假或夸耀，因为所有这些细节都证明他比现今所有另立学说的人更资深，且他并非不认识真理。他不只是宗徒的追随者，更是同工，尤其是保禄的同工，保禄本人也在书信中宣告说：\BibleQuote{德玛斯离弃了我，……往得撒洛尼去了；克勒斯先往迦拉达去；弟铎往达尔玛提亚去。只有路加在我这里。} \BibleRef{弟后 4:10} 由此表明他始终与保禄相连，不可分离。他又在致哥罗森人的书信中说：\BibleQuote{可爱的医生路加问候你们。} \BibleRef{哥 4:14} 但倘若路加常与保禄一同传道，被保禄称为\BibleQuote{可爱的}，并与保禄一同履行福音传道者的工作，受托将福音传给我们，他所学的不异于保禄（正如从他的话语所指出的），那么这些从未依附保禄的人，怎能夸口他们学到了隐藏而不可言传的奥秘呢？\par

2. 但保禄以纯朴教导他所知道的，不仅向与他一同工作的人，也向听他讲道的人，他自己清楚显明了。因为当从厄弗所和其他邻近城市来的主教与长老聚集在米肋托时，他本人正赶往耶路撒冷过五旬节，在向他们作多项证言并宣告他在耶路撒冷必要遭遇的事之后，他接着说：\BibleQuote{我知道你们以后不会再见到我的面了。因此我今日向你们作证：我是纯洁的，与众人的血无关。因为我从来没有避讳而不把天主的全部旨意传报给你们。你们应当留心自己，又留心整个羊群，因为圣神立你们为监督，牧养主的教会，这是祂用自己的血所获得的。} \BibleRef{宗 20:25-28} 然后，他论及将要兴起的恶导师，说：\BibleQuote{我知道在我离去之后，将有凶暴的豺狼进到你们中间，不怜惜羊群。就是在你们中间，也要有人起来讲说谬论，引诱门徒跟从他们。} \BibleRef{宗 20:29-30} \BibleQuote{我从来没有避讳}，他说，\BibleQuote{而不把天主的全部旨意传报给你们。} \BibleRef{宗 20:27} 宗徒们就是这样单纯地、不看情面地将他们从主所学的一切传授给众人。同样，路加也不看情面，将他从他们所学到的传授给我们，正如他自己所证言的：\BibleQuote{正如那些从起初亲眼看见并为圣言的仆役所传授给我们的。} \BibleRef{路 1:2}\par

3. 如今，若有人弃绝路加，认为他不知真理，此人便显然拒绝了他自称是门徒的那部福音。因为借着路加，我们得知了福音中许多重要部分，例如：若翰的出生、匝加利亚的事迹、天使拜访玛利亚、依撒伯尔的呼喊、天使降临牧人处、他们所说的话、亚纳和西默盎关于基督的见证、祂十二岁时留在耶路撒冷、若翰的洗礼、主受洗时的年龄、以及此事发生在提庇留凯撒执政第十五年。在祂为师的职务中，祂对富人说：\BibleQuote{你们富有的人有祸了，因为你们已得了你们的安慰；}\BibleRef{路 6:24}以及\BibleQuote{你们饱足的人有祸了，因为你们将要饥饿；你们现在欢笑的人有祸了，因为你们将要哀恸哭泣；}还有\BibleQuote{众人夸赞你们的时候，你们就有祸了，因为他们的祖先也同样对待过假先知。}以下这类事情，我们唯独从路加得知（且通过他，我们得知了主的许多作为，这些也为所有圣史所提及）：伯多禄的同伴在按主命令撒网时所捕获的大量鱼群；\EditorialNote{路 5}患病十八年的妇人，在安息日获得痊愈；\EditorialNote{路 13}患水臌的人，主在安息日治愈了他，并为自己在安息日行医辩护；祂教导门徒不要追求上座；我们应邀请那些无力回报的贫穷软弱者；半夜敲门求饼的人，因切求而得到；\EditorialNote{路 11}当主与一位法利塞人同席时，一个罪妇亲吻祂的脚，用香膏涂抹，以及主为她对西满所说关于两个债务人的话；\EditorialNote{路 7}还有关于那积蓄财物的富人的比喻，他对那人说：\BibleQuote{今夜就要索回你的灵魂，你所预备的一切将归谁呢？}\BibleRef{路 12:20}类似的还有那身穿紫红袍、奢华宴乐的富人，以及贫穷的拉匝禄；\EditorialNote{路 16}还有祂对门徒说\BibleQuote{请增加我们的信德}时的回答；\BibleRef{路 17:5}祂与税吏匝凯的谈话；\EditorialNote{路 19}以及同时祈祷的法利塞人与税吏；\EditorialNote{路 18}还有在路上同时被祂洁净的十个癞病人；\EditorialNote{路 17}以及祂命令从街巷中召集瘸腿和瞎眼的来参加婚宴；\EditorialNote{路 18}还有那不敬畏天主的法官的比喻，因寡妇的切求而为她伸冤；\EditorialNote{路 13}以及葡萄园中不结果的无花果树。还有许多其他细节只见于路加，玛尔基翁和瓦伦提努都加以利用。此外，还有复活后主在路上对门徒所说的话，以及他们如何在擘饼时认出祂来。\EditorialNote{路 24}\par

4. 因此，这些人必然要么接受路加其余记述，要么也拒绝这些部分。因为凡有常识者，都不会容许他们把路加所记某些事当作真实接受，而把另一些事搁置一边，仿佛他不知真理似的。如果玛尔基翁的门徒拒绝这些，他们将没有福音；因为如我所说，他们删节了路加福音，却还夸耀拥有福音。而瓦伦提努的门徒则必须放弃他们那些全然虚妄的言论，因为他们从这部福音中获取了许多借以发展自己臆测的机会，对路加所言本善的事作了恶意的解释。另一方面，如果他们不得不接受其余部分，那么通过研究完整的福音和宗徒的教导，他们将发现必须悔改，才能从面临的危险中得救。\par

\SourceRecord{Translated by Alexander Roberts and William Rambaut.；Ante-Nicene Fathers；Vol. 1.；Edited by Alexander Roberts, James Donaldson, and A. Cleveland Coxe.；Buffalo, NY: Christian Literature Publishing Co.,；1885.；<http://www.newadvent.org/fathers/0103314.htm>.}

% structure: p0001 source=h1 target=section reason=page-title

\section{\WorkTitle{驳斥异端（第三卷，第十五章）}}

1. 然而，我们同样指控那些不承认\Person{保禄}为宗徒的人：他们要么拒绝接受我们仅仅通过\Person{路加}得知的福音的其他言论，并且不使用它们；要么，如果他们接受所有这些，他们必须也承认那关于\Person{保禄}的见证，即\Person{路加}告诉我们的：主起初从天上对他说：\Quote{扫禄，扫禄，你为什么迫害我？我是耶稣基督，你所迫害的。} \BibleRef{宗 22:8} 然后又对\Person{阿纳尼雅}说，关于他：\Quote{你去吧！因为他是我所拣选的器皿，为把我的名字带到外邦人、君王和以色列子民面前。我要指示他，从今以后，他必须为我的名字受多少苦。} \BibleRef{宗 9:15} 因此，那些不接受他（作为导师）的人，他是天主为这目的所拣选的，要勇敢地奉他的名，被派到前述的民族那里，他们轻视天主的拣选，并自绝于宗徒的团体。因为他们既不能争辩说\Person{保禄}不是宗徒，因他为此被拣选；也不能证明\Person{路加}说谎，因为他殷勤地向我们宣讲真理。的确，也许正是为此，天主通过\Person{路加}的工具设定了许多福音真理，所有人都应认为有必要使用它们，以便所有人，追随他后来关于宗徒的行为和教义的见证，持守不掺杂的真理准则，可以得救。因此，他的见证是真实的，宗徒的教义是公开而坚定的，毫无保留；他们也没有私下教导一套教义，而在公开场合教导另一套。\par

2. 因为这是假人、邪恶的诱惑者和伪君子的托词，正如那些来自\OriginalTerm{瓦伦廷}{Valentinus}的人所做的。这些人向群众谈论那些属于教会的人，他们自己称他们为\Quote{粗俗的}和\Quote{教会的}。用这些话，他们诱捕较单纯的人，并引诱他们，模仿我们的措辞，以便这些（受骗者）更频繁地听他们；然后这些人被问及关于我们，为什么他们持有与我们相似的教义，我们却毫无理由地远离他们的团体；以及当他们说同样的话，持有同样的教义时，我们称他们为异端？当他们通过问题推翻任何人的信仰，并使他们成为自己无反驳的听众后，他们私下向他们描述他们\OriginalTerm{圆满}{Pleroma}的不可言喻的奥秘。但那些想象他们可以从异端引用的圣经文本中学到那些他们的言辞貌似教导的（教义）的人，完全是被欺骗了。因为错误是似是而非的，与真理相似，但需要伪装；而真理没有伪装，因此被托付给小孩子。如果他们的听众中有人确实要求解释，或对他们提出异议，他们断言他是不配接受真理的人，并且没有从上面得到来自他们母亲的种子；因此实际上不给他回答，而只是宣称他属于中间领域，即属于动物本性。但如果有人像一只小羊一样把自己交托给他们，并跟随他们的做法和他们的\Quote{\OriginalTerm{救赎}{redemption}}，这样的人就会膨胀到认为自己既不在天上也不在地上，而是已经进入了圆满之内；并且已经拥抱了他的天使，他蹒跚而行，面带傲慢，拥有一只公鸡的所有浮夸气派。他们中间有些人断言，那从上而来的人应当遵循良好的行为方式；因此他们也假装一种严肃（的态度），带着某种傲慢。然而，大多数人成了嘲弄者，仿佛已经完美，生活毫不关心（外表），甚至蔑视（良善），称自己为\Quote{属灵的人}，并声称他们已经熟悉了他们的圆满之内的安息之地。\par

3. 但让我们回到（迄今为止所遵循的）同一条论证线索。因为当已经明确宣告，那些真理的宣讲者和自由的宗徒，除了唯一的真天主圣父和他的圣言——他在万有中居首位——之外，没有称任何别者为天主，或称他为主；那么将清楚地证明，他们（宗徒）承认那创造天地、也曾与\Person{梅瑟}说话、赐予他法律的安排、并召叫先祖的那一位为主天主；并且他们不认识任何其他者。因此，宗徒和那些从他们的话学习的人（\Person{马尔谷}和\Person{路加}）关于天主的意见已经显明了。\par

\SourceRecord{Translated by Alexander Roberts and William Rambaut.；Ante-Nicene Fathers；Vol. 1.；Edited by Alexander Roberts, James Donaldson, and A. Cleveland Coxe.；Buffalo, NY: Christian Literature Publishing Co.,；1885.；<http://www.newadvent.org/fathers/0103315.htm>.}

% structure: p0001 source=h1 target=section reason=page-title

\section{驳斥异端（第三卷·第十六章）}

1. 然而有些人说，耶稣只是基督的容器，基督如鸽子般从高处降下在他身上，当基督宣告了那不可名状的圣父之后，就以不可理解且无形的方式进入了\OriginalTerm{圆满}{Pleroma}：因为他不被理解，不仅不被人类理解，甚至不被天上那些能力和德能理解；而耶稣是圣子，但基督是圣父，并且基督的父是天主；另一些人则说他只是在表面上受苦，本性是不受苦的。而瓦伦提尼安派则坚持认为，经世的耶稣就是那经过玛利亚的同一个，那个来自更高领域的救主降在他身上，这救主也被称为\OriginalTerm{万有}{Pan}，因为他拥有所有产生他者的\OriginalTerm{名称}{vocabula}；但这位救主与经世的耶稣分享了他的能力和名字，以便借此消灭死亡，而圣父则通过那从高处降下的救主被认识，他们也称这救主本身是基督和整个圆满的容器；他们确实在口头上承认一位基督耶稣，但在实际意见上却是分裂的：因为如前所述，这些人惯于说有一位基督，是由\OriginalTerm{独生子}{Monogenes}产生的，为了圆满的坚固；但另一位，即救主，被派遣为光荣圣父；还有第三位，即经世的那位，他们声称他受了苦，他也怀有基督，也就是那返回圆满的救主。因此我认为有必要考虑宗徒们关于我主耶稣基督的全部心意，表明他们不仅从未持有任何此类意见，而且更进一步，他们通过圣神宣布，传讲这类教义的人是撒殚的差役，被派来倾覆某些人的信德，并将他们引离生命。\par

2. 关于若望知道同一个且唯一的圣言，祂是独生子，并且祂为了我们的救恩降生成人，即耶稣基督我们的主，我已从若望自己的话中充分证明。玛窦也同样承认同一个且唯一的耶稣基督，展示祂作为人由童贞女而生，如同天主曾应许达味，要从他身体的果实中兴起一位永恒之王，并早先就向亚巴郎作出了同样的应许，他说：\Quote{亚巴郎之子达味之子耶稣基督的族谱。}\BibleRef{玛 1:1} 然后，为使我们心中对若瑟的疑虑消除，他说：\Quote{基督的诞生是这样的：祂的母亲玛利亚已许配给若瑟，但在他们同居之前，她因圣神有了身孕。} 然后，当若瑟考虑休退玛利亚，因为她已怀孕时，［玛窦告诉我们］天主的使者站在他身旁，说：\Quote{不要害怕娶你的妻子玛利亚，因为那在她内受孕的是出于圣神。她要生一个儿子，你要给祂起名叫耶稣，因为祂要把自己的民族从他们的罪恶中拯救出来。这一切发生，是为应验主借先知所说的话：看，一位贞女将怀孕生子，人要称祂的名字为\OriginalTerm{厄玛奴耳}{Emmanuel}，意思是：天主与我们同在。} 这清楚地表明，对祖先所作的应许已经实现：天主子由童贞女所生，并且祂自己就是先知们所预言的基督救主；而不是像这些人所断言的那样，耶稣是由玛利亚所生的那一位，而基督是从高处降下的那一位。玛窦本来可以说：\Quote{现在\Emph{耶稣}的诞生是这样的；} 但圣神预见了［真理的］败坏者，并预先防范他们的欺骗，便通过玛窦说：\Quote{但\Emph{基督}的诞生是这样的；} 又表明祂是\OriginalTerm{厄玛奴耳}{Emmanuel}，免得我们或许把祂看作一个普通人：因为\Quote{不是从血肉的意愿，也不是从人的意愿，而是从天主的意愿，圣言成了血肉；} 叫我们不要想象耶稣是一位，基督是另一位，而应知道他们是同一位。\par

3. \Person{保禄}在致罗马人书时，已经解释了这一点：\Quote{\Person{保禄}，耶稣基督的宗徒，预定为天主的福音，这福音是天主先前藉祂的众先知在圣经中所预许的，论及祂的圣子，祂按肉身是达味的后裔，按圣德之神，藉着从死者中复活，被预定为具有大能的天主子，我们的主耶稣基督。}\BibleRef{罗 1:1} 他又论及以色列而致罗马人书，说：\Quote{圣祖们是他们的，并且按肉身说，基督也是从他们来的，祂是在万有之上的天主，永受赞美。}\BibleRef{罗 9:5} 又在他的迦拉达书里说：\Quote{但时期一满，天主就派遣了祂的圣子，由女人所生，生于法律之下，为把在法律之下的人赎出来，使我们获得义子的名分；}\BibleRef{迦 4:4} 这明确显示了一位天主，祂曾藉先知预许了圣子，与一位主耶稣基督，祂按由玛利亚的出生，是达味的后裔；并且耶稣基督按圣德之神，藉着从死者中复活，被立为具有大能的天主子，作为一切受造物中的首生者；\BibleRef{哥 1:14} 圣子成为了人子，好使我们藉着祂获得义子的名分——人性承载、接纳并拥抱了天主子。因此\Person{马尔谷}也说：\Quote{天主圣子耶稣基督福音的开始，正如先知书上所记载的。}\BibleRef{谷 1:1} 我们认识同一天主子耶稣基督，祂由先知所预报，从达味的后裔而来，是\OriginalTerm{厄玛奴耳}{Emmanuel}，\Quote{父的伟大谋士的使者}；藉着祂，天主使晨曦和义人从达味家兴起，并为达味兴起了一个拯救的角，\Quote{并在雅各伯中建立了见证}；\BibleRef{路 1:69} 正如\Person{达味}在论及祂的出生的缘由时所说：\Quote{祂在以色列中颁定了法律，使另一世代知道，这些将要出生的子孙，他们起来，要对自己的子女宣讲，使他们寄望于天主，并寻求祂的诫命。} 此外，天使在向\Person{玛利亚}报喜时说：\Quote{祂将伟大，被称为至高者的圣子；上主天主要把祂祖先达味的御座赐给祂；}\BibleRef{路 1:32} 承认那至高者的圣子，祂自己也是达味之子。\Person{达味}藉着圣神知道这一位来临的安排，藉此祂超越一切生者与死者，便承认祂为主，坐在至高圣父的右边。\par

4. 然而\Person{西默盎}——他曾从圣神得到启示，知道自己未见到基督耶稣以前不会看见死亡——将那位童贞女的首生者抱在手中，赞美天主，并说：\Quote{主啊！现在可照祢的话，让祢的仆人平安离去；因为我的眼睛已看见了祢的救恩，即祢在万民面前所预备的：为作外邦人的光明，和祢子民以色列的荣耀。}\BibleRef{路 2:29} 这样承认，他手中所抱的婴孩耶稣，由\Person{玛利亚}所生，正是基督自己，天主子，万民的光明，以色列的荣耀，以及安眠者的平安与安慰。因为祂已经剥夺了人，消除他们的无知，赐予他们自己的知识，并驱散那些认识祂的人，正如\Person{依撒意亚}所说：\Quote{称祂的名字为快取掠物，急分赃物。}\BibleRef{依 8:3} 这些便是基督的作为。所以，祂本身就是基督，\Person{西默盎}抱着［在怀中］祝福了至高者；牧童看见祂就光荣了天主；\Person{若翰}还在母胎中，而祂（基督）在\Person{玛利亚}的胎中，认出是主，就跳跃致礼；贤士们看见祂，便朝拜了祂，并献上礼物［给祂］，正如我已经说过的，他们俯伏在地朝拜了永世的君王，然后从另一条路离去，不再从亚述的路回去。\Quote{因为在这孩子还未晓得叫父亲或母亲以前，祂必得着大马士革的能力和撒玛黎雅的财物，对抗亚述王。}\BibleRef{依 8:4} 这确实以奥秘的方式，但有力地宣告：主曾以隐藏的手与阿玛肋克战斗。为此，祂也突然除去了\Person{达味}家那些有福此时出生的孩子，使他们先进入祂的王国；祂自己既是婴孩，便这样安排，使人类的婴孩成为殉道者，按圣经所说，为基督而被杀，基督生于犹大的白冷，即\Person{达味}城。\BibleRef{玛 2:16}\par

5. 因此，主在复活后也对祂的门徒说：\Quote{无知的人哪！对先知们所说的一切话，你们的心信得太迟钝了。基督不是必须受这些苦，然后进入祂的光荣吗？} \BibleRef{路 24:25} 祂又对他们说：\Quote{这就是我从前与你们同在时所告诉你们的话：凡梅瑟法律、先知和圣咏上所记载关于我的事，都必须应验。} 于是祂开启了他们的明悟，使他们理解圣经，并对他们说：\Quote{经上这样记载，基督必须受苦，第三天从死者中复活；并且必须从耶路撒冷开始，因祂的名向万邦宣讲悔改以得罪的赦免。} \BibleRef{路 24:44} 这就是那位由玛利亚所生的；因为祂说：\Quote{人子必须受许多苦，被弃绝，被钉在十字架上，并在第三天复活。} 因此，福音所认识的人子，除了那位出自玛利亚、也确实受苦的以外，没有别的；也没有在受苦之前就离开耶稣的基督；福音只知道那被生的是天主之子耶稣基督，而这位同一的主受苦并复活了，正如主的门徒若望所证实的，说：\Quote{但这些事是记载的，为叫你们相信耶稣是基督，天主子；并使你们信的人，因祂的名获得永生。} \BibleRef{若 20:31} ——这是预见那些亵渎性的体系，它们尽其所能地分裂主，说祂是由两种不同本体形成的。为此，若望也在他的书信中这样为我们作证：\Quote{孩子们，现在是末时了；你们曾听说过反基督要来，如今已经有许多反基督出现了；由此我们晓得现在是末时了。他们从我们中间出去，却不属于我们；因为他们若属于我们，就必会留在我们中间；但他们出去了，为要显明他们都不是属于我们的。所以你们知道，一切谎言都是从外面来的，不是出于真理。谁是说谎者呢？岂不是那否认耶稣是基督的吗？这就是反基督。}\par

6. 但是，所有上述那些人，虽然口里确实承认一位耶稣基督，却愚弄自己，想的是一套，说的是另一套；因为他们的假设各有不同，正如我已经表明的，他们声称：一位存在受苦并被生，这就是耶稣；但另一位降临在祂身上，这就是基督，祂也升上去了；他们还争辩说，出自\OriginalTerm{德谬哥}{Demiurge}的，或出自经世安排的，或出自若瑟的那一位，是那受苦的存在；但那位从不可见且不可言喻之处降临的，是前者，他们断言祂是不可理解、不可见、不受苦的。这样他们就偏离了真理，因为他们的教义离开了那位真实的天主，他们不知道天主的独生子圣言，祂始终与人类同在，按圣父的旨意与祂自己的受造物结合并相混，并成了血肉，祂就是耶稣基督我们的主，祂为我们受苦，并为我们复活，祂还要在圣父的光荣中再来，复活所有血肉，并彰显救恩，对祂所造的一切施行公正审判的准则。因此，正如我已经指出的，只有一位天主圣父，和一位基督耶稣，祂藉着整个经世安排而来，并在祂内总括万物。\BibleRef{弗 1:10} 而且，在各方面，祂也是人，是天主的杰作；这样，祂就将人收纳于祂内：不可见的成了可见的，不可理解的成为可理解的，不受苦的成为能受苦的，圣言成了人，在祂内总括一切：以致正如在超天、属神、不可见的事物中，天主的圣言是至高的，同样在可见和形体的事物中，祂也拥有至高权威，并为自己取得首位，立自己为教会的头，在适当的时候将万物吸引归向祂。\par

7. 在祂内没有不完整或不合时宜的事，正如在圣父内没有不协调的事。因为这一切都是圣父预知的；但圣子在适当的时候，以完美的秩序和次序予以成就。这就是为什么当\Person{玛利亚}催促祂行那关于酒的奇妙奇迹，并在时间未到之前就想饮用那具有象征意义的杯时，主阻止她不及时的催促，说：\Quote{女人，这于我和你有什么关系？我的时辰尚未来到。} \BibleRef{若 2:4} ——祂等待着圣父预知的那个时辰。也是因此，当人们屡次想捉拿祂时，经上说：\Quote{没有人下手捉拿祂，因为祂的时辰还没有到；} \BibleRef{若 7:30} 也没有到祂受难的时刻，那是圣父预知的；正如先知\Person{哈巴谷}也说：\Quote{藉此你将被人认识，当岁月临近时；当时候来到，你将被展示；因为我的灵魂因愤怒而不安，你将记起你的仁慈。} \BibleRef{哈 3:2} \Person{保禄}也说：\Quote{但时期一满，天主就派遣了自己的儿子。} \BibleRef{迦 4:4} 由此可见，圣父所预知的一切，我们的主都按其秩序、季节和时辰予以成就，这些是预知和适宜的，祂确实是同一而丰富伟大的。因为祂履行了祂圣父那慷慨而全面的旨意，祂本身就是得救者的救主，治下者的主，以及一切受造之物的天主，圣父的独生子，那被预报的基督，天主的圣言，祂在时期一满时降生成人，那时天主之子必须成为人子。\par

8. 凡借知识之名，以耶稣为一、基督为另一、独生子为另一（由此又有圣言）、救主为另一——这些错谬之徒宣称救主是由那些堕落为永世者所产生——所有这些人都置身于基督信仰的救恩计划之外。他们外表是羊，因为他们公开说的话看起来像我们，重复我们所说的话；但内里却是狼。他们的教义是杀人的，因为它捏造了众多神祇，模拟许多父亲，却以多种方式贬低并分割天主子。主早已预先警告我们要提防这些人；祂的宗徒在他已经提到的书信中命令我们要避开他们，说：\BibleQuote{因为有许多迷惑者进入了世界，他们不承认耶稣基督是成了肉身而来的。这就是迷惑者，就是假基督。你们要提防他们，免得丧失你们所劳苦的。}他在同一书信中又说：\BibleQuote{有许多假先知已到世界上来了。凭此你们可以认出天主的神：凡神承认耶稣基督是成了肉身而来的，就是出于天主；凡神不承认耶稣基督的，就不是出于天主，而是出于假基督。}这些话与福音中所说的\BibleQuote{圣言成了血肉，居住在我们中间}是一致的。因此他又在他的书信中宣称：\BibleQuote{凡信耶稣是基督的，都是由天主而生的}\BibleRef{若一 5:1}，因为他知道耶稣基督是同一个，天国之门因祂取了肉身而为其敞开；祂也要以祂受苦时所同有的肉身再来，显示圣父的光荣。\par

9. 保禄与这些陈述一致，对罗马人宣告：\BibleQuote{何况那些领受丰富恩宠和正义以得永生的人，更要藉着一个人，耶稣基督，为王而活。}\BibleRef{罗 5:17}由此可见，他并不知道那个从耶稣那里飞走的基督；也不知道他们以为不受苦的那位至高救主。因为，如果实际上一个受苦，另一个保持不受苦；一个生，另一个降到已生者身上又离开他，那么所显出的就不是一个，而是两个。但宗徒确实知道祂是一位，既是生又是受苦的那一位，即基督耶稣。他在同一书信中又说：\BibleQuote{难道你们不知道，我们受洗归于基督耶稣的人，就是受洗归于祂的死吗？正如基督从死者中复活，我们也应该在新生活中行走。}\BibleRef{罗 6:3}但为表明基督确实受苦，并且祂就是天主子，为我们而死，在预定之时用自己的血救赎了我们，他又说：\BibleQuote{当我们还在软弱的时候，基督就在预定的时候为不虔敬的人死了。但天主以自己的爱证明给我们：当我们还是罪人的时候，基督为我们死了。那么，现在我们既然因祂的血成义，更要藉着祂从义怒中得救。因为假如我们还在为仇敌的时候，曾藉着祂圣子的死与天主和好，那么，既然和好了，岂不更要藉着祂的生命得救吗？}祂以最清楚的方式宣告，那个被捉拿、受苦、为我们流血的同一位，既是基督又是天主子，祂也确实复活了，被提升天，正如他自己[保禄]所说：\BibleQuote{但不仅如此，基督死了，而且复活了，现今在天主右边}\BibleRef{罗 8:34}又说：\BibleQuote{我们知道，基督既从死者中复活，就不再死}\BibleRef{罗 6:9}因为他自己藉着圣神预见了恶教师们在[主的位格]上的分裂，渴望为他们除掉一切挑剔的借口，便说了已提及的话，[还宣告：]\BibleQuote{如果那使耶稣从死者中复活者的圣神住在你们内，那么，那使基督从死者中复活的，也必要藉着那住在你们内的圣神，使你们有死的身体复活。}\BibleRef{罗 8:11}他不仅对愿意听的人说，也对所有人说：不要错误：耶稣基督，天主子，是同一的，祂藉着受苦使我们与天主和好，并从死者中复活；祂坐在圣父右边，在一切事上完美；\BibleQuote{祂受辱骂不还口，受苦时不恐吓}\BibleRef{伯前 2:23}当祂遭受暴政时，祂祈求圣父赦免那些钉死祂的人。因为祂自己确实带来了救恩：祂本身就是天主的圣言，本身就是圣父的独生子，我们的主基督耶稣。\par

\SourceRecord{Translated by Alexander Roberts and William Rambaut.；Ante-Nicene Fathers；Vol. 1.；Edited by Alexander Roberts, James Donaldson, and A. Cleveland Coxe.；Buffalo, NY: Christian Literature Publishing Co.,；1885.；<http://www.newadvent.org/fathers/0103316.htm>.}

% structure: p0001 source=h1 target=section reason=page-title

\section{驳斥异端（第三卷，第十七章）}

宗徒们当然能够宣告基督降临于耶稣，或所谓的更高救主降临于那经世者，或来自不可见之处的那一位降临于来自\OriginalTerm{得穆革}{Demiurge}的那一位；但他们既不知道，也没有说过这类事：因为，倘若他们知道，他们必定也会陈述出来。但他们所记录的，正是实际情形，即：天主的圣神仿佛鸽子降临在祂身上；这圣神，依撒意亚曾论及说：\BibleQuote{天主的圣神将安息在祂身上，}\BibleRef{依 11:2}正如我已说过的。又说：\BibleQuote{上主的神临于我，因为祂给我傅了油。}\BibleRef{依 61:1}那就是主所宣告的圣神：\BibleQuote{因为说话的不是你们，而是你们父的圣神在你们内说话。}\BibleRef{玛 10:20}再者，主赐给门徒们归于天主的\OriginalTerm{再生}{regeneration}权能，对他们说：\BibleQuote{你们要去使万民成为门徒，因父、及子、及圣神之名给他们授洗。}\BibleRef{玛 28:19}因为天主曾应许，在末期祂要将圣神倾注在仆人和使女身上，使他们说预言；因此圣神也降临在天主子身上，祂成了人子，习惯于与祂共融而居住在人类中，安息在人间，并住在上主的工程中，在他们内承行父的旨意，并使他们从旧习更新到基督的新生。\par

这圣神是达味为人类所祈求的，说：\BibleQuote{并以祢统御万有的圣神坚固我；}这圣神又照路加所记载，在主升天后五旬节日降临在门徒身上，具有使万民进入生命之门并开启新盟约的权能；因此，他们同声合意用各种语言颂扬天主，圣神使远方的部落合而为一，并将万民的首熟之果献给圣父。因此主也应许要派遣\OriginalTerm{护慰者}{Comforter}，\BibleRef{若 16:7}祂要使我们与天主结合。因为正如一块紧实的面团，没有液体就干麦子就无法成形，面包也无法拥有合一；同样，我们这许多人，若没有来自天上的水，也不能在基督耶稣内成为一体。又正如干土若不得到水分就不能生发；同样，我们原本是一棵枯树，若无那从上面自愿降下的雨水，也绝不能结出生命的果实。因为我们的身体藉着那引入\OriginalTerm{不朽}{incorruption}的洗礼而彼此合一；而我们的灵魂则藉着圣神合一。因此两者都是必要的，因为两者都有分子于天主的生命。我们的主怜悯那迷途的撒玛黎雅妇人——她没有守住一个丈夫，而是嫁过多次而犯奸淫——便指出并应许活水给她，使她不再口渴，也不用再劳累获取可提神的泉水，因她内里有涌到永生的水泉。主从父接受了这份恩赐，祂也亲自将它赐予那些分享祂的人，将圣神派遣到普世。\par

基德红，\BibleRef{民 6:37}那个天主拣选来拯救以色列子民脱离外邦人势力的以色列人，预见这恩赐，便改变他的请求，预言羊毛（人民的预像）上会有干涸，起初只有羊毛上有露水；从而指明他们将不再从天主那里拥有圣神，正如依撒意亚所说：\BibleQuote{我也要命云彩不再雨在它上面，}\BibleRef{依 5:6}但那露水，即天主的神，曾降临在主身上，却要遍及全地，就是\BibleQuote{智慧和聪敏的神，超见和刚毅的神，明达和孝爱之神，敬畏天主的神。}\BibleRef{依 11:2}这圣神，祂再次赐予了教会，从天上派遣\OriginalTerm{护慰者}{Comforter}到普世；正因如此，主告诉我们，魔鬼像闪电一样被摔下来。\BibleRef{路 10:18}因此我们需要天主的露水，免得被火焚烧，也不致不结果实，并且我们有了控告者之处，也有了一位\OriginalTerm{辩护者}{Advocate}，\BibleRef{若壹 2:1}主将祂自己的人交给了圣神，这人曾落在强盗手中，\BibleRef{路 10:35}主自己怜悯了他，包扎了他的伤口，给了两个王家的\OriginalTerm{罗马钱币}{denaria}。如此，我们藉着圣神领受了父和圣子的肖像与名号，就能使委托给我们的钱币结果实，将盈余算给主。\BibleRef{玛 25:14}\par

4. 因此，圣神按预定的时期降下，天主子、独生子，也是圣父的圣言，在时期一满时来临，为了人而成为人，取了人性，并满足了人性的一切条件。我们的主耶稣基督是同一而不可分的，正如主亲自作证，宗徒们承认，先知们所宣告的——所有这些人的教义，他们发明了所谓的\OriginalTerm{八神体}{Ogdoads}和\OriginalTerm{四神体}{Tetrads}，并想象出（对主的位格）的划分，已被证明是谬误。事实上，这些人完全摒弃了圣神；他们认为基督是一位，耶稣是另一位；并教导说不是一位基督，而是有许多位。即使他们说二者合一，他们却又将其分开：因为他们指出，一位确实受苦，而另一位则保持\OriginalTerm{不动情}{impassible}；一位确实升到\OriginalTerm{圆满}{Pleroma}之中，而另一位则留在中间之地；一位确实在不可见、无可名状之处宴乐，而另一位则与\OriginalTerm{德谬哥}{Demiurge}同坐，使他失去能力。因此，你们和所有留意这篇著作、关心自身救恩的人，有责任不要轻易认同他们在外所散布的言论。因为，如我已指出的，他们所说的听起来与信徒的教义相似，但他们所持的见解不仅不同，而且完全相反，处处充满亵渎，借以毁灭那些因言语相似而吸收不适合他们体质之毒药的人。正如有人将石灰与水调和当作牛奶，用颜色的相似来误导人。一位比我高明的人曾就所有以任何方式败坏天主之事、掺杂真理的人说过：\Quote{罪恶地将石灰掺入天主的奶中。}\par

\SourceRecord{Translated by Alexander Roberts and William Rambaut.；Ante-Nicene Fathers；Vol. 1.；Edited by Alexander Roberts, James Donaldson, and A. Cleveland Coxe.；Buffalo, NY: Christian Literature Publishing Co.,；1885.；<http://www.newadvent.org/fathers/0103317.htm>.}

% structure: p0001 source=h1 target=section reason=page-title

\section{\WorkTitle{驳斥异端}（卷三、第十八章）}

1. 既然已经清楚证明：那在起初与天主同在的圣言，万有都是藉祂而造成的，祂也一直与人类同在，在这末世，按照圣父所定的时间，与祂自己的工程结合，成为可受苦的人，那么，所有那些说\Quote{如果我们的主在那时出生，基督就因此没有先存}之人的异议便被推翻了。因为我已经表明，天主子并非在那时才开始存在，祂从起初就与圣父同在；但当祂降生成人、成为人时，祂开启了人类的长久谱系，并以简要而包罗万有的方式赐予我们救恩；好使我们在亚当内所失去的——即按照天主的肖像和模样——得以在基督耶稣内恢复。\par

2. 因为那曾一次被征服、因悖逆而毁灭的人，不可能改造自己并获得胜利的奖赏；那已落在罪恶权势之下的人，也不可能获得救恩——圣子成就了这两件事，祂是天主的圣言，从天父降下，降生成人，自谦至死，并完成为我们预定的救恩计划。关于祂，保禄劝我们要毫不犹疑地相信，又说：\BibleQuote{谁能升到天上去？那就是领基督下来；谁能下到深渊里去？那就是再次把基督从死里领上来。}\BibleRef{罗 10:6}接着他继续说：\BibleQuote{你若口里承认主耶稣，心里相信天主使祂从死者中复活了，你必得救。}\BibleRef{罗 10:9}他又说明天主子行这些事的原因，说：\BibleQuote{因为基督为此而生，又死而复活了，为作死人和活人的主。}\BibleRef{罗 14:9}又写信给格林多人时，他宣布：\BibleQuote{我们却是宣扬被钉在十字架上的基督耶稣；}\BibleRef{格前 1:23}并补充说：\BibleQuote{我们所祝福的那祝福之杯，岂不是共结合于基督的血吗？}\BibleRef{格前 10:16}\par

3. 但谁曾在饮食上与我们相通呢？是那些人所认为的在上面的基督，祂透过\OriginalTerm{界限}{Horos}延伸自己，并赋予他们的母亲一个形体；还是生于童贞女的那位厄玛奴耳，祂吃了奶油和蜂蜜，\BibleRef{依 8:14}先知曾论及祂说：\BibleQuote{祂也是人，谁能认识祂？}\BibleRef{耶 17:9}保禄也同样宣讲祂：\BibleQuote{因为我首先传给你们的是：基督照经上记载的，为我们的罪死了，被埋葬了，且照经上记载的，第三天复活了。}\BibleRef{格前 15:3}那么，显然保禄不认识任何别的基督，只认识那一位受苦、埋葬、又复活、也诞生的基督，并且他称祂为人。因为他说了\BibleQuote{如果基督被传报了，祂从死者中复活了}\BibleRef{格前 15:12}之后，接着说明祂降生成人的原因：\BibleQuote{因为死是由一人而来，死者的复活也是由一人而来。}每当论及我们主的苦难、祂的人性以及祂屈服于死亡时，他处处使用基督的名字，例如：\BibleQuote{你的食物不可使那基督为他而死的人丧亡。}\BibleRef{罗 14:15}又说：\BibleQuote{但是现在在基督内，你们从前远离的人，靠着基督的血成为亲近的了。}\BibleRef{弗 2:13}又说：\BibleQuote{基督赎了我们脱离法律的诅咒，为我们成了可咒骂的，因为经上记载：凡挂在树上的，是可咒骂的。}\BibleRef{迦 3:13}又说：\BibleQuote{因你的知识，那软弱的弟兄，基督为他而死的，就要丧亡。}\BibleRef{格前 8:11}这表明那不能受苦的基督并没有降临到耶稣身上，而是祂自己，因为祂是耶稣基督，为我们受苦；祂躺过坟墓，又复活了，祂下降又上升——天主子成了人子，正如这名本身所宣告的。因为在基督这名内，隐含了那傅油的、被傅油的和那傅油本身，即是祂所受的傅油。圣父是那傅油的，圣子是被圣神傅油的，圣神就是那傅油，正如圣言藉依撒意亚所宣告的：\BibleQuote{上主的神临到我身上，因为祂给我傅了油}\BibleRef{依 61:1}——指出了傅油的圣父、被傅油的圣子和那傅油，即圣神。\par

4. 主自己也显明了那受苦的是谁；因为当祂问门徒：\BibleQuote{人们说人子是谁？}\BibleRef{玛 16:13}伯多禄回答：\BibleQuote{祢是基督，永生天主之子。}当祂称赞他（说）：\BibleQuote{不是血肉启示了你，而是我在天之父}时，祂就清楚表明：祂，人子，就是基督，永生天主之子。因为\BibleQuote{从那时起，据说，祂开始向门徒显示，祂必须上耶路撒冷去，受司祭长许多的苦，并被弃绝，被钉在十字架上，并在第三日复活。}\BibleRef{玛 16:21}那位被伯多禄承认是基督的，称他为有福的，因为父将永生天主子启示给了他；祂说祂自己必须受许多苦，并被钉十字架；然后祂责斥伯多禄，因伯多禄以为祂是像一般人所以为的基督，并且不愿祂受苦；祂对门徒说：\BibleQuote{谁若愿意跟随我，该弃绝自己，背起自己的十字架来跟随我。因为谁若愿意救自己的生命，必丧失生命；但谁若为我的缘故丧失自己的生命，必获得生命。}\BibleRef{玛 16:24}因为基督公开说了这些事，祂自己正是那些因承认祂而被交付死亡并丧失生命的救主。\par

5. 然而，如果祂自己并不受苦，而是应从耶稣身上逃离，那么祂为何劝诫门徒背起十字架跟隨祂呢？——那十字架正是这些人声称祂未曾背起的，并且［他们将祂说成是］放弃了受苦的计划。因为祂说这话并非指承认上文的\OriginalTerm{十字架}{Stauros}，如他们中有些人胆敢解释的那样，而是指祂自己将要承受的苦难，以及门徒应忍受的苦难。祂说：\Quote{因为凡要救自己生命的，必丧掉生命；凡为祂丧掉生命的，必得着生命。}祂对犹太人说：\BibleQuote{看，我打发先知、贤士和经师到你们这里来；你们要杀害、钉死他们中的一些。}\BibleRef{玛 23:24}祂常对门徒说：\BibleQuote{你们要为我的缘故被带到总督和君王面前；他们要鞭打你们中的一些人，杀害你们，从这城迫害到那城。}\BibleRef{玛 10:17}因此，祂知道哪些人会受迫害，也知道哪些人会因祂而被鞭打、被杀害；祂所说的并非别的十字架，而是祂自己首先承受、门徒随后当受的苦难。为此祂劝勉他们说：\BibleQuote{不要怕那些杀害身体却不能杀害灵魂的；但要怕那能把灵魂和身体都投入地狱的。}\BibleRef{玛 10:28}［以此劝勉他们］持守对祂所作的信仰宣认。因为祂应许要在祂的父面前承认那些在人前承认祂名的人；却说要否认那些否认祂的人，并要以那些羞于承认祂的人为耻。尽管这些事情如此，这些人中有些竟胆大妄为到藐视殉道者，辱骂那些因承认主而被杀并承受主所预言一切的人——他们在这一点上努力追随主苦难的足迹，成为那受苦者的殉道者；我们也将他们列入殉道者之列。因为当追讨他们的血债，他们获得光荣时，一切曾诋毁他们殉道的人，都要被基督所羞辱。祂在十字架上呼喊：\BibleQuote{父啊，赦免他们，因为他们不知道自己做的是什么。}\BibleRef{路 23:34}由此显明基督的宽忍、忍耐、怜悯和良善，因为祂受苦，且亲自为虐待祂的人辩护。因为天主的圣言曾对我们说：\BibleQuote{爱你们的仇敌，为那仇恨你们的人祈祷。}\BibleRef{玛 5:44}祂自己在十字架上正是这样行了；祂如此爱人类，甚至为处死祂的人祈祷。然而，如果有人假设有两位基督，并对他们作出判断，那么那位［基督］将被视为更优秀、更忍耐、真正良善的一位——他在自己的创伤、鞭伤和所受的其他苦难中仍施恩惠，不念所遭受的冤屈——而那位逃遁、未受伤也未受辱的则不然。\par

6. 这一点同样也针对那些主张祂只是貌似受苦的人。因为如果祂并非真实受苦，那就毫无功劳可言，因为根本没有苦难；当我们真正开始受苦时，祂就好像欺骗了我们，劝我们忍受殴打并\BibleQuote{转过另一面脸来}\BibleRef{玛 5:39}，而祂自己并没有在我们之前真实地同样受苦；并且正如祂以貌似而非真实的形象误导了他们，祂也同样误导我们，劝我们忍受祂自己所没有忍受的。［那样的话］我们甚至超过了师傅，因为我们忍受和担当了我们的主从未承担或忍受的苦。然而，正如我们的主独自是真正的师傅，同样，天主子是真正良善和忍耐的，天父的圣言已变成了人子。因为祂争战并得胜了；因为祂是人，为父辈们而战，藉着服从彻底消除了不服从：因为祂\BibleQuote{捆绑了壮士}\BibleRef{玛 12:29}，释放了弱者，通过摧毁罪恶，赋予祂亲手所造的救恩。因为祂是至圣至慈的主，热爱人类。\par

7. 因此，正如我所说的，祂使人性依附并成为与天主合一。因为除非人战胜了人的仇敌，否则仇敌就不会被合法地击败。而且，除非是天主无偿赐予救恩，我们绝不能稳固地拥有它。并且，除非人已经与天主结合，否则他绝不能分享不朽。因为天主与人之间的中保，由于祂与双方的关系，必须使双方达到友谊与和谐，将人呈献给天主，同时向人启示天主。因为，我们若未曾通过圣子从祂那里领受那关乎祂自身的共融，除非祂的圣言降生成人，进入与我们共融，我们怎能分享义子名分呢？因此，祂也经历了生命的每个阶段，为所有人恢复与天主的共融。所以，那些声称祂是幻现地显现，既非在肉身中生，亦非真正成为人的人，仍处于旧有的定罪之下，为罪提供庇护；因为按照他们的说法，死亡尚未被征服，它\BibleQuote{从亚当直到梅瑟统治了那些没有像亚当的违命一样犯罪的人。}\BibleRef{罗 5:14} 但是，藉着梅瑟赐下的法律来临，为罪作证判定它为罪人，确实夺去了（死亡的）王国，显示出它不是君王，而是强盗；并揭示它是凶手。然而，它给人加上了沉重的负担，因为人自身有罪，表明他该受死亡。因为法律是属神的，它只是把罪凸显出来，却没有摧毁它。因为罪不是统治那属神的，而是统治人。因为那将要摧毁罪、救赎处于死亡权势之下的人，祂自己必须成为那被罪牵引进入束缚、被死亡所控制的人所是的，也就是说，人；以便罪被人摧毁，而人从死亡中走出去。因为正如 \BibleQuote{因那最初从处女泥土塑造出来的一个人的不服从，许多人成了罪人}\BibleRef{罗 5:19}，并丧失了生命；同样，因那最初从童贞女出生的一人的服从，许多人将得以成义并接受救恩。这样，天主的圣言成了人，正如梅瑟也说：\BibleQuote{天主，祂的作为是真实的。}\BibleRef{申 32:4} 但如果没有成为肉身，祂只是看起来像肉身，那么祂的工作就不是真实的。而祂显现为什么，祂就是什么：天主在祂内重演了人古老的形态，为杀死罪，剥夺死亡的权势，并使人生活；因此，祂的作为是真实的。\par

\SourceRecord{Translated by Alexander Roberts and William Rambaut.；Ante-Nicene Fathers；Vol. 1.；Edited by Alexander Roberts, James Donaldson, and A. Cleveland Coxe.；Buffalo, NY: Christian Literature Publishing Co.,；1885.；<http://www.newadvent.org/fathers/0103318.htm>.}

% structure: p0001 source=h1 target=section reason=page-title

\section{驳斥异端（第三卷，第十九章）}

1. 但是，那些断言祂只是一个由若瑟所生的凡人，仍处在旧悖逆的奴役中的人，是处于死亡的状态，因为他们尚未与天主圣父的圣言结合，也没有通过圣子获得自由，正如祂亲自宣称的：\BibleQuote{如果圣子使你们自由，你们就真的自由了。}\BibleRef{若 8:36} 但是，由于他们不认识那位由童贞女而来的厄玛奴耳，便失去了祂的恩赐，即永生；\BibleRef{罗 6:23} 并且，因不接受那不朽的圣言，他们仍处于必死的肉身中，成为死亡的债务人，未能获得生命的解毒剂。圣言对他们说，提及祂自己的恩宠恩赐：\Quote{我曾说：你们都是至高者的儿子，都是神；但你们要像人一样死去。} 毫无疑问，祂这些话是对那些没有接受义子恩赐的人说的，他们蔑视天主圣言纯洁生育的降生，剥夺了人性提升为天主之子的可能性，并证明自己对那为他们而成为血肉的天主圣言忘恩负义。因为正是为此目的，天主圣言成了人，那天主圣子成为人子，为使人被纳入圣言，接受义子名分，而成为天主子。因为我们不可能通过其他方式达到不朽与不死，除非我们与不朽和不死结合。但我们如何能与不朽和不死结合，除非首先不朽和不死成为我们也是的样式，好使可朽的被不朽所吞没，必死的被不死所吞没，为使我们能接受儿子的义子名分？\par

2. 为此，（经上说：）\BibleQuote{谁能述说祂的世代？}\BibleRef{依 53:8} 因为\BibleQuote{祂是一个人，谁能认识祂？}\BibleRef{耶 17:9} 但那受天上之父启示的人，\BibleRef{玛 16:16} 认识祂，以致明白那\BibleQuote{不是由血肉的意愿，也不是由人的意愿所生}\BibleRef{若 1:13} 的人子，就是基督，永生天主之子。因为我已经从圣经中证明，亚当的众子中没有一个在一切方面绝对地被称为天主或称为主。但是，祂本身，超越所有曾活过的人，是天主、是主、是永恒的王、是降生成人的圣言，为众先知、宗徒和圣神自己所宣告，凡获得一点点真理的人都能看见。圣经不会为祂作这些见证，如果祂像其他人一样仅是一个凡人。但是，祂超越所有其他人，在自己内拥有那来自至高圣父的卓绝出生，也经历了那来自童贞女的卓绝生育，\BibleRef{依 7:14} 神圣的圣经在两方面都为祂作证：祂是一个没有仪容、且能受苦的人；\BibleRef{依 53:2} 祂骑在驴驹上；\BibleRef{匝 9:9} 祂喝醋和苦胆；祂在百姓中被藐视，自己谦抑至死；并且祂是圣主、奇妙、策士、容貌秀美、全能的天主，\BibleRef{依 9:6} 驾云而来作所有人的审判者，\BibleRef{达 7:13} ——这一切都是圣经对祂的预言。\par

3. 因为正如祂成为人为了受试探，同样祂也是圣言，为要受光荣；圣言保持静默，为能受试探、受凌辱、被钉十字架并承受死亡，但人性被天主性吞没，当人性得胜、坚忍、行善、复活并被接升天时。因此，我们的主，天主子，既是圣父的圣言，又是人子，因为祂的人性出自玛利亚——她出自人类，本身也是人——祂成了人子。\BibleRef{依 7:13} 因此，主自己也给了我们一个记号，在地的深处，在天的高处，这是人所没有求的，因为人从未指望一个童贞女能怀孕，或者可能一个保持童贞的女子能生子，并且所生的乃是\BibleQuote{\Emph{厄玛奴耳，天主与我们同在}}，并下降到地下的事物中，寻找那迷失的羊——那羊正是祂自己的独特手工作品——并升到高处，将那被寻获的\OriginalTerm{人性}{hominem}献给并交托给祂的圣父，在自己身上作了人复活的初果；使元首从死者中复活，同样，身体的其余部分——每个活着的人的身体——当那因悖逆而存在的判决之时满足后，也能复活，通过关节和韧带\BibleRef{弗 4:16} 在天主的增长中交织并坚固，各肢体在身体中各有其适当的位置。因为在圣父的家里有许多住处，\BibleRef{若 14:2} 因为身体上也有许多肢体。\par

\SourceRecord{Translated by Alexander Roberts and William Rambaut.；Ante-Nicene Fathers；Vol. 1.；Edited by Alexander Roberts, James Donaldson, and A. Cleveland Coxe.；Buffalo, NY: Christian Literature Publishing Co.,；1885.；<http://www.newadvent.org/fathers/0103319.htm>.}

% structure: p0001 source=h1 target=section reason=page-title

\section{驳斥异端（第三卷，第二十章）}

1. 因此，当人成为违约者时，天主忍耐了，因为祂预见了那要藉着圣言赐予人的胜利。因为，当能力在软弱中成了完全的\BibleRef{格后 12:9}，就显示了天主的仁慈和超越的大能。正如祂容忍约纳被鲸鱼吞下，不是要他完全被吞灭而消亡，而是当他再次被吐出后，能更服从天主，更能光荣那赐予他如此意想不到的救恩的那一位，并能带领尼尼微人持久的悔改，使他们归向那拯救他们脱离死亡的主，因约纳身上所行的这个预兆而畏惧，正如圣经论及他们所说的：\BibleQuote{他们就各自离开自己的邪道，以及他们手中的不义，说：谁知道天主是否会后悔，转离祂的怒气，使我们不致灭亡呢？}\BibleRef{约纳 3:8}——同样地，从起初，天主就允许人被那大鱼（即过犯的创始者）吞下，不是要他被吞入时完全灭亡，而是安排并预备那藉着圣言并通过约纳的预象而完成的救恩计划，为那些与约纳持相同意见、论及主，并且承认说：\BibleQuote{我是主的仆人，我敬拜天上的主天主，是祂造了海和旱地。}\BibleRef{约纳 1:9}的人。这是为要使人从天主获得意想不到的救恩，从死者中复活，光荣天主，并重复那由约纳预言所说的话：\BibleQuote{我在患难中呼求我的天主上主，祂就从阴府的深处垂听了我；}\BibleRef{约纳 2:2}并且使他能不断光荣天主，不停地感谢那从他获得的救恩，\BibleQuote{为使没有血肉之躯能在主前自夸；}\BibleRef{格前 1:29}并且使人决不可对天主持有相反的看法，假设属于他的不朽是他天生就有的，因而不握持真理，以空洞的傲慢自夸，好像他天生就如天主一样。因为牠（撒殚）就使人更忘恩负义，对待造物主，遮蔽了天主对人的爱，蒙蔽了他的心智，使他不能觉察什么合乎天主，将自己与天主比较，并自认为与天主同等。\par

2. 因此，这就是天主的忍耐的目的：使人经历万事，获得道德纪律的知识，然后达到从死者中的复活，并由经验学习他得救的根源，以便常常生活在感恩主的状态，从祂那里领受不朽的恩赐，使他更爱祂；因为\BibleQuote{那多得宽恕的，爱得也多；}\BibleRef{路 7:43}并且使他认识自己，是多么可朽和软弱；同时他也了解天主，知道祂是不朽和全能的，足以将不朽赐给可朽的，将永恒赐给短暂的；也了解天主向他显示的其他属性，藉此受了教导，能合乎天主的伟大来思想天主。因为人的光荣在于天主，但天主的化工是天主的光荣；而收纳祂全部智慧和能力的容器是人。正如医生藉病人得到证明，天主也藉着人而显示。因此保禄宣告：\BibleQuote{因为天主把众人都禁锢在不信服之中，是为了要怜悯众人；}\BibleRef{罗 11:32}他不是指着属灵的永世说这句话，而是指着那曾悖逆天主、被剥夺了不朽的人，然后得到了怜悯，通过天主子接受了由祂自身完成的义子名分。因为那没有骄傲和夸耀地持有对受造物和造物主——祂是万有的全能天主，并赐予一切存在——的真实光荣（看法）的人，[这样的人]继续存留在祂的爱内\BibleRef{若 15:9}和顺服中，并感恩，也必将从祂领受更大的晋升的光荣，期待那时刻，他将相似那为他而死的那一位，因为祂也同样\BibleQuote{成了罪恶肉身的模样，}\BibleRef{罗 8:3}，为定断罪恶，并将罪恶如今作为已被定断之物逐出肉身之外，而是为要召唤人进入祂自己的相似，指派他作为[祂自己的]效法者归向天主，并将祂父的法律加给他，使他能看见天主，并赐给他权力接纳父；[祂本是]居住在人心中的天主的圣言，并成了人子，为要使人习惯于接纳天主，并使天主习惯于居住在人间，按父的美意。\par

3. 因此，主自己，就是从童贞女而来的厄玛奴耳\BibleRef{依 7:4}，是我们救恩的预兆，因为正是主自己拯救了他们，因为他们不能靠自己的手段得救；因此，当保禄陈述人的软弱时，他说：\Quote{我知道在我的肉身内，没有善；}\BibleRef{罗 7:18}表明我们救恩的\Quote{善}不是出于我们，而是出于天主。又说：\Quote{我真不幸！谁能救我脱离这死亡的身体呢？}\BibleRef{罗 7:24}然后他引入了拯救者，说：\Quote{我们主耶稣基督的恩宠。}依撒意亚也宣布了这一点，[说：]\BibleQuote{你们这些下垂的手、无力的膝，要坚强；你们这些心怯的人，要鼓起勇气，不要害怕；看，我们的天主已施行审判，带着报应，祂要酬报；祂必亲自来拯救我们。}\BibleRef{依 25:3}由此我们看到，我们得救不是靠自己，而是靠天主的帮助。\par

4. 再者，为使我们不被一个纯粹的人拯救，也不是一个没有血肉的——因为天使是没有血肉的——那位同一位先知宣布说：\BibleQuote{既不是长者，也不是天使，而是主亲自拯救他们，因为祂爱他们，并且宽恕他们；祂亲自解救他们。} \BibleRef{依 63:9} 并且祂自己将要成为真实的人，有形的，当祂是赐予救恩的圣言时，依撒意亚又说：\BibleQuote{看，熙雍城：你的眼睛将看见我们的救恩。} \BibleRef{依 33:20} 而且那为我们而死的不是一个纯粹的人，依撒意亚说：\BibleQuote{圣主记念祂死去的以色列，那些曾安睡在坟墓之地的人；祂降临向他们宣讲祂的救恩，为拯救他们。} 亚毛斯（米该亚）先知也同样宣布：\BibleQuote{祂将回转，并怜悯我们；祂将消灭我们的不义，将我们的罪投入深海。} \BibleRef{米 7:9} 再者，指明祂来临的地方，他说：\BibleQuote{主从熙雍发言，从耶路撒冷发出祂的声音。} \BibleRef{岳 3:16} 而且，从犹大产业南方的那一地区，天主子将要来临，祂是天主，且来自白冷，主在那里诞生，祂的赞美将传遍全地，哈巴谷先知这样说：\BibleQuote{天主要从南方而来，圣者从厄弗辣因山而来。祂的威能遮蔽诸天，大地充满祂的赞美。在祂面前，圣言将行进，祂的脚要在平原上前进。} 这样，祂清楚地表明祂是天主，祂的降临是在白冷，从位于产业南方的厄弗辣因山，并且祂是人。因为他说：\BibleQuote{祂的脚要在平原上前进：} 这恰好是人的标志。\par

\SourceRecord{Translated by Alexander Roberts and William Rambaut.；Ante-Nicene Fathers；Vol. 1.；Edited by Alexander Roberts, James Donaldson, and A. Cleveland Coxe.；Buffalo, NY: Christian Literature Publishing Co.,；1885.；<http://www.newadvent.org/fathers/0103320.htm>.}

% structure: p0001 source=h1 target=section reason=page-title

\section{\WorkTitle{驳斥异端}（第三卷，第21章）}

1. 天主于是成了人，主亲自拯救了我们，赐给我们童贞女的记号。但并非如某些人（那些现今胆敢解释圣经的人）所声称的：\BibleQuote{看，一位童贞女要怀孕生子，}\BibleRef{依 7:14}——如同厄弗所的\Person{特敖多提翁}，以及本都的\Person{阿基拉}（两人都是犹太归化者）所翻译的。根据他们，\Person{厄俾约尼特人}断言主是由\Person{若瑟}所生；这样，他们尽其可能地摧毁了天主如此奇妙的救恩计划，并且弃绝了来自天主的先知们的见证。确实，这个预言是在人民被掳到\Place{巴比伦}之前发出的；也就是说，早于玛待人和波斯人获得霸权。但它是由犹太人自己译成希腊文的，远在我们的主降临之前，为的是使人无法怀疑犹太人可能是在迎合我们的趣味而这样解释这些言辞。事实上，如果犹太人预知到我们将来会存在，并且我们会使用这些来自圣经的证明，他们自己会毫不犹豫地焚烧自己的圣经——这些圣经确实宣告所有外邦民族都分享永生，并表明那些以\Person{雅各伯}家和以色列民自夸的人被剥夺了天主的恩宠。\par

2. 因为在罗马人统治他们的王国之前，当马其顿人仍占据\Place{亚细亚}的时候，\Person{托勒密}（\Person{拉古斯}之子）渴望用所有有价值的人的著作来装饰他在\Place{亚历山大}建立的图书馆，于是向\Place{耶路撒冷}居民提出请求，要他们将圣经译成希腊文。他们——因为那时他们仍臣服于马其顿人——就派了七十位精通圣经和两种语言的长老到\Person{托勒密}那里，去完成他所期望的工作。但\Person{托勒密}想要个别地测试他们，并担心他们或许会通过商议而用他们的翻译掩盖圣经中的真理，就把他们彼此分开，命令他们都写出相同的译文。他对所有的书籍都这样做。但当他们一同来到\Person{托勒密}面前，每个人都把自己的译文与其他人的比较时，天主确实受到光荣，圣经也被承认真是天主的。因为他们所有人都从头到尾用完全相同的词语和名称朗读出共同的译本（他们准备好的），以至于在场的甚至外邦人也看出圣经是由天主的默感翻译的。天主这样做并不奇怪——祂在\Person{拿步高}统治下人民被掳、圣经被毁坏，七十年后犹太人返回故土时，在波斯王\Person{阿尔塔薛西斯}时代，感动了肋未支派的司祭\Person{厄斯德拉}，让他重写所有先知的言语，并重新向人民确立\Person{梅瑟}的律法。\par

3. 既然圣经以如此的信实并通过天主的恩宠被翻译，并且天主从这些圣经中预备并重新塑造了我们对祂圣子的信德，又将无掺杂的圣经保存在\Place{埃及}——在那里\Person{雅各伯}家因躲避\Place{客纳罕}的饥荒而兴盛；在那里我们的主逃避\Person{黑落德}发起的迫害时也被保存；并且这些圣经的翻译是在我们的主降世之前完成的，是在基督徒出现之前就存在的——因为我们的主大约在\Person{奥古斯都}在位的第四十一年诞生；而\Person{托勒密}（在他在位时圣经被翻译）要早得多——[既然这些事如此，我断言，]这些人确实被证明是无耻和放肆的，他们现在竟想作出不同的翻译，而我们用这些圣经驳斥他们，并将他们限制在对天主圣子来临的信仰之中。但\Emph{我们的}信德是坚定的、无伪的、唯一真实的，有从这些圣经中得来的明确证明，这些圣经按我所述的方式被翻译；而教会的宣讲没有增添。因为宗徒们——他们比所有这些异端者年代更古老——都赞同上述译本；并且这译本与宗徒的传统和谐一致。因为\Person{伯多禄}、\Person{若望}、\Person{玛窦}、\Person{保禄}以及其余的宗徒和他们的继承者，都宣示了所有先知性的宣告，正如同长老们的译本所包含的那样。\par

4. 因为同一天主之神，曾借着先知预告主的来临是怎样的，也借着这些长老真正预言的正确解释；并且祂自己曾借着宗徒宣布，义子之期的圆满已来到，天国临近了，且祂自己正居住在那些相信由童贞女所生的厄玛奴耳的人心中。为此，他们作证说，在若瑟与玛利亚同居之前，她尚保持童贞时，\BibleQuote{她因圣神怀孕了}\BibleRef{玛 1:18}；天使加俾额尔对她说：\BibleQuote{圣神要临于你，至高者的能力要庇荫你，因此那要诞生的圣者，将称为天主子。}\BibleRef{路 1:35}；天使在梦中告诉若瑟说：\BibleQuote{这一切事的发生，是为应验上主借先知所说的话：看，一位贞女将怀孕生子。}\BibleRef{玛 1:23}。但长老们这样解释依撒意亚所说的：\BibleQuote{上主又对阿哈次说：你向你的天主上主求一个征兆，或来自阴府深处，或来自上天高处。阿哈次说：我不求，我不试探上主。依撒意亚说：你们使人厌烦还算小事，还要使我的天主厌烦吗？因此吾主要给你们一个征兆：看，一位贞女要怀孕生子，给他起名叫厄玛奴耳。他要吃奶酪和蜂蜜，直到他知道弃恶择善的时候，因为孩子尚不知道善恶，他就不愿作恶，以便选择善。}\BibleRef{依 7:10}。因此，圣神借上述的话仔细指出了祂由贞女所生，以及祂的本质，即祂是天主（因为\OriginalTerm{厄玛奴耳}{Emmanuel}这名表明这一点）。祂又表明祂是人，当祂说：\Quote{他要吃奶酪和蜂蜜}，而且称祂为婴孩，说：\Quote{在他知道善恶之前}，因为这些都是人类婴孩的标记。但祂\Quote{不愿作恶，以便选择善}——这却是天主所固有的；这样，因祂要吃奶酪和蜂蜜，我们不应以为祂仅是一个纯粹的人，另一方面，从厄玛奴耳这名，也不应怀疑祂是无血肉的天主。\par

5. 当祂说：\BibleQuote{达味家，请听！}\BibleRef{依 7:13}，祂所做的是表明天主应许达味从他自己腹中所生的后裔中兴起一位永远的君王，正是那位由童贞女所生、而童贞女本身出自达味家族的。因此，天主应许这位君王\Quote{从他腹中的果实}（原文为\Emph{腹}的果实），这是适合童贞女怀孕的说法，而不是\Quote{从他腰间的果实}或\Quote{从他肾间的果实}，因为后者适合男人生育以及女人由男人怀孕。因此，在此应许中，圣经排除了任何男性的影响；然而，它当然没有说那所生者不是出于人的意愿。但它固定并确立了\Quote{腹中的果实}，以表明那一位由童贞女所生的诞生，正如依撒伯尔充满圣神时所证明的，她对玛利亚说：\BibleQuote{你在妇女中应受赞美，你腹中的胎儿也应受赞美。}\BibleRef{路 1:42}圣神向愿意听的人指出，天主所作的应许——从[达味的]腹中兴起一位君王——是在童贞女的诞生中实现的，即出自玛利亚。因此，那些把依撒意亚的这段话篡改为\Quote{看，一位青年女子要怀孕}，并认为他是若瑟之子的人，也应该篡改天主给达味的应许的形式，即天主应许他从腹中兴起基督君王的角。但他们不明白，否则他们也会胆敢篡改这一节经文。\par

6. 但依撒意亚所说的：\BibleQuote{或来自阴府深处，或来自上天高处}\BibleRef{依 7:11}，是要表明\BibleQuote{那降下的，正是那升上的}\BibleRef{厄 4:10}。而他说\Quote{上主自己要给你们一个征兆}时，他宣布了关于祂的诞生的一件意想不到的事，这件事除非由万有的天主上主亲自给达味家一个征兆，否则无法以其他方式完成。因为一个青年女子由男人怀孕生子——这是所有生子的妇女都会发生的事，有什么伟大或什么征兆呢？但由于天主的帮助要为人提供意想不到的救恩，所以也成就了由童贞女意想不到的诞生；天主给了这个征兆，而人并没有成就它。\par

7. 为此，达尼尔\BibleRef{达 2:34}预见了祂的来临，说有一块非人手凿成的石头来到了这个世界。因为\Quote{非人手}的意思，就是祂来到这个世界不是借着人手的操作，即那些习惯凿石头的人的手；也就是说，若瑟完全没有参与，只有玛利亚单独配合那预先安排的计划。因为这块石头从大地获得存在，是出自天主的德能和智慧。因此依撒意亚也说：\BibleQuote{上主这样说：看，我在熙雍奠放一块石头，一块宝贵的、精选的、首要的、角石，应受尊重。}\BibleRef{依 28:16}因此，我们明白祂的人性来临不是出于人的意愿，而是出于天主的意愿。\par

8. 因此，梅瑟也立了一个预像，将他的杖投在地上\BibleRef{出 7:9}，为使那杖藉着成为血肉，揭露并吞没埃及人所有反对天主预定计划的对抗\BibleRef{出 8:19}，好让埃及人自己作证：是天主的手指为百姓施行救恩，而非若瑟的儿子。因为如果他是若瑟的儿子，怎能大于撒罗满，或大于约纳\BibleRef{玛 12:41}，或大于达味\BibleRef{玛 22:43}呢？他是从同样的血统而生，是这些人的后裔。又如何他称伯多禄为有福的，因为伯多禄承认他是永生天主之子？\BibleRef{玛 16:17}\par

9. 然而，此外，如果他真是若瑟的儿子，按耶肋米亚所说，他就不能是君王或继承人了。因为若瑟被证明是约雅金和耶苛尼雅之子，正如玛窦在其家谱中所陈述的\BibleRef{玛 1:12}。但耶苛尼雅和他所有的后裔都被剥夺了王位的继承权；耶肋米亚如此宣告说：\Quote{我指着我的永生起誓，上主说：如果犹大王约雅金的儿子耶苛尼雅曾是我右手的印玺，我也必将你从那里拔除，交在寻索你性命的人手中。}\BibleRef{耶 22:24}又说：\Quote{耶苛尼雅像一个不中用器皿，被抛弃了，因为他被抛到一个他不认识的地方。大地，听上主的话吧！你应当记录：这个人是一个被剥夺继承权的人，因为他的后裔中没有一人能坐在达味的宝座上，或在犹大中作首领。}\BibleRef{耶 22:28}此外，天主论及他的父亲约雅金说：\Quote{因此，上主关于犹大王约雅金如此说：他必没有人在达味的宝座上坐位；他的尸首必被抛掷在日间的炎热和夜间的寒霜中。我必惩罚他和他的儿子，并他们和耶路撒冷的居民以及犹大地上的人，降与他们所说的一切灾祸。}\BibleRef{耶 36:30}因此，那些说他是从若瑟所生，并在他身上寄予希望的人，是使自己被剥夺了王国的继承权，落在针对耶苛尼雅及其后裔的诅咒和斥责之下。因为关于耶苛尼雅说了这些话，是圣神预知了那些邪恶教师的教训，好叫他们知道：他必不从他的后裔——即从若瑟——而生，但按照天主的应许，从达味的腹中兴起永恒的王，他在自己内总括万有，把古时的人性收聚于自身。\par

10. 因为正如因一人的悖逆，罪恶进入了世界，藉着罪恶死亡也得了地位\BibleRef{罗 5:19}，照样，因一人的服从，正义被引入，将使生命在那些过去已死的人身上结出果实。又正如元祖亚当自己，他的质料是从未耕种的、尚是童贞的土壤而来（\Quote{因为天主还没有降雨在地上，也没有人耕种田地}——见：\BibleRef{创 2:5}），并由天主的手形成，即由天主的圣言形成，因为\Quote{万物是藉着他造成的}\BibleRef{若 1:3}，是主从地上取尘土形成了人；同样，那圣言，在自己内总括亚当，理应从尚是童贞的玛利亚接受一个出生，使他能收聚亚当归于自己。如果第一个亚当有一个人作为父亲，并由人种而生，那么说第二个亚当是从若瑟所生也合情合理。但如果前者是从尘土所取，而天主是他的造主，那么后者也必须在自身内进行总括，由天主按人形形成，在起源上与前者保持类似。那么，为何天主不再取尘土，却使形成出自玛利亚呢？这是为了不使另一个受造物被召唤出来，也没有另一个需要被拯救，而是使同一个受造物被总括在基督内，如曾在亚当内一样，相似性得以保存。\par

\SourceRecord{Translated by Alexander Roberts and William Rambaut.；Ante-Nicene Fathers；Vol. 1.；Edited by Alexander Roberts, James Donaldson, and A. Cleveland Coxe.；Buffalo, NY: Christian Literature Publishing Co.,；1885.；<http://www.newadvent.org/fathers/0103321.htm>.}

% structure: p0001 source=h1 target=section reason=page-title

\section{\WorkTitle{驳斥异端（第三卷，第二十二章）}}

1. 因此，那些声称祂从童贞玛利亚一无所取的人大大错误；为了抛弃肉身的遗产，他们也拒绝了〔祂与亚当之间的〕类比。因为若那出于尘土的一位确实由天主的手和化工所形成并具有实体，而另一位却非出于天主的手和化工，那么那位按前者的肖像和模样受造者，便未保存人的类比，祂必显为一件不一致的作品，没有可显示其智慧之处。这就是说，祂也仅仅外表上显现为人而并非真人，且祂成为人时并未从人取任何事物。因为若祂未从人取得肉身的实体，祂便既未成为人，也未成为人子；若祂未成为我们之所是，祂在所受的苦和所忍受的事上便未成就大事。然而，人人都将承认我们有取自尘土的身体和由天主接受神明的灵魂。因此，天主圣言成为这样，在祂自己内\OriginalTerm{总归}{recapitulatio}了祂自己的工程；因此祂承认自己是人子，并祝福\BibleQuote{温良的人，因为他们要承受土地。}\BibleRef{玛 5:5} 此外，宗徒保禄在致迦拉达人书中明言：\BibleQuote{天主派遣了自己的儿子，生于女人。}\BibleRef{迦 4:4} 又，在致罗马人书中，他说：\BibleQuote{论到祂的儿子，按肉身说，是从达味的后裔生的，按圣德的神性说，由于从死者中复活，被立为具有威能的天主之子，我们的主耶稣基督。}\BibleRef{罗 1:3}\par

2. 在这种情况下，祂降至玛利亚也是多余的；因为若祂不从她取任何事物，为何要进入她内？再者，若祂未从玛利亚取任何事物，祂绝不会享用那些由土地而来的食物，因为那由土地所取的身体正是靠这些营养；祂也不会在禁食四十天时感到饥饿，如同梅瑟和厄里亚那样，除非祂的身体渴求合宜的营养；祂的门徒若望在写到祂时也不会说：\BibleQuote{但耶稣因行路疲倦，就坐着歇息}；\BibleRef{若 4:6} 达味也不会预先宣告说：\Quote{他们加重了我伤口的疼痛}；祂也不会为拉匝禄流泪，也不会汗滴如血；也不会说：\BibleQuote{我的心灵忧闷得要死}；\BibleRef{玛 26:38} 当祂的肋旁被刺透时，也不会流出血和水。因为这一切都是那由土地所取之肉身的标记，祂将之\OriginalTerm{总归}{recapitulatus}在祂自己内，为祂自己的工程带来救恩。\par

3. 因此，路加指出追溯我们主世代直至亚当的家谱包含七十二代，将终始相连，暗示正是祂在自己内总括了从亚当以来分散的一切民族、一切语言和世代，连同亚当本人。因此，亚当也被保禄称为\BibleQuote{那要来的那位的预像}，\BibleRef{罗 5:14} 因为圣言，万物的创造者，预先为祂自己规定了与天主之子相连的人类未来的安排；天主预定了第一个人应具有动物性的本性，目的在于他可被属神者所拯救。因为既然祂作为拯救者先已存在，那么那可能被拯救者也应被召叫存在，免得拯救者徒然存在。\par

4. 与此计划相符，童贞玛利亚显出顺服，说：\BibleQuote{看！主的婢女，愿照你的话成就在我身上！}\BibleRef{路 1:38} 但厄娃却不服从；因为她还是童贞时就没有听从。事实上，她虽有丈夫亚当，却仍为童贞（因为在乐园里\BibleQuote{二人赤身露体，并不害羞}，\BibleRef{创 2:25} 因为他们刚受造不久，尚不知生育之事：他们必须先达成年，然后自那时起繁衍），她由于不服从，成为死亡的原因，既对自己也对着整个人类；同样，玛利亚虽已有未婚夫，却仍为童贞，由于顺服，成为救恩的原因，既对自己也对着整个人类。因此，法律称许配给丈夫的女子为他的妻子，尽管她仍是童贞；这暗示从玛利亚回溯至厄娃，因为联合在一起的事物除非通过翻转其结合过程的顺序，否则无法分离；这样，后者的纽带取消前者，后者使前者重新得释。事实上，这已发生：第一个结合因第二个纽带而解开，而第二个纽带取代了已被取消的第一个的位置。为此，主宣称\BibleQuote{为首的将成为最后的，最后的将成为首先的。}\BibleRef{玛 19:30} 先知也指明同一件事，说：\Quote{代替你的父亲们，为你生了儿子们。}因为主，生于死者中成为\BibleQuote{\OriginalTerm{首生者}{primogenitus}}，\BibleRef{默 1:5} 将古圣祖们拥入怀中，使他们重生进入天主的生命，祂自己成为生活者之始，正如亚当成为死者之始。\BibleRef{格前 15:20} 因此，路加从主开始家谱，一直追溯到亚当，暗示正是祂使他们重生进入生命的福音，而非他们使祂重生。同样，厄娃的\OriginalTerm{悖逆之结}{nodus inoboedientiae}也因玛利亚的顺服而解开；因为童贞厄娃因不信所束缚的，童贞玛利亚因信德而释放了。\par

\SourceRecord{Translated by Alexander Roberts and William Rambaut.；Ante-Nicene Fathers；Vol. 1.；Edited by Alexander Roberts, James Donaldson, and A. Cleveland Coxe.；Buffalo, NY: Christian Literature Publishing Co.,；1885.；<http://www.newadvent.org/fathers/0103322.htm>.}

% structure: p0001 source=h1 target=section reason=page-title

\section{驳斥异端（第三卷，第二十三章）}

1. 因此，主来到那迷失的羊那里，进行如此全面安排的总归，并寻找祂亲手所造的工程，就有必要拯救那个按照祂的肖像和模样受造的人，即亚当，填满因不顺从而招致的定罪时期——这些时期是\BibleQuote{圣父以祂自己的权能所指定的。}\BibleRef{宗 1:7}这也是必要的，因为关乎人的整个救恩计划都是按照圣父的美意进行的，好使天主不被战胜，祂的智慧也不被（受造物眼中）贬低。因为如果那个天主为使他活着而创造的人，在被那腐蚀他的蛇伤害后失去了生命，再也不能返回生命，而是完全（永远）被抛弃于死亡，那么天主（在这种情况下）就被战胜了，蛇的邪恶就胜过了天主的旨意。但由于天主是不可战胜且长久忍耐的，祂在纠正人和试验一切的事上确实显出了长久忍耐，正如我已经指出的；并且藉着第二个人，祂捆绑了那壮士，抢走了他的财物，\BibleRef{玛 12:29}废除了死亡，使那处于死亡中的人活了过来。因为起初亚当成了撒殚手中的器皿，撒殚也确实用权势抓住他，即不义地将罪带到他身上，在永生的假象下将死亡带给他。因为蛇许诺他们将像神一样，而这是人绝不可能做到的，他却在他们身上造成了死亡；因此，那掳掠人的，被天主公义地捕获了；但那被掳掠的人，从定罪的捆绑中被释放了。\par

2. 但这个人是亚当，倘若应说真话，就是最初被造的人，关于他圣经说，主曾说：\BibleQuote{让我们按我们的肖像和模样造人}\BibleRef{创 1:26}；我们都从他而出，既然出自他，因此我们也都继承了他的称号。但既然人得了救，那被造为原始的人也应得救，才是合适的。因为坚持认为那被敌人深深伤害并首先遭受囚禁的人，没有被那战胜敌人的主所解救，而他的孩子——那些在同样囚禁中所生的——却得解救，这是过于荒谬的。如果旧有的掠物仍留在敌人手中，敌人也不算是已被战胜。举一个例子：如果一支敌军打败了一些（敌人），捆绑了他们，将他们掳去，并长期奴役他们，以致他们之中生下了孩子；然后有人怜悯那些做了奴隶的人，打败了同一支敌军；他若释放了那些被掳者的孩子，使他们摆脱奴役他们父亲的势力，却让那些遭受被掳之苦的人本身仍受敌人辖制——那些人正是他为报复的缘故而攻击的对象——而孩子却因替父亲报仇而获得自由，但父亲本身却未被解救（仍在捆绑中），那他行事就不公平。因为天主既非无能也非不义，祂帮助了人，使人恢复祂自己的自由。\par

3. 也是出于这个原因，在亚当犯罪之后，正如圣经所载，祂没有亲自诅咒亚当，而是因他的工作诅咒了土地，正如古人中某人所观察的：\Quote{天主确实将诅咒转移到土地上，以免它停留在人身上。}\BibleRef{创 3:16}但人因犯罪所受的惩罚是：耕种土地的劳苦，要汗流满面地吃面包，并且归于他所出自的尘土。同样，女人也（受了）劳苦、辛劳、叹息、生产的阵痛以及隶属的状态，即她应服侍丈夫；这样他们既不会因天主的诅咒而完全灭亡，也不会因未受责罚而变得藐视天主。但全然的诅咒落在那欺骗了他们的蛇身上。经上记载：\BibleQuote{天主对蛇说：因为你做了这事，你在一切牲畜和地上一切的野兽中是可诅咒的。}\BibleRef{创 3:14}同样，主也在福音中对那些在左边的人说：\Quote{离开我，你们这些可诅咒的，进入那为魔鬼和他的天使所预备的永火里去！}这表明永火原不是为人预备的，而是为那个欺骗了人并使人犯罪的——就是那背叛之首，以及和他一同成为背叛者的天使们；的确，那些像他一样坚持作恶、不知悔改、不肯回头的人，也将公义地感受到这火。\par

4. 这些人行事如同加音，当天主劝他安静，因为他没有公平地分给他兄弟应得的那一份，反而怀着嫉妒和恶意，以为可以压制他，他不但不顺服，反而罪上加罪，以行为表明了他的心态。他所计划的，他也付诸实行：他暴虐地杀了他的兄弟；天主使义者服从不义者，为叫义者因所受的苦难被证明是义人，而不义者因所行的被揭露是不义者。即便如此，他也没有心软，也没有停止那恶行；但当被问及他兄弟在哪里时，他说：\Quote{我不知道；难道我是我兄弟的看守者吗？}以此回答扩大并加重了他的邪恶。因为杀害兄弟已是邪恶，更何况这样傲慢无礼地回答全知的天主，仿佛能与祂争战。因此他自身背负了诅咒，因为他无故献上了罪恶的祭品，对天主毫无敬畏，也不因杀弟之行为而羞愧。\par

5. 然而，亚当的情况与此并不相似，而是全然不同。因为他被另一人以不死为借口所引诱，立刻陷入恐惧，并躲藏起来；并非因为他能逃避天主，而是因违背了祂的命令而慌乱，自觉不配在天主面前出现并与祂交谈。\BibleQuote{敬畏上主是智慧的开端；}\BibleRef{箴 1:7}对罪的认识引向悔改，天主向悔改者施予慈悲。因为亚当以行动表明了他的悔改，他用了腰带，用无花果叶遮盖自己，尽管有许多其他叶子，对身体刺激较小。然而他选择了与他的悖逆相称的服装，因敬畏天主而战兢；他抗拒了肉体的错谬和情欲的倾向（因为他失去了天然的本性和孩童般的心智，认识了恶事），为自己和妻子系上了节制的缰绳，敬畏天主，等待祂的来临，并仿佛表明了这样的事：他说，既然因悖逆失去了从圣神而来的圣洁外袍，我现在也承认我配得上这种不提供愉悦，反而咬啮、刺激身体的遮盖。如果慈悲的天主没有用皮衣代替无花果叶给他们穿上，他无疑会永远保留这样的衣服，以此自卑。为此，天主也审问他们，为叫罪责落到女人身上；又审问她，为叫她将罪责归给蛇。因为她讲述了所发生的事。\BibleQuote{蛇欺骗了我，我就吃了。}\BibleRef{创 3:13}但祂没有向蛇提问，因为祂知道他是罪行的始作俑者；但祂首先向他宣布诅咒，为叫这诅咒以较轻的责备落到人身上。因为天主憎恨那引人误入歧途的，但对那被欺骗的，祂渐渐、逐渐地显示慈悲。\par

6. 因此，祂也将他逐出乐园，使他远离生命树，并非如有些人妄言的那样，是祂嫉妒生命树，而是因为怜悯他，不愿他永远作罪人，也不愿环绕他的罪成为不朽的，邪恶成为无终的、不可补救的。祂为他的罪过设定了界限，藉着介入死亡，使罪止息，\BibleRef{罗 6:7}藉着肉体在尘世中的解体，使罪终结，这样，人终于停止为罪而活，死于罪，开始为天主而活。\par

7. 为此，祂在蛇与女人以及她的后裔之间放置了敌意，他们互相保持这敌意：祂的脚跟被咬，却有能力践踏仇敌的头；而另一方则咬伤、杀害、阻碍人的脚步，直到那被指定践踏其头的后裔来到——那后裔由童贞玛利亚所生，先知论及祂说：\Quote{你要践踏毒蛇和蝮蛇，你要蹂躏狮子和毒龙。}——这表明那针对人而设立、蔓延、使人服于死亡的罪，连同统治人的死亡，将被剥夺权能；而狮子，即敌基督，在末世猖獗攻击人类，将被祂践踏；祂要捆绑\BibleQuote{那龙，那古蛇}\BibleRef{默 20:2}，将他置于曾被征服的人的权下，\BibleRef{路 10:19}使他的全部势力被践踏。现在亚当已被征服，一切生命被夺去；因此，当仇敌反过来被征服时，亚当获得了新生命；最后的仇敌——死亡——被毁灭，\BibleRef{格前 15:26}这死亡起初占有了人。因此，当人被释放时，\BibleQuote{经上所记载的将应验：死亡被胜利吞灭了。死亡啊，你的刺在哪里？}\BibleRef{格前 15:54}这话若那首先被死亡辖制的人不得释放，就不能公正地说出。因为他的救恩就是死亡的毁灭。所以，当主使人，即亚当，活过来时，死亡同时被毁灭。\par

8. 因此，凡否认（亚当的）救恩的人，说的都是谎言，他们将自己永远排除在生命之外，因为他们不信那只迷失的羊已被找到。\BibleRef{路 15:4} 因为若它未被找到，全人类仍陷于丧亡状态。所以，那个首先提出这个想法——更确切地说，这种无知和盲目——的人，即\OriginalTerm{塔提安}{Tatian}，是虚假的。正如我已经指出的，此人已自陷于所有异端者之中。然而，这条教义是他自己杜撰的，为的是通过引入某种与其余人不同的新奇事物，并说虚妄之言，为自己赢得缺乏信德的听众，装作被尊为教师，并时常努力使用保罗常引用的话：\BibleQuote{在亚当内众人都死了；}\BibleRef{格前 15:22} 然而却不知道\BibleQuote{罪恶在那里越多，恩宠在那里也格外丰富。}\BibleRef{罗 5:20} 既然这事已清楚显明，就让他的所有门徒都感到羞愧吧，让他们为亚当争辩，仿佛若他不被拯救，他们就会得到莫大的好处；然而他们并未因此获益更多，正如蛇在诱惑人（犯罪）时也没有获益，只是证明了人是违命者，把人作为他自己背道的初果。但他不知道天主的权能。同样，那些否认亚当救恩的人也没有收获，只是使自己成为异端者和背离真理者，并显明自己是蛇和死亡的庇护者。\par

\SourceRecord{Translated by Alexander Roberts and William Rambaut.；Ante-Nicene Fathers；Vol. 1.；Edited by Alexander Roberts, James Donaldson, and A. Cleveland Coxe.；Buffalo, NY: Christian Literature Publishing Co.,；1885.；<http://www.newadvent.org/fathers/0103323.htm>.}

% structure: p0001 source=h1 target=section reason=page-title

\section{驳斥异端（第三卷，第二十四章）}

1. 如此，所有引入不敬神学、谈论我们的造物主和塑造者（祂也形成了这世界，除祂以外别无天主）的人都被揭露了；那些教导关于我们主的本质以及祂为祂的受造物——人——而完成的救恩计划的谎言的人，也被自己的论据推翻了。相反地，教会的宣讲处处一致，持续平稳地前进，并得到先知、宗徒和所有门徒的见证——正如我所证明的——通过起初、中间和终末，以及天主的整个救恩计划和那稳固的、导向人救恩的系统，即我们的信德。这信德从教会领受，我们予以保存，并且它常由天主之神更新青春，犹如一件精美器皿中的珍贵宝库，也使那盛放它的器皿本身更新青春。因为天主的这恩赐已托付给教会，如同气息托付给最初受造的人，为这目的：使所有领受它的肢体都得以活起来；而基督内的共融之媒介已遍布其中，即圣神，不朽的保证，坚固我们信德的媒介，以及上升天主的阶梯。\BibleQuote{因为在教会内，}经上说，\BibleQuote{天主设立了宗徒、先知、教师，}\BibleRef{格前 12:28}以及圣神所运作的所有其他媒介；凡不加入教会的人都不能分受这些，反而因自己的乖谬意见和可耻行为剥夺自己的生命。因为哪里教会，哪里就有天主之神；哪里天主之神，哪里就有教会和各种恩宠；但圣神是真理。因此，凡不分受祂的人，既不能从母怀中得滋养而得生命，也不能享受那从基督身体涌出的极清澈之泉源；反而为自己挖掘地上坑中的破裂水池\BibleRef{耶 2:13}，饮用泥沼中的臭水，逃避教会的信德以免被定罪；并拒绝圣神，以免受教导。\par

2. 这样疏远真理，他们理所当然在一切错误中打滚，被它抛来抛去，对同一事物在不同时候持不同看法，从未达到稳固的认识，更热衷于做辞藻的诡辩家而非真理的门徒。因为他们不是建立在唯一的盘石上，而是建立在沙土上，沙土本身包含许多石头。因此他们也想象出许多神，并且总有寻找真理的借口（因为他们瞎了眼），却从未成功找到。因为他们亵渎造物主——那位真实的天主，祂也提供找到真理的能力——却想象自己发现了另一个超越天主的神，或另一个\OriginalTerm{圆满}{Pleroma}，或另一种救恩安排。因此，来自天主的光也不照亮他们，因为他们羞辱和轻视天主，把祂看得很小，因为祂出于自己的爱和无限良善，进入了人类知识的范围（然而，知识不是关于祂的伟大或本质——因为没有人曾度量或把握过——而是这样的：我们应当知道，那制造、塑造并吹入生命气息，借着受造物滋养我们，以祂的圣言建立万有，以祂的\OriginalTerm{智慧}{Wisdom}维系万有的那位，即是唯一真实的天主）；但他们却梦想一个非实有者高于祂，以便被视为发现了那位伟大的神，而他们声称无人能认识这位神与人类交流或管理尘世事务；也就是说，他们发现了\Person{伊壁鸠鲁}的神，这神既不为自己也不为他人做任何事，即根本不施行眷顾。\par

\SourceRecord{Translated by Alexander Roberts and William Rambaut.；Ante-Nicene Fathers；Vol. 1.；Edited by Alexander Roberts, James Donaldson, and A. Cleveland Coxe.；Buffalo, NY: Christian Literature Publishing Co.,；1885.；<http://www.newadvent.org/fathers/0103324.htm>.}

% structure: p0001 source=h1 target=section reason=page-title

\section{\WorkTitle{驳斥异端} 第三卷 第二十五章}

1. 天主确实对万物施行\OriginalTerm{天主眷顾}{providence}，因此祂也赐予劝言；当赐予劝言时，祂与那些留意道德纪律的人同在。由此自然得出，凡受看管和治理的事物应认识它们的统治者；这些事物并非无理或虚妄的，它们由天主的眷顾获得理解。因此，一些外邦人——较少沉迷于感官的诱惑和逸乐，也未受偶像崇拜的迷信所引导——虽被祂的眷顾轻微触动，仍确信他们应称这宇宙的造物主为父，祂对万物施行眷顾，并安排我们世界的事务。\par

2. 再者，为了将责斥和审判的权能从圣父那里移除，他们以为这不配于天主，并认为他们发现了一位既无忿怒又[仅]是良善的天主，便声称一位[天主]审判，另一位拯救，不知不觉中剥夺了两位神祇的智慧与正义。因为如果那位\OriginalTerm{审判者天主}{judicial one}不是同样良善，不能施恩于应得者，也不能向需要者赐予斥责，他就不显为公正或智慧的审判者。另一方面，那\OriginalTerm{良善的天主}{good God}，如果祂仅是良善，而不考验那些祂将施以良善的人，祂就超出了正义和良善的范围；祂的良善也将显得不完全，如同不拯救所有人；[因为若没有审判相随，]本应如此。\par

3. 因此，马尔强自己将天主分为两位，主张一位是良善的，另一位是审判的，这实际从两方面都取消了神性。因为那位审判者，若不善，就不是天主，因为良善缺失者根本不是天主；同样，那位良善者，若没有审判权能，也遭受与前一位相同的[损失]，被剥夺其神性特征。如果他们认为圣父是明智的，怎能不将\OriginalTerm{审判能力}{judicial faculty}归于祂？因为如果祂是明智的，祂也是考验[他人]者；但审判权能属于考验者，公义随从审判能力，以达至公正的结论；公义引出审判，而审判在公义中施行时，将通向智慧。因此，圣父的智慧将超越一切人类和天使的智慧，因为祂是主、审判者、\OriginalTerm{公义者}{Just One}，并掌管万有。因为祂是良善、慈悲、忍耐的，拯救祂所当拯救的；在施行公义时，良善并未离开祂，祂的智慧也未减少；因为祂拯救祂所应拯救的，审判那些应受审判的。祂也从不显出不慈悲的公义；因为祂的良善无疑是先行并占优先的。\par

4. 因此，那位慈善地使祂的太阳升起照临众人\BibleRef{玛 5:45}，降雨给义人和不义的人的天主，要审判那些享受祂平均分施的恩慈却生活不配祂慷慨之尊严的人；他们虚度光阴于放荡和奢靡之中，违背祂的慈善，甚至还亵渎了那位施予他们如此巨大恩惠的天主。\par

5. 柏拉图被证明比这些人更虔诚，因为他承认同一天主既是公义又是良善的，掌管万有，亲自施行审判，这样说道：\Quote{诚然，天主，作为祂也是\OriginalTerm{远古的圣言}{ancient Word}，拥有万物的开端、终末和中际，按着万物的本性环绕万物而行，公正地运行一切；但报应的正义总是跟随祂，对付那些违犯神圣律法的人。}接着，他又指出宇宙的\OriginalTerm{造物主和塑造者}{Maker and Framer}是良善的。他说：\Quote{对于良善者，}他说，\Quote{任何事物都不会生出嫉妒。}从而确立了天主的良善，作为世界创造的开端和原因，而非无知，亦非迷途的\OriginalTerm{永世}{Æon}，亦非缺陷的后果，亦非哭泣哀叹的\OriginalTerm{母亲}{Mother}，亦非另一位神明或圣父。\par

6. 他们的\OriginalTerm{母亲}{Mother}哀哭他们，实在是应该的，因为他们竟能构思并发明这类事情，可说是他们合情合理地宣告了那反对自己的谎言：他们的母亲超出\OriginalTerm{圆满}{Pleroma}，即超出对天主的认识，且他们整个群众成了无形且粗劣的流产：因为它不解真理；它落入空虚和黑暗：因为他们的\OriginalTerm{智慧}{Sophia}是空虚的，并被黑暗包裹；\OriginalTerm{界限}{Horos}不准她进入圆满；因为那神（\OriginalTerm{亚卡莫特}{Achamoth}）没有接纳他们进入安息之所。因为他们的父亲，通过生养无知，在他们身上造成了死亡的苦楚。我们并没有歪曲这些[观点]；但他们自己证实，自己教导，并以此夸耀，他们幻想关于他们母亲的一个崇高[奥迹]，说她是在没有父亲（即没有天主）的情况下受生的，是由女性生出女性，即是由错误生出败坏。\par

7. 我们的确祈祷，这些人不要留在自己所挖掘的坑中，而要脱离这样一位\OriginalTerm{母亲}{Mother}，远离\OriginalTerm{深渊}{Bythus}，离开虚空，放弃阴影；并愿他们归向天主的教会，合法地重生，并愿基督在他们内形成，使他们认识这宇宙的创造者与制造者，那唯一真实的天主和万有之主。我们为他们祈祷这些事，爱他们胜过他们似乎爱自己。因为我们的爱，既然是真诚的，若他们愿意接受，便为他们有益。它好比一种严厉的药剂，切除伤口上骄傲的腐肉；因为它终结他们的骄傲和自大。因此，我们不会厌倦，尽我们所能向他们伸出援手。除了已经陈述的以外，我已推迟到下一卷书引用主的话；如果通过基督本身的教导，能说服他们中的一些人，使他们放弃这样的错误，不再亵渎他们的创造者——祂既是唯一的天主，也是我们的主耶稣基督的父。阿们。\par

\SourceRecord{Translated by Alexander Roberts and William Rambaut.；Ante-Nicene Fathers；Vol. 1.；Edited by Alexander Roberts, James Donaldson, and A. Cleveland Coxe.；Buffalo, NY: Christian Literature Publishing Co.,；1885.；<http://www.newadvent.org/fathers/0103325.htm>.}

% structure: p0001 source=h1 target=section reason=page-title

\section{\WorkTitle{驳斥异端（第四卷，序言）}}

1. 我至爱的朋友，通过将这部名为\WorkTitle{揭露与驳斥假知识}的第四卷传递给你，正如我所应许的，我将借助主的话来加重我已提出的论点；这样，你也能如你所请求的，从我这里获得驳斥各处所有异端者的方法，不让他们处处被击退后进一步陷入错误的深渊，也不让他们淹没在无知的海洋中；而是你引导他们进入真理的港湾，使他们获得救恩。\par

2. 然而，想要归化他们的人，必须准确了解他们的体系或学说方案。因为如果不知道病人的疾病，任何人都无法医治病人。这就是我的前辈们——远比我优秀的人——尽管未能满意地驳斥瓦伦廷派的原因，因为他们不了解这些人的体系；我已在第一卷中悉心向你传递了这一体系，并在其中表明他们的学说是一切异端的总汇。正因为如此，在第二卷中，我们如同在镜中看到了他们彻底溃败的景象。因为那些以正确方法反对这些人（瓦伦廷派）的人，实际上是在反对所有心怀恶意的人；而那些推翻他们的人，其实推翻了每一种异端。\par

3. 因为他们的体系是亵渎神明超过一切[其他异端]，因为他们声称那创造者与形成者，即独一天主，如我所表明的，是从一个缺陷或背道中产生的。他们还亵渎我们的主，将耶稣与基督分割开来，又将基督与救主分割开来，再将救主与圣言分割开来，又将圣言与\OriginalTerm{独生子}{Only-begotten}分割开来。由于他们断言创造者起源于缺陷或背道，他们也教导说基督和圣神是为了这个缺陷而发出的，并且救主是从那些由缺陷产生的\OriginalTerm{永世}{Æons}中出来的；因此在他们中间只有亵渎。在前一卷中，关于所有这些点的宗徒思想已经阐明，[表明]他们，即\BibleQuote{从起初就是圣言的眼见者和服务者}\BibleRef{路 1:2}（真理的），不仅没有持有这些观点，而且还教导我们避开这些学说，\BibleRef{弟后 2:23}，藉着圣神预见了那些软弱的、将要被引入歧途的人。\par

4. 因为正如蛇诱骗了厄娃，应许给她自己所没有的东西，\BibleRef{伯后 2:19}，同样这些人也假装[拥有]高超的知识，并[通晓]不可言喻的奥秘；并且应许他们所谈论的进入\OriginalTerm{圆满}{Pleroma}之内的接纳，使那些相信他们的人陷入死亡，使他们成为背弃创造他们者的叛徒。那时，确实，那背道的天使藉着蛇使人类不服从，以为逃脱了主的注意；因此天主赋予他[蛇]的形状和名称。但现在，末期已到，邪恶在人类中蔓延，不仅使他们成为背教者，而且[魔鬼]藉着各种诡计，即藉着所有上述的异端者，激起对创造者的亵渎。因为所有这些，尽管从不同的地域出来，宣扬不同的[意见]，却都同意同一个亵渎的意图，伤人至死，通过教导亵渎我们的创造者和维系者天主，并贬低人的救恩。人是由灵魂和肉体混合而成的组织，是按天主的肖像造成的，并由祂的手所塑造，即由圣子和圣神，祂对他们说：\BibleQuote{让我们造人。}\BibleRef{创 1:26}。因此，那嫉妒我们生命者的目的，就是使人们不信自己的救恩，并亵渎天主创造者。因为所有异端者可能极其庄严地提出的，最终归结为亵渎创造者，并否认天主作品的救恩，即肉体；为此，我已用多种方式证明天主圣子完成了整个[仁慈的]安排，并表明圣经中称为天主的，除了万有之父、圣子以及那些拥有\OriginalTerm{义子身份}{adoption}的人之外，别无其他。\par

\SourceRecord{Translated by Alexander Roberts and William Rambaut.；Ante-Nicene Fathers；Vol. 1.；Edited by Alexander Roberts, James Donaldson, and A. Cleveland Coxe.；Buffalo, NY: Christian Literature Publishing Co.,；1885.；<http://www.newadvent.org/fathers/0103400.htm>.}

% structure: p0001 source=h1 target=section reason=page-title

\section{驳斥异端（第四卷，第一章）}

1. 既然这事是确实而坚定的，即除了那作为天主统管万有的、连同祂的圣言，以及那些领受义子的圣神（即相信唯一真天主和祂的子耶稣基督的人）的那一位之外，圣神并未宣告任何别的天主或主；同样，宗徒们自己也不称任何别人为天主，或名为［别的］主；而且，更为重要的是，［既然这是真的］我们的主［也如此行］，祂曾命令我们只承认在天上的那一位为圣父，祂是唯一的天主和唯一的圣父——那么，这些欺骗者和最邪恶的诡辩家所提出的主张显然就是虚假的，他们坚持说，他们自己所虚构的那一位按其本性既是天主也是父；而巨匠造物主按其本性既不是天主也不是父，只是出于礼貌（\OriginalTerm{字面上}{verbo tenus}）才这样称呼，因为祂统治受造物，这些邪恶的神话编造者如此说，以他们的思想反对天主；并且，他们抛弃基督的教训，自己臆造谎言，反对天主的整个救恩计划。因为他们坚持说，他们的永世、神祇、父和主，还被进一步称为诸天，连同他们的母亲，他们也称她为\Quote{地}和\Quote{耶路撒冷}，还给她取了其他许多名字。\par

2. 现在，有谁不清楚，如果主曾知道许多父和神，祂就不会教导祂的门徒认识［唯一］天主\BibleRef{若 17:3}并只呼祂为圣父了？然而，祂宁可把那些仅在口头上（\OriginalTerm{字面上}{verbo tenus}）被称为神的，与那真正为天主的区分开来，使他们不致在祂的教训上犯错，也不致将一位误认为另一位。如果祂确实教导我们称一位为圣父和天主，而祂自己有时又以同样的意义承认别的父和神，那么祂就会显得是在给门徒规定一条与祂自己所行不同的道路。但这样的行为并不表明是一位好老师，而是一位误导和嫉妒的老师。按照这些人的说法，宗徒们也被证明违犯了诫命，因为他们承认造物主为天主、主和圣父，正如我已指出的——如果祂不是唯一的天主和圣父。那么，耶稣对于他们来说，就会是这种违犯的创始者和老师，因为祂命令只称一位为圣父\BibleRef{玛 23:9}，从而迫使他们不得不承认造物主为他们的圣父，正如已经指出的。\par

\SourceRecord{Translated by Alexander Roberts and William Rambaut.；Ante-Nicene Fathers；Vol. 1.；Edited by Alexander Roberts, James Donaldson, and A. Cleveland Coxe.；Buffalo, NY: Christian Literature Publishing Co.,；1885.；<http://www.newadvent.org/fathers/0103401.htm>.}

% structure: p0001 source=h1 target=section reason=page-title

\section{《驳斥异端》（第四卷，第二章）}

1. 梅瑟在《申命纪》中重述他从\OriginalTerm{造物主}{Demiurge}所领受的全部法律时说：\BibleQuote{诸天，请侧耳，我要发言；大地，请听我口中的言语。}\BibleRef{申 32:1} 达味又说他的援助来自上主，断言：\Quote{我的援助来自上主，他造了天地。}依撒意亚也承认那创造天地并治理万有的天主发言了。他说：\BibleQuote{诸天，请听！大地，请侧耳！因为上主说了话。}\BibleRef{依 1:2} 又说：\BibleQuote{那创造诸天并展开它们，奠定大地及其上的万物，赐气息给地上的万民，赐生命给行走其上者的上主天主如此说：}\BibleRef{依 42:5}\par

2. 再者，我们的主耶稣基督承认这同一位是他的父，他说：\BibleQuote{父啊，天地的主宰，我称谢你。}\BibleRef{玛 11:25} 那些极其乖谬的潘多拉诡辩家要我们理解哪一位父呢？是他们凭幻想虚构的\OriginalTerm{深渊}{Bythus}，还是他们的母亲，或是独生子？或是马西翁派及其他派别所捏造的神（我早已充分证明那根本不是神）？还是实际上那位天地的创造者，先知们也曾宣告过的，基督也承认他为父，法律也宣布说：\BibleQuote{以色列，请听！上主你的天主是唯一的天主。}\BibleRef{申 6:4}\par

3. 但既然梅瑟的\Emph{著作}是基督的话，他自己就向犹太人声明，正如若望在福音中所记载：\BibleQuote{如果你们相信梅瑟，就会相信我，因为他写的是关于我的。但如果你们不相信他的著作，也不会相信我的话。}\BibleRef{若 5:46} 他以最清楚的方式指明梅瑟的著作就是他的话。那么，如果梅瑟如此，其他先知的话毫无疑问也是他的话，正如我已指出的。再者，主亲自显示亚巴郎对那富人说，关于所有还活着的人：\BibleQuote{如果他们不听从梅瑟和先知，即使有人从死者中复活去他们那里，他们也不会相信。}\BibleRef{路 16:31}\par

4. 现在，他不仅给我们讲了一个关于穷人和富人的故事，而且教训我们：首先，没有人应当过奢侈的生活，也不要沉溺于世俗享乐和不断的宴乐，成为自己情欲的奴隶而忘记天主。他说：\BibleQuote{有一个富人，身穿紫红袍和细麻衣，日日奢华宴乐。}\BibleRef{路 16:19}\par

5. 关于这样的人，圣神也曾藉依撒意亚说过：\BibleQuote{他们伴着竖琴、手鼓、瑟和笛饮酒，却不留意天主的作为，也不顾念他手所作的。}\BibleRef{依 5:12} 因此，为避免我们遭受与这些人相同的惩罚，主向我们揭示了他们的结局；同时表明，如果他们听从梅瑟和先知，就会相信他们所宣讲的那一位，即从死者中复活、赐给我们生命的天主子；他又显示所有的人都来自同一本质，即亚巴郎、梅瑟、先知，以及主自己——他从死者中复活，许多受割礼的人都信他，他们也听从梅瑟和先知预言天主子的来临。但那些嘲笑真理的人断言这些人出自另一本质，他们不认识那从死者中首生者，把基督理解为一位不同的存在，似乎他是不能受苦的，而受难的耶稣则完全分离于他。\par

6. 因为他们不从父领受关于子的知识，也不从子学习父是谁，子清楚地、不用比喻地教导那真实的天主。他说：\BibleQuote{你们不可发誓；不可指天发誓，因为天是天主的宝座；不可指地发誓，因为地是他的脚凳；也不可指耶路撒冷发誓，因为它是大王的城。}\BibleRef{玛 5:34} 这些话显然是针对创造者说的，正如依撒意亚所说：\BibleQuote{天是我的宝座，地是我的脚凳。}\BibleRef{依 66:1} 除他之外没有别的天主；否则主就不会称他为\Quote{天主}或\Quote{大王}；因为能被这样描述的存在，不允许有其他存在比他或在他之上。因为任何在他之上者、在他权下者，绝不能被称为\Quote{天主}或\Quote{大王}。\par

7. 但这些人也无法坚持认为这些话是以讽刺的方式说出的，因为从这些话本身就可以向他们证明，它们是认真的。因为说这些话的那一位是真理，祂也确实为自己的房屋伸张正义，将那些买卖兑换银钱的人从屋里赶出去，对他们说：\Quote{经上记载：我的房屋应称为祈祷之所，但你们竟把它变成了贼窝。}\BibleRef{玛 21:13} 祂如此行、如此说，为祂的房屋伸张正义，如果祂宣讲的是另一位天主，那有什么理由呢？但祂这样做，是为了指出祂父法律的违犯者；因为祂既没有控告那房屋，也没有指责那法律——祂来正是为成全法律——而是责备那些滥用祂房屋的人和那些违犯法律的人。因此，从法律时期就开始轻视天主的经师和法利塞人，没有接受祂的话，也就是说，他们不信基督。关于这些人，依撒意亚说：\Quote{你的首领是叛逆的，与盗贼为伍，喜爱贿赂，追逐报酬，不审判孤儿，不顾寡妇的冤情。}\BibleRef{依 1:23} 耶肋米亚也同样说：\Quote{他们，}他说，\Quote{治理我子民的人不认识我；他们是愚昧无知的儿女；他们善于作恶，却不知行善。}\BibleRef{耶 4:22}\par

8. 但凡敬畏天主、关心祂法律的人，都奔向基督，并都得救了。因为祂对祂的门徒说：\Quote{你们往以色列家迷失的羊那里去。}\BibleRef{玛 10:6} 据说，当主在撒玛黎雅人中停留了两天后，更多的撒玛黎雅人\Quote{因祂的话而相信了，并对那妇人说：现在我们信，不是因为你的话，而是因为我们亲自听见了祂，并且知道这确实是世界的救主。}\BibleRef{若 4:41} 同样，保禄也宣告：\Quote{这样，全以色列必得救；}\BibleRef{罗 11:26} 但他也说过，法律是我们的启蒙师，引领我们归向基督耶稣。\BibleRef{迦 3:24} 因此，他们不要把某些人不信归咎于法律。因为法律从未阻止他们相信天主之子；相反，它还劝勉他们如此行，\BibleRef{户 21:8} 说人若不信那一位——祂以罪恶肉身的形状，被高举离开地面到殉道之木上，吸引万有归向自己，\BibleRef{若 12:32} 并使死者复生——就无法从蛇的旧伤中得救。\par

\SourceRecord{Translated by Alexander Roberts and William Rambaut.；Ante-Nicene Fathers；Vol. 1.；Edited by Alexander Roberts, James Donaldson, and A. Cleveland Coxe.；Buffalo, NY: Christian Literature Publishing Co.,；1885.；<http://www.newadvent.org/fathers/0103402.htm>.}

% structure: p0001 source=h1 target=section reason=page-title

\section{驳斥异端（第四卷，第三章）}

1. 此外，对于他们恶意断言：如果天确实是天主的宝座，地是他的脚凳，并且经上宣告天地将要消逝，那么当这些消逝时，坐在上面的天主也必然消逝，因此他不能是统御万有的天主；首先，他们不明白这句话的意思，即天是他的宝座，地是他的脚凳。因为他们不知道天主是什么，却想象他像人一样坐着，被界限所包容，而并不包容。他们也不明白天地消逝的含义；但\Person{保禄}对此并非无知，他宣告说：\Quote{因为这世界的局面正在逝去。}\BibleRef{格前 7:31}其次，\Person{达味}解释了他们的疑问，他说当这世界的局面逝去时，不仅天主将存留，他的仆人也一样，他在第一〇一篇圣咏中这样表达：\Quote{起初，上主，你奠定了大地，诸天是你手的化工。它们将要毁灭，你却永存；万物都要像衣服一样破旧；你更换它们如同更换衣裳，它们就被更换。但是你却永远不变，你的岁月没有穷尽。你仆人的子孙将长存，他们的后裔在你面前永远坚立。}这就清楚地指明了哪些事物会消逝，以及谁永远存留——天主，以及他的仆人。同样，\Person{依撒意亚}说：\Quote{你们要仰观诸天，俯察大地，因为天如同烟雾消散，地像衣服破旧，地上的居民也要同样死亡。但我的救恩却永远长存，我的正义永不消灭。}\BibleRef{依 51:6}\par

\SourceRecord{Translated by Alexander Roberts and William Rambaut.；Ante-Nicene Fathers；Vol. 1.；Edited by Alexander Roberts, James Donaldson, and A. Cleveland Coxe.；Buffalo, NY: Christian Literature Publishing Co.,；1885.；<http://www.newadvent.org/fathers/0103403.htm>.}

% structure: p0001 source=h1 target=section reason=page-title

\section{驳斥异端（第四卷，第四章）}

1. 此外，关于耶路撒冷和主，他们竟敢断言，如果耶路撒冷曾是\BibleQuote{伟大君王的城，}\BibleRef{玛 5:35}就不会被遗弃。这就如同有人说，如果稻草是天主的受造物，就绝不会与麦子分离；葡萄藤的枝条，如果是天主所造的，就绝不会被剪掉而失去果实。但是，这些枝条原本不是为自身而造，而是为了它们上面的果实；这果实一旦成熟并被摘取，枝条就被留下；那些不结果实的枝条则被完全剪掉。同样，耶路撒冷也是如此——她曾背负着奴役的轭（人曾被置于这轭下，从前死亡作王时，人不服从天主，被征服后，才适于获得自由）；当自由的果实来临、成熟、被收割并储存在仓里，当那些有能力结果实的已从她（耶路撒冷）那里被带走，并分散到全世界时，正如\Person{依撒意亚}所说，\BibleQuote{雅各伯的子孙必将生根，以色列必将开花，全世界必充满他的果实。}\BibleRef{依 27:6}因此，当果实已播种到全世界后，她（耶路撒冷）理应被遗弃，那些从前丰结果实的东西也被带走；因为从这些，按照肉身，基督和宗徒得以结果实。但现在这些东西对结果实不再有用。因为一切在时间中有开端的事物，当然也必须在时间中有终结。\par

2. 那么，既然法律始于\Person{梅瑟}，它必然终于\Person{若翰}。基督来是要成全法律，因此\BibleQuote{法律和先知}与他们同在\BibleQuote{直到若翰。}\BibleRef{路 16:16}所以，耶路撒冷始于\Person{达味}，满了它的时候，当新盟约被启示出来时，必须有立法的终结。因为天主以度量和秩序行一切事；祂没有什么是不度量的，因为没有什么是不合秩序的。那说不可测量的圣父在圣子内受到度量的人说得很好；因为圣子是圣父的尺度，因为圣子也包含圣父。但关于他们（犹太人）的管理是暂时的，\Person{依撒意亚}说：\BibleQuote{熙雍的女儿必剩下，像葡萄园中的草棚，像瓜田中的茅舍。}\BibleRef{依 1:8}什么时候这些事被遗弃？岂不是当果实被取走，只剩叶子，而叶子不再有能力结果实的时候吗？\par

3. 但我们何必提耶路撒冷？既然全世界的样式也必过去，当它消失的时刻来到之时，为的是让果实被收进仓里，而剩下的糠秕被火烧尽？\BibleQuote{因为主的日子来临，如同燃烧的火炉，所有罪人，即那些作恶的人，都是碎秸，那日子要烧毁他们。}\BibleRef{拉 4:1}谁是那带来这样的一天的主呢？\Person{若翰洗者}指出，论到基督时说：\BibleQuote{祂要用圣神和火给你们施行洗礼；祂手里拿着簸箕，要扬净祂的禾场，将麦粒收在仓里，将糠秕用不灭的火烧尽。}\BibleRef{玛 3:11}因为使糠秕和使麦子的不是不同的人，而是同一位，祂审判他们，即把他们分开。但麦子和糠秕，无生命无理智，按其本性被造成这样。但人，被赋予了理智，在这方面相似天主，具有自由意志，自己掌权，就是他自己有时成为麦子、有时成为糠秕的原因。因此，他也应受公正的定罪，因为他虽被造为理性存在，却失去了真正的理性，无理地生活，反对天主的正义，将自己交付给每个属地的灵，服事一切情欲；正如先知所说：\BibleQuote{人纵然受尊崇，却不能领悟；他竟与无理性的牲畜相似，变成与它们一样。}\par

\SourceRecord{Translated by Alexander Roberts and William Rambaut.；Ante-Nicene Fathers；Vol. 1.；Edited by Alexander Roberts, James Donaldson, and A. Cleveland Coxe.；Buffalo, NY: Christian Literature Publishing Co.,；1885.；<http://www.newadvent.org/fathers/0103404.htm>.}

% structure: p0001 source=h1 target=section reason=page-title

\section{驳斥异端（第四卷，第五章）}

1. 天主因此是唯一且同一的，祂卷起诸天如同书卷，更新地面；祂为人创造了时间性的事物，使人得以在其中成熟，结出不朽的果实；并因祂的良善，也将永恒的事物赐予［人］，\BibleQuote{为在未来的世代，显示祂恩宠的超越丰盛。}\BibleRef{厄 2:7} 祂为法律与先知所宣告，基督宣认祂为自己的圣父。如今祂是创造者，祂就是统御万有的天主，正如依撒意亚所说：\BibleQuote{我是见证——上主天主说——也是我所拣选的仆人，为使你们知道、相信并明白我就是那一位。在我以前没有别的神，在我以后也不会再有。我是天主，除我之外没有救主。我已宣告，我已拯救。}\BibleRef{依 43:10} 又说：\BibleQuote{我自身是首先的天主，我也凌驾于将来的事之上。}\BibleRef{依 12:4} 因为祂说这些话，并非以模糊、傲慢或夸耀的方式，而是既然没有天主就无法认识天主，祂便藉着自己的圣言教导人认识天主。因此，对于那些对此无知、并因此以为发现了另一位父的人，有人公正地说：\BibleQuote{你们错了，既不明白圣经，也不明了天主的能力。}\BibleRef{玛 22:29}\par

2. 因为我们的主与师傅，在回答撒杜塞人——他们说不存在复活，并因此羞辱天主，贬低法律的威信——时，既指示了复活，又启示了天主，对他们说：\Quote{你们错了，既不明白圣经，也不明了天主的能力。} \Quote{关于死人复活，}祂说，\Quote{你们没有读过天主对你们所说的话吗？他说：我是亚巴郎的天主，依撒格的天主，雅各伯的天主。}祂又加上：\Quote{祂不是死人的天主，而是活人的天主；因为众人都为祂而活。}藉这些论证，祂无疑清楚地表明，那位曾在荆棘中向梅瑟说话、并宣称自己为祖先的天主者，就是活人的天主。因为谁会是活人的天主，除非那身为天主、且在他之上别无其他天主者？达尼尔先知也——当波斯王居鲁士对他说：\Quote{你为何不朝拜贝耳？}——宣告说：\Quote{因为我不朝拜人手所造的偶像，而是朝拜生活的天主，祂建立了天与地，统御所有血肉。}他又说：\Quote{我要朝拜上主我的天主，因为祂是生活的天主。}因此，那被先知们朝拜为生活天主的那一位，就是活人的天主；而祂的圣言，就是那位也曾向梅瑟说话、也曾使撒杜塞人无言、也曾赐予复活恩赐的那一位，从而向那些盲目的人启示了双重真理，即复活和天主［的真实特性］。因为如果祂不是死人的天主，而是活人的天主，却仍被称为那安眠祖先的天主，那么他们无疑是为天主而活着，并未归于无有，因为他们是复活之子。但我们的主自己就是复活，正如祂自己所宣告的：\BibleQuote{我是复活和生命。}\BibleRef{若 11:25} 但祖先们是祂的子女；因为先知说过：\Quote{代替你的祖先，你的子女被赐给你。}因此，基督自己与圣父一起，就是活人的天主，祂曾向梅瑟说话，也曾显现给祖先。\par

3. 教导这同一真理时，祂对犹太人说：\BibleQuote{你们的祖先亚巴郎曾欢欣地希望看见我的日子，他看见了，且欢喜快乐。}\BibleRef{若 8:56} 这意味着什么？\BibleQuote{亚巴郎相信了天主，这就算为他的正义。}\BibleRef{罗 4:3} 首先，他相信祂是天地万物的创造者，唯一的天主；其次，相信祂要使他的后裔像天上的星辰。这就是保禄所说的意思：\BibleQuote{如同世上的光。}\BibleRef{斐 2:15} 因此，他正义地撇下了地上的亲族，跟随了天主的圣言，作为朝圣者与圣言同行，以便日后能与圣言同住。\par

4. 正义地，宗徒们——亚巴郎的后裔——也撇下了船和父亲，跟随了圣言。正义地，我们这些拥有与亚巴郎同样信德的人，也背起十字架，如同依撒格背起木柴一样，\BibleRef{创 22:6}跟随祂。因为在亚巴郎身上，人预先学习并习惯了跟随天主的圣言。因为亚巴郎凭着他的信德，遵行了天主圣言的命令，以乐意的心将他的独生圣子献上作为祭献，为的是天主也乐意为了他所有的后裔，献上自己的独生圣子作为我们救赎的祭献。\par

5. 因此，既然亚巴郎是一位先知，并在圣神中看见了主来临的日子，以及祂受苦的时期，藉着祂，他自己以及所有以他的信德为榜样、信赖天主的人，都必得救，他就极其欢喜。因此，主对亚巴郎并非未知，亚巴郎曾渴望看见祂的日子；\BibleRef{若 8:56} 同样，主的圣父也非未知，因为他从主的圣言学习并相信了祂；因此，这为他被主算为正义。因为对天主的信德使人成义；因此他说：\BibleQuote{我要向至高者天主——创造天地的天主——举手起誓。}\BibleRef{创 14:22} 然而，所有持歪曲观点的人，都试图推翻所有这些真理，因为他们错误地理解了一处经文。\par

\SourceRecord{Translated by Alexander Roberts and William Rambaut.；Ante-Nicene Fathers；Vol. 1.；Edited by Alexander Roberts, James Donaldson, and A. Cleveland Coxe.；Buffalo, NY: Christian Literature Publishing Co.,；1885.；<http://www.newadvent.org/fathers/0103405.htm>.}

% structure: p0001 source=h1 target=section reason=page-title

\section{\WorkTitle{驳斥异端}（第四卷，第六章）}

1. 因为主向祂的门徒启示，表明祂自己就是那赐予圣父知识的圣言，并谴责那些自以为认识天主、却拒绝祂的圣言（天主藉着圣言被认识）的犹太人，祂说：\BibleQuote{除了圣父外，没有人认识圣子；除了圣子及圣子所愿意启示的人外，也没有人认识圣父。}\BibleRef{玛 11:27}玛窦是这样记载的，路加类似，马尔谷也一样；若望省略了这段经文。然而，那些自以为比宗徒更聪明的人，却如此书写[这节经文]：\Quote{没有人\Emph{认识}圣父，除了圣子；也没有人认识圣子，除了圣父，以及圣子所愿意启示的人。}他们解释这话，仿佛在我主来临之前，无人认识真天主；他们断言那由先知所预言的天主并非基督的圣父。\par

2. 但如果基督只是在祂作为人来到[世界]时才开始存在，而圣父在提庇留·凯撒的时代才记起要供应[人所需]，并且祂的圣言被证明并非总是与祂的受造物共存，[可以指出]那么，也无必要宣告另一位天主，反而[更应当]调查祂如此大意和疏忽的原因。因为不应该让这样的问题产生并积聚力量，以致实际上既改变了天主，又摧毁了我们对那位藉着受造物支持我们的造物主的信德。因为正如我们将信德指向圣子，我们也应当拥有对圣父坚定不移的爱德。在他反对马尔强的书中，犹斯定说得很好：\Quote{我甚至不会相信主自己，如果祂宣告了另一位而不是那一位——祂是我们的塑造者、造作者和养育者。但因为独生子从独一天主来到我们这里，祂既创造了这个世界，也塑造了我们，并包含和管理万有，将祂自己的手工作品总归在祂自己内，我对祂的信德是坚定的，我对圣父的爱德是不可动摇的，天主将两者赐予我们。}\par

3. 因为除非藉着天主的圣言，即除非藉着圣子的启示，没有人能认识圣父；同样，除非藉着圣父的善意，也没有人能认识圣子。但圣子实行圣父的善意；因为圣父派遣，圣子被派遣，并来了。祂的圣言知道，祂的圣父就我们而言是不可见且无限的；既然祂不能被[任何人]宣告，祂亲自向我们宣告祂；另一方面，只有圣父认识祂自己的圣言。我们的主宣布了这两个真理。因此，圣子藉着祂自己的显现揭示对圣父的认识。因为圣子的显现就是对圣父的认识；因为万有都是藉着圣言而显现的。因此，为了让我们知道，那来的圣子就是那赐予信祂之人对圣父认识的这一位，祂对门徒说：\BibleQuote{除了圣父外，没有人认识圣子；除了圣子和圣子所愿意启示的人外，也没有人认识圣父；}\BibleRef{玛 11:27}从而将自己和圣父呈现为祂[实际]所是，使我们不接受任何别的圣父，除了圣子所启示的那一位。\par

4. 但这位[圣父]是天地的创造者，正如祂的言语所示；并不是那假父，即由马尔强、或华伦提努、或巴西理得、或卡珀克拉特斯、或西门，以及其余所谓\Quote{诺斯替派}的人所捏造的。这些人中没有一个是天主子；而是我们的主基督耶稣[是]，他们将自己的教导与之对立，并胆敢宣扬一位未知的天主。但他们应当听到这话来反对自己：那位被他们认识的天主，怎么可能是未知的？因为凡是甚至被少数人所认识的，就不是未知的。但主并没有说圣父和圣子\OriginalTerm{完全}{in totum}不能认识，因为那样的话，祂的来临就是多余的。祂为何来到这里？难道是要对我们说：\Quote{不要寻求天主；因为祂是未知的，你们不会找到祂；}正如华伦提努的门徒错误地宣称基督曾对他们的\OriginalTerm{永世}{Æons}所说的话？但这确实是徒劳的。因为主教导我们，除非被天主教导，没有人能够认识天主；也就是说，没有天主就不能认识天主；但这是圣父明确的旨意，即天主应当被认识。因为凡圣子所启示的人，他们将认识祂。\par

5. 圣父为此目的启示了圣子，好藉着祂，祂能向所有人显示自己，并接纳那些相信祂的义人进入不朽和永久的福乐（信祂就是承行祂的旨意）；但祂将会公义地将那些不相信、因而回避祂的光的人，关入他们为自己选择的黑暗之中。因此，圣父藉着使祂的圣言为众人可见，而将自己启示给众人；反之，圣言也向众人宣告了圣父和圣子，因为祂已成为众人可见的。因此，天主的公义审判[将落在]所有那些像其他人一样看见、却不像其他人一样相信的人身上。\par

6. 因为藉着受造界本身，圣言启示了造物主天主；藉着世界，圣言宣告了主、世界的创造者；藉着人的形成，圣言宣告了形成他的工匠；并藉着圣子，启示了生圣子的圣父。这些事确实以相同的方式向所有人说话，但并非所有人都以相同的方式相信它们。然而藉着法律和先知，圣言同样向所有人宣讲了祂自己和圣父；全体人民都同样听见了祂，但并非所有人都同样相信。并且透过成了可见可触的圣言自己，圣父被显示出来，虽然并非所有人都同样相信祂；但所有人在子内看见了父：因为圣父是圣子的不可见者，圣子是圣父的可见者。正因如此，所有人当基督（在地上）显现时都与祂说话，并称祂为天主。是的，甚至魔鬼看见圣子时也呼喊说：\Quote{我们知道你是谁，你是天主的圣者。}\BibleRef{谷 1:24} 魔鬼注视着祂，试探祂说：\Quote{你若是天主子，}\BibleRef{玛 4:3}——所有这些人确实看见了并谈论圣子和圣父，但并非所有人都相信。\par

7. 因为真理应当从众人获得见证，并成为审判：为相信的人乃是救恩，但为不相信的人乃是定罪；好使众人受公平审判，并使对圣父和圣子的信德被所有人认可，也就是说，应被所有人确立为唯一的救恩途径，从众人领受见证，既从属于它的人（因为是它的朋友），也从与它无关的人（虽为它的仇敌）。因为那证据是真实的，无可辩驳，它甚至从敌手那里激起对其有利的惊人见证；他们因着对眼前之事的直接默观而被说服，为之作证并宣扬它。但过些时候，他们便转而敌视，成为（他们曾认可的）控告者，并希望他们自己的见证不被视为真实。因此，那位被认识者，与那位宣告\Quote{没有人认识父}的是同一位，并非不同的存在，而是同一的，圣父使万有服从于祂；而祂从众人领受见证，证明祂是真人，又是真天主：来自圣父、来自圣神、来自天使、来自受造界本身、来自人、来自背道的邪灵和魔鬼、来自仇敌，最后来自死亡本身。然而圣子为圣父管理万有，从起初直到终结工作，没有祂，没有人能获得对天主的认识。因为圣子是对圣父的认识，而对圣子的认识是在圣父内，并透过圣子启示了；这就是主为何宣称\Quote{除了父外，没有人认识子；除了子外，也没有人认识父，以及子所愿意启示的人。}因为\Quote{将启示}一词并非仅指未来，好像圣言在由玛利亚诞生时才开始启示圣父，而是普遍适用于所有时代。因为圣子从起初就临在于祂亲手所作的工程，向众人启示圣父；祂愿意启示谁，何时启示，并如同圣父所愿意的。因此，在万有中并透过万有，只有一位天主圣父，一位圣言，一位圣子，一位圣神，以及一个救恩，给所有相信祂的人。\par

\SourceRecord{Translated by Alexander Roberts and William Rambaut.；Ante-Nicene Fathers；Vol. 1.；Edited by Alexander Roberts, James Donaldson, and A. Cleveland Coxe.；Buffalo, NY: Christian Literature Publishing Co.,；1885.；<http://www.newadvent.org/fathers/0103406.htm>.}

% structure: p0001 source=h1 target=section reason=page-title

\section{驳斥异端（第四卷，第七章）}

1. 因此，亚巴郎也藉着那创造天地的\OriginalTerm{圣言}{\GreekText{Λόγος}}认识了圣父，并宣认祂为天主；又因着（对他所作的）宣告得知天主子将作为人而居于人之中，藉祂的来临，他的后裔将如天上星辰那样繁多，他便渴望见到那一日，好使他也能亲自拥抱基督；而藉着预言之神见到那日子，他就喜乐了。\BibleRef{创 17:17} 因此，他的后裔之一西默盎也完美地实现了族长的喜乐，说道：\BibleQuote{主啊！现在可照祢的话，让祢的仆人平安去吧。因为我亲眼看见了祢的救恩，即祢在万民面前所预备的：为启示外邦人的光明，和祢子民以色列的荣耀。} \BibleRef{路 2:29} 同样，天使们也向夜间守牧的牧羊人报告了大喜的信息。\BibleRef{路 2:8} 此外，玛利亚说：\BibleQuote{我的灵魂颂扬主，我的心神欢跃于天主，我的救主。} \BibleRef{路 1:46} ——亚巴郎的喜乐降临到他的子孙身上，即那些警醒、看见基督并信祂的人；另一方面，有一种互惠的喜乐从子女回溯到亚巴郎，他也曾渴望看见基督来临的日子。因此，我们的主正确地为他作证说：\BibleQuote{你们的父亲亚巴郎曾欢欣仰望我的日子，他看见了，且喜乐了。}\par

2. 因为祂说这些话不仅是为亚巴郎的缘故，也是要指出，凡从起初认识天主并预言基督来临的人，都是从圣子本身领受了启示；圣子在末期成为可见的、可受苦的，并与人交谈，为要从石头中给亚巴郎兴起子孙，并实现天主应许给他的诺言，使他的后裔如同天上的星辰一样繁多，\BibleRef{创 15:5} 正如若翰洗者所说：\BibleQuote{天主能从这些石头给亚巴郎兴起子孙来。} \BibleRef{玛 3:9} 如今，耶稣正是这样做的：祂使我们离开对石头的崇拜，从坚硬且不结果实的思虑中引领我们出来，并使我们建立如同亚巴郎一样的信德。正如保禄也作证说，我们因信德的相似和继承的应许而成为亚巴郎的子女。\BibleRef{罗 4:12}\par

3. 因此，呼召亚巴郎并给他应许的，是同一、唯一的天主。那创造者，也就是通过基督在世界上预备光明者（即从外邦人中信从的人）。祂说：\BibleQuote{你们是世界的光；} \BibleRef{玛 5:14} 也就是说，如同天上的星辰。因此，我正确地指出了，祂不被任何人认识，除非藉着圣子，并藉着圣子所愿意启示的人。但圣子向祂愿意使祂被认识的一切人启示圣父；没有圣父的美意，也没有圣子的工作，无人能认识天主。因此，主对祂的门徒说：\BibleQuote{我是道路、真理、生命；除非经过我，谁也不能到父那里去。你们若认识我，也就必然认识我父；现在你们已认识祂，并且已看见祂。} \BibleRef{若 14:6} 从这些话明显可知，祂是因圣子，即\OriginalTerm{圣言}{\GreekText{Λόγος}}，而被认识的。\par

4. 因此，犹太人因不接受圣言，就离开了天主，他们妄想无需圣言（即圣子）而独自认识圣父；他们不知道那位以人形对亚巴郎说话的天主，\BibleRef{创 18:1} 又对梅瑟说：\BibleQuote{我看见我的百姓在埃及所受的痛苦，我下来要拯救他们。} \BibleRef{出 3:7} 因为那圣言，即天主的圣子，从起初就预先安排了这一切——圣父不需要天使来创造万物和塑造人，万物也是为了人而被造的；圣父也不需要任何工具来创造受造之物或安排与人有关的事；同时，祂拥有不可胜数的仆役。因为祂的\OriginalTerm{后裔}{offspring}和\OriginalTerm{肖像}{similitude}在各方面为祂服务，即圣子和圣神、圣言和\OriginalTerm{智慧}{\GreekText{Σοφία}}；所有天使都服事他们，并服从他们。因此，那些因那句\BibleQuote{除了圣子，没有人认识圣父} \BibleRef{玛 11:27} 而引入另一位未知的圣父的人，是徒劳的。\par

\SourceRecord{Translated by Alexander Roberts and William Rambaut.；Ante-Nicene Fathers；Vol. 1.；Edited by Alexander Roberts, James Donaldson, and A. Cleveland Coxe.；Buffalo, NY: Christian Literature Publishing Co.,；1885.；<http://www.newadvent.org/fathers/0103407.htm>.}

% structure: p0001 source=h1 target=section reason=page-title

\section{驳斥异端（卷四，第八章）}

1. 马尔奇安及其追随者徒然［试图］将亚巴郎排除在继承之外，圣神藉多人、如今又藉保禄为亚巴郎作证说：\Quote{亚巴郎信了天主，天主就以此算为他的正义。}\BibleRef{罗 4:3} 主［也为亚巴郎作证，］首先，从石头中给亚巴郎兴起子孙，并使他的后裔如同天上的星辰，说：\Quote{将有许多人从东方和西方，从北方和南方而来，与亚巴郎、依撒格和雅各伯在天国里一同坐席；}\BibleRef{玛 8:11} 然后又对犹太人说：\Quote{当你们看见亚巴郎、依撒格、雅各伯及众先知在天国里，而你们却被弃在外时。}\BibleRef{路 13:28} 因此，这一点是清楚的：那些否认亚巴郎的救恩，并构想另一位神、而非那向亚巴郎作出许诺的那一位的人，是在天国之外，被剥夺了不朽的［恩赐］，他们藐视并亵渎天主，而天主正是藉耶稣基督将亚巴郎及其后裔——即教会——领进天国，教会也获得了那应许给亚巴郎的义子和继承权。\par

2. 主藉着解救亚巴郎的后裔脱离束缚并召叫他们获得救恩，来维护他们，正如祂在治愈那妇人的事例中所行的，祂公开对没有亚巴郎那样信德的人说：\Quote{假善人！你们每个人在安息日，难道不解开自己的牛或驴，牵去饮水吗？这个女人原是亚巴郎的女儿，被撒殚捆绑了十八年，在安息日就不该解开她的束缚吗？}\BibleRef{路 13:15} 因此很明显，祂解救了并赋予生命给那些像亚巴郎一样信从祂的人，祂在安息日治病并没有作任何违反法律的事。因为法律并没有禁止人在安息日被治愈；［相反，］法律甚至在安息日给人行割损，并命令司祭为人民举行礼仪；是的，法律甚至不禁止医治哑巴牲畜。祂在史罗亚池和以后多次在安息日施行治愈；因此，许多人在安息日都到祂那里去。因为法律命令他们戒避一切劳碌的工作，即戒避通过贸易和其他世俗事务牟取财富的一切贪婪；但法律劝勉他们从事灵魂的操练，即反思，并进行有益于邻人的训诲。因此，主责备那些不义地指责祂在安息日治病的人。因为祂并未废除法律，而是成全了法律，履行了大司祭的职务，为人类向天主赎罪，并洁净癞病人，医治病人，并且亲自受死，使被放逐的人得以脱离定罪，并能无所畏惧地归回自己的产业。\par

3. 再者，法律并不禁止那些在安息日饥饿的人取用现成的食物：但它确实禁止他们收割和收进仓里。因此，主对那些指责祂的门徒掐了麦穗搓着吃的人说：\Quote{你们没有读过这事吗？达味在饥饿时做了什么？他进了天主的殿，吃了供饼，还给了同他一起的人；这饼只准司祭吃，别人不可吃。}\BibleRef{路 6:3} 祂用法律的话为祂的门徒辩护，指出司祭可以自由行动。因为达味虽然被撒乌耳迫害，却已被天主立为司祭。因为所有的义人都拥有司祭的品位。主的所有宗徒都是司祭，他们在此世不继承土地或房屋，而是不停地事奉天主和祭台。关于他们，梅瑟在\WorkTitle{申命纪}祝福肋未时也说：\Quote{他对自己的父亲和母亲说：我不认识你们；他也不承认自己的兄弟，他抛弃了自己的儿子：他遵守了你的诫命，持守了你的盟约。}\BibleRef{申 33:9} 但谁离开了父亲和母亲，并因天主的话和祂的盟约告别了所有的邻人，如果不是主的门徒，还有谁呢？关于他们，梅瑟又说：\Quote{他们不可有产业，因为上主自己就是他们的产业。}\BibleRef{户 18:20} 又说：\Quote{司祭肋未人不可在肋未全支派中，也不可在以色列人中拥有份或产业；他们的产业是上主的祭品（\OriginalTerm{果实}{fructifications}）：他们要以此为食。}\BibleRef{申 18:1} 因此保禄也说：\Quote{我不是寻求礼物，而是寻求果实。}\BibleRef{斐 4:17} 祂对祂的门徒——那些拥有上主司祭职的人，饥饿时吃麦穗是合法的——说：\Quote{工人自当得他的食物。}\BibleRef{玛 10:10} 而圣殿里的司祭违犯了安息日，却不算为有罪。那么，他们为什么不算为有罪呢？因为他们在圣殿里，从事的不是世俗事务，而是事奉上主，成全法律，却没有逾越法律，如同那人，他擅自将干柴带回天主的营中，因而被正义地处以石刑。\BibleRef{户 15:32} \Quote{凡不结好果子的树，都要砍倒，丢在火里；}\BibleRef{玛 3:10} 并且\Quote{谁若败坏天主的殿，天主必要败坏那人。}\BibleRef{格前 3:17}\par

\SourceRecord{Translated by Alexander Roberts and William Rambaut.；Ante-Nicene Fathers；Vol. 1.；Edited by Alexander Roberts, James Donaldson, and A. Cleveland Coxe.；Buffalo, NY: Christian Literature Publishing Co.,；1885.；<http://www.newadvent.org/fathers/0103408.htm>.}

% structure: p0001 source=h1 target=section reason=page-title

\section{\WorkTitle{驳斥异端}（第四卷，第九章）}

1. 因此，万物都来自同一实体，即来自同一天主；正如主对门徒所说：\BibleQuote{所以，凡成为天国门徒的经师，就像一个家主，从他的宝库里取出新的和旧的东西。}\BibleRef{玛 13:52} 祂并非教导说，取出旧的和取出的新的是不同的人，而是同一位。因为主是良善的家主，管理祂父的整个家；祂颁布适合奴隶和尚未受管教之人的法律；也向那些因信德而成义的自由人赐予合宜的诫命，并将自己的产业向儿子们敞开。祂称门徒为\Quote{经师}和\Quote{天国教师}；关于他们，祂又对犹太人说：\BibleQuote{看，我派遣智者、经师和教师给你们；其中有些你们要杀死，有些要迫害，从这城到那城。}\BibleRef{玛 23:34} 毫无疑问，祂所说从宝库里取出的新和旧，指的是两个盟约：旧的，是以前颁布的法律；新的，是福音所要求的生活方式，达味对此说：\Quote{向上主唱一首新歌。}依撒意亚说：\Quote{向上主唱一首新的赞美诗。祂的\OriginalTerm{开始}{initium}，祂的名从地极被颂扬：众海岛宣扬祂的大能。}耶肋米亚说：\BibleQuote{看，我将订立新约，不像我同你们祖先所订立的}\BibleRef{耶 31:31}在曷勒布山上。然而，同一位家主产生了两个盟约，就是天主圣言、我们的主耶稣基督，祂曾与亚巴郎和梅瑟说话，并重新赐给我们自由，增多了出自祂的恩宠。\par

2. 祂宣告：\BibleQuote{在这里有比圣殿更大的一位。}\BibleRef{玛 12:6} 但\Emph{更大}和\Emph{更小}不适用于那些彼此之间毫无共同之处、性质相反且相互排斥的事物；而是用在具有同一实体、拥有共同属性、仅在数量和大小上有所区别的事物上，例如水与水、光与光、恩宠与恩宠。因此，为自由而颁布的立法比为奴役而颁布的立法更伟大；因此它也不仅传播给一个民族，而是遍布全世界。因为同一位主，比圣殿更大，比撒罗满更大，比约纳更大，祂赐予人类恩赐，即祂自己的临在和死者复活；但祂并没有改变天主，也没有宣告另一位父，而是同一位，祂总是有更多的东西分赐给祂家中的人。随着他们对天主的爱德增长，祂赐予更多更大的恩赐；正如主对门徒所说：\BibleQuote{你要看见比这些更大的事。}\BibleRef{若 1:50} 保禄也宣称：\Quote{这不是说我已经达到成义，或已经成全。因为我们现在所知道的，只是局部的；我们作先知所讲的，也只是局部的；但当那圆满的一来到，局部的，就必要消逝。} 因此，当那圆满的一来到，我们不会看见另一位父，而是我们现在渴望看见的那一位（因为\BibleQuote{心里洁净的人是有福的，因为他们要看见天主}\BibleRef{玛 5:8}）；我们也不会寻求另一位基督和天主之子，而是那位由童贞玛利亚所生、也受苦的，我们所信靠和爱慕的；正如依撒意亚说：\BibleQuote{到那日，人要说：看，这是我们的天主，我们所信赖的，我们因祂的救恩而欢欣。}\BibleRef{依 25:9} 伯多禄在他的书信中说：\BibleQuote{你们虽然没有见过祂，却爱慕祂；虽然现在看不见祂，却相信祂，以不可言喻的喜乐欢欣。}\BibleRef{伯前 1:8} 我们也不会领受另一位圣神，而是那与我们同在、呼喊\Quote{阿爸，父啊！}的那一位；\BibleRef{罗 8:15} 我们将在同样的事物上增长并进步（如同现在一样），以致不再透过镜子或谜语，而是面对面地享受天主的恩赐——同样，现在既然接受了比圣殿和撒罗满更大的那一位，即天主圣子的来临，我们就没有被教导另一位天主，而是那位从起初就指明给我们的万有创造者和制造者；也不是另一位基督、天主之子，而是那位由先知所预言的那一位。\par

3. 因为新约既已被先知们认识和宣讲，那要按圣父的美意实行新约的一位也被宣讲了，按天主所喜悦的向人启示了；使他们因信祂而不断进步，并藉着（相继的）盟约，逐步达致完美的救恩。因为只有一个救恩、一位天主；但塑造人的规诫众多，引领人归向天主的步骤不少。一个属世暂时的君王，虽是人，尚可有时赐予臣民更大的恩惠：那么，天主既是永远同一，又总是愿意赐予人类更大（程度的）恩宠，并不断以许多恩赐尊荣那些中悦祂的人，难道这是不允许的吗？但如果进步就是这样：找出另一位父，不同于那从起初所传的；然后又找出第三位，不同于那第二位想象中发现的那位——那么这人的进步就在于他从第三位进行到第四位；再从第四位到另一位，又一位：这样，凡以为自己在总是作这种进步的人，将永不能安息于唯一天主。因为，从那真正是【天主】者被赶走，而后退，他将永远寻求，却永不得寻见天主；\BibleRef{弟后 3:7}却不断地在无底的深渊中漂浮，除非他藉悔改归正，返回他被逐出的地方，明认唯一天主、圣父、创造者，并相信那由法律和先知所宣告、由基督所见证的，正如祂亲自向那些控告祂的门徒不守长老传统的人所宣告的：\Quote{为什么你们因你们的传统而废弃天主的法律？因为天主说过：\InnerQuote{应孝敬你的父亲和母亲；}又说：\InnerQuote{咒骂父亲或母亲的，应处以死刑。}} \BibleRef{玛 15:3} 祂第二次又对他们说：\Quote{你们因你们的传统，废弃了天主的话。} 基督以最明确的方式承认那位曾在法律中说过：\Quote{应孝敬你的父亲和母亲，好使你事事顺遂。} 的，就是父和天主。因为真天主承认法律的诫命是天主的话，除了自己的父以外，没有称别人为天主。\par

\SourceRecord{Translated by Alexander Roberts and William Rambaut.；Ante-Nicene Fathers；Vol. 1.；Edited by Alexander Roberts, James Donaldson, and A. Cleveland Coxe.；Buffalo, NY: Christian Literature Publishing Co.,；1885.；<http://www.newadvent.org/fathers/0103409.htm>.}

% structure: p0001 source=h1 target=section reason=page-title

\section{\WorkTitle{驳异端}（第四卷，第十章）}

1. 因此，若望也恰当地记述了主对犹太人说：\BibleQuote{你们查考圣经，因为你们以为其中有永生；正是这些为我作证。你们却不愿意到我这里来，为获得生命。}\BibleRef{若 5:39} 那么，\WorkTitle{圣经}如何为祂作证呢？除非它们来自同一个圣父，预先教导人有关祂圣子的来临，并预言由祂带来的救恩。\BibleQuote{如果你们相信梅瑟，你们也必会相信我，因为他曾论及我写了。}\BibleRef{若 5:46} 这无疑是因为天主圣子处处贯穿于他的著作中。的确，有时祂与亚巴郎说话，正要与他一同进食；有时与诺厄说话，将方舟的尺寸赐给他；有时查问亚当；有时对索多玛人降下审判；又有一次，祂显现出来，指引雅各伯的旅程，并从荆棘丛中与梅瑟说话。\BibleRef{出 3:4} 若要一一数算梅瑟所显示的天主圣子的事例，那是说不尽的。关于祂受难的日期，梅瑟也并非不知道；相反，他以象征的方式，通过逾越节的名称预言了祂；我们的主正是在那同一节日——梅瑟在很久以前就预定了的节日——受难，从而应验了逾越节。他不仅描述了那日子，还有那地点，以及苦难结束时的时辰和日落的现象，说：\BibleQuote{你不可在你任何城里——那上主你的天主赐给你的城中——随便献逾越节祭；只应在上主你的天主所选定、为立祂名的地方，你才能在傍晚、日落时献逾越节祭。}\BibleRef{申 16:5}\par

2. 并且梅瑟早已预告了祂的来临，说：\BibleQuote{犹大的权杖不会失去，领袖不会从他的腰间消失，直到那应得的人来到，祂是万民的期望；祂将驴驹拴在葡萄树上，将母驴的驹子拴在蔓延的常春藤上。祂在酒中洗涤外袍，在葡萄血中洗涤衣裳；祂的眼睛比酒更欢愉，祂的牙齿比奶更洁白。}\BibleRef{创 49:10} 因为，让那些以研究一切闻名的人去查考，何时犹大失去了王子和领袖？谁是万民的期望？葡萄树是谁？那属于祂的驴驹是什么？衣服、眼睛、牙齿、酒各指什么？让他们逐一查考这些要点，他们就会发现所预告的正是我们的主基督耶稣，而非别人。因此，梅瑟在斥责百姓的忘恩负义时，说：\BibleQuote{愚昧无知的人民，你们就这样报上主吗？}\BibleRef{申 32:6} 他又指出，那位从起初就建立并创造他们的圣言，也在末期救赎并赐予我们生命者，被展示为挂在树上，他们却不相信祂。因为他（梅瑟）说：\BibleQuote{你的性命必悬在你眼前，你却不信你的性命。}\BibleRef{申 28:66} 又说：\BibleQuote{这同一一位，难道不是你的父亲拥有你、造了你、创造了你吗？}\BibleRef{申 32:6}\par

\SourceRecord{Translated by Alexander Roberts and William Rambaut.；Ante-Nicene Fathers；Vol. 1.；Edited by Alexander Roberts, James Donaldson, and A. Cleveland Coxe.；Buffalo, NY: Christian Literature Publishing Co.,；1885.；<http://www.newadvent.org/fathers/0103410.htm>.}

% structure: p0001 source=h1 target=section reason=page-title

\section{驳斥异端（卷四，第十一章）}

1. 但是，主已显明，不仅是先知和许多义人，藉着圣神预见祂的来临，祈求他们能到达那个时期，得以面对面看见他们的主，并听见祂的话语——当祂对祂的门徒说：\BibleQuote{许多先知和义人曾渴望看见你们所看见的，却没有看见；并渴望听见你们所听见的，却没有听见。}\BibleRef{玛 13:17} 那么，他们如何渴望听见和看见呢？除非他们预先知道祂将来的来临。但他们又怎能预先知道呢？除非他们先前从祂自己那里得到了预知。圣经又如何为祂作证呢？除非所有的事都曾由同一天主藉着\OriginalTerm{圣言}{Word}向信徒启示和显示；祂时而与祂的受造物交谈，时而颁布祂的律法；时而责备，时而劝勉，然后释放祂的仆人，并收养他\OriginalTerm{作为儿子}{in filium}；并在适当的时候，赐予不朽的产业，为使人达于成全？因为祂造人是为了成长和增进，如圣经所说：\BibleQuote{增长繁衍。}\BibleRef{创 1:28}\par

2. 在这方面，天主与人的不同之处在于：天主确实创造，而人则被造；那创造者常是同一的，而被造者必须接受开始、中间、增添和增长。天主确实技巧地创造，而就人而言，他\Emph{是}被技巧地创造的。天主也真正在一切事上完美，自身相等且相似于自己，因为祂全是光、全是理智、全是实体、及一切美善的源泉；但人向天主接受进步和增长。因为正如天主常是同一的，同样，当人在天主内被发现时，他将永远迈向天主。因为天主从不停顿地施恩惠于人，或丰富人；人也从不停顿地接受恩惠，被天主所丰富。因为祂仁慈的容器和祂光荣化的工具，是那感谢祂造主的人；再者，祂公义审判的容器，是那忘恩负义的人，他既藐视他的造主，又不顺从祂的圣言；祂曾应许，要给那些常结果实的人很多，而对那些拥有主人银钱的，还要更多。\BibleQuote{做得好，}祂说，\BibleQuote{良善而忠信的仆人：因为你在小事上忠信，我要委派你管理许多事；进入你主人的福乐吧。}\BibleRef{玛 25:21} 主祂自己如此应许了很多。\par

3. 因此，正如祂应许要给那些现在结果实的人很多，是按照祂恩宠的恩赐，而不是按那变化不定的\Quote{\OriginalTerm{知识}{knowledge}；}（因为主保持不变，同一圣父被启示）同样，那同一的主藉着祂的来临，赐予后期的人比旧约经济下的人更大的恩宠恩赐。因为那些在旧约的人确实藉着[祂的]仆人听见君王要来，并因而欢喜到某种程度，因为他们盼望祂的来临；但那些亲眼看见祂亲临，获得了自由，并成为祂恩赐的分享者的人，拥有更大的恩宠和更高的喜乐，因君王的到来而欢欣：正如达味也说过：\Quote{我的灵魂将因主而喜乐；它将因祂的救恩而欢欣。} 为此，当祂进入耶路撒冷时，路上所有的人都认出了他们心怀忧伤的达味王，便铺开自己的衣服，并用青绿的树枝装饰道路，大声欢呼，充满喜乐：\BibleQuote{贺三纳于达味之子；奉主名而来的当受赞美；至高之处的贺三纳。}\BibleRef{玛 21:8} 但对于那些嫉妒的邪恶管家，他们欺诈下属，统治那些没有大智慧的人，并因此不愿君王来临，且对祂说：\Quote{你听见这些人在说什么？}主回答说：\Quote{你们没有读过：从婴孩和吃奶者的口中你完成了赞美？}——从而指出，那藉达味论及天主之子的话，已应验在祂自己身上；并指出他们确实无知于圣经的意义和天主的安排；但宣告祂自己就是那由先知所预言的基督，祂的名在普世受赞美，并从婴孩和吃奶者的口中向祂的圣父完成赞美；因此，祂的荣耀也被高举于诸天之上。\par

4. 因此，如果那由先知所预言的同一位，我们的主耶稣基督，现在已经临在，并且祂的来临给那些接受祂的人带来了更充实的恩宠和更大的恩赐，那么显然，圣父也是那由先知所宣告的同一圣父，而圣子来临时，并未传播另一位圣父的知识，而是传播那从起初就被宣讲的同一位圣父的知识；从祂那里，祂也给那些以合法的方式、自愿的心、并全心侍奉祂的人带来了自由；而对于嘲笑者、不顺服天主的人、以及那些为得人称赞而追求外在洁净的人（这些仪式曾被赐予作为未来事物的预像——律法仿佛以影儿预表某些事物，以暂时的事物描绘永恒的事物，以属地的事物描绘属天的事物），并对于那些假装自己遵守比规定更多之事的人，似乎把他们自己的热心置于天主本身之上，而内在却充满虚伪、贪婪及一切邪恶的人——祂为他们指定了永久的灭亡，将他们从生命中剪除。\par

\SourceRecord{Translated by Alexander Roberts and William Rambaut.；Ante-Nicene Fathers；Vol. 1.；Edited by Alexander Roberts, James Donaldson, and A. Cleveland Coxe.；Buffalo, NY: Christian Literature Publishing Co.,；1885.；<http://www.newadvent.org/fathers/0103411.htm>.}

% structure: p0001 source=h1 target=section reason=page-title

\section{\WorkTitle{驳斥异端}（第四卷，第十二章）}

1. 长老们自己所传的遗训——他们佯装按照法律遵守它——是与梅瑟所赐的法律相反的。因此依撒意亚也宣告说：\Quote{你们的商贾将酒兑水}\BibleRef{依 1:22}，表明长老们习惯将兑水的遗训与天主的简明诫命混合；也就是说，他们设立了一条伪造的、与［真正的］法律相反的法律。正如主也明确指出的，祂对他们说：\Quote{你们为何为了你们的遗训而违犯天主的诫命？}\BibleRef{玛 15:3}因为不仅通过实际违犯他们藐视天主的法律，将酒兑水；他们还制定了他们自己的法律与之对立，这法律直至今日仍被称为法利塞人的法律。在这［法律］中，他们压制某些事，添加其他事，又按自己的私意解释另一些事，他们的老师各自使用这些；并且渴望维护这些遗训，他们不愿服从天主的法律，而这法律是为他们预备基督的来临。他们甚至责备主在安息日治病，正如我已经指出的，法律并未禁止此事。因为他们自己在某种意义上也在安息日施行治病，当他们在［那日］给人行割损时；但他们没有因通过遗训和前述的法利塞人法律违犯天主的诫命，以及没有遵守法律的诫命——即爱天主——而责备自己。\par

2. 但是，这是第一条也是最大的诫命，其次［关于］爱近人，主已经教导了，当祂说全部法律和先知都系于这两条诫命。此外，祂并没有从天上带来任何比这更大的诫命，而是将这同一条诫命更新给祂的门徒，命令他们全心爱天主，并爱人如己。但如果祂是从另一位父降下，祂绝不会使用法律的第一条和最大诫命；而无疑会千方百计从完美的父那里带来一条比这更大的，以便不使用那由法律之神所赐的诫命。保禄同样宣告说：\Quote{爱德是法律的满全}\BibleRef{罗 13:10}；并且［他宣告］当其他一切都被废除时，将存留\Quote{信德、望德、爱德；但其中最大的是爱德}\BibleRef{格前 13:13}；并且没有爱天主的爱德，知识也无益\BibleRef{格前 13:2}，对奥迹的理解、信德、先知话也都无益；但没有爱德，一切都是空洞和虚妄的；再者，爱德使人成全；并且爱天主的人在现世和来世都是成全的。因为我们从不停止爱天主；而是随着我们继续默观祂，我们就越爱祂。\par

3. 因此，在法律和福音中，第一条也是最大的诫命是全心爱上主天主，其次与之相似的是爱人如己；如此，法律和福音的作者被显明是同一的。因为绝对完美生活的规诫，既然它们在每个盟约中都是相同的，已向我们指出了同一位天主，祂确实为每个盟约颁布了特别的律法；但那些更突出和最大的［诫命］，没有它们救恩无法［获得］，祂在两个盟约中同样劝勉我们遵守。\par

4. 主也没有废除这位［天主］，当祂表明法律并非来自另一位神时，祂向那些受祂教导的人、群众和祂的门徒这样表达：\Quote{经师和法利塞人坐在梅瑟的座位上。所以，凡他们所吩咐你们遵守的，你们要遵守并实行；但不要效法他们的行为：因为他们说而不做。他们把沉重难担的担子捆起来，放在人的肩上；但自己连一个指头也不肯动它们。}\BibleRef{玛 23:2}因此，当耶路撒冷尚且安全时，祂并没有指责那由梅瑟所赐的法律，反而劝勉遵守它；但祂\Emph{确实}指责那些人，因为他们虽重复法律的话，却没有爱德。因此他们被视为对天主和对邻人都行不义。正如依撒意亚也说：\Quote{这百姓用嘴唇尊敬我，心却远离我；他们敬拜我也是枉然，教导人的教条和吩咐。}\BibleRef{依 29:13}祂并没有称那由梅瑟所赐的法律为人的吩咐，而是长老们自己所发明的遗训，他们维护这些遗训，便使天主的法律失效，并因此也不服从祂的圣言。因为这就是保禄论及这些人所说的：\Quote{因为他们不了解天主的正义，企图确立自己的正义，就不服从天主的正义。因为基督是法律的终向，使一切相信的人获得正义。}\BibleRef{罗 10:3}基督怎能是法律的终向，如果祂不是法律的终因呢？因为祂带来了终向，祂也亲自开启了开端；正是祂亲自对梅瑟说：\Quote{我确实看见了我的人民在埃及所受的苦难，我下来要拯救他们。}\BibleRef{出 3:7}因为这从起初就是天主圣言的习惯，上上下下，为拯救那些在苦难中的人。\par

5. 法律预先教训人类必须跟随基督，这一点主亲自显明了。当有人问祂该做什么才能获得永生时，主回答说：\BibleQuote{你若要进入生命，就该遵守诫命。}\BibleRef{玛 19:17} 另一个人问：\Quote{哪些诫命？} 主又回答说：\BibleQuote{不可奸淫，不可杀人，不可偷盗，不可作假见证，孝敬父母，并要爱人如己，}——祂为那些愿意跟随祂的人，将法律的诫命作为进入生命的\OriginalTerm{阶梯}{velut gradus}置于他们面前；那时祂对一个人说的，就是对所有人说的。但那人说：\BibleQuote{这一切我都遵守了}（很可能他并没有遵守，因为如果遵守了，主就不会对他说\BibleQuote{遵守诫命}了），主揭露了他的贪心，对他说：\BibleQuote{你若愿意成为成全的，去，变卖你所有的，分给穷人，然后来跟随我。} 祂向那些如此行动的人应许了属于宗徒的份\OriginalTerm{宗徒的份}{apostolorum partem}。祂并没有向跟随者宣讲另一位圣父，有别于法律从起初所宣告的那位；也没有另一位圣子；也没有母亲，即那在痛苦与背逆中存在的\OriginalTerm{伊昂}{Æon}的\OriginalTerm{意旨}{enthymesis}；也没有那三十个\OriginalTerm{伊昂}{Æon}的\OriginalTerm{圆满}{Pleroma}，这已被证明是虚妄而不可信的；也没有其他异端者所编造的故事。祂教导他们应遵守天主从起初所命令的诫命，以善行消除先前的贪心，并跟随基督。分施财产给穷人能消除先前的贪心，这由匝凯所作的表白显明了，他说：\BibleQuote{主，你看，我把我财物的一半施舍给穷人；我如果欺骗过谁，我就以四倍赔偿。}\BibleRef{路 19:8}\par

\SourceRecord{Translated by Alexander Roberts and William Rambaut.；Ante-Nicene Fathers；Vol. 1.；Edited by Alexander Roberts, James Donaldson, and A. Cleveland Coxe.；Buffalo, NY: Christian Literature Publishing Co.,；1885.；<http://www.newadvent.org/fathers/0103412.htm>.}

% structure: p0001 source=h1 target=section reason=page-title

\section{\WorkTitle{驳斥异端（第四卷 第十三章）}}

1. 主并未废除使人成义的法律的自然规诫——那些因信德而成义、中悦天主的人在赐下法律之前也曾遵守这些规诫——而是扩展并成全了它们，这从祂的话可以显明。因为祂说：\BibleQuote{古时曾对你们说：不可奸淫。我却对你们说：凡注视妇女而贪恋她的，已在心中与她犯了奸淫。}\BibleRef{玛 5:27} 又说：\BibleQuote{曾有话说：不可杀人。我却对你们说：凡向弟兄发怒的（无故），应受裁判。}\BibleRef{玛 5:21} 以及：\BibleQuote{曾有话说：不可背誓。我却对你们说：什么誓都不可起；你们说话，是就说是，非就说非。}\BibleRef{玛 5:33} 以及其他类似的教训。所有这些都不是对过去规诫的反对或推翻——如马尔基昂的追随者们竭力主张的那样——而是成全与扩展，正如祂亲自宣告的：\BibleQuote{除非你们的义德超过经师和法利塞人的义德，你们决不能进天国。}\BibleRef{玛 5:20} 因为这\Quote{超过}是什么意思？首先，我们不仅要相信圣父，也要相信如今启示的圣子；因为正是祂引导人进入与天主的共融与合一。其次，我们不仅要说，而且要去做；因为他们说了，却没有去做。并且，我们不仅要戒避恶行，甚至连对恶行的贪欲也要戒避。如今祂教导我们这些，并非反对法律，而是为成全法律，并在我们内栽植法律的各样义德。倘若祂命令门徒做法律所禁止的事，那便是违反法律。但祂所命令的——即不仅戒避法律所禁止的事，甚至戒避对这些事的贪求——并不违反法律，正如我已指出的，这并非废除法律者的言论，而是成全、扩展并赋予法律更广阔范围者的言论。\par

2. 法律既是为那些在奴役之下的人设立的，便藉那些外在的、属于形体的对象教训灵魂，如同用一条绳索牵引它遵守诫命，使人学会服事天主。但\OriginalTerm{圣言}{Verbum}解放了灵魂，并教导灵魂藉着它，身体应心甘情愿地被净化。这既已完成，便自然随之而来：奴役的锁链——人已惯于其中——应被除去，人应无捆缚地跟随天主；此外，自由的法规应被扩展，对君王的顺服应增加，以致于没有一位皈依者看来不配他的解放者，而应有的孝爱和服从——即对家主应尽的——应同样由仆人和儿女献上；儿女却拥有更大的信心，因为自由的工作比在奴役状态下所行的服从更伟大、更光荣。\par

3. 为此，主以禁止\OriginalTerm{贪欲}{concupiscentia}取代了那条诫命——\BibleQuote{不可奸淫}；以禁止发怒取代了这条——\BibleQuote{不可杀人}；并且，取代法律关于缴纳什一之物的规定，祂教导我们将一切所有分施给穷人\BibleRef{玛 19:21}；不仅要爱近人，甚至要爱仇敌；不仅慷慨施舍馈赠，甚至对拿走我们财物的人也要白白给予。因为祂说：\BibleQuote{拿走你外衣的人，连内衣也给他；拿走你财物的，不要再讨回；你们愿意人怎样待你们，你们也要怎样待人。}\BibleRef{路 6:29} 这样我们便不像那些不愿被欺骗的人那样忧伤，而像那些乐意施予的人一样欢喜，且是向近人施恩而非出于无奈。祂又说：\BibleQuote{若有人强迫你走一里路，就同他走两里；}\BibleRef{玛 5:41} 这样你不再是作为奴隶跟随他，而是作为自由人走在他前面，在一切事上对近人显示善意和帮助，不顾念他们的恶意，而履行善举，使自己肖似天父，\BibleQuote{祂使太阳升起，照耀恶人也照耀善人，降雨给义人也给不义的人。}\BibleRef{玛 5:45} 如今我既已指出，所有这些规诫都不是废除法律者的命令，而是成全、扩展并扩大法律在我们中间者的命令；正如有人会说，自由更广泛的运行意味着在我们内植入了对\OriginalTerm{解放者}{Liberator}更完全的服从与热爱。因为祂解放我们，并非为使我们离开祂——的确，当人不在主恩惠的范围内时，无人有能力为自己获取救恩的途径——而是为使我们越接受祂的恩宠，就越爱祂。我们越爱祂，在常与圣父同在时，从祂所得的光荣也越大。\par

4. 由于所有自然诫命对我们和他们都（犹太人）是共同的，他们确实在其中获得了开端和起源；但在我们身上，它们获得了成长和完成。因为赞同天主、跟随祂的圣言、爱祂超过一切、并爱邻人如同自己（现在人就是邻人）、以及戒除一切恶行，以及所有其他类似性质、共同属于两个盟约的事，都启示了同一位天主。但这就是我们的主，天主的圣言，祂起初的确将奴隶带到天主面前，但后来祂释放了那些服从于祂的人，正如祂亲自向门徒宣告的：\BibleQuote{我不再称你们为仆人，因为仆人不知道他主人所做的事；但我称你们为朋友，因为凡我从我父那里听到的一切，我都告诉了你们。}\BibleRef{若 15:15} 因为当祂说：\BibleQuote{我不再称你们为仆人}，祂极其明确地指出，正是祂自己最初借着法律为人类设立了那在天主面前的奴役，然后后来又赐予他们自由。而祂说：\BibleQuote{因为仆人不知道他主人所做的事}，祂借着祂自己的到来，指出了处于奴役状态之人的无知。但当祂称祂的门徒为\BibleQuote{天主的朋友}时，祂明确宣告自己就是天主的圣言，亚巴郎也自愿地、\OriginalTerm{不受束缚地}{sine vinculis}跟随了祂，因为他的信德的高贵天性，从而成为了\BibleQuote{天主的朋友。}\BibleRef{雅 2:23} 但天主的圣言接受亚巴郎的友谊，并非因为祂需要它，因为祂从起初就是完美的（祂说：\BibleQuote{在亚巴郎出现以前，我就存在}\BibleRef{若 8:58}），而是因为祂的良善，要将永生赐予亚巴郎本人，因为天主的友谊使那些拥抱它的人获得不朽。\par

\SourceRecord{Translated by Alexander Roberts and William Rambaut.；Ante-Nicene Fathers；Vol. 1.；Edited by Alexander Roberts, James Donaldson, and A. Cleveland Coxe.；Buffalo, NY: Christian Literature Publishing Co.,；1885.；<http://www.newadvent.org/fathers/0103413.htm>.}

% structure: p0001 source=h1 target=section reason=page-title

\section{反异端论（第四卷，第十四章）}

1. 因此，在起初，天主造了\Person{亚当}，并非因为祂需要人，而是要有一个【某个人】可施予恩惠的对象。因为不仅先于亚当，而且在一切受造之前，圣言已光荣了祂的圣父，常居于祂内；祂自己也由圣父所光荣，正如祂自己所宣告的：\BibleQuote{父啊，在创世以前，我同祢所享有的光荣，求祢现在藉我在祢面前光荣我。}\BibleRef{若 17:5} 祂也不需要我们的服务而吩咐我们跟随祂；祂乃是借此将救恩赐予我们。因为跟随救主就是分享救恩，跟随光就是接受光。但在光中的人并不光照光，而是被光照耀并被光显示：他们确实对光毫无贡献，反而接受恩惠，被光照亮。同样，对天主的服务确实无益于天主，天主也不需要人的服从；但祂赐给那些跟随并事奉祂的人生命、不朽和永恒的光荣，施恩给那些事奉祂的人，因为他们事奉祂，并给那些跟随祂的人，因为他们跟随祂；但祂不从他们那里得到任何好处：因为祂是富足的、完美的、一无所缺的。但为此原因，天主要求人的服务，为的是，既然祂是良善仁慈的，祂可以恩待那些持续事奉祂的人。因为，天主一无所缺，人却多么需要与天主共融。因为人的光荣就在于持续并永远留在天主的服务中。因此，主也对祂的门徒说：\BibleQuote{不是你们拣选了我，而是我拣选了你们；}\BibleRef{若 15:16} 表明他们跟随祂时并没有光荣祂；反而，在跟随圣子时，他们被祂光荣了。又说：\BibleQuote{我愿我在哪里，他们也同我在一起，为能看见我的光荣；}\BibleRef{若 17:24} 并非因此虚夸，而是渴望祂的门徒分享祂的光荣：关于他们，\Person{依撒意亚}也说：\BibleQuote{我要从东方领回你的后裔，从西方把你聚集；我要对北方说：交出来！对南方说：不要阻拦！从远方领回我的儿子们，从地极领回我的女儿们；凡以我名被召的，都是我为自己光荣所创造、所形成的。}\BibleRef{依 43:5} 既然如此，\BibleQuote{无论在哪里有尸体，老鹰也必聚集在那里，}\BibleRef{玛 24:28} 我们就参与主的光荣，祂形成了我们，并为我们预备了这一切，使我们当与祂同在时，能分享祂的光荣。\par

2. 同样，天主起初造人，也是因为祂的慷慨；祂为了救恩的缘故而拣选了圣祖；祂预备了一个民族，教导那顽固者跟随天主；祂在地上兴起先知，使习惯于承载祂的圣神【在身内】，并与天主共融：祂自己确实一无所缺，却将共融赐予那些需要它的人，并像建筑师一样，为凡祂所喜悦的人描绘出救恩的蓝图。祂亲自给那些在埃及没有看见祂的人提供引导，而对那些在旷野中变得不顺从的人颁布了非常适合【他们的】法律。然后，祂赐予进入福地的百姓一份高贵的产业；祂为那些回头归向圣父的人杀了肥犊，并给他们穿上最好的袍子。\BibleRef{路 15:22} 这样，祂以多种方式调整人类，使之与救恩协调一致。为此，\Person{若望}在\WorkTitle{默示录}中也宣告说：\BibleQuote{祂的声音犹如大水的响声。}\BibleRef{默 1:15} 因为【天主的】圣神确实是【如】大水，因为圣父既富足又伟大。圣言经过所有这些【人】，以书面形式颁布了适用于每一类人的法律，慷慨地赐予祂的臣民恩惠。\par

3. 同样，祂给【犹太】百姓规定了建造会幕、修建圣殿、拣选肋未人、祭献、供物、法律训诫，以及法律的一切其他服务。祂自己确实不需要这些，因为祂永远充满一切美善，在\Person{梅瑟}存在之前，祂就拥有全部仁慈的香气和一切馨香的芬芳。此外，祂教导那易于转向偶像的百姓，通过反复的呼吁来劝导他们坚持并事奉天主，以次要之事召叫他们走向首要之事；即借助预象达到真实，借助暂时达到永恒，借助肉体达到精神，借助地上达到天上；正如曾对\Person{梅瑟}所说的：\BibleQuote{你要按照在山上所显示的样式制造一切。}\BibleRef{出 25:40} 因为四十天之久，他学习牢记天主的言语、天上的样式、精神的图像和未来之事的预像；正如\Person{保禄}所说：\BibleQuote{因为他们饮了那跟随他们的磐石中的水；那磐石就是\Person{基督}。}\BibleRef{格前 10:11} 又说，在首先提到法律中所包含的内容之后，他继续说：\BibleQuote{这些事发生在他们身上，是为作鉴戒，并且写下来，是为劝戒我们这些生活在末世的人。} 因为藉着预像，他们学会了敬畏天主，并持续地事奉祂。\par

\SourceRecord{Translated by Alexander Roberts and William Rambaut.；Ante-Nicene Fathers；Vol. 1.；Edited by Alexander Roberts, James Donaldson, and A. Cleveland Coxe.；Buffalo, NY: Christian Literature Publishing Co.,；1885.；<http://www.newadvent.org/fathers/0103414.htm>.}

% structure: p0001 source=h1 target=section reason=page-title

\section{驳斥异端（第四卷，第十五章）}

1. 他们（犹太人）因此拥有法律、训练课程和未来事物的预言。因为天主起初确实通过自然的诫命（即从起初就植入人类中的十诫）警告他们，如果任何人没有遵守，就没有救恩，那时并没有向他们要求更多。正如梅瑟在\BibleBook{申命}{纪}中说：\BibleQuote{以下是上主在西乃山上向以色列全体会众所说的一切话，并没有再多说；并将这些话写在两块石版上，交给了我。}\BibleRef{申 5:22} 正因为如此，祂这样做，使那些愿意跟随祂的人能遵守这些诫命。但当他们转向铸造牛犊，心中回到埃及，宁愿作奴隶而不作自由人时，他们就被置于一种适合他们愿望的奴役状态——这种奴役并没有使他们与天主隔绝，而是使他们服从束缚的轭；正如先知厄则克耳在说明颁布这种法律的原因时宣称：\BibleQuote{他们的眼睛随从了心中的欲望，我赐给他们不良好的法律，和不能藉以生活的规则。}\BibleRef{则 20:24} 路加也记载了斯德望——他是宗徒们首先选入执事职，并第一个为基督的见证而殉道的人——关于梅瑟说了以下的话：\Quote{这人确实接受了活天主的命令要传给你们，但你们的祖先不服从，反而\InnerQuote{拒绝祂}，心中转回埃及，对亚郎说：\InnerQuote{给我们造神在我们前面引路，因为不知道那领我们出埃及地的梅瑟遭遇了什么事。}那时他们造了一个牛犊，向偶像献祭，并因他们双手的作为而欢欣。但天主转身，任他们去敬拜天上的万军，正如先知书上所记载的：\BibleQuote{以色列家啊，你们在旷野四十年，曾向我献过牺牲和祭品吗？你们抬着摩洛的帐幕和楞番神的星，就是你们为敬拜所造的像。}\BibleRef{亚 5:25}\BibleRef{宗 7:38}} 清楚地指出，这样的法律并非由另一位神赐给他们，而是适应他们奴役的状况，来自同一位天主。因此，在\BibleBook{出谷}{纪}中也对梅瑟说：\BibleQuote{我要派遣我的使者在你前面，因为我不同你上去，因为你是固执的百姓。}\BibleRef{出 33:2}\par

2. 不仅如此，主也表明，某些法律是由梅瑟为他们制定的，因为他们的心硬，以及不愿意服从。当他们问祂：\Quote{为什么梅瑟却命令给休书休妻呢？}祂对他们说：\Quote{因为你们的心硬，他允许了你们这样作，但从起初并不是这样。}\BibleRef{玛 19:7} 这样，祂为梅瑟辩护，称他为忠信的仆人，同时承认一位天主，祂从起初就造了男人和女人，并责备他们心硬和不顺从。因此，他们从梅瑟那里接受了休妻的法律，适应他们刚硬的本性。但我为何在旧约中提及这些事呢？因为在新约中，宗徒们也是基于同样的理由这样做的，保禄明确地说：\BibleQuote{这是我说，不是主说。}\BibleRef{格前 7:12} 又说：\BibleQuote{我说这话，是出于特许，不是出于命令。}\BibleRef{格前 7:6} 又说：\BibleQuote{关于童女，我没有主的命令，但我既蒙主怜悯，成为忠信的，就提供我的意见。}\BibleRef{格前 7:25} 此外，在另一处他说：\BibleQuote{免得撒殚因你们不能节制而诱惑你们。}\BibleRef{格前 7:5} 因此，即使在新约中，宗徒们考虑到人性的软弱，因为有些人的不能节制，以免这些人变得顽固，对自己的救恩完全绝望，而成为背叛天主的人，就赐予某些规则——那么，在旧约中，同一位天主为祂的子民的利益允许类似的宽容，通过已经提到的规诫引导他们，使他们因遵守十诫而获得救恩的恩赐，并受到祂的约束，不转向偶像崇拜，不背叛天主，而是学习全心爱祂，就不足为奇了。如果有人因为那些不服从和灭亡的以色列人，而断言法律的赐予者（\OriginalTerm{导师}{doctor}）能力有限，他们将在我们的时代看到\BibleQuote{被召的人多，被选的人少}\BibleRef{玛 20:16}，而且有些人内心是狼，却在世人眼中（\OriginalTerm{外面}{foris}）披着羊皮；天主始终保留人的自由和自治的能力，同时发出祂自己的劝诫，为的是使那些不服从祂的人因没有服从而正义地受审判，而那些服从并相信祂的人则应被尊崇，获得永生。\par

\SourceRecord{Translated by Alexander Roberts and William Rambaut.；Ante-Nicene Fathers；Vol. 1.；Edited by Alexander Roberts, James Donaldson, and A. Cleveland Coxe.；Buffalo, NY: Christian Literature Publishing Co.,；1885.；<http://www.newadvent.org/fathers/0103415.htm>.}

% structure: p0001 source=h1 target=section reason=page-title

\section{驳斥异端（卷四，第十六章）}

1. 此外，我们从圣经本身得知，天主赐予割礼，并非作为义德的成全者，而是作为标记，使亚巴郎的种族保持可辨认。因为圣经宣告：\Quote{天主对亚巴郎说：你们中的每一个男子都应受割礼；你们应割去包皮，作为我与你们之间盟约的记号。}\BibleRef{创 17:9} 厄则克耳先知论及安息日也说：\Quote{我也赐给他们我的安息日，作为我与他们之间的记号，使他们知道我是祝圣他们的上主。}\BibleRef{则 20:12} 在出谷纪中，天主对梅瑟说：\Quote{你们应遵守我的安息日，因为它在我与你们之间，世世代代是一个记号。}\BibleRef{出 21:13} 这些事乃是作为记号赐下的；但这些记号并非没有象征意义，即既非无意义，也非徒然，因为是由智慧的工匠赐下的；但血肉的割礼预表了属神的割礼。因为宗徒说：\Quote{我们，}宗徒说，\Quote{受了非人手所行的割礼。}\BibleRef{哥 2:11} 先知也宣告：\Quote{要割去你们心的硬壳。}但安息日教导我们应当日复一日地事奉天主。\Quote{因为我们被看做，}宗徒保禄说，\Quote{终日待宰的群羊，}\BibleRef{罗 8:36} 也就是说，被祝圣［归于天主］，并不断服务于我们的信德，持守它，戒除一切贪婪，不在地上积攒或拥有财宝。\BibleRef{玛 6:19} 此外，\OriginalTerm{天主的安息}{requietio Dei}，即天国，可以说由受造物所指示；在这个［国度］中，那坚持\OriginalTerm{侍奉天主}{Deo assistere}的人，将在安息中分享天主的宴席。\par

2. 而且，人并非因这些事而成义，它们只是作为记号赐给百姓，这一事实表明——亚巴郎本人，没有割礼，也不守安息日，\Quote{信了天主，天主就以此算为他的义德；他被称为天主的朋友。}\BibleRef{雅 2:23} 此外，罗特没有受割礼，却被从索多玛领出来，从天主那里获得了救恩。同样，诺厄悦乐天主，虽未受割礼，却领受了［方舟］的尺寸，那是第二［人类］世界的尺寸。哈诺客也悦乐天主，未受割礼，却履行了天主在天使中的使节职务，尽管他是人，并被提去，至今保存着，作为天主公正审判的见证，因为天使们犯罪后被罚堕地受审，而那悦乐［天主］的人却被提升以得救恩。此外，所有亚巴郎以前的义人，以及梅瑟之前的圣祖们，都是独立于上述之事，并在梅瑟法律之外成义的。正如梅瑟本人在申命纪中对百姓说：\Quote{上主你们的天主在曷勒布订立了盟约。上主不是与你们的祖先订立这盟约，而是与你们。}\BibleRef{申 5:2}\par

3. 那么，主为何没有与祖先订立盟约？因为\Quote{法律不是为义人设立的。}\BibleRef{弟前 1:9} 但是义德的先祖们将十诫的意义写在他们的心和灵魂里，也就是说，他们爱那创造他们的天主，并不伤害邻人。因此，他们没有理由被警告性的诫命（\OriginalTerm{警告性的文字}{correptoriis literis}）所警戒，因为他们自身拥有法律的正直。但当这种正直和对天主的爱在埃及被遗忘并熄灭时，天主因对人类极大的善意，必然以声音显现自己，并以大能带领百姓出埃及，为使人再次成为天主的门徒和追随者；祂惩罚不顺服的人，使他们不致藐视创造者；并以玛纳喂养他们，使他们为灵魂领受食物（\OriginalTerm{接受理性的食物}{uti rationalem acciperent escam}）；正如梅瑟在申命纪所说：\Quote{祂用玛纳喂养了你，这是你祖先不知道的，为使你明白，人生活不只靠饼，也靠天主口中所发的一切言语。}\BibleRef{申 8:3} 它命令爱天主，并教导我们以公正对待邻人，使我们既不不义，也不配不起天主，祂通过十诫预备人获得祂的友谊，同样也预备人与邻人的和好——这些事确实有益于人本身；然而，天主并不需要从人得到任何东西。\par

4. 因此圣经说：\Quote{这些话是上主在山上对以色列全会众说的，并没有再加添；}\BibleRef{申 5:22} 因为正如我已经指出的，祂不需要从他们那里得到任何东西。梅瑟又说道：\Quote{以色列！现在上主你的天主向你要求什么？只是要求你敬畏上主你的天主，履行祂的一切道路，爱祂，全心全灵事奉上主你的天主。}\BibleRef{申 10:12} 这些事确实使人光荣，补足了人所欠缺的，即天主的友谊；但对天主并无益处，因为天主根本不需要人的爱。因为天主的荣耀是人所欠缺的，而除了事奉天主，人无法以任何其他方式获得这荣耀。因此，梅瑟又对他们说：\Quote{你要选择生命，好叫你和你后裔得以存活，爱慕上主你的天主，听从祂的声音，依附于祂；因为这就是你的生命，是你长寿的日子。}\BibleRef{申 30:19} 为预备人进入这生命，主亲自以祂自己的声音向众人说了十诫的话；因此，同样地，它们永久地与我们同在，借着祂血肉的来临得到扩展和增长，而非废除。\par

5. 然而，束缚的法典是被梅瑟逐一颁布给百姓的，适合他们的教导或惩罚，正如梅瑟自己所宣告的：\BibleQuote{那时主吩咐我教导你们法令和典章。} \BibleRef{申 4:14} 因此，这些为束缚而赐下、并作为他们记号的事物，祂藉着自由的新约取消了。但祂却增加了并扩展了那些自然的、高贵的、众人共有的法律，藉着义子身份，宽宏大量地、毫无吝啬地赐予人，使他们认识天主圣父，全心爱祂，坚定不移地遵行祂的圣言，同时不仅戒绝恶行，甚至戒绝恶念。但祂也增加了敬畏之情；因为儿子对父亲应有比奴隶更多的尊敬和更大的爱。因此，主说：\BibleQuote{至于人所说的每一句虚言，在审判之日，他们都要交账。} \BibleRef{玛 12:36} 并且，\BibleQuote{凡注视妇女而有意贪恋她的，他心里已经犯了奸淫；} \BibleRef{玛 5:28} 以及，\BibleQuote{凡向弟兄发怒而无故的，应受裁判的危险。} \BibleRef{玛 5:22} [这一切都说明，] 使我们知道，我们向天主交账，不仅在于行为，如同奴隶，而且在于言语和思想，如同真正接受了自由之权的人，在这种自由中，人受到更严格的考验，看他是否尊敬、敬畏并爱慕主。为此，伯多禄说：\BibleQuote{我们不可将自由当作邪恶的掩饰，} \BibleRef{伯前 2:16} 而是当作考验和证明信德的方法。\par

\SourceRecord{Translated by Alexander Roberts and William Rambaut.；Ante-Nicene Fathers；Vol. 1.；Edited by Alexander Roberts, James Donaldson, and A. Cleveland Coxe.；Buffalo, NY: Christian Literature Publishing Co.,；1885.；<http://www.newadvent.org/fathers/0103416.htm>.}

% structure: p0001 source=h1 target=section reason=page-title

\section{\WorkTitle{驳斥异端}（第四卷，第十七章）}

1. 此外，先知们最充分地指出，天主并不需要他们奴性的服从，而是为了他们自己的缘故，他才在法律中规定某些仪式。再者，主清楚地教导说，天主并不需要他们的祭品，而[只是要求它]，是为了献祭者本人的缘故，正如我已经指出的。因为当他看见他们忽略正义，远离天主的爱，并幻想天主能被祭献和其他象征性的仪式所平息时，撒慕尔就这样对他们说：\BibleQuote{天主不喜全燔祭和祭献，却愿他的声音被听从。看，欣然服从胜过祭献，听命胜过公羊的脂油。}\BibleRef{撒上 15:22} 达味也说：\BibleQuote{祭献和祭品祢不渴望，但祢开了我的耳朵；赎罪的全燔祭祢也不要求。} 他以此教导他们，天主渴望服从，这使他们安全，而非祭献和全燔祭，这些对正义毫无益处；[藉此声明]他同时预言了新约。在第五十篇圣咏中，他更清楚地论及这些事：\BibleQuote{因为祢若渴望祭献，我必会奉上；祢不会喜悦全燔祭。天主的祭献是痛悔的精神；忧伤和痛悔的心，上主必不轻看。} 因此，因为天主一无所缺，他在前一篇圣咏中声明：\BibleQuote{我不从你家取牛犊，也不从你圈中取公山羊。因为地上的一切走兽，山上的牛羊都是我的；我知道天上的一切飞鸟，田野的各种族类都是我的。我若饿了，不必告诉你，因为世界和其中的万物都是我的。难道我要吃公牛之肉，喝山羊之血吗？} 然后，免得有人以为他是因愤怒而拒绝这些事，他继续给人（人）劝告：\BibleQuote{向天主奉献赞颂的祭献，向至高者还你的愿，并在你困苦之日呼求我，我必拯救你，你必光荣我。} 他确实拒绝了罪人幻想能以此平息天主的事物，并显示他自己一无所缺；但他劝诫并建议他们做那些使人成义并亲近天主的事。同样的声明，依撒意亚也作了：\BibleQuote{你们向我献上大量祭献有何益处？上主说。我已饱足。}\BibleRef{依 1:11} 当他拒绝了全燔祭、祭献、祭品，以及新月、安息日、节期和所有这些伴随的仪式后，他继续劝勉他们做属于救恩的事：\BibleQuote{你们要洗涤，自洁，从我眼前除掉你们心里的邪恶：停止作恶，学习行善，寻求正义，解救受压迫者，为孤儿伸冤，为寡妇辩护；然后来，让我们彼此辩论，上主说。}\par

2. 并非因为他像人一样发怒，如同许多人敢于声称的那样，而拒绝了他们的祭献；而是出于对他们盲目的怜悯，并为向他们指明真正的祭献，通过奉献这祭献他们能平息天主，好从祂那里获得生命。正如他在别处声明：\BibleQuote{向天主的祭献是忧伤的心：一个赞美造物主的心是上主的馨香。} 因为如果他因愤怒而拒绝了他们的这些祭献，好像他们是不配获得他怜悯的人，他肯定不会催促他们做那些能使他们得救的事。但天主是仁慈的，他没有断绝他们的善劝。因为他在耶肋米亚说过：\BibleQuote{你们从舍巴带来乳香，从远方带来肉桂，有何益处？你们的全燔祭和祭献，我不接纳。}\BibleRef{耶 6:20} 他接着说：\BibleQuote{所有犹大，听上主的话。上主，以色列的天主这样说：修正你们的行径和行为，我必使你们安居在此。不要信赖谎言，它们对你们毫无益处，说：上主的圣殿，上主的圣殿，它[在此]。}\BibleRef{耶 7:2}\par

3. 再者，当祂指出祂领他们出埃及并非为此目的，即要他们向祂献祭，而是要他们忘记埃及人的偶像崇拜，能够听从主的声音——这声音为他们乃是救恩与荣耀时，祂通过同一\Person{耶肋米亚}宣告说：\Quote{上主这样说：将你们的全燔祭与祭献收集起来，吃你们的肉罢！因为在我领你们出埃及的那一天，我并没有论及全燔祭与牺牲的事，也没有吩咐你们。我只命令你们这事说：要听从我的声音，使我成为你们的天主，你们成为我的子民；并行走在我所命令你们的一切道路上，好使你们获得幸福。但他们没有听从，也没有侧耳倾听，反而随从自己邪恶心中的计谋行走，且向后退而不向前。}\BibleRef{耶 7:21} 又在同一先知说：\Quote{但凡要夸耀的，只可在这一点上夸耀：明悟并认识我是上主，我在世上施行慈爱、公义和审判；}\BibleRef{耶 9:24} 祂接着说：\Quote{因为我喜爱这些事——上主宣告——}\Quote{并非喜爱祭献、全燔祭或供物。}因为百姓并没有把这些诫命视为首要的\OriginalTerm{首要地}{principaliter}，而是视为次要的，且出于前述的理由，正如\Person{依撒意亚}又说：\Quote{你没有将你的全燔祭的羊带给我，也没有用你的祭献荣耀我；你没有用祭献服事我，也没有为乳香费心劳力；你没有用银钱为我买乳香，也没有贪图你祭献的油脂；你却以你的罪恶和不义站在我面前。}\BibleRef{依 43:23} 因此祂说：\Quote{这样的人才是我所垂顾的：谦卑、温良、对我的话战兢的人。}\BibleRef{依 46:2} \Quote{因为肥肉和脂油决不能除去你的不义。}\BibleRef{耶 11:15} \Quote{我所中意的斋戒是这样的：解开邪恶的锁链，废除暴虐的轭，释放受压迫者，折断一切不义的横木。把你的食粮分给饥饿的人，将无家可归的穷人领到你家中。你若见到裸体的人，就给他衣服穿，不要轻慢你的骨肉至亲\OriginalTerm{你本族的人}{domesticos seminis tui}。那时你的光明将如晨光迸发，你的伤口将迅速愈合；你的公义必在你前面行走，上主的荣耀必作你的后盾；你还在呼求时，我就应允：我在这里。}\BibleRef{依 58:6} 十二先知中的\Person{匝加利亚}也将天主的旨意指示给百姓说：\Quote{万有的上主这样说：要执行公义的审判，彼此施行仁慈和怜悯。不可欺压寡妇、孤儿、外方人和穷人；不可心里图谋恶事陷害弟兄。}\BibleRef{匝 7:9} 他又说：\Quote{你们应当讲的这些话是：你们每人要对自己的邻人讲真话，在你们的城门口要施行公义的审判。谁也不能心中图谋恶事陷害弟兄，你们不可喜爱假誓；因为这一切都是我所憎恨的——万有的上主说。}\BibleRef{匝 8:16} 此外，\Person{达味}也相似地说：\Quote{谁愿享受生命，愿享见幸福的日子？就要谨守口舌，不说坏话，谨守嘴唇，不说欺诈；躲避邪恶，努力行善，寻求和平，追随不辍。}\par

4. 从这一切清楚可见，天主并不要求他们的祭献和全燔祭，而是要求信德、服从和正义，为了他们的救恩。正如天主在\Person{欧瑟亚}先知中教导他们祂的旨意时所说：\Quote{我喜爱仁慈，胜过祭献；喜爱认识天主，胜过全燔祭。}\BibleRef{欧 6:6} 此外，我们的主也劝勉他们同样的事，祂说：\Quote{如果你们知道\InnerQuote{我喜欢仁慈胜过祭献}这句话的意思，就决不会判断无罪的人了。}\BibleRef{玛 12:7} 如此，祂为先知作证，证明他们宣讲了真理；同时斥责这些人（祂的听众）因自己的过错而愚昧无知。\par

5. 再者，当祂指示祂的门徒，向天主献上祂受造物的初果——并非因为祂需要它们，而是为使门徒自己不致不结果实或不感恩——祂拿起受造物中的饼，祝谢了，并说：\Quote{这是我的身体。}\BibleRef{玛 26:26} 同样，那杯——也是我们所归属的受造物的一部分——祂承认是自己的血，并教导了那新盟约的新祭献。这祭献由教会从宗徒们领受，在全世界上献给天主，即那位在\WorkTitle{新约}中赐予我们祂恩赐的初果作为生活之需的天主。关于这祭献，十二先知中的\Person{玛拉基亚}曾这样预言：\Quote{我不喜欢你们——万有的上主说——也不接受你们手中的祭献。因为从太阳升起到日落之处，我的名在外邦人中受尊崇；在各处，我的名受焚香和纯洁的祭献。因为我的名在外邦中尊大——万有的上主说。}\BibleRef{拉 1:10}——这些话极其清楚地指明：那先前的百姓（犹太人）确实将不再向天主献祭，但祭献——而且是纯洁的祭献——将在各处献给祂；祂的名在外邦人中受尊崇。\par

6. 但是，除了我们的主的名——圣父藉着这名受光荣，人也藉着这名受光荣——之外，在外邦人中还有什么别的名被光荣呢？因为是祂自己的圣子的名，而圣子是由祂造成的人，所以祂称这名是祂自己的。正如一位君王，若他亲自画了一幅他儿子的肖像，他理所当然地称这幅肖像是他自己的，这有两个原因：因为这是他的儿子的肖像，并且是他自己创作的作品。同样，圣父也承认耶稣基督的名是祂自己的，这名称在普世教会中被光荣，既因为这是祂圣子的名，也因为描述这名的祂将这子赐予人，为拯救人类。因此，既然圣子的名是属于圣父的，并且教会是在全能的天主内通过耶稣基督献上祭献，那么祂基于这两点说得好：\Quote{在各地都有人向我的名献上馨香和纯洁的祭献}。现在，若望在\BibleBook{若望}{默示录}中声明，\Quote{香}就是\Quote{圣人们的祈祷}。\par

\SourceRecord{Translated by Alexander Roberts and William Rambaut.；Ante-Nicene Fathers；Vol. 1.；Edited by Alexander Roberts, James Donaldson, and A. Cleveland Coxe.；Buffalo, NY: Christian Literature Publishing Co.,；1885.；<http://www.newadvent.org/fathers/0103417.htm>.}

% structure: p0001 source=h1 target=section reason=page-title

\section{驳斥异端（第四卷，第十八章）}

1. 因此，主命令在全天下献上的教会的祭品，在天主看为纯洁的祭献，并蒙祂悦纳；并非祂需要我们的祭献，而是献上者若蒙悦纳，他本人便因所献之物而得光荣。因为通过礼物，向君王显出尊荣和爱意；主愿意我们以纯朴无邪之心献上，便如此说：\BibleQuote{所以，你在祭台上献礼物时，若想起你的弟兄有什么怨你的事，就把礼物留在祭台前，先去与你的弟兄和好，然后再来献礼物。}\BibleRef{玛 5:23} 因此，我们有必要将受造物的初果献给天主，正如梅瑟也说：\BibleQuote{你不可空手出现在上主你的天主面前。}\BibleRef{申 16:16} 这样，人因他所显示感恩之事而被视为知恩者，便可得那从祂而来的荣耀。\par

2. 各种祭献的类别并未被废除；因为从前（在犹太人当中）有祭品，如今（在基督徒当中）也有祭品。从前在百姓中有祭祀，如今在教会中也有祭祀；只是种类有所改变，因为现今的奉献不是由奴隶，而是由自由人献上。因为主是（永远）同一的；但奴隶式的祭献有其（自身）特性，自由人的祭献也有其特性，为的是藉着祭献本身，显出自由的标记。因为在祂那里，没有任何事物是无目的的，也没有无意义或无意旨的。为此缘故，他们（犹太人）的确将什一之物奉献给祂，但那些获得自由的人，却将所有财产为主的事工而置，喜乐而慷慨地献出，并非只献次要部分，因为他们怀有（来世）更好事物的希望；如同那穷寡妇所作的一样，将她全部生活费投入天主的库中。\par

3. 因为在起初，天主看中了亚伯尔的礼物，因为他以纯朴和正直之心献上；但没有看中加音的奉献，因为他的心被嫉妒和恶意所分裂，就是他对兄弟所怀的，正如天主在责备他隐藏的（念头）时所说：\BibleQuote{你若献得正当，却分得不正，岂不犯了罪？安静吧！}\BibleRef{创 4:7} 因为天主并不被祭物所安抚。因为若有人只在外表上以无懈可击、按部就班地献祭，而心中却不给予邻人应有的正确交往，也不敬畏天主——如此心怀隐罪的人，并不能以外表正确的祭献欺瞒天主；这样的祭献对他毫无益处，只有放弃那在他内滋生的邪恶，才能使罪不致因假善的行为，使他成为自己的毁灭者。因此主也宣告：\BibleQuote{祸哉，你们经师和法利塞人，假善人啊！因为你们好像刷白的坟墓；外面好看，里面却装满死人的骨头和各样污秽。同样，你们外面在人前显出正义，里面却充满邪恶和假善。}\BibleRef{玛 23:27} 因为他们虽然在外表上被认为献得正确，内心却怀有类似加音的嫉妒；因此他们杀了那正义者，藐视圣言的劝告，如同加音一样。因为（天主）对他说：\BibleQuote{安静吧！}但他没有听从。那么，\Quote{安静吧}除了放弃预谋的暴力之外，还能是什么呢？祂对这些人说了类似的话，宣告说：\BibleQuote{瞎眼的法利塞人啊！先洗净杯子的里面，好使外面也干净。}\BibleRef{玛 23:26} 他们却不听从祂。因为耶肋米亚说：\BibleQuote{看，你的眼和你的心都不正，只顾贪图不义之财，流无辜人的血，施行欺压和杀人，好去行恶。}\BibleRef{耶 22:17} 依撒意亚又说：\BibleQuote{你们筹划，却不出于我；你们立约，却不出于我的神。}\BibleRef{依 30:1} 因此，为了使他们内心的意愿和思想被揭露，以显示天主无可指责，且不作恶——这位天主揭示（内心）隐藏的事，却不作恶——当加音丝毫没有安静下来时，祂对他说：\BibleQuote{你要制服它，你要管辖它。}\BibleRef{创 4:7} 祂也这样对比拉多说：\BibleQuote{你对我没有任何权柄，除非从上头赐给你。}\BibleRef{若 19:11} 天主总是将义者（在此生）交出（受苦），使他经过所受和所忍的考验，得以（最终）蒙悦纳；而作恶者，因他所行的行为受审判，而被弃绝。因此，祭献并不能圣化人，因为天主并不需要祭献；乃是献祭者的良心，当它纯洁时，圣化祭献，从而感动天主（接纳）这（奉献）如同来自朋友。\BibleQuote{但罪人，}祂说，\BibleQuote{凡向我献牛犊为祭的，就如杀了狗一般。}\BibleRef{依 66:3}\par

4. 既然教会以纯一之心奉献，她的礼品就应被算作在天主前纯洁的祭献。如保禄也对斐理伯人说：\BibleQuote{我一切都满足，因为我从\Person{厄帕夫洛狄托}收到了你们所送来的，那馨香的芬芳，可悦纳的祭献，中悦天主的。}\BibleRef{斐 4:18} 因为我们理当向天主奉献祭献，并在一切事上以纯洁的心灵、无伪的信德、稳固的望德、炽热的爱德，献上祂所造万物的初果，以感恩的心对待我们的造物主天主。唯有教会向造物主奉献这纯洁的祭献，怀着感恩之心，将祂所造的万物奉献给祂。但犹太人并不如此奉献：因为他们的手沾满了血；因为他们没有领受那圣言——借着祂，祭献才献给天主。同样，异端者的任何\OriginalTerm{秘密集会}{synagogae}也不献此祭。因为有些人主张父不同于造物主，当他们向父奉献属于我们这受造界之物时，就把父描绘成贪图他人财产、希求非己之物的人。另一些人主张我们周围的事物源于背弃、无知和情欲，当他们向父奉献无知、情欲和背弃的果实之时，反而得罪了他们的父，是对父的侮辱而非感恩。然而，他们怎能自圆其说呢？因为他们说那祝谢过的饼是主的身体，那杯是祂的血，却不承认祂自己是世界造物主之子，即祂的圣言——借着祂，树木结果，泉源涌流，大地生出\BibleQuote{先有苗，后有穗，然后穗上结满饱满的麦粒。}\BibleRef{谷 4:28}\par

5. 再者，他们怎能说那以主的身体和祂的血所滋养的肉体会腐朽，并不得分享生命呢？因此，他们要么改变意见，要么停止奉献上述之物。但我们的意见与圣体圣事相符，而圣体圣事也反过来证实我们的意见。因为我们向祂奉献祂自己的万物，始终宣告血肉与圣神的共融与结合。因为正如那由大地所出的饼，当它接受了天主的呼求之后，就不再是寻常的饼，而是圣体圣事，由两重实在组成——地上的和天上的；同样，我们的身体在领受圣体圣事后，也不再腐朽，而是怀着复活的希望而达于永生。\par

6. 如今我们向祂奉献，并非因为祂需要什么，而是为了感恩祂的恩赐，并借此圣化受造之物。因为正如天主不需要我们的财产，我们也需要向天主奉献；正如撒罗满说：\BibleQuote{怜悯穷人的，就是借给上主。}\BibleRef{箴 19:17} 因为一无所缺的天主，为了使我们获得祂美善之物的赏报，将我们的善工纳为己有，正如我们的主说：\BibleQuote{我父所祝福的，你们来吧！承受那为你们预备的国度吧！因为我饿了，你们给了我吃的；我渴了，你们给了我喝的；我作客，你们收留了我；我赤身露体，你们给了我穿的；我患病，你们看顾了我；我在监里，你们来探望了我。}\BibleRef{玛 25:34} 因此，正如祂不需要这些服务，却仍愿意我们为了自己的益处而实行，免得我们不结果实；同样，圣言也向百姓颁赐了关于奉献祭献的诫命，虽然祂并不需要它们，为使他们学习事奉天主：因此，同样也是祂的旨意，要我们常而不辍地在祭台上献礼。祭台则在天上（因为我们的祈祷和奉献都朝向那里）；圣殿也同样在那里，正如若望在默示录中说：\BibleQuote{天主的圣殿敞开了。}\BibleRef{默 11:19} 帐幕也在那里：\Quote{看，}祂说，\Quote{天主的帐幕，祂要与人同住。}\par

\SourceRecord{Translated by Alexander Roberts and William Rambaut.；Ante-Nicene Fathers；Vol. 1.；Edited by Alexander Roberts, James Donaldson, and A. Cleveland Coxe.；Buffalo, NY: Christian Literature Publishing Co.,；1885.；<http://www.newadvent.org/fathers/0103418.htm>.}

% structure: p0001 source=h1 target=section reason=page-title

\section{\WorkTitle{驳斥异端，第四卷，第十九章}}

1. 现在，那些礼品、供物和所有祭献，人民是在预像中领受的，正如在山上指示给梅瑟的那位唯一同一的天主，如今祂的名字在外邦人中的教会里受到光荣。但那些确实散布在我们周围的尘世之物，适合作为天上之物的预像，它们两者都是由同一天主所创造的。因为祂没有别的方法能使属灵事物的形象[适合于我们的理解]。但若有人断言那些超天界和属灵的、就我们而言不可见且不可言说的事物，反过来又是天上之物和另一个\OriginalTerm{圆满}{Pleroma}的预像，并且说天主是另一位圣父的肖像，那便是扮演真理的流浪者，以及全然愚昧和愚蠢之人的角色。因为，如我多次指出的，这样的人将不得不不断地找出预像的预像、肖像的肖像，并且永远无法将他们的心智确立于唯一真天主之上。因为他们的想象超越天主，他们在心中已超过主自己，确实在观念上自高自大，凌驾于祂之上，但事实上却背离了真天主。\par

2. 对这些人，可以公正地说（正如圣经本身所暗示的）：你们把想象提到天主之上多远啊，啊，你们这些鲁莽自高自大的人？你们曾听见说：\BibleQuote{诸天曾被祂的手掌量过：}\BibleRef{依 40:12} 告诉我那尺度，数算那无尽的肘尺，为我解释那测量的丰满、宽度、长度、高度、起点和终点——这些是人心不能理解，也无法领悟的事。因为天上的宝库确实伟大：天主不能在心中被测量，在理智中不可理解；祂用掌心托住大地。谁能领悟祂右手的尺度？谁知道祂的手指？或者谁理解祂的手——那量度无限的手；那以自身尺度展开诸天尺度的手，那在掌中容纳大地和深渊的手；那在自己内包含整个受造界的宽度、长度、下面的深度和上面的高度的手；那被看见、被听见、被理解，却又不可见的手？因此，天主\BibleQuote{高于一切率领者、掌权者、异能者，以及一切可称呼的名号，}\BibleRef{弗 1:21} 超过一切受造和设立之物。祂充满诸天，俯视深渊，也与我们每个人同在。因为祂说：\BibleQuote{我岂是近处的天主，而不是远处的天主？人若藏在隐秘处，我岂看不见他？}\BibleRef{耶 23:23} 因为祂的手抓住万有，这手照亮诸天，也光照天之下的事物，察验肺腑和心肠，也临在于隐秘之处和我们的秘密思想中，公开地养育和保护我们。\par

3. 但若人不理解祂手的丰满和伟大，谁能在心中领悟或认识如此伟大的天主呢？然而，仿佛他们已经测量并彻底调查了祂，从各方面探索了祂，他们便臆测在祂之外存在着另一个由众\OriginalTerm{永世}{Æons}构成的\OriginalTerm{圆满}{Pleroma}和另一位圣父。他们确实不仰望天上之物，而是真正陷入疯狂的\OriginalTerm{深渊}{Bythus}之中；声称他们的圣父只延伸到那些在圆满之外的事物的边界，而另一方面，\OriginalTerm{造物主}{Demiurge}则达不到圆满；这样他们就把两者都描述为不完全且不透知万物的。因为前者在圆满之外形成的整个世界方面有缺陷，而后者在圆满之内形成的那个[理想]世界方面有缺陷；因此[两者]都不能成为万有的天主。但是，无人能通过天主所造之物完全宣示天主的良善，这一点对所有人都是显而易见的。而且，祂的尊威并无缺陷，反而包含万有，甚至延伸到我们，与我们同在，凡对天主怀有相称观念的人都会承认这一点。\par

\SourceRecord{Translated by Alexander Roberts and William Rambaut.；Ante-Nicene Fathers；Vol. 1.；Edited by Alexander Roberts, James Donaldson, and A. Cleveland Coxe.；Buffalo, NY: Christian Literature Publishing Co.,；1885.；<http://www.newadvent.org/fathers/0103419.htm>.}

% structure: p0001 source=h1 target=section reason=page-title

\section{反驳异端（第四卷，第二十章）}

1. 就祂的伟大而言，人不可能认识天主，因为父是不可测量的；但就祂的爱（这是藉祂的圣言引领我们归向天主的）而言，当我们顺服祂时，我们常学习到有如此伟大的天主，是祂亲自建立、拣选、装饰并包含万有；在万有之中，也包括我们自己和这个世界。我们也是与祂所包含的事物一起被造的。这就是圣经所说：\BibleQuote{天主用地上的泥土形成了人，将生命的气息吹在他的脸上。}\BibleRef{创 2:7}因此，不是天使造了我们，也不是天使形成了我们，天使也没有能力制造天主的肖像，除主的圣言外，也没有任何其他远隔万有之父的能力。因为天主并不需要这些〔存有〕来完成祂自己事前决定要做的事，好像祂没有自己的双手似的。因为\OriginalTerm{圣言}{Word}与\OriginalTerm{上智}{Wisdom}，即圣子和圣神，常与祂同在；祂藉着他们并在他内，自由而自发地创造了万有，祂也对他们说：\BibleQuote{让我们按我们的肖像和模样造人。}\BibleRef{创 1:26}祂从自己身上取出受造物的实体、受造物的样式和世上一切装饰的典型。\par

2. 确实，圣经宣告说：\Quote{首先当相信只有一位天主，祂建立了万有，并使万有得以完成，并使万有从无中生出。}祂包含万有，自己却不被任何人包含。玛拉基亚先知也正确地说：\BibleQuote{不是一位天主创造了我们吗？我们不是同有一位父吗？}\BibleRef{拉 2:10}与此一致，宗徒也说：\BibleQuote{只有一个天主，就是父，祂超越万有，并在我们众人之内。}\BibleRef{弗 4:6}同样，主也说：\BibleQuote{一切都由我父交付给我。}\BibleRef{玛 11:27}明显是由那创造万有的父交付；因为祂交给祂的不是别人的东西，而是自己的。但\Emph{万有}意味着没有〔任何事物〕被保留〔不给祂〕，因此同一人是生者与死者的审判者；\BibleQuote{他拿着达味的钥匙：他开了，没有人能关；他关了，没有人能开。}\BibleRef{默 3:7}因为无论在天上、地上或地下的，无人能打开父的书卷，或得见父，除了被宰杀的羔羊，祂用自己的血赎回了我们，从同一天主（祂藉圣言创造了万有，并藉〔祂的〕上智装饰了万有）那里获得了驾驭万有的权柄，那时\BibleQuote{圣言成了血肉}；为的是天主圣言在天上拥有主权，同样也在世上拥有主权，因为〔祂是〕义人，\BibleQuote{祂没有犯罪，口中也没有诡诈。}\BibleRef{伯前 2:23}祂要超越那些在地下的，祂自己成为\BibleQuote{死者中的首生者}\BibleRef{哥 1:18}；为的是万有，如我已说过的，能得见他们的君王；为的是父的光明能与我们的主的肉身相遇并安息，并从祂光辉的肉身来到我们这里，这样人就能穿上父的光明而获得不朽。\par

3. 我也已充分证明，圣言，即圣子，常与父同在；而上智，即圣神，也在万有创世之前就已与祂同在，祂藉撒罗满宣告说：\BibleQuote{天主以上智建立了大地，以智慧奠定了高天。因祂的知识，深渊裂开，云彩滴下露珠。}\BibleRef{箴 3:19}又说：\BibleQuote{上主在祂造化的起头，在祂工作的开始就造了我。在永恒之初，在起初，在造地之先，在深渊出现之前，在水泉涌出之前，在群山奠定之前，在诸山冈之前，祂就生了我在世之初。}又说：\BibleQuote{当祂预备高天时，我已在祂那里，当祂奠定深渊的泉源时，当祂使大地的基础坚强时，我与祂一同预备。我日日为祂所喜悦，常在祂面前欢乐，当祂因完成世界而欢乐时，欣喜于人子。}\BibleRef{箴 8:27}\par

4. 因此，只有一位天主，祂藉圣言和上智创造并安排了万有；这位造物主（Demiurge）将这个世界赐予人类，就祂的伟大而言，祂对所有被造者都是未知的（因为没有人能探究祂的高度，无论是古代已安息的人，还是任何现今活着的人）；但就祂的爱而言，祂常通过祂所借以安排万有的那一位而为人所知。这一位就是祂的圣言，我们的主耶稣基督，祂在末世成了人中间的人，为的是将终结与开端相连，即人与天主相连。为此，先知们从同一圣言领受了先知之恩，宣告了祂按肉身的来临；藉此，天主与人的结合与共融得以按照父的美意实现；天主的圣言从起初就预言，天主将被人看见，并在世上与人交谈，与人商议，并与祂自己的受造物同在，拯救它，成为可被它感知的，使我们摆脱所有恨我们之人的手，即脱离一切邪恶之灵；并使我们一生在圣德和正义中事奉祂\BibleRef{路 1:71}，为使人，既拥抱了天主的圣神，能进入父的光荣之中。\par

5. 先知们以预言的方式宣告了这些事；但他们并未像有些人所声称的那样，宣告被先知看见的那位是不同的［天主］，而万有的圣父是不可见的。然而这正是那些完全不懂预言本质的［异端者］所宣称的。因为预言是对未来之事的预告，即预先陈明后来将要发生的事。因此，先知们预先指明天主将被人看见；正如主也说：\Quote{心里洁净的人是有福的，因为他们要看见天主。}\BibleRef{玛 5:8} 但就他的伟大和他奇妙的荣耀而言，\Quote{没有人看见天主而能存活，}\BibleRef{出 33:20} 因为圣父是不可领悟的；但就他的爱、仁慈和他无限的权能而言，他甚至将此赐给那些爱他的人，即看见天主，这正是先知们也预言过的事。\Quote{因为在人不可能的，在天主却可能。}\BibleRef{路 18:27} 因为人不是凭自己的能力看见天主；而是当他愿意时，他就被那些他所愿意、在他所愿意的时候、照他所愿意的方式的人看见。因为天主在一切事上都是大能的，当时藉着圣神以预言的方式被看见，也藉着圣子以义子的方式被看见；而且他将来还要在天国中以父的方式被看见，圣神确实在天主之子内预备人，圣子引领人归向圣父，而圣父也赐予人不朽以得永生，永生来自每人看见天主的事实。因为正如那些看见光的人在光内，并分享其光辉；同样，那些看见天主的人在天主内，并领受他的荣光。但他的荣光赐予他们生命；因此，看见天主的人就领受生命。正因如此，他虽超越理解、无限且不可见，却使自己成为可见的、可理解的、并在信徒的能力范围之内，为的是使他能赐生命给那些藉着信德接受并看见他的人。因为他的伟大不可穷究，同样，他的良善不可言喻；藉此良善，他被人看见时，就将生命赐给那些看见他的人。离开生命就不能生存，而生命之资在于与天主的共融；但与天主的共融就是认识天主，并享受他的良善。\par

6. 人因此将看见天主，为能生活，藉那看见成为不朽的，甚至达到天主那里；正如我已经说过，先知们以预象的方式宣告过，即天主将被那些在他内拥有圣神并始终耐心等待他来的人看见。正如梅瑟在\BibleBook{申命}{纪}中所说：\Quote{我们今日看见天主与世人说话，而人仍能存活。}\BibleRef{申 5:24} 因为这些人中某些人曾看见预言的圣神及其为各样恩赐而倾注的活跃影响；另一些人［看见］主的来临，以及那从起初就有的、藉此他成就了圣父对天上和地下之事的旨意的经世；还有一些人［看见］与时代、与当时看见和听见的人以及后来一切将要听见的人相应的父的荣耀。因此，天主就是这样被启示出来的；因为天主圣父藉着所有这些［运作］显示出来：圣神确实在工作，圣子在服务，圣父在赞许，人的救恩得以完成。正如他通过欧瑟亚先知所宣告的：\Quote{我，}他说，\Quote{多加了异象，并藉先知们的手（\Emph{in manibus}）用了比喻。}\BibleRef{欧 12:10} 但宗徒解释了这一处经文，他说：\Quote{恩赐各有区别，却是同一圣神；职分各有区别，却是同一主；功效各有区别，却是同一天主，在一切人身上行一切。圣神显示在每人身上，是为人的益处。}\BibleRef{格前 12:4} 然而，作为在一切内行一切的天主，就他的本性和伟大而言，［天主］对他所造的一切事物是不可见、不可描述的，但他绝非不可认识：因为万物藉着他的圣言知道只有一个天主圣父，他包罗万有，并赐予万有存在，正如福音书上所写：\Quote{从来没有人看见天主，只有那在父怀里的独生子，他给我们详述了。}\BibleRef{若 1:18}\par

7. 因此，圣父的圣子从起初就宣告［他］，因为他从起初就与圣父同在，他也按着适当的次序和联系，在合适的时候，为［人类］的益处，向人类显示了预言的异象、各样的恩赐、他自己的服务以及圣父的荣耀。因为哪里有了承继的次序，哪里就有固定性；哪里有固定性，哪里就有时代的适应性；哪里有适应性，哪里就有益处。正因如此，圣言成了为人类利益而分施父的恩宠的管家，他为人类行了如此伟大的经世，一方面将天主启示给人，另一方面将人呈献给天主，同时保全了圣父的不可见性，免得人随时轻慢天主，并使人总有一些可进取的目标；但另一方面，他又通过多种经世将天主启示给人，免得人完全远离天主而归于无有。因为天主的光荣是活生生的人；而人的生命在于瞻仰天主。因为那藉着受造物所作的天主的显示，给地上一切生物带来生命，那么，那藉着圣言而来的圣父的启示，更给那些看见天主的人带来生命。\par

8. 既然天主之神藉先知预示未来的事，预先塑造和调整我们，使我们服从天主；但人藉圣神的恩惠看见[天主]仍是未来的事，那么，那些藉以宣告未来之事的人，必须看见他们所预告的人将看见的天主；如此，天主、天主之子、圣子、圣父不仅被先知性地宣告，而且也要被祂所有成圣并在天主之事上受教的肢体看见，使人预先受训并事先操练，以领受那日后将在爱天主的人身上显示的荣耀。因为先知们不仅以言语预言，也以异象、生活方式和按圣神默感所行的事工预言。因此，他们以这种无形的方式看见了天主，正如依撒意亚所说：\Quote{我亲眼看见了君王，万军的上主。}\BibleRef{依 6:5}指出人将用眼睛看见天主，并听见祂的声音。同样，他们也看见了天主之子作为与人相处的人，预言将要发生的事，说那尚未来临的已经临在，宣布那不可受难者\OriginalTerm{不可受难者}{impassible}应受苦难，宣告那在天上的已降入死亡的尘土中。此外，关于祂所要完成的总结的其他安排，有些他们通过异象看见，有些他们用言语宣告，有些他们藉[外在]行动象征性地指示，用眼看见那些将要被看见的事；用口宣告那些将要被听见的事；用实际操演将要发生的事；但同时也都先知性地宣告一切。因此，梅瑟也宣布天主对违犯法律的百姓实在是以火\OriginalTerm{火}{igneum}吞噬的天主\BibleRef{申 4:24}，并警告天主将带来火的日子；但对敬畏天主的人，他说：\Quote{上主天主是慈悲宽仁的，缓于发怒，富于慈爱和忠诚，对千万人保持慈爱，赦免过犯、罪行和罪恶。}\BibleRef{出 34:6}\par

9. 圣言向梅瑟说话，在他面前显现，\Quote{好像人与朋友谈话一样。}\BibleRef{户 12:8}但梅瑟渴望面对面看见那与他说话者，却被如此告知：\Quote{站在这磐石的深处，我将用手遮盖你。当我的荣耀经过时，你将看见我的背面，但你不能看见我的面，因为没有人看见我而能存活。}\BibleRef{出 33:20}这里表明了两种事实：人不可能看见天主；以及，通过天主的智慧，人将在末世在磐石的深处看见祂，即祂以人的身份来临。为此，祂[主]在山上与他面对面交谈，厄里亚也在场，正如福音所记\BibleRef{玛 17:3}，祂如此在末后实现了古老的应许。\par

10. 因此，先知们并未公开看见天主的真正面容，而是看见了那些以后人将看见天主的安排和奥秘。正如对厄里亚所说：\Quote{你明天要出去，站在上主面前；看，上主从那里经过。在上主面前，有强烈的大风，裂山碎石，但上主不在风中；风后有地震，但上主不在震中；地震后有火，但上主不在火中；火后有细微的声音}\OriginalTerm{细微的声音}{vox auræ tenuis}。\BibleRef{列上 19:11}因为藉此，那位因百姓的过犯和先知被杀而十分愤慨的先知，被教导行事要更加温和；同时，主作为人的来临被指明，它将在梅瑟所赐的法律之后，温柔而宁静，在此祂不折断压伤的芦苇，也不吹灭将熄的灯心。\BibleRef{依 42:3}祂王国的温和与平安的安息也被指明。因为，在裂山的大风之后，在地震之后，在火之后，祂王国的宁静和平安时代随之而来，在此天主之神以最温柔的方式使人活过来并成长。这一点由厄则克耳更清楚地显明：先知们部分地看见了天主的安排，而非天主本身。因为当此人看见了天主的异象\BibleRef{则 1:1}、革鲁宾、轮子，并叙述了整个前进的奥秘，看见了它们之上的宝座样式，宝座上有人形模样的样式，腰间以上有琥珀色的样式，以下有如火的形状，并陈述了宝座异象的其余部分，免得有人以为在这些[异象]中他真看见了天主，他补充说：\Quote{这就是天主荣耀的形像的样式。}\BibleRef{则 2:1}\par

11. 那么，如果梅瑟、厄里亚、厄则克耳，这些人都曾见过许多天上的异象，却都没有看见天主；而他们所见的只是主荣耀的相似物和对未来之事的预言，那么显然，圣父确实不可见，主也曾论到祂说：\BibleQuote{从来没有人看见过天主。}\BibleRef{若 1:18} 但是祂的圣言，照祂自己的旨意，并为了那些看见之人的益处，显示了圣父的荣光，并说明了祂的旨意（主也如此说：\BibleQuote{那在圣父怀里的独生天主，祂已解释出来；}祂自己也诠释圣父的圣言，因其丰盛而伟大）；祂不是以一种形象，也不是以一种特征，向那些看见祂的人显现，而是按照祂各项安排的目的和效果，如同在\WorkTitle{达内尔书}所记。因为有一次，祂与阿纳尼雅、阿匝黎雅、米沙耳周围的人一同显现，在火窑的火焰中与他们同在，保护他们免受火的伤害：经上说：\BibleQuote{那第四位的面貌好像天主子。}\BibleRef{达 3:26} 另一次，祂如同一块\BibleQuote{非人手所凿的石头从山上滚下，}\BibleRef{达 7:13} 击碎所有暂时的国度，并将它们\OriginalTerm{吹散}{ventilans ea}，而祂自己充满全地。然后，同样这位被看见为人子乘着天上的云彩而来，走近万古常存者，从祂领受一切权能、荣耀和国度的：经上说：\BibleQuote{祂的统治是永远的统治，祂的国度永不灭亡。}\BibleRef{达 7:4} 主的门徒若望，当看见祂王国司祭和荣耀的来临时，在\WorkTitle{默示录}中说：\BibleQuote{我转过身来，要看那同我说话的声音；我转过来，看见七盏金灯台；在灯台中间有一位相似人子，身穿长袍，胸间束着金带；祂的头和头发洁白如羊毛、如雪；眼目如同火焰；脚像在炉中锻炼过的光铜；声音如同大水的声音；祂右手拿着七星；口中吐出一把双刃利剑；面貌如同烈日发光。}\BibleRef{默 1:12} 因为在这话中，祂显示了自圣父所得荣耀的一部分，例如头部；也显示了有关司祭职分的事，如那垂到脚上的长袍。这也是梅瑟如此装饰大司祭的原因。也暗示了终结（万物终局），如光铜在火中锻炼，象征信德的力量和在祈祷中的恒心，因为世界末了将要来的焚火。但当若望不能忍受这景象（因他说：\BibleQuote{我就在祂脚前像死了一样；}\BibleRef{默 1:17} 为应验所记的：\BibleQuote{没有人看见天主而能存活}\BibleRef{出 33:20}），而圣言使他苏醒，并提醒他这就是那位他在晚宴上曾依靠胸口、询问谁将出卖祂的主，宣告说：\BibleQuote{我是首先的、末后的，是那活着、曾死过，如今永远活着，并且拿着死亡和地狱的钥匙。}在这之后，他在第二个异象中看见同一位主，说：\BibleQuote{我又看见在宝座和四个活物中间，并在长老们中间，有一只羔羊站立，像是被杀过的，有七角和七眼，就是天主的七神，奉差遣往普天下去。}\BibleRef{默 5:6} 他又论到同一位羔羊说：\BibleQuote{我又看见一匹白马，骑在马上的称为忠信和真实，祂按公义审判和争战。祂的眼睛如同火焰，头上戴着许多冠冕；又有写着的名字，除了祂自己没有人知道；祂穿着溅了血的衣服；祂的名称为天主圣言。天上的众军骑着白马，穿着洁白的细麻衣跟随祂。有利剑从祂口中出来，可以击杀列国；祂必用铁杖\OriginalTerm{牧放}{pascet}他们；并要踹全能天主烈怒的酒榨。在祂衣服和大腿上写着名号：万王之王，万主之主。}\BibleRef{默 19:11} 如此，天主的圣言始终保持着未来之事的轮廓，犹如向人指出圣父各项安排的种种\OriginalTerm{形式}{species}，教导我们关于天主的事。\par

12. 然而，天主不仅借着所见的神视和所传的话语，也借着实际的事工，让先知看见祂，为的是借他们预表和显示未来之事。为此，先知欧瑟亚娶了\Quote{淫妇}，借这行动预言\BibleQuote{这地要犯奸淫，离弃上主}\BibleRef{欧 1:2}，即地上的人；上主乐意从这种人中选出一群教会\BibleRef{宗 15:14}，借着与祂圣子的共融而成圣，如同那女人因与先知结合而成圣一样。为此，保禄宣告\Quote{不信主的妻子因信主的丈夫而成圣}。再者，先知给儿女起名叫\BibleQuote{未获怜悯}和\BibleQuote{非我民}\BibleRef{欧 1:6}，为的是如宗徒所说：\BibleQuote{那非我民的要成为我民；未获怜悯的要获得怜悯。从前在什么地方对他们说：你们不是我的人民，将来就在那里称他们为永生天主的儿女。}\BibleRef{罗 9:25} 先知借着行动所预表的事，宗徒证明基督已在教会中真实成就了。同样，梅瑟也娶了一位雇士女子，使她成为以色列人，预先表明野橄榄枝被接在栽培的橄榄树上，同沾肥脂。因为按血肉出生的基督确实曾被人民追捕以杀害，但在埃及——即外邦人中——得救，圣化那里尚处幼年的人，并由此在那一方完成祂的教会（因为埃及从起初就是外邦之地，雇士亦然）；为此，借着梅瑟的婚姻，预表了圣言的婚姻；借着雇士新娘，显明了从外邦人中取来的教会；凡诋毁、控告、嘲笑她的人，必不洁净，因为他们要满身癞病，被逐出义人的营外。同样，妓女辣哈布虽自认有罪，作为外邦人，犯有一切罪过，却接待了三个侦察全地的侦探，将他们藏在家中；这三人无疑是圣父、圣子及圣神的预象。当整个城市在七支号角声中倒塌时，妓女辣哈布借着朱红绳的信德，保全了自己全家\Emph{in ultimis}；如同主对那些不接受祂来临的人（无疑就是法利塞人）所宣布的——他们否定了朱红绳的标记，这标记原指逾越节、人民的救赎和出离埃及——主说：\BibleQuote{税吏和娼妓要在你们以先进入天国}\BibleRef{玛 21:31}。\par

\SourceRecord{Translated by Alexander Roberts and William Rambaut.；Ante-Nicene Fathers；Vol. 1.；Edited by Alexander Roberts, James Donaldson, and A. Cleveland Coxe.；Buffalo, NY: Christian Literature Publishing Co.,；1885.；<http://www.newadvent.org/fathers/0103420.htm>.}

% structure: p0001 source=h1 target=section reason=page-title

\section{驳斥异端（第四卷，第二十一章）}

1. 但宗徒在致迦拉达人书里充分教导说，我们的\OriginalTerm{信德}{faith}也在亚巴郎身上预先被表明，而且他是我们\OriginalTerm{信德}{faith}之父，可以说是我们\OriginalTerm{信德}{faith}的先知。他说：\BibleQuote{那赐与你们圣神，并在你们中间施行奇迹的，是由于遵行法律，还是由于听\OriginalTerm{信德}{faith}呢？正如亚巴郎信了天主，这事就为他算作\OriginalTerm{成义}{righteousness}。所以你们要知道，具有\OriginalTerm{信德}{faith}的人，才是亚巴郎的子孙。圣经预先看见，天主将使外邦人因\OriginalTerm{信德}{faith}而成义，就预先向亚巴郎传报说：\InnerQuote{万民都要因你获得祝福。}因此，具有\OriginalTerm{信德}{faith}的人，必与有\OriginalTerm{信德}{faith}的亚巴郎一同蒙受祝福。}\BibleRef{迦 3:5}为此，他宣称这人不仅是\OriginalTerm{信德}{faith}的先知，也是那些从外邦人中相信耶稣基督者的父亲，因为他的\OriginalTerm{信德}{faith}和我们的\OriginalTerm{信德}{faith}是同一个：他因天主的应许，相信未来的事如同已经成就；同样，我们也因天主的应许，藉着\OriginalTerm{信德}{faith}看见那在将来王国中为我们存留的[产业]。\par

2. 依撒格的历史也不是没有\OriginalTerm{预象特征}{symbolic character}的。因为在致罗马人书中，宗徒宣告说：\BibleQuote{再者，当黎贝加由我们的祖先依撒格一人怀了孕}，她从圣言得到了回答，\BibleQuote{为使天主选拔的计划坚定不移，不因行为，而因召选者，就有话对她说：\InnerQuote{两个民族在你腹中，两种人民要从你身内分出；这一民族要强过那一民族，大的要服事小的。}}\BibleRef{罗 9:10}由此明显可见，不仅祖先有\OriginalTerm{预言}{prophecy}，而且黎贝加所生的孩子也是对两个民族的\OriginalTerm{预言}{prediction}；一个确实较大，另一个较小；一个受束缚，另一个得自由；但两者都出自同一位父亲。我们的天主，唯一的天主，也是他们的天主，祂知道隐藏的事，在事情发生之前早已知道一切；因此祂说：\BibleQuote{我爱了雅各伯，却恨了厄撒乌。}\BibleRef{罗 9:13}\par

3. 若有人再细察雅各伯的行为，就会发现它们并非没有含义，而是对\OriginalTerm{救恩计划}{dispensations}充满意义。首先，在他出生时，因他抓住了他哥哥的脚跟，\BibleRef{创 25:26}他就被称为雅各伯，即\Emph{\OriginalTerm{取代者}{supplanter}}——抓住却未被抓住；绑住脚却未被绑住；奋斗并得胜；手中紧握对手的脚跟，即胜利。因为主为此而生，他预先表明了他的出生方式，关于这位主，若望在默示录中说：\BibleQuote{祂骑了出去，胜了又胜。}\BibleRef{默 6:2}其次，雅各伯在他哥哥轻视长子的权利时获得了长子权；同样，年轻的民族在年长的民族拒绝祂，说\BibleQuote{除了凯撒，我们没有君王}\BibleRef{若 19:15}时，接受了祂，基督，首生子。但在基督里一切祝福都聚集了，因此后来的民族从父亲那里夺走了前者的祝福，正如雅各伯夺走了厄撒乌的祝福。因此，他的兄弟遭受了兄弟的迫害和阴谋，正如教会从犹太人那里遭受同样的事。在异乡，十二支派诞生了，以色列民族，因为基督也是在异乡要生出教会的十二柱石基础。各种颜色的羊被分配给雅各伯作为工资；基督的工资是来自各种不同民族聚集到同一\OriginalTerm{信德}{faith}队伍中的人，正如圣父应许祂说：\BibleQuote{你向我请求，我必将万民赐你作产业，将八极赐你作领土。}正如从他的众多儿子中，主[后来]的先知兴起，雅各伯必须从两个姊妹生子，正如基督从同一圣父的两条法律生子；同样也从婢女生子，表明基督要从自由人和按肉体为奴的人中生出天主的子女，同样赐给所有人赋予生命的圣神之恩赐。但他（雅各伯）所做的一切都是为了那有秀丽眼睛的少女辣黑耳，她预象了教会，基督为她忍耐；那时，祂藉着祂的祖先和先知，预先预表和宣告未来的事，在\OriginalTerm{救恩计划}{dispensations}中提前完成祂的部分，并训练祂的子民服从天主，在世上如同在旅途中漂泊，跟随祂的圣言，并预先指明未来的事。因为在天主那里，没有一事是没有意义或恰当象征的。\par

\SourceRecord{Translated by Alexander Roberts and William Rambaut.；Ante-Nicene Fathers；Vol. 1.；Edited by Alexander Roberts, James Donaldson, and A. Cleveland Coxe.；Buffalo, NY: Christian Literature Publishing Co.,；1885.；<http://www.newadvent.org/fathers/0103421.htm>.}

% structure: p0001 source=h1 target=section reason=page-title

\section{驳斥异端（卷四，第二十二章）}

1. 在末世，当自由时期的圆满到来时，圣言亲自用自己的手洗了门徒的脚，以此自身\BibleQuote{涤除熙雍女子的污秽}\BibleRef{依 4:4}，\BibleRef{若 13:5}。因为这就是人类继承天主的终向：正如起初，藉着我们的原祖父母，我们都因沦为死亡的奴隶而被奴役；同样，在末期，藉着新人，所有从起初就是祂门徒的人，既已被洗净和涤除了属于死亡的事物，就应进入天主的生命。因为那位洗了门徒的脚的人圣化了整个身体，并使之洁净。为此，祂也以斜倚的姿势给他们食物，表明那些躺卧在尘土中的人正是祂来赋予生命的人。正如耶肋米亚宣布：\Quote{圣洁的上主记起了祂死去的以色列人，他们沉睡在坟墓之地；祂下到他们那里，向他们宣告祂的救恩，好使他们得救。}也正为此，当基督的苦难临近时，门徒的眼睛沉重；起初，主发现他们睡着时，祂任凭他们——以此表明天主对人类沉睡状态的忍耐；但祂第二次来临时，唤醒了他们，使他们站起，以此象征祂的苦难是唤醒沉睡的门徒，为了他们的缘故，\BibleQuote{祂也下到地下深处}\BibleRef{弗 4:9}，亲眼看见那些从劳苦中安息的人的状况。关于这些人，祂也对门徒说：\BibleQuote{许多先知和义人渴望看见你们所看见的，听见你们所听见的。}\BibleRef{玛 13:17}\par

2. 因为基督来临，不仅仅是为了那些在提庇留凯撒时代信仰祂的人，也并非圣父只为现在活着的人行使祂的眷顾，而是为了从起初以来所有世代中，按各自的能力敬畏并爱慕天主、对邻人行正义与虔敬、并切望看见基督、听见祂声音的所有人。因此，在祂第二次来临时，祂将首先唤醒所有这类人，使他们从睡眠中起身，连同其余将要受审判的人一起，并赐给他们祂国度中的位置。因为确实\Quote{是天主}引导圣祖们走向祂的救恩计划，且\BibleQuote{使受割礼者因信德成义，也使未受割礼者藉着信德成义。}\BibleRef{罗 3:30}因为如同我们被预表在起初的他们身上，同样，他们也被代表在我们身上，即在教会中，并获得他们所作所为的报偿。\par

\SourceRecord{Translated by Alexander Roberts and William Rambaut.；Ante-Nicene Fathers；Vol. 1.；Edited by Alexander Roberts, James Donaldson, and A. Cleveland Coxe.；Buffalo, NY: Christian Literature Publishing Co.,；1885.；<http://www.newadvent.org/fathers/0103422.htm>.}

% structure: p0001 source=h1 target=section reason=page-title

\section{驳斥异端（第四卷，第二十三章）}

1. 为此，主对门徒说：\BibleQuote{我告诉你们：举目观看\OriginalTerm{田地}{regiones}，它们已经白了，可以收割了。收割的人得工价，积蓄五谷到永生，叫撒种的和收割的一同快乐。在这事上，那句\InnerQuote{那人撒种，这人收割}的话是真的。我差你们去收割你们没有劳苦的；别人劳苦了，你们享受他们所劳苦的。}\BibleRef{若 4:35}那么，谁曾劳苦，并协助了天主的救世计划呢？显然是那些圣祖和先知，他们甚至预表了我们的信德，并通过大地传播了天主之子的来临——祂是谁，祂应如何——使得后代，拥有对天主的敬畏，藉着先知的教导，能轻易接受基督的来临。为此，当若瑟得知玛利亚怀孕，并打算暗中休退她时，天使在梦中向他说：\BibleQuote{不要怕，娶你的妻子玛利亚来，因那在她内受孕的是出于圣神。她要生一个儿子，你要给祂起名叫耶稣，因为祂要把自己的民族从他们的罪恶中拯救出来。}\BibleRef{玛 1:20}并劝勉他，补充说：\BibleQuote{这一切事的发生，是为应验主藉先知所说的话：\InnerQuote{看，一位贞女将怀孕生子，人要称祂的名字为厄玛奴耳。}}从而藉先知的话说服他，免去玛利亚的责难，指出她正是从前依撒意亚所预言的那位将生厄玛奴耳的贞女。因此，当若瑟完全信服后，他不仅娶了玛利亚，而且喜乐地服从了关于基督抚养的一切其余事宜，前往埃及再返回，然后迁居纳匝肋。【为此，】那些既不知道圣经，也不知道天主的应许，更不知道基督的救恩计划的人，最后竟称他为孩子的父亲。为此，主自己也曾在葛法翁诵读依撒意亚的预言：\BibleRef{路 4:18}\BibleQuote{主的神在我身上，因为祂用油膏了我，叫我传福音给贫穷的人，差遣我医好伤心的人，报告被掳的得释放，瞎眼的得看见。}\BibleRef{依 61:1}同时，表明祂就是依撒意亚先知所预言的那位，祂对他们说：\BibleQuote{今天这经应验在你们耳中了。}\par

2. 为此，斐理伯遇见埃塞俄比亚女王的太监正在诵读这段经文：\BibleQuote{祂被牵到宰杀之地，又像羊羔在剪毛者前无声，祂也不开口；在祂屈辱时，祂的审判被夺去；}\BibleRef{宗 8:27}以及先知接着所说的关于祂的苦难和祂以肉身来临的一切事，以及祂如何被那些不信祂的人所羞辱；很容易就说服他相信，祂就是基督耶稣，在般雀比拉多手下被钉十字架，承受了先知所预言的一切，并且祂是天主子，赐予人永生。当斐理伯给他付洗后，就立刻离开了他。因为对这位已经受过先知教导的人来说，所缺的只是圣洗圣事：他并非不认识天主圣父，也不缺乏有关正常生活方式的规诫，只是不知道天主之子的来临；当他得知这一点后，在短时间内，他就欢欢喜喜地上路，到埃塞俄比亚去宣讲基督的来临。因此，斐理伯在此人身上并没有费什么大功夫，因为他已藉着先知在敬畏天主中预备好了。为此，宗徒们召集以色列家迷失的羊，从圣经与他们谈论，证明这位被钉死的耶稣就是基督，永生天主之子；他们说服了一大群人，这些人本来已怀有对天主的敬畏。于是在一天之内，有三千、四千、五千人受了洗。\BibleRef{宗 2:41}\par

\SourceRecord{Translated by Alexander Roberts and William Rambaut.；Ante-Nicene Fathers；Vol. 1.；Edited by Alexander Roberts, James Donaldson, and A. Cleveland Coxe.；Buffalo, NY: Christian Literature Publishing Co.,；1885.；<http://www.newadvent.org/fathers/0103423.htm>.}

% structure: p0001 source=h1 target=section reason=page-title

\section{\WorkTitle{驳斥异端}（第四卷，第24章）}

因此，\Person{保禄}既是外邦人的宗徒，就说：\Quote{我比他们众人更劳苦。} \BibleRef{格前 15:10} 对前者的教导[即犹太人]是容易的，因为他们可以从圣经中引证，并且他们素常听\Person{梅瑟}和先知，也很容易接受死者中的首生者和天主生命之元首——祂藉着伸开双手消灭了\Place{阿玛肋克}，并藉着对祂的信德使被蛇咬伤的人活过来。正如我在前卷中所指出的，宗徒首先教导外邦人离开偶像的迷信，去崇拜唯一的天主，即天地万物的创造者、整个受造界的塑造者；并教导祂的圣子就是祂的圣言，藉祂创造了万有；祂在末世成为人中人，革新了人类，摧毁并征服了人类的仇敌，赐予祂手中的工程战胜敌对者的胜利。虽然受割损的人仍然不顺从天主的话，因为他们是藐视者，但他们先前受过教导：不可奸淫、不可行淫乱、不可偷窃、不可欺诈；凡损害邻人的事都是邪恶的，为天主所憎恶。因此他们也乐意戒除这些事，因为他们受过这样的教导。\par

但他们也必须教导外邦人同样的事：即这类行为是邪恶的、有害的、无益的，并对行这些事的人具有破坏性。因此，那领受了向外邦人的宗徒职务的人，比那些在受割损的人中宣讲天主之子的人更为劳苦。因为前者有圣经的帮助，主来临时既应验又成全了圣经所预言的；但是在这里，[就外邦人而言]，却有某种陌生的学识和一种新的教义[要接受，即]：万民的神祇不仅根本不是神，甚至还是魔鬼的偶像；并且只有一位天主，祂是\BibleQuote{超越一切率领者、掌权者、能力者、宰制者，以及一切可称呼的名号} \BibleRef{弗 1:21}；祂的圣言，本性不可见，却在人中间成为可触摸、可见的，并且下降\BibleQuote{致死，且死在十字架上} \BibleRef{斐 2:8}；此外，凡信祂的人将是不朽的、不受苦难的，并将承受天国。这些事也是藉言语传给了外邦人，没有[圣经]的帮助；因此，在外邦人中宣讲的人经历了更重的劳苦。但另一方面，外邦人的信德被证明更为高贵，因为他们没有来自[圣经]的教导而遵循了天主的话\OriginalTerm{没有文字的教导}{sine instructione literarum}。\par

\SourceRecord{Translated by Alexander Roberts and William Rambaut.；Ante-Nicene Fathers；Vol. 1.；Edited by Alexander Roberts, James Donaldson, and A. Cleveland Coxe.；Buffalo, NY: Christian Literature Publishing Co.,；1885.；<http://www.newadvent.org/fathers/0103424.htm>.}

% structure: p0001 source=h1 target=section reason=page-title

\section{《驳斥异端》（卷四，第25章）}

1. 因为亚巴郎的子孙本当如此，就是天主从这些石头中为亚巴郎兴起的子孙\BibleRef{玛 3:9}，并使他们与那成为我们信德的首领和先驱者（他在尚于未受割损时领受了因信德而成的义，之后才领受了割损的盟约，好叫两个盟约在他身上得以预表，使他成为所有追随天主圣言、在此世过着朝圣生活者的父亲，即那些出自割损和未受割损而有信德的人，正如 \BibleQuote{基督\BibleRef{弗 2:20}是主要的角石} 支撑万物）的人为邻；祂将那些从两个盟约中堪当为天主建筑的人，聚集于亚巴郎的唯一信德中。但这未受割损的信德，因其将终始相连，已成为最先的和最后的。因为，如我已表明的，它在亚巴郎受割损之前就已存在，在其余所有中悦天主的义人身上亦然；而在这些末期，它又藉着主的来临在人类中重新兴起。但割损和行为的律法则占据了中间时期。\par

2. 这件事确实由许多其他事件表明，但尤以犹大的儿媳塔玛尔的事迹为预表。\BibleRef{创 38:28}因为当她怀了双胎时，其中一个先伸出手来；接生妇以为他是长子，就将一条朱红线拴在他手上。但这之后，他缩回手去，他的兄弟培勒兹却先出来；然后，那有朱红线的匝辣才出生为次子：圣经清楚指明那拥有朱红记号的人民，即处于割损状态的信德，这信德最初在族长们身上已预先显明；但之后被收回，好让他的兄弟出生；同样，那身为长者的，后来才出生，却因系在他身上的朱红线而受区别，即义人的苦难，这苦难在起初已由亚伯尔预表，由先知们描述，但在末期由天主子完成。\par

3. 因为必须有一些事由祖先以父亲的方式预先宣布，另一些由先知以律法的方式预表，还有一些由那些获得义子身份的人以基督的样式描述；而所有的事都在唯一天主内显明。因为虽然亚巴郎只是一位，他却在自己身上预表了两个盟约，其中有些人播种，另一些人收割；因为经上说：\Quote{这话是真的：\InnerQuote{一个民族}播种，另一个民族却要收割。}\BibleRef{若 4:37}但却是同一天主赐予两者合宜之物——给播种者种子，给收割者食粮。正如那栽种的和浇灌的是一样，但使之生长的是天主。\BibleRef{格前 3:7}因为族长和先知们播种了关于基督的圣言，教会则收割，即领受了果实。为此，这些（先知）自己也祈求在教会中有一居所，如耶肋米亚说：\Quote{谁使我在旷野中有一最后居所？}好使播种者和收割者都能在基督的国度中一同欢乐，基督与所有从起初就被天主认可、蒙祂赐予圣言与他们同在的人同在。\par

\SourceRecord{Translated by Alexander Roberts and William Rambaut.；Ante-Nicene Fathers；Vol. 1.；Edited by Alexander Roberts, James Donaldson, and A. Cleveland Coxe.；Buffalo, NY: Christian Literature Publishing Co.,；1885.；<http://www.newadvent.org/fathers/0103425.htm>.}

% structure: p0001 source=h1 target=section reason=page-title

\section{反驳异端（第四卷，第二十六章）}

1. 因此，任何人若用心阅读圣经，必在其中找到关于基督的记述以及\OriginalTerm{新召叫}{vocationis}的预象。因为基督就是藏在田地里的宝藏\BibleRef{玛 13:44}，那田地就是这世界（因为\BibleQuote{田地就是世界}\BibleRef{玛 13:38}）；但藏在圣经里的宝藏就是基督，祂藉着预象和比喻被预示出来。因此，在那些预言之事的完成，即基督来临之前，祂的人性是无法被理解的。所以，先知达尼尔曾被告知：\BibleQuote{你要封闭这些话，将书卷封起，直到终结之时，使多人学习，知识得以圆满。因为那时，当分散完成，他们必明白这一切。}\BibleRef{达 12:4}耶肋米亚也说：\BibleQuote{在最后的日子，他们必明白这些事。}\BibleRef{耶 23:20}因为一切预言在实现之前，对人而言如同谜语和歧义。但当时刻到来，预言应验之后，预言便有了清晰而确定的解释。正因如此，如今当法律读给犹太人听时，它就像一则寓言；因为他们不掌握关于天主圣子降生成人之事的一切解释；但基督徒阅读时，它便成了宝藏，确实藏在田里，却藉着基督的十字架显明出来，并得到解释，既充实人的悟性，又彰显天主的智慧，宣告祂对人的计划，预构基督的王国，预先宣告圣耶路撒冷的产业，预告那爱天主的人将达到如此崇高的境界，甚至能看见天主，听见祂的话语，并因聆听祂的言论而荣耀至极，以致他人不能目睹他面容的光荣，正如达尼尔所说：\BibleQuote{明慧者必发光如同穹苍的光辉，许多义人必如星辰，直到永远。}\BibleRef{达 12:3}因此，我已表明，若有人用心阅读圣经，即是如此。因为主在复活后就是这样教导门徒，从圣经本身向他们证明\BibleQuote{基督必须受苦，进入祂的光荣，并且因祂的名，将赦罪之道传遍普世。}于是门徒得以成全，如同家主，\BibleQuote{从自己的宝库里取出新旧的东西。}\par

2. 因此，必须服从教会中的司铎——那些如我所表明的，从宗徒继承下来的司铎；他们与主教职位继承一起，按照圣父的美意，领受了真理的确定恩赐。但[也应当]怀疑那些背离原始传承，在任意地方聚集的人，[视他们]或为心怀邪僻的异端者，或为骄傲自满的分裂者，或为贪图利益和虚荣而行此事的伪善者。因为所有这些人都偏离了真理。异端者将异火带到天主的祭台上——即异端邪说——必被天上降下的火烧灭，正如纳达布和阿彼乌\BibleRef{肋 10:1}。但那些起来反对真理，并劝诱他人反抗天主教会的人，[将]留在\OriginalTerm{地狱}{apud inferos}之中，被地裂吞噬，如科辣黑、达堂和阿彼兰同伙\BibleRef{户 16:33}。而那些分裂并割裂教会合一的人，[将]受到来自天主与雅洛贝罕相同的惩罚\BibleRef{列上 14:10}。\par

3. 然而，那些被许多人视为司铎，却顺从自己的私欲，不将天主的敬畏置于心中，反而轻蔑他人，因占据首位而骄傲自负，并在暗中行恶，说\Quote{没有人看见我们}的人，将被那不看\OriginalTerm{外貌}{secundum gloriam}、不凭面容而鉴察内心的圣言定罪；他们将听到达尼尔先知的话：\BibleQuote{你这客纳罕的种子，而非犹大的种子，美貌欺骗了你，情欲败坏了你的心。你在罪恶中衰老，如今你从前所犯的罪行暴露了；因为你作出了不义的判决，习惯定无辜者有罪，释放有罪者，尽管主说：无辜和正义的人，你不可杀害。}主也曾论及他们说：\BibleQuote{但若那恶仆心里说：我的主人必来得迟，便开始殴打仆婢，也吃也喝，也沉醉；那仆人的主人要在想不到的日子，不知道的时辰来到，把他斩成两段，分派他与不信者同列。}\par

4. 因此，我们应当远离所有这样的人，而依附那些——正如我已指出的——持守宗徒教训，并连同\OriginalTerm{长老品秩}{presbyterii ordine}，以纯全的言论和无瑕的行为来巩固和矫正他人的人。这样，梅瑟——被托付如此领导之职的人——凭著良心的无愧，在天主面前为自己辩护说：\BibleQuote{我没有贪图这些人的任何东西，也没有对他们行过恶。}\BibleRef{户 16:15}同样，撒慕尔——他审判百姓多年，无骄傲地治理以色列——最终也为自己辩护说：\BibleQuote{我从幼年直到今日，走在你们面前：请在天主面前和\OriginalTerm{祂的受傅者}{Christi ejus}面前答复我；我拿了谁的牛？拿了谁的驴？我压迫过谁？亏负过谁？若我从任何人手中接受过贿赂或哪怕一只鞋，请当面控告我，我必偿还给你们。}\BibleRef{撒上 12:3}当百姓对他说：\BibleQuote{你没有压迫我们，没有亏负我们，也没有从任何人手中拿过任何东西}，他便呼求主作证说：\BibleQuote{主今日作见证，祂的受傅者也作见证，你们在我手中没有找到任何东西。他们便对他说：祂是见证。}同样，宗徒保禄，因怀有无愧的良心，对格林多人说：\BibleQuote{因为我们不像许多人那样，混淆天主的圣言；而是出于真诚，出于天主，在基督内当着天主的面讲论；}\BibleRef{格后 2:17}\BibleQuote{我们没有亏负过任何人，没有败坏过任何人，也没有欺骗过任何人。}\BibleRef{格后 7:2}\par

5. 教会养育这样的长老，先知论到他们说：\BibleQuote{我要赐给你和平的统治者，正义的主教。}\BibleRef{依 60:17}主也曾论到他们说：\BibleQuote{那么，谁是那忠信又精明的\OriginalTerm{管家}{actor}，主人派他管理自己的家仆，按时分粮给他们呢？主人来时，看见他这样做，那仆人才是有福的。}\BibleRef{玛 24:45}保禄教导我们在何处能找到这样的人，说：\BibleQuote{天主在教会中设立的：第一是宗徒，第二是先知，第三是教师。}\BibleRef{格前 12:28}因此，主的恩赐放在何处，我们就当在那里学习真理，即从那些拥有从宗徒传下来的教会继承者，且在他们中间有行为纯正无瑕、言论纯全不朽的人学习。因为这些人也保持我们对于创造万有的独一天主的信德；他们增进我们对天主圣子的爱，祂为我们成就了如此奇妙的救世计划；他们毫无危险地为我们诠释圣经，既不亵渎天主，也不侮辱圣祖，更不轻视先知。\par

\SourceRecord{Translated by Alexander Roberts and William Rambaut.；Ante-Nicene Fathers；Vol. 1.；Edited by Alexander Roberts, James Donaldson, and A. Cleveland Coxe.；Buffalo, NY: Christian Literature Publishing Co.,；1885.；<http://www.newadvent.org/fathers/0103426.htm>.}

% structure: p0001 source=h1 target=section reason=page-title

\section{驳斥异端（第四卷，第27章）}

1. 我从一位长老那里听说，他从见过宗徒的人以及宗徒的门徒那里听说，圣经中所宣告的惩罚，对于古人在未受圣神引导时所行的事，已足够了。因为天主不偏待人，祂对不悦乐祂的行为施以适当的惩罚。正如达味的情况，\EditorialNote{撒慕尔纪上18}他为了义而受撒乌耳的迫害，逃离撒乌耳王，不肯报复他的仇敌；他歌颂基督的来临，以智慧教训万民，凡事随从圣神的引导，中悦天主。但当他的情欲促使他娶了乌黎雅的妻子巴特舍巴时，圣经论及他说：\Quote{现在，达味所行的\OriginalTerm{那事}{sermo}，上主眼中看为恶。}\BibleRef{撒下11:27}于是纳堂先知被派到他那里，向他指明他的罪行，为使他自我定罪，从而获得基督的怜悯和宽恕：\BibleQuote{纳堂对他说：在一座城里有两个人，一富一贫。富人有许多牛羊；穷人除了一只小母羊羔外，一无所有。他喂养它，它与他和他的儿女一同长大，吃他的食物，喝他的杯，睡在他怀中，如同他的女儿一样。有一个客人来到富人那里，他舍不得从自己的羊群中取一只预备给客人，却取了穷人的母羊羔，款待了来客。达味就十分恼怒那人，对纳堂说：上主永在，行这事的人\OriginalTerm{该死}{filius mortis est}，并且赔偿四倍，因为他行了这事，且没有怜恤穷人。纳堂对达味说：你就是那人。}\BibleRef{撒下12:1}接着他继续责备他，历数天主对他的恩惠，指明他的行为何等令上主不悦。因为这类行为不中悦天主，反使他的家遭受大忿怒。达味听了这话，心中懊悔，喊道：\Quote{我得罪了上主。}于是他唱了一首忏悔的圣咏，等候主的来临，祂要洗涤并洁净那被罪恶捆绑的人。同样，撒罗满也是如此。当他秉公行义，宣扬天主的智慧，建造了作为真理预像的圣殿，彰显天主的光荣，宣告将要临到万民的和平，预像基督的国度，讲论关于主来临的三千比喻，写作赞美天主的五千首歌，阐释天主在造化中的智慧，论及各种树木、草本植物、飞禽、走兽和鱼类；他说：\BibleQuote{天与天上的天尚且容不下天主，祂岂真愿与人同住在地上？}\BibleRef{列上8:27}他中悦天主，为众人所钦佩；地上的众王都\OriginalTerm{寻求与他见面}{quærebant faciem ejus}，为听天主赐予他的智慧。\BibleRef{列上4:34}南方的女王也从地极而来，为明了其中的智慧。\BibleRef{列上10:1}主也曾提到她，说在审判时她要与那些听见祂的话却不相信的世代一同起来定他们的罪，因为她顺从了由天主的仆人传报的智慧，而这些人却轻看那直接出于天主圣子的智慧。因为撒罗满是仆人，基督却是天主子，是撒罗满的主。因此，当他无过地事奉天主，为祂的救恩计划尽职时，他就获得光荣；但当他从万民中娶妻，容许她们在以色列中设立偶像时，圣经就如此论及他说：\BibleQuote{撒罗满王爱许多外邦女子，并娶了她们。当撒罗满年老时，他的心没有完全归于上主他的天主。外邦女子诱惑了他的心去敬拜别的神。撒罗满行了上主眼中看为恶的事，没有像他父亲达味那样全心随从上主。上主向撒罗满发怒，因为他的心没有完全归于上主，不像他父亲达味的心。}\BibleRef{列上11:1}圣经就这样充分责备了他，正如那位长老所说的，为叫凡有血肉的，不得在天主面前自夸。\par

2. 主也下降到地下区域，在那里宣讲祂的来临，并宣告凡信祂的人都得到罪过的赦免。\BibleRef{伯前 3:19} 凡是对祂怀着希望的人，就是那些宣讲祂来临、顺服祂救恩计划的人，即义人、先知和族长们，都相信了祂。祂同样赦免了他们的罪，如同赦免我们一样；我们不应将罪过归咎于他们，免得我们轻视天主的恩宠。因为正如这些人没有将我们（外邦人）在基督显现于我们中间之前所犯的过犯归咎于我们，同样我们也不应将过错加于那些在基督来临前犯罪的人。因为\BibleQuote{众人都不及天主的荣耀}，他们并非因自己而成义，而是因主的来临而成义——就是那些目光专注地朝向祂之光的人。他们的行为被记录下来是为了教训我们，首先要知道，我们的天主和他们的天主是同一位，罪过不蒙祂喜悦，即使是名人所犯的罪；其次，我们应远离邪恶。因为如果这些古时的先贤，在所赐的恩赐上领先于我们，而天主子尚未为他们受苦，当他们犯罪、顺从肉欲时，就蒙受如此耻辱，那么现今轻视主的来临、成为自己情欲奴隶的人将遭受何等痛苦呢？诚然，主的死亡成为前者医治和罪赦的途径，但基督不会为现今犯罪的人再死，因为死亡不再统治祂；而是圣子要在父的光荣中来临，向祂的管家和分施者连本带利索回委托给他们的钱财；且祂向谁给的多，向谁要的也多。因此，正如那位长老所评论的，我们不应骄傲，也不应苛责古人，而应自己恐惧，免得在认识基督之后，若做了天主不悦的事，就不再得到罪的赦免，反而被摒于祂的王国之外。因此保禄说：\BibleQuote{因为天主若不怜惜原来的枝条，也要当心，祂也不怜惜你；你原是野橄榄，被接在橄榄树的肥汁中，成了它肥汁的分享者。}\par

3. 你也要注意，类似地，平民的过犯也被记载下来，并非为了当时犯罪之人的缘故，而是为教训我们，使我们明白，这些人所冒犯的天主，与现今那些自称信祂的人中犯罪者所冒犯的是同一位。但这一点，正如那位长老所说的，保禄也在\WorkTitle{致格林多人书}中极为清楚地宣告了，他说：\BibleQuote{弟兄们，我不愿意你们不知道，我们的祖先都曾在云柱下，都从海中经过，都曾在云中和海中受洗归于梅瑟，都吃过同样的神粮，都喝过同样的神饮；因为饮于跟随他们的神石，那神石就是基督。但在他们中，天主并不喜欢多数人，因而他们倒毙在旷野。这些事作为我们的\Emph{\OriginalTerm{预像}{in figuram nostri}}，为叫我们不要贪恋恶事，像他们贪恋过一样；也不要拜偶像，像他们中有些人，正如经上所记载：\BibleRef{出 32:6}\InnerQuote{百姓坐下吃喝，起来玩乐。}我们也不要行奸淫，像他们中有些人行过，一天就倒毙了两万三千人；也不要试探基督，像他们中有些人试探过，就被蛇所灭；也不要抱怨，像他们中有些人抱怨过，就被毁灭者所灭。这一切发生在他们身上，是作为预像，并写下来为警戒我们这到了\Emph{\OriginalTerm{世代}{sæculorum}}的终极的人。所以，自以为站得稳的，要小心，免得跌倒。} \BibleRef{格前 10:1}\par

4. 因此，既然宗徒毫不怀疑、毫无矛盾地表明，有同一的天主，祂既曾审判那些先前的事，也查究现今的事，并指出为何这些事被记载下来；那么，那些因古时之人的过犯和众多之人的悖逆，而声称这些人的天主有一位，且祂是世界的造主，处于堕落状态；但还有另一位圣父由基督宣告，而这位神祇是每个人凭自己的心志所构想出来的；他们不明白，正如在前一种情况下，天主在许多事例中对犯罪者表示不悦，同样在后一种情况下，\BibleQuote{被召的人多，被选的人少。}\BibleRef{玛 20:16}正如那时不义的人、拜偶像的人和行淫的人灭亡了，现在也是如此：因为主宣称，这样的人要被投入永火；\BibleRef{玛 25:41}宗徒也说：\BibleQuote{你们不知道不义的人不得继承天主的国吗？不要被迷惑：无论是淫乱的、拜偶像的、奸淫的、作娈童的、与男人苟合的、偷窃的、贪婪的、酗酒的、辱骂的、勒索的，都不能继承天主的国。}\BibleRef{格前 6:9}正如这些话他并不是对外面的人说的，而是对我们说的，免得我们因行这样的事而被逐出天主的国，他接着说：\BibleQuote{你们从前也的确是这样，但你们已被洗净，因主耶稣基督的名和因我们天主的圣神，成了圣洁的。}正如那时，那些行为恶劣、引人误入歧途的人被定罪并被逐出，同样现在，犯罪的眼、脚、手也要被砍掉，免得其余肢体也照样丧亡。\BibleRef{玛 18:8}我们还有这样的训诫：\BibleQuote{若有人称为弟兄，却是淫乱的、贪婪的、拜偶像的、辱骂的、酗酒的、勒索的，这样的人，不可与他一同吃饭。}\BibleRef{格前 5:11}宗徒又说：\BibleQuote{不要让任何人以空言欺骗你们，因为为了这些事，天主的义怒临到悖逆之子身上。所以你们不要与他们同流合污。}\BibleRef{弗 5:6}正如那时罪人的定罪也涉及那些赞同他们、与他们交往的人，现在同样如此：\BibleQuote{一点面酵能使整个面团发起来。}\BibleRef{格前 5:6}正如那时天主的义怒降在不义的人身上，同样，宗徒也说：\BibleQuote{天主的义怒从天上显明在一切不虔敬和不义的人身上，就是那些以不义压制真理的人。}\BibleRef{罗 1:18}正如那时，天主的报复临到那些不义地惩罚以色列的埃及人身上，现在也是如此：主真实地说：\BibleQuote{天主岂不替祂所拣选、日夜向祂呼号的人伸冤吗？我告诉你们，祂必要快快为他们伸冤。}\BibleRef{路 18:7}同样，宗徒在写给得撒洛尼人的书信中也说：\BibleQuote{既然天主是公义的，就必将患难报复给那加患难给你们的人；并将平安赐给你们这些受患难的人，以及我们，当主耶稣从天上同祂有能力的使者显现时，在火焰中报复那些不认识天主和那不听从我们主耶稣福音的人。他们要受惩罚，永远丧亡，远离主的面和祂威能的光荣，当祂在圣徒身上受光荣，在一切信祂的人身上受赞美的那一天。}\BibleRef{得后 1:6}\par

\SourceRecord{Translated by Alexander Roberts and William Rambaut.；Ante-Nicene Fathers；Vol. 1.；Edited by Alexander Roberts, James Donaldson, and A. Cleveland Coxe.；Buffalo, NY: Christian Literature Publishing Co.,；1885.；<http://www.newadvent.org/fathers/0103427.htm>.}

% structure: p0001 source=h1 target=section reason=page-title

\section{《驳斥异端》（第四卷，第二十八章）}

1. 因此，既然在两约中，当天主施行报复时，显示的是同一天主的正义，在一方是象征性的、暂时的、较为温和的，而在另一方则是真实的、持久的、更为严厉的；因为火是永恒的，以及那将从天上、从主的面容（正如达味也说：\BibleQuote{但是主的面容敌对行恶的人，要由地上剪除他们的记念}）彰显的天主义怒，给那些招致它的人带来更重的惩罚——长老们指出，那些缺乏理智的人，他们根据从前不服从天主之人所遭遇的事，试图引入另一位父，以主在来临时为拯救那些接受祂、怜悯他们的人所行的伟大之事来对抗这些惩罚，却对祂的审判保持沉默；以及所有那些将要临到听见祂的话却不实行的人身上的事，并且他们不出生倒好\BibleRef{玛 26:24}，在审判时，索多玛和哈摩辣所受的比那城更容易忍受，因为那城没有接受祂门徒的话\BibleRef{玛 10:15}。\par

2. 因为在\WorkTitle{新约}中，人对于天主的信德得以增强，加之于已经启示之上，圣子被赐予，使人也能分享天主；因此，我们的生活方式也需要更加谨慎，我们被吩咐不仅戒绝恶行，还要戒绝恶念、闲言、空谈和诽谤的话语；同样，那些不信从天主圣言、轻视祂来临、向后转身之人的惩罚也加重了，不再是暂时的，而是成为永恒的。因为主将对谁说什么：\BibleQuote{可咒骂的，离开我，到永火里去}\BibleRef{玛 25:41}，这些人将永远被罚；主将对谁说什么：\BibleQuote{我父所祝福的，来承受为你们预备的国}\BibleRef{玛 25:34}，这些人将永远承受这国，并在其中不断进步；因为同一天主圣父，以及祂的圣言，始终与人类同在，通过各种不同的安排，成就了许多事，并从一开始就拯救了被拯救者（因为他们爱天主，并按他们的种类跟随天主的圣言），并审判了被审判者，即那些忘记天主、亵渎并违反祂圣言的人。\par

3. 因为我们已经提及的那同一批异端者，因控告他们自称相信的主，而脱离了自身。因为他们所关注的那些关于天主当时对不信者施以暂时惩罚、击打埃及人、同时拯救顺从者的事实，同样也将在主身上重复，祂永恒地审判祂所审判的人，永恒地释放祂所释放的人；并且根据这些人所用的语言，祂将被发现是那些对祂下手、刺透祂之人的最严重罪的起因。因为如果祂没有这样来，那么这些人就不可能成为他们主的杀害者；如果祂没有派遣先知到他们那里，他们肯定也不能杀害先知，也不能杀害宗徒。因此，对攻击我们的人说：如果埃及人没有遭受灾祸，并在追赶以色列时淹死在海中，天主就不能拯救祂的子民，我们可以这样回答——除非犹太人成为主的杀害者（这确实从他们那里夺去了永生），并通过杀害宗徒和迫害教会而陷入愤怒的深渊，我们就不能得救。因为正如他们藉着埃及人的盲目而得救，同样我们藉着犹太人的盲目而得救，如果确实主被钉十字架是那些将祂钉在十字架上、不信祂来临之人的定罪，却是相信祂之人的救恩。因为宗徒在\WorkTitle{格林多后书}中也说：\BibleQuote{因为我们就是基督的馨香，在得救的人中，也于丧亡的人中：对后者，这是致死的香气致死；对前者，这是生命的香气致生。}\BibleRef{格后 2:15}那么，对谁有致死的香气致死，若不是那些不信从也不顺从天主圣言的人？甚至那时将自己交付死亡的是谁？无疑是不信从、不顺从天主的人。再者，得救并承受产业的是谁？无疑是那些信从天主并持守祂爱德的人，如耶孚乃的儿子加肋布和农的儿子若苏厄\BibleRef{户 14:30}，以及无辜的孩童，他们不知恶。如今得救并承受永生的是谁？岂不就是那些爱天主、相信祂的应许，并且\BibleQuote{在邪恶上成为像小孩子一样}\BibleRef{格前 14:20}的人吗？\par

\SourceRecord{Translated by Alexander Roberts and William Rambaut.；Ante-Nicene Fathers；Vol. 1.；Edited by Alexander Roberts, James Donaldson, and A. Cleveland Coxe.；Buffalo, NY: Christian Literature Publishing Co.,；1885.；<http://www.newadvent.org/fathers/0103428.htm>.}

% structure: p0001 source=h1 target=section reason=page-title

\section{\WorkTitle{驳斥异端（第四卷，第二十九章）}}

1. \Quote{但是}，他们说，\BibleQuote{天主使法郎和他臣民的心硬。}\BibleRef{出 9:35} 那么，提出这类难题的人，并没有读到福音中主回答门徒的那段话，当时他们问祂：\BibleQuote{你为什么用比喻对他们讲话？}——\BibleQuote{因为天国的\OriginalTerm{奥迹}{mystery}，是给你们知道的，但给他们，我用比喻，叫他们看是看见，却不明白；听是听见，却不领悟；免得他们悔改，而得赦免。}\BibleRef{玛 13:11} 为应验依撒意亚关于他们的预言，说：\BibleQuote{你们听是听见，却不领悟；看是看见，却不明白。} 因为天主使这百姓的心变硬，使他们的耳朵发沉，使他们的眼睛昏迷，免得他们眼睛看见，耳朵听见，心里明白，而悔改，获得痊愈。 因为同一位天主［祝福别人］使那些不信的人、轻视祂的人瞎眼；正如祂所造的太阳，对于那些因眼睛有某种弱点而不能看见其光的人，也是如此；但对于那些相信祂、跟随祂的人，祂赐予更丰沛、更明亮的心智光照。因此，按照这句话，宗徒在致格林多人后书中说：\BibleQuote{今世的神已蒙蔽了不信者的心思，免得基督荣耀福音的光照著他们。}\BibleRef{格后 4:4} 又在致罗马人书中说：\BibleQuote{他们既然不肯认真认识天主，天主就任凭他们陷于邪僻的心思，去行那些不正当的事。}\BibleRef{罗 1:28} 说到假基督时，他在致得撒洛尼人后书中清楚地写道：\BibleQuote{为此，天主使一种错误的效力在他们身上，叫他们相信谎言，为使一切不信真理，而喜爱不义的人，都被定罪。}\BibleRef{得后 2:11}\par

2. 因此，即使在现今，天主知道那些将不信的人的人数，因为祂预知一切，祂就把他们交于不信，并向这类人转面不顾，把他们留在自己选择的黑暗中，那么，祂在那时也把从不相信的法郎及其随从交于他们的不信，又有什么奇怪呢？正如圣言从荆棘中对\Person{梅瑟}所说：\BibleQuote{我知道，除非用强力，埃及王不会让你们走。}\BibleRef{出 3:19} 正因为主用比喻讲话，并使以色列人眼瞎，使他们看是看见，却不明白，因为祂知道他们心中的不信，同样，祂也使法郎的心硬；为的是，法郎虽看见是天主的手指引领百姓出来，却仍不信，而沉入不信的海洋，执迷于认为这些［以色列人］的出走是靠魔术的能力成就的，而不是天主的作为；红海为百姓开辟道路，也不是由于天主的运作，而是\OriginalTerm{自然如此}{sed naturaliter sic se habere}的结果。\par

\SourceRecord{Translated by Alexander Roberts and William Rambaut.；Ante-Nicene Fathers；Vol. 1.；Edited by Alexander Roberts, James Donaldson, and A. Cleveland Coxe.；Buffalo, NY: Christian Literature Publishing Co.,；1885.；<http://www.newadvent.org/fathers/0103429.htm>.}

% structure: p0001 source=h1 target=section reason=page-title

\section{驳斥异端（卷四，第三十章）}

1. 有些人吹毛求疵，指责民众在天主命令下，于出发前夕从埃及人那里拿了各种器皿和衣服，并以此离去了；并且后来在旷野中用这些战利品建造了会幕。这些人表明自己对天主公义的作为和他的计划一无所知，正如长老所指出的：因为如果天主没有在这\OriginalTerm{预像性的出谷}{typical exodus}中如此安排，就没有人能在我们真正的出谷中得救；也就是说，在我们所被建立的信德中，并藉此信德我们从外邦人中被领出来。因为在我们中，有人跟随少量，有人跟随大量财产，这些财产是从\OriginalTerm{不义之财}{mammon of unrighteousness}中获得的。因为我们居住的房屋、穿着的衣服、使用的器皿，以及每日生活所需的一切，若非来自我们身为外邦人时以贪婪获取的，或从我们外邦的父母、亲属或朋友（他们不义地获得了这些）那里继承的，还能从哪里得来呢？更不用说我们现在在信德中也取得这些东西。因为谁做买卖而不想从买者那里获利呢？谁买东西而不想从卖者那里得到好处呢？谁经营生意而不为谋生呢？至于在王宫中的那些信徒，他们使用的器皿难道不是来自属于\Person{凯撒}的财产吗？而对于那些没有的，这些基督徒中的每一个人不是按自己的能力给予吗？埃及人不仅欠了以色列子民的财产，也欠了他们的生命，因为先祖\Person{若瑟}昔日的恩惠；但是外邦人在哪方面亏欠我们，以致我们从中得到利益和好处呢？他们劳苦积累的一切，我们却在信德中不劳而获地享用。\par

2. 直到那时，以色列子民仍处于对埃及人最卑微的奴役中，正如圣经所说：\BibleQuote{埃及人严厉地对待以色列子民，使他们因苦工而生活痛苦，就是和泥、做砖，以及田间各样的工作，凡他们所做的工，都严严地待他们。}\BibleRef{出 1:13}他们以极大的劳力为埃及人建造了设防的城市，在漫长的岁月中增加了这些人的财富，并透过各种奴役方式；而这些人不仅忘恩负义，甚至图谋彻底消灭他们。那么，以色列子民从众多物品中只取少量，难道有什么不义吗？他们若不曾服侍，本可拥有许多财产，并可能富足地离去；而实际上，他们只收到对自己沉重奴役的极其微不足道的补偿，便贫穷地离去了。这就好像一个自由人被另一个人强行掳走，服侍多年，增加了那个人的财富，最终得到一些支持时，却被人认为占有了主人财产的一小部分，但实际上他只是从自己巨大的劳动中得到了很少，而从巨大的已得财产中几乎一无所得地离去，而这竟被人拿来指控他行为不当。这个指控者反而显得是对那个被强行掳为奴隶的人不公正的审判者。这些人也是如此，他们指责以色列子民从众多物品中取了少量，却不对那些没有回报祖先服务的人提出任何指控；甚至，他们使以色列子民陷入最繁重的奴役，并从中获取了最大的利益。这些反对者声称以色列子民行为不诚实，因为他们拿走了少许未刻印的金银，正如我所说，装在几件器皿中作为他们劳动的报酬；而他们自己（让我们说实话，尽管有人可能觉得可笑）却声称行为诚实，因为他们从别人的劳动中，将刻有\Person{凯撒}的刻印和像的铸币金、银、铜，系在腰带上带走。\par

3. 然而，倘若在我们和他们之间作一比较，哪一方似乎更公平地接受了（他们的世物）？是（犹太）民族，他们（取了）埃及人（在各方面都是他们的债务人）的财物；还是我们，（从罗马人及其他民族那里接受财产），而这些民族对我们并无类似的义务？是的，此外，通过他们的工具，世界得享和平，我们无惧地行走在大道上，随意航行。因此，针对这种人（即异端者），吾主的话是适用的，它说：\BibleQuote{假善人，先从你眼中取出大梁，然后你才能清楚地取出你兄弟眼中的木屑。}\BibleRef{玛 7:5} 因为如果那个将这些事归咎于你、并自夸于自己智慧的人，已与外邦人的团体分离，且毫无（取自）他人财物的东西，而是实实在在地赤身、赤足，居无定所于山间，如同那些食草的野兽一样，那么他（说这样的话）就可以原谅，因为他对我们生活方式的需求一无所知。但若他也享受在人们看来属于他人的财物，并且（同时）贬低他们的典型，他就证明自己极其不义，将这种指责转向自己。因为他会被发现携带不属于自己的财物，并贪图不属于他的东西。因此主说：\BibleQuote{你们不要判断，免得你们受判断：因为你们用什么判断来判断，你们也要受判断。}\BibleRef{玛 7:1}（意思）当然不是说我们不应指责罪人，也不是要我们同意作恶的人，而是说我们不应不公正地判断天主的安排，因为祂自己已预作安排，使一切事以符合正义的方式变为美好。因为，祂知道我们会善加运用我们从他人那里接受的财产，祂说：\BibleQuote{有两件衣裳的，要分给那没有的；有食物的，也要照样做。}\BibleRef{路 3:11} 又说：\BibleQuote{因为我饿了，你们给了我吃的；我渴了，你们给了我喝的；我赤身露体，你们给了我穿的。}\BibleRef{玛 25:35} 又说：\BibleQuote{当你施舍时，不要让你的左手知道右手所做的。}\BibleRef{玛 6:3} 我们因任何其他善行而被证明是义人，如同从陌生人的手中赎回我们的财产。但我这样说：\Quote{从陌生人的手中，}并非因为世界不属于天主的财产，而是说我们有这类的恩赐，并从他人那里接受它们，正如这些人从不认识天主的埃及人那里得到它们一样；并借着这些同样的恩赐，我们在自己内建立了天主的帐幕：因为天主住在行为正直的人内，如主所说：\BibleQuote{你们要用不义的钱财结交朋友，以便在你们失势时，他们接待你们进入永恒的帐幕。}\BibleRef{路 16:9} 因为凡我们在外教人时从不义中获取的，当我们成为信徒时，因着把它们用于主的益处，就被证明是义人。\par

4. 因此，很自然地，这些事是预先在预象中实行的，并由此建造了天主的帐幕；那些人公平地接受了它们，正如我已经表明的，而我们在它们中被预先指出——我们将来要借着别人的东西事奉天主。因为全体子民出离埃及（这是在神圣引导下发生的）是整个教会将来从外邦人中的出离的预象和形象；为此原因，祂最后把他们从这个世界领出，进入祂自己的产业，那产业是天主仆人\Person{梅瑟}所没有（赐予）的，而是天主子\Person{耶稣}将作为产业赐予的。如果有人仔细注意先知们关于（时期的）终结所说的那些事，以及主的门徒\Person{若望}在\WorkTitle{默示录}中所见的那些事，他将发现万民普遍要遭受同样的灾祸，如同当时埃及特别遭受的一样。\par

\SourceRecord{Translated by Alexander Roberts and William Rambaut.；Ante-Nicene Fathers；Vol. 1.；Edited by Alexander Roberts, James Donaldson, and A. Cleveland Coxe.；Buffalo, NY: Christian Literature Publishing Co.,；1885.；<http://www.newadvent.org/fathers/0103430.htm>.}

% structure: p0001 source=h1 target=section reason=page-title

\section{驳异端（卷四，第三十一章）}

1. 当提及某些关于古人的这类事情时，（前面提到的）长老常教导我们说：\Quote{对于圣经自身责备族长和先知过的那些不端行为，我们不应痛斥他们，也不应成为像含那样嘲笑父亲羞耻而招致诅咒的人；但我们应该[反而]为他们的缘故感谢天主，因为他们的罪过已通过我们主的来临获得了赦免；因为他说，他们曾[为我们]感恩，并因我们的救恩而夸耀。至于那些圣经未加指责、仅单纯记载[作为已经发生的事]的事情，我们不应成为[行此事者]的控告者，因为我们并不比天主更精确，也不能胜过我们的主；反而应在[它们中]寻求预象。因为圣经中凡未受谴责而记载下来的事，没有一件是没有意义的。} 例子见于罗特的事例，他把女儿们从索多玛领出来，她们随后由自己的父亲怀孕；他又将妻子留在（那地的）境内，她直到今日仍是一根盐柱。因为罗特并非出于自己的意志冲动，也不是由于肉欲的驱使，更没有丝毫此类知识或念头，而确实成就了（未来事件的）预象。正如圣经所说：\BibleQuote{那夜，长女进去与她父亲同睡；罗特不知道她躺下，也不知道她起来。} \BibleRef{创 19:33} 同样的事也发生在次女身上：\BibleQuote{他不知她与他同睡，也不知她起来。} \BibleRef{创 19:35} 既然罗特不知道（他所做的事），也不受情欲（在他行为中）的支配，（天主所计划的）安排就实现了，由此指明两个女儿（即两个教会）从同一位父亲所生，不受肉欲（影响）而生子。因为没有别的人（如她们所想的）能赋予她们生命的种子和生育子女的能力，正如经上所记：\BibleQuote{长女对次女说：地上没有一个人按世间的常规来亲近我们；来，我们让父亲喝酒喝醉，与他同睡，好由父亲得子。} \BibleRef{创 19:31}\par

2. 因此，罗特的女儿们怀着淳朴与天真说了这样的话，以为全人类都已灭亡，如同索多玛人所遭遇的，并且天主的愤怒已降临到整个大地。因此她们也应被谅解，因为她们以为只有她们和父亲被留下以保存人类；正因如此，她们欺骗了父亲。此外，她们所说的话指明了一个事实：除了我们的父以外，没有其他人能赐予长幼两个教会生育子女的[能力]。现在，人类之父是天主圣言，正如梅瑟指出时所说：\Quote{祂岂不是获得你[借着生育]、形成你、创造你的父吗？} 那么，祂是在何时将赐生命的种子——即赦罪的圣神，借由祂我们得以生活——倾注在人类身上的呢？岂不是当祂在地上与人一同吃饭、饮酒的时候吗？因为经上说：\BibleQuote{人子来了，也吃也喝。} \BibleRef{玛 11:19} 并且当祂躺下时，祂睡着了并休息了。正如祂自己在达味中所说的：\BibleQuote{我睡了，也休息了。} 又因祂住在我们中间、与我们同住时常这样做，祂又说：\BibleQuote{我的睡眠对我成为甘饴。} \BibleRef{耶 31:26} 这一切都通过罗特被预示：万有的父的种子——即创造万有的天主圣神——与肉身（即祂自己的化工）相混合并合一；借此混合与合一，两个会堂——即两个教会——从她们自己的父亲产生了活生生的儿子，归于生活的天主。\par

3. 当这些事发生时，他的妻子留在[索多玛境内]，不再是可朽坏的肉身，而是一根永远存在的盐柱；借着那些属于人类的自然过程，表明教会——她是地上的盐\BibleRef{玛 5:13}——也被留在大地的界限内，受人类痛苦的支配；并且当整个肢体经常从她身上被取走时，盐柱仍然存留，如此预表了信德的基础，这信德使人坚强，并将子女送到他们的父面前。\par

\SourceRecord{Translated by Alexander Roberts and William Rambaut.；Ante-Nicene Fathers；Vol. 1.；Edited by Alexander Roberts, James Donaldson, and A. Cleveland Coxe.；Buffalo, NY: Christian Literature Publishing Co.,；1885.；<http://www.newadvent.org/fathers/0103431.htm>.}

% structure: p0001 source=h1 target=section reason=page-title

\section{驳斥异端（第四卷，第三十二章）}

1. 同样，一位宗徒的门徒、一位长老，也曾就两个盟约进行推理，证明二者确实来自同一天主。因为他主张，除创造并塑造我们的那一位以外，再没有别的神；那些人声称我们的这个世界是由天使、或任何其他能力、或另一位神所造，他们的言论毫无根据。因为人一旦离开万物的创造者，并承认我们所属的这个受造界是由任何其他者（或通过任何其他者）所形成，他就必然陷入诸多不一致和此类矛盾，且无法提出任何可能或真实的解释。因此，那些引进其他教义的人，向我们隐瞒他们自己对天主的看法，因为他们知道自己教义的站不住脚和荒谬性，并且害怕一旦被击败，就难以顺利逃脱。但如果有人相信独一天主，祂也藉着圣言创造了万物，正如梅瑟所说：\BibleQuote{天主说：\InnerQuote{要有光}，就有了光} \BibleRef{创 1:3}，以及我们在福音中读到：\BibleQuote{万物是藉着祂造的；凡受造的，没有一样不是藉着祂造的} \BibleRef{若 1:3}，同样，宗徒保禄也说：\BibleQuote{只有一个主，一个信德，一个圣洗，一个天主和圣父，祂在万有之上，贯通万有，且在你们众人之内} \BibleRef{弗 4:5}——这人首先\BibleQuote{持守元首，整个身体都靠祂联络得合适，借着各关节的供应，按各部分的功能，使身体增长，在爱中建立自己} \BibleRef{弗 4:16}。那么，一切话语在他眼中也将显得一致，只要他本人与教会中的长老们一同勤奋研读圣经，宗徒的教导就在他们中间，正如我已指出的。\par

2. 因为所有宗徒都教导，两个民族之间确实有两个盟约；但同一天主为了那些将要相信天主的人的利益而制定了二者（这些盟约正是为他们而赐下），这一点我已从宗徒的教导本身在第三卷中证明了。并且，第一个盟约不是无理由、无目的或偶然地赐下的；相反，它使领受者顺服于事奉天主，为他们有益（因为天主不需要人的事奉），并展示了天上事物的预像，因为人当时还不能通过直接看见而认识天主的事；它还预表了那些现今确实存在于教会中的事物的图像，以使我们的信德得以坚固；并包含了对未来事物的预言，使人知道天主预知一切。\par

\SourceRecord{Translated by Alexander Roberts and William Rambaut.；Ante-Nicene Fathers；Vol. 1.；Edited by Alexander Roberts, James Donaldson, and A. Cleveland Coxe.；Buffalo, NY: Christian Literature Publishing Co.,；1885.；<http://www.newadvent.org/fathers/0103432.htm>.}

% structure: p0001 source=h1 target=section reason=page-title

\section{驳斥异端（第四卷，第三十三章）}

一位这样领受了天主之神——祂从起初就在天主的一切时期中与人同在，宣告未来之事，启示现在之事，讲述过去之事——的属神门徒，确实\BibleQuote{判断一切人，但自己却不被任何人判断}。因为他判断外邦人，\BibleQuote{他们崇拜受造物超过崇拜造物主}\BibleRef{罗 1:21}，并且以败坏的心思将所有劳力浪费于虚妄。他也判断犹太人，他们不接受自由之言，也不愿出来得自由，虽然救主就在他们面前；他们却在不适合的时候，假装以超出律法要求的敬奉事奉一无所缺的天主，并且不承认基督为拯救人类而完成的来临，也不愿明白所有先知都预言了祂的两次来临：第一次，祂成为忍受鞭打、知道软弱是什么的人，\BibleRef{依 53:3}骑在驴驹上，\BibleRef{匝 9:9}成为匠人所弃的石头，被牵到宰杀之地如羊，\BibleRef{依 53:7}并借着伸开双手消灭了阿玛肋克；\BibleRef{出 17:11}同时，祂从地极聚集流散的儿女进入父的羊圈，\BibleRef{依 11:12}并记念祂以前睡了的死人，降到他们那里以拯救他们：第二次，祂将驾云来临，\BibleRef{达 7:13}带来那如火炉燃烧的日子，\BibleRef{拉 4:1}以口中的言语击打大地，\BibleRef{依 11:4}以嘴唇的气息杀死不虔敬的人，手里拿着簸箕，扬净禾场，把麦子收进仓里，把糠秕用不灭的火烧掉。\BibleRef{玛 3:12}\par

此外，他还将审查马尔强的教义，查明他如何主张两个神彼此相隔无限距离。或者，那位从他处带走不属于他的人，并召唤他们进入自己王国的人，怎能是善的呢？为什么他的良善不拯救所有人，反有瑕疵？再者，为什么他对人似乎是善的，但对创造人的那一位却极不公正，因为他剥夺了祂的产业？此外，主若属于另一位父亲，怎能公正地从我们所属的受造界拿起饼，承认那是祂的身体，并肯定调和之杯是祂的血？如果祂没有经历属于人的出生，又怎能承认自己是人子？再者，祂怎能赦免我们对造物主和天主所负的那些罪债？假设祂不是血肉之躯，而只是外表的人，祂怎能被钉在十字架上，并从被刺的肋旁流出血和水？\BibleRef{若 19:34}埋葬祂的人又放入了什么样的身体到坟墓里？又从死者中复活的是什么？\par

这位属神的人也将判断所有瓦伦廷的追随者，因为他们固然用舌头承认独一天主圣父，并承认万有都来自祂，但同时主张形成万有的那一位是背道或缺陷的产物。同样，他们固然用舌头承认独一主耶稣基督、天主子，但在他们的教义体系中却将一种自己的产物归于独生子，另一种自己的产物归于圣言，另一种归于基督，又另一种归于救主；因此，按照他们的说法，这一切存在者固然在圣经中被说成好像一个，但他们却主张每一个都应被理解为与其他分开，并根据其特殊的结合而有自己独特的起源。这样看来，他们的舌头恰恰承认了合一，而他们的意见和理性（由于他们探究深奥的习惯）却已经偏离了合一，接受了多神的观念——当他们受基督审查他们所发明的那几点教义时，这点必然显明。此外，他们还断言祂是在众永世之圆满之后出生的，并且祂的产生是在堕落或背道之后发生的；他们还主张由于索菲亚所经历的受难，他们自己才得以诞生。但是他们自己的特殊先知荷马，听了他的话他们发明了这些教义，将亲自责备他们，因为他这样说：\par

\Quote{可憎如哈得斯之门的人，\newline{}心里想一样，嘴上说另一样。}\par

这位属神的人也将判断那些乖谬的诺斯替派的虚妄言论，指明他们是西满·玛古斯的门徒。\par

4. 祂也要审判厄彼翁派；[因为]若非天主在地上完成了他们的救恩，他们如何能得救呢？人又如何能进入天主，除非天主[先]进入人？人又如何逃脱那必死的世代，若非藉着一种新而奇妙的、出乎意料的（但作为救恩的标记）由天主所赐的世代——即由童贞女藉着信德而涌出的重生？他们若仍留在这种[属]于这世界之人的本性世代中，又如何能从天主获得义子之名？祂（基督）怎能比撒罗满更大\BibleRef{玛 12:41}，或比约纳更大，或是达味的主\BibleRef{玛 22:43}，而祂与他们原是同一性体？祂又怎能制伏\BibleRef{玛 22:29}那比人更强壮者，那不仅征服了人，还将人控制在自己权下，并击败了那征服者，同时释放了被征服的人类，除非祂比那如此被击败的人更大？但那按天主的肖像受造的人，除了按祂的肖像造人的天主子，有谁比那人更优越、更崇高？为此，在末世祂显示了相似；[因为]天主子成了人，将祂手中的古老创造收入自己的本性，正如我在前一卷书中所表明的。\par

5. 祂也要审判那些认为基督只是[在]外表[成人]的人。因为如果他们的导师只是幻想的存在，他们怎能想象自己进行真正的讨论？如果祂只是想象的存在而非真实，他们又怎能从祂领受任何稳固之物？如果祂显为纯粹想象的存在，这些人在祂内所信靠的，又怎能真正分享救恩？因此，与这些人相关的一切都是虚幻的，没有任何真实的特性；在这种情况下，可能会产生疑问：是否（因为或许他们自己同样并非人类，而是哑巴动物）他们在大多数情况下仅仅呈现一个人的幻影？\par

6. 祂也要审判假先知，这些人没有从天主领受预言的恩赐，也没有天主的敬畏之心，而是为了虚荣，或为了某种个人利益，或在邪灵的影响下行其他事，假装说预言，却始终在天主前说谎。\par

7. 祂也要审判那些制造分裂的人，他们缺乏对天主的爱，只顾自己的特殊利益，而不顾教会的合一；他们因微不足道的理由，或任何想到的理由，切割并分裂基督那伟大而荣耀的身体，并尽己所能[真正]摧毁它——这些人侈谈和平，却引起战争，真是滤出蚊虫，吞下骆驼。\BibleRef{玛 23:24}因为不能由他们实现如此重大的改革，足以补偿他们分裂所造成的损害。祂也要审判所有真理之外的人，即教会之外的人；但祂自己不受任何人审判。因为在祂一切一致：祂完全信仰独一全能的天主，万物由祂而有；信仰天主子耶稣基督我们的主，万物藉祂而有，并信仰与祂相关的诸般安排，藉此天主子成了人；并坚定相信圣神，祂赐给我们真理的知识，并展示了父与子的安排，按父的旨意，祂与每一世代的人同住。\par

8. 真知识[在于]宗徒的教导，以及普世教会从起初的宪章，基督的身体藉主教的传承而有的独特彰显，藉此他们将那存在于各地并传到我们这里的教会传递下来，这教会藉着非常完整的教导体系，受到守护和保存，没有任何篡改圣经，既不增添也不删减[她所信仰的真理]；并且在于无伪地阅读[天主圣言]，以及合法而勤勉地按圣经解释，既无危险也无亵渎；[尤其在于]爱德的卓越恩赐，\BibleRef{格后 8:1}它比知识更珍贵，比预言更荣耀，超越一切[天主的]其他恩赐。\par

9. 因此，教会因她对天主的爱德，在各处，历世历代，不断将众多殉道者送到父那里；而所有其他人不仅自己中间没有此类事迹可指，甚至声称这种见证完全非必要，因为他们的教义体系即为真见证[为基督]，唯一例外是，或许从主显现于世以来，他们中间有一两个偶尔与我们的殉道者一同为这名号受辱（好像他也[异端者]得了怜悯），并与他们一同被带去[受死]，如同赐给他们的随从。因为唯独教会纯洁地承受那些为义而受迫害、忍受各种惩罚、因爱天主及宣认祂的圣子而被处死之人的凌辱；她常被削弱，却立即增长其成员，并恢复完整，正如她的预像，罗特的妻子，成了盐柱。同样，她经历与古先知相似的经历，如主所说：\BibleQuote{因为他们在你们以前也这样迫害过先知；}\BibleRef{玛 5:12}因为她确实是以后来的方式，受那些不接受天主圣言者的迫害，而同一圣神停在她身上，\BibleRef{伯前 4:14}[如同在那些古先知身上]。\par

10. 事实上，先知们除了预言其他事物外，也预言了这一点：凡圣神将临于其身、并服从圣父之言、按各自能力事奉祂的人，都必遭受迫害，被石头砸死，被杀害。先知们因着对天主的爱，并为祂的圣言的缘故，在他们自身中预表了这一切。因为他们自己是基督的肢体，每一个都在其作为肢体的位置上，据此预表了〔所托付给他们的〕预言；他们全体虽多，却只预表一位（基督），并宣告那属乎一（基督）的事。正如全部身体的活动是通过我们的肢体来展现，而完整的人的形像并非由一个肢体显示，而是通过所有肢体合在一起，同样，所有先知都预表了那一位（基督）；而他们中的每一位，在其特定的作为肢体的位置上，据此完成了〔所确立的〕救恩计划，并预先影射了与那肢体相关的基督的特定行动。\par

11. 因为他们中的一些人，看见祂在荣耀中，看见了祂在圣父右手的荣耀\OriginalTerm{生活}{conversationem}；另一些人看见祂驾着云彩而来，如同人子；\BibleRef{达 7:13}而那些论及祂说\BibleQuote{他们要瞻望他们所刺透的那一位}的人，\BibleRef{匝 12:10}预示了祂的第二次来临，关于此事祂亲自说：\BibleQuote{你想人子来临时，能在世上找到信德吗？}保禄也论及此事说：\BibleQuote{如果天主是公义的，祂必以患难报复那些加患难于你们的人，却将安宁赐予你们这些受难的人，与我们一起，当主耶稣从天上同祂大能的天使、在火焰中显现的时候。}\BibleRef{得后 1:6}另有其他人论及祂如同法官，并〔提及〕仿佛炽热的火炉，即主的日子，祂\BibleQuote{将麦子收入自己的仓里，却用不灭的火烧尽糠秱，}\BibleRef{玛 3:12}以此威胁那些不信的人，关于他们主也亲自宣称：\BibleQuote{可咒骂的，离开我，到那为魔鬼和他的使者预备的永火里去吧！}\BibleRef{玛 25:41}宗徒也同样论及他们说：\BibleQuote{他们要受惩罚，永远沉沦，离开主的面和祂大能的荣耀，当祂来临时，要在祂的圣徒身上受光荣，在一切信者身上受赞美。}\BibleRef{得后 1:9}还有一些人宣称：\BibleQuote{你比世人更美；}以及\BibleQuote{天主，你的天主，用喜乐之油膏你胜过你的同伴；}以及\BibleQuote{大能者啊，愿你腰间佩刀，带着你的威严和荣美，为真理、谦卑、公义前进，得胜，统治。}以及其它类似关于祂的言论，这些表明祂国度的美善与光辉，以及祂权下所有被高举的超凡卓越〔属于〕他们，使听见的人渴望在那里被寻见，做天主喜悦的事。此外，还有人说：\BibleQuote{祂是一个人，谁能认识祂？}以及\BibleQuote{我进到女先知那里，她生了一个儿子，祂的名字称为奇妙、策士、全能的天主；}以及那些宣扬祂为厄玛奴耳，〔生于〕童贞女的人，表明了圣言与祂自己的创造物的结合，〔宣告〕圣言成了血肉，天主子成为人子（纯洁者纯洁地打开了那纯洁的、使人重生归于天主的子宫，祂亲自使其纯洁）；祂成为和我们一样的人，却〔仍然是〕全能的天主，拥有不可言说的生育。还有一些人说：\BibleQuote{主在熙雍发言，从耶路撒冷发出声音；}\BibleRef{岳 3:16}以及\BibleQuote{天主在犹大为人所认识；}——这些表明祂在犹太地发生的来临。又有宣告\BibleQuote{天主从南方来到，从茂密的山岭而来}的人，\BibleRef{哈 3:3}宣告祂在白冷的来临，正如我在前一卷中指出的。从那地方，那统治并牧养祂父的百姓的那一位也来了。又有宣告在祂来临时\BibleQuote{瘸子必跳跃如鹿，哑巴的舌头必说话清晰，瞎子的眼必睁开，聋子的耳必听见，}\BibleRef{依 35:5}以及\BibleQuote{下垂的手、无力的膝必得坚固，}\BibleRef{依 35:3}以及\BibleQuote{坟墓里的死人必要复活，}\BibleRef{依 26:19}以及祂亲自\BibleQuote{担当我们的病患，背负我们的忧愁}\BibleRef{依 53:4}——〔所有这些〕宣告了由祂所完成的医治之工。\par

12. 其中有些人——当他们预言祂要如同一个软弱无荣的人，一个知道承担病苦的人，\BibleRef{依 53:3}，并骑着驴驹进入耶路撒冷，\BibleRef{匝 9:9}；祂要将背转向鞭打的人，\BibleRef{依 50:6}，将自己的面颊转向掌掴之人；祂要像一只羊被牵到宰杀之地，\BibleRef{依 53:7}；他们要给祂喝醋和苦胆；祂要被祂的朋友和近人离弃；祂要终日伸手，\BibleRef{依 65:2}；祂要被观看的人嘲笑和毁谤；祂的衣服要被分开，为祂的衣袍拈阄；祂要被带到死亡的尘土之中，以及其他一切类似的事——预言了祂以人的形象进入耶路撒冷，在那里，祂藉着受难和被钉十字架，经历了以上所说的一切。另有些人，当他们说：\Quote{神圣的上主记念了祂那些睡在尘土中的死者，并下到他们那里，为要叫他们复活，好拯救他们}，便为我们提供了祂忍受这一切苦难的原因。再者，那些说：\Quote{在那日，上主说：太阳将在正午落下，大地将在晴朗的白昼变为黑暗；我要把你们的庆节变为悲哀，把你们的一切歌曲变为哀歌}的人，\BibleRef{亚 8:9}，清楚地宣告了在祂被钉十字架时，从第六时辰起发生的日蚀，以及此后，他们按法律所守的那些节日和歌曲，在交于外邦人时，将变为忧伤和哀哭。耶肋米亚在论及耶路撒冷时，更清楚地说明了这一点，他这样说：\Quote{那生了七个子女的，如今衰弱了；她的灵魂已厌倦；她的太阳在正午已落下；她受了羞辱，遭受了凌辱；他们中剩下的人，我要在敌人面前交与刀剑。}\BibleRef{耶 15:9}\par

13. 此外，那些论及祂曾打盹、安眠，以及祂因上主扶持而复起的人，并吩咐天上的掌权者打开永久的门，好让光荣的君王进入的人，预先宣告了祂藉着圣父的能力从死者中复活，以及被接入天上。当他们这样说：\Quote{祂的升起从高天之上，祂的回归直到至高之天；没有一人能在祂的热力前隐藏}时，便宣告了这一真理，即祂被提回祂降临之处，并且没有一人能逃脱祂公义的审判。那些说：\Quote{上主为王了，愿万民战栗；祂坐在革鲁宾之上，愿大地震动}的人，部分预言了祂升天后，万民对相信祂的人所发的愤怒，以及全地反对教会的震动；部分预言了当祂随同祂的众天使从天降临时，全地将被震撼，正如祂自己所说：\Quote{将有大灾难，是从起初直到如今没有过的。}\BibleRef{玛 24:21} 再者，当一位说：\Quote{凡受审判的，要站到对面；凡成义的，要就近天主的仆人}；以及\Quote{祸哉，你们！因为你们要像衣服一样变旧，蛀虫要吃掉你们}；以及\Quote{凡有血肉的都将谦卑，惟有上主在至高之处被尊崇}，\BibleRef{依 2:17}——这便显示了，在祂受难和升天之后，天主要使一切反对祂的人屈服于祂的脚下，祂要超越万有，没有人能成义或与祂相比。\par

14. 至于那些宣称天主将与人订立新盟约\BibleRef{耶 31:31}的，不是像在曷勒布山上与祖先所立的那样，并将赐给人一颗新心和一个新灵；\BibleRef{则 36:26} 以及又说：\Quote{不要追念旧事；看，我要行新事，现在就要发生，你们要知道；我必在旷野中开辟道路，在干地上开凿江河，为给我的选民，我的人民饮水，他们是我为自己所获取的，为叫他们传扬我的赞美。}\BibleRef{依 43:19}——这便清楚地宣告了那区别新盟约的自由，以及新酒装入新皮囊之中的事，\BibleRef{玛 9:17}，那就在基督内的信德，祂藉此宣布了在旷野中产生的正义之路，以及圣神在干地上的河流，为给天主所拣选的人民，即祂所获取的人民饮水，好叫他们传扬祂的赞美，而不是亵渎那行这些事的天主。\par

15. 所有那些其他要点，即我已通过如此长篇的圣经系列证明先知们所宣告的，真正属神的人将会藉此指明：对于每一件所说的事，它涉及主救恩计划的哪一特定点，并[藉此展示]天主圣子工作的整个体系，他始终认识同一的天主，始终承认同一的天主圣言——尽管祂[仅仅]如今才向我们显现；也始终承认同一的天主圣神——尽管祂在这末世以一种新的方式倾注在我们身上，[知道祂从世界的创造直到终结，临于人类全体]，从祂那里，那些相信天主并跟随祂圣言的人，领受那源于祂的救恩。另一方面，那些离开祂、藐视祂的诫命、以行为羞辱造他们的主、以意见亵渎养育他们的主的人，为自己堆积最公正的审判。见：\BibleRef{罗 2:5}因此，他（即属神的人）洞察并考验一切，但他自己却不受任何人考验：见：\BibleRef{格前 2:15}他既不亵渎圣父，也不废弃祂的救恩计划，也不攻击先祖，也不羞辱先知——不主张他们[被差遣]来自另一个天主[与他所崇拜的不同]，也不主张他们的预言来自不同源头。\par

\SourceRecord{Translated by Alexander Roberts and William Rambaut.；Ante-Nicene Fathers；Vol. 1.；Edited by Alexander Roberts, James Donaldson, and A. Cleveland Coxe.；Buffalo, NY: Christian Literature Publishing Co.,；1885.；<http://www.newadvent.org/fathers/0103433.htm>.}

% structure: p0001 source=h1 target=section reason=page-title

\section{驳斥异端（第四卷，第三十四章）}

1. 现在我要简单地说，反对所有异端者，尤其是反对马尔强的追随者及类似他们的人，他们主张先知来自另一位天主（而非福音中所宣告的那位）：请认真阅读那由宗徒传给我们的福音，并认真阅读先知书，你们将会发现，我们主的一切行为、一切训诲、一切苦难，都通过先知预言了。但是，如果你心中产生这样一种想法，问说：主藉祂的来临给我们带来了什么呢？要知道，祂带来了全部（可能的）\Quote{新}，因为祂亲自带来了那位曾被预告的那一位。因为这件事早已被预告：将有新事来到，以更新和振兴人类。因为君王的来临，早由那些被派遣的仆人预先宣告，为预备和装备那些要迎接他们主的人。但当君王真正来临，祂的臣民充满了预先宣告的喜乐，获得了祂所赐的自由，得见祂的容颜，聆听祂的话语，享受祂施予的恩赐时，任何有理智的人都不会再问：君王除开那些宣告祂来临者所预告的之外，还带来了什么新鲜事？因为祂亲自带来了自己，并将那些预先宣告的恩赐赐予人类，这些恩赐正是天使所切望窥察的。\BibleRef{伯前 1:12}\par

2. 但是，如果基督在来临之时，被发现完全如先前所预告的那样，却没有实现他们的话，那么那些仆人就会被证明是虚假的，并非主所派遣。因此祂说：\BibleQuote{你们不要以为我来是废除法律或先知；我来不是为废除，而是为成全。我实在告诉你们：即使天地过去了，法律和先知的一笔一划也决不会过去，直到一切都实现。}\BibleRef{玛 5:17-18} 因为祂藉着来临亲自成全了一切，并且仍在教会中成全由法律所预告的新盟约，直至万事的终结。为此，祂的宗徒保禄在致罗马人书中也说：\BibleQuote{但是如今，在法律之外，天主的义德已彰显出来，有法律和先知为证据；因为义人将因信德而生活。}\BibleRef{罗 3:21}\BibleRef{哈 2:4} 而这\BibleQuote{义人将因信德而生活}的事实，早已由先知预告了。\BibleRef{哈 2:4}\par

3. 但是，先知们从哪里获得能力预言君王的来临，预先宣讲祂所赐的自由，并预先宣告基督所行的一切事——祂的话语、祂的作为、祂的苦难，以及预言新盟约呢？如果他们是从另一位天主（而非福音中启示的那位）获得先知灵感，而你又说他们不知道那不可言喻的圣父、祂的王国和祂的救恩安排——这些正是天主圣子在这末世来到世上时所实现的——那么他们又怎能做到这些呢？你也不能说这些事是偶然发生的，仿佛先知们论及的是另外一个人，而类似的事件却发生在主身上。因为所有先知都预言了同样的事，但这些事从未在古代任何人身上应验。如果这些事曾在他们当中的某人身上发生过，那么后来的先知肯定不会再预言这些事将在末世发生。此外，事实上，在先祖、先知或古代君王中，没有一个人身上曾确切而具体地发生任何一件这样的事。因为他们都曾预言基督的苦难，但他们自己远远没有遭受与预言相似的苦难。而且，主受难时所发生的各项细节，在其他任何情况下都未曾实现。因为古代任何人的死亡都不曾伴随正午日食，圣殿的帐幔未曾撕裂，大地未曾震动，岩石未曾崩裂，死人未曾复活，也没有任何一个古人曾在第三日复活并被接升天，在他升天时诸天未曾打开，外邦人未曾相信任何其他名字；也没有一个人从死里复活而开启自由的新盟约。因此，先知们所说的不是别人，正是主，因为所有这些预兆都集中在他身上。\par

4. 若有人为犹太人辩护，主张这新盟约在于从巴比伦流放后、在则鲁巴贝耳领导下所重建的那座圣殿，以及百姓七十年后从那里离开，那么，让他知道：用石头建造的圣殿当时确实重建了（因为那时仍遵守刻在石版上的法律），但并没有赐下新盟约，他们一直使用梅瑟法律直到主的来临；但从主降临起，那带来和平的新盟约和赐予生命的法律已传遍普世，正如先知们所说：\BibleQuote{因为法律将出自熙雍，上主的话将出自耶路撒冷；祂要责斥许多民族，他们要把刀剑铸成犁头，把枪矛打成镰刀，不再学习战事。}\BibleRef{依 2:3} 因此，如果另一道出自耶路撒冷、外邦人接受的法律和圣言带来了这样的和平，并藉着他们使许多民族认识到自己的愚昧，那么先知们似乎是在谈论另一个人。但如果那自由的法律，即天主的圣言，由宗徒们（他们从耶路撒冷出发）传遍全地，引起了如此大的改变，以至这些民族把刀剑和战矛打成犁头和镰刀用来收割谷物，即用于和平目的的工具，并且他们如今已不习惯打仗，反而挨打时也献上另一边面颊，\BibleRef{玛 5:39} 那么先知们所说的就不是另一个人，而是成就这些事的那一位。这一位就是我们的主，在这位身上这预言得以应验；因为正是祂亲自制作了犁，引入了镰刀，即人的初次播种，这播种是在亚当身上显示的创造，以及末后由圣言收集的果实；为此缘故，祂将起初与末后连结，是两者的主，最终祂展示了犁，即木头与铁结合，以此清洁了祂的土地；因为圣言与肉体坚固结合，并用钉子固定在其构造中，收服了蛮荒的大地。起初，祂藉着亚伯预表了镰刀，指出应有一支正义的人种被收集。祂说：\BibleQuote{看，义人怎样灭亡，无人放在心上；虔诚的人被除去，无人介意。}\BibleRef{依 57:1} 这些事先在亚伯身上预演，也由先知们预先宣告，却在主身上实现；同样的事也发生在我们身上，身体跟随元首的榜样。\par

5. 这些就是适合反驳那些主张先知们受另一神感动、我们的主来自另一位父之人的论证，但愿这些异端者最终能停止这种极端的愚妄。我竭力引证这些圣经证据，正是为了藉着这些经文，尽我所能驳斥他们，阻止他们如此严重的亵渎和疯狂捏造众多神祇。\par

\SourceRecord{Translated by Alexander Roberts and William Rambaut.；Ante-Nicene Fathers；Vol. 1.；Edited by Alexander Roberts, James Donaldson, and A. Cleveland Coxe.；Buffalo, NY: Christian Literature Publishing Co.,；1885.；<http://www.newadvent.org/fathers/0103434.htm>.}

% structure: p0001 source=h1 target=section reason=page-title

\section{\WorkTitle{驳斥异端（第四卷，第三十五章）}}

1. 再者，为反驳\OriginalTerm{瓦伦廷派}{Valentinians}及其他自称\OriginalTerm{诺斯替派}{Gnostics}的人，他们主张圣经的某些部分一时来自\OriginalTerm{普累若麻}{Pleroma}之巅，藉着来自那处的种子（\Emph{种子}）而发言，一时又来自中间居所，藉着那鲁莽的母亲\OriginalTerm{普鲁尼卡}{Prunica}而发言，但许多部分则出自世界的创造者，先知们也正是从祂领受使命；我们要说，将宇宙之父贬低到如此窘迫的境地，以致祂不能拥有自己合用的工具，藉以完美地宣告普累若麻中的事，这是全然不合理的。因为祂惧怕谁，以致不能按自己的方式、独立、自由、不受那生于堕落与无知之灵的影响，而启示自己的旨意呢？是担心当更多人听到纯正的真理时，会有许多人得救吗？或者，反过来，祂无力为自己预备那些宣告救主来临的人呢？\par

2. 然而，如果救主来到此世时，派遣祂的宗徒进入世界，准确地宣告祂的来临，并教导圣父的旨意，与外邦人和犹太人的教训毫无共同之处；那么，当祂尚在普累若麻内时，岂不更要任命自己的先驱，宣告祂将来要降临此世，并与那些源自\OriginalTerm{德穆革}{Demiurge}的预言毫无共同之处吗？但若祂在普累若麻内时，曾使用那些在律法之下的先知，并藉着他们宣告自己的事；那么，当祂降临此地时，岂不更会使用同样的教师，藉着他们向我们传福音吗？因此，他们不能再声称伯多禄、保禄及其他宗徒传讲了真理，而经师、法利塞人以及其他颁布律法的人却传讲了谬误。因为，当祂降临时，祂派遣自己的宗徒在真理的灵内，而非在错谬的灵内；对先知也是如此，因为天主圣言始终是同一的。并且，按照这些人的体系，若来自普累若麻的灵是光明的灵、真理的灵、完美的灵、知识的灵，而来自德穆革的灵是愚昧、堕落、错谬的灵，且是黑暗的产物；那么，在同一位存在者身上，怎能同时有完美与缺陷、知识与愚昧、错谬与真理、光明与黑暗呢？但若这种事在先知身上不可能发生，因为他们传讲了来自独一真天主的主的圣言，并宣告了祂圣子的来临；那么，主自己更绝不会一时从高处说话，一时又从低下的堕落中说话，从而同时成为知识与愚昧的老师；祂也绝不会一时荣耀世界创造者为父，一时又荣耀超越祂的那一位为父，正如祂自己所宣称的：\BibleQuote{没有人把新衣服的一块补在旧衣服上，也没有人把新酒装在旧皮袋里。}\BibleRef{路 5:36} 因此，这些人要么完全与先知无关，如同对待古人一般，不再声称他们被德穆革预先差遣，并在那属于普累若麻的新影响下说了某些话；要么，他们应被我们的主说服，因祂声明新酒不能装在旧皮袋里。\par

3. 然而，他们母亲的后裔从何处获得关于普累若麻内奥秘的知识，以及论说它们的能力呢？假设母亲在普累若麻之外确实生下了这个后裔；但他们所说的普累若麻之外，乃是知识之外，即愚昧。那么，那在愚昧中受孕的种子，怎能具有宣告知识的能力呢？或者，母亲本身，一个无形无状、被当作流产丢弃的东西，虽在外部被塑造、被赋予形态，并被\OriginalTerm{霍洛斯}{Horos}禁止进入内部，直到（万物的）终结仍留在普累若麻之外，即在知识之外，她又如何获得关于普累若麻内奥秘的知识呢？再者，当他们说主的受难是上方基督的延伸的预像——祂藉着霍洛斯完成此事，并因而赋予他们的母亲以形态——他们在主受难的其他细节上被驳倒，因为关于这些细节，他们没有任何类似预像可展示。因为上方的基督何时曾被人给醋和苦胆喝？何时他的衣服被分开？何时他被刺透，流出血和水？何时他汗珠如大血点滴下？至于主所遭遇的其他细节，先知们也曾预言，同样的问题也适用。那么，母亲或她的后裔，怎能预言那些尚未发生、但以后必将发生的事呢？\par

4. 他们断言，此外还有一些东西是从普累若麻发出的，但这些都被圣经中关于基督来临的话所驳斥。然而对于这些（从普累若麻发出的东西）是什么，他们并不一致，而是给出不同的答案。因为如果有人想考验他们，逐一询问他们的首领关于某段经文，他会发现其中一人将这段经文指向元父——即拜托斯；另一人归之于阿耳刻——即独生子；另一人归之于万有之父——即圣言；还有人会说这是关于那个由普累若麻中的移涌共同形成的移涌说的；另一些人（认为这段经文）是指基督，而另一个人（指）救主。又一个比这些人更精明的人，在长时间沉默之后，宣称这是关于霍洛斯说的；另一个说它意指着在普累若麻内的索菲亚；另一个说它宣告了普累若麻外的母亲；而另一个人会提到那造世界的天主（德谬哥）。他们对同一（经文）就有如此多的分歧，对同一圣经持有不一致的意见；当同一段经文被读出时，他们全都开始皱起眉毛，摇头，并且说他们确实能说出高超的言辞，但并非所有人都能理解其中隐含的思想的伟大；因此，在智慧人中沉默是首要之事。那至高的\OriginalTerm{沉默}{Sige}必须由他们所保持的那种沉默来表征。这样一来，他们所有的人，各自怀有对同一件事如此多的意见，并且暗自保守他们聪明的见解，便都（彼此）分离。因此，当他们自己之间对圣经所预言的事情达成一致时，他们也就会被我们驳倒。因为尽管持有错误的见解，但他们同时却定了自己的罪，因为他们对同一话语意见不一。但我们以唯一真的天主为我们的导师，并以祂的话语作为真理的准则，我们大家都在同一事上说同样的话，只知道一位天主，这个宇宙的创造者，祂派遣了先知，带领百姓出埃及地，在这末期内显示了自己的圣子，为要使不信者羞愧，并寻找义德的果实。\par

\SourceRecord{Translated by Alexander Roberts and William Rambaut.；Ante-Nicene Fathers；Vol. 1.；Edited by Alexander Roberts, James Donaldson, and A. Cleveland Coxe.；Buffalo, NY: Christian Literature Publishing Co.,；1885.；<http://www.newadvent.org/fathers/0103435.htm>.}

% structure: p0001 source=h1 target=section reason=page-title

\section{反异端论（第四卷，第三十六章）}

1. 此一位天主，主并不拒绝，也不说先知们是出于不同于祂父的另一位神所说的话；也不是出于任何别的本质，而是出于同一位同一的父；也不是说任何别的存在者创造了世界上的万物，唯独除了祂自己的父，当祂在教导中如此说时：\Quote{有一个家主，他栽了一个葡萄园，周围围上篱笆，园中挖了一个压酒池，建了一座楼，租给园户，就往远方去了。到了收果子的时候，他打发仆人到园户那里，去收果子。园户拿住仆人，打了一个，杀了一个，又用石头打死了一个。主人又打发别的仆人，比先前的更多，园户还是照样待他们。最后，他打发自己的儿子到他们那里去，说：\InnerQuote{他们必会尊敬我的儿子。}但园户看见那儿子，就彼此说：\InnerQuote{这是继承人；来吧，我们杀了他，占他的产业。}于是他们拿住他，把他推出葡萄园外，杀了。这样，葡萄园的主人来的时候，要怎样处治这些园户呢？}他们回答说：\InnerQuote{要恶狠狠地消灭这些恶人，把葡萄园另租给按时交果子的园户。} \BibleRef{玛 21:33} 主又说道：\Quote{你们没有念过吗？\InnerQuote{匠人所弃的石头，已成了屋角的基石；那是主的作为，在我们眼中奇妙莫测。}因此我告诉你们：天主的国必从你们手中夺去，赐给那能结果子的民族。} \BibleRef{玛 21:42} 这些话清楚地给祂的门徒指出同一位家主——即一位天主父，祂自己创造了一切；同时也指明有各种不同的园户：一些是顽梗、骄傲、无用、杀害主人的人，另一些则以完全的服从按时交纳果子；并且是同一位家主，有时差遣仆人，有时差遣祂的儿子。因此，那派遣儿子被派到杀害祂的园户那里去的父，也就是派遣仆人的那位。但儿子，由于来自父，拥有\OriginalTerm{最高的权威}{principali auctoritate}，常如此说：\Quote{我却对你们说。} \BibleRef{玛 5:22} 而仆人们，由于是来自他们的主，以仆人的方式说话（传递信息），所以他们常说：\Quote{上主这样说。}\par

2. 这些人（先知们）曾向不信者宣讲为主的那一位，正是基督教导那些服从祂的人所认识的那一位；那召叫前世代的人的天主，也正是收纳后来世代的人的那一位。换言之，那起初使用导致奴役的律法者，也就是后来藉着义子之名召唤祂子民的那一位。因为天主栽种了人类这个葡萄园：起初祂塑造了亚当，并拣选了列祖；然后祂把葡萄园租给园户，当祂设立梅瑟律法时；祂围上篱笆，即关于他们的敬拜给了特别的指示；祂建了一座楼，即拣选了耶路撒冷；祂挖了一个压酒池，即预备了接收先知圣神的器皿。于是，在流亡巴比伦之前，祂派遣了先知，然后在那事件之后，又派遣了更多比先前更多的先知，去收取果子，这样对他们（犹太人）说：\Quote{上主这样说：要洁净你们的行为和作为，施行公正的审判，各人要以怜悯和同情看待自己的兄弟；不要欺压寡妇、孤儿、外方人和穷人；你们心中不可图谋恶事害弟兄，也不可喜爱假起誓。你们要洗濯，自洁，从你们心中除掉恶念，学习行善，寻求公正，保护受压者，为\OriginalTerm{孤儿}{pupillo}伸冤，替寡妇辩护；然后来，我们彼此辩论，上主说。} 又说：\Quote{要禁止舌头不出恶言，嘴唇不说诡诈的话；远离恶，行善；寻求和平，追求和平。} 先知们宣讲这些事，寻求正义的果子。最后，祂派遣自己的儿子，我们的主耶稣基督，到那些不信者那里，却被恶园户赶出葡萄园并杀害了。因此，上主天主甚至把葡萄园（不再围上篱笆，而是向全世界敞开）给了别的园户，他们按时交纳果子——那美丽的被拣选的楼也处处升起。因为卓越的教会如今处处可见，压酒池也处处挖好：因为那些接受圣神的人处处都是。由于从前那些人拒绝了天主子，并杀害祂，把祂赶出葡萄园，天主便公义地弃绝了他们，把葡萄园的出产交给了葡萄园外的外邦人。这符合耶肋米亚所说的：\Quote{上主弃绝并抛开了做这些事的民族；因为犹大的子孙在我眼中行了恶，上主说。} \BibleRef{耶 7:29} 同样，耶肋米亚又这样说：\Quote{我给你们派了守望者；你们要听号角声；他们却说：\InnerQuote{我们不听从。}因此外邦人听见了，那些在其中牧放羊群的人也听见了。} \BibleRef{耶 6:17} 因此，栽种葡萄园、带领百姓、派遣先知、派遣自己的儿子、又把葡萄园交给那些按时交纳果子的别的园户的，是同一位父。\par

3. 因此，主对祂的门徒们说话，为使我们成为好的工人：\BibleQuote{你们应当谨慎，并要时刻醒寤，免得你们的心为宴饮沉醉及今生的挂虑所累，那日子就忽然临到你们；因为它要像网罗一样临到所有住在地上的人。}\BibleRef{路 21:34}\BibleQuote{你们要把腰束起，把灯点着，应当如同那些等候自己的主人由婚宴回来的人。}\BibleRef{路 12:35}\BibleQuote{因为在\Person{诺厄}的日子里，人们吃喝、买卖、嫁娶，直到\Person{诺厄}进了方舟，洪水来了，把他们全都消灭了；同样，在\Person{罗特}的日子里，人们吃喝、买卖、种植、建造，直到\Person{罗特}离开\Place{索多玛}，天降火与硫磺，消灭了他们全体：人子来临时也要如此。}\BibleRef{路 17:26}\BibleQuote{所以，你们要醒寤，因为你们不知道你们的主在哪一天要来。}\BibleRef{玛 24:42}[在这些话中] 祂宣告是同一的主，曾在\Person{诺厄}的时代因人的悖逆而降下洪水，也曾在\Person{罗特}的日子因\Place{索多玛}人众多的罪恶而从天降火，并因同样的悖逆和类似的罪恶，将在时间的末后（\OriginalTerm{在末后}{in novissimo}）带来审判的日子；在那个日子，祂宣告对\Place{索多玛}和\Place{哥摩辣}比那不接受祂宗徒之言的城和家更容易忍受。\BibleQuote{至于你，\Place{葛法翁}，}祂说，\BibleQuote{莫非你要被高举到天上吗？你必被降到地狱！因为如果在你们那里所行的异能行在\Place{索多玛}，它必存留到今天。我实在告诉你们：在审判的日子，\Place{索多玛}所受的惩罚比你们还容易忍受。}\BibleRef{玛 11:23}\par

4. 既然天主子永远是同一的，祂赐给相信祂的人涌流到永生的水泉\BibleRef{若 4:14}[涌出]到永生，却使不结果的无花果树立即干枯；在\Person{诺厄}的日子，祂公义地降下洪水，为灭绝当时那种最邪恶的人类，他们不能为天主结出果实，因为犯罪的天使与他们混杂在一起，并且祂这样做是为了抑制这些人的罪恶，但同时保存了\Person{亚当}形成的\OriginalTerm{原型}{archetype}。又是祂在\Person{罗特}的日子，从天上降下火与硫磺，落在\Place{索多玛}和\Place{哥摩辣}，作为\Quote{天主公义审判的榜样}，为使众人知道，\BibleQuote{凡不结好果子的树，都要被砍倒，投入火中。}\BibleRef{玛 3:10}又是祂用这些话，说在公审判中，\Place{索多玛}所受的惩罚比那些看见祂的奇能却不相信祂、不接受祂教训的人更容易忍受\BibleRef{玛 11:24}。因为正如祂借着来临赐给那些相信祂并承行祂旨意的人更大的特权，同样祂也指明那些不相信祂的人在审判中应受更严厉的惩罚；这样，祂给予众人等量的公正，且要向他们中的一些要求更多，因祂给予他们更多；然而，更多并非因为祂启示另一位父的知识，正如我如此充分且反复说明的，而是因为祂借着来临，将父的恩宠的更大恩赐倾注在人类身上。\par

5. 如果我所陈述的仍不足以说服任何人，相信先知们是由同一位圣父所派遣，我们的主也是由这位圣父所派遣，那么，就让这样的人打开他心灵的口，呼求主基督耶稣这位导师，听祂所说的话：\BibleQuote{天国好像一个国王，为他的儿子举办婚宴，他打发仆人去召叫那些被请的人来赴婚宴。}当他们不肯听从时，祂继续说：\BibleQuote{他又打发别的仆人，说：告诉那些被请的人，来，我已预备了我的午餐，我的公牛和肥畜都已宰了，一切都齐备了，来赴婚宴吧。但他们漠不关心，就走了，一个到自己的田里，另一个去做自己的买卖；其余的拿住仆人，有的凌辱，有的杀死了。国王听了，就发怒，派军队消灭了那些凶手，焚毁了他们的城市，然后对仆人们说：婚宴已经预备好了，但那些被请的人都不配。所以你们到大路上，凡是你们遇到的，都请来赴婚宴。仆人们就出去，在大路上把所遇到的，不论坏人好人，都聚集了来，婚宴上就坐满了客人。国王进来巡视坐席的客人，看见那里有一个人没有穿婚宴礼服，便对他说：朋友，你没有穿婚宴礼服，怎么到这里来了？那人无言以对。于是国王对仆人们说：捆起他的手和脚，把他丢到外面的黑暗中：在那里要有哀号和切齿。因为被召的人多，被选的人少。}\BibleRef{玛 22:1} 现在，主通过祂的这些话清楚地表明所有[这些点，即：]有一位万有的父、君王及主，祂曾先前说过：\BibleQuote{你不可指着耶路撒冷起誓，因为它是大王的城；}并且祂从起初为祂的儿子预备了婚宴，并用最大的仁慈，藉着祂的仆人，召叫先前的时代的人来赴婚宴；当他们不肯听从时，祂又打发别的仆人去邀请他们，然而他们仍不听从，反而用石头砸死并杀死了那些带邀请消息的人。祂于是派出军队消灭了他们，焚毁了他们的城市；但祂从所有的大道上，即从万民中，召集[宾客]来参加祂儿子的婚宴，正如祂藉着\Person{耶肋米亚}所说：\BibleQuote{我也曾打发我的仆人先知们到你们那里去，说：现在你们各人当从极恶的道路上回转，改正你们的行为。}\BibleRef{耶 35:15} 祂又藉着同一位先知说：\BibleQuote{我也曾打发我的仆人先知们到你们那里去，终日从早到晚；但他们没有听从我，也没有侧耳向我。你要将这话告诉他们：这是一个不听从主的声音、也不接受纠正的民族；信德已从他们口中消失了。}\BibleRef{耶 7:25} 因此，那位藉着宗徒们到处召叫我们的主，正是从前藉着先知们召叫他们的主，正如主的话所显示的；虽然他们向不同的民族宣讲，但先知们并非来自一位天主，宗徒们来自另一位；而是[出于]同一位，他们中有些人宣告了主，有些人宣讲了圣父，其他人则预言了天主子来临，而另一些人则向当时远方的人宣布祂已经来临。\par

6. 此外，祂还明明显示：我们在蒙召之后，也应当以义行装饰自己，以便天主的圣神安息在我们身上；因为这正是婚宴礼服，宗徒也提到它：\BibleQuote{我们并不愿意脱去身体，而是愿意穿上另一层，为使那有死的被生命所吞灭。}\BibleRef{格后 5:4} 但那些确实被召赴天主的宴席，却因恶行而没有领受圣神的人，祂宣告说，他们\BibleQuote{将被投入外面的黑暗中。}\BibleRef{玛 22:13} 祂清楚表明：正是同一位君王，从各处聚集忠信者来赴祂儿子的婚宴，并赐予他们不朽的宴席，［也］下令将没有穿婚宴礼服的人，即轻视它的人，投入外面的黑暗。因为正如在旧盟约中，\BibleQuote{他们许多人不得天主喜悦；}\BibleRef{格前 10:5} 同样，此处也是\BibleQuote{被召的人多，但被选的人少。}\BibleRef{玛 22:14} 因此，并非一位天主审判，另一位父召叫我们得救；也不是一位，当然，赐予永恒的光明，而另一位却下令将没有穿婚宴礼服的人投入外面的黑暗。而是同一天主，我们主的父，先知们也由祂领受使命，祂确实藉着无限仁慈召叫不配的人，但祂查验那些被召的人，［察看］他们是否穿上合宜且适合祂儿子婚宴的礼服，因为任何不当或邪恶的事都不令祂喜悦。这与主对那被治好的人所说的话相符：\BibleQuote{看，你已痊愈了；不要再犯罪，免得你遭遇更坏的事。}\BibleRef{若 5:14} 因为那良善、公义、纯洁、无瑕的祂，决不容许任何邪恶、不义、可憎之事进入祂的婚宴厅。这就是我们主的父，万物藉祂的眷顾存立，万有凭祂的命令治理；祂将自己的恩赐白白赐予那些应当［领受］的人；但至公义的报应者最公正地按他们的功过施报，给忘恩负义和漠视祂仁慈的人；因此祂说：\BibleQuote{祂派遣了军队，消灭了那些凶手，并烧毁了他们的城。}\BibleRef{玛 22:7} 祂在此说\BibleQuote{祂的军队}，因为所有人都是天主的财产。因为\BibleQuote{大地和其中的万物，世界和居住其间的，都属于上主。} 因此，宗徒\Person{保禄}也在\BibleBook{罗马}{书}中说：\BibleQuote{因为没有权柄不是出于天主的；掌权的都是天主所立的。所以，抗拒掌权的，就是抗拒天主的命令；抗拒的必自取刑罚。因为官长不是使行善的惧怕，而是使作恶的惧怕。你愿意不怕掌权的吗？你只管行善，就可以从他得称赞；因为他是天主的仆人，是为你的益处。但你若作恶，就该惧怕；因为他不是徒然佩剑；他是天主的仆人，是伸冤的，刑罚那作恶的。所以你们必须顺服，不但是因为刑罚，也是因为良心。为此你们也要纳税；因为他们是天主的差役，常常专管这事。}\BibleRef{罗 13:1} 因此，主和宗徒都宣告独一天主圣父，祂是赐予法律者，派遣先知者，创造万有者；因此祂说：\BibleQuote{祂派遣了军队，} 因为每个人，就其为人而言，都是祂的造物，尽管他可能不认识他的天主。因为祂赐给众人存在；祂，\BibleQuote{使太阳上升，光照恶人，也光照善人；降雨给义人，也给不义的人。}\BibleRef{玛 5:45}\par

7. 而且，不只藉着所述之事，也藉着两个儿子的比喻——其中幼子因与娼妓挥霍无度而耗尽了家产——主教导了同一位圣父，祂连一只小山羊都不给长子；但对于那失落的，即他的幼子，却吩咐人宰杀肥牛犊，并给他穿上最好的袍子。\BibleRef{路 15:11} 此外，藉着在不同时辰被派到葡萄园工作的工人的比喻，也宣告了同一天主\BibleRef{玛 20:1} 在起初，当世界最初被创造时，召叫了一些人；之后又召叫了另一些人；在中间时期召叫了一些人；在漫长岁月之后又召叫了一些人；而在时间的终结又召叫了另一些人；因此，在各自的时代有很多工人，但只有一位家主召集他们。因为只有一个葡萄园，既然只有一种义，也只有一位安排者，因为只有一位天主的圣神安排万有；同样也只有一份工资，因为每人都领了一个德纳，上面印有［君王的］肖像和名号，即对天主圣子的认识，那是不朽的。因此，祂开始把工资给那些最后［被雇］的人，因为在末世，当主被显示时，祂把自己呈现给众人［作为他们的赏报］。\par

8. 随后，关于那在祈祷上胜过法利塞人的税吏，[我们发现]他之所以从主那里得到证言说他是成义的，而不是[另一个人]，并不是因为他敬拜了另一位父；而是因为他以极大的谦卑，脱离了一切夸耀和骄傲，向同一天主作了宣认。\BibleRef{路 18:10} 两个儿子的比喻也是如此：那些被派往葡萄园的人，其中一个确实违抗了他的父亲，但后来悔改了，当悔改对他毫无益处时；另一个却答应去，立即向父亲保证，但他没有去（因为\BibleQuote{每个人都是说谎者；} \BibleQuote{愿意是有的，但行不出来} \BibleRef{罗 7:18}）——[我所说的这个比喻]指出了同一位父。接着，这一真理又由无花果树的比喻清楚地显示出来，关于这比喻，主说：\BibleQuote{看，这些年来我在这无花果树上寻找果子，但找不到} \BibleRef{路 13:6}（藉着先知指向祂的来临，祂不断前来，在那里寻求公义的果子，却没有找到），也由于前面已提及的理由，无花果树应该被砍掉。不用比喻，主对\Place{耶路撒冷}说：\BibleQuote{\Place{耶路撒冷}，\Place{耶路撒冷}，你杀害先知，又用石头砸死那些被派遣到你这里来的人；我多少次愿意聚集你的子女，好像母鸡把自己的小鸡聚集在翅膀下，但你却不愿意！看，你的房屋将为你留下一片荒凉。} \BibleRef{路 13:34} 因为在比喻中所说的\BibleQuote{看，我三年以来寻找果子，}以及明确的话中又[主说]：\BibleQuote{我多少次愿意聚集你的子女，} 若是我们不理解祂藉先知所[预告]的来临——事实上，祂只到他们那里一次，而且是第一次——那么这些话便是[虚假的]。但那拣选了圣祖和那些[活在第一个盟约之下的人]的，是同一圣言，祂曾藉先知之神眷顾了他们，也眷顾了我们这些藉祂的来临从四方被召集起来的人；除了已经说过的，祂确实宣告说：\BibleQuote{将有许多人从东方和西方来，与\Person{亚巴郎}、\Person{依撒格}和\Person{雅各伯}在天国里一同坐席；但本国的子民将被赶到外面的黑暗里，那里将有哀号和切齿。} \BibleRef{玛 8:11} 那么，如果那些因宗徒们的宣讲而在东方和西方信从祂的人，将与\Person{亚巴郎}、\Person{依撒格}和\Person{雅各伯}在天国里一同坐席，与他们共享[天上的]筵席，那么所显明的就是同一天主，祂确实拣选了圣祖，也眷顾了百姓，并召叫了外邦人。\par

\SourceRecord{Translated by Alexander Roberts and William Rambaut.；Ante-Nicene Fathers；Vol. 1.；Edited by Alexander Roberts, James Donaldson, and A. Cleveland Coxe.；Buffalo, NY: Christian Literature Publishing Co.,；1885.；<http://www.newadvent.org/fathers/0103436.htm>.}

% structure: p0001 source=h1 target=section reason=page-title

\section{\WorkTitle{驳斥异端（第四卷，第三十七章）}}

1. 吾主这句话：\Quote{我多少次愿意聚集你的子女，……而你却不 愿意！}\BibleRef{玛 23:37} 阐明了人类自由的古老法则，因为天主从起初就 造人自由，人拥有自己的权能，正如他拥有自己的灵魂，能自愿地、而非受天主 强制地服从天主的命令。因为天主没有强制，而是永远以善意待我们。因此，祂 向所有人赐予善意的劝告。祂将选择的能力置于人以及天使之中（因为天使是理 性的受造物），以便那些服从者能公正地拥有那善——这善确由天主赐予，却由 他们自己保存；反之，那些不服从者将公正地不得拥有这善，并受到相称的惩罚 ：因为天主仁慈地将善赐予他们，但他们自己并未殷勤持守，也未视之为珍宝， 反而藐视祂那至高的仁慈。因此，他们弃绝了善，如同将其吐出，必将理所应当 地遭受天主公义的审判，这也正是宗徒保禄在致\BibleBook{罗马}{书}中所 见证的，他说：\Quote{难道你轻视祂丰厚的仁慈、宽容和忍耐，却不晓得天主 的仁慈是引你悔改吗？你凭着你那顽固而不悔改的心，反而在震怒之日，即显示 天主公义审判的日子，为你自己累积愤怒。}\Quote{但是荣耀和尊崇，}他说， \Quote{是给一切行善的人。}因此，天主赐予了那善，正如宗徒在这封书信中所 告诉我们的；行善的人将获得荣耀和尊崇，因为他们本可以不行善，却行了善； 而那些不行善的人将受到天主公义的审判，因为他们本可以行善，却没有行善。\par

2. 倘若有些人天生是恶的，另一些人天生是善的，那么后者行善就不 值得称赞，因为他们被造如此；前者也不应受责备，因为他们原本就被造成那样 。但既然所有人本性相同，既能持守并实行善，也能抛弃善而不实行——有些人 就在遵守善法的人中（更不用说在天主面前）公正地受到称赞，并获得对其普遍 择善并持守的应有见证；而另一些人则因弃绝公平与善而受责备，并遭受公义的 惩罚。因此，先知们曾劝勉人行善、行正义、实行仁义，正如我已大量证明的， 因为这样做是在我们的能力之内，也因为我们可能因过分疏忽而遗忘，从而需要 那善天主通过先知赐予我们的善意劝告。\par

3. 因此，主也说：\Quote{你们的光当在人前照耀，好使他们看见你 们的善行，光荣你们在天之父。}\BibleRef{玛 5:16} 又说：\Quote{你们要谨 慎，免得你们的心为宴饮沉醉及人生的挂虑所累。}\BibleRef{路 21:34} 又说 ：\Quote{你们要把腰束紧，把灯点着，如同仆人等候自己的主人由婚筵回来， 为的是主人来到敲门时，立刻给他开门。主人来到时，看见那仆人如此行事，那 仆人才是有福的。}\BibleRef{路 12:35-36} 又说：\Quote{那知道主人的旨意 ，而不准备，也不照他旨意行的仆人，必受许多鞭打。}\BibleRef{路 12:47} 又说：\Quote{你们为什么称呼我：主啊！主啊！而不行我所吩咐的呢？} \BibleRef{路 6:46} 又说：\Quote{如果那仆人心裡说：我的主人必来迟，便开 始拷打仆人和使女，也吃喝醉酒，那仆人的主人要在他不期待的日子，在他料想 不到的时刻来到，把他斩绝，使他与假善人遭受同样的命运。}\BibleRef{路 12:45} 所有这些段落都证明了人的独立意志，同时也表明了天主赐予他的劝告 ，借此祂劝勉我们服从祂，并试图使我们远离对祂的不信之罪，然而祂绝不强制 我们。\par

4. 毫无疑问，如果有人不愿跟随福音本身，他有权拒绝，但这并不 有益。因为人有权不服从天主，并丧失那善，但这样的行为带来不小的伤害和祸 患。因此，保禄说：\Quote{凡事我都可行，但不全有益}\BibleRef{格前 6:12} ，这既指人的自由——就此而言\Quote{凡事都可行}，天主对人毫不强制；也藉 \Quote{不有益}指出我们\Quote{不可用自由作恶事的掩饰}\BibleRef{伯前 2:16} ，因为这不有益。他又说：\Quote{你们每人应说实话，与邻居来往} \BibleRef{弗 4:25}；\Quote{一切坏话都不可出自你们的口；但只按需要说造就 人的话，使听者获得恩宠。不要叫天主的圣神忧郁}\BibleRef{弗 4:29-30}； \Quote{你们从前原是黑暗，但现在在主内却是光明；生活行事就该像光明之子 ，不可在宴乐中醉酒，不可在色欲和放荡中，不可在争执和嫉妒中。你们中从前 也有些人如此，但你们已经洗净了，已经圣化了，已因我们主的名成了义人} \BibleRef{格前 6:11}。倘若行这些事或不行这些事不在我们的能力之内，那么 宗徒（更不用说主自己）为何要劝告我们行某些事、戒绝另一些事呢？但正因为 人起初就拥有自由意志，而按天主肖像造人的天主也拥有自由意志，所以人总是 受到劝告，要坚守那善，而这正是通过对天主的服从来实现的。\par

5. 不仅在工作上，而且在信德上，天主也保存了人的意志自由并由人自己掌控，祂说：\BibleQuote{按你们的信德，给你们成就吧。}\BibleRef{玛 9:29} 由此表明，有一种信德是特别属于人的，因为人有自己特有的意见。又说：\BibleQuote{为信的人，一切都是可能的。}\BibleRef{谷 9:23} 以及：\BibleQuote{去吧！照你所信的，给你成就吧。}\BibleRef{玛 8:13} 所有这些说法都表明，人对于信德有自主权。为此，\BibleQuote{信子的人，有永生；不信子的人，不得见永生，天主的义怒常在他身上。}\BibleRef{若 3:36} 同样，主既显示自己的良善，又表明人拥有自由意志和自己的权能，对\Place{耶路撒冷}说：\BibleQuote{我多少次愿意聚集你的子女，有如母鸡把自己的雏鸡聚集在翅膀下，但你们不愿意！看吧，你们的房屋必给你们留下一片荒凉。}\BibleRef{玛 23:37}\par

6. 而那些持相反[结论]的人，他们自己将主呈现为无能为力的，仿佛祂不能完成祂所愿的；或者，另一方面，仿佛祂不知道他们本性是\Quote{\OriginalTerm{物质的}{material}}，如这些人所说，是不能接受祂的不朽的。\Quote{但祂不应该}，他们说，\Quote{创造出具有背叛能力的天使，以及立刻对祂忘恩负义的人；因为他们被造为理性受造物，赋有审查和判断的能力，不是像无理性或纯动物性的东西那样，它们不能凭自己的意志行事，而是由必然性和强制力牵引到善，在这些事物中只有一种思想、一种用途，机械地沿着同一轨道运作\OriginalTerm{不灵活且无判断}{inflexibiles et sine judicio}，他们不可能成为被造以外的任何东西。}但根据此假设，善既不会令他们感恩，与天主的共融也不会珍贵，善也不会受到极大追求，因为它无需他们自己的努力、关心或研究就能呈现，而是自动地、不经他们操心就植入他们内。如此，他们的善将毫无意义，因为他们是出于本性而非出于意志而为善，是自动地而非因选择而拥有善；因此他们不会明白善是美好之事，也不会以此为乐。因为那些对善无知的人怎能享受善？那些未曾追求善的人有何功劳？那些未曾像竞赛得胜者那样追求善的人，有何冠冕？\par

7. 为此，主也断言天国是属于\Quote{猛力的人}；祂说：\BibleQuote{猛力的人夺取了它。}\BibleRef{玛 11:12} 也就是说，那些以力量和认真努力时刻准备夺去的人。为此，\Person{保禄}宗徒也对\Place{格林多人}说：\BibleQuote{你们不知道吗？在赛跑场上赛跑的，大家都跑，但只有一个得奖赏。你们也该这样跑，好能得到奖赏。凡参加竞赛的，在一切事上都要节制；他们为得到可朽坏的冠冕，而我们为得到不朽的冠冕。但我跑，不像无定向的；我打拳，不像打空气的；我使我身体青肿，并带之服从，免得我给别人宣讲，自己反被摒弃。}\BibleRef{格前 9:24} 因此，这位能干的角力者劝谏我们为不朽而努力，使我们得以被加冕，并视冠冕为珍贵，即藉我们的努力所获得，而非自动环绕我们的\OriginalTerm{并非自然生成}{sed non ultro coalitam}。而越努力，它就越宝贵；越宝贵，我们就越当珍视它。实际上，自动而来的东西，不如经过许多焦虑注意达到的东西那样受珍视。既然这能力已赐予我们，主教导了，宗徒也嘱咐我们要更爱天主，以便我们藉追求为自己到达这[奖品]。否则，无疑我们的善将是[近乎]无理性的，因为不是考验的结果。再者，视觉不会显得如此令人向往，除非我们知道失去视觉的损失；健康也因对疾病的了解而更显可贵；光明也因与黑暗对比；生命也因与死亡对比。同样，天国对于那些知道地上王国的人是尊贵的。但越尊贵，我们就越珍视它；若我们更珍视它，我们在天主面前将更光荣。主因此为我们忍受了这一切，为使我们藉这一切受教导，将来在各方面谨慎，并因理性地受教爱天主，而能继续在祂完美的爱中：因为天主在人的背道上显示了忍耐；而人藉此受教导，正如先知所说：\BibleQuote{你自己的背道要医治你。}\BibleRef{耶 2:19} 天主因而预定了万事，为使人类达至成全，为建立他们，并为启示祂的安排，使善得以彰显，正义得以成全，教会得以按照祂圣子的肖像被塑造，人类最终得以在将来某个时候成熟，藉着这些恩宠而成熟，得以看见并领悟天主。\par

\SourceRecord{Translated by Alexander Roberts and William Rambaut.；Ante-Nicene Fathers；Vol. 1.；Edited by Alexander Roberts, James Donaldson, and A. Cleveland Coxe.；Buffalo, NY: Christian Literature Publishing Co.,；1885.；<http://www.newadvent.org/fathers/0103437.htm>.}

% structure: p0001 source=h1 target=section reason=page-title

\section{驳斥异端（第四卷，第三十八章）}

1. 可是，若有人说：\Quote{那么，天主难道不能从一开始就使人成为成全的吗？}就该知道：就天主本身而言，他确实是永远同一、非受生的，凡事对他都是可能的。但受造物必然不如创造他们的那一位，这恰恰是由于他们较晚的起源；因为新近受造的事物不可能是不受造的。既然它们不是不受造的，正因如此，它们就达不到成全。因为正如这些事物年代较晚，它们就是婴孩般的；它们尚未习惯、未在成全的纪律中受过操练。诚然，一位母亲有能力给她的婴孩吃固体食物，但她并不这样做，因为孩子尚不能接受更实在的滋养；同样，天主自己本来可以一开始就使人成为成全的，但人还不能接受这成全，因为他尚是婴孩。为此，我们的主在末期，当他将万有总归于自己时，来到我们中间，并非照他可能来的方式，而是照我们能够看见他的方式。他原可以轻易地带着他不朽的荣耀来到我们中间，但在那种情况下，我们绝不能承受那荣耀的伟大；因此，他，那圣父的成全之粮，将自己作为奶献给我们，因为我们如同婴孩。当他以人的形像出现时，他这样做了，好使我们，好似从他肉身的乳房得到滋养，并通过这样一段奶的滋养，习惯于饮食天主圣言，便也能在自己内怀有那不朽之粮，即圣父的圣神。\par

2. 为此，保禄向格林多人宣告说：\BibleQuote{我曾用奶喂养你们，没有用肉，因为那时你们还不能吃。}\BibleRef{格前 3:2}这就是说，你们确实学到了我们的主作为人的来临；然而，由于你们的软弱，圣父的圣神尚未停留在你们身上。他又说：\BibleQuote{因为你们中间有嫉妒、纷争和分裂，你们岂不是属血肉的，按照人的方式行事吗？}\BibleRef{格前 3:3}这就是说，圣父的圣神尚未与他们同在，是由于他们生活行事的不成全和缺欠。因此，正如宗徒有能力给他们吃硬食——因为凡宗徒们覆手的人都领受了圣神，他是生命（永生）之粮——但他们却不能够接受，因为他们灵魂的感知官能尚软弱，在有关天主的事上未受操练；同样，天主在起初有能力赐给人成全；但人既然刚被创造，他不可能接受它，即使接受了，也不能保有它，或者既保有它，也不能持守它。为此，圣子，虽然他是成全的，却与其余人类一样度过了婴孩阶段，这样作并非为了他自己的益处，而是为了人的婴孩般生存阶段的益处，为使人能够接受他。因此，人不是不受造的受造物，这在天主方面没有任何不可能或欠缺之处；但这仅仅适用于那新近受造的，即人。\par

3. 在天主那里，同时显示出能力、智慧和良善。他的能力和良善在此彰显：他凭自己的旨意，使从未存在的事物得以存在并塑造它们；他的智慧在于他使受造之物成为和谐一致的整体的一部分；那些因他超卓的仁慈而获得成长和长久存续的事物，反映出非受造者的荣耀，即那位毫不吝惜地施予美善的天主。因为从这些事物受造的事实本身，就得出它们不是非受造的；但藉着它们在整个漫长年代中持续存在，它们将通过天主无偿地赐予它们永恒存在而获得非受造者的能力。因此，在一切事物中，天主居于首位，唯有他是非受造的，是万物的元始，是万物存在的首要原因，而其他一切事物都处于天主的管辖之下。但处于天主的管辖之下就是不朽的延续，而不朽就是非受造者的荣耀。因此，通过这一安排、这些和谐以及这种性质的过程，人，一个受造、有组织的存在，被塑造成非受造天主的肖像和模样——圣父妥善地计划一切并发出命令，圣子执行这些命令并施行创造之工，圣神滋养并增长（所造之物），而人则日复一日地进步，向着成全上升，即接近非受造者。因为非受造者就是成全，即天主。现在，人必须首先受造；受造之后，应获得成长；获得成长后，应得以坚强；得以坚强后，应丰盛；丰盛后，应康复（脱离罪的疾病）；康复后，应受光荣；受光荣后，应看见他的主。因为天主是那将要被看见的，而看见天主产生不朽，但不朽使人亲近天主。\par

4. 因此，那些不等候成长的时间，反而将他们本性的软弱归咎于天主的人，在各方面都是不合理的。这样的人既不认识天主，也不认识自己，贪得无厌、忘恩负义，不愿起初就成为受造之物——即有情感的人；反而超越人类的法则，在成为人之前，就希望现在就如同他们的造主天主一样；而他们比哑巴动物更缺乏理性，[坚持]认为未受造的天主与今日受造的人之间没有区别。因为这些[哑巴动物]并不因为没有使它们成为人而指责天主；相反，每个动物都按其所受造的样子，感谢自己受造。因为我们归咎于祂，是因为我们起初没有被造为神，而先是仅仅为人，最终才成为神；尽管天主出于纯粹的仁慈采取了这一过程，使无人能归咎于祂嫉妒或吝啬。祂宣告说：\BibleQuote{我曾说：你们是神，都是至高者的儿子。}但由于我们不能承受神性的能力，祂接着说：\BibleQuote{但你们要像人一样死亡。}以此陈述两个真理——祂白白恩赐的仁慈、我们的软弱，以及我们拥有自主的能力。因为祂在极大的仁慈之后，恩赐了良善[给我们]，并使人类相似于祂自己[即在于他们自己的权力]；同时，祂凭着预知知道人类的软弱及其后果；但藉着[祂的]爱和[祂的]能力，祂将克服受造本性的实质。因为首先必须展示本性；然后，有死的被不朽所征服并吞没，可朽坏的被不可朽坏所吞没，而人应按照天主的肖像和模样被造，并拥有对善恶的知识。\par

\SourceRecord{Translated by Alexander Roberts and William Rambaut.；Ante-Nicene Fathers；Vol. 1.；Edited by Alexander Roberts, James Donaldson, and A. Cleveland Coxe.；Buffalo, NY: Christian Literature Publishing Co.,；1885.；<http://www.newadvent.org/fathers/0103438.htm>.}

% structure: p0001 source=h1 target=section reason=page-title

\section{\WorkTitle{驳斥异端}（第四卷，第三十九章）}

1. 人领受了善恶的知识。顺服天主、信靠祂、遵守祂的诫命是善的，这就是人的生命；不顺服天主是恶的，这就是人的死亡。因此，天主既赐予人这种\OriginalTerm{心智能力}{magnanimitatem}，人便知道了顺服之善与不顺服之恶，使心灵之眼借着对二者的经验，能有判断地选择更善的事；并且永不怠惰或疏忽天主的命令；更借着经验得知是什么坏事夺去了他的生命——即不顺服天主——从而根本不试图去做；反而知道保全生命的事——即顺服天主——是善的，便以全部热忱勤谨持守。为此，人也有了两方面的经验，兼具两类知识，好能以训练有素的心智选择更善的事物。然而，若人对相反事物一无所知，他又怎能获得关于善事的教导呢？因为如此，我们对所领受之事的把握，便比单凭意见的揣测更为确实、无疑。正如舌头借着品尝体验甜与苦，眼睛借着观看辨别黑与白，耳朵借着听闻识别声音的差异；同样，心灵借着对二者的经验得悉何为善，便因顺服天主而更坚定地持守它：首先，借着悔改丢弃不顺服，如同讨厌和令人作呕之物；然后，真正明白不顺服是什么，即它违背善与甘饴，使心灵根本连尝试不顺服天主的念头都不曾起。但若有人避免对这两类事物的知识，避免双重的认识，他便不知不觉地丧失人性。\par

2. 那么，他还未成为人，怎能成为神呢？刚被造不久，怎能是成全的呢？再者，他若在必死的本性中不顺服造物主，又怎能是不死的呢？因为你必须首先持有人的地位，然后才能分享天主的荣耀。因为不是你造了天主，而是天主造了你。那么，如果你是天主的化工，就当等候你的创造者的手，祂按时创造一切；就你而言，你的受造正在进行之中，正是时候。你要将你的心向祂呈上，柔软而易顺从，并保存创造主塑造你的形状，使你自己内有湿润，免得因硬化而失去祂手指的印痕。但若保存这个框架，你必升达成全之境，因为你里面的湿土已被天主的化工隐藏。祂的手塑造了你的质料；祂还要用纯金纯银在你内外覆蔽你，如此装饰你，甚至\Quote{君王必喜爱你的美貌}。但若你顽固硬化，拒绝祂的技巧之工，并对祂忘恩负义，仅仅因为你被造成一个人——你如此对天主忘恩，便立刻既失去了祂的化工，也失去了生命。因为创造是美善天主的属性，而被造是人的本性。那么，你若向祂交付你所有的一切，即对祂的信德和顺服，你必领受祂的化工，成为天主成全的作品。\par

3. 然而，你若不信祂，且逃避祂的手，不完全的原因必在于你，因为你没有顺服，而不在于那召唤你的祂。因为祂派遣使者召叫人们赴婚宴，但那些不顺服祂的人使自己丧失了王的盛宴。\BibleRef{玛 22:3}因此，天主的技巧并不欠缺，因为祂有能力从石头中给亚巴郎兴起子孙来\BibleRef{玛 3:9}；但那人若未得到它，他就是自己不完全的原因。同样，光也不因那些弄瞎自己的人而缺失；光依旧如常，而那些弄瞎自己的人却因自己的过犯陷入黑暗。光从不靠必然性奴役任何人；同样，天主也不勉强任何不愿意接受祂技巧之工的人。因此，那些背离圣父所赐之光、违反自由律法的人，是因自己的过犯而如此行，因为他们被造为自由的、拥有自主之权的人。\par

4. 但天主预知一切，为两者准备了适宜的居所，仁慈地赐予那渴望之光给那些寻求不朽之光并投奔它的人；但对于那些躲避并背弃这光、如同自盲的轻蔑者和嘲弄者，祂预备了适合这些反对光明者的黑暗，并对那些试图逃避服从祂的人施加了适当的惩罚。服从天主是永恒的安息，因此，那些躲避光明的人拥有一个与其逃避相称的地方；而那些逃避永恒安息的人，则有一个与其逃离相应的住所。既然一切美善都与天主同在，那些因自己的决定而逃离天主的人，便剥夺了自己的一切美善；被剥夺了与天主相关的一切美善后，他们必将落入天主的公义审判之下。因为那些逃避安息的人理应受到惩罚，而那些躲避光明的人理应住在黑暗中。正如这暂时的光，躲避它的人将自己交付给黑暗，以至于他们自己成为自己缺乏光明并居住于黑暗的原因；而正如我已经指出的，光并不是他们这种不幸生存状况的原因；同样，那些逃避天主永恒之光（此光包含一切美善）的人，他们自己成为自己居住在永恒黑暗、缺乏一切美善的原因，他们自己成为这类住所的原因。\par

\SourceRecord{Translated by Alexander Roberts and William Rambaut.；Ante-Nicene Fathers；Vol. 1.；Edited by Alexander Roberts, James Donaldson, and A. Cleveland Coxe.；Buffalo, NY: Christian Literature Publishing Co.,；1885.；<http://www.newadvent.org/fathers/0103439.htm>.}

% structure: p0001 source=h1 target=section reason=page-title

\section{\WorkTitle{驳斥异端（第四卷，第四十章）}}

1. 因此，是同一天主圣父为那些渴望与祂共融并服从祂的人预备了美善之事，也为背道的魁首——魔鬼及其同谋者预备了永火；主曾宣告\BibleRef{玛 25:41}那些被分别站在祂左边的人将被投入那火。这也是先知所说的：\BibleQuote{我是嫉妒的天主，制造和平，也制造灾祸}\BibleRef{依 45:7}；祂这样与那些悔改归向祂的人和好，并引领他们合一，却为那执迷不悟、逃避光明的人预备了永火和外面的黑暗，这些对陷入其中的人确实是灾祸。\par

2. 然而，如果真有一位父赐予安息，而另一位神预备了火，那么他们的儿子也必然互不相同：一位将人送入父的国，另一位则送入永火。但既然同一位主指明，全人类将在审判时被分开，\BibleQuote{好像牧人分开绵羊和山羊一样}\BibleRef{玛 25:32}，并要对一些人说：\BibleQuote{我父所祝福的，你们来吧！承受那为你们预备了的国度}\BibleRef{玛 25:34}，而对另一些人说：\BibleQuote{可咒骂的，离开我，到那为魔鬼和他的使者预备的永火里去吧！}\BibleRef{玛 25:41}，那么显然，同一圣父就是经文所宣告的\BibleQuote{制造和平，也制造灾祸}的那一位，为双方预备了适宜的事；同样，一位审判者将他们送往适宜的地方，正如主在莠子和麦子的比喻中所说：\BibleQuote{就如把莠子收集起来，用火焚烧，在末日也将是这样。人子要差遣祂的天使，把一切使人跌倒的事和作恶的人从祂的国里收集起来，扔到火窑里；在那里要有哀号和切齿。那时，义人在他们父的国里，要发光如同太阳}\BibleRef{玛 13:40}。因此，那位为义人预备了国度、并由圣子接纳配得之人进入的父，也是预备了火窑的那一位，人子差遣的天使将根据天主的命令，把那些应得的人投入其中。\par

3. 主确实在自己的田里撒了好种子，他说：\BibleQuote{田地就是世界。}但人们睡觉的时候，仇敌来了，\BibleQuote{在麦子中间撒上莠子，就走了}\BibleRef{玛 13:28}。由此我们知道，这就是那背叛的天使和仇敌，因为他嫉妒天主的创造，并企图使这创造与天主为敌。因此，天主将他从自己面前驱逐出去，因为他自己偷偷地撒了莠子，即带来违命的那一位；但天主怜悯了人，人虽然因疏忽（当然，但也因另一方邪恶）而陷入悖逆，但天主将魔鬼本想使人成为天主的敌人的敌意转向了其作者，将祂的愤怒从人身上移开，转向蛇。正如经上告诉我们的，天主对蛇说：\BibleQuote{我要把仇恨放在你和女人之间，和你的后裔和她的后裔之间；她的后裔要踏碎你的头颅，你要咬伤他的脚跟}\BibleRef{创 3:15}。主亲自将这仇恨实现在自己身上，当他由女人而成为人，并踏碎了蛇的头颅时，正如我在前书中所指出的。\par

\SourceRecord{Translated by Alexander Roberts and William Rambaut.；Ante-Nicene Fathers；Vol. 1.；Edited by Alexander Roberts, James Donaldson, and A. Cleveland Coxe.；Buffalo, NY: Christian Literature Publishing Co.,；1885.；<http://www.newadvent.org/fathers/0103440.htm>.}

% structure: p0001 source=h1 target=section reason=page-title

\section{驳斥异端（第四卷，第四十一章）}

1. 主既说有些天使——即魔鬼的天使——为他们预备了永火；又说及莠子时，祂声明说：\Quote{莠子是恶者的子女}，\BibleRef{玛 13:38}，因此必须肯定：祂将所有背弃者都归属于这悖逆的首领。但祂并没有使天使或人天生如此。因为我们并未发现魔鬼创造了任何东西，因为他自己也是如同其他天使一样，是天主的受造物。因为天主造成了万物，正如达味论及此类一切时所说：\Quote{祂一发言，万物造成；祂一命令，万物创成。}\par

2. 既然万物都由天主造成，而魔鬼成了他自己及他人的背弃原因，那么圣经通常将那些常处于背弃状态的人称为\Quote{魔鬼之子}和\Quote{恶者的天使}（\OriginalTerm{恶者}{maligni}），这是公义的。因为据在我之前的一位所说，\Quote{儿子}一词有双重意义：一是在自然秩序中，因生而为子；二是在被造就方面，被称为子，尽管生而为子和被造就为子之间有区别。因为前者确实是从所提及的人生出；但后者是由他造成的，无论就创造而言，还是就其教义的传授而言。凡从别人口中受教者，就被称为教导他的人的\Quote{儿子}，而后者被称为他的父亲。因此，就本性而言——即就创造而言，可以这样说——我们都是天主的儿子，因为我们都由天主所造。但就服从和教义而言，我们并不都是天主的儿子：只有那些信祂并承行祂旨意的人才是。而那些不信、不顺服祂旨意的人，则是魔鬼的儿子和天使，因为他们作魔鬼的作为。祂在依撒意亚书中宣告了这一点：\Quote{我养育了子女，使他们长大，他们却背叛了我。}\BibleRef{依 1:2} 又说，这些子女是外邦的：\Quote{外邦的子女向我说谎。}所以，按照本性，他们是天主的子女，因为他们是如此被造的；但就他们的作为而言，他们不是祂的子女。\par

3. 因为在人当中，那些不服从父亲的儿子，虽被剥夺继承权，但在本性上仍是他们的儿子，但在法律上却被剥夺继承权，因为他们不能成为生身父母的继承人；同样，在天主那里也是如此——那些不服从祂的人，被祂剥夺继承权后，就不再是祂的儿子。因此，他们不能承受祂的产业：如达味所说：\Quote{罪人一出母胎就与祢疏远；他们怒气如同蛇的怒气。}所以主称祂所知道的人类后代为\Quote{毒蛇的种类}，\BibleRef{玛 23:33} 因为他们像这些动物一样诡计多端，伤害别人。因为祂说：\Quote{你们要谨慎，防备法利塞人和撒杜塞人的酵母。}\BibleRef{玛 16:6} 谈到黑落德时，祂也说：\Quote{你们去告诉那只狐狸，}\BibleRef{路 13:32} 意指他的狡猾和欺骗。因此先知达味说：\Quote{人虽然在尊荣中，却如同牲畜一样。} 耶肋米亚又说：\Quote{他们像发情的牡马，各向自己邻舍的妻子嘶鸣。}\BibleRef{耶 5:8} 依撒意亚在犹太宣讲、与以色列辩论时，称他们为\Quote{索多玛的首脑}和\Quote{哈摩辣的人民}，\BibleRef{依 1:10} 暗示他们在邪恶方面如同索多玛人，同样的罪恶在他们中间流行，并因他们行为的相似而用同一名称称呼他们。既然他们并非天生被天主造成这样，而是也有能力行善，同一位先知给他们好劝告说：\Quote{你们洗濯，自洁，从我眼前除掉你们的恶行，停止作恶。}\BibleRef{依 1:16} 无疑，既然他们以同样的方式犯罪和逾越，也就受到如同索多玛人同样的责备。但当他们悔改、归向，并停止作恶时，他们就有能力成为天主的儿子，并承受由祂所赐予的不朽产业。因此，祂称那些听从魔鬼、作魔鬼作为的人为\Quote{魔鬼的天使}和\Quote{恶者的子女}。\BibleRef{玛 25:41} 然而，这些人同时又都是同一天主所造。当他们相信并服从天主，继续持守祂的教导时，他们是天主的儿子；但若他们背弃并陷入过犯，就被归属于他们的首领——魔鬼，即那首先成了自己、而后又成了他人背弃的原因者。\par

既然主的话语众多，而它们都宣告同一位圣父——这个世界的创造者，我也应当为了他们自己的益处，用许多论据来反驳那些陷入许多错误的人，若有可能，当他们被许多证据驳倒时，他们可以皈依真理而得救。但是，在接下来的部分中，有必要在本著作之后附加保禄在主的言语之后的教导，审查这个人的观点，阐述这位宗徒，并解释所有被异端者作其他解释的经文（他们完全误解了保禄所说的话），并指出他们疯狂观点的愚蠢；还要从同一位保禄——他们从他的著作中向我们提出难题——那里证明，他们确实是说谎者，而这位宗徒是真理的宣讲者，他教导的一切都与真理的宣讲相符；即同一位天主圣父曾对亚巴郎说话，颁赐法律，预先派遣先知，在末期派遣了祂的圣子，并将救恩赐予祂自己的作品——即肉身的实质。因此，在另一本书中，我将整理主关于圣父的其余话语——这些话语不是以比喻，而是以明显的字面意义表达（\Emph{sed simpliciter ipsis dictionibus}）——以及受祝福的宗徒的书信的阐述，在天主的助佑下，我将向您提供揭露和驳斥那伪称为知识之学问的完整著作；这样，在这五本书中，我和您都可操练如何反对所有异端者。\par

\SourceRecord{Translated by Alexander Roberts and William Rambaut.；Ante-Nicene Fathers；Vol. 1.；Edited by Alexander Roberts, James Donaldson, and A. Cleveland Coxe.；Buffalo, NY: Christian Literature Publishing Co.,；1885.；<http://www.newadvent.org/fathers/0103441.htm>.}

% structure: p0001 source=h1 target=section reason=page-title

\section{反驳异端（第五卷，序言）}

在前四卷书中，我至亲爱的朋友，我已将一切异端者揭露，使他们的教义大白于天下，并驳斥了那些捏造不虔敬见解的人。我藉着引用各人特有且留于其著作中的教义，并运用更普遍且适用于他们全体的论证，完成了此事。然后，我指出了真理，并展示了教会的宣讲，这宣讲是先知们所预言的（正如我已证明的），但由基督使之完美，宗徒们传承下来；教会从宗徒接受了这些真理，并独自在普世完整地（\Emph{bene}）保存它们，传递给她的子女。此外——在解决了异端者向我们提出的一切问题，解释了宗徒的教义，并清楚阐述了主以比喻所说和所做的许多事情之后——我将在全书第五卷（本卷旨在揭露和驳斥那伪称为知识的）中，努力展示主其余教义和书信中的证据：如此满足你的要求，正如你所请求我的（因为我确实在圣言的职务中得到了一席之地）；并且，尽我所能者，努力为你提供大量的帮助，以对抗异端者的反驳，同时挽回迷途者，引领他们归向天主的教会，并坚定初信者的心志，使他们能坚守所接受的信仰，由教会完整地保存，丝毫不被那些教唆谬论、引人背离真理的人所败坏。然而，你以及所有可能读到这篇著作的人，务必以极大的注意细读我已说过的，好使你们明白我所争论的对象。因为如此，你们才能以正当的方式驳斥他们，并准备好接受对他们的反驳证据，借着属天的信德，将他们的教义视如尘芥；但要跟随唯一真实而坚定的师傅，天主的圣言，我们的主耶稣基督，祂藉着超越的爱，成了我们现在的样子，为带领我们成为祂自己那样。\par

\SourceRecord{Translated by Alexander Roberts and William Rambaut.；Ante-Nicene Fathers；Vol. 1.；Edited by Alexander Roberts, James Donaldson, and A. Cleveland Coxe.；Buffalo, NY: Christian Literature Publishing Co.,；1885.；<http://www.newadvent.org/fathers/0103500.htm>.}

% structure: p0001 source=h1 target=section reason=page-title

\section{驳斥异端（第五卷·第一章）}

1. 因为我们没有其他方法可以得知天主的事，除非我们的师傅，即那作为圣言而存在的，降生成人。因为除了祂自己的圣言，没有任何受造物有能力向我们启示圣父的事。还有谁能\BibleQuote{知道主的心意}，或是谁\BibleQuote{做了祂的参谋}？\BibleRef{罗 11:34} 再者，我们也没有其他方法可以学习，唯有亲眼看见我们的师傅，亲耳听见祂的声音，使我们成为祂工作的效法者以及祂言语的实行者，从而与祂共融，从那完美的一位并获得增长，从那先于一切受造者的一位那里获得增长。我们这些不过是最近才由那唯一良善者——那具有不死之恩的——按祂的肖像被塑造（按照圣父的先见被预定，我们这些尚未存在者得以进入存在），并成为受造界初果的，已在预知的时代，按照那在一切事上完美的圣言的职分，领受了救恩的福分。祂是强有力的圣言，也是真实的人，藉着祂自己的血按合理的方式救赎了我们，将自己作为赎价赐予那曾被掳去的人。既然那背叛不公地统治了我们，虽我们本性属于全能的天主，它却违反本性使我们疏离，使我们成为它的门徒，天主的圣言——在一切事上充满能力，且不欠缺自己的正义——便正义地对抗那背叛，并从它手中赎回自己的财产，不是用暴力手段，如同那背叛起初统治我们时，它贪得无厌地攫取不属于自己的东西，而是用说服的方式，正如一位愿意谋划的天主，祂不用暴力获取祂所愿的；如此，正义既不受损害，天主古老的化工也不致毁灭。主既藉着自己的血这样救赎了我们，将祂的灵魂给予我们的灵魂，将祂的肉身给予我们的肉身，也倾注了圣父的圣神，为使人神合一与共融，藉着圣神将天主赐给人，另一方面，藉着祂的降生成人将人联系于天主，并在祂来临时藉着与天主的共融，持久而真实地赐予我们不死的生命——异端所有教义便全然崩溃。\par

2. 那些声称祂只是外表显现的人，实在是虚妄的。因为这些事情不是只在表面上发生，而是真实地发生。但如果祂表面上像人而实际上不是人，那么圣神也不会降临在祂身上——这件事确曾发生——因为圣神是不可见的；那样，祂里面就没有任何真理，因为祂并非祂所显现的那样。但我已经说过，\Person{亚巴郎}和其他先知以先知的方式看见祂，在异象中预言将要发生的事。那么，如果这样的存在如今以与祂实际不同的外表显现，那就出现了一个给人的先知性异象；而且必须期待祂另一次来临，在那来临中，祂将是如同现在在先知方式中被看见的那样。我已经证明，说祂只是外表显现，和断言祂没有从\Person{玛利亚}取得什么，是同一回事。因为除非祂在自己内总结亚当的古老形成，祂就不会真真实实拥有那救赎我们的血肉了。因此，\Person{瓦伦廷}的门徒是虚妄的，他们提出这个意见，为要将血肉排斥于救恩之外，并抛弃天主所塑造的。\par

3. 厄彼翁派也是虚妄的，他们不以信德接受神人结合于自己灵魂内，却停留在（本性的）出生的旧酵母中，不愿明白圣神曾降临于\Person{玛利亚}，至高者的能力曾荫庇她：\BibleRef{路 1:35} 因此，所生的也是圣的，是至高者天主之子、万有的圣父之子，祂成就了这存在的降生成人，并显示了一种新的（样式）的诞生；好使我们藉着这新诞生继承生命，如同藉着前一次诞生继承死亡。因此，这些人拒绝天上之酒的\OriginalTerm{混合}{commixture}，只愿它是世上的水，不接受天主与自己结合，却停留在那被击败并逐出乐园的\Person{亚当}内；他们没有考虑，正如在我们于\Person{亚当}内形成之初，那出自天主的生命气息与所塑造的结合，使人有生命，并显示他为理性的存在；同样，在（时期的）末后，圣父的圣言和天主的圣神，与\Person{亚当}形成的古老质料结合，使人成为活而完美的，能接纳完美的圣父，好使如在天性的（\Person{亚当}）内我们都死了，照样在属神方面我们都要活。\BibleRef{格前 15:22} 因为\Person{亚当}从未逃脱天主的\Emph{手}，圣父曾对祂说：\BibleQuote{让我们照我们的肖像，按我们的模样造人。}因此，在末后的时期（\Emph{终结}），不是凭血气，也不是凭人意，而是凭圣父的旨意，\BibleRef{若 1:13} 祂的手形成了活人，好使\Person{亚当}按天主的肖像和模样被（重新）创造。\par

\SourceRecord{Translated by Alexander Roberts and William Rambaut.；Ante-Nicene Fathers；Vol. 1.；Edited by Alexander Roberts, James Donaldson, and A. Cleveland Coxe.；Buffalo, NY: Christian Literature Publishing Co.,；1885.；<http://www.newadvent.org/fathers/0103501.htm>.}

% structure: p0001 source=h1 target=section reason=page-title

\section{\WorkTitle{驳斥异端}（第五卷，第二章）}

1. 同样，那些说天主来到不属于祂的事物中，仿佛觊觎他人财产的人，也是虚妄的；他们声称天主是为了将那个由另一个人所创造的人，交付给那位既未创造也未形成任何东西的天主，而这位天主从起初就被剥夺了祂自己正当的人性塑造。因此，这些人所主张的降临到他人事物中的那一位，其来临并非正义；祂也并非真正以自己的血救赎了我们，如果祂没有真正成为人，将起初关于人的话——即人是按照天主的肖像和模样造的——归还原属于祂亲手的工作；不是用诡计夺走他人的财产，而是以正义和仁慈的方式占有自己的财产。就背叛一事而言，祂确实以自己的血正义地救赎了我们；但就我们这些被救赎的人而言，祂是仁慈地行事。因为我们先前并没有给祂什么，祂也不像需要什么那样渴望从我们得到什么；但我们确实需要与祂共融。正是为此，祂仁慈地倾注了自己，为将我们聚集到圣父的怀中。\par

2. 但是，那些轻视天主的整个计划、否认肉身的救恩、蔑视其重生、主张肉身不能承受不朽的人，在各方面都是虚妄的。如果肉身确实不能获得救恩，那么主就没有用祂的血救赎我们，圣体圣事之杯也不是祂血的共融，我们所擘的饼也不是祂身体的共融。\BibleRef{格前 10:16} 因为血只能来自血管和肉身，以及构成人本质的其他一切——天主圣言正是成了这样的实体。祂以自己的血救赎了我们，正如祂的宗徒所宣告的：\BibleQuote{在祂内我们藉祂的血获得了救赎，罪过的赦免。} \BibleRef{哥 1:14} 既然我们是祂的肢体，我们也藉受造物得到滋养（祂亲自将受造物赐给我们，因为祂使太阳升起，按祂的意愿降雨 \BibleRef{玛 5:45}）。祂承认那杯（受造物的一部分）为自己的血，从中祂浇灌我们的血；那饼（也是受造物的一部分）祂立为自己的身体，从中祂使我们身体增长。\par

3. 因此，当调和之杯和制成的饼接受了天主圣言，成为了基督的血和身体的圣体圣事——我们肉身的实体由此得到增长和支持——他们怎能断言肉身不能承受天主的恩赐，即永生？这肉身正是由主的身体和血所滋养，并且是祂的肢体！——正如蒙福的保禄在致厄弗所人书中所宣告的：\BibleQuote{我们是祂身体的肢体，是祂的肉和祂的骨。} \BibleRef{弗 5:30} 他这些话不是指某个属灵的、无形的人，因为神灵没有骨也没有肉 \BibleRef{路 24:39}；而是指那计划（主成为真实的人，由肉、神经和骨组成）——那肉身由作为祂血的杯所滋养，并因作为祂身体的饼而增长。正如从葡萄树剪下的枝条栽种在地里，按季节结果；又如一粒麦子落在地里腐烂，藉着包含万物的天主圣神，以多倍的收成升起，然后藉着天主的智慧供给人使用，并接受了天主圣言，成为圣体圣事，即基督的身体和血；同样，我们的身体被它滋养，埋葬在地里，在那里腐烂，将在指定的时刻复活，天主圣言赐予它们复活，归于天主、即圣父的光荣，祂白白地将不朽赐予这必死的，将不朽坏赐予这可朽坏的 \BibleRef{格前 15:53}，因为天主的德能在软弱中才全显出来 \BibleRef{格后 12:3}，为使我们永不骄傲，仿佛我们靠自己拥有生命，并自高反对天主，我们的心变得忘恩负义；而是藉经验知道，我们永恒的生命来自这存有的超越能力，而非出自我们自己的本性，这样我们既不会轻视环绕天主的荣耀，也不会无知于自己的本性，而是要知道天主能做什么，人领受什么益处，从而永不偏离对事物真相的正确理解，即对天主和对人的理解。而且，正如我早已指出的，岂不可能是天主为这目的而允许我们分解入死亡之普通尘土，好让我们藉各种方式受到教导，将来在一切事上都准确无误，既不明于天主，也不明于我们自己？\par

\SourceRecord{Translated by Alexander Roberts and William Rambaut.；Ante-Nicene Fathers；Vol. 1.；Edited by Alexander Roberts, James Donaldson, and A. Cleveland Coxe.；Buffalo, NY: Christian Literature Publishing Co.,；1885.；<http://www.newadvent.org/fathers/0103502.htm>.}

% structure: p0001 source=h1 target=section reason=page-title

\section{反驳异端（第五卷，第三章）}

1. 宗徒\Person{保禄}以最清晰的方式指出，人已被交付于自身的软弱，以免他因高举而远离真理。因此他在\WorkTitle{致格林多人后书}中说：\BibleQuote{免得我因那高超的启示而骄傲，就有一根刺在我身上，就是撒殚的使者来击打我。关于这事，我曾三次求主使它脱离我；但他对我说：\InnerQuote{我的恩宠为你够了，因为能力在软弱中才得以成全。}所以我甘心情愿夸耀我的软弱，好使基督的德能常在我身上。} \BibleRef{格后12:7} 那么，这是怎么回事？（有人或许会惊呼：）主是否希望他的宗徒如此遭受击打，并忍受这样的软弱？正是如此；圣言这样说。因为能力在软弱中才得以成全，使那因自己的软弱而认识天主能力的人成为更好的人。因为一个人若非借经验了解两者，又怎能知道他自己是软弱的、本性的可朽者，而天主是不朽和全能的呢？借忍耐学习自己的软弱并非坏事；相反，它甚至有益，能防止他对自己的本性产生不当的见解（\Emph{不偏离其本性}）。但高举自己反抗天主，将祂的光荣归于自己，使人忘恩负义，已为他带来许多灾祸。[因此，我说，人必须借经验学习这两件事]，使他对自己和造物主都不缺乏真理与爱德。而两者的经验赋予他关于天主和人的真知识，并增加他对天主的爱。哪里爱德增长，哪里就有更大的光荣借天主的德能为那些爱祂的人成就。\par

2. 因此，那些人置天主的德能于不顾，也不思考圣言所宣告的，当他们专注于肉身的软弱时，却不考虑那从死者中复活肉身者的能力。因为如果祂不赋予可朽者生命，不使可腐朽者归于不朽，祂就不是全能的天主。但祂在这一切方面都是有能力的，我们应当从我们的起源得知，因为天主从地中取泥土造了人。确实，从未存在的骨骸、神经、血管以及人身体的其他部分，使这一切成为现实，并造出一个有生命有理性的受造物，这比重新整合那已受造而后分解归土（由于上述原因）之物要困难和不可思议得多，因为后者已进入了先前不存在的人所形成的那些元素。因为那在起初使尚不存在者存在的主，按祂的意愿，更会重新恢复那些先前存在的人，当祂愿意他们承受祂所赐的生命时。而且，那起初接受了天主精巧塑造的肉身，也将被证明适于并能够接受天主的德能；以致一部分成为眼睛用于看，另一部分成为耳朵用于听，另一部分成为手用于触摸和工作，另一部分成为筋腱遍布各处以连结肢体，另一部分成为动脉和静脉作为血液和空气的通道，另一部分成为各种内脏，另一部分成为血液，它是灵魂与身体结合的纽带。但何必继续呢？数字无法表达人体结构部位的多样性，它除了天主的极大智慧外别无造成。但那些分享天主技巧和智慧的事物，也分享祂的德能。\par

3. 因此，肉身并不缺乏对天主创造智慧和德能的参与。但如果那赐予生命者的德能在软弱中——即肉身中——才得以成全，让他们告诉我们，当他们坚持肉身无能接受天主所赐的生命时，他们是否像现在活着的人、分享生命的人那样说这些话，还是承认自己完全没有生命，此刻是死人。如果他们真是死人，他们怎能走动、说话，并履行那些不属于死人而属于活人的功能呢？但如果他们现在活着，并且全身分享生命，他们怎能妄称肉身没有资格分享生命，而他们却承认自己此刻拥有生命呢？这好比有人拿起一块吸满水的海绵，或一支燃烧的火炬，却宣称海绵不可能分享水，或火炬不可能分享火。正是这样，那些人声称自己活着并在肢体中承载生命，随后却又说这些肢体不能接受生命，自相矛盾。但如果现世暂时的生命——其本性远逊于永生——尚且能这样做，使我们的可朽肢体活起来，那么永恒的生命——远比这更有能力——岂不更能赋予那已经与生命打交道并习惯于维持生命的肉身以生命呢？因为肉身确实能分享生命，这由它活着的事实所证明；只要天主愿意，它就继续活着。同样明显的是，天主有能力赐予它生命，因为祂赐生命给我们这些存在的人。因此，既然主有能力将生命注入祂所塑造的，既然肉身能被赋予生命，还有什么能阻止它分享不朽——即天主所赐的极乐无尽的生命呢？\par

\SourceRecord{Translated by Alexander Roberts and William Rambaut.；Ante-Nicene Fathers；Vol. 1.；Edited by Alexander Roberts, James Donaldson, and A. Cleveland Coxe.；Buffalo, NY: Christian Literature Publishing Co.,；1885.；<http://www.newadvent.org/fathers/0103503.htm>.}

% structure: p0001 source=h1 target=section reason=page-title

\section{\WorkTitle{驳斥异端}（第五卷，第四章）}

1. 那些虚构在造化者之外另有一位父亲，并称其为善的天主的人，是在自欺；因为他们把祂描述为一位软弱、无价值、疏忽的存在，甚至可说是恶毒且充满嫉妒，因为他们断言我们的身体并非由祂赋予生命。当他们声称那些明显永生的事物——如灵魂与精神之类——是由那位父亲赋予生命，而另一事物（即身体）——它唯有通过天主的恩赐才能获得生命——却被剥夺了生命时，他们要么承认这表明他们的父亲是软弱无力的，要么便是嫉妒与恶毒的。因为既然造化者在此世就赋予我们必朽的身体以生命，并像我所指出的那样，通过先知应许他们复活，那么谁显得更有能力、更强大、或真正良善呢？是那使整个的人活着的造化者，还是他们那被误称为父亲的那位？祂假装赋予那些本性不朽、生命常存的事物以生命，却并不仁慈地赋予那些需要祂的帮助才能存活的事物以生命，而是漫不经心地任凭它们落入死亡的权势之下。那么，究竟是他们的那位父亲有能力赐予生命却不赐予呢，还是没有能力？如果是因为祂不能，那么按照这个假设，祂就不是一位有能力的存在，也不比造化者更完美；因为造化者所赐予的，我们必须看到，是\Emph{祂}自己所不能给予的。但如果是因为祂有能力却不赐予，那就证明祂不是一位良善的父亲，而是嫉妒与恶毒的。\par

2. 再者，如果他们为其父亲不赋予身体生命而提出任何原因，那么那个原因必然显得高于父亲，因为它阻止了祂施行仁慈；而祂的仁慈将因他们所提出的那个原因而显得软弱。人人都必须明白，身体是能够接受生命的。因为只要天主愿意，它们就能活着；既然如此，（异端者们）便不能主张这些身体完全无能力接受生命。因此，如果由于必然性或任何其他原因，那些能够分享生命的事物未被赋予生命，那么他们的父亲便成了必然性和那个原因的奴隶，而非自由者，不能自主地掌控自己的意愿。\par

\SourceRecord{Translated by Alexander Roberts and William Rambaut.；Ante-Nicene Fathers；Vol. 1.；Edited by Alexander Roberts, James Donaldson, and A. Cleveland Coxe.；Buffalo, NY: Christian Literature Publishing Co.,；1885.；<http://www.newadvent.org/fathers/0103504.htm>.}

% structure: p0001 source=h1 target=section reason=page-title

\section{\WorkTitle{反异端论}（第五卷第五章）}

一、【为要学习】身体确实持续存在了很长一段时间，只要天主乐意它们繁荣，让【这些异端者】阅读圣经，他们就会发现我们的前人在年龄上超过了七百、八百和九百岁；他们的身体与漫长的日子同步，并参与生命，直到天主愿意他们存活。但我为什么提到这些人呢？因为\Person{哈诺客}当他取悦天主时，就在同一身体内被提升，他确实取悦了祂，从而预先指出了义人的提升。\Person{厄里亚}也被接去，【当时他还在】【自然】形体的实质中；这样在预言中展示了那些属神者的被提升，并且表明他们的身体被提升和接去并没有任何障碍。因为正是借着起初塑造他们的同一双手，他们才获得这种提升和接去。因为在亚当身上，天主的手已习惯于安排、管理和维持祂自己的创造，并把它带到和放置在祂喜欢的地方。那么，第一个人被安置在哪里？当然在乐园里，正如圣经所宣告的：\BibleQuote{天主在东方的伊甸栽种了一个乐园（\Emph{paradisum}），并在那里安置了祂所造的人。}\BibleRef{创 2:8}然后后来当【人】证明不顺服时，他便从那里被逐出到这个世上。因此，作为宗徒弟子的长老们告诉我们，那些被提升的人被转移到了那个地方（因为乐园已为义人预备好了，就是那些拥有圣神的人；\Person{保禄}宗徒被提升到那里时，听到了在我们目前的状况下对人而言不可言传的话语\BibleRef{格后 12:4}），并且那些被提升的人将留在那里直到【万物的】终结，作为不朽的先声。\par

二、然而，如果有人想象人不可能存活如此长的时间，并且\Person{厄里亚}不是带着肉身被提升，而是他的肉身在火战车中被烧尽了，那么让他考虑一下\Person{约纳}，当他被抛入深渊，被吞入鲸鱼腹中时，却因天主的命令再次被安全地抛到陆地上。\BibleRef{纳 2:11}然后，当\Person{阿纳尼雅}、\Person{阿匝黎雅}和\Person{米沙耳}被抛入加热七倍的火窑中时，他们丝毫未受伤害，连火的气味也没有在他们身上闻到。正如天主的手与他们同在，在他们身上成就了奇妙的事——这些事是人的本性所不可能做到的——那么，如果也在那些被提升的人身上，天主的手依照天父的旨意成就了一些奇妙的事，这又有什么可奇怪的呢？这手就是天主子，正如圣经所记载的\Person{拿步高}王所说：\BibleQuote{我们不是把三个人捆着抛进火窑里吗？看，我看见四个人在火中行走，那第四个好像天主子。}\BibleRef{达 3:19}因此，任何受造物的本性，或肉体的软弱，都不能胜过天主的旨意。因为天主并不服从于受造物，而是受造物服从于天主；万物都服从祂的旨意。因此主也宣告：\BibleQuote{在人不可能的，在天主却可能。}\BibleRef{路 18:27}因此，正如在今日那些不认识天主安排的人看来，有人能活这么多年似乎不可思议且不可能，但那些在我们之前的人确实活了【这么久】，并且那些被提升的人确实活着作为未来长寿的保证；以及【同样似乎不可能】从鲸鱼腹中和火窑中的人毫发无损地出来，然而他们确实如此，仿佛是借着天主的手被领出来，为要宣告祂的大能：同样现在，虽然有些人不知道天主的能力和应许，可能反对自己的救恩，认为那使死人复活的天主不可能有能力赐予他们永恒，但这一等人的怀疑并不能使天主的信实归于无效。\par

\SourceRecord{Translated by Alexander Roberts and William Rambaut.；Ante-Nicene Fathers；Vol. 1.；Edited by Alexander Roberts, James Donaldson, and A. Cleveland Coxe.；Buffalo, NY: Christian Literature Publishing Co.,；1885.；<http://www.newadvent.org/fathers/0103505.htm>.}

% structure: p0001 source=h1 target=section reason=page-title

\section{驳斥异端（第五卷，第六章）}

1. 天主将因祂的工程而受光荣，祂要使之符合并仿效祂自己的圣子。因为藉圣父的双手，即藉圣子和圣神，人——而非人的一部分——是按照天主的肖像造的。灵魂与神魂固然是人的\Emph{一部分}，但绝不是\Emph{其人}本身；因为成全的人在于灵魂领受父的圣神，与那按天主的肖像塑造的血肉本性相混合并结合。为此，宗徒宣告：\BibleQuote{我们在成全的人中也讲智慧} \BibleRef{格前 2:6}，称那些领受了天主圣神的人为\Quote{成全的}，他们藉天主的圣神说各种语言，正如他本人也这样说。同样，我们在教会中也听到许多弟兄，他们拥有先知之恩，藉圣神说各种语言，为公共益处显明人隐藏的事，并宣示天主的奥秘。宗徒也称他们为\Quote{属神的}，他们之所以属神，是因为他们分受了圣神，而非因为他们的血肉被剥去、被移除，并成为纯粹属神的。因为若有人除去血肉的实体，即天主的工程，而理解那纯粹属神的，那么这样就不是属神的人，而是人的神魂，或天主的圣神。但当这神魂与灵魂混合，并与天主的工程结合时，人便因圣神的倾注而成为属神的和成全的，这就是那按天主的肖像和模样所造的人。但若灵魂缺少圣神，这样的人确实是属血气的，留在血肉中，便是一个不完全的存在，在他的形成中（\OriginalTerm{在形成中}{in plasmate}）确实拥有天主的肖像，但未藉圣神接受模样；因此这个存在是不完全的。同样，若有人除去肖像，舍弃工程，他就不能将其理解为一个人，而是要么如我所说的一部分，要么是别的什么东西。因为那被塑造的血肉本身并非成全的人，而是人的身体，人的一部分。灵魂本身，若单独考虑，也不是人，而是人的灵魂，人的一部分。神魂也不是人，因为它被称为神魂，而非人；但所有这些的混合与结合构成了成全的人。为此，宗徒解释自己，表明得救的人既是完全的人，也是属神的人，在致得撒洛尼人前书中这样说：\BibleQuote{愿和平的天主亲自\OriginalTerm{成全}{perfectos}你们，并愿你们的整个神魂、灵魂和身体，在我们的主耶稣基督来临时，保持无可指摘。}他祈祷这三者——即灵魂、身体和神魂——在主来临时得以保存，其目的是什么？除非他意识到这三者将来要重新整合与结合，并继承同一个救恩。为此，他也宣告那些向主呈献这三部分而无过犯的人就是\Quote{成全的人}。因此，那些使天主的圣神常住在他们内，并保持灵魂和身体无瑕，持守天主的信德（即那朝向天主的信德），并维持与邻人正义往来的人，就是成全的人。\par

2. 因此他也说，这工程是\BibleQuote{天主的圣殿}，如此宣告：\BibleQuote{你们不知道，你们是天主的圣殿，天主圣神住在你们内吗？若有人毁坏天主的圣殿，天主必要毁灭那人，因为天主的圣殿是圣的，这圣殿就是你们。} \BibleRef{格前 3:16}这里他明确宣布身体是圣神居住的圣殿。正如主论及自己说：\BibleQuote{你们拆毁这座圣殿，三天之内我要把它重建起来。}但据说他这话\BibleQuote{是指他身体的圣殿。} \BibleRef{若 2:19}不仅宗徒承认我们的身体是圣殿，甚至是基督的圣殿，他对格林多人说：\BibleQuote{你们不知道你们的身体是基督的肢体吗？那么，我能拿基督的肢体作为娼妓的肢体吗？} \BibleRef{格前 3:17}他说这些话，并非指某个其他属神的人；因为这种本性的存在与娼妓毫无关系；但他宣告\Quote{我们的身体}，即那继续在圣洁和纯洁中的血肉，是\Quote{基督的肢体}；但若与娼妓成为一体，就成为娼妓的肢体。因此他说：\BibleQuote{若有人毁坏天主的圣殿，天主必要毁灭那人。}那么，声称天主的圣殿（圣父的圣神住在其中）和基督的肢体不能分享救恩，反而归于丧亡，岂不是最大的亵渎吗？此外，我们的身体复活不是出自其本身的实体，而是藉天主的能力，他对格林多人说：\BibleQuote{身体不是为淫乱，而是为主，主也是为身体。天主既使主复活了，也要以自己的能力使我们复活。} \BibleRef{格前 6:13}\par

\SourceRecord{Translated by Alexander Roberts and William Rambaut.；Ante-Nicene Fathers；Vol. 1.；Edited by Alexander Roberts, James Donaldson, and A. Cleveland Coxe.；Buffalo, NY: Christian Literature Publishing Co.,；1885.；<http://www.newadvent.org/fathers/0103506.htm>.}

% structure: p0001 source=h1 target=section reason=page-title

\section{\WorkTitle{反异端论}第五卷，第七章}

1. 因此，正如基督以\OriginalTerm{肉身的实体}{substance of flesh}复活了，并向祂的门徒指出钉痕和肋旁的伤口（这些是从死者中复活的那肉身的标记），同样，经上说：\Quote{祂也要}藉祂自己的德能\Quote{使我们复活}。\BibleRef{格前 6:14} 他又对罗马人说：\Quote{然而，如果那使耶稣从死者中复活者的圣神住在你们内，那使基督从死者中复活的，也必藉那住在你们内的圣神，使你们有死的身体活过来。}\BibleRef{罗 8:11} 那么，什么是\OriginalTerm{有死的身体}{mortal bodies}？难道可以是灵魂吗？不，因为灵魂与有死的身体相比是\OriginalTerm{非物质的}{incorporeal}；因为天主\Quote{向人的面上吹了一口\OriginalTerm{生命的气息}{breath of life}，人就成了一个\OriginalTerm{有灵的生灵}{living soul}}。这生命的气息是非物质的东西。他们当然不能坚持说这生命气息本身是可死的。因此达味说：\Quote{我的灵魂也要为祂而生活}，仿佛它的实体是不朽的。另一方面，他们也不能说圣神就是有死的身体。那么，还有什么可被称为\Quote{有死的身体}，除非是那被塑造的，即肉身，经上也说天主将使这肉身\OriginalTerm{活起来}{vivify}？因为这是那死去并\OriginalTerm{分解}{decompose}的，而不是灵魂或圣神。因为死就是失去生命能力，从此成为无气息、无生命、无活动，并分解为那些构成其实体的成分。但这既不发生在灵魂身上，因为它是生命的气息；也不发生在圣神身上，因为圣神是单纯而非复合的，故不能被分解，且祂本身就是那些接受祂者的生命。因此我们必须断定，死亡是就肉身而言的；这肉身在灵魂离开后，变得无气息、无生命，并逐渐分解为它所从出的尘土。这就是可朽的。也正是这个，他说：\Quote{也必使你们有死的身体活过来。}因此，他在（第一封）致格林多人书中论及它说：\Quote{死人的复活也是这样：播种的是可朽坏的，复活起来的是\OriginalTerm{不朽}{incorruption}的。}\BibleRef{格前 15:42} 因为他声明：\Quote{你所播种的，若不先死去，就不能活起来。}\BibleRef{格前 15:36}\par

2. 那如同麦粒被撒在地里而腐朽的，若非被埋在地里的尸体，种子也撒在其上，又是什么呢？为此他说：\Quote{播种的是在\OriginalTerm{羞辱}{dishonour}中，复活起来的是在\OriginalTerm{荣耀}{glory}中。}\BibleRef{格前 15:43} 因为有什么比死去的肉身更可耻？另一方面，有什么比它复活并分享不朽时更光荣？\Quote{播种的是在\OriginalTerm{软弱}{weakness}中，复活起来的是在\OriginalTerm{德能}{power}中：}\BibleRef{格前 15:43} 当然是在它自己的软弱中，因为既然是土，就归于土；但藉天主的德能，祂使它从死者中复活。\Quote{播种的是\OriginalTerm{属生灵的身体}{animal body}，复活起来的是\OriginalTerm{属神的身体}{spiritual body}。}\BibleRef{格前 15:44} 他无疑教导说，这样的话并非指着灵魂或圣神说的，而是指着成为\OriginalTerm{尸体}{corpses}的身体。因为它们是属生灵的身体，即分享生命的身体；当它们失去生命时，就陷入死亡；然后藉着圣神的媒介复活，成为属神的身体，从而藉圣神拥有\OriginalTerm{永恒的生命}{perpetual life}。\Quote{因为如今，}他说，\Quote{我们认识是局部的，我们讲预言也是局部的；但到那时，就要面对面了。}\BibleRef{格前 13:9} 这也是伯多禄所说的：\Quote{你们虽然没有见过祂，却爱祂；如今虽然看不见祂，却相信祂；且相信时，要以不可言传的喜乐而欢跃。}\BibleRef{伯前 1:8} 因为我们的脸将要看见主的脸，并以不可言传的喜乐而欢跃——也就是说，当它看见自己的喜悦时。\par

\SourceRecord{Translated by Alexander Roberts and William Rambaut.；Ante-Nicene Fathers；Vol. 1.；Edited by Alexander Roberts, James Donaldson, and A. Cleveland Coxe.；Buffalo, NY: Christian Literature Publishing Co.,；1885.；<http://www.newadvent.org/fathers/0103507.htm>.}

% structure: p0001 source=h1 target=section reason=page-title

\section{\WorkTitle{驳斥异端（第五卷，第八章）}}

1. 然而我们现在确实领受了祂圣神的一部分，这趋向于完美，预备我们进入不朽，使我们逐渐习惯于领受并承载天主；宗徒将此称为\Quote{抵押}，即天主所应许给我们荣耀的一部份；他在致厄弗所人的书信中说：\BibleQuote{在祂内你们也听了真理之言，即你们救恩的福音，相信了，就领受了所应许的圣神为印记，这圣神是我们嗣业的抵押。}\BibleRef{弗 1:13}因此，这抵押既住在我们内，即使现在也使我们成为\OriginalTerm{属神的}{spiritual}，并且必死的被不朽所吞没。\BibleRef{格后 5:4}\BibleQuote{你们已不在肉性内，而在圣神内，只要天主的圣神住在你们内。}\BibleRef{罗 8:9}然而，这事的发生并非藉着抛弃血肉，而是藉着圣神的传授。因为他所写信的对象并非没有血肉，而是那些领受了天主圣神的人，\BibleQuote{藉着祂我们呼喊：阿爸，父啊！}\BibleRef{罗 8:15}所以，如果现在，拥有抵押，我们就呼喊\Quote{阿爸，父啊}，那么当我们复活后面对面看见祂时，所有肢体迸发出持续的凯旋赞歌，荣耀那位使他们从死者中复活、并赐予永生恩赐的那一位，又将是何等情形呢？因为如果抵押将人聚集于自身，现在已使他呼喊\Quote{阿爸，父啊}，那么由天主赐予人的那圆满的圣神恩宠，将成就何事呢？它将使我们肖似祂，并满全圣父的旨意；因为它将按天主的肖像与模样造人。\par

2. 那么，那些拥有圣神抵押、不受肉性私欲奴役、而是顺服于圣神、并在一切事上依理智之光而行的人，宗徒恰当地称之为\Quote{属神的}，因为天主的圣神住在他们内。然而，属神的人并非无形无体的神灵；而是我们的实体，即血肉与精神的结合，领受天主圣神，构成了属神的人。但那些确实拒绝圣神的引导、成为肉性私欲的奴隶、过着违反理性的生活、毫无节制地投身于自己的欲望、对神圣的圣神毫无渴望的人，像猪狗一样生活；这些人，[我说]，宗徒非常恰当地称之为\OriginalTerm{属肉性的}{carnal}，因为他们除了肉性的事外别无他念。\par

3. 同样，先知也因他们行为的不理性而将他们比作无理性的动物，说：\BibleQuote{他们成了发情的公马，各向邻舍的妻子嘶鸣。}\BibleRef{耶 5:3}又说：\Quote{人在尊贵时被比作牲畜。}这表明，由于他自身的过失，他因效仿无理性的生活而成了畜类。我们也按惯例，将这类人称为畜类和无理性的野兽。\par

4. 法律曾以象征方式预言了这一切，用各种动物来描绘人：\BibleRef{肋 11:2}凡分蹄和反刍的，[经上]说，都定为洁净的；而那些不具备其中一项或两项的，则被排除在外，视为不洁。那么，谁是洁净的呢？那些藉着信德坚定地走向父和子的人；因为这由分蹄的稳定性所象征；他们昼夜默思天主的话，好能以善行装饰自己：这就是反刍的含义。而不洁净的，则是那些既不分蹄也不反刍的；即那些对天主既无信德、也不默思祂话语的人：这就是外邦人的可憎之物。至于那些确实反刍却没有分蹄、因而本身不洁的动物，我们从中看到对犹太人的象征描述：他们口中确实有天主的话，却未能将根基的坚定锚定在父和子内；因此他们是不稳定的世代。因为那些蹄子连成一片的动物容易滑倒；而那些分蹄的则更为稳健，它们的裂蹄在前进中交替支撑。同样，那些分蹄却不反刍的也是不洁的：这明显指示一切异端者，以及那些不默思天主的话、也不以正义行为装饰自己的人；主也对他们说：\BibleQuote{你们为什么称呼我主啊主啊，却不做我所吩咐的事呢？}\BibleRef{路 6:46}因为这类人确实说他们相信父和子，却从不按应当的方式默思天主的事，也不以正义行为装饰自己；反而，如我先前所述，他们效法猪狗的生活，沉溺于污秽、贪食和各种放纵。因此，宗徒正确地称所有这些人为\Quote{属肉性的}和\OriginalTerm{属生灵的}{animal}——即所有那些因自己的不信和放纵而不接受圣神、并用各种方式将\OriginalTerm{生命圣言}{life-giving Word}从自身逐出、愚昧地随从私欲的人；先知也称他们为驮畜和野兽；习俗也将他们视为牛马和无理性的生物；法律则定他们为不洁。\par

\SourceRecord{Translated by Alexander Roberts and William Rambaut.；Ante-Nicene Fathers；Vol. 1.；Edited by Alexander Roberts, James Donaldson, and A. Cleveland Coxe.；Buffalo, NY: Christian Literature Publishing Co.,；1885.；<http://www.newadvent.org/fathers/0103508.htm>.}

% structure: p0001 source=h1 target=section reason=page-title

\section{\WorkTitle{反异端}（第五卷第九章）}

1. 宗徒所宣布的其他真理中，也有这一项：\BibleQuote{血肉不能继承天主的国。}\BibleRef{格前 15:50} 这段经文被所有的异端者用来支持他们的愚昧，试图困扰我们，并指出天主的受造物不得救恩。他们不考虑这一事实：如我所指出的，完整的人由三部分组成——肉身、灵魂和灵。其中一样确实保存并塑造人——这就是灵；而另一样被结合和形成——这就是肉身；然后介于二者之间的——这就是灵魂，它有时随从灵而被提升，但有时同情肉身而陷入肉欲。因此，凡没有那保存并塑造我们进入永生之物的人，将被称为、也将是血肉；因为他们在自身内没有天主的圣神。因此，这类人被主称为\BibleQuote{死人}；因为他说：\BibleQuote{让死人埋葬他们的死人，}\BibleRef{路 10:60} 因为他们没有那使人生活的圣神。\par

2. 另一方面，凡敬畏天主并信赖祂圣子来临的人，并通过信德在天主心中确立了天主的圣神——这样的人将要正当地被称为\Quote{纯洁的}、\Quote{属神的}和\Quote{活着面向天主的}，因为他们拥有那净化人并将人提升到天主生命的圣父之神。因为主曾证明\BibleQuote{肉身是软弱的}，也同样说\BibleQuote{心神是切愿的}。\BibleRef{玛 26:41} 因为心神能实践自己的意愿。因此，若有人将圣神的乐意倾向，犹如刺激，掺入肉身的软弱，则必然的结果是，强壮的将压倒软弱的，以致肉身的软弱被圣神的强力所吸收；那发生这种情况的人，不能是属血肉的，而是属神的，因为与圣神的共融。因此，殉道者们作证，轻看死亡，不是出于肉身的软弱，而是由于圣神的乐意。因为当肉身的软弱被吸收时，它显示出圣神的强大；反之，当圣神吸收软弱的肉身时，它便拥有肉身作为自己的产业，从两者中形成一个活着的人——活着，是因为他分享了圣神；是人，是因为有肉身的实体。\par

3. 因此，肉身若缺少天主的圣神，便是死的，没有生命，不能承受天主的国：如同无理性的血，像水倾倒在地。因此他说：\BibleQuote{那属于土的怎样，凡属于土的也怎样。}\BibleRef{格前 15:48} 但在哪里有天主圣父的圣神，哪里就有活着的人；有理性的血被天主保存，以报复流它的人；肉身被圣神所拥有，确实忘记自身所属，而采纳圣神的品质，与天主的圣言相符。因此，他（宗徒）宣告：\BibleQuote{我们怎样带了那属于土的肖像，也要怎样带那属于天上的肖像。}\BibleRef{格前 15:49} 那么，什么是属土的？那被塑造的。什么是属天的？圣神。因此，正如他说，当我们缺少天上的圣神时，我们从前在肉身的旧生活里行走，不顺从天主；同样，现在我们领受了圣神，愿我们在生命的新颖中行走，顺从天主。因此，既然没有天主的圣神我们不能得救，宗徒劝勉我们通过信德和贞洁的言行，保存天主的圣神，免得我们因不参与天主的圣神而失去天国；他大声呼喊：肉身和血本身不能承受天主的国。\par

4. 然而，若我们必须严格地说，【我们当说】肉身\Emph{不}继承，而是\Emph{被}继承；正如主也宣称：\BibleQuote{温良的人是有福的，因为他们要承受土地。}\BibleRef{玛 5:5}，仿佛在【未来的】国度中，土地——我们肉身的质料从那里而来——要藉继承而承受。这正是祂愿圣殿（即肉身）洁净的理由，为使天主的圣神得以在其中喜悦，如同新郎喜悦新娘。因此，正如新娘不能【说】出嫁，而是被娶，当新郎来迎娶她时；同样，肉身不能凭自己继承天主的国，但它可以被取\Emph{为}继承进入天主的国。因为活着的人继承死者的产业；继承是一回事，被继承是另一回事。前者统治、管理并按其意志支配被继承的事物；而后者处于从属、被管理、被获得继承的人所统治的状态。那么，活着的是什么？无疑是天主的圣神。而死者的产业又是什么？当然是人的各部分，它们在土中腐朽。但当它们被迁移到天主的国时，这些就被圣神所继承。为此，基督也死了，为使福音的盟约向全世界显现并为人所知，首先释放祂的奴隶；然后，正如我已表明的，当圣神藉继承拥有他们时，使他们成为祂产业的嗣子。因为活着的人继承，而肉身被继承。为使我们不因失去那拥有我们的圣神而丧失生命，宗徒劝勉我们共融于圣神，合理地说了前面引述的话：\BibleQuote{血肉不能继承天主的国。}正如他所说：\Quote{不要错误；因为除非天主的圣言住在你们内，圣父的圣神也在你们内，而你们竟只如同这单纯的血肉般轻率粗心地生活，你们就不能继承天主的国。}\par

\SourceRecord{Translated by Alexander Roberts and William Rambaut.；Ante-Nicene Fathers；Vol. 1.；Edited by Alexander Roberts, James Donaldson, and A. Cleveland Coxe.；Buffalo, NY: Christian Literature Publishing Co.,；1885.；<http://www.newadvent.org/fathers/0103509.htm>.}

% structure: p0001 source=h1 target=section reason=page-title

\section{\WorkTitle{反对异端（第五卷，第十章）}}

1. 因此，[他宣称]这一真理，为的是我们不可在纵容肉身的同时，拒绝圣神的接枝。他说：\Quote{但你原是野橄榄，}\Quote{已被接在好橄榄树上，并得以分享橄榄树的肥汁。}\BibleRef{罗 11:17} 因此，正如野橄榄被接枝后，若仍保持原状，即仍是一棵野橄榄，它就被\Quote{砍下来，丢在火里；}\BibleRef{玛 7:19} 但若它顺利接受接枝，改变成好橄榄树，就成为结果实的橄榄树，仿佛栽植在君王的园囿（\OriginalTerm{乐园}{paradiso}）中：同样，人若真正因信德而向着更美好的事物前进，接受天主的圣神，并结出圣神的果实，将成为属神的人，如同栽植在天主的乐园中。但若他们抛弃圣神，仍保持原状，宁愿属于肉身而不属于圣神，那么，对这等人正当地说：\Quote{血肉不能承受天主的国；}\BibleRef{格前 15:50} 犹如有人说野橄榄不能进入天主的乐园一样。因此，宗徒在谈论血肉和野橄榄时，奇妙地展示了我们的本性以及天主普世的安排。因为正如好橄榄树，若一段时间被忽略，任其野生生长，就会自行变成野橄榄；又如野橄榄若被精心照料并接枝，就自然会回复到先前结果实的状态：同样，人若变得粗心，结出肉身的私欲如同木质产物，就因自己的过错而在正义上不结果实。因为当人睡觉时，仇敌撒下莠子的材料；\BibleRef{玛 13:25} 为此，主命令祂的门徒要警醒。再者，那些不结出正义果实，仿佛被荆棘覆盖而迷失的人，若勤勉努力，并接受天主的圣言作为接枝，\BibleRef{雅 1:21} 就达到人的原始本性——即按天主的肖像和模样受造的本性。\par

2. 然而，正如被接枝的野橄榄并未失去其木质实体，却改变了果实的品质，并得到另一个名字，如今不再是野橄榄，而是结果实的橄榄树，并被如此称呼；同样，当人因信德被接枝并接受天主的圣神时，他确实没有失去肉身的实体，却改变了果实（即他的行为）的品质，并得到另一个名字，\BibleRef{默 2:17} 表明他已变得更好，如今不再是单纯的血肉，而是属神的人，并被如此称呼。此外，再次如前，野橄榄若不被接枝，因其木质特性而对主人无用，就被砍下作为不结果实的树，丢在火里；同样，人若因信德不接受圣神的接枝，就留在旧的状态中，作为单纯的血肉，不能承受天主的国。因此，宗徒正确地宣告：\Quote{血肉不能承受天主的国；}\BibleRef{格前 15:50} 以及\Quote{属于血肉的人不能中悦天主。}\BibleRef{罗 8:8} 他并非藉这些话否认肉身的实体，而是表明圣神必须注入肉身之中。为此，他说：\Quote{这必朽坏的必须穿上不朽，这必死的必须穿上不死。}\BibleRef{格前 15:53} 他又宣告：\Quote{你们不属于血肉，而属于圣神，只要天主的圣神住在你们内。}\BibleRef{罗 8:9} 他更清楚地阐明这一点，说：\Quote{身体固然因罪而死亡，但神魂却因正义而生活。再者，如果那使耶稣从死者中复活者的圣神住在你们内，那使基督从死者中复活的，也必要藉那住在你们内的圣神，使你们必死的身体活过来。}\BibleRef{罗 8:10} 他还在\WorkTitle{罗马书}中说：\Quote{如果你们随从肉性生活，必要死亡。}\BibleRef{罗 8:13} 他藉这些话并非禁止他们在肉身中生活，因为他自己写信时也在肉身中；而是斩断肉身的私欲，那些引人死亡的东西。为此他接着说：\Quote{但你们若因圣神治死肉身的作为，必要生活。因为凡受天主圣神引导的，都是天主的子女。}\par

\SourceRecord{Translated by Alexander Roberts and William Rambaut.；Ante-Nicene Fathers；Vol. 1.；Edited by Alexander Roberts, James Donaldson, and A. Cleveland Coxe.；Buffalo, NY: Christian Literature Publishing Co.,；1885.；<http://www.newadvent.org/fathers/0103510.htm>.}

% structure: p0001 source=h1 target=section reason=page-title

\section{\WorkTitle{驳斥异端}（第五卷，第十一章）}

1. 宗徒预见不信者邪恶的言论，特别列举了他称为肉性的事；他解释自己，以免那些不诚实地曲解他意思的人有任何怀疑的余地，在\WorkTitle{致迦拉达人书}中这样说：\BibleQuote{本性私欲的作为是显而易见的：即淫乱、不洁、放荡、崇拜偶像、施行邪法、仇恨、竞争、嫉妒、忿怒、争吵、不睦、分党、妒忌、醉酒、宴乐，以及类似的事。我从前警告过你们，如今再警告你们：做这种事的人，决不能承受天主的国。}\BibleRef{迦 5:19} 这样他便更明确地向听众指出他所说的\BibleQuote{肉和血不能承受天主的国}是什么意思。因为做这些事的人，既然确实随从肉性，就没有能力活于天主。然后，他又继续告诉我们使人活的神性行动，即圣神的嫁接，如此说：\BibleQuote{圣神的效果是：仁爱、喜乐、平安、忍耐、良善、温和、忠信、柔和、节制：这样的事，没有法律禁止。}\BibleRef{迦 5:22} 因此，正如那已走向更好之事、结出圣神果实的人，因着与圣神的共融而完全得救；同样，那坚持上述肉性之事的人，真正被视为肉性的，因为他没有接受天主的圣神，就没有能力承受天国。又如同一宗徒向格林多人作证说：\BibleQuote{你们岂不知道，不义的人不得承继天主的国吗？你们不要自欺：无论是淫秽的、拜偶像的、犯奸淫的、做娈童的、搞同性恋的、偷窃的、贪婪的、酗酒的、辱骂的、勒索的，都不能承继天主的国。你们中有些人从前也这样；但是你们因主耶稣基督之名，并因我们天主的圣神，已经洗净了，已经圣洁了，已经成义了。}\BibleRef{格前 6:9} 他明明地指出，人是因什么而走向毁灭，如果他继续随从肉性生活；另一方面，他又指出人是因什么而得救。他说那救人的是我们的主耶稣基督之名，以及我们天主的圣神。\par

2. 既然在那段话中，他列举了那些没有圣神的、给[行这些事的人]带来死亡的肉性之事，他就在书信的结尾，按照他已经宣告的，喊道：\BibleQuote{就如我们怎样带了那属于土的肖像，也要怎样带那属于天上的肖像。弟兄们，我告诉你们：肉和血不能承受天主的国。}\BibleRef{格前 15:49-50} 现在他所说的\BibleQuote{就如我们怎样带了那属于土的肖像}，与已宣告的\BibleQuote{你们中有些人从前也这样；但是你们因主耶稣基督之名，并因我们天主的圣神，已经洗净了，已经圣洁了，已经成义了}相似。那么，我们什么时候带了那属于土的肖像呢？无疑就是当那些被称为\Quote{肉性之事}的行为在我们内发生时。然后，又是什么时候[我们带了]那属于天上的肖像呢？无疑就是当他说\BibleQuote{你们已经洗净了}，相信主的名，并接受祂的圣神时。现在我们已经洗去了，不是我们身体的实体，也不是我们[原初]形成的肖像，而是先前虚妄的生活。因此，在这些肢体中，我们曾因行败坏之事而走向死亡，正是在这些肢体中，我们因行圣神之事而得以活过来。\par

\SourceRecord{Translated by Alexander Roberts and William Rambaut.；Ante-Nicene Fathers；Vol. 1.；Edited by Alexander Roberts, James Donaldson, and A. Cleveland Coxe.；Buffalo, NY: Christian Literature Publishing Co.,；1885.；<http://www.newadvent.org/fathers/0103511.htm>.}

% structure: p0001 source=h1 target=section reason=page-title

\section{驳斥异端（第五卷，第十二章）}

1. 因为正如肉体能朽坏，它也能不朽坏；正如它能死亡，它也能获得生命。这两者互相让位，不能同时存在于同一处所，而是其一被另一驱逐，一者的在场摧毁另一者。那么，如果死亡占据了人，它就把生命从他身上驱走，证明他是死的，那么生命一旦获得对人的权能，岂不更赶走死亡，使他作为活人归于天主？因为如果死亡带来必死性，那么生命来临时，为何不该使人活起来？正如依撒意亚先知所说：\BibleQuote{死亡吞噬，因为它得了胜。} 又说：\BibleQuote{天主擦去了每张脸上的每一滴眼泪。} 因此，先前的生命被驱逐，因为它不是由圣神所赐，而是由气息所赐。\par

2. 因为生命的气息——这使人为有灵的生物——是一回事，而赋予生命的圣神——这又使人成为属神的——是另一回事。为此，依撒意亚说：\BibleQuote{上主如此说，祂创造了天并奠定了它，奠定了地及其中的万物，将气息赐给地上的百姓，将圣神赐给行在地上的人；}\BibleRef{依 42:5} 这告诉我们，气息确实普遍赐给地上所有的人，但圣神只属于那些践踏世俗欲望的人。因此，依撒意亚自己区分上述事物，再次呼喊：\BibleQuote{因为圣神将由我而出，我造了每口气息。}\BibleRef{依 57:16} 这样，他将圣神视为天主独有的，在末期藉着义子收养倾注于人类；但气息是普遍受造的，他指出它是受造之物。受造之物与造物者不同。因此，气息是暂时的，而圣神是永恒的。气息亦在短时间内增长，存续一段时间；之后它离去，使先前的居所缺欠气息。但当圣神充满人内外时，因它常驻，永不离开他。\BibleQuote{但属神的不是在先，} 宗徒说，这话仿佛是论及我们人类；\BibleQuote{而是属生灵的在前，以后才是属神的，}\BibleRef{格前 15:46} 这是合理的。因为首先必须造人，受造者必须接受灵魂；之后它才这样接受圣神的共融。因此，\BibleQuote{第一亚当由主造成一个活的灵魂，第二亚当是赋予生命的神。}\BibleRef{格前 15:45} 那么，正如那成为活灵魂的人因转向邪恶而丧失生命，同样，同一个人当他转向良善并接受赋予生命的圣神时，必将获得生命。\par

3. 因为死的并非一个东西，活的又是另一个，正如失落的不是一个，被寻回的又是另一个，而是主来寻找那同一只失落的羊。那么，什么死了呢？无疑是肉体的实体；同样，它也失去了生命的气息，成为无气息而死的。因此，主来使之复活的就是这个，为使我们如同在亚当内因属生灵的本性都死了，在基督内因成为属神的都活起来，不是撇弃天主的工程，而是撇弃肉体的私欲，并接受圣神；正如宗徒在\WorkTitle{致哥罗森人书}中所说：\BibleQuote{所以，你们要致死你们在地上的肢体。} 它们是什么，他自己解释：\BibleQuote{淫乱、不洁、邪情、恶欲；贪婪，即崇拜偶像。}\BibleRef{哥 3:5} 宗徒所宣讲的正是撇弃这些；他声明行这些事的人，由于仅仅是血肉，不能承受天国。因为他们的灵魂趋向更坏的事，堕入地上的情欲，成了与这些情欲（即\Quote{属地的}）相同名称的分享者；宗徒命令我们撇弃这些时，在同一书信中说：\BibleQuote{脱去旧人和他的行为。}\BibleRef{哥 3:9} 但他说这话时，并未除去人古老的构造；因为那样我们就得藉自杀摆脱它的陪伴了。\par

4. 但宗徒本人也一样，他曾成胎于母腹并由此而出，写信给我们，在\WorkTitle{致斐理伯人书}中承认\BibleQuote{活在肉身中，是[他]工作的果实；}\BibleRef{斐 1:22} 他这样表达。圣神工作的最终果实就是肉身的得救。因为那无形的圣神还有什么别的可见果实，而不是使肉身成熟并堪当不朽呢？那么，如果[他说]：\Quote{活在这肉身中，对我是工作的结果，} 他当然不是在说\BibleQuote{脱去旧人和他的行为}那句话时轻蔑肉身的实体，\BibleRef{哥 3:9} 而是指出我们应脱去从前的行为，那衰老而败坏的行为；为此他接着说：\BibleQuote{穿上新人，那按创造者的肖像在知识上更新的。}\BibleRef{哥 3:10} 因此，他说\Quote{在知识上更新}时，他表明同一人，过去无知，即对天主的无知，藉着那关于祂的知识而被更新。因为认识天主更新人。当他说\Quote{按创造者的肖像}时，他陈述了同一人的\OriginalTerm{重演}{recapitulation}，这人在起初是按天主的肖像而造的。\par

5. 而他，那位宗徒，正是那位从母胎出生、即由古代的肉体实质所生的人，他本人在\WorkTitle{迦拉达书}中声明说：\BibleQuote{但是，从母胎中已选拔了我，以恩宠召叫我的天主，却决意将他的圣子启示给我，叫我在外邦人中传扬他，}\BibleRef{迦 1:15} 并非如我已指出的，一位是从母胎出生的，另一位是宣讲天主圣子的福音者；而是同一个人，他先前无知，曾迫害教会，但当从天上获得启示、主与他交谈之后——如我在第三卷中所指出的——便宣讲了在般雀·比拉多下被钉死的耶稣基督、天主圣子的福音；他先前的无知被后来的知识驱除：正如主所治好的瞎子们，他们确实失去了瞎眼，但接受了他们眼睛的完好本质，并在他们先前看不见的同一双眼睛中获得了视觉能力；黑暗仅被视力驱散，而眼睛的本质被保留，为的是通过他们曾看不见的那双眼睛，重新运用视觉能力，感谢那位使他们重见光明者。同样，那位手枯萎而被治好的人，以及所有被治好的人，通常并没有改变他们出生时从母胎出来的那些身体部分，而是仅仅重新获得了它们健康的状态。\par

6. 因为万有的创造者、天主的圣言，他从起初就塑造人，当他发现自己的工程被邪恶损害时，便在其上施行各种治疗。有时［祂这样做］，是针对每一个别的肢体，如同在他自己的工程中所发现的；另一次他一次性将人完全复原，在所有方面健全无损，预备他完善无缺，为他自己，直到复活。因为他治疗肉身的部分、并将它们恢复至原初状态，其目的是什么？如果那些被他治好的部分不能获得救恩？因为如果他所赐予的［仅仅］是暂时的益处，那么他对那些他治疗的对象就没有赐予什么重要的东西。或者，他们怎能坚持说肉身不能接受从他涌出的生命，既然它从他接受了治疗？因为生命是通过治疗带来的，而不朽是通过生命带来的。因此，那赐予治疗的，也同样赐予生命；而那位［赐予］生命的，也以不朽环绕他自己的工程。\par

\SourceRecord{Translated by Alexander Roberts and William Rambaut.；Ante-Nicene Fathers；Vol. 1.；Edited by Alexander Roberts, James Donaldson, and A. Cleveland Coxe.；Buffalo, NY: Christian Literature Publishing Co.,；1885.；<http://www.newadvent.org/fathers/0103512.htm>.}

% structure: p0001 source=h1 target=section reason=page-title

\section{\WorkTitle{驳斥异端（第五卷，第十三章）}}

1. 让我们的对手——即那些反对自己得救的人——告诉我们这一点：大司祭的女儿死了；寡妇的儿子死了，被抬出城门口附近；\BibleRef{路 7:12} 以及躺在坟墓里已四天的拉匝禄，\BibleRef{若 9:30} ——他们是在什么身体里复活的？无疑是在他们死去时的同一个身体里。因为如果不是完全相同的身体，那么死去的人本身就没有复活。因为（圣经）说：\Quote{主拉着死人的手，对他说：\InnerQuote{青年人，我对你说，起来。}那死人便坐起来，主就吩咐给他一些东西吃，并把他交给他母亲。}\newline{} 祂又大声喊叫拉匝禄说：\Quote{拉匝禄，出来！那死人便出来了，手脚裹着布带。}这是那被罪恶捆绑之人的预象。因此主说：\Quote{解开他，让他去。}\newline{} 因此，正如那些被治愈的人在从前受苦的肢体上得到痊愈，死者在相同的身体中复活，他们的肢体和身体得到健康，以及主所赐予的生命——主借暂时之事预表永恒之事，并表明祂自己能够将医治和生命赐予祂的创造物，好使祂关于（未来）复活的言语得以被相信——同样，在终结之时，当主借\Quote{最后的号角}\BibleRef{格前 15:52} 发出声音时，死人将要复活，正如祂自己所宣告的：\Quote{时候要到，凡在坟墓里的，都要听见人子的声音，就出来；行善的，复活进入生命；作恶的，复活受审判。}\BibleRef{若 5:28}\par

2. 因此，那些不愿意看见如此明显清晰之事、反而逃避真理之光、像悲剧中的俄狄浦斯一样盲目自己的人，实在是虚空且可悲。正如那些不熟练摔跤的人，与别人角力时，死死抓住对手身体的某一部分，却因所抓之处反而跌倒；而跌倒时，却以为自己得胜，因为固执地抓住了起初抓住的部位，除了跌倒之外还成了笑柄；同样，关于异端者那句（喜爱的）话：\Quote{血肉不能承受天主的国}——他们取用保禄的两句话，却没有理解宗徒的含义，也没有仔细考究词语的力量，只是死死抓住词语本身，因而被它们所害（\OriginalTerm{围绕它们}{\GreekText{περὶ} \GreekText{αὐτάς}}），尽可能颠覆了天主的整个计划。\par

3. 因为他们会这样主张，说这段经文指严格意义上的血肉，而不是血肉的行为，正如我所指出的，这样就把宗徒说成自相矛盾。因为紧接着在同一封书信中，他明确地说，关于血肉这样说：\Quote{因为这可朽坏的必须穿上不可朽坏的，这可死的必须穿上不可死的。当这可死的穿上不可死的时候，就应验了经上的话：\InnerQuote{死亡被吞灭于胜利中。死亡，你的胜利在哪里？死亡，你的刺在哪里？}}\BibleRef{格前 15:53} 这些话将适时地说出，当这可死可朽坏的血肉——它屈服于死亡，也被某种死亡的统治所压迫——进入生命，穿上不可朽坏和不可死的时候。因为那时，死亡确实被战胜了，当那被它压制的血肉从它的统治下走出来时。\newline{} 又对斐理伯人说：\Quote{但我们的家乡是在天上，我们也等待救主，主耶稣基督，祂将改变我们卑贱的身体，相似祂光荣的身体，依祂所能（\OriginalTerm{如此能}{ita ut possit}）按祂大能的运行。}\BibleRef{斐 3:29} 那么，这\Quote{卑贱的身体}是什么，主将改变它，相似\Quote{祂光荣的身体}？显然，这是由血肉组成的身体，当它落入泥土时确实是卑贱的。它的改变（如此发生）：它本是可死可朽坏的，变成不可死不可朽坏的，不是按它自己的本质，而是按主的大能运行，祂能将可死的穿上不可死的，将可朽坏的穿上不可朽坏的。因此他说：\Quote{使可死的被生命吞灭。那为我们完成这事的天主，也赐给我们圣神作抵押。}\BibleRef{格后 5:4} 他这些话极其明显地指血肉；因为灵魂不是可死的，圣神也不是。现在，可死的被生命吞灭，那时血肉不再死去，而是活着且不朽坏，颂扬那为我们完成这事的天主。因此，为使我们得以完成这事，他恰当地对格林多人说：\Quote{你们要在身体上光荣天主。}\BibleRef{格前 6:20} 而天主就是那赐予不朽的。\par

4. 他向格林多人清楚、无疑且毫不含糊地声明，他这些话是指肉躯的身体，而非别的：\BibleQuote{我们身上时常带着耶稣的死，为使耶稣基督的生命也彰显在我们身上。因为我们活着的人，为耶稣的缘故常被交于死，为使耶稣的生命也彰显在我们可朽的肉身上。}\BibleRef{格后 4:10} 关于圣神抓住肉躯，他在同一封书信中说：\BibleQuote{你们就是基督的书信，由我们供职所写，不是用墨，而是用生活天主的圣神；不是写在石版上，而是写在血肉的心版上。}\BibleRef{格后 3:3} 因此，如果现今血肉的心已分享圣神，那么在复活时，它们领受由圣神赐予的生命，又有何惊奇呢？关于这复活，宗徒在致斐理伯人书信中说：\BibleQuote{使我相似他的死，我希望我无论如何能到达从死者中的复活。}\BibleRef{斐 3:11} 那么，除了那因承认天主而被置于死的实体之外，还能在别的可朽肉身上理解生命被彰显呢？——正如他本人所宣布的：\BibleQuote{如果我在厄弗所像一个普通人那样与野兽搏斗，于我有什么益处呢？因为如果死人不复活，基督也就没有复活。如果基督没有复活，我们的宣讲便是空的，你们的信德也是空的。并且，我们也被发现是为天主作假见证，因为我们曾为天主作证说祂使基督复活了，但[根据那种假设]祂并没有使祂复活。因为如果死人不复活，基督也就没有复活。但如果基督没有复活，你们的信德便是空的，你们还在你们的罪中。那么，那些在基督内长眠的人也丧亡了。如果我们只在今生对基督抱有希望，我们就是众人中最可怜的了。但是，基督实在从死者中复活了，做了长眠者的初果，因为死亡既借着一人而来，死者的复活也借着一人而来。}\BibleRef{格前 15:13}\par

5. 因此，在所有这些段落中，正如我已经说过的，这些人必须要么声称宗徒在关于\BibleQuote{肉和血不能承受天主的国}这句话上自相矛盾；要么，另一方面，他们将被强迫对所有段落作出歪曲和扭曲的解释，以颠覆和改变词语的意义。因为他们若试图以别种方式解释他所写的\BibleQuote{因为这可朽坏的必须穿上不可朽坏的，这必死的穿上不死}\BibleRef{格前 15:53}，以及\BibleQuote{为使耶稣的生命彰显在我们可朽的肉身上}\BibleRef{格后 4:11}，以及所有其他宗徒明确清楚地宣告肉躯复活与不朽的段落时，他们还能说出什么有理智的话呢？如此，那些不愿正确理解的人，将被逼对这些段落作出错误的解释。\par

\SourceRecord{Translated by Alexander Roberts and William Rambaut.；Ante-Nicene Fathers；Vol. 1.；Edited by Alexander Roberts, James Donaldson, and A. Cleveland Coxe.；Buffalo, NY: Christian Literature Publishing Co.,；1885.；<http://www.newadvent.org/fathers/0103513.htm>.}

% structure: p0001 source=h1 target=section reason=page-title

\section{驳斥异端（卷五，第十四章）}

1. 既然宗徒并没有断言血肉的本质本身不能承受天主的国，那么这位宗徒在论及主耶稣基督时，也到处使用了\Quote{血肉}一词。一方面是为确立祂的人性（因为祂自称为人子），另一方面是为确认我们血肉的救恩。因为如果血肉不能得救，天主的圣言就不会成为血肉。如果义人的血不被追问，主当然就不会有血（在祂的组成中）。但既然从（世界）起初血就发出呼号（\Emph{vocalis est}），天主在\Person{加音}杀死他的弟弟时，就对\Person{加音}说：\Quote{你弟弟的血的声音从地里向我哀告。} \BibleRef{创 4:10} 而且，因为他们的血将要被追问，祂对与\Person{诺厄}同船的人说：\Quote{我要向一切野兽追讨你们灵魂的血；} 又说：\Quote{凡流人血的，他的血也要为人所流。} 同样，主也对那些将要流祂血的人说：\Quote{凡地上所流的义人的血，从义人\Person{亚伯尔}的血起，直到你们在圣所与祭坛之间所杀的\Person{巴辣基雅}的儿子\Person{匝加利亚}的血，都要归到你们身上。我实在告诉你们：这一切都要归于这一代。} 祂由此指出，从起初所有义人和先知的血的倾流，将在祂自己的身上发生\OriginalTerm{重演}{recapitulation}，并且藉着祂自己，他们的血将被追讨。这血若不具有被救的能力，就不能被追讨；主也不会将这些事总结在祂自己身上，除非祂是按照起初受造的方式成了血肉，在末后在自己的身上拯救了那在起初于\Person{亚当}内所丧失的。\par

2. 但如果主为了别的秩序而成了血肉，并取了另一种本质的血肉，那么祂就没有在自己的身上总结人性，那样祂也不能被称为血肉。因为血肉真是由那最初从尘土所塑造之物的传递而形成的。但如果祂必须从另一种本质取得（祂身体的）材料，圣父起初就会从另一种本质（而不是实际所用的）塑造（血肉的）材料了。但如今情况是，圣言拯救了那确实（受造的）——即那已丧失的人性，藉着自身实现与它的共融，并寻求它的救恩。但那已丧失的拥有血肉。因为主从地上取尘土塑造了人；所有主来临的救恩计划都是为了他。因此，主自己拥有血肉，在祂自己身上\OriginalTerm{总结}{recapitulating}的并非别的，而是圣父那最初的创造工程，寻求那已丧失的。为此，宗徒在致哥罗森人的书信中说：\Quote{你们从前也与天主隔绝，并因邪恶的行为与祂为敌，但现在祂藉着祂血肉的身体，藉着死亡，使你们与自己和好，好使你们成为圣洁、无瑕、无可指责地呈现在祂面前。} \BibleRef{哥 1:21} 他说\Quote{你们在祂血肉的身体里与祂和好}，因为义的血肉使那在罪中被束缚的血肉与天主和好，并使之与天主成为朋友。\par

3. 因此，如果有人断言主的血肉在这方面与我们的不同——即它确实没有犯罪，祂的灵魂里也没有诡诈，而我们却是罪人——他说的倒是事实。但如果他声称主拥有另一种血肉的本质，那么关于和好的话就不会适用于那个人了。因为那和好的东西从前曾是处于敌对的。如果主取了另一种本质的血肉，祂就不会藉此将那因过犯而成为仇敌的与天主和好。但如今，主藉着与自己的共融，使人与天主圣父和好，藉着祂自己血肉的身体使我们与祂自己和好，并用祂自己的血救赎了我们，正如宗徒致厄弗所人所说：\Quote{在祂内，我们藉着祂的血得蒙救赎，过犯得到赦免；} \BibleRef{弗 1:7} 又对他们说：\Quote{你们从前远离的人，如今在基督耶稣内，藉着祂的血，成为亲近的了；} \BibleRef{弗 2:13} 又说：\Quote{以祂的肉身，废掉了那含规条命令的法律，所导致的仇视。} \BibleRef{弗 2:15} 宗徒在每一封书信中都明确作证，我们是藉着我们主的血肉和祂的血而得救的。\par

4. 因此，如果血肉是那些为我们获取生命的事物，那么，并非在\OriginalTerm{字面意义}{proprie}上说血肉不能承受天主的国；而是指那些已经提及的属肉体的行为，它们将人扭曲入罪，剥夺其生命。为此，他在\WorkTitle{罗马书}中说：\BibleQuote{因此，不要让罪在你们必死的身体上为王，使你们顺从其控制；也不要把你们的肢体献给罪，作不义的武器；但要献上你们自己给天主，如同从死者中复活的人，并把你们的肢体献给天主，作正义的武器。}\BibleRef{罗 6:12}因此，正是在这些同样的肢体中，我们曾用来服事罪，结出死亡的果子，他愿意我们服从正义，以便结出生命的果子。所以，我亲爱的朋友，请记住，你已因我们主的血肉被救赎，因他的血被重建；并\BibleQuote{持守着头，全身——教会——既已配合紧密，就从他得到增长}\BibleRef{哥 2:19}——即，承认天主圣子降生成人的来临及其\OriginalTerm{神性}{deum}，并恒常期待他的\OriginalTerm{人性}{hominem}，同时善用这些从圣经中得出的证据——如我所指出的，你就能轻易推翻那些后来编造出来的异端者的一切观念。\par

\SourceRecord{Translated by Alexander Roberts and William Rambaut.；Ante-Nicene Fathers；Vol. 1.；Edited by Alexander Roberts, James Donaldson, and A. Cleveland Coxe.；Buffalo, NY: Christian Literature Publishing Co.,；1885.；<http://www.newadvent.org/fathers/0103514.htm>.}

% structure: p0001 source=h1 target=section reason=page-title

\section{驳斥异端（第五卷，第十五章）}

1. 现在，那在起初创造人的那一位，曾应许人在死后要得到第二次出生，依撒意亚这样宣告说：\BibleQuote{死人要复活，在坟墓中的要起来，在地下的人要欢欣。因为从你而来的甘露是他们的健康。}\BibleRef{依 26:19} 又说：\BibleQuote{我要安慰你们，你们要在耶路撒冷得到安慰；你们要看见，你们的心要欢欣，你们的骨骸要像青草一样茂盛；上主的手要为敬拜祂的人所认识。}\BibleRef{依 66:13} 厄则克耳这样说：\BibleQuote{上主的手临于我，上主藉圣神领我出去，把我放在平原中间，那地方满是骨骸。祂使我绕着它们走，看啊，平原上有很多骨骸，非常干枯。祂对我说：人子，这些骨骸能活吗？我说：上主，是祢知道。祂对我说：你向这些骨骸说预言，对他们说：干枯的骨骸，听上主的话！上主对这些骨骸这样说：看，我要使生命的气息进入你们里面，我要给你们加上筋，使肉长在你们身上，并披上皮肤，将我的圣神放在你们内，你们就活了；你们要知道我是上主。我遂遵命说预言。正在我说预言时，忽有地震，骨骸互相结合起来，每个骨头都连到自己的关节上。我看看，看啊，有筋和肉生在它们上面，皮肤也包在上面，但它们内里还没有气息。祂对我说：你向气息说预言，人子，对气息说：上主这样说：气息，从四方而来，\OriginalTerm{风}{spiritibus}，吹在这些死人身上，使他们活过来。我遵命说预言，气息就进入他们内，他们便活了，站起来，成为一支极庞大的军队。}\BibleRef{则 37:1} 他又说：\BibleQuote{上主这样说：看，我要打开你们的坟墓，使你们从坟墓中出来，带你们进入以色列地；你们要知道我是上主，当我打开你们的坟墓，为领我的百姓从坟墓中出来时，我要将我的圣神放在你们内，你们就活了；我要安置你们在你们的地上，你们要知道我是上主。我说了，必要实行——上主说。}\BibleRef{则 37:12} 我们立刻觉察到，此处所描述的创造者\OriginalTerm{工匠之神}{Demiurgo}使我们死亡的身体复活，应许它们复活并从坟墓和墓葬中苏醒，并赐予它们不朽（祂说：\BibleQuote{因为他们的日子要像生命树的日子一样}\BibleRef{依 65:22}），祂被证明是成就这些事的唯一天主，并且祂本身就是良善的圣父，慈爱地将生命赐给那些本身没有生命的人。\newline{}\par

2. 为此，主极其明确地向祂的门徒显示了自己和圣父，免得他们也许在祂以外——那位造人并赐给人生命气息的主——寻找另一位天主，也免得人疯狂到在创造者以上捏造另一位父。祂也用一句话治愈了所有因罪而软弱的人，祂对他们说：\BibleQuote{看，你已痊愈了，不要再犯罪，免得遭遇更不幸的事。}\BibleRef{若 5:14} 这指明因着悖逆之罪，疾病临到了人。然而，对于那个生来瞎眼的人，祂不是用一句话，而是通过外在行动赐他看见；这样做并非无目的或偶然，而是为了显示天主的手——那在起初塑造人的手。因此，当祂的门徒问祂，那人为何生来瞎眼，是他自己的罪还是他父母的罪，祂回答说：\BibleQuote{不是他犯了罪，也不是他的父母，而是为叫天主的工作在他身上显出来。}\BibleRef{若 9:3} 天主的工作就是塑造人。因为圣经说，祂通过一种过程造了人：\BibleQuote{上主用地上的尘土形成了人。}\BibleRef{创 2:7} 因此，主也在地上吐唾沫，和泥，抹在瞎子的眼睛上，表明了起初的塑造过程是如何完成的，并向那些能理解的人显示了天主的手——人是从尘土中被那只手所造的。因为那位工匠圣言在母腹中所未完成的部分（即瞎子的眼睛），祂后来在公开场合补全了，为叫天主的工作在他身上显出来，使我们在被造时不再寻找另一只手，也不再寻找另一位父；我们明白，那当初塑造我们、并在母腹中塑造我们的天主之手，在末期寻找我们这些迷失的人，领回祂自己的羊，将亡羊扛在肩上，欢欢喜喜地带回生命的羊栈。\newline{}\par

3. 如今，天主的圣言在母腹中塑造我们，祂对耶肋米亚说：\BibleQuote{我还没有在母腹中形成你以前，就已认识了你；你还没有从母胎出生以前，我已祝圣了你，并立你为万民的先知。}\BibleRef{耶 1:5} 保禄也同样说：\BibleQuote{但是，当那位从母胎中已选拔了我，并借祂的恩宠召叫我的天主，愿意将祂的子启示给我，叫我在外邦人中传扬祂时，}\BibleRef{迦 1:15} 因此，既然我们是由圣言在母腹中塑造的，这同一圣言也在那生来瞎眼的人身上塑造了视觉的能力；公开显示了那位在隐秘中塑造我们的那一位，因为圣言自己已显示给人类；并宣告了亚当的最初形成，以及他受造的方式，和由哪一只手塑造，从部分指出整体。因为形成视觉能力的主，就是创造整个人的那位，他实现圣父的旨意。而由于人，就那在亚当之后的形成而言，已陷入过犯，需要重生的圣洗，主（在给他涂上泥以后）对他说：\BibleQuote{去史罗亚洗洗罢；}\BibleRef{若 9:7} 这样，他恢复了双重的祝圣，以及那藉由圣洗而产生的重生。因此，当他洗后，他看见而来，为使他能认识那塑造他的主，也为使人能认识那位赐予他生命的主。\par

4. 因此，所有瓦伦廷的追随者都输掉了他们的案子，因为他们说人不是由这泥土塑造的，而是由一种流动且分散的物质。因为，主为那人形成眼睛所用的泥土，显然人原先也是由这同一种泥土塑造的。因为眼睛由一种来源形成而身体其余部分由另一种来源形成，是不相容的；同样，由一位塑造身体，而由另一位塑造眼睛，也是不相容的。但是，那位在起初塑造亚当的同一主，圣父也曾与祂说话：\BibleQuote{让我们按我们的肖像和模样造人，}\BibleRef{创 1:25} 在这末世祂向人类启示自己，为那（从亚当所承受的身体中）瞎眼的人形成了视觉器官（\OriginalTerm{视觉}{visionem}）。因此，圣经指出将要发生的事，说：当亚当因不服从而躲藏时，主在傍晚来到他那里，呼唤他，并说：\BibleQuote{你在哪里？}\BibleRef{创 3:9} 这意味着在末世，这同一圣言来呼唤人，提醒他的行为，在此行为中他曾躲避主。因为正如那时天主在傍晚对亚当说话，寻找他；同样在末世，藉着同一声音，寻找他的后代，祂已访问了他们。\par

\SourceRecord{Translated by Alexander Roberts and William Rambaut.；Ante-Nicene Fathers；Vol. 1.；Edited by Alexander Roberts, James Donaldson, and A. Cleveland Coxe.；Buffalo, NY: Christian Literature Publishing Co.,；1885.；<http://www.newadvent.org/fathers/0103515.htm>.}

% structure: p0001 source=h1 target=section reason=page-title

\section{驳斥异端（第五卷，第十六章）}

1. 既然亚当是从我们所属的这泥土所塑造，圣经告诉我们，天主向他说：\Quote{你必须汗流满面，才有饭吃，直到你归于你由之所出的尘土。}\BibleRef{创 3:19} 如果死后我们的身体归于任何其他物质，那么它们也必是从那物质而来。但如果它恰恰归于这［尘土］，那么显然人的形体也是从它而被造的；正如主清楚地显示，当祂从这同一物质为那［被祂赐予视力的］人塑造了眼睛时。如此，天主的手清楚地显明了，亚当藉此被塑造，我们也被塑造；既然有同一圣父，祂的声音从起初到终结始终临在于祂的工程，而我们被塑造的物质也藉福音清楚地宣告，我们就不该寻求祂以外的另一位父，也不该［寻求］我们被塑造的另一物质，除了那预先提及并由主所显示的；也不该寻求另一只天主的手，除了那从起初到终结塑造我们、预备我们得生命、临在于祂的工程、并按照天主的肖像和样式使之完美的那一只。\par

2. 再者，这圣言得以彰显，是在天主的圣言成了人时，祂使自己相似于人，也使人相似于祂自己，好使人因与圣子相似而成为父所珍视的。因为在远古时代，人只是被\Emph{说}是按天主的肖像所造，但并未被\Emph{显示}出来；因为圣言尚不可见，人是按祂的肖像被造的。因此人也轻易失去了这相似性。然而，当天主的圣言成了血肉时，祂确认了这两点：祂既真实地显示了肖像，因为祂自己成了那肖像；又按确定的方式重建了相似性，藉可见的圣言使人相似于不可见的圣父。\par

3. 主不仅藉上述之事彰显了自己，也藉祂的苦难［彰显了自己］。因为消除那起初因一棵树而发生的人类的悖逆，祂\Quote{成了服从至死，且死在十字架上}\BibleRef{斐 2:8}；藉着那在［十字］木上的服从，矫正了因一棵树而发生的悖逆。如果祂宣告另一位父，祂就不会藉着同一［肖像］来消除那对我们造物主所犯的悖逆。然而，正因为我们藉着这些事物悖逆了天主，没有相信祂的话，所以祂也藉着同样的事物带来了对祂圣言的服从与认同；藉着这些事物，祂清楚地显示了天主自己，那位我们确实在首先的亚当中冒犯了，因他没有遵守祂的诫命。在第二个亚当中，我们却被和好，成为服从至死的。因为我们只欠那一位的债，就是我们起初违犯了其诫命的那一位。\par

\SourceRecord{Translated by Alexander Roberts and William Rambaut.；Ante-Nicene Fathers；Vol. 1.；Edited by Alexander Roberts, James Donaldson, and A. Cleveland Coxe.；Buffalo, NY: Christian Literature Publishing Co.,；1885.；<http://www.newadvent.org/fathers/0103516.htm>.}

% structure: p0001 source=h1 target=section reason=page-title

\section{驳斥异端（第五卷，第十七章）}

1. 现在，这个存有就是\OriginalTerm{创造者}{Demiurgus}，他因他的爱而被称为圣父，因他的权能而被称为主，因他的智慧而被称为我们的造物主和形成者；我们因违犯了他的诫命而成为他的仇敌。因此，在末期，主藉着他的降生成人使我们重新和好，成为\BibleQuote{天主与人之间的中保} \BibleRef{弟前 2:5}，确实为我们平息了我们曾冒犯的圣父，并藉着他自己的服从，\OriginalTerm{取消}{consolatus}了我们的悖逆；还赐予我们与造物主共融和服从的恩赐。为此，他也教导我们在祈祷中说：\BibleQuote{宽免我们的罪债} \BibleRef{玛 6:12}，因为他确实是我们欠了债的圣父，我们违犯了他的诫命。但这存有是谁？是某位未知者，是那位不向任何人颁布诫命的圣父吗？还是那位在圣经中被宣告的天主，我们欠了他的债，违犯了他的诫命？因为诫命是由圣言赐给人类的。据说，亚当\BibleQuote{听见了上主天主的声音} \BibleRef{创 3:8}。因此，他的圣言正当地对人说：\BibleQuote{你的罪赦了} \BibleRef{玛 9:2}；那在起初我们所冒犯的同一位置，在终结时赐予罪过的赦免。但如果我们确实违背了另一位神的命令，而不同的神说：\BibleQuote{你的罪赦了} \BibleRef{玛 9:2}，那么这位神既不善良，也不真实，更不正义。因为他若不给出属于他自己的东西，怎能是善良的？或他若掠夺别人的财物，怎能是正义的？罪过又怎能真正被宽恕，除非我们所冒犯的那一位亲自藉着\BibleQuote{我们天主慈悲的心肠}来赐予宽恕，在其中\BibleQuote{他眷顾了我们} \BibleRef{路 1:78}，藉着他的圣子？\par

2. 因此，当他治好那患瘫痪的人时，[圣史]说：\BibleQuote{群众看见，就光荣了天主，因为他赐给了人们如此大的权柄} \BibleRef{玛 9:8}。那么，旁观者光荣了哪一位天主？果真是异端者们所杜撰的那位未知的父吗？他们怎能光荣那位完全不为他们所知的神呢？所以，显然，以色列子民光荣了那位藉着法律和先知被宣告为天主的，他也是我们主之圣父；因此，他藉着所行的那些标记，通过人的感官证据，教导人光荣天主。然而，如果他本人来自另一位父，而人看见他的奇迹时光荣了另一位不同的父，那么他就使他们对于那位派遣了治愈恩赐的圣父忘恩负义。但由於独生子从那位天主那里来为拯救人类，他既藉着他一向所行的奇迹激发不信的人光荣圣父；又对法利塞人说——他们不承认他圣子的来临，因而不相信藉着他所赐予的罪赦：\BibleQuote{为叫你们知道，人子在地上有赦罪的权柄} \BibleRef{玛 9:6}。他这样说了，就命令那瘫痪的人拿起他所躺的床，回家去。他以此行为使不信者困惑，并表明他自己就是天主的声音，人藉此领受了诫命，却违背了它而成为罪人；因为瘫痪是罪的后果。\par

3. 因此，藉着赦免罪过，他确实治愈了人，同时也显示了他自己是谁。因为如果除了天主独自以外，无人能赦免罪过，而主赦免了罪并治愈了人，那么显然，他自己就是成为人子的天主圣言，从圣父那里领受了赦罪的权柄；因为他是人，因为他是天主，为的是：作为人，他为我们受苦，作为天主，他怜悯我们，并宽免我们的罪债，我们因这罪债而成为欠债于天主我们创造者的人。因此，达味预先说：\Quote{罪恶得赦免，过犯得遮掩的人，是有福的。上主不归咎于他的人，是有福的。} 从而指出那随着他来临而来的罪赦，藉此他\BibleQuote{销毁了字据}我们欠债的，并将它\BibleQuote{钉在十字架上} \BibleRef{哥 2:14}，好使我们藉着一块木头向天主欠债，同样也藉着一块木头获得债务的赦免。\par

4. 许多其他人，尤其是藉着先知厄里叟，已经显著地表明了这一点。因为当他的先知同伴们砍伐木材建造帐幕时，铁斧头从斧柄上脱落，掉进了约但河，他们找不到它。厄里叟来到那地方，得知所发生的事后，他将一些木头扔进水里。他这样做之后，斧头的铁部分浮了上来，他们从水面上捞起了先前失去的东西。\BibleRef{列下 6:6} 藉此行动，先知指出，天主的确定圣言，我们曾因一棵树而疏忽地失去，并且无法再找到它，我们将藉着一棵树的安排（即基督的十字架）重新领受它。因为天主圣言被比作斧头，洗者若翰就此宣告说：\Quote{现在斧子已经放在树根上。}\BibleRef{玛 3:10} 耶肋米亚也说了同样的话：\Quote{天主圣言如斧头劈开岩石。}\BibleRef{耶 23:29} 因此，正如我已经指出的，这向我们隐藏的圣言，藉着树的安排彰显出来。因为正如我们藉着一棵树失去了它，同样藉着一棵树它再次向众人彰显，在其自身中显示了高度、长度、宽度和深度；并且，正如我们的一位前辈所观察到的：\Quote{藉着一位神圣人物双手的伸展，将两个民族聚集于同一天主之前。} 因为这是两只手，因为有两个民族分散到地极；但中间有一个头，因为只有一位天主，祂超越万有，贯乎万有，也在我们众人之内。\par

\SourceRecord{Translated by Alexander Roberts and William Rambaut.；Ante-Nicene Fathers；Vol. 1.；Edited by Alexander Roberts, James Donaldson, and A. Cleveland Coxe.；Buffalo, NY: Christian Literature Publishing Co.,；1885.；<http://www.newadvent.org/fathers/0103517.htm>.}

% structure: p0001 source=h1 target=section reason=page-title

\section{\WorkTitle{反驳异端（第五卷，第十八章）}}

1. 天主并未藉他人的受造物，而是藉自己的受造物，成就了如此伟大的一项救恩计划；并非藉那些出于无知与缺陷的产物，而是藉那些由其圣父的智慧与能力所孕育的受造物。因为祂既非不义，以致贪图他人的产业；也非匮乏，以致不能凭自己的能力将自己的生命赋予自己的受造物，并藉自己的受造物为人类带来救恩。的确，受造物无法支撑祂（在十字架上），如果祂所派遣的（仅凭委托）是无知与缺陷的产物。我们已屡次指出，降生成人的天主圣言曾被悬挂在木头上，连异端分子自己也承认祂被钉在十字架上。那么，无知与缺陷的产物如何能支撑那包含万有知识、真实而完美的天主圣言呢？或者，那向父隐藏、远离父的受造物如何能承受祂的圣言呢？如果这个世界是由天使造的（我们无需假定他们对至高天主是无知还是知晓），当主宣告说：\BibleQuote{因为我在父内，父也在我内} \BibleRef{若 14:11}，那么这天使的作品如何能同时承受父与子呢？再者，那超越普累若麻的受造物如何能容纳那包含整个普累若麻的圣言呢？既然这一切都是不可能且无法证明的，那么唯有教会的宣讲是真实的：祂自己的受造物承受了祂，这受造物赖天主的德能、技能和智慧而存在；它虽以无形的方式为父所承载，却以有形的方式承载了祂的圣言——这就是真实的圣言。\par

2. 因为父同时承载着受造物和祂自己的圣言，而被父承载的圣言按父的意愿将圣神赐予众人。祂赐给一些人如同造物那样所造之物；另一些人则如同义子那样，即从天主而来的，即重生。因此宣告了独一的天主圣父，祂是超越万有、贯乎万有、并在万有之内。圣父确实超越万有，且是基督的头；但圣言是贯乎万有，且是教会的头；而圣神在我们众人之内，祂是活水，\BibleRef{若 7:39}，是主赐给那些正确信祂、爱祂，并知道\BibleQuote{只有一位父，祂是超越万有、贯乎万有、并在我们众人之内} \BibleRef{弗 4:6}的人。主的门徒若望也在福音中为此作证，如此说：\BibleQuote{在起初已有圣言，圣言与天主同在，圣言就是天主。这道在起初与天主同在。万有是藉着祂造的，凡被造的，没有一样不是藉着祂造的。} \BibleRef{若 1:1} 然后他论及圣言本身说：\BibleQuote{祂在世界，世界是藉着祂造的，世界却不认识祂。祂来到自己的地方，自己的人却不接受祂。凡接受祂的，祂赐他们权能成为天主的儿女，就是信祂名的人。} \BibleRef{若 1:10} 再者，若望显示关于祂人性的救恩计划，说：\BibleQuote{圣言成了血肉，居住在我们中间。} \BibleRef{若 1:14} 接着他说：\Quote{我们见过祂的光荣，正如父独生者的光荣，满溢恩宠和真理。} 他如此清楚地向那些愿意听的人，即有耳朵的人指明：只有一位天主，即超越万有的圣父，和一位天主的圣言，祂是贯乎万有的，万有是藉着祂造的；这世界属于祂，并按照父的旨意由祂造，而非由天使；也非由背道、缺陷和无知；也非由普纽尼刻的任何权能——他们中有些人称她为\Quote{母亲}；也非由任何不认识父的造物主。\par

3. 因为世界的创造者确是天主的圣言：这就是我们的主，在末世成了人，存在于这个世界，并以无形的方式包含一切受造物，且内在于整个受造界，因为天主的圣言治理并安排万有；因此祂以有形的方式来到自己的地方，成了血肉，并悬挂在木头上，为将万有总归于自己。\Quote{祂自己的百姓并未接受祂}，正如梅瑟在百姓中宣告此事：\BibleQuote{你的性命必悬在你眼前，你必不相信你的性命。} \BibleRef{申 28:66} 因此，那些不接受祂的，就没有接受生命。\BibleQuote{凡接受祂的，祂赐他们权能成为天主的儿女。} \BibleRef{若 1:12} 因为祂从父获得掌管万有的权能，祂是天主的圣言，也是真人，以理智的方式与无形者相通，并为感官设立可遵守的法律，使万有各守其序；祂显然掌管有形可见的、与人有关的事物；并施行公义和相称的审判于众人；正如达味清楚指明此事说：\Quote{我们的天主必公开降临，决不缄默。} 然后他也显示由祂带来的审判，说：\Quote{烈火在祂面前燃烧，暴风在祂四周旋转。祂呼唤上天与大地，为审判祂的百姓。}\par

\SourceRecord{Translated by Alexander Roberts and William Rambaut.；Ante-Nicene Fathers；Vol. 1.；Edited by Alexander Roberts, James Donaldson, and A. Cleveland Coxe.；Buffalo, NY: Christian Literature Publishing Co.,；1885.；<http://www.newadvent.org/fathers/0103518.htm>.}

% structure: p0001 source=h1 target=section reason=page-title

\section{\WorkTitle{Against Heresies (Book V, Chapter 19)}}

1. 主显然正来到属于祂自己的事物那里，并借着由祂自己所支撑的受造界维系它们，同时藉着[祂悬在]树上的服从，重演那与树相关的悖逆；那曾使已许配给一个男子的童贞厄娃不幸受骗的欺骗，其[效果]也因着天使向（也）许配给一个男子的童贞玛利亚传达真理[之言]而得蒙喜报。因为正如前者因天使的话语而被引入歧途，以致她违抗天主的言语而逃离天主；同样，后者因天使的传报，喜闻她将怀有（\Emph{portaret}）天主，并服从了祂的言语。倘若前者确实违抗了天主，那么后者却被说服服从了天主，为使童贞玛利亚成为童贞厄娃的\OriginalTerm{中保}{advocata}。如此，正如人类因一位童贞而陷入死亡的枷锁，同样也因一位童贞而被解救；童贞的悖逆由童贞的服从在天平另一侧所平衡。因为\OriginalTerm{首先被造的人}{protoplasti}的罪同样由首生者的纠正而得弥补，蛇的来临被鸽子的无害所征服，那曾将我们牢牢束缚于死亡的锁链被解开了。\par

2. 所有异端者既未受过教育，也不认识天主的安排，且不熟悉祂取人性身的那一救恩计划（\Emph{inscii ejus quæ est secundum hominem dispensationis}），因为他们对真理自己瞎了眼，实则是在反对自己的救恩。他们中有些人引入另一位父，而非创造者；另有些人说世界及其质料是由某些天使创造的；还有些人[主张]它被\OriginalTerm{界限}{Horos}与他们所宣称的那位父远远分开——它自行生出（\Emph{floruisse}），并自自身而生。此外，另有一些[异端者]断言，它从缺陷和无知中，在父所包含的事物中获得了实体；还有些人鄙视主可感知的来临，因为他们不承认祂的降生成人；而另一些人忽略[祂应]由童贞女[所生]的安排，主张祂是由若瑟所生。更有甚者，有些人肯定既不是他们的灵魂也不是他们的身体能领受永生，而仅仅是那内在的人。而且，他们主张这[内在的人]就是他们里面的理解力（\Emph{sensum}），并规定只有它上升到\Quote{那完美的}。另有些人，正如我在第一卷中所说的，主张当灵魂得救时，他们的身体并不分享那来自天主的救恩；在那卷书中，我也陈述了所有这些人的假设，并在第二卷中指出他们的软弱和矛盾。\par

\SourceRecord{Translated by Alexander Roberts and William Rambaut.；Ante-Nicene Fathers；Vol. 1.；Edited by Alexander Roberts, James Donaldson, and A. Cleveland Coxe.；Buffalo, NY: Christian Literature Publishing Co.,；1885.；<http://www.newadvent.org/fathers/0103519.htm>.}

% structure: p0001 source=h1 target=section reason=page-title

\section{\WorkTitle{驳斥异端（第五卷，第二十章）}}

1. 如今所有这些（异端者）都比宗徒托付教会之主教晚得多；这一点我在第三卷中已竭尽全力加以证明。由此自然得出，上述异端者既盲目于真理，偏离（正确的）道路，必将行走于各种歧途；因此他们学说的足迹散落各处，毫无一致与联系。但属于教会之人的道路环绕整个世界，因其拥有来自宗徒的可靠传统，并让我们看到众人的信德是同一个且相同的，因为众人都接受同一个天主圣父，相信同一项关于天主圣子降生成人的安排，认识同一圣神的恩赐，遵守同一诫命，保持同一教会体制形式，期待同一位主的来临，等候同一个完整的人（即灵魂与身体）的救恩。毫无疑问，教会的宣讲是真实而稳固的，其中在全世界显示了同一条救恩之路。因为天主的真光托付给了她；因此天主的\Quote{智慧}，藉此她拯救众人，\BibleQuote{在出发时被宣告；在街市上忠信地发出（其声音），在城墙上被宣讲，在城门口不断说话。}\BibleRef{箴 1:20} 因为教会处处宣讲真理，她便是那承载基督真光的七枝灯台。\par

2. 因此，那些离弃教会宣讲的人，质疑圣善长老的知识，却不考虑一个虔诚的人（即便在俗世中）比一个亵渎而无耻的诡辩者重要得多。如今，所有异端者都是如此，还有那些想象自己发现了某种超越真理之事物的人，以致于追随上述那些论调，以多样、不和谐且愚昧的方式前行，对同一事物并非始终持同一见解，正如盲人领盲人，他们必将掉进路上无知的坑中，永寻真理而不得。\BibleRef{弟后 3:7} 因此，我们应当避开他们的学说，小心谨慎，以免受其伤害；而应逃向教会，在她的怀抱中被养育，用主的圣经来滋养。因为教会已如同一个\OriginalTerm{乐园}{paradisus}被栽种在这世界上；故此圣神说：\BibleQuote{你可以随意吃园中所有树上的果子，}\BibleRef{创 2:16} 也就是说，可吃主的所有圣经；但不可用高傲的心去吃，也不可接触任何异端的不和谐。因为这些人自称拥有善恶的知识，并将自己不敬的心置于创造他们的天主之上。因此他们对超出理解范围的事物妄下论断。为此，宗徒也说：\BibleQuote{不要自作聪明，而要明智地谨慎行事，}\BibleRef{罗 12:3} 免得我们因吃这些人的\Quote{知识}（那知道过多之事的知识）而被逐出生命的乐园。主将那些听从祂召唤的人引入这乐园，\BibleQuote{使天上和地上的一切总归于自己，}\BibleRef{弗 1:10} 天上的事物是属神的，地上的事物则构成\OriginalTerm{按照人性而有的安排}{secundum hominem est dispositio}。因此，祂将这些总归于自己：将人结合于圣神，并使圣神住在人内，祂自己成了圣神的头，并将圣神赐给人作头：因为藉着祂（圣神），我们观看、听闻并说话。\par

\SourceRecord{Translated by Alexander Roberts and William Rambaut.；Ante-Nicene Fathers；Vol. 1.；Edited by Alexander Roberts, James Donaldson, and A. Cleveland Coxe.；Buffalo, NY: Christian Literature Publishing Co.,；1885.；<http://www.newadvent.org/fathers/0103520.htm>.}

% structure: p0001 source=h1 target=section reason=page-title

\section{驳斥异端（第五卷，第二十一章）}

1. 因此，祂在恢复万有的工程中，总括了一切，既向我们的仇敌作战，又粉碎了那起初在亚当内俘虏我们的人，并践踏了他的头，正如你在\BibleBook{创世}{纪}中所见天主对蛇说：\BibleQuote{我要把仇恨放在你和女人之间，放在你的后裔和她的后裔之间；祂要\OriginalTerm{窥视}{observabit}你的头，你要窥视祂的脚跟。}\BibleRef{创 3:15}因为从那时起，那要由女人——即由童贞女——所生、相似亚当的那一位，被传报为要窥视蛇的头。这就是宗徒在\BibleBook{迦拉达}{书}中所说的后裔：\BibleQuote{行为法律得以确立，直到那应许的后裔来到。}\BibleRef{迦 3:19}这事实在同一书信中更清楚地显明，他如此说：\BibleQuote{但时期一满，天主就派遣了自己的儿子，生于女人。}\BibleRef{迦 4:4}因为仇敌若不被他所征服的——一个由女人所生的人——所战胜，就不能算是公平地被击败。因为起初他是藉一个女人占了人的上风，自命为人的对手。因此，主自称为人子，在祂自己内包含了那原初的人——那女人是从他而受造的（\OriginalTerm{从他所塑造的女人}{ex quo ea quae secundum mulierem est plasmatio facta est}）——为使我们这一族类如何因一个被征服的人而堕入死亡，也同样如何因一个得胜者而上升至生命；并且，正如死亡藉一个人赢得了针对我们的胜利之棕枝，同样地，我们也藉一个人赢得针对死亡的棕枝。\par

2. 主若来自另一位父，就不会在自己内重新恢复那对蛇的古老而原始的仇恨，以履行\OriginalTerm{造物主}{Demiurgi}的应许，并执行祂的命令。但祂是同一位，既在起初造了我们，又在末期派遣了自己的儿子，主确是由女人所生，执行了祂的命令，既毁灭了我们的仇敌，又使人按照天主的肖像和模样而得以成全。为此，祂没有从别处，而是从法律的话中取得挫败他的方法，并利用父的命令作为帮助，以毁灭和混乱那叛逃的天使。祂像梅瑟和厄里亚一样，禁食四十天后，就饿了。首先，为使我们认识祂是一个真实而实在的人——因为人禁食时感到饥饿是自然的；其次，为使祂的对手有机会攻击祂。因为正如起初仇敌藉食物，虽然不饿，却诱使人违背天主的诫命，同样，最终他没有成功诱使那饥饿的祂取用那出于天主的食物。因为当他试探祂时，他说：\Quote{你若是天主子，就命这些石头变成饼吧！}\BibleRef{玛 4:3}但主以法律的诫命击退了他，说：\BibleQuote{经上记载：人生活不只靠饼。}\BibleRef{申 8:3}至于他所说的\Quote{你若是天主子}，主没有回应；但藉此承认了祂的人性，祂挫败了仇敌，并以祂父的话耗尽了他第一次攻击的力量。因此，发生在乐园里因（原祖父母）二人的吃而带来的人的败坏，藉主在这世上的缺乏食物而被除去了。但他既被法律战败，又试图自己引用一条法律诫命来发起攻击。因为他把祂带到圣殿的最高处，对祂说：\Quote{你若是天主子，就跳下去，因为经上记载：\BibleQuote{他必为你吩咐他的天使，用手托着你，免得你的脚碰在石头上。}}这样，他以圣经的外表掩盖谎言，正如一切异端者所为。因为那句话确实被记载了，即\Quote{他必为你吩咐他的天使}；但\Quote{从这里跳下去}没有任何圣经论及祂：这类劝诱是魔鬼从自己产生的。主因此以法律驳斥了他，说：\BibleQuote{经上又记载：你不可试探上主你的天主。}\BibleRef{申 6:16}指由法律的话指出人之本分，即不该试探天主；并关于祂自己，因祂以人形显现，声明祂不会试探上主、祂的天主。因此，蛇内的理性骄傲被（基督）这人所有的谦卑所消灭，现在魔鬼两次从圣经被征服，因为他被查出建议违反天主诫命的事，并由他的心思显示出是天主的仇敌。然后，他这样被决定性地击败，接着，如集中兵力，调集他所有可用力量以行欺骗，第三次\Quote{把天下万国及其荣华指给祂看}\BibleRef{路 4:6}，按路加记载，说：\Quote{这一切我都要给你——因为已交付给我；我愿给谁就给谁——你若俯伏朝拜我。}主随即揭露他的真面目，说：\Quote{去吧，撒殚！因为经上记载：\BibleQuote{你要朝拜上主你的天主，惟独事奉他。}}\BibleRef{玛 4:10}祂既以此名揭露了他，也（同时）显明了祂自己是谁。因为希伯来语\OriginalTerm{撒殚}{Satan}一词意为叛逃者。如此，第三次战胜他后，祂最终将他作为被法律征服者而驱离；于是，亚当内所发生的违反天主诫命之事，藉人子——那位不违反天主诫命者——所遵守的法律规条而被除去了。\par

3. 那么，基督为之作证的这位主天主是谁？谁也不能试探他，所有人都应朝拜他，且惟独事奉他？毫无疑问，就是那位也赐予法律的天主。因为这些事情已在法律中预言过了，并且主藉着法律的\OriginalTerm{话}{sententiam}指明，法律确实宣告了从圣父而来的天主的圣言；而天主的背命天使被它的话语摧毁，露出真面目，并被遵守天主诫命的人子所击败。因为在起初，他诱使人违犯造物主的法律，从而将人置于自己的权下；但他的权势在于违犯和背命，并且他用这些捆绑了人[归向自己]；所以再一方面，有必要的是，藉着人本身，当他被击败时，应被捆绑在捆绑了人的同一锁链中，为使人获得释放，能回到自己的主那里，将那些曾经束缚自己的锁链——即罪恶——留给他（撒殚）。因为当撒殚被捆绑时，人便获得释放；因为\BibleQuote{除非先将壮士捆住，谁也不能进入壮士的家，抢劫他的财物。}因此，主揭露他说话相反那创造万有的天主的话语，并藉着诫命制伏他。如今法律就是天主的诫命。这人证明他是法律的逃犯和违犯者，也是背弃天主者。在[这人做了这事]之后，圣言将他作为自己的逃犯牢牢捆绑，并夺取了他的财物——即那些被他束缚且被他无理利用的人。确实，他理应被俘虏，因为他曾无理地奴役人；而过去曾被俘的人，依照圣父天主的慈爱被从他掌握者手中救出，圣父怜悯自己的亲手所造，赐给他救恩，藉着圣言——即基督——恢复它，为使人们从实际经验得知，他所接受的不朽性并非出于自己，而是出于天主的白白恩赐。\par

\SourceRecord{Translated by Alexander Roberts and William Rambaut.；Ante-Nicene Fathers；Vol. 1.；Edited by Alexander Roberts, James Donaldson, and A. Cleveland Coxe.；Buffalo, NY: Christian Literature Publishing Co.,；1885.；<http://www.newadvent.org/fathers/0103521.htm>.}

% structure: p0001 source=h1 target=section reason=page-title

\section{反对异端（第五卷，第二十二章）}

1. 如此，主明确显示：法律所揭示的正是真主和唯一的天主；因为法律所宣示为天主的，就是基督所指明为圣父的那一位，唯独基督的门徒应当事奉祂。祂借着法律的话使我们的仇敌彻底混乱；而法律指示我们要赞美天主\OriginalTerm{造物主}{Demiurgum}，并唯独事奉祂。既然如此，我们不应在祂之外或之上寻求另一位圣父，因为只有一位天主，祂使受割损的因信德成义，也使未受割损的因信德成义。\BibleRef{罗 3:30} 因为如果在祂之上还有另一位成全的圣父，那么基督绝不会借着祂的话和诫命推翻撒殚。因为一个无知不能借着另一个无知被除去，一个缺陷也不能借着另一个缺陷被除去。因此，如果法律出于无知和缺陷，那么其中所包含的话怎能消除魔鬼的无知并战胜那壮士呢？因为壮士不能被一个较弱的或同等的人战胜，只能被具有更大能力的人战胜。然而，天主的圣言是超越万有的，祂就是在法律中大声宣告的那一位：\BibleQuote{以色列啊，你要听：上主你的天主是唯一的天主；}以及\BibleQuote{你当全心爱上主你的天主；}以及\BibleQuote{你当朝拜祂，唯独事奉祂。}接着，在福音中，祂借着这些言辞摧毁了背道，既以祂圣父的声音战胜了壮士，也承认法律的诫命表达了祂自己的心意，当祂说：\BibleQuote{你不可试探上主你的天主。}\BibleRef{玛 4:7} 因为祂不是借着别人的话，乃是借着属于祂自己圣父的话使仇敌困惑，如此战胜了壮士。\par

2. 祂藉诫命教导我们：我们既已获得自由，当饥饿时，应取用天主所赐的食物；当我们被提升至可领受的一切恩宠的高位时，不可因信赖正义的行为，或身披卓越服务的神恩而骄傲自大，也不可试探天主，反应在一切事上心怀谦卑，并随时记着这话：\BibleQuote{你不可试探上主你的天主。}\BibleRef{申 6:16} 正如宗徒教导说：\BibleQuote{不可图谋高傲的事，应降就卑微的事；}\BibleRef{罗 12:16} 好叫我们不被财富、世俗的光荣或现世的幻想所迷惑，而应知道我们必须\BibleQuote{朝拜上主你的天主，唯独事奉祂}，并不可听从那虚假应许不属于自己之物的人，他说：\BibleQuote{这一切我都要给你，只要你俯伏朝拜我。}因为他自己承认：朝拜他并遵行他的旨意，就是从天主的荣耀中堕落。那堕落的人还能享受什么愉快或美善的事呢？这样的人除了死亡还能期望或盼望什么？因为死亡紧邻着堕落的人。因此，他也不会赐下他所应许的。因为他怎能施恩给那堕落的人呢？此外，既然天主统治人类和他自己，且没有我们在天之父的旨意，连一只麻雀也不会掉在地上，\BibleRef{玛 10:29} 那么他的宣告——\BibleQuote{因为这一切都交给了我，我愿意给谁就给谁}——显然出自他骄傲自大。因为受造界并不服从他的权能，他自己也不过是受造物之一。他也不会将统治人的权力交给别人；相反，所有其他事物和一切人类事务，都按天主圣父的安排而处理。此外，主宣称\BibleQuote{魔鬼从起初就是撒谎者，真理不在他内。}\BibleRef{若 8:44} 那么，如果他是个撒谎者，真理不在他内，他所说的\BibleQuote{因为这一切都交给我了，我愿意给谁就给谁}\BibleRef{路 4:6} 必定是谎言，而非真理。\par

\SourceRecord{Translated by Alexander Roberts and William Rambaut.；Ante-Nicene Fathers；Vol. 1.；Edited by Alexander Roberts, James Donaldson, and A. Cleveland Coxe.；Buffalo, NY: Christian Literature Publishing Co.,；1885.；<http://www.newadvent.org/fathers/0103522.htm>.}

% structure: p0001 source=h1 target=section reason=page-title

\section{\WorkTitle{驳斥异端（卷五，第二十三章）}}

1. 他早已惯于为迷惑人而说谎反对天主。因为在起初，天主赐给人多种食物，却只命令他不可吃一棵树上的果子，正如圣经告诉我们天主对亚当说：\Quote{园中各种树上的果子，你都可吃；只有知善恶树上的果子，你不可吃，因为那一天你吃了，必定要死。}\BibleRef{创 2:16} 那时他说谎反对主，引诱人，正如圣经说蛇对女人说：\Quote{天主真说了，你们不可吃园中任何树上的果子吗？}\BibleRef{创 3:1} 当她揭露了谎言，并单纯地重述了命令，如他所说的：\Quote{园中树上的果子，我们都可以吃；只有园中央那棵树上的果子，天主说过，你们不可吃，也不可摸，免得死亡。}\BibleRef{创 3:2} 他从女人那里听说了天主的命令后，便施展诡计，最终用谎言欺骗了她，说：\Quote{你们决不会死！因为天主知道，你们吃了那果子，眼睛就会开了，将如同天主一样知道善恶。}\BibleRef{创 3:4} 首先，在天主的园子里，他辩论关于天主的事，好像天主不在那里，因为他不知道天主的伟大；其次，他从女人那里得知天主说过，如果他们尝了上述的树，就会死，他便张口说出第三个谎言：\Quote{你们决不会死。} 但结果证明了天主是真实的，蛇是说谎者：死亡临到了那些吃了的人身上。因为他们与果子一同落入了死亡的权势，因为他们因不顺从而吃；不顺从天主带来死亡。因此，他们成为死亡的负债者，从那一刻起，他们就被交给了死亡。\par

2. 这样，在他们吃的那一天，他们实在是死了，成为了死亡的负债者，因为那一天是创造的一日。经上说：\Quote{过了晚上，过了早晨，这是第一天。} 就在他们吃的这一天，他们也死了。但按照日子的循环和进程——其中一个称为第一，另一个第二，又一个第三——如果有人勤勉地探究亚当是在七天中的哪一天死的，他通过考查主的救世计划就会发现。因为主在祂自己内从起初到终结总结了全人类，也总结了死亡。由此可知，主为服从祂的圣父而受死，正是在亚当因不服从而死的那一天。亚当是在他吃的那一天死的，因为天主说：\Quote{你吃的那一天，必定要死。} 因此，主在自己内重演这一天，在安息日的前一日，即创造的第六天受难，这也是人被创造的日子；祂藉着苦难为人再造，即从死亡中再造。还有一些人把亚当的死推迟到第一千年，因为\Quote{主看一日如千年}\BibleRef{伯后 3:8}，他没有越过一千年，而是在其中死了，如此证实了对其罪恶的判决。因此，无论关于不顺从（即死亡），即因不顺从而被交于死亡、成为死亡的负债者；无论关于他们在同一天吃、也在同一天死（因为那是创造的一日）；无论关于按此日子的循环，他们在吃的那一天死，即预备日，称为\Quote{纯洁的晚餐}，即节期的第六天，主也在那天受难显示了这一点；还是关于他（亚当）没有越过一千年，而是在其限度内死了——从所有这些含义来看，天主确实是真实的。因为尝了那树的人都死了；蛇被证明是撒谎者和杀人者，正如主论及它说：\Quote{他从起初就是杀人者，不站在真理上，因为他心里没有真理。}\BibleRef{若 8:44}\par

\SourceRecord{Translated by Alexander Roberts and William Rambaut.；Ante-Nicene Fathers；Vol. 1.；Edited by Alexander Roberts, James Donaldson, and A. Cleveland Coxe.；Buffalo, NY: Christian Literature Publishing Co.,；1885.；<http://www.newadvent.org/fathers/0103523.htm>.}

% structure: p0001 source=h1 target=section reason=page-title

\section{驳斥异端（第五卷，第二十四章）}

1. 因此，正如魔鬼起初说谎，它也在末时说谎，当它说：\Quote{这一切都交付给了我，我愿意给谁，就给谁。}\BibleRef{玛 4:9} 因为不是它，而是天主设立了这世上的国度；因为\Quote{君王的心在天主手中。}\BibleRef{箴 21:1} 圣言也藉著\Person{撒罗满}说：\Quote{藉着我，君王执政，首领秉公行义。藉着我，权贵被兴起，藉着我，君王统治大地。}\BibleRef{箴 8:15} 宗徒\Person{保禄}也就同一主题说：\Quote{要服从所有更高的权力，因为没有权力不是出于天主的：现今存在的，都是由天主所指定的。}\BibleRef{罗 13:1} 再次，论到他们，他说：\Quote{因为他并非徒然佩剑；他是天主的仆役，是为向作恶的人发怒施以报复的。}\BibleRef{罗 13:4} 他说这些话，不是指天使的权力，也不是指不可见的统治者——如有些人胆敢解释这段经文——而是指实际的人类权威，[他显示这一点]当他说：\Quote{为此你们也该纳税，因为他们是天主的仆役，专为这事服役。}\BibleRef{罗 13:6} 主也证实了这一点，当他没有做魔鬼引诱他做的事，反而吩咐为他自己和\Person{伯多禄}向税吏纳税；\BibleRef{玛 17:27} 因为\Quote{他们是天主的仆役，专为这事服役。}\par

2. 因为人既远离了天主，竟至狂怒到视自己的兄弟为仇敌，并无所畏惧地陷入各种不安的行为、谋杀和贪婪；天主就将对人的恐惧加于人类，因为他们不承认对天主的敬畏，好使他们受制于人的权威，并被他们的法律所约束，从而能达到某种程度的正义，并因悬在眼前的刀剑的畏惧而彼此容忍，正如宗徒所说：\Quote{因为他并非徒然佩剑；他是天主的仆役，是为向作恶的人发怒施以报复的。}为此，官吏们本身，当他们在公正和合法的方式下行事时，他们拥有法律作为正义的衣袍，他们的行为就不会受到质问，也不会受罚。但若他们为了颠覆正义，而不义地、不敬地、非法地、暴虐地行事，在这些事上他们也要灭亡；因为天主公义的审判同样临到众人，绝无欠缺。因此，地上的统治是由天主为万民的利益而设立的，并非由那永不安宁、且不愿看见万民安静行事的魔鬼所设，好使人在对人的统治的畏惧之下，不致像鱼一样互相吞噬；而是藉着法律的建立，遏制万民中过度的邪恶。从这个观点来看，向我们征税的人乃是\Quote{天主的仆役，专为这事服役}。\par

3. 所以，既然\Quote{在上的权力都是天主所指定的}，那么魔鬼说\Quote{这一切都交付给了我，我愿意给谁，就给谁}，显然是谎言。因为藉着那召叫人类存在的同一位天主的法律，君王也被指定，适应于当时处于他们治理之下的人。这些统治者中，有些是为纠正和造福其臣民、并为维护正义而赐予的；另一些则为了恐惧、惩罚和责备的目的；还有些，照[臣民]所应得的，是为了欺骗、耻辱和骄傲；然而天主公义的审判，如我先前所说，同样临到众人。魔鬼，作为叛逃的天使，只能做到这种地步，正如起初所做的，[即]欺骗和引诱人的心智，使之不服从天主的诫命，并逐渐使那些试图服事他的人的心昏暗，以至于忘记真天主，反而崇拜他自己为神。\par

4. 正如有人，身为叛徒，以敌对的方式夺取他人的领土，骚扰其居民，以便在对他的叛逃和抢夺一无所知的人中，为自己窃取君王的荣耀；同样，魔鬼，作为那些被置于空中之神的天使之一，正如宗徒\Person{保禄}在致\Place{厄弗所}人书信中所宣告的，\BibleRef{弗 2:2} 因嫉妒人而成为神圣法律的叛逃者：因为嫉妒是与天主不相干的事。而由于他的叛逃被人揭露，人成了探究他思想的[工具]（\OriginalTerm{人成了他思想的审查}{et examinatio sententiæ ejus, homo factus est}），他就越发坚决地敌对众人，嫉妒人的生命，并希望把人卷入他自己叛逃的权势中。然而，天主的圣言，万物的创造者，藉着人性战胜了他，并显明他是叛逃者，反而将他置于人的权下。因为他说：\Quote{看，我赐予你们践踏蛇和蝎子，以及敌人一切权柄的能力，}\BibleRef{路 10:19} 好使他如何藉着叛逃获得对人的统治，同样，他的叛逃也如何藉着人转回天主而被剥夺力量。\par

\SourceRecord{Translated by Alexander Roberts and William Rambaut.；Ante-Nicene Fathers；Vol. 1.；Edited by Alexander Roberts, James Donaldson, and A. Cleveland Coxe.；Buffalo, NY: Christian Literature Publishing Co.,；1885.；<http://www.newadvent.org/fathers/0103524.htm>.}

% structure: p0001 source=h1 target=section reason=page-title

\section{《驳斥异端》（第五卷，第二十五章）}

1. 不仅通过已提及的细节，而且通过将在假基督时期发生的事件，也表明他作为背教者和强盗，渴望被当作天主来敬拜；尽管他只是一个奴隶，却希望自己被宣布为王。因为他（假基督）具有魔鬼的全部能力，将来到，不是作为一个正义的君王，也不是作为一个合法的君王（即顺从天主的君王），而是作为一个不敬虔、不义、无法无天的君王；作为一个背教者、邪恶和凶杀的；作为一个强盗，在他身上集中了撒殚的一切背教行为，并废弃偶像，以劝服人他自己就是天主，把自己提升为唯一的偶像，在他身上拥有其他偶像的多种错误。他这样做，是为了使那些如今通过许多可憎之事敬拜魔鬼的人，可以通过这一个偶像来事奉他自己。关于此人，宗徒在\BibleBook{得撒洛尼}{后书}中这样说：\BibleQuote{除非先有背叛的事，那罪恶之人，即丧亡之子，将被显示出来，他反对并高举自己于一切称为神或受人崇拜者之上，以致他坐在天主的殿中，显示自己好像是天主。} 因此，宗徒清楚地指出他的背教，以及他被高举于一切称为神或受人崇拜者之上——即高于每一个偶像——因为这些固然被人称为神，但并不是真正的神；并且他将以暴虐的方式试图把自己展现为天主。\par

2. 此外，宗徒也指出了这一点，即我已经以多种方式表明的，耶路撒冷的圣殿是按照真天主的吩咐建造的。因为宗徒亲自以他自己的口吻清楚地称它为天主的殿。我在第三卷中已经表明，当宗徒们以自己身份说话时，没有人被称为神，除非那真正是神的那一位，即我们主的父，耶路撒冷的圣殿是按照祂的吩咐而建造的，为的是我已经提到过的那些目的；敌人将坐在那殿中，企图显示自己是基督，正如主也宣布的：\BibleQuote{但当你们看见招致荒凉的可憎之物，即由达内尔先知所说过的，站在圣所（诵读的应明白）时，那时在犹太的，要逃到山上；在屋顶上的，不要下来取家中的东西；因为那时将有大灾难，是从世界开始直到如今从未有过的，将来也绝不会再有。}\par

3. 达内尔也前瞻最后一个王国的结局，即最后的十个国王，这些人的王国将被分割，而丧亡之子将降临在他们身上。他宣称从兽身上将长出十角，另有一小角将从中升起，它将从它面前拔除三个先前的角。他说：\BibleQuote{我观看，见这角有人眼，有一张说大话的嘴，它的容貌比它的同伴更傲慢。我观看，见这角与圣者作战，战胜了他们，直到万古常存者来到，将审判授予至高天主的圣者；时候到了，圣者就取得了王国。} \BibleRef{达 7:8} 然后，在异象的解释中，又对他说：\BibleQuote{第四只兽将是地上的第四个王国，它将胜过所有其他的王国，并吞食全地，把它践踏碎。它的十角是十位将要兴起的国王；在他们之后，将兴起另一位，他将在恶行上超过所有以前的人，并推翻三位国王；他将说亵渎至高天主的话，磨灭至高天主的圣者，并企图改变节期和法律；一切将被交在他手中，直到一个时期、两个时期和半个时期，} \BibleRef{达 7:23} 即三年六个月。在这期间，当他来时，他将在全地掌权。关于他，保禄宗徒在\BibleBook{得撒洛尼}{后书}中又说道，同时宣告他来临的原因：\BibleQuote{那时，那个邪恶者将被显露出来，主耶稣要以口中的气息杀死他，并以祂来临的显现消灭他。那邪恶者的来临是基于撒殚的运作，具有各种能力、征兆和虚假的奇迹，以及各种不义的欺骗，为那些丧亡的人；因为他们没有接受爱真理的心，以致得救。为此，天主将给他们一种错误的作用，叫他们相信谎言；使一切不信真理、反而喜爱不义的人都被定罪。} \BibleRef{得后 2:8}\par

4. 主也这样对那不信祂的人说：\Quote{我因我父的名而来，你们却不接待我；若另有一位凭自己的名而来，你们倒要接待他。}\BibleRef{若 5:43}此处称假基督为\Quote{那另一位}，因为他与主隔绝了。这也正是主所提及的不义的审判官，\Quote{不敬畏天主，也不尊重人，}\BibleRef{路 18:2}那寡妇因忘记天主——即属地的耶路撒冷——而投奔于他，求他伸冤。在假基督的王国时期，他也要这样做：他将他的王国迁往那城，并坐在天主的殿中，迷惑那些崇拜他的人，仿佛他就是基督。为此，\Person{达尼尔}又说道：\Quote{他要亵渎圣所；罪被当作祭品献上，正义被抛弃在地，他行了事（\Emph{fecit}），并且顺利成就。}\BibleRef{达 8:12}天使\Person{加俾额尔}在解释这异象时论及此人说：\Quote{到了他们国度的末期，将兴起一位面貌凶恶的君王，明白隐谜，极有力量，充满奇事；他要败坏、引导、影响（\Emph{faciet}），使强人屈服，同样也对付圣民；他的轭如同花环般套在颈上；他手中有欺诈，心高气傲；他要用诡诈毁灭多人，领多人丧亡，将他们如同鸡蛋般捏碎。}\BibleRef{达 8:23}接着，他指出他暴政持续的时间，在此时期圣徒们将被迫逃亡，那些向天主献上纯洁祭品的人：\Quote{在一周之半，}他说，\Quote{祭品和奠祭将被废除，那招致荒凉的可憎之物被安置在圣殿中，直到最后终结，荒凉才告结束。}\BibleRef{达 9:27}三年半正好是半周。\par

5. 从所有这些段落中，不仅启示了我们关于背道以及那聚集一切撒殚错误者的具体细节，而且也启示了：同一天主圣父，曾由先知们所预言，却由基督所彰显。因为如果\Person{达尼尔}所预言的终结已被主所确认，当祂说：\Quote{几时你们看见那招致荒凉的可憎之物，即由\Person{达尼尔}先知所说的，}\BibleRef{玛 24:15}（并且天使\Person{加俾额尔}曾向\Person{达尼尔}解释异象，他是创造者的总领天使，也曾向玛利亚宣告基督可见的来临与降生成人），那么，同一天主就被最清楚地指明：祂派遣了先知，预许了圣子，并召叫我们认识祂。\par

\SourceRecord{Translated by Alexander Roberts and William Rambaut.；Ante-Nicene Fathers；Vol. 1.；Edited by Alexander Roberts, James Donaldson, and A. Cleveland Coxe.；Buffalo, NY: Christian Literature Publishing Co.,；1885.；<http://www.newadvent.org/fathers/0103525.htm>.}

% structure: p0001 source=h1 target=section reason=page-title

\section{\WorkTitle{驳斥异端（第五卷，第二十六章）}}

1. 若望在《默示录》中更清楚地为主的门徒指示了末世将发生的事，以及那时将兴起的十位君王，现今统治（世界）的帝国将在这十位君王中被分割。他教导我们，达尼尔所见的十角是什么，告诉我们曾如此对他说：\Quote{你看见的那十角，就是十位君王，他们还没有得到王国，但将与那兽一同得到权柄，如同君王一小时之久。他们同心合意，将自己的能力和权柄交给那兽。他们要与羔羊交战，羔羊将战胜他们，因为祂是万主之主，万王之王。}\BibleRef{默 17:12} 由此可见，在这些（权柄）中，那将要来的（君王）将杀死三个，将余下的置于自己的权下，而他本人将是他们中的第八位。他们将使巴比伦变为荒场，用火焚烧她，并将他们的王国交给那兽，使教会逃散。此后，他们将因我主的来临而被消灭。因为王国必须分裂，并因此灭亡，主（在祂）说时（宣告了这一点）：\Quote{凡一国自相纷争，必成荒场；凡一城或一家自相纷争，必不得站立。}\BibleRef{玛 12:25} 因此，王国、城和家必须分裂为十；为此祂已预先暗示了（将要发生的）分裂和分割。达尼尔也特别说到，第四王国的终结在于拿步高所见像的脚趾，有非人手所凿的石头打在其上；正如他自己所说：\Quote{脚，一部分是铁，一部分是泥，直到非人手凿出的石头击打那像的铁泥脚，把它们砸得粉碎，直至尽头。}\BibleRef{达 2:33} 随后在解释这一点时，他说：\Quote{你看见那脚和脚趾，一部分是窑匠的泥，一部分是铁，那王国必要分裂，其中却有铁的坚硬，正如你看见铁与窑泥混合。脚趾，一部分是铁，一部分是泥。}\BibleRef{达 2:41} 因此，这十个脚趾就是这十位君王，王国将在他们中间被分割，其中一些将是强壮而活跃或有力的；另一些则将是迟钝无用的，并且不会和睦；正如达尼尔也说：\Quote{那王国一部分强，一部分脆弱。你看见铁与窑泥混合，他们必有人种的混杂，却不能彼此相合，正如铁不能与陶器相熔合一样。}\BibleRef{达 2:42} 既然有一个终结，他说：\Quote{在这些君王的日子，天上的天主必要兴起一个永不灭亡的王国，祂的王国必不留给其他民族。它要粉碎并消灭一切王国，而它自己将永远屹立。正如你看见那石头非人手从山上凿出，打碎了窑泥、铁、铜、银和金，天主已将此后必成的事指示给君王；这梦是真实的，这解释是可信的。}\BibleRef{达 2:44}\par

2. 因此，如果伟大的天主藉达尼尔启示了未来的事，并藉祂的圣子予以确认；如果基督就是那非人手凿出的石头，祂将要毁灭暂时的王国，并引入一个永恒的王国，即义人的复活；正如祂所宣称的：\Quote{天上的天主必要兴起一个永不灭亡的王国。}——那么，那些拒绝\OriginalTerm{创造者}{Demiurgum}的人，不承认先知们是从同一位父——主也是从祂而来——被预先派遣的，反而声称预言源于不同的权能，这样的人既已被驳斥，就应当醒悟过来。因为创造者藉所有先知所预言的一切，基督在末时都成就了，服膺于祂父的旨意，完成了祂对人类所安排的救恩计划。因此，那些亵渎创造者的人，无论是像玛尔基翁的门徒那样以公开的言语，还是像华伦提诺的追随者和所有所谓诺斯替派那样以曲解（圣经）意义的方式，都应当被所有敬拜天主的人认作撒殚的代理人；撒殚正是藉着他们，现在（而不是以前）被看见反抗天主，甚至反抗那为各种背道准备了永火的天主。因为他自己不敢公开地亵渎他的主；正如起初他藉着蛇的媒介引诱人，仿佛在天主面前隐藏自己。犹斯定确实说过：在主显现之前，撒殚从不敢亵渎天主，因为他尚不知道自己的判决，因它隐藏在比喻和寓言中；但在主显现之后，当他从基督和祂宗徒的话中清楚得知，永火已为他预备，因为他自愿背离天主，也同样为所有不悔改而继续背道的人预备了，他现在就藉着这样的人亵渎那带来审判的主（因他已被定罪），并将他背道的罪责归于他的创造主，而非他自己自由意志的选择。正如那些违法者，当受到惩罚时，他们归咎于立法者，而不归咎于自己。同样，这些充满撒殚精神的人，对我们的创造主——祂不仅赐予我们生命之神，还建立了适用于万民的法律——提出无数的指控，却不肯承认天主的审判是公义的。因此，他们还开始设想另一位父，祂既不关心也不治理我们的事务，甚至赞同一切罪恶。\par

\SourceRecord{Translated by Alexander Roberts and William Rambaut.；Ante-Nicene Fathers；Vol. 1.；Edited by Alexander Roberts, James Donaldson, and A. Cleveland Coxe.；Buffalo, NY: Christian Literature Publishing Co.,；1885.；<http://www.newadvent.org/fathers/0103526.htm>.}

% structure: p0001 source=h1 target=section reason=page-title

\section{\WorkTitle{驳斥异端}（第五卷，第二十七章）}

1. 如果圣父不施行审判，那么审判就不属于祂，或者祂同意所发生的一切行为；如果祂不审判，那么所有人就都是平等的，并被视为处于相同状况。基督的来临就没有了目的，甚至荒谬，因为祂不行使审判权。因为\BibleQuote{祂来是要使儿子与父亲分离，使女儿与母亲分离，使儿媳与婆婆分离；}\BibleRef{玛 10:25} \BibleQuote{当两个人同在一张床上，要取一个而留下另一个；两个女人一同推磨，要取一个而留下另一个；}\BibleRef{路 17:34} \BibleQuote{在末期，要命令收割的人先将莠子收集起来，捆成捆，用不灭的火焚烧，却将麦子收入仓内；}\BibleRef{玛 13:30} \BibleQuote{并将绵羊领进为他们预备的王国，却将山羊投入为魔鬼及其使者所预备的永火里。}\BibleRef{玛 25:33} 这是为什么？圣言来是为了使许多人跌倒和复活吗？当然，是为了使那些不信祂的人跌倒，祂曾警告他们在审判之日将遭受比索多玛和哈摩辣更重的处罚；\BibleRef{路 10:12} 但为了信徒和遵行祂在天之父旨意的人的复活。如果圣子的来临确实同样临到所有人，但目的是为了审判，并将信者与不信者分开，既然信者自愿遵行祂的旨意，而不信者也自愿不顺从祂的教导，那么显然，祂的圣父使所有人处于相同状况，每个人都有自己的选择和自由理解；祂眷顾一切，施行照料一切，\BibleQuote{使太阳上升，光照恶人也光照善人；降雨给义人也给不义的人。}\BibleRef{玛 5:45}\par

2. 凡继续留在对天主的爱中的人，祂就赐予他们与祂的共融。但与天主的共融就是生命和光，以及享受祂所储备的一切恩惠。但那些按自己的选择离开天主的人，祂便施加他们自己选择的分离。与天主分离就是死亡，与光分离就是黑暗；与天主分离就在于丧失祂所储备的一切恩惠。因此，那些因背教而抛弃上述事物的人，实际上丧失了一切善，便经历各种惩罚。然而，天主并非直接亲自惩罚他们，而是该惩罚落到他们身上，因为他们丧失了一切善。既然善的事物在天主那里是永恒无终的，故此丧失这些善也是永恒无尽的。这正如大量光芒的情况：那些弄瞎自己眼睛或被别人弄瞎的人，永远被剥夺了享受光的乐趣。不是光照施加了瞎眼的惩罚，而是瞎眼本身给他们带来了灾祸：因此主宣称：\BibleQuote{信我的人不受审判}\BibleRef{若 3:18}，也就是说，不与天主分离，因为他藉着信德与天主结合。另一方面，祂说：\BibleQuote{不信的人已受了审判，因为他不信天主独生子的名字；}也就是说，他自动与天主分离。\BibleQuote{审判就在于此：光明来到世界上，世人却爱黑暗甚于光明。因为凡作恶的恨光明，也不来就光明，怕他的行为被显露；但履行真理的却来就光明，为显示他的行为是在天主内完成的。}\par

\SourceRecord{Translated by Alexander Roberts and William Rambaut.；Ante-Nicene Fathers；Vol. 1.；Edited by Alexander Roberts, James Donaldson, and A. Cleveland Coxe.；Buffalo, NY: Christian Literature Publishing Co.,；1885.；<http://www.newadvent.org/fathers/0103527.htm>.}

% structure: p0001 source=h1 target=section reason=page-title

\section{\WorkTitle{驳异端}（第五卷，第二十八章）}

1. 既然如此，在此世界上（\OriginalTerm{世界}{\GreekText{αἰῶνι}}）有些人投奔光明，并因信德与天主结合，但另一些人躲避光明，使自己与天主分离，天主的圣言便来为两者预备恰当的居所。因为那些在光明中的人，为使他们享受光明及其中所含的美善；而那些在黑暗中的人，为使他们分担其中的灾祸。因此主说，右边的人被召进天国，但左边的人祂要投入永火，因为他们已自绝于一切美善。\par

2. 为此，宗徒说：\BibleQuote{因为他们没有接受爱天主的心，为能得救，所以天主也给他们一种错误的能力，叫他们相信谎言，使一切不信真理而喜爱不义的人都被判断。}\BibleRef{得后 2:10} 因为当那\OriginalTerm{假基督}{Antichrist}来临时，他自行将背教集中于一身，按照自己的意愿和选择成就一切，并坐在天主的圣殿里，使他的受骗者敬拜他为基督；因此他也应当被投入火湖中。\BibleQuote{被投入火湖中}\BibleRef{默 19:20}（这是按天主的安排），天主凭其预知预见这一切，并在适当的时候差遣这样的人，\BibleQuote{使他们相信谎言，使一切不信真理而喜爱不义的人都被判断。}。若望在\WorkTitle{默示录}中如此描述他的来临：\BibleQuote{我所看见的兽，相似豹子，脚像熊的脚，口像狮子的口；那龙把自己的能力、宝座和大权柄都交给了它。它的头中，有一个似乎受了致命伤，但那死伤却治好了；全地的人都惊奇，跟随那兽。他们敬拜那龙，因为龙把权柄给了那兽；也敬拜那兽，说：\InnerQuote{谁能比得上这兽？谁能与它交战？}又赐给它一个口，说大话和亵渎的话，赐给它权柄任凭四十二个月。它开口向天主说亵渎的话，亵渎他的名和他的帐幕，以及那些住在天上的。又赐给它权柄，管辖各族、各民、各方、各国。凡住在地上的人，名字从创世以来没有被宰杀的羔羊的生命册上记着的，都要拜它。凡有耳的，就应当听。如果人掳掠人，他也必被掳掠；如果用刀杀人，他也必被刀杀。圣徒的恒心和信德就在于此。}\BibleRef{默 13:2} 此后，他又描述他的副手，也称他为假先知：\BibleQuote{它说话好像龙，在头一个兽面前施行它所有的权柄，叫地和住在地上的人都拜那死伤医治好的头一个兽。又行大奇事，甚至在人面前叫火从天降到地上。它迷惑住在地上的人。}\BibleRef{默 13:11} 不要以为这些奇事是出于天主的能力，而是出于巫术的作用。我们也不必惊讶，既然魔鬼和背教的神体为他服务，他藉着他们行奇事，迷惑住在地上的人。若望还说：\BibleQuote{它吩咐给兽造一个像，又给那像有气息，叫那像能说话，并使那些不肯拜兽像的人被杀害。}又说：\BibleQuote{它又叫人，无论大小贫富，都在右手上或额上受一个印记；除了那有兽的印记或兽名或兽名数字的人，谁也不能买或卖。那数字是六百六十六。}\BibleRef{默 13:14} 即六个一百、六个十、六个一。这是总结整个六千年来的背教。\par

3. 因为世界创造了多少日，也要在多少千年中结束。为此，圣经说：\BibleQuote{天地和它们的一切点缀都完成了。到第六天，天主完成了祂所造的工作；第七天，天主停止了祂所做的一切工作，安息了。}\BibleRef{创 2:2} 这是对先前受造之物的记述，也是对将来的预言。因为主的日子如一千年；\BibleRef{伯后 3:8} 在六天内，创造之物得以完成；因此，显然它们将在第六千年终结。\par

4. 因此，在整个时间内，从起初由天主的手——即圣子和圣神——所塑造的人，是按天主的肖像和模样造的；那糠秕，即背教，被抛弃；而麦子，即那些因信德向天主结出果实的人，被收入仓库。为此，为那些得救的人，患难是必需的，好使他们经过某种方式的粉碎、筛细，由天主圣言的坚忍所浇灌，并经火炼（为得净化），被配得王的筵席。正如我们中间有一位，因他为天主作证而被判给野兽时所说：\Quote{我是基督的麦子，被野兽的牙齿磨碎，好成为天主纯洁的饼。}\par

\SourceRecord{Translated by Alexander Roberts and William Rambaut.；Ante-Nicene Fathers；Vol. 1.；Edited by Alexander Roberts, James Donaldson, and A. Cleveland Coxe.；Buffalo, NY: Christian Literature Publishing Co.,；1885.；<http://www.newadvent.org/fathers/0103528.htm>.}

% structure: p0001 source=h1 target=section reason=page-title

\section{\WorkTitle{反异端（第五卷，第二十九章）}}

1. 在前几卷中，我已陈述了天主允许这一切受造的原因，并指出所有这一切都是为那得救的人性益处而创造的，使那拥有自由意志和自身能力的（人性）成熟以得永生，并预备它，使之更适于永恒地服从天主。因此，受造界是适合（人的需要）的；因为人不是为受造界而造，而是受造界为人而造。然而，那些不自行举目望天、不感谢造物主、不愿看见真理之光，而像藏在愚昧深处的瞎眼老鼠的民族，圣言正当地将他们算着\BibleQuote{如同从槽里流出的废水，如同天秤上的转砝——实为虚无}\BibleRef{依 40:15}；他们对于义人而言，有用和有益之处，就如秸秆有助于麦子生长，其稻草通过焚烧有助于熔炼金子。因此，当末后教会忽然被从此世提升时，经上说：\BibleQuote{那时必有大灾难，是从宇宙开始直到如今从未有过的，也决不会再有。}\BibleRef{玛 24:21} 因为这是义人最后的争战，他们得胜时，将戴上不朽的冠冕。\par

2. 因此，在这兽来临时，它身上综合了所有的不义和一切欺骗，为要使所有背教的力量都流入并封闭在它内，从而投入火炉。所以，它的名字理所当然地拥有六百六十六这个数目，因为它本身综合了诺厄时代以前因天使背教所发生的一切邪恶混合。因为当洪水来到地上时，诺厄已是六百岁，那洪水为那在诺厄时代的可憎世代而扫荡了背叛的世界。而且（敌基督）也综合了洪水以后所发明的一切偶像错误，连同杀害先知和铲除义人。因为拿步高所立的像，高度为六十肘，宽度为六肘；为此，阿纳尼雅、阿匝黎雅和米沙耳因不敬拜那像而被投入火炉，这藉着他们所遭遇的事，预言性地指出了末世将临到义人身上的愤怒。因为那像整体而言，是这人之来临的预表，宣告他无疑应独自被所有人崇拜。因此，诺厄的六百年（洪水因背教而发生在诺厄时代）以及那像的肘数（为此这些义人被投入火炉），都指示了那人之名字的数目，在他身上集中了全部六千年来的背教、不义、邪恶、假预言和欺骗；为此，也将有火焰的洪流（降在地上）。\par

\SourceRecord{Translated by Alexander Roberts and William Rambaut.；Ante-Nicene Fathers；Vol. 1.；Edited by Alexander Roberts, James Donaldson, and A. Cleveland Coxe.；Buffalo, NY: Christian Literature Publishing Co.,；1885.；<http://www.newadvent.org/fathers/0103529.htm>.}

% structure: p0001 source=h1 target=section reason=page-title

\section{\WorkTitle{驳斥异端（第五卷，第三十章）}}

1. 既如此，这个数字出现在所有最权威和最古老的\WorkTitle{默示录}抄本中，并且那些曾亲眼见过\Person{若望}的人也为此（抄本）作证；同时，理性也引导我们得出结论，那兽的名字的数字，（若）按照希腊计算方法，依据其包含的字母（数值），应等于六百六十六；也就是说，十位数应与百位数相等，百位数应与个位数相等（因为那个始终体现数字六的数值，表明了那背教的全部重演，它发生于起初、中间时期，并将在终结时发生）——我不知道为何有些人竟随从通常的说法而犯错，歪曲了名字中的中间数字，从中减去五十，以至于他们所得的不是六个十，而是一个十。[我倾向于认为这是由于抄写者的失误造成的，这是常有的事，因为数字也用字母表示；因此，表示数字六十的希腊字母容易被扩展成希腊字母伊奥塔（\GreekText{Ι}）。]其他人于是不加查考地接受了这个读法；有些人出于天真，并凭自己的主见，使用了这个表示一个十的数字；而另一些人，因缺乏经验，竟冒险去寻找一个含有那错误和虚假数字的名字。至于那些出于天真并无恶意而这样做的人，我们可以假定天主会宽恕他们。但是那些为了虚荣而断然主张含有那虚假数字的名字应当被接受，并断言他们自己猜中的这个名字就是将要来的那一位的名字的人，这样的人必然不会没有损失，因为他们既使自己陷入了错误，也使信赖他们的人受到了误导。首先，偏离真理并误以为非为是，本身就是一种损失；其次，凡对圣经有所增减的人，必受不轻的惩罚，\BibleRef{默 22:19}这样的人必然要落入那种惩罚之下。此外，另一个绝非微不足道的危险将降临到那些妄自以为自己知道假基督名字的人身上。因为如果这些人假设了一个（数字），而当这位（假基督）带着另一个数字来临时，他们就会轻易被他诱骗，以为他不是那位必须防备的期待者。\par

2. 因此，这些人应当了解（真正的情况），并回到那名字的真实数字上来，以免被列入假先知之中。但是，既然他们知道了圣经所宣告的确切数字，即六百六十六，就让他们首先等候那王国分裂成十个；其次，当这些君王在位，开始整顿事务、推进其王国时，（让他们学会）承认那位前来为自己夺取王国、并恐吓我们之前所说的人的，他的名字包含上述数字，那一位就是真正致人荒凉的可憎之物。宗徒也同样断定：\Quote{当人们说\InnerQuote{平安无事}时，毁灭会突然临到他们。} \BibleRef{得前 5:3} \Person{耶肋米亚}不仅指出他的突然来临，甚至还指明他将出自哪个支派，他说：\Quote{我们必从丹听到他快马的声音；他骏马嘶鸣的声音震动全地；他必来吞灭大地和其中的所有，也吞灭城邑和其中的居民。} \BibleRef{耶 8:16} 这也是为什么在\WorkTitle{默示录}中，这个支派没有被列入得救的支派之列的原因。\par

3. 因此，等候预言的应验比猜测和搜寻可能出现的名字更为确定且较少风险，因为可以找到许多名字含有所述的数字；而同样的问题终将无法解决。因为如果找到许多名字含有这个数字，就会问将要来的那人该用哪一个。我说这话并非因为缺少包含那名字数字的名字，而是出于对天主的敬畏和对真理的热忱：因为名字\Emph{\OriginalTerm{埃凡塔斯}{\GreekText{ΕΥΑΝΘΑΣ}}}含有所需的数字，但我对此不作论断。此外，\Emph{\OriginalTerm{拉泰诺斯}{\GreekText{ΛΑΤΕΙΝΟΣ}}}也有六百六十六这个数字；而且这很可能就是答案，因为这是那最后王国（但以理所见的四个王国之一）的名字。因为拉丁人就是现今统治的人；但我不会为这个巧合自夸。\Emph{\OriginalTerm{泰坦}{\GreekText{ΤΕΙΤΑΝ}}}（第一个音节由两个希腊元音\OriginalTerm{\GreekText{ε}}{\GreekText{ε}}和\OriginalTerm{\GreekText{ι}}{\GreekText{ι}}写成）在我们所找到的所有名字中，也颇有可信之处。因为它本身包含预言中的数字，由六个字母组成，每个音节三个字母；而且这个词本身是古老的，不常用于世俗；因为在我们的君王中，没有一位名叫泰坦，无论希腊人还是外邦人所公开敬拜的偶像中也没有这个称号。在许多人心目中，这个名字被视为神圣，以至于现今掌权的人甚至称太阳为\Quote{泰坦}。这个词也带有某种外在的报应和施加应得惩罚的意味，因为他（假基督）假装为受压迫者伸冤。此外，这是一个古老的名字，可信，具有王室尊严，更进一步说，这是一个暴君的名字。既然\Quote{泰坦}这个名字有如此多可取之处，那么有很大的可能性，我们可以从许多名字中推断，那将要来的人或许被称为\Quote{泰坦}。但我们不愿冒险关于假基督的名字作明确断言；因为如果他的名字必须在此刻明确启示，那么看到默示录异象的那位本应宣告出来。因为那异象在不久之前，几乎在我们的时代，即\Person{多米提安}统治末期才被看见。\par

4. 但祂如今指示了名字的数字，以便当此人到来时，我们能认出他是谁而避开他；名字却被隐藏了，因为不值得由圣神宣告。因为如果由祂宣告，他（假基督）可能会持续一段很长的时间。但如今，正如\BibleQuote{他曾经在，如今不在，将要从深渊中上来，并走向灭亡}\BibleRef{默 17:8}，如同一个不存在的人；同样，他的名字也未宣告，因为那不存在的名字不被宣告。但当这假基督在世上肆虐一切之后，他将统治三年六个月，并坐在耶路撒冷的圣殿里；然后主将从天上的云中降临，带着圣父的荣耀，将这人和跟随他的人投入火湖；但为义人带来国度的时期，即安息，圣化的第七日；并将应许的产业归还给亚巴郎，在这国度中，主曾宣告：\BibleQuote{有许多人要从东方和西方来，与亚巴郎、依撒格、雅各伯一同坐席。}\BibleRef{玛 8:11}\par

\SourceRecord{Translated by Alexander Roberts and William Rambaut.；Ante-Nicene Fathers；Vol. 1.；Edited by Alexander Roberts, James Donaldson, and A. Cleveland Coxe.；Buffalo, NY: Christian Literature Publishing Co.,；1885.；<http://www.newadvent.org/fathers/0103530.htm>.}

% structure: p0001 source=h1 target=section reason=page-title

\section{\WorkTitle{驳异端（第五卷，第三十一章）}}

1. 再者，有些被视为正统的人，为了高举义人而超越预定的计划，且不晓得他们在不朽之前受训练的方法，因而持异端意见。因为异端者藐视天主的化工，不承认他们肉身的救恩，同时也轻蔑天主的应许，在他们所形成的思想中完全超越天主，断言他们死后立刻升到诸天和\OriginalTerm{德缪哥}{Demiurge}之上，并到母亲\OriginalTerm{阿卡莫特}{Achamoth}或他们所虚构的父那里去。因此，那些否认涉及整个人类的复活\OriginalTerm{完全否认复活}{universam reprobant resurrectionem}，并尽其所能将其从[基督信仰计划]中移除的人，如果他们再次对复活计划一无所知，又有什么可惊讶的呢？因为他们不愿明白，如果事情如他们所说，那主自己——他们自称所信的那位——就不会在第三日复活；而是在他断气在十字架上时，无疑立刻升到高处，将身体留在地上。但事实是，他在死者所在之处住了三天，正如先知论到他所说：\Quote{主记念他那先前睡在坟墓之地中的圣死者；他降临到他们那里，为要拯救并救出他们。}主自己说：\BibleQuote{正如约纳在大鱼腹中三天三夜，人子也要在地心三天三夜。}\BibleRef{玛 11:40}然后宗徒也说：\BibleQuote{但他升了天，若不是说他先降到地底下，是什么呢？}\BibleRef{弗 4:9}达味预言他时也这样说：\Quote{你从极深的地狱中救了我的灵魂；}在他第三日复活后，他对首先看见并朝拜他的玛利亚说：\BibleQuote{你别摸我，因为我还没有升到父那里；但你去到门徒那里，对他们说：我升到我的父和你们的父那里去。}\BibleRef{若 20:17}\par

2. 那么，如果主遵守了死者的律法，为成为从死者中首先复生的，并停留到第三日\BibleQuote{在地底下}\BibleRef{弗 4:9}；然后肉身复活，甚至向门徒显示钉痕，这样才升到父那里——[如果这一切都发生了，我说]，这些人怎能不羞愧呢？他们声称\Quote{地底下}指的是我们这个尘世，而他们的内心之人离开这里的身体，升到天上之地之上。因为主\Quote{进入死荫之中}，那里是死者灵魂所在之处，之后却肉身复活，复活后被接升[天]，那么显然，主的门徒的灵魂——主为了他们而经历这些——也要去到天主为他们所指定的不可见之处，在那里等待复活，期待那时刻；然后领回他们的身体，完整地复活，即肉身复活，正如主复活一样，他们就这样来到天主面前。\BibleQuote{因为门徒不能高过老师，凡成全的，必像他的老师一样。}\BibleRef{路 6:40}因此，正如我们的主没有立即飞升[天堂]，而是等待父所规定的复活时刻——这也是通过约纳所预表的——并在三天后复活，然后被接升[天]；同样，我们也应当等待天主所规定、众先知所预言的我们复活的时候，然后复活，被接升天——凡主认为配得此[特恩]的人，都要如此。\par

\SourceRecord{Translated by Alexander Roberts and William Rambaut.；Ante-Nicene Fathers；Vol. 1.；Edited by Alexander Roberts, James Donaldson, and A. Cleveland Coxe.；Buffalo, NY: Christian Literature Publishing Co.,；1885.；<http://www.newadvent.org/fathers/0103531.htm>.}

% structure: p0001 source=h1 target=section reason=page-title

\section{《反异端论》卷五，第三十二章}

1. 因此，既然某些［正统］人士的见解来源于异端言论，他们既不明了天主的经纶，也不明了义人复活和［属世的］国度的奥秘——这国度是不朽的开端，藉着这国度，那些配得的人逐渐习惯于分享天主性（\OriginalTerm{承受天主}{capere Deum}）；必须就这些事告诉他们：义人理当首先承受天主应许给先祖的产业的应许，并在其中为王，当他们复活、在这被更新的受造物中看见天主时，而后审判才会发生。因为公义的是，正是在他们劳苦或受难、在各样苦难中被试炼的这同一个受造物中，他们应得他们苦难的赏报；正是在他们因爱天主而被杀害的这受造物中，他们应再次复活；正是在他们忍受奴役的这受造物中，他们应统治。因为天主在万有中是富有的，万有都属于他。因此，受造物本身既恢复其原始状态，理当毫无约束地处于义人的统治之下；宗徒在\WorkTitle{致罗马人书}中已清楚说明这一点，当他这样说：\BibleQuote{因为受造之物切望等待天主子女的显扬。因为受造之物被屈伏于虚幻之下，不是出于自愿，而是出于那屈伏者，因着希望；因为受造之物本身也要从腐败的束缚中被释放，进入天主子女光荣的自由。}\BibleRef{罗 8:19}\par

2. 这样，天主给亚巴郎的应许就坚定不移。因为天主这样说：\BibleQuote{举起你的眼，从你现在的地方向南北东西观看。因为你所看见的这一切地，我都要赐给你和你的后裔，直到永远。}\BibleRef{创 13:13} 他又说：\BibleQuote{起来，纵横走遍这地，因为我必赐给你。}\BibleRef{创 13:17} 然而他并没有在其中得到产业，连一脚之地也没有得到，却始终在那里作客旅和寄居的。\BibleRef{宗 7:5} 在他妻子撒辣去世时，赫特人愿意赠他一块葬地，他谢绝了礼物，却从赫特人左哈尔的儿子厄斐龙那里买了坟地（付出了四百银塔冷通）。\BibleRef{创 23:11} 他就是这样耐心等候天主的应许，不愿看似从人手中接受天主应许要赐给他的东西，当天主又对他这样说：\BibleQuote{我要将这地赐给你的后裔，从埃及河直到幼发拉底大河。}\BibleRef{创 15:13} 那么，如果天主应许他将承受这地为产业，而他在这地寄居的全部时间并未得到，就必定是：他要和他的后裔，即敬畏天主并信靠他的人，在义人复活时一同承受。因为他的后裔就是教会，教会藉着主获得天主的义子名分，正如洗者若翰所说：\BibleQuote{天主能从这些石头中给亚巴郎兴起子孙来。}\BibleRef{路 3:8} 宗徒在\WorkTitle{致迦拉达人书}中也说：\BibleQuote{可是，弟兄们，你们像依撒格一样，是恩许的子女。}\BibleRef{迦 4:28} 在同一封书信中，他又清楚宣告说，凡信靠基督的人就得到基督，即给亚巴郎的应许，这样说：\BibleQuote{所应许的话原是向亚巴郎和他后裔说的。他没有说\InnerQuote{和后裔们}，好像是许多人，而是说\InnerQuote{和你那后裔}，指一个人，就是基督。}\BibleRef{迦 3:16} 他又确认先前的话说：\BibleQuote{正如亚巴郎信了天主，天主就以此算为他的正义。所以你们要知道：凡出于信德的，都是亚巴郎的子孙。圣经预见天主要使外邦人因信德成义，就预先把福音传给亚巴郎说：\InnerQuote{万民都要因你而得福。}这样，凡出于信德的，都和有信德的亚巴郎同受祝福。}\BibleRef{迦 3:6} 这样，凡出于信德的，都和有信德的亚巴郎同受祝福，这些人就是亚巴郎的子孙。天主曾将大地应许给亚巴郎和他的后裔；然而亚巴郎和他的后裔，即因信德成义的人，现在并没有得到其中的任何产业；他们将在义人复活时得到它。因为天主是真实而信实的；正因如此，他说：\BibleQuote{温良的人是有福的，因为他们要承受土地。}\BibleRef{玛 5:5}\par

\SourceRecord{Translated by Alexander Roberts and William Rambaut.；Ante-Nicene Fathers；Vol. 1.；Edited by Alexander Roberts, James Donaldson, and A. Cleveland Coxe.；Buffalo, NY: Christian Literature Publishing Co.,；1885.；<http://www.newadvent.org/fathers/0103532.htm>.}

% structure: p0001 source=h1 target=section reason=page-title

\section{\WorkTitle{驳斥异端}（第五卷，第33章）}

1. 因此，当祂即将受难时，为了向亚伯拉罕和与他在一起的人宣布继承已开放的喜讯，[基督]在拿着杯祝谢后，饮了杯，并递给门徒，对他们说：\Quote{你们都要喝这杯，因为这是我的血，新盟约的血，将为多人流出，以赦免罪过。但我告诉你们：从今以后，我不再喝这葡萄树的果实，直到我在我父的国里与你们同喝新酒的那一天。}\BibleRef{玛 26:27}因此，祂将亲自更新大地的继承，并重新组织[祂]儿子们荣耀的奥秘；正如达味所说：\Quote{祂更新了地面。}祂应许与门徒同饮葡萄树的果实，从而表明这两点：大地的继承，其中新葡萄树的果实被饮用；以及祂门徒在肉身上的复活。因为那新肉身复活，正是那接受新杯的肉身。祂绝不可能是与门徒定居在超天的某处时饮用葡萄树的果实；同样，那些饮用它的人也不是没有肉身的，因为饮用那从葡萄树流出的东西属于肉身，而不属于灵。\par

2. 为此，主宣告说：\Quote{当你设午宴或晚宴时，不要请你的朋友、邻居或亲戚，免得他们回请你，报答你。宁可请瘸腿的、瞎眼的、贫穷的，你就有福了，因为他们无力报答你；但在义人复活时，你必得到报答。}\BibleRef{路 14:12}祂又说：\Quote{凡为我的缘故舍弃田地、房屋、父母、兄弟、子女的，今生必得百倍，来世必得永生。}\BibleRef{玛 19:29}因为在今世那百倍[赏报]是什么？不就是款待穷人，以及那些有回报的晚餐吗？这些[要发生]在王国时期，即在第七天，那已被祝圣的日子，天主歇了祂所创造的一切工作，这就是义人真正的安息日，他们在这日不再从事任何世俗工作，但天主为他们预备了桌子，供应各种佳肴。\par

3. 依撒格祝福他的次子雅各伯的祝福具有同样的含义，他说：\Quote{看，我儿子的香气，犹如主所祝福的丰田的香气。}\BibleRef{创 27:27}但\Quote{田地就是世界。}\BibleRef{玛 13:38}因此他又加上：\Quote{愿天主赐给你天上的甘露，土地的肥脂，五谷和美酒的丰盈。愿万民事奉你，万王向你下拜；愿你作你兄弟的主人，你父亲的儿子们向你下拜；诅咒你的，必受诅咒；祝福你的，必蒙祝福。}\BibleRef{创 27:28}如果任何人拒绝接受这些事是指向所定的王国，他必陷入许多矛盾和冲突，正如犹太人那样，他们被完全困惑。因为不仅今生的万民没有事奉这位雅各伯，甚至在他得到祝福之后，他自己离开[家]去服事他的舅父叙利亚人拉班二十年；\BibleRef{创 31:41}他不但没有作他兄弟的主人，反而从美索不达米亚回到父亲那里时，在他兄弟厄撒乌面前下拜，并送给他许多礼物。\BibleRef{创 33:3}再者，他如何在这里继承了许多五谷和美酒？他因居住之地遭受饥荒而移民到埃及，并臣服于当时统治埃及的法郎。因此，那预言的祝福无疑属于王国时期，那时义人将从死者中复活而掌权；那时受造界也被更新并释放，从天上的甘露和土地的丰饶中结出丰盛的食物：正如那些见过主的门徒若望的长老们所传述的，他们从若望那里听到主如何教导关于这些时期，并说：\Quote{日子将到，那时葡萄树将长出，每棵有一万枝，每枝有一万条嫩枝，每条实枝有一万新芽，每个新芽有一万串，每串有一万颗葡萄，每颗葡萄压榨后将出产二十五梅特来得斯葡萄酒。当任何一位圣人抓住一串时，另一串会喊道：\InnerQuote{我是一串更好的，拿我吧；通过我赞美主。}}\Quote{同样，[主宣称]一粒麦子将结出一万穗，每穗有一万粒，每粒麦子将产出十磅（\OriginalTerm{五比勒布斯}{quinque bilibres}）洁净、纯细的面粉；以及其他所有果树、种子和草，都将以相应的比例（\OriginalTerm{相应比例}{secundum congruentiam iis consequentem}）出产；并且所有只吃地上出产的动物，[在那时]将彼此和平和谐，完全顺服于人。}\par

4. 这些事由若望的听者、波利卡普的同伴\Person{帕皮亚}在其第四卷书中书面作证；因为他汇编了\OriginalTerm{汇编}{\GreekText{συντεταγμένα}}五卷书。而且他补充说：\Quote{现在这些事对信徒是可信的。}而且他说，\Quote{当叛徒\Person{犹达斯}不信这些事，并且发问说：\InnerQuote{这些将要如此丰盛地产生的事，主怎能成就呢？}主就宣告说：\InnerQuote{那些将要来到这些时期的人将会看见。}}因此，\Person{依撒意亚}预言这些时期时说道：\BibleQuote{豺狼将与羔羊同食，豹子将与山羊羔同卧；牛犊、幼狮和肥畜必一同吃食，小孩子要牵引它们。牛与熊必同食，它们的幼崽必一起躺卧；狮子必吃草如牛。吃奶的婴儿必在虺蛇的洞口玩耍，断奶的幼儿必按手在毒蛇的穴上。在我整个圣山上，它们都不再伤人，也不再害物。}他在\OriginalTerm{总结}{recapitulation}中又说：\BibleQuote{那时，豺狼和羔羊将一起觅食，狮子要像牛一样吃草，蛇要以尘土为食；在我整个圣山上，它们不再伤害或骚扰任何事物，上主说。}\BibleRef{依 40:6}我很清楚，有些人努力将这些话解释为来自不同民族、各种习性的野蛮人，他们来相信，并在相信后与义人和谐相处。但尽管目前这确实适用于某些来自不同民族、归向信德和谐的人，然而在义人的复活中，这些话也将应用于那些提到的动物。因为天主在一切事上都是富有的。当受造界复原时，这是合宜的：所有动物都应服从人并受制于人，回归天主最初赐予的食物——因为它们原本在服从\Person{亚当}的情况下受制——即地上的产物。但需要另找机会，而非现在，来证明狮子那时将吃草。这表明果实的大小和丰富品质。因为如果那动物，狮子，在那个时期吃草，那用作狮子合适食物的麦子本身该是何等品质呢？\par

\SourceRecord{Translated by Alexander Roberts and William Rambaut.；Ante-Nicene Fathers；Vol. 1.；Edited by Alexander Roberts, James Donaldson, and A. Cleveland Coxe.；Buffalo, NY: Christian Literature Publishing Co.,；1885.；<http://www.newadvent.org/fathers/0103533.htm>.}

% structure: p0001 source=h1 target=section reason=page-title

\section{驳斥异端（第五卷，第三十四章）}

1. 再者，依撒意亚本人也明明白白地宣告：在义人复活时，将有这等喜乐，他说：\Quote{死人要复活，那些在坟墓里的也要兴起，那些在地下的要欢乐，因为从祢来的甘露是他们的健康。}\BibleRef{依 26:19} 而厄则克耳也同样说：\Quote{看，我要打开你们的坟墓，把你们从坟墓中领出来，当我把我的百姓从坟墓中领出来时，我要在你们内注入气息，你们便要生活，我要把你们安置在你们自己的土地上，那时你们便知道我是上主。}\BibleRef{则 37:12} 同一位先知又这样说：\Quote{上主这样说：我要从异民中聚集以色列，从他们被驱逐所到的各地聚集他们，我要在众异民眼前，在他们身上显圣，他们便住在自己的土地上，就是我赐给我的仆人雅各伯的土地。他们要安居在那里，建造房屋，栽种葡萄园，怀着希望居住，那时我要向一切羞辱他们的，和四周包围他们的人施行审判，他们便知道我是上主，他们的天主，也是他们祖先的天主。}\BibleRef{则 28:25} 我不久前已指出，教会就是亚巴郎的后裔；为此，为使我们知道，那在新约中\Quote{能从石头中给亚巴郎兴起子孙来}\BibleRef{玛 3:9} 的，就是那要按旧约从万民中聚集得救者，耶肋米亚说：\Quote{看，时日将到，上主说，人不再说：那领以色列子民从北方和从他们所被驱逐的各地出来的上主是永生的，祂要使他们回到自己的土地，就是我赐给他们祖先的土地。}\BibleRef{耶 23:6}\par

2. 按照天主的旨意，整个受造界将大大增加，以便结出并维持〔如上所述〕的果实，依撒意亚宣告说：\Quote{到那一天，在一切高山和一切高岗上，必有水流；当许多人灭亡，城墙倒塌时。月光将如日光，日光将增加七倍，如同七天的光，那时上主将医治祂百姓的创伤，并治愈祂打击的伤痕。}\BibleRef{依 30:25} 所谓\Quote{打击的伤痕}，是指起初在亚当内加给悖逆之人的创伤，即死亡；上主在我们复活时，要治愈这创伤，并恢复祖先的产业，正如依撒意亚又说：\Quote{你要因上主而自信，祂必使你走遍大地，并以你祖先雅各伯的产业喂养你。}\BibleRef{依 58:14} 这正是主所宣布的：\Quote{那些主人来到时，发现警醒着的仆人，是有福的。我实在告诉你们：主人要束上腰，请他们坐席，亲自伺候他们。他若在第二更，或在第三更来到，发现他们如此，那些仆人才是有福的。}\BibleRef{路 12:37} 同样，若望在默示录中也说了同样的话：\Quote{凡参与第一次复活的人是有福的，是圣洁的。}\BibleRef{默 20:6} 此外，依撒意亚也说明了这些事发生的时候，他说：\Quote{我要遂说：上主，要到几时呢？直到城邑荒凉，无人居住，房屋空寂无人，土地变成荒野。此后，上主将使我们远离人类（\Emph{longe nos faciet Deus homines}），那些存留的人将在地上繁衍。}\BibleRef{依 6:11} 达尼尔也同样说：\Quote{国度、权柄和天下万国的尊荣，都归于至高者的圣民；祂的国是永远的，一切权柄都要事奉并服从祂。}\BibleRef{达 7:27} 为避免人们将这应许理解为指今世，先知被告知：\Quote{你要去等候你的终局，安息吧！到末日你必起来承受你的产业。}\BibleRef{达 12:13}\par

3. 这些应许不仅给予先知和祖先，也给予从万民中与这些〔祖先〕联合的教会，圣神称这些教会为\Quote{岛屿}（因为它们建立在风浪之中，忍受亵渎的风暴，成为遇险者的安全港，也是那些爱慕高处〔天堂〕并努力避免\OriginalTerm{深渊}{Bythus}，即错误深渊者的避难所），耶肋米亚如此宣告：\Quote{万民哪，请听上主的话，在遥远的海岛宣扬说：那驱散以色列的，要再聚集他，看守他，像牧人看守自己的羊群。因为上主救赎了雅各伯，从他强于自己的手中将他救出。他们要前来，在熙雍山上欢呼，奔向美物，奔向谷、酒、果品、牛羊之地；他们的灵魂像结实的果树，不再饥饿。那时，处女在青年团体中欢乐，老年人也欢欣，我要将他们的悲哀转为喜乐，使他们欢乐，并增加他们；我要使司祭肋未子孙的内心满足，我的百姓要因我的美物而满足。}\BibleRef{耶 31:10} 在前一卷书中，我已指出，主的所有门徒都是肋未人和司祭，他们曾在圣殿里亵渎安息日，却是无可指摘的。\BibleRef{玛 12:5}\par

4. 再者，提及\Place{耶路撒冷}和祂在那里为王的事，\Person{依撒意亚}宣告说：\Quote{上主这样说：凡在\Place{熙雍}有后裔，在\Place{耶路撒冷}有仆人的人，是有福的。看啊，一位正义的君王要统治，首领们将以公义治理。}\BibleRef{依 31:9}至于重建的根基，他说：\Quote{看啊，我要为你安置红玉作基石，以蓝宝石作根基；我要以碧玉造你的城垛，以水晶造你的城门，以精选的宝石造你的城墙。你所有的儿女都必受天主教诲，你的儿女必大享平安；你必在公义上被建造。}\BibleRef{依 54:11}他又再次说同样的事：\Quote{看啊，我要使\Place{耶路撒冷}成为欢乐，使我的百姓[成为欢乐]；在她内再也听不到哭泣的声音，也听不到哀号的声音。那里没有夭折的人，也没有不滿寿数的老人；因为百岁的人还算青年，而罪人虽活百岁仍被诅咒。他们要建造房屋，自己居住；要栽种葡萄园，吃所结的果实，并喝酒。他们建造的，别人不得居住；他们栽种的，别人不得享用。因为在你内的百姓的日子，必如生命树的日子；他们手所作的，必存留。}\BibleRef{依 65:18}\par

\SourceRecord{Translated by Alexander Roberts and William Rambaut.；Ante-Nicene Fathers；Vol. 1.；Edited by Alexander Roberts, James Donaldson, and A. Cleveland Coxe.；Buffalo, NY: Christian Literature Publishing Co.,；1885.；<http://www.newadvent.org/fathers/0103534.htm>.}

% structure: p0001 source=h1 target=section reason=page-title

\section{驳斥异端（卷五，第三十五章）}

1. 然而，如果有人试图以寓意来解释这类预言，他们将发现在所有点上自相矛盾，并被这些表达本身的教导所驳斥。例如：\BibleQuote{当城市} 外邦人 \BibleQuote{成为荒凉，无人居住，房屋也无人居住，土地变为荒废} \BibleRef{依 6:11} \BibleQuote{因为，看啊，} \Person{依撒意亚}说，\BibleQuote{主的日子来临，无可救药，充满烈怒和忿怒，要使大地成为荒凉，从中铲除罪人。} \BibleRef{依 13:9} 他又说：\BibleQuote{让他被带走，免得他看见天主的荣耀。} \BibleRef{依 26:10} 当这些事完成时，他说：\BibleQuote{天主将使人远离，留下的人要在地上增多。} \BibleRef{依 6:12} \BibleQuote{他们要建造房屋，自己居住；栽种葡萄园，自己吃用。} \BibleRef{依 65:21} 因为所有这些话和其他话，无疑都是指着义人的复活说的，这复活发生在假基督来临之后，在他统治下所有民族被毁灭之后；在那时候，义人将在世上为王，因看见主而日益坚强；并藉着祂，他们将习惯于分享天主父的荣耀，在王国中与圣天使交往共融，与属灵的存在结合；并且关乎那些主将在肉身中寻见、从天上等候祂、既受逼迫又逃脱了那恶者之手的人。因为先知正是指着他们说：\BibleQuote{那留下的人要在地上增多。} \Person{耶肋米亚}先知也指明，天主为此目的所预备的信徒，为使留下的人在地上增多，他们既要受圣者统治，为这耶路撒冷服务，且祂的王国要在其中，说：\BibleQuote{你向耶路撒冷四周观望东部，看那从天主亲自来到你那里的喜乐。看，你所放逐的儿子们要来了，他们要从东方直到西方成群而来，因圣者的话，因你的天主所赐的光辉而喜乐。耶路撒冷啊！脱去你哀悼和苦难的长袍，穿上你天主所赐的永恒光辉的美丽。束上从你天主而来的正义的双重外衣，将永恒荣耀的冠冕戴在你的头上。因为天主要在普天之下显扬你的荣耀。因为你的名字要被天主自己永远称为：\InnerQuote{正义的平安}和\InnerQuote{敬畏天主者的荣耀}。耶路撒冷啊！起来，站在高处，向东方观望，看你的儿子们从日出之地直到西方，因圣者的话，因对天主的纪念而喜乐。因为步兵从你那里出发，被敌人掳去；天主却要使他们荣耀地归来，如同王国的宝座。因为天主已命定，一切高山都要降低，永恒的丘陵也要降低，山谷要填满，使地面平坦，使以色列——天主的荣耀——安全行走。森林也要提供荫凉之处，各样香气的树木要因天主的命令为以色列自己生长。因为天主要在祂的光辉中喜乐地前行，带着从祂而来的怜悯和正义。}\par

2. 现在，凡此种种既已如此，便不能理解为指向超天界之事；因为经上说：\BibleQuote{天主要向普天之下显示你的荣耀。} 但在王国时期，大地被基督重新召回到原初状态，耶路撒冷按上层的耶路撒冷的样式得以重建，关于这城，依撒意亚先知说：\BibleQuote{看，我已将你的城墙绘在我掌上，你常在我眼前。} \BibleRef{依 49:16} 宗徒在写给迦拉达人的书信中，也同样说：\BibleQuote{但在上层的耶路撒冷是自由的，她是我们的母亲。} \BibleRef{迦 4:26} 他说这话，并非想到一个漂移的\OriginalTerm{永世}{Æon}，或某个从\OriginalTerm{普累若麻}{Pleroma}流出的其他权能，或\OriginalTerm{普鲁尼科斯}{Prunicus}，而是指那已被描绘在\Person{天主}手上的耶路撒冷。在 \BibleBook{}{默示录} 中，\Person{若望}看见这新的耶路撒冷降在新地上。\BibleRef{默 21:2} 因为在王国时期之后，他说：\BibleQuote{我看见一个白色的大宝座，和坐在上面的那位；天地从他面前逃避，再也找不到它们的位置了。} \BibleRef{默 20:11} 他又陈述有关普遍复活和审判的事，提到\BibleQuote{死人，无论大小}。他说：\BibleQuote{海交出其中的死者，死亡和阴府也交出其中的死者；书卷就展开了。} 他又说：\BibleQuote{生命册也展开了，死者按书卷上所记载的，照他们的行为受审判；死亡和阴府被扔进火湖，这是第二次死亡。} \BibleRef{默 20:12} 这就是所谓的\OriginalTerm{革哈纳}{Gehenna}，\Person{主}称之为永火。\BibleRef{玛 25:41} 经上说：\BibleQuote{若有人名字没有记在生命册上，他就被扔进火湖。} \BibleRef{默 20:15} 此后，他说：\BibleQuote{我看见一个新天新地，因为先前的天地已经过去了，海也不再有了。我又看见圣城新耶路撒冷，由天主那里从天而降，好像一位新娘装饰好了，等候她的丈夫。} 经上说：\BibleQuote{我听见有大声音从宝座出来说：\InnerQuote{看，天主的帐幕在人间，他要与人同住；他们要作他的子民，天主要亲自与他们同在，作他们的天主。他要擦去他们眼中一切的眼泪；不再有死亡，也不再有悲哀、哭号、疼痛，因为先前的事都过去了。}} \BibleRef{默 21:1} 依撒意亚也宣布同样的事：\BibleQuote{因为要有新天新地；先前的事不再被记念，也不再涌上心头；但他们要因我所造的永远欢喜快乐。} \BibleRef{依 65:17} 这就是宗徒所说的：\BibleQuote{因为这世界的样式正在逝去。} \BibleRef{格前 7:31} 同样的，主也宣告：\BibleQuote{天地要逝去。} \BibleRef{玛 26:35} 因此，当这些事在天地之上逝去时，主的门徒\Person{若望}说，上层的、新的耶路撒冷将像新娘装饰好等候丈夫那样降下；这就是天主的帐幕，天主将在其中与人同住。这耶路撒冷是那座耶路撒冷的预像——那先前大地上的耶路撒冷，在其中义人预先被训练以达致不朽，并为救恩作准备。关于这帐幕，\Person{梅瑟}在山上领受了样式；\BibleRef{出 25:40} 没有一样是可以寓意化的，而一切都是坚定、真实、实在的，由天主造给义人享用。因为正如天主真正使人复活，人也真正从死者中复活，而非寓意式的，如我多次指明的。正如他确实复活，他也确实在王国时期预先被训练以达致不朽，并前进、繁盛，以便能承受\Person{圣父}的荣耀。然后，当万物更新时，他将真正住在\Person{天主}的城里。因为经上说：\BibleQuote{坐在宝座上的那位说：\InnerQuote{看，我更新一切。} 主说：\InnerQuote{你要写下这一切，因为这些话是可信靠的、真实的。} 他对我说：\InnerQuote{已成全了。}} \BibleRef{默 21:5} 这就是事情的真相。\par

\SourceRecord{Translated by Alexander Roberts and William Rambaut.；Ante-Nicene Fathers；Vol. 1.；Edited by Alexander Roberts, James Donaldson, and A. Cleveland Coxe.；Buffalo, NY: Christian Literature Publishing Co.,；1885.；<http://www.newadvent.org/fathers/0103535.htm>.}

% structure: p0001 source=h1 target=section reason=page-title

\section{驳斥异端（第五卷，第三十六章）}

1. 既然有真实的人，就必有真实的\OriginalTerm{建立}{plantationem}，免得他们消失于虚无之物中，而在实存者中进步。因为受造的实质和本质绝不会被消灭（建立它的那位是信实而真实的），但\Quote{这世界的\Emph{样式}将要过去}\BibleRef{格前 7:31}；也就是说，在其中发生过过犯的那些事物将要过去，因为人在其中已变得衰老。因此，这个[现时的]样式被造成暂时的，天主预知一切；正如我在前一卷中所指出的，也尽可能说明了这暂时之物世界的受造原因。但当这个[现时的]样式过去，人得以更新，在不朽的状态中繁荣，从而不可能再变得衰老时，便将有新天新地，新的人将[永远]居住在当中，常常与天主保持新鲜的交谈。而且（或：既然）这些事将永远持续，没有终结，依撒意亚宣告说：\Quote{正如我所造的新天新地，在我面前常存，上主说，你们的后裔和你们的名字也必常存。}\BibleRef{依 66:22}如同长老们所说：那时，那些堪当天上居所的人将去那里，另一些人将享受乐园的愉悦，还有一些人将拥有城市的光辉；因为救主将随处被看见，按照看见他的人所堪当的程度。\par

2. [他们还说]，结百倍果实者的居所、结六十倍者的居所、结三十倍者的居所之间有这样的区别：第一种将被提到天上，第二种将住在乐园中，最后一种将居住在城市里；主为此曾宣告说：\Quote{在我父的家里有许多住处。}\BibleRef{若 14:2}因为万物都属于天主，祂赐给每人适宜的住所；正如祂的圣言所说，圣父照着各人是否或将成为堪当的，分配分额。这就是宾客们应邀赴婚宴时所靠的卧榻\BibleRef{玛 22:10}。宗徒的门徒们——长老们肯定说，这就是得救者的等级和安排，他们经由这样的步骤升进；也即，他们通过圣神上升到圣子，并通过圣子上升到圣父，到了适当的时机，圣子将祂的工程交给圣父，正如宗徒所说：\Quote{因为祂必须为王，直到把一切仇敌放在祂脚下。最后被毁灭的仇敌是死亡。}\BibleRef{格前 15:25}因为在王国时期，在地上的义人那时将忘记死亡。\Quote{但祂说万物都要屈服在祂脚下时，显然那使万物屈服于祂的除外。等到万物都屈服于祂，那时子自己也要屈服于那使万物屈服于祂的，好叫天主成为万物之中的万有。}\BibleRef{格前 15:27}\par

3. 因此，若望清楚地预见了第一次\Quote{义人的复活}\BibleRef{路 14:14}和地上王国的继承；先知们关于这事所预言的也与之[他的异象]一致。因为主也教导了这些事，当祂应许要在王国中与祂的门徒一起新喝那混合的杯。宗徒也承认受造物将从败坏的束缚中获得自由，[进入]天主儿女的光荣自由\BibleRef{罗 8:21}。在这一切事中，并借着这一切，同一个天主圣父被彰显出来：祂造了人，赐给祖先地上的继承应许，在义人复活时领出受造物（脱离束缚），并为祂圣子的王国履行应许；随后以父亲的方式赏赐那些眼所未见、耳所未闻、人心也未曾想到的事\BibleRef{格前 2:9}。因为只有一个圣子，完成了祂圣父的旨意；也只有一个人类，天主的奥迹在其中运作，\Quote{天使也切愿察看}\BibleRef{伯前 1:12}；他们不能穷究天主的智慧，借这智慧，祂亲手的工作被确认并与祂的圣子结合，达到完美：使祂的首生者圣言降下到\OriginalTerm{受造物}{facturam}，即到那被\OriginalTerm{塑造之物}{plasma}，并被祂容纳；另一方面，受造物应容纳圣言，并上升到祂那里，超越天使，成为按照天主的肖像和模样。\par

\SourceRecord{Translated by Alexander Roberts and William Rambaut.；Ante-Nicene Fathers；Vol. 1.；Edited by Alexander Roberts, James Donaldson, and A. Cleveland Coxe.；Buffalo, NY: Christian Literature Publishing Co.,；1885.；<http://www.newadvent.org/fathers/0103536.htm>.}
